% !TEX program = xelatex
% 中文自动译本：Daniele L. Marchisio and Rodney O. Fox,
% Computational Models for Polydisperse Particulate and Multiphase Systems
% (Cambridge University Press, 2013).
\documentclass[11pt,a4paper,openany,fontset=none]{ctexbook}
\usepackage[margin=2.55cm,headheight=15pt]{geometry}
\usepackage{microtype}
\usepackage{fancyhdr}
\usepackage{hyperref}
\usepackage{xurl}
\usepackage{titlesec}
\usepackage{amsmath,amssymb,amsfonts,bm,mathtools}
\usepackage[export]{adjustbox}
\usepackage{enumitem}
\usepackage{graphicx}
\usepackage{caption}
\usepackage{longtable,booktabs,array}
\setmainfont[AutoFakeBold=2,AutoFakeSlant=.2]{Arial Unicode MS}
\setCJKmainfont{Songti SC}
\setCJKsansfont{Songti SC}
\setCJKmonofont{Songti SC}
\setlength{\parindent}{2em}
\setlength{\parskip}{0.65em}
\setlength{\emergencystretch}{3em}
\linespread{1.30}\selectfont
\xeCJKsetup{CJKglue = \hskip 0.075em plus 0.04em minus 0.03em}
\setlist[itemize]{leftmargin=2.8em,itemsep=0.3em,topsep=0.45em}
\titleformat{\chapter}[display]{\bfseries\Huge\centering}{第\thechapter 章}{0.75em}{}
\titleformat{\section}{\large\bfseries}{}{0pt}{}
\titleformat{\subsection}{\normalsize\bfseries}{}{0pt}{}
\titleformat{\paragraph}[block]{\normalsize\bfseries}{}{0pt}{}
\titlespacing*{\section}{0pt}{1.45em}{0.7em}
\titlespacing*{\subsection}{0pt}{1.2em}{0.55em}
\titlespacing*{\paragraph}{0pt}{1.0em}{0.4em}
\pagestyle{fancy}
\fancyhf{}
\fancyhead[C]{\small\nouppercase{\leftmark}}
\fancyfoot[C]{\thepage}
\renewcommand{\headrulewidth}{0.3pt}
\hypersetup{
  colorlinks=true,
  linkcolor=black,
  urlcolor=blue,
  pdftitle={多分散颗粒与多相系统的计算模型},
  pdfauthor={Daniele L. Marchisio, Rodney O. Fox}
}
\graphicspath{{./}}
\title{多分散颗粒与多相系统的计算模型\\[0.5em]\Large 中文自动译文}
\author{Daniele L. Marchisio、Rodney O. Fox 著\\
\small 译文由自动化工具生成，供学习与检索使用}
\date{原著：Cambridge University Press，2013\\生成日期：\today}
% 统一译本版式：避免公式、图片和题注跨页或侵入页眉、页脚。
\usepackage{needspace}
\clubpenalty=10000
\widowpenalty=10000
\displaywidowpenalty=10000
\raggedbottom
\setlength{\abovedisplayskip}{0.8\baselineskip plus 0.2\baselineskip minus 0.1\baselineskip}
\setlength{\belowdisplayskip}{0.8\baselineskip plus 0.2\baselineskip minus 0.1\baselineskip}
\setlength{\abovedisplayshortskip}{0.55\baselineskip}
\setlength{\belowdisplayshortskip}{0.55\baselineskip}
\newenvironment{sourceblock}
  {\par\addvspace{0.65\baselineskip}\noindent\begin{minipage}{\linewidth}\centering}
  {\end{minipage}\par\addvspace{0.45\baselineskip}}
\captionsetup{font=small,labelfont=bf,justification=centering,singlelinecheck=false,skip=0.45\baselineskip}
\begin{document}
\frontmatter
\maketitle
\chapter*{说明}
\markboth{说明}{说明}
本文件是 Daniele L. Marchisio 与 Rodney O. Fox 所著
\emph{Computational Models for Polydisperse Particulate and Multiphase Systems}
的中文自动译本。正文按章节与小节的正常阅读顺序重新排版，不与原书逐页对齐。

原 PDF 使用专用数学字体。为保证积分号、矩阵、上下标及求积节点等数学符号可靠，
本译本把显示公式、图表、代码清单、符号表和参考文献作为原 PDF 的矢量裁片嵌入；
正文和标题则重新排版为中文。涉及公式推导、图表数据和算法实现时，请同时核对原著。
原著版权归作者与 Cambridge University Press 所有。

{\small\tableofcontents}
\clearpage
\makeatletter\@mainmattertrue\makeatother
\pagenumbering{arabic}
\begingroup\let\clearpage\relax
\chapter*{前言}\addcontentsline{toc}{chapter}{前言}
\markboth{前言}{前言}
\endgroup
本书面向对多分散型多相流模拟感兴趣的科学和工程不同分支（即化学、机械、环境、能量学等）的研究生，以及已经在该领域工作的科学家和工程师。事实上，这本书对控制多分散多相系统的基本理论进行了系统且一致的讨论，适合新手，并提出了一类用于实际模拟的特定计算方法，这可能会让更有经验的学者感兴趣。正如整本书所解释的，分散多相系统的特点是多相，其中一相是连续的，其他相是分散的（即以不同的颗粒、液滴或气泡的形式）。在本文中使用术语“多分散”来指定表征分散相元素的相关属性，例如质量、动量或能量，随着元素的不同而变化，从而产生通常所说的分布。通常用作多相系统的特征签名的典型分布例如是晶体尺寸分布（CSD）、颗粒尺寸分布（PSD）和颗粒速度分布。不同的科学界以多种方式处理描述这些分布的演变（在空间和时间上）的问题，重点关注与其社区最相关的方面。例如，在结晶和沉淀领域，用晶体或颗粒尺寸来描述问题（通常忽略空间不均匀性），由此产生的控制方程称为群体平衡方程（PBE）。在蒸发（和非蒸发）喷雾领域，问题是根据颗粒表面积来表述的，控制方程称为威廉姆斯-玻尔兹曼方程。在这个领域和其他领域，人们非常重视所研究的系统在空间上不均匀的事实。气溶胶和超细颗粒通常用颗粒质量来描述，最终的控制方程称为颗粒动力学方程。相反，颗粒流中涉及的颗粒系统仅根据颗粒速度进行研究，控制方程是著名的玻尔兹曼方程（BE）对多相系统的非弹性扩展，该方程用于描述气体动力学中的分子速度分布。尽管这些明显不同的理论框架有不同的名称，但其基本理论（以经典统计力学为基础）是完全相同的。这也产生了大量用于求解控制方程的数值方法，这些方法通常共享许多共同元素，但通常只专门关注问题的一部分。例如，在 PBE 中，表示构成多相系统的元素的分布通常被离散成类或部分，生成所谓的离散群体平衡方程（DPBE）。在开发的众多方法中，计算流体动力学 (CFD) 从业者广泛使用的一种方法是多尺寸组 (MUSIG) 方法。这种方法的基本思想类似于所谓的离散速度方法 (DVM) 中对 BE 进行的离散化。 类似地，矩方法 (MOM) 已用于求解 PBE 和 BE，但通过在两种情况下采用不同的策略来克服由此产生的闭合问题。就 BE 而言，最流行的矩闭包是 Grad 提出的矩闭包，它基于 13 或 26 个矩子集的解，以及速度分布的假定函数形式。相反，在 PBE 的情况下，闭合策略通常涉及已知矩之间的插值（如插值闭合矩方法 MOMIC）。鉴于方法过多，对于新手来说通常不可能看出不同社区所采用的方法之间的联系。

本书基于一般理论框架对这些问题进行了一致的处理。这又源于广义群体平衡方程（GPBE），其中包括前面提到的所有其他控制方程（例如 PBE 和 BE）作为特殊情况。在讨论了该方程的起源之后，提出了其数值解的不同计算模型。本书的结构如下。

• 第1 章通过使用两个例子介绍了关键概念，例如流态和相关无量纲数：细颗粒的PBE 和气体颗粒流的KE。随后，详细解释了全书使用的介观建模方法，特别关注与微观和宏观模型的关系以及由此产生的闭合问题。

• 第2 章简要介绍了多分散系统的介观描述。本章介绍了许多可能的数密度函数 (NDF)，这些函数是用不同的内坐标选择来表示的，然后介绍了各种形式的 PBE。本章最后简短讨论了与 PBE 相关的矩输运方程与湍流理论中的系综平均所产生的矩输运方程之间的差异。

• 第3 章介绍了高斯求积和基于求积矩方法(QBMM) 中使用的矩反演算法。本章讨论了经典单变量矩求积法 (QMOM) 所采用的乘积差分 (PD) 和 Wheeler 算法，以及针对多元问题开发的强力算法、张量积算法和条件 QMOM 算法。本章最后讨论了扩展求积矩方法 (EQMOM) 和直接矩求积法 (DQMOM)。

• 在第4 章中推导了GPBE，强调了从微观模型到介观模型的过渡必须引入的封闭。本章还概述了从 GPBE 导出 NDF 矩的输运方程所需的数学步骤。由此产生的矩闭合问题也得到了深入讨论。

• 第5章重点介绍从文献中选择的关键物理和化学过程的介观模型。本章首先对介观尺度建模原理及其数学框架进行一般性讨论。由于文献中提出的介观尺度模型数量巨大，本章的目标是介绍物理空间和相空间中的平流和扩散以及零阶、一阶和二阶点过程（例如成核、破裂和聚集）的模型示例。

• 第6 章专门讨论QBMM 背景下的硬球碰撞模型（以及相关的更简单的动力学模型）主题。特别是，在具有不同直径/质量的两个粒子之间的非弹性二元碰撞的情况下，推导了由于碰撞而产生的整数矩的精确源项，并说明了使用 QBMM 来克服封闭问题。

• 第7章专门讨论空间齐次GPBE的求解方法，包括类法和截面法、MOM和QBMM以及蒙特卡罗方法。本章最后给出了几个比较所选同质 PBE 求解方法的示例。

• 第8 章重点介绍使用矩方法求解空间非齐次GPBE。空间非均匀系统的关键问题是矩可实现性和损坏（由于数值平流和扩散算子）以及粒子轨迹穿越（PTC）的存在。在介绍基于动理学的有限体积方法之后，通过提出能够保留矩可实现性的数值方案并通过实际例子证明 QBMM 非常适合捕获 PTC，对这些问题进行了讨论。本章最后给出了一些空间一维数值示例。

• 为了使本书完整，包含四个附录。附录 A 包含第 3 章中介绍的最常见矩反演算法的 Matlab 脚本。附录 B 更详细地讨论了第 8 章中介绍的基于动理学的有限体积方法。最后，相空间中 PTC 的关键问题（发生在远离碰撞平衡的系统中）以及某些 QBMM 的矩守恒问题分别在附录 C 和附录 D 中讨论。

作者非常感谢以不同方式为本书的完成做出贡献的许多人。本书的核心是 Robert L. McGraw 博士的开创性研究，他是第一个开发 QMOM 和 Jacobian 矩阵变换（这是 DQMOM 的基础）来解决 PBE 的人，并引起我们注意使用矩方法时矩损坏和可实现性的重要性。因此，作者特别感谢 Daniel E. Rosner 教授，他在 1999 年将他们的注意力转向了 McGraw 博士新发表的关于 QMOM 的著作。他们还要感谢 R. Dennis Vigil 教授认可 QMOM 解决聚集和破损问题的能力，感谢 Prakash Vedula 教授提供了用于计算第 6 章中报告的硬球碰撞矩源项的数学框架。本书中描述的解决方法的中心主题是在数值近似中保持矩集可实现性的重要性。在这一点上，作者特别感谢 Marc Massot 教授，他启发了他们关于双曲系统的基于动力学的方法的微妙性以及粒子轨迹穿越的一般主题。得益于 Olivier Desjardins 教授出色的数值分析能力以及 Philippe Villedieu 博士在湍流研究中心 2006 年暑期项目期间提出的重要建议，Massot 教授的言论最终为我们指明了第 8 章中描述的可实现的有限体积方案的方向。在这方面，我们还要感谢 Z. J. Wang 教授在高阶有限体积方案领域的关键贡献以及 Varun Vikas 博士的开发以及附录 B 中描述的可实现的准高阶方案的实现。

与剑桥大学出版社合作出版本书的想法是 Massimo Morbidelli 教授对该主题表现出兴趣的结果。我们还感谢许多其他同事的贡献，其中包括 Antonello A. Barresi、Marco Vanni、Giancarlo Baldi、Miroslav Soos、Jan Sefcik、Christophe Chalons、Fr ́ed ́erique Laurent、Hailiang Liu、Alberto Passalacqua、Venkat Raman、Julien Reveillon 和 Shankar Subramaniam。过去十年来，作者指导的所有研究生和博士后研究人员对本书的研究结果做出了贡献，在此致以谢意，并详细引用了他们的具体贡献。本书背后的研究工作得到了许多机构的资助，其中值得一提的是欧盟委员会（DLM）、意大利教育、大学和研究部（DLM）、ISI基金会（DLM、ROF）、美国国家科学基金会（ROF）、美国能源部（ROF）、巴黎中央理工学院（ROF）和斯坦福大学湍流研究中心（ROF）。我们还深深感谢众多工业合作伙伴（意大利埃尼集团、德国巴斯夫集团、美国 BP 化学公司、美国康菲石油公司和美国 Univation Technologies）的持续刺激和财政支持。

\chapter*{符号表}\addcontentsline{toc}{chapter}{符号表}
\markboth{符号表}{符号表}
下表给出群体平衡与求积矩方法中最常用符号和缩写的中文释义。符号结构保持原书写法；
后续完整原书符号表用于核对全部上下标、条件平均和模型专用系数。

\begingroup\small
\renewcommand{\arraystretch}{1.18}
\begin{longtable}{>{\centering\arraybackslash}p{0.25\textwidth}p{0.65\textwidth}}
\toprule 符号或缩写 & 中文释义 \\ \midrule
\endfirsthead
\toprule 符号或缩写 & 中文释义 \\ \midrule
\endhead
$n(t,\boldsymbol{x},\boldsymbol{v},\boldsymbol{\xi})$ & 数密度函数（NDF） \\
$\boldsymbol{x},\boldsymbol{v},\boldsymbol{\xi}$ & 物理位置、颗粒速度和内坐标矢量 \\
$m_{\boldsymbol{k}}$ & 多指标 $\boldsymbol{k}$ 对应的 NDF 矩 \\
$w_\alpha,\ \boldsymbol{\xi}_\alpha$ & 第 $\alpha$ 个求积节点的权重与节点位置 \\
$A_p,\ A_f$ & 颗粒加速度和流体加速度 \\
$A_{fp},\ A_{pf}$ & 流体--颗粒及颗粒--流体动量交换引起的加速度 \\
$B(\boldsymbol{g},\boldsymbol{x})$ & 硬球碰撞核 \\
$B_i^{\mathrm{agg}},\ B_i^{\mathrm{break}}$ & 聚并与破碎造成的第 $i$ 组颗粒数变化率 \\
$d_p,\ A_p,\ V_p$ & 颗粒直径、表面积和体积 \\
$\rho_p,\rho_f,\alpha_p$ & 颗粒密度、流体密度和颗粒相体积分数 \\
$e_{\alpha\beta}$ & $\alpha$、$\beta$ 类颗粒碰撞的恢复系数 \\
PBE / GPBE & 群体平衡方程 / 广义群体平衡方程 \\
MOM / QMOM & 矩方法 / 求积矩方法 \\
DQMOM / EQMOM & 直接求积矩方法 / 扩展求积矩方法 \\
DPBE / DVM & 离散群体平衡方程 / 离散速度方法 \\
PTC & 颗粒轨迹交叉 \\
\bottomrule
\end{longtable}
\endgroup
\clearpage
\begin{sourceblock}
\includegraphics[page=19,pagebox=mediabox,trim={58.00bp 50.16bp 58.23bp 185.00bp},clip,width=\textwidth,max height=0.82\textheight]{Marchisio2013.pdf}
\end{sourceblock}
\clearpage
\begin{sourceblock}
\includegraphics[page=20,pagebox=mediabox,trim={58.00bp 50.16bp 58.23bp 55.00bp},clip,width=\textwidth,max height=0.82\textheight]{Marchisio2013.pdf}
\end{sourceblock}
\clearpage
\begin{sourceblock}
\includegraphics[page=21,pagebox=mediabox,trim={58.00bp 50.16bp 58.23bp 55.00bp},clip,width=\textwidth,max height=0.82\textheight]{Marchisio2013.pdf}
\end{sourceblock}
\clearpage
\begin{sourceblock}
\includegraphics[page=22,pagebox=mediabox,trim={58.00bp 50.16bp 58.23bp 55.00bp},clip,width=\textwidth,max height=0.82\textheight]{Marchisio2013.pdf}
\end{sourceblock}
\clearpage
\begin{sourceblock}
\includegraphics[page=23,pagebox=mediabox,trim={58.00bp 50.16bp 58.23bp 55.00bp},clip,width=\textwidth,max height=0.82\textheight]{Marchisio2013.pdf}
\end{sourceblock}
\clearpage
\begin{sourceblock}
\includegraphics[page=24,pagebox=mediabox,trim={58.00bp 50.16bp 58.23bp 55.00bp},clip,width=\textwidth,max height=0.82\textheight]{Marchisio2013.pdf}
\end{sourceblock}
\clearpage
\begin{sourceblock}
\includegraphics[page=25,pagebox=mediabox,trim={58.00bp 50.16bp 58.23bp 55.00bp},clip,width=\textwidth,max height=0.82\textheight]{Marchisio2013.pdf}
\end{sourceblock}
\clearpage
\begin{sourceblock}
\includegraphics[page=26,pagebox=mediabox,trim={58.00bp 50.16bp 58.23bp 55.00bp},clip,width=\textwidth,max height=0.82\textheight]{Marchisio2013.pdf}
\end{sourceblock}
\clearpage
\begin{sourceblock}
\includegraphics[page=27,pagebox=mediabox,trim={58.00bp 50.16bp 58.23bp 55.00bp},clip,width=\textwidth,max height=0.82\textheight]{Marchisio2013.pdf}
\end{sourceblock}
\clearpage
\begin{sourceblock}
\includegraphics[page=28,pagebox=mediabox,trim={58.00bp 50.16bp 58.23bp 55.00bp},clip,width=\textwidth,max height=0.82\textheight]{Marchisio2013.pdf}
\end{sourceblock}
\clearpage
\begin{sourceblock}
\includegraphics[page=29,pagebox=mediabox,trim={58.00bp 50.16bp 58.23bp 55.00bp},clip,width=\textwidth,max height=0.82\textheight]{Marchisio2013.pdf}
\end{sourceblock}
\clearpage
\begin{sourceblock}
\includegraphics[page=30,pagebox=mediabox,trim={58.00bp 50.16bp 58.23bp 55.00bp},clip,width=\textwidth,max height=0.82\textheight]{Marchisio2013.pdf}
\end{sourceblock}
\clearpage
\begin{sourceblock}
\includegraphics[page=31,pagebox=mediabox,trim={58.00bp 50.16bp 58.23bp 55.00bp},clip,width=\textwidth,max height=0.82\textheight]{Marchisio2013.pdf}
\end{sourceblock}
\clearpage
\begin{sourceblock}
\includegraphics[page=32,pagebox=mediabox,trim={58.00bp 50.16bp 58.23bp 55.00bp},clip,width=\textwidth,max height=0.82\textheight]{Marchisio2013.pdf}
\end{sourceblock}
\clearpage
\begin{sourceblock}
\includegraphics[page=33,pagebox=mediabox,trim={58.00bp 50.16bp 58.23bp 55.00bp},clip,width=\textwidth,max height=0.82\textheight]{Marchisio2013.pdf}
\end{sourceblock}
\clearpage
\begin{sourceblock}
\includegraphics[page=34,pagebox=mediabox,trim={58.00bp 50.16bp 58.23bp 55.00bp},clip,width=\textwidth,max height=0.82\textheight]{Marchisio2013.pdf}
\end{sourceblock}
\clearpage
\begin{sourceblock}
\includegraphics[page=35,pagebox=mediabox,trim={58.00bp 50.16bp 58.23bp 55.00bp},clip,width=\textwidth,max height=0.82\textheight]{Marchisio2013.pdf}
\end{sourceblock}
\clearpage
\begin{sourceblock}
\includegraphics[page=36,pagebox=mediabox,trim={58.00bp 50.16bp 58.23bp 55.00bp},clip,width=\textwidth,max height=0.82\textheight]{Marchisio2013.pdf}
\end{sourceblock}
\clearpage
\begin{sourceblock}
\includegraphics[page=37,pagebox=mediabox,trim={58.00bp 50.16bp 58.23bp 55.00bp},clip,width=\textwidth,max height=0.82\textheight]{Marchisio2013.pdf}
\end{sourceblock}
\clearpage
\begin{sourceblock}
\includegraphics[page=38,pagebox=mediabox,trim={58.00bp 50.16bp 58.23bp 55.00bp},clip,width=\textwidth,max height=0.82\textheight]{Marchisio2013.pdf}
\end{sourceblock}
\clearpage
\begin{sourceblock}
\includegraphics[page=39,pagebox=mediabox,trim={58.00bp 50.16bp 58.23bp 55.00bp},clip,width=\textwidth,max height=0.82\textheight]{Marchisio2013.pdf}
\end{sourceblock}
\clearpage
\begin{sourceblock}
\includegraphics[page=40,pagebox=mediabox,trim={58.00bp 50.16bp 58.23bp 55.00bp},clip,width=\textwidth,max height=0.82\textheight]{Marchisio2013.pdf}
\end{sourceblock}
\chapter{引言}
\section*{1.1 分散型多相流}\phantomsection\addcontentsline{toc}{section}{1.1 分散型多相流}
化学加工工业中使用的大多数设备都采用多相流。鼓泡塔、流化床、火焰反应器以及液液萃取、固体干燥和尺寸放大或缩小的设备是常见的例子。为了有效地设计、优化和扩大工业系统，模拟多相流的计算工具非常重要。多分散型多相流在其他领域也很常见，例如汽车和飞机发动机中的燃油喷雾、航空航天器中的电力不足条件以及环境中发生的颗粒流。尽管乍一看，多种多样的工业和环境多相流似乎彼此非常不同，但它们有一个非常重要的共同元素：可以识别连续相和分散相（通常具有特征“粒度”的分布）。从历史上看，分散型多相流的理论框架和计算模型的发展集中在两个不同的方面：（i）分散相的演化（例如气泡或液滴的破裂和聚结、颗粒与颗粒的碰撞等）和（ii）多相流体动力学。第一类模型主要关注分散相的描述，并基于空间均匀群体平衡方程（PBE）的求解。 PBE 是用数密度函数 (NDF) 编写的连续性陈述，将在第 2 章中详细描述。NDF 包含有关由分散相的离散元素构成的总体如何分布在某些特性上的信息，这些特性决定了产品质量等。例如，在结晶过程中，晶体的最终质量通常取决于晶体尺寸分布 (CSD)。 CSD（或 NDF）定义了晶体总体在晶体尺寸上的分布情况，PBE 是一个偏积分微分方程，定义了 NDF 的演化，如 Ramkrishna (2000) 和 Randolph \& Larson (1971) 所描述的。一般来说，这些模型由 PBE 与空间均匀质量和能量平衡相结合组成，可有效地用于工艺设备的建模、设计和放大，以及开发

\footnotetext[1]{在更复杂的处理中，空间不均匀性是通过假设为均匀的连接“区域”或空间区域来建模的。由此产生的群体平衡方程 (PBE) 并未明确考虑流场的局部变化。}
颗粒系统的控制策略（Chiu \& Christofides, 1999; Crowley et al., 2000; Nagy, 2009）。然而，在许多工业和环境应用中，分散相和连续相之间的流体动力学相互作用也非常重要，并且许多研究工作都集中在问题的这些方面（例如Delnoij等人，1997年；Fox，2012年；Gavi等人，2007年；Gerber＆Mousavi，2007年；Lain等人，2002年；Laurent＆Massot，2001年；Marchisio） \& Fox，2007；Prat \& Ducoste 等，2006；Sanyal 等，2002；例如，在气固系统中，气体动力学理论在颗粒流中的应用以及多相计算模型的发展使人们对这些问题有了更深入的理解（Gidaspow，1994；Goldhirsch，2003；Jenkins \& Mancini，1989）。一般来说，这些基于输运现象的模型由分散相和连续相之间的空间不均匀质量、动量和能量（热和颗粒）平衡组成，并且在计算流体动力学（CFD）的背景下通过使用有限体积方法（Leveque，2002）来处理空间依赖性。此类流的重要工业实例包括流化床和提升管流、浆料流反应器和鼓泡塔。在所有这些例子中，由于分散相和连续相之间的质量、动量和能量交换而产生的耦合导致了与单相流中观察到的明显不同的流动动力学。因此，用于描述分散型多相流的流体动力学的 CFD 模型通常涉及通过相相互作用项紧密耦合的多个连续性和动量方程（Drew \& Passman，1999；van der Hoef 等，2008；Ishii，1975；Portela \& Oliemans，2006）。对于多分散型多相流，情况变得更加复杂，因为需要描述分散相的“粒径”分布（De Chaise-martin et al., 2009; Fan et al., 2004; Fox, 2007; Riber et al., 2009），以及不同尺寸颗粒与连续相的耦合。当前工作的主要目标之一是提出一个系统的建模框架来完成这项任务。一般来说，上述第一类模型能够描述表征分散相的 NDF 的演化。例如，如果将此方法应用于结晶问题，则可以描述在某些操作条件下工作的结晶器中 CSD 的演变，同时考虑所有可能的物理和化学过程，例如成核、分子生长、聚集和断裂。然而，这些模型无法考虑性质的空间梯度以及它们与连续相或分散相的关系、晶体与连续相之间的流体动力学相互作用以及晶体本身之间的相互作用。因此，第一类模型无法预测局部流动相互作用对 CSD 演化的影响。另一方面，第二类模型能够描述详细的流体动力学相互作用，但为了简单起见，通常假设恒定（甚至单分散）CSD。 本书的主要目的是介绍共同的基础理论，并通过引入适当的计算方法，在两种方法之间架起一座桥梁。由此产生的基于 CFD 的计算模型可用于解决各种多分散颗粒和多相流问题。

\paragraph*{多分散型多相流的工作定义}
{\small 就本书而言，明确定义多相流中分散和多分散的含义将很有用。所谓分散，是指一个或多个相由清晰可辨的离散实体组成，例如固体颗粒、液滴或气泡。通过多分散，我们的意思是\par}
每个实体的分散相实体可以不同（例如具有不同质量、成分或温度的颗粒）。例如，燃料喷雾在喷嘴附近有一个液体射流不分散的区域，随后是主射流破碎后由单独的液滴组成的区域。后一个区域将被视为多分散型多相流，可以使用本书中描述的方法进行建模。相反，例如，气泡柱流可以完全单分散，因为所有气泡都具有（大约）相同的属性（例如气泡直径）。在实践中，单分散型多相流相对较少，因此，拥有一个自然地考虑多分散性的建模框架非常重要。 “多分散性”的另一个重要表现是存在不同速度的分散相实体（即使是具有相同物理性质的实体）。熟悉分子气体动力学理论的读者会认识到这种类型的多分散性导致了速度分布函数，它在多分散型多相流的输运理论中发挥着重要作用。与分子动力学理论一样，单分散型多相流的介观描述可以用速度分布函数的动力学方程来表达。然而，与标准温度和压力下的分子气体不同，分散相实体通常由于碰撞而很少相互作用，因此在碰撞主导的情况下有效的标准流体动力学近似不再准确。与分子气体类似，这种多相流表现为稀薄的颗粒气体，其中除了碰撞之外的过程（例如与连续相的动量交换）在确定流态方面占主导地位。因此，通常需要保留介观描述作为描述分散相流动动力学的起点。总之，多分散型多相流由一种（或多种）分散相组成，这些分散相具有可能不同的物理性质和速度。用于描述此类流动的介观尺度建模方法（第 1.3 节中更详细地描述）是本书的主要焦点。对于不由明显不同的分散相组成的多相流，必须遵循其他建模方法（例如，解决分离相的界面动力学的方法，或体积平均方法以及相相互作用和传输项的伴随的临时闭包）。正如读者在后续章节中将更加清楚的那样，介观建模方法允许对宏观尺度输运方程进行严格推导，从而确保在模拟真实多分散型多相流时所得到的 CFD 模型尽可能准确。

\section*{1.2 两个示例系统}\phantomsection\addcontentsline{toc}{section}{1.2 两个示例系统}
在本节中，我们提供了两个多分散多相系统的示例。在第一个示例中，假设颗粒非常小，因此它们的惯性小到相对于流体相的惯性可以忽略不计。因此，粒度分布 (PSD) 由 PBE 控制，并因聚集和破碎而发生变化。在第二个例子中，假设粒子足够大，相对于周围流体具有有限的惯性，我们让所有粒子相同并考虑粒子速度的分布，这由动力学方程（KE）描述。

\subsection*{1.2.1 细颗粒物的群体平衡方程}\phantomsection\addcontentsline{toc}{subsection}{1.2.1 细颗粒物的群体平衡方程}
{\small 作为第一个例子，我们考虑悬浮在液体中的固体颗粒群。我们假设粒子的密度和大小使得它们与\par}
液体。流体中的局部速度梯度（通常由剪切率量化）会引起颗粒聚集和破碎，从而改变 PSD（由相应的 PBE 决定）。在这种情况下，我们将假设（胶体系统通常是这种情况）聚集和破碎是完全可逆的，这意味着，当颗粒聚集时，形成的聚集体随后可以通过破碎而破碎，而破碎形成的碎片又可以聚集在一起。接下来，我们将分析 PBE 并引入一组无量纲数和特征时间尺度，这些将在多相系统的研究中发挥作用。

\paragraph*{相关无量纲数字和时间尺度}
如第 4 章所述，以下 PBE 控制 PSD 的演化，代表固体骨料的状态：

\begin{sourceblock}
\includegraphics[page=44,pagebox=mediabox,trim={199.50bp 453.01bp 65.07bp 220.27bp},clip,max width=0.94\textwidth,max height=0.38\textheight]{Marchisio2013.pdf}
\end{sourceblock}
其中 vf 是颗粒平流输送的已知流体速度，Г 是扩散系数。为了简单起见，在本例中，仅考虑一个空间坐标x，并且扩散系数Г是恒定的。 PSD（用 n 表示）是一个数密度函数 (NDF)，表示体积在 V 和 V + dV 之间的粒子的数浓度。 PSD 是时间 t、空间坐标 x 和聚集体体积 V 的函数。聚集体由具有固定直径 dp 的球形初级颗粒组成。如果聚集体中仅包含一个初级粒子，则 V = πd3 p/6，而对于由两个初级粒子组成的聚集体，V = πd3 p/3，对于较大的聚集体，依此类推。如第 5 章所述，由于聚合和破坏而产生的源项是

\begin{sourceblock}
\includegraphics[page=44,pagebox=mediabox,trim={102.61bp 275.49bp 65.07bp 371.95bp},clip,max width=0.94\textwidth,max height=0.38\textheight]{Marchisio2013.pdf}
\end{sourceblock}
其中 a(V, V') 和 b(V) 是聚集和破坏核，N(V|V') 是子分布函数，为了清楚起见，省略了时间和空间依赖性。正如将在后面的章节中更详细地描述的，聚集和破碎内核表达了颗粒聚集和破碎的趋势。一般来说，这些核强烈依赖于空间坐标 x，因为真实的多相系统的特征是聚集和破碎以非常不同的速率发生的区域。例如，悬浮在搅拌罐中的颗粒可能会在叶轮附近经历强烈破碎，并在其他区域经历强烈聚集。子分布函数给出的是单个聚集体破碎所产生的颗粒的尺寸分布。为了突出方程所展示的不同状态。 (1.1) 以及更适合其在每个方案中的解决方案的不同处理，PBE 必须标准化并无量纲化。为此可以使用不同的量，例如零阶矩和 PSD 之一：

\begin{sourceblock}
\includegraphics[page=44,pagebox=mediabox,trim={190.23bp 50.16bp 65.07bp 590.59bp},clip,max width=0.94\textwidth,max height=0.38\textheight]{Marchisio2013.pdf}
\end{sourceblock}
零阶矩确定总聚集体数量密度（即每单位体积的聚集体总数），而一阶矩是分散相体积分数，更常见地用 αp 表示。对于这个简单的两相系统，如果用αf来表示流体相体积分数，则定义为αp + αf = 1。通过使用这些矩，可以定义对应于一阶矩和零阶矩之比的平均颗粒体积，此外，还可以定义无量纲PSD ψ(ψ) = n (t, x, V) m1(t, x)/m2 0(t, x)，其中ψ = Vm0/m1。无量纲 PSD 具有等于 1 的一阶和二阶矩，在搜索表征纯聚集、纯破坏以及同时聚集和破坏问题的自相似解时非常有用（Ramkrishna，2000）。 PSD的二阶矩，

\begin{sourceblock}
\includegraphics[page=45,pagebox=mediabox,trim={184.92bp 483.98bp 67.56bp 182.22bp},clip,max width=0.94\textwidth,max height=0.38\textheight]{Marchisio2013.pdf}
\end{sourceblock}
相反，它可用于定义 PSD 的标准差，它等于 m2 −m2 1/m0。这里是 PSD 和 Eq。通过引入特征长度 L、速度 U、骨料体积 ζ 和总骨料数密度 Nt，使式（1.1）无量纲化。例如，对于悬浮在搅拌罐中的颗粒，L 可以是罐直径，U 可以是叶轮尖端速度，而 z 和 Nt 可以通过稳定状态下容器中的体积平均 PSD 来提取。最后两个量又可用于定义特征分散相体积分数： φp = ζNt。了解这些特征量足以计算聚合 τa = [a(z, z)Nt]−1 和断裂 τb = [b(z)]−1 的特征时间尺度。这两个时间尺度表示两个后续聚集或破坏事件之间的平均时间间隔。当这些时间尺度非常短时，聚集和破坏事件非常频繁，导致 PSD 快速演化。事实上，当粒子在域中移动和扩散时，会发生许多聚集和破碎事件，导致 PSD 快速适应其局部“平衡”值 neq(t, x, V)，该值由聚集和破碎内核的局部值决定。这里的平衡一词是指聚集和破坏相互平衡时达到的稳定状态。2由于聚集和破坏可逆的假设，这种平衡在点与点之间会有所不同，因为它取决于聚集和破坏核的局部值。平衡解对应于生成的 PSD，在粒子无法远离 x 点的情况下，但根据 x 点处的局部聚集和破碎内核随时间演化。相反，当特征时间尺度很长时，聚集和破坏事件不是很频繁，PSD的演化变得相当缓慢。因此，当颗粒在域中移动和扩散时，由于聚集和破碎频率较低以及 PSD 变化缓慢，只能感知到平均速率。 PSD 的演变由体积平均内核决定，而不是由它们的局部值决定。仔细观察聚集时间尺度可知，它与特征总聚集数密度Nt成反比，而特征总聚集数密度Nt与特征分散相体积分数φp有关。当系统稀释时（即 φp ≪1），Nt 为
2 根据聚合和破坏内核的功能形式，可能出现不同的情况。当聚合盛行时，系统可以进行广泛的聚合。由于当颗粒聚集时，其数量减少，因此可以达到胶凝极限，其中颗粒总数趋于零。相反，当破碎普遍存在时，颗粒总数可能达到无穷大，从而给出相反的极限，即破碎。当聚集和破碎相互平衡时，粒子群可以演化到稳定状态。这种稳定状态可以通过设置方程中的碰撞源项来计算。 (1.2) 等于零。我们将此稳态解称为平衡 PSD。

小，因此聚合时间尺度很大。在这些条件下，聚合事件被很长的时间间隔分开。事实上，对于颗粒很少的稀释系统，需要更长的时间间隔来观察两个颗粒之间的碰撞。相反，随着系统变得更加集中（即随着 Nt 和 φp 增加），两个聚集事件之间的时间间隔变得非常短，因为随着更多的颗粒悬浮在流体中，将会发生更多的碰撞（和聚集事件）。特征分散相体积分数也会影响破碎时间尺度；事实上，随着聚合事件变得更加频繁，所得到的特征聚合体积 ζ 也会增加。由于破碎核通常随着颗粒尺寸的增加而增加（即较大的聚集体比较小的聚集体更容易破碎），因此破碎时间尺度也变短。总之，稀系统的特点是聚集和破坏时间尺度大，而对于稠密系统，这些时间尺度非常短。只有通过分析归一化和无量纲 PBE，才能对稀系统和稠密系统进行更定量的定义。通过令 ε = V/ze、n*= nze/Nt、x*= x/L、v* f = vf/U 和 t*= tΓ/L2，可获得以下无量纲 PBE：

\begin{sourceblock}
\includegraphics[page=46,pagebox=mediabox,trim={87.08bp 368.59bp 65.07bp 241.18bp},clip,max width=0.94\textwidth,max height=0.38\textheight]{Marchisio2013.pdf}
\end{sourceblock}
其中归一化的聚集和破碎核由 a*(xi, xi′) = a(V, V′)/ a(ze, ze) 和 b*(ze) = b(V)/b(ze) 给出。 P ́eclet 数，定义为

\begin{sourceblock}
\includegraphics[page=46,pagebox=mediabox,trim={209.64bp 304.92bp 65.07bp 368.47bp},clip,max width=0.94\textwidth,max height=0.38\textheight]{Marchisio2013.pdf}
\end{sourceblock}
是粒子扩散和粒子平流时间尺度之间的比率。聚合 Damk ̈ohler 数，定义为

\begin{sourceblock}
\includegraphics[page=46,pagebox=mediabox,trim={176.11bp 241.16bp 65.07bp 431.33bp},clip,max width=0.94\textwidth,max height=0.38\textheight]{Marchisio2013.pdf}
\end{sourceblock}
是粒子扩散和粒子聚集时间尺度之间的比率。破损 Damk ̈ohler 数，定义为

\begin{sourceblock}
\includegraphics[page=46,pagebox=mediabox,trim={197.52bp 178.63bp 65.07bp 494.65bp},clip,max width=0.94\textwidth,max height=0.38\textheight]{Marchisio2013.pdf}
\end{sourceblock}
是粒子扩散和破碎时间尺度之间的比率。根据 Pe、Daa 和 Dab 的值，可以识别不同的状态。例如，当 Pe 变得远小于 1 时，扩散粒子传输变得比平流传输更重要。在这种情况下，根据 Daa 和 Dab 的值，仍然可能出现不同的情况。当两个数字都远小于单位时，扩散比聚集和破碎快得多，而当两个数字都远大于单位时，聚集和破碎比粒子扩散快。这里重要的是要记住，粒子扩散往往会使空间域中的 PSD 均质化，而由于内核通常包含某种形式的空间依赖性，因此聚集和破坏往往会在 PSD 中产生梯度。

\paragraph*{主要流动状态}
从上面的讨论中，可以从等式中识别出三种主要的操作方式。 (1.5)。

(i)在φp<1的稀释条件下，聚集和破碎比粒子扩散慢，或者换句话说，Daa<1和Dab<1。 PSD 的演变不是由聚合和破碎内核的局部值控制，而是由它们的体积平均值控制。在稀释条件下，由于扩散混合颗粒的速度比颗粒聚集和破碎的速度快，因此两相系统通常被认为是空间均匀的或混合良好的（即域不同点的 PSD 相同）。混合良好的系统 ˆn(t, V) 的 PSD 演化由体积平均 PBE 控制：

\begin{sourceblock}
\includegraphics[page=47,pagebox=mediabox,trim={110.85bp 446.57bp 67.56bp 200.86bp},clip,max width=0.94\textwidth,max height=0.38\textheight]{Marchisio2013.pdf}
\end{sourceblock}
其中 ˆa 和 ˆb 是体积平均聚集和破坏内核，由于它们对空间坐标的非线性依赖，通常与局部内核 a 和 b 有很大不同。当 Daa 和 Dab 小于 10−2 时，系统通常被认为是稀释的（因此混合良好）。根据问题的类型以及聚集和破碎内核的函数形式，这些条件会导致对特征分散相体积分数的不同约束。对于大多数应用，当 φp < 10−3 时，将运行状态视为稀释状态是合理的。

(ii) 在中等密度条件下，Daa 和 Dab 大于 10−2 但小于 102。聚集和破碎的特征时间尺度与颗粒扩散的特征时间尺度相当。在这些条件下，系统不能被认为是良好混合的，并且 PSD 在空间域上会发生很大的变化。因此，物理空间中粒子的运动及其聚集和破碎无法解耦，两相系统的演化可以通过求解方程（1）来发现。 (1.5)。

(iii) 在致密条件下，Daa 和 Dab 都大于 102，因此与颗粒扩散相比，聚集和破碎非常快（或几乎瞬时）。由于聚集和破碎的可逆性假设，当粒子在域中扩散时，它们会立即适应局部平衡解，即 neq。在这种情况下，PSD 的演化完全由局部聚集和破坏内核以及由此产生的局部稳态决定。

总之，对于经历聚集和破碎的细颗粒，应针对上述每种情况使用不同的建模方法。在稀释条件下，系统可以近似为充分混合的系统，因为 PSD 中的空间梯度不相关。系统的演化由方程控制。 (1.9) 使用体积平均聚合和破碎内核。在中等密度条件下，PSD 的空间梯度不可忽略，因此系统的演化由方程 1 控制。 (1.5)。在稠密条件下，系统根据通过设置方程右侧找到的平衡解演化。 (1.2) 等于零。

\subsection*{1.2.2 气-颗粒流动的动理学方程}\phantomsection\addcontentsline{toc}{subsection}{1.2.2 气-颗粒流动的动理学方程}
作为第二个例子，我们考虑悬浮在恒定密度气相中的单分散等温固体颗粒的动力学方程 (KE)。为了清楚起见，我们假设颗粒材料的密度明显大于气体的密度，因此仅需要流体阻力和浮力项来解释两相之间的动量交换（Maxey＆Riley，1983）。在此示例中，颗粒足够大，具有有限的惯性，因此它们的演化速度可能与气相的速度截然不同。

\paragraph*{相关无量纲数}
与输运现象的所有数学描述一样，多分散型多相流理论引入了一组与描述流动行为相关的无量纲数。根据流动的复杂性（例如，由于化学反应、碰撞等导致的物理属性的变化），无量纲数的集合可能非常大。 （动量交换物理模型的详细信息将在第 5 章中给出。）正如第 4 章中将详细描述的，可以导出分散相速度 n(t, x, v) 的数密度函数 (NDF) 的动力学方程。同样在这个例子中，为了清楚起见，我们将假设问题只有一个粒子速度分量 v，并且在物理空间中是一维的，在时间 t 处坐标为 x。此外，我们假设 NDF 已经标准化（通过将其乘以粒子的体积），使得前三个速度矩为

\begin{sourceblock}
\includegraphics[page=48,pagebox=mediabox,trim={189.78bp 261.77bp 65.07bp 338.35bp},clip,max width=0.94\textwidth,max height=0.38\textheight]{Marchisio2013.pdf}
\end{sourceblock}
其中 αp、Up 和 θp 分别是分散相的体积分数、平均速度和颗粒温度。示例动力学方程由下式给出

\begin{sourceblock}
\includegraphics[page=48,pagebox=mediabox,trim={202.25bp 193.66bp 65.07bp 480.69bp},clip,max width=0.94\textwidth,max height=0.38\textheight]{Marchisio2013.pdf}
\end{sourceblock}
其中加速度（由于浮力、重力和阻力）和碰撞项分别定义为

\begin{sourceblock}
\includegraphics[page=48,pagebox=mediabox,trim={181.16bp 97.36bp 65.07bp 548.09bp},clip,max width=0.94\textwidth,max height=0.38\textheight]{Marchisio2013.pdf}
\end{sourceblock}
{\small 在这些表达式中，gx 是重力的 x 分量，ρp 是颗粒密度，pg 是气相压力。 τp 和 τc 分别是阻力和碰撞的特征时间尺度。 Ug 是气相速度，neq 是平衡\par}
分布。为简单起见，我们使用线性碰撞模型3，其中 neq 的前两个矩与 n 的相同，第三个是

\begin{sourceblock}
\includegraphics[page=49,pagebox=mediabox,trim={171.67bp 577.69bp 67.56bp 88.77bp},clip,max width=0.94\textwidth,max height=0.38\textheight]{Marchisio2013.pdf}
\end{sourceblock}
0 ≤e ≤1 是粒子-粒子碰撞的恢复系数。对于弹性碰撞，e = 1，而 e < 1 对应于非弹性碰撞。分散相动理学方程与连续相的连续性和动量方程耦合，分别由下式给出

\begin{sourceblock}
\includegraphics[page=49,pagebox=mediabox,trim={88.00bp 460.48bp 67.56bp 185.41bp},clip,max width=0.94\textwidth,max height=0.38\textheight]{Marchisio2013.pdf}
\end{sourceblock}
其中 αg、ρg 和 νg 分别是气相体积分数、密度和运动粘度，4。根据定义，αp + αg = 1。 （1.14）无因次，我们引入特征气相速度U和特征长度L。则无因次时间为t*= tU/L。注意到 αg 是无量纲的，并令 x*= x/L、U* g = Ug/U 和 U* p = Up/U，则我们有

\begin{sourceblock}
\includegraphics[page=49,pagebox=mediabox,trim={84.28bp 322.67bp 67.56bp 317.90bp},clip,max width=0.94\textwidth,max height=0.38\textheight]{Marchisio2013.pdf}
\end{sourceblock}
\begin{sourceblock}
\includegraphics[page=49,pagebox=mediabox,trim={98.32bp 298.78bp 307.56bp 385.73bp},clip,max width=0.94\textwidth,max height=0.38\textheight]{Marchisio2013.pdf}
\end{sourceblock}
其中 p* x = gx/g。气相雷诺数和弗劳德数分别由 Reg = UL/νg 和 Frg = U/(Lg)1/2 定义，其中 g 是重力加速度常数。斯托克斯数由 St = Uτp/L 定义，相质量比由 φ2 = αpρp/(αgρg) 定义。后者是一个重要参数，因为它决定气相是否“看到”分散相（即当 φ2 <1 时动量耦合可以忽略不计）。方程给出的动力学方程。通过定义 v*= v/U 和 n*= Un 使 (1.11) 变得无量纲：

\begin{sourceblock}
\includegraphics[page=49,pagebox=mediabox,trim={181.78bp 193.82bp 67.56bp 479.31bp},clip,max width=0.94\textwidth,max height=0.38\textheight]{Marchisio2013.pdf}
\end{sourceblock}
其中 n*(t*, x*, v*) 是无量纲速度分布函数，5 和

\begin{sourceblock}
\includegraphics[page=49,pagebox=mediabox,trim={171.61bp 119.88bp 67.56bp 529.26bp},clip,max width=0.94\textwidth,max height=0.38\textheight]{Marchisio2013.pdf}
\end{sourceblock}
Knp

3 有关碰撞模型的更多详细信息，请参阅第 6 章。 4对于一维流动，粘度项通常为零。然而，我们将粘度项作为全 3D 情况的占位符，以显示雷诺数依赖性。

\footnotetext[5]{请注意，n 和 neq 被无量纲化，使得 αp =}
{\small n*dv*。由于动力学方程在 n 中是线性的，因此该过程不会生成新的无量纲数。\par}
相密度比由 φ1 = ρp/ρg 定义，并且由于 φ1 ≫1，浮力项可以忽略不计。此过程中生成的两个新的无量纲数是相速度比 φ3 = U† p/U,6 和分散相努森数 Knp = U† pτc/L。此外，方程的无量纲形式。 (1.10),

\begin{sourceblock}
\includegraphics[page=50,pagebox=mediabox,trim={171.34bp 497.83bp 65.07bp 110.53bp},clip,max width=0.94\textwidth,max height=0.38\textheight]{Marchisio2013.pdf}
\end{sourceblock}
\begin{sourceblock}
\includegraphics[page=50,pagebox=mediabox,trim={283.15bp 474.49bp 182.04bp 209.91bp},clip,max width=0.94\textwidth,max height=0.38\textheight]{Marchisio2013.pdf}
\end{sourceblock}
引入分散相马赫数 Map = U† p。与可压缩气体流类比，Map 是分散相中特征平均粒子速度 U† p 与“声速”(θ1/2 p ) 的比率。因此，对于 Map ≪1 来说，分散相中的输运主要是由于速度波动造成的，而对于 Map ≫1 来说，则是由于平均平流造成的。按照气体动力学中使用的惯例，我们将根据分散相是亚音速还是超音速来不同地定义速度比和分散相努森数：

\begin{sourceblock}
\includegraphics[page=50,pagebox=mediabox,trim={132.03bp 320.55bp 65.07bp 319.28bp},clip,max width=0.94\textwidth,max height=0.38\textheight]{Marchisio2013.pdf}
\end{sourceblock}
这里 dp 是粒径，7 g0(αp/α* p) 是径向分布函数，α* p ≈ 0.63 是密堆积时的最大体积分数。对于 αp ≪α* p，g0 ≈1。然而，当 αp 接近 α* p 时，g0 发散至无穷大。请注意，比率 φ3 将很大程度上取决于斯托克斯数，

\begin{sourceblock}
\includegraphics[page=50,pagebox=mediabox,trim={202.18bp 225.56bp 65.07bp 445.19bp},clip,max width=0.94\textwidth,max height=0.38\textheight]{Marchisio2013.pdf}
\end{sourceblock}
其中 τp 的表达式对于小粒子雷诺数有效，定义为

\begin{sourceblock}
\includegraphics[page=50,pagebox=mediabox,trim={207.66bp 170.38bp 65.07bp 502.30bp},clip,max width=0.94\textwidth,max height=0.38\textheight]{Marchisio2013.pdf}
\end{sourceblock}
当 St ≥ 1 时，动力学方程将与气相分离，并且分散相将表现为颗粒气体。在相反的极限条件下，St ≪1，θp ≈0，U† p ≈U，因此分散相马赫数将非常大，并且 φ3 ≈1。在斯托克斯数的中间值，可以观察到取决于无量纲参数的所有值的丰富多样的流动现象。

\footnotetext[6]{因为 φ3 和 Knp 一起出现在式（1）中。 （1.17），它们可以组合成一维无量纲数。我们将它们分开，以便 Knp 和 Map 仅依赖于分散相变量。}
\footnotetext[7]{在固定体积分数 αp 下，减小 dp 会增加碰撞率，因为颗粒总数浓度 N 显著增加。}
\clearpage
\begin{center}\small 表1.1。气体-颗粒流的相关无量纲数，其中 Rep ≪1 且 ρg ≪ρp\end{center}
\begin{sourceblock}
\includegraphics[page=51,pagebox=mediabox,trim={58.00bp 301.96bp 58.23bp 94.25bp},clip,width=\textwidth,max height=0.72\textheight]{Marchisio2013.pdf}
\end{sourceblock}
为方便起见，本节中导出的气体-颗粒流的相关无量纲数收集在表 1.1 中。在实践中，必须针对特定问题选择合适的U和L值。例如，它们可以由入口和/或边界条件确定。然而，一种特别令人感兴趣的情况是落在无界域中的粒子，对此方便的选择是 L = dp 和 U = U† p = |Up −Ug| = gτp|g* x| （即沉降速度）。对于这种情况，没有 θp 的源项，因此由于阻力，它在稳态下松弛为零。8 因此，分散相马赫数变得无穷大。对于沉降问题，通常使用粒子阿基米德数（参见表 1.1）来代替弗劳德数。读者可能会惊讶地发现没有专门为分散相定义的雷诺数。这是因为分散相粘度仅在 Knp <1（即碰撞主导或流体动力学状态）下得到明确定义。在此限制下，νp ∝θpτc ∝ Knp/Map，因此当 Map ≤1 时，分散相雷诺数将与 Map/Knp 成正比。然而，在许多气体-颗粒流中，即使对于αp∼0.1，分散相努森数也不会很小，因为阻力和非弹性碰撞会大大降低颗粒温度（以及由此产生的碰撞频率）。相比之下，标准温度和压力下的分子气体的 Knp <1，即使分子所占的体积分数约为 0.001。这个事实可以是

\footnotetext[8]{流体动力相互作用（未包含在本例的动力学方程中）可能导致有限的 θp。然而，在气体-颗粒流中，分散相马赫数通常非常大。}
这是因为分子的高均方根 (RMS) 速度。事实上，如果将典型气体-颗粒流中的颗粒温度转换为开尔文温度，那么与室温相比，它会非常小。

\paragraph*{主要流态}
鉴于表 1.1 中的无量纲参数数量相对较多，因此在气体-颗粒流中观察到的不同流态的数量也很大也就不足为奇了。对于其他多分散型多相流来说通常也是如此。例如，如果我们考虑斯托克斯数非常小的流动（St ≪1），使得分散相速度几乎与连续相速度相同，那么我们可以期望观察到与 φ1 ≈1 时恒定密度单相流观察到的流态相似的流态。然而，如果相密度非常不同，并且 φ2 存在重要的空间变化（例如，分层流的顶部有一个大的 φ2（相对于重力），而小 φ2 的顶部），那么人们就会观察到典型的变密度单相流的流态。尽管存在大量可能的流态，但仍可以确定一些具有分散型多相流特征的主要流态。因此，本节的目的是向读者介绍多相流文献中常见的主要流态。在许多关于分散型多相流的讨论中，分散相体积分数 αp 的值用于区分主要流态（Balachandar \& Eaton，2010）。从表1.1中，我们可以观察到αp出现在φ2、Knp和g0中。如前所述，φ2 与分散相和连续相之间的动量耦合有关，并且在所有其他参数固定的情况下，可以通过降低 αp 来减小 φ2。分散相努森数和径向分布函数通过碰撞项进入控制动力学方程。正如第 6 章详细讨论的，碰撞项对多分散系统的平均动量方程有两个贡献：（1）“纯”点状碰撞项（对于单分散系统为零）和（2）由于粒子的有限直径而产生的碰撞通量项。当 g0 ≈1（即 αp ≪α* p）时，第二个贡献可以忽略不计。因此，一般来说，Knp 的值表示碰撞的相对重要性（即当 Knp <1 时，流动是碰撞主导的，而当 Knp = ∞ 时，则是无碰撞的），并且增加 αp 通常会通过降低 Knp 来增加碰撞率。然而，值得注意的是，当恢复系数 e 明显小于 1 时，颗粒温度也会受到 αp 的影响。事实上，由于碰撞过程中的能量损失，颗粒温度可能会降至零，在这种情况下，分散相努森数将趋于无穷大。 9 因此，一般来说，αp 的值不一定是碰撞速率的可靠指标，因为后者也取决于颗粒温度。另一个常用于区分流态的参数是斯托克斯数 St。在 Stokes 数非常大的情况下，粒子加速度（方程（1.17））不依赖于连续相速度，分散相表现为颗粒气体。在小 St 的限制下，存在两个由临界斯托克斯数 Stc 分隔的重要状态（Jabin，2002）。对于 St ≤ Stc，任何两个粒子的轨迹永远不会相互交叉，这意味着在给定时间流中的任何点只能观察到一个粒子速度。 这种流动称为单动流动（Massot，2007）或平衡流动（Ferry \& Balachandar，2001），并且由于速度分布是 delta 函数，因此粒状流动

\footnotetext[9]{由于初始条件，对于致密弹性系统，分散相克努森数也可能是无穷大。例如，如果释放出最初在真空中静止的密集粒子集合，这些粒子将由于重力而加速，使得它们都具有相同的速度。对于这种情况，颗粒温度为零且 Knp = ∞。}
这种流动中的温度为零。10 对于 St > Stc，颗粒轨迹交叉是可能的，因此即使最初为零，也会局部生成非零颗粒温度。就 αp 和 St 而言，文献中讨论的一些主要流态如下。

(i) 非常稀。对于 αp <1 和 φ2 <1，与连续相的碰撞和耦合都不重要。这种机制通常被称为单向耦合。非常稀的区域可以根据斯托克斯数进一步细分。

(a) St ≪Stc。颗粒具有与连续相相同的速度。这种流动状态通常被称为含尘气体。

(b)St≤Stc。粒子的速度与连续相略有不同，但局部是唯一的。

(c) St > Stc。由于粒子的有限惯性，粒子速度分布的动力学状态是多值的。

(d) St ≫1。无碰撞颗粒气体。

(ii) 稀释。对于 αp <1 且 φ2 ≥1，与连续相的碰撞和/或耦合是不可忽略的。当碰撞可以忽略不计时（Knp≫1），这种状态被称为双向耦合。当碰撞很严重时（Knp ≤ 1），这种状态称为四向耦合。碰撞占主导地位的情况（Knp <1）被称为流体动力学状态。稀释区域可以使用斯托克斯数进一步细分，其结果与非常稀释区域的结果相似，但有一个重要的区别。由于 φ2 ≥1，动量耦合可能导致新的流动不稳定性，而这种不稳定性在非常稀薄的区域中是观察不到的。例如，稀流在雷诺数 Reg 下可能会变成湍流，而非常稀的流则为层流（Passalacqua 等，2010）。

(iii) 中等密度。当 αp < α* p 且 g0 > 1 时，碰撞通量项不可忽略。对于流化床等气体颗粒流，当 0.02 < αp < α* p 时，就会出现这种情况。 φ1 ≈1（例如浆料）和 Reg ≪1 的液体-颗粒流通常处于中等密度状态，但由于长程流体动力相互作用，颗粒-颗粒相互作用比简单的二元碰撞更复杂。例如，流体动力相互作用改变了双粒子径向分布函数的函数形式（Ladd，1993）。

（四）密集。当 αp ≈α* p 时，粒子与粒子之间的接触是持续的，不能用瞬时碰撞来表示。然后，二元碰撞模型必须替换为一个模型，该模型可以解释“力链”中发生的粒子之间的长程相互作用（Aarons 等，2010）。

上面讨论的流态通常用于描述 φ1 <1 的系统。当 φ1 ≫1 时，会出现其他流态（例如气泡流），因为额外的力模型（例如增加的质量）引入了新的无量纲参数，并且因为连续相的动量变得占主导地位。例如，在气泡流中，当 Rep > 1 时，湍流液体在气泡后面苏醒，导致“伪湍流”（Mudde 等人，2009 年；Riboux 等人，2010 年；Sato 和 Sekoguchi，1975 年），从而改变了连续相中气泡与气泡相互作用的性质。

\footnotetext[10]{请注意，如果由于颗粒直径为亚微米而斯托克斯数非常小，则布朗运动引起的速度波动将很重要。}
最后，我们还应该注意到分散型多相流可能是湍流。事实上，工业和环境应用中遇到的许多多相流都是湍流。虽然多相湍流模型不是本书的重点，但它们是具有相当实际重要性的主题。一般来说，在上述非常稀的区域中，连续相的湍流模型将与单相流完全相同。然而，有必要修改阻力项中出现的流体速度，以解决未解决的流体速度波动的影响（Buyevich，1972；Derevich＆Zaichik，1990；Minier＆Peirano，2001；Pandya＆Mashayek，2003a；Reeks，1977，1991；Zaichik等人，2003， 2004 年，2008 年；扎奇克和阿利普琴科夫，1998 年）。另一方面，在稀状态下，连续相的单相湍流模型必须进行修改，以包括相之间动量交换（平均和波动）的影响（Sommerfeld \& Ho，2003）。更一般地，当 φ2 ≥1 时，由于体积分数梯度，可能会出现新的湍流流态，而体积分数梯度是由流动动力学自发产生的。这些真正的多相湍流流态表现出单相流中不存在的湍流产生机制，这使得它们更难以建模。例如，由于阻力定律对分散相体积分数的非线性依赖性，气体-颗粒流可能会变得湍流。由于这种不稳定性首先发生在相对较小的尺度上，然后发展到更大的尺度，因此无法成功地将其建模为从积分尺度到耗散尺度的简单湍流能量级联。因此，至少有必要在湍流模型中添加新项，以解释多相流中存在的附加物理现象。在这项工作中，我们将重点关注分散型多相流基本方程的推导，这些方程随后可以用作使用经典雷诺平均和大涡模拟方法推导湍流模型的起点（Fox，2012；Pope，2000）。然而，随着我们的进展，我们将指出模型方程与纳维-斯托克斯方程相比在形式上的关键差异，并评论这些差异将如何影响湍流建模过程。

\section*{1.3 介观尺度建模方法}\phantomsection\addcontentsline{toc}{section}{1.3 介观尺度建模方法}
我们将用于推导多分散型多相流的基本控制方程的建模方法将处于介观尺度水平。建模框架的示意图如图 1.1 所示。我们将使用术语“微尺度”来描述在单个粒子水平上完全解析所有相关长度和时间尺度的建模方法。例如，流体-颗粒流的微观模型将解析连续相中的所有流动结构，包括颗粒表面周围的边界层，以及系统中所有颗粒对的颗粒-颗粒碰撞。在单相流的背景下，微尺度模型被称为直接数值模拟（DNS）。对于多相系统，微尺度模型对每个相使用 DNS，并完全解析相间耦合项。术语宏观尺度将用于表示采用分散相的流体动力学描述的多相模型。此类模型也称为多流体模型（因为分散相被视为有效流体）或欧拉-欧拉模型。后者的名称来自于分散相的数值处理（即固定网格上的离散化），而不是欧拉-拉格朗日模型，其中分散相在拉格朗日框架中作为离散实体进行跟踪。我们应该注意到，在

\begin{sourceblock}
\includegraphics[page=55,pagebox=mediabox,trim={58.00bp 394.80bp 58.23bp 59.00bp},clip,width=\textwidth,max height=0.72\textheight]{Marchisio2013.pdf}
\par\vspace{0.45\baselineskip}
\small 图1.1。分散流建模框架。
\end{sourceblock}
在介观模型的背景下，分散相的拉格朗日处理本质上是动力学方程解的随机表示（例如，相当于动力学理论中使用的直接模拟蒙特卡罗（DSMC）方法（Bird，1994））。这可以与基于分子动力学 (MD) 模拟的颗粒流微观模型形成对比。对于后者，不能保证闭合的动力学方程（由于碰撞项中存在双粒子分布函数），而前者在动力学方程闭合后应用。一般来说，可以使用体积平均或系综平均方法从微观模型中找到宏观模型（Drew，1971；Drew \& Passman，1999），或者首先推导介观模型，然后评估所得动力学方程的矩。第一步（推导介观模型）主要涉及物理建模和基于物理的简化，以获得封闭的动力学方程。第二步（从介观模型导出宏观模型）涉及数学近似以封闭矩输运方程。换句话说，如果一个定义良好的宏观模型再现了从介观尺度模型的精确解中找到的矩，则该模型被认为是准确的。相反，介观模型的准确性取决于从微观模型的解导出的物理模型的保真度。介观尺度建模方法中存在的物理和数学建模误差的这种划分在体积平均和集合平均方法中不存在。在后一种方法中，人们直接从微尺度输运方程导出矩输运方程（尽管它们通常不是指矩，而是平均质量和平均动量等），并直接对出现在矩输运方程中的未知项建立闭包（或“本构”方程）。物理误差和数学误差之间缺乏分离使得很难使用平均方法提出准确且适定的宏观模型。相比之下，通过首先推导介观模型，可以通过与微观模型的解决方案直接比较来隔离物理建模错误。

正如我们将在后面的章节中看到的，对于多分散型多相流，建模误差的分离对于获得适定的宏观模型非常有价值。类似的情况尤其出现在湍流反应流（Fox，2003）和一般湍流（Pope，2000）中。通过采用概率密度函数（PDF）方法，可以写出速度和成分变量的联合 PDF 的输运方程，这相当于动力学方程。 PDF 输运方程的闭合涉及分子混合（或微混合）等物理过程的模型。因此，PDF输运方程是湍流的介观尺度模型，其矩是经典湍流模型中出现的雷诺平均量。对于湍流反应流，化学源项在介观模型中是封闭的，但在宏观模型中是不封闭的。后者需要数学闭包，这在本质上不同于微混合模型所使用的物理闭包。由于多相流中存在额外的物理场，导出宏观方程的直接闭包比单相湍流更困难（Fox，2012）。湍流 PDF 方法与本书中描述的介观尺度模型之间的另一个重要区别是，前者处理由湍流引起的系综平均值，而后者是非湍流多相系统的动力学理论模型。事实上，我们可以从本书中导出的基本宏观模型方程出发，导出湍流多相流的 PDF 方法。

\subsection*{1.3.1 与微尺度模型的关系}\phantomsection\addcontentsline{toc}{subsection}{1.3.1 与微尺度模型的关系}
如图 1.2 所示，介观尺度方法的关键建模步骤是 (1) 定义介观尺度变量，以及 (2) 推导一个模型，了解一个粒子的介观尺度变量如何因涉及所有粒子的微观物理现象而变化。在某些情况下，介观变量的选择很简单。例如，由于质量和

\begin{sourceblock}
\includegraphics[page=56,pagebox=mediabox,trim={58.00bp 69.07bp 58.23bp 408.84bp},clip,width=\textwidth,max height=0.72\textheight]{Marchisio2013.pdf}
\par\vspace{0.45\baselineskip}
\small 图1.2。从微观描述到介观模型。
\end{sourceblock}
动量守恒、颗粒质量和速度通常用作介观变量，以及用于传质和/或化学反应系统的颗粒温度和化学成分。在某些情况下，可能还需要定义考虑颗粒性质（例如形状、内部结构等）的介观尺度变量。在任何情况下，介观尺度变量的选择与推导这些变量如何针对单个粒子变化的模型的第二步密切相关。对于质量和动量变化，开发此类模型通常很简单，但对于其他类型的介观尺度变量而言并非总是如此。正如第 5 章所讨论的，此类模型通常包含一个随机分量来表示微尺度模型中存在的未解析尺度。一个典型的例子是由流动中相邻粒子的湍流尾流产生的粒子速度的随机波动。介观建模工作的主要目标是在单个粒子的介观模型方程中尽可能多地捕获微观物理现象。介观尺度模型的验证可以使用微尺度模型的数值解来进行（Tenneti et al., 2010），其方式与 DNS 用于湍流单相流模型验证的方式大致相同。典型的介观建模策略由四个步骤组成。

1. 利用基于微观尺度描述和对物理学的深入理解的见解，开发介观模型。

2. 设计一个微观“数值实验”，测试介观尺度模型开发中的关键假设。

3.利用微观模拟结果验证介观模型的预测。

4. 返回步骤 1，通过应用从验证过程中发现的差异中获得的新见解来改进介观模型。

请注意，这些步骤的顺序很重要。事实上，成功的介观尺度模型很少是通过简单地研究微尺度模拟的结果而产生的，而微尺度模拟并不是专门为测试特定的假设或假说而设计的。一般来说，如果在每个时间步都存储所有流动变量，微尺度模拟将产生大量数据，并且在没有对物理学的先验理解的情况下期望这些数据集将“揭示”介观模型的正确形式是不现实的。另一方面，微尺度模拟提供了一种宝贵且通常独特的工具，用于整理所提出的介观尺度模型的有效性，因为它们包含所有微尺度变量，其中许多变量在实验室实验中无法测量。总之，微观尺度描述提供了介观尺度模型开发所需的两条重要信息。首先，微尺度模型的数学公式包括完整描述分散型多相流所需的所有相关物理学，为了解需要哪些介观变量以及这些变量如何在介观尺度上相互作用提供了有价值的见解。这些见解用于制定介观模型。其次，微尺度模型的详细数值解直接用于验证所提出的介观尺度模型。当观察到介观模型预测与微观模拟之间存在显著偏差时，这些差异会导致介观模型的重新表述，以改进物理描述。请注意，重要的是要记住，此验证步骤应通过将介观尺度模型的精确解与微观尺度结果进行比较来完成，而不是通过求解宏观模型得到的近似解来完成。如前所述，宏观尺度模型包括介观尺度水平上不存在的额外闭合步骤，因此微观尺度和宏观尺度模型之间的分歧可能仅仅是由于这种闭合的不足。

\subsection*{1.3.2 数密度函数}\phantomsection\addcontentsline{toc}{subsection}{1.3.2 数密度函数}
在介观尺度模型的制定中，数密度函数（NDF）起着关键作用。因此，我们在第 2 章中详细讨论了 NDF 的属性。换言之，NDF 是在给定的介观尺度变量值集下每单位体积的粒子数。由于在任何时刻，微尺度粒子都会具有一组独特的微尺度变量，因此 NDF 也称为单粒子 NDF。一般来说，单粒子 NDF 仅对于介观变量的可实现值才非零。换句话说，可实现的介观尺度值是在微观尺度模拟中出现的所有粒子的集合中观察到的值。相比之下，在微观尺度模拟中从未观察到的一组介观尺度值是无法实现的。可实现性限制可能因多种原因而发生，例如由于质量、动量、能量等守恒，是微尺度模型的固有属性。还需要注意的是，介观尺度值通常具有很强的相关性。我们的意思是，任何两个介观变量的 NDF 都不能根据每个变量单独的 NDF 的知识来重建。因此，通过构造，单粒子 NDF 包含仅一个粒子的介观尺度变量之间的所有潜在相关性。正如第 2 章所讨论的，单粒子 NDF 通常不提供微尺度系统的完整描述。例如，包含 N 个粒子的微尺度系统将由 N 粒子 NDF 完全描述。这是因为原则上任何一个粒子的介观尺度变量都会受到所有 N 个粒子的介观尺度变量的影响。或者，换句话说，N组介观尺度变量可以彼此相关。例如，通过二元碰撞相互作用的粒子系统表现出碰撞前后两个粒子的速度之间的相关性。因此，由于碰撞，单粒子 NDF 速度的时间演化将涉及双粒子 NDF。在介观尺度建模方法中，主要的物理建模步骤涉及通过单粒子 NDF 的函数来近似 N 粒子 NDF（即精确的微尺度模型）。一个典型的例子是通过用两个单粒子 NDF 的乘积来近似二粒子 NDF 来封闭碰撞项（参见第 6 章）。对于多分散型多相流，单粒子 NDF 包含介观尺度变量之间的所有相关性，这一事实使得介观尺度建模方法对于推导宏观尺度模型非常有用。例如，气体-颗粒流中的颗粒尺寸对斯托克斯数有很大的影响，因此由于与气相的耦合而对颗粒的加速度有很大的影响。此外，如果颗粒尺寸由于聚集或破碎而发生变化，则颗粒尺寸和速度之间的耦合将相当复杂。由于 NDF 可以准确描述这种复杂的相关性，因此介观尺度模型将能够预测平均颗粒尺寸和平均颗粒速度的演变。相反，如果试图直接导出平均颗粒尺寸的宏观模型，则很难使用临时闭包捕获与颗粒速度的强耦合。

{\small 总之，描述多分散型多相流的介观变量之间的相关性的复杂性要求我们采用能够准确描述这种相关性的数学公式。单粒子 NDF 包含所有\par}
一个粒子中介观尺度变量之间的相关性，因此非常适合开发介观尺度建模框架。

\subsection*{1.3.3 分散相动理学方程}\phantomsection\addcontentsline{toc}{subsection}{1.3.3 分散相动理学方程}
从数学上来说，单粒子 NDF 由动力学方程控制。然而，对于特定的应用，控制方程有其他名称。例如，在动力学理论中，速度 NDF 由玻尔兹曼方程控制（Boltzmann，1872；Cercignani 等，1994），而在气溶胶科学中，相关动力学方程被称为一般粒子动力学方程（Friedlander，2000；Gelbard \& Sein-feld，1978），而在喷雾领域，它被标记为威廉姆斯-玻尔兹曼方程（Williams， 1985）。在这项工作中，我们将不包括速度但包括粒子尺寸等介观尺度变量的动力学方程称为群体平衡方程（PBE），将包括速度和“内部”介观尺度变量的动力学方程称为广义群体平衡方程（GPBE）。虽然从数学角度来看这种划分有些任意，但它具有强调基础分散相的多分散性质的优点。在介观变量仅限于粒子速度的示例中，我们将遵循动力学理论中使用的术语。在介观尺度水平上，NDF中出现的自变量是时间、物理位置和介观尺度变量。为了区分物理点和“介观”点，我们引入相空间的概念。因此，物理点在物理空间中呈现出值，而介观尺度点在相空间中呈现出值。然而，两种类型的变量都作为自变量（与时间一起）出现在动力学方程中，因此 NDF 是自变量的函数。相比之下，宏观模型只有时间和空间作为自变量。因此，可以通过在相空间上积分来从动力学方程中找到宏观模型，以消除对介观变量的依赖性。这个“在相空间上积分”的过程与评估 NDF 的矩有关。例如，在1.2.1节中，我们使用了粒径NDF的前两个矩与代表总聚集体数密度和分散相体积分数的宏观变量之间的关系；在 1.2.2 节中，我们使用速度 NDF 的前三个矩与代表分散相体积分数、平均颗粒速度和颗粒温度的宏观变量之间的关系。介观变量的动力学方程包含影响分散相的每个微观物理过程的项。例如，在第 1.2.2 节中，气体-颗粒流的动力学方程包含由于浮力、重力和流体阻力引起的加速度以及由于碰撞引起的速度变化的项。在介观尺度建模的背景下，这些术语中的每一个都是底层微尺度物理的模型，每一个都可以通过与微尺度模拟进行比较来验证。一般来说，当添加新的介观变量（例如粒径）时，新的项将添加到动力学方程中以解释该变量的变化。重要的是要认识到，如果动力学方程中的项是根据自变量制定的并且是单粒子 NDF 的函数，则称它们是封闭的。例如，独立作用于每个粒子的重力等过程是 NDF 的线性函数。 介观建模方法的最终目标是得到一个封闭的动力学方程，准确地描述微尺度物理对介观变量的影响。

\subsection*{1.3.4 介观尺度层面的闭合}\phantomsection\addcontentsline{toc}{subsection}{1.3.4 介观尺度层面的闭合}
在大多数情况下，动力学方程中项的闭合需要先了解介观尺度变量如何受到基础物理的影响。以流体阻力项为例，最简单的模型具有以下形式

\begin{sourceblock}
\includegraphics[page=60,pagebox=mediabox,trim={211.89bp 549.25bp 65.07bp 124.59bp},clip,max width=0.94\textwidth,max height=0.38\textheight]{Marchisio2013.pdf}
\end{sourceblock}
其中，对于无限域 (αp = 0) 中单个球体周围斯托克斯流 (Rep = 0) 的特殊情况，τp 是常数（参见方程 (1.20)）。对于有限 Rep，τp 取决于阻力系数，该系数是 |Ug−v| 的函数。同样，对于 αp > 0，阻力模型取决于分散相体积分数（即 Ad 随着 αp 的增加而增加）。改变粒子形状（或其他介观变量）通常也会影响阻力模型。因此，利用微观模拟的结果，有必要制定类似于方程（1）的阻力模型。 （1.22）这取决于（可能）所有介观尺度变量。 （有关阻力模型的更详细讨论，请参阅第 5 章。）从微尺度方程制定介观尺度模型的过程广泛用于输运现象（Ferziger \& Kaper，1972）。例如，分散相和流体之间的传热取决于努塞尔数，传质取决于舍伍德数。努塞尔特数和舍伍德数与介观尺度变量和 NDF 矩（例如平均颗粒温度和平均颗粒浓度）的相关性可在文献中找到。随着微尺度模拟变得越来越复杂，基于微尺度结果的修正相关性将变得越来越普遍（Beetstra 等，2007；Holloway 等，2010；Tenneti 等，2010）。请注意，由于动力学方程需要在相空间中局部有效的介观尺度模型（即对于一组特定的介观尺度变量），而不是从宏观尺度变量中找到平均相关性，因此微观尺度模型的直接数值模拟可能是获得彻底验证此类模型所需数据的唯一方法。例如，宏观模型将取决于平均阻力，其表示为

\begin{sourceblock}
\includegraphics[page=60,pagebox=mediabox,trim={168.59bp 241.74bp 65.07bp 432.09bp},clip,max width=0.94\textwidth,max height=0.38\textheight]{Marchisio2013.pdf}
\end{sourceblock}
其中 〈τp〉 是仅取决于平均量（例如 |Ug −Up|）的宏观阻力时间尺度。相反，τp 取决于介观变量，因此可能包含对流体阻力的更完整的描述。同样的评论也适用于 NDF 动力学方程中出现的其他介观项。

\subsection*{1.3.5 与宏观模型的关系}\phantomsection\addcontentsline{toc}{subsection}{1.3.5 与宏观模型的关系}
如上所述，宏观模型是根据 NDF 低阶矩的输运方程编写的。不同类型的矩将在第 2 章和第 4 章中讨论。然而，单分散颗粒宏观模型中通常出现的低阶矩是分散相体积分数、分散相平均速度和分散相颗粒温度。当颗粒是多分散的时，PSD 的描述（至少）需要颗粒尺寸的平均值和标准偏差，或者换句话说，PSD 的前三个矩。然而，对 PSD 的更完整的描述将需要更大的粒子尺寸矩集，当斯托克斯数非零时，其输运方程通常以非平凡的方式与速度矩耦合。

宏观模型中出现的输运方程可以使用感兴趣矩的定义从动力学方程导出。例如，如果感兴趣的时刻是分散相体积分数，那么对介观变量进行积分就足够了。 （有关此过程的详细讨论，请参阅第 4.3 节。）使用第 1.2.2 节中的速度分布函数作为示例，此过程产生

\begin{sourceblock}
\includegraphics[page=61,pagebox=mediabox,trim={72.53bp 519.30bp 67.56bp 146.65bp},clip,max width=0.94\textwidth,max height=0.38\textheight]{Marchisio2013.pdf}
\end{sourceblock}
或者，使用等式中的定义。 (1.10),

\begin{sourceblock}
\includegraphics[page=61,pagebox=mediabox,trim={76.19bp 433.92bp 67.56bp 205.35bp},clip,max width=0.94\textwidth,max height=0.38\textheight]{Marchisio2013.pdf}
\end{sourceblock}
由于碰撞期间质量守恒，涉及 C 的积分为零。使用数密度函数可积并因此有界于 |v| 的事实来简化涉及 A 的积分。 → 无穷大：

\begin{sourceblock}
\includegraphics[page=61,pagebox=mediabox,trim={76.65bp 342.74bp 67.56bp 313.39bp},clip,max width=0.94\textwidth,max height=0.38\textheight]{Marchisio2013.pdf}
\end{sourceblock}
因此，等式。 (1.25) 简化为分散相体积分数的连续性方程。然而，分散相平均速度 Up 是未知的，必须通过求解单独的传输方程来找到。由速度分布函数的一阶矩可得分散相平均动量方程：

\begin{sourceblock}
\includegraphics[page=61,pagebox=mediabox,trim={76.64bp 223.59bp 67.56bp 431.33bp},clip,max width=0.94\textwidth,max height=0.38\textheight]{Marchisio2013.pdf}
\end{sourceblock}
或者，使用等式中的定义。 (1.10),

\begin{sourceblock}
\includegraphics[page=61,pagebox=mediabox,trim={76.19bp 136.29bp 67.56bp 502.97bp},clip,max width=0.94\textwidth,max height=0.38\textheight]{Marchisio2013.pdf}
\end{sourceblock}
由于碰撞期间动量守恒，涉及 C 的积分再次为零。使用分部积分来简化涉及 A 的积分：

\begin{sourceblock}
\includegraphics[page=61,pagebox=mediabox,trim={80.42bp 56.53bp 67.56bp 596.61bp},clip,max width=0.94\textwidth,max height=0.38\textheight]{Marchisio2013.pdf}
\end{sourceblock}
因此，等式。 (1.28) 简化为分散相平均动量的输运方程：

\begin{sourceblock}
\includegraphics[page=62,pagebox=mediabox,trim={99.03bp 588.79bp 65.07bp 83.07bp},clip,max width=0.94\textwidth,max height=0.38\textheight]{Marchisio2013.pdf}
\end{sourceblock}
然而，分散相颗粒温度 θp 和平均加速度 〈A〉 未知，必须使用矩闭包进行建模。这个例子说明了使用宏观模型时面临的经典问题（Struchtrup，2005）。无论矩的输运方程如何导出，它们始终包含依赖于高阶矩的未封闭项（例如 Up 取决于 θp 等）。相比之下，NDF 动力学方程的解包含有关所有可能矩的信息。换句话说，如果我们可以直接计算 n(t, x, v)，则无需使用宏观模型方程。那么显而易见的问题就出现了：“为什么我们不简单地求解介观模型的动力学方程，而是使用宏观模型呢？”这个问题的答案主要是由于存在大量自变量而导致求解动力学方程的计算成本。在 3D 速度分布函数 n(t, x, v) 的最简单示例中，自变量的数量为 1 + 3 + 3 = 7。然而，对于多分散型多相流，介观尺度变量的数量可能远大于 3。相比之下，矩输运方程涉及四个自变量（物理空间和时间）。此外，矩输运方程的形式使得它们可以很容易地集成到标准计算流体动力学（CFD）代码中。动力学方程的直接求解器构建起来要困难得多，如果要获得准确的结果，则需要专门的数值方法（Filbet \& Russo，2003）。例如，对于直接求解器，有必要对所有相空间进行离散化，因为先验 n 的非零值的位置是未知的，当相空间无界时，这可能会非常昂贵。在实践中，模拟介观尺度模型最常用的方法是采用统计描述，其中 NDF 通过计算“地块”的有限集合来近似（Bird，1994）。在多相流文献中，此类方法被称为欧拉-拉格朗日模型，因为连续相使用欧拉公式处理，而离散相则使用拉格朗日粒子跟踪方法。对于相间耦合较弱或不存在的情况（例如非常稀的流动），通过增加地块数量以最小化统计误差，可以使欧拉-拉格朗日模拟变得相当准确。然而，当相之间的耦合很强时，即使分散相模拟中很小的统计误差也可能导致与从动力学方程的直接解中观察到的流态性质不同的流态（Passalacqua 等人，2010）。此外，重要的是要认识到，在地块数量固定的情况下，欧拉-拉格朗日模拟计算出的统计数据的准确性将随着介观尺度变量数量的增加而降低。因此，尽管在欧拉-拉格朗日模拟代码中增加介观变量的数量通常很简单，但这并不意味着统计精度不成问题。欧拉-拉格朗日模拟的另一个挑战是需要对连续相输运方程中出现的耦合项进行时间精确的统计（而不仅仅是时间平均值）。 由于所有这些原因，人们有强烈的动机去寻找精确的方法来闭合矩输运方程。

\section*{1.4 矩输运方程的闭合方法}\phantomsection\addcontentsline{toc}{section}{1.4 矩输运方程的闭合方法}
如上所述，对于复杂的工业和环境流，直接求解 NDF 的介观尺度模型几乎是不可能的。相反，我们希望导出精确的宏观模型，其解与从介观模型解中找到的矩一致。由于宏观模型是根据 NDF 矩的传递方程制定的，因此我们需要找到矩输运方程的精确闭包，最好仅使用尽可能小的矩集。从动力学理论（Chapman \& Cowling，1961），我们知道，对于碰撞时间尺度远小于所有其他过程的情况（即当分散相努森数非常小时），可以正式推导仅使用质量、动量和总能量的五矩方程的流体动力学模型。然而，当碰撞时间尺度不小时，有必要开发对所有分散相努森数值都有效的矩闭包（Grad，1949b；Struchtrup，2005）。在本节中，我们简要概述这两种闭包方法背后的概念基础。后面的章节将使用具体示例介绍更多详细信息。为简单起见，我们将在本节中将注意力集中在速度分布函数和相应的玻尔兹曼方程上。

\subsection*{1.4.1 流体动力学模型}\phantomsection\addcontentsline{toc}{subsection}{1.4.1 流体动力学模型}
流体动力学模型源自介观尺度模型（例如玻尔兹曼方程），使用努森数幂的 Chapman-Enskog 展开式（Bardos 等，1991；Cercignani 等，1994；Chapman \& Cowling，1961；Ferziger \& Kaper，1972；Jenk-ins \& Mancini，1989）。基本思想是，碰撞项将驱动速度分布 n 趋向平衡函数 neq（即 C(neq) = 0 的解），因此与平衡的偏差可以近似为 n ≈neq + Knn1。从线性化玻尔兹曼方程中，可以求解 n1，从而闭合 neq 和 n1 定义中出现的守恒量 (αp, Up, θp) 的矩方程。因此，根据结构，当 Kn <1 时，流体动力学模型最准确。图 1.3 显示了玻尔兹曼方程的不同流态。两个相应的流体动力学模型的有效范围对于欧拉方程为 Kn = 0，对于纳维-斯托克斯 (NS) 方程为 Kn ≤0.1。在某些情况下，对于稍大的努森数，通过调整边界条件（即部分滑移），可以使用 NS 方程获得合理的预测。然而，对于较大的努森数，流动将不可避免地开始表现出非平衡行为11，仅使用五个矩方程无法捕获这种行为（Grad，1949b）。然而，请记住，玻尔兹曼方程（即介观模型）的精确解对于所有克努德森数都有效。因此，流体动力学模型的失败是用于闭合矩方程的数学近似的结果，而不是介观建模方法的固有属性。对于气体-粒子流，非平衡行为最明显的表现是 Knp = ∞ 时的粒子轨迹穿越 (PTC)（即无碰撞）。 PTC 发生在极稀流状态 (αp <1)，并且当颗粒温度为零 (Map = ∞) 时最容易识别。 PTC 的一个例子如图 1.4 所示。面板上的

\footnotetext[11]{平衡是指速度分布函数的形式。当 n = neq 时，流动处于平衡状态。当 n ≈neq + Knn1 时，流动接近平衡。否则，流动是非平衡的，必须通过求解介观尺度模型来找到 f。}
\begin{sourceblock}
\includegraphics[page=64,pagebox=mediabox,trim={58.00bp 525.80bp 58.23bp 59.00bp},clip,width=\textwidth,max height=0.72\textheight]{Marchisio2013.pdf}
\par\vspace{0.45\baselineskip}
\small 图1.3。以分散相克努森数表示的速度动力学方程的流体动力学模型的有效性范围。 NS，纳维-斯托克斯。经牛津大学出版社许可改编自 Bird (1994)。
\end{sourceblock}
\begin{sourceblock}
\includegraphics[page=64,pagebox=mediabox,trim={58.00bp 369.40bp 58.23bp 225.78bp},clip,width=\textwidth,max height=0.72\textheight]{Marchisio2013.pdf}
\par\vspace{0.45\baselineskip}
\small 图1.4。 Knp = ∞ 的高度非平衡流的体积分数场示例。左：流体动力学模型预测粒子云无法交叉，并且所有粒子最终的 y 动量为零。右：介观尺度模型预测粒子云将在没有碰撞的情况下交叉。改编自 Freret 等人。 （2008）。
\end{sourceblock}
左边显示了求解流体动力学模型的结果（尽管远远超出了其有效性范围！），右边的面板显示了介观模型的解（Kah 等人，2010）。流体动力学模型中 y = 0 处的质量累积在文献中称为 Delta 冲击（Chen \& Liu，2003）。从数学上讲，这是由于在流体动力学模型中，每个空间位置仅存在一个介观尺度速度值（因为颗粒温度以及速度波动为零）。当粒子云试图相互交叉时，局部速度被平均速度取代，平均速度为零。相比之下，介观尺度模型允许每个空间位置有多个速度，因此 PTC 可以不受数学近似的阻碍而发生。对于有限的 Knp，碰撞将导致速度为多值（但不一定是各向同性），并且粒子云将分散。只有在碰撞相对于传输速度较快的情况下，流体动力学模型才能提供气体-颗粒流的准确预测。原则上，可以通过扩展到更高阶（例如n ≈neq + Knn1 + Kn2n2等）来提高流体动力学模型的有效性范围。对于稀薄气体（例如玻尔兹曼方程），该过程取得了一些成功，但并没有极大地改善可以精确模拟的努森数范围（Struchtrup，2005）。对于具有多个无量纲参数（例如 Knp、St 等）的更复杂的介观模型，一致的高阶展开式能否提供比一阶流体动力学模型更好的预测并不明显。事实上，在许多多分散型多相流中，观察到的分散相努森数范围可以在同一流内以不稳定的方式变化很大（Passalacqua＆Fox，2011）。由于这些原因，需要开发一种对介观模型中出现的无量纲参数的所有可能值都有效的近似方法。目前，只有矩方法提供了这样的灵活性。

\subsection*{1.4.2 矩方法}\phantomsection\addcontentsline{toc}{subsection}{1.4.2 矩方法}
矩方法有许多不同的变体，但总体思路是增加传递矩的数量（超出流体动力学变量），以改进对非平衡行为的描述。如前所述，矩输运方程通常对于任何有限的矩集都不封闭。因此，任何矩方法的第一步都是对截断的矩方程组应用闭包过程。一般来说，这可以通过两种方式之一来完成。

1. 直接处理矩输运方程中的未封闭项，找到封闭它们的函数形式（Struchtrup，2005）。例如，涉及矩 Mk+1 的空间通量可以使用涉及矩 Mk 的梯度扩散模型来闭合。

2. 使用传递矩重构未知密度函数，然后计算矩输运方程中的未闭项。例如，在 Grad 的 13 矩方法（Grad，1949b）中，重建的密度函数使用高达三阶的矩，并通过积分找到四阶矩。

对于一般的多元介观尺度模型，直接闭合矩方程可能相当困难。例如，无法保证此类闭包会产生可实现的矩（即对应于非负 NDF 的矩）。由于可实现矩的空间具有相当复杂的边界，并且就其本质而言，高度非平衡问题的矩接近边界，因此向矩输运方程添加闭包很容易导致矩变得不可实现。因此，直接闭包矩方程通常仅适用于距离平衡不太远的问题。当用于远离平衡的问题时，此类方法的主要缺点不仅在于它们不准确，而且在于方程的数值属性可能会发生变化，从而导致求解器灾难性地失败。为了避免此类困难，有必要实现不一定与流动的基础物理一致的矩校正算法。因此，从物理学和数值的角度来看，非常希望仅使用保证产生可实现矩的矩闭包。后者最容易通过重建密度函数来实现。基于 NDF 重建的矩方法背后的基本思想如图 1.5 所示。从动力学方程开始，直接动力学求解器将使用 6D 计算网格上的时间步长来计算密度 n(t, x, v)，然后在相空间上积分以找到矩 Mk(t, x)。如前所述，通过首先在相空间上积分以找到未封闭的矩输运方程，可以避免直接求解器的高昂计算成本。原则上，可以使用 3D 有限体积网格上的时间步长来求解这些方程，以找到矩。然而，为了闭合矩输运方程，必须根据每个网格单元中的矩值在该单元中重建密度函数。重建的密度函数用 n*(t, x, v) 表示。一旦 n* 已知，矩输运方程为

\begin{sourceblock}
\includegraphics[page=66,pagebox=mediabox,trim={58.00bp 381.80bp 58.23bp 59.00bp},clip,width=\textwidth,max height=0.72\textheight]{Marchisio2013.pdf}
\par\vspace{0.45\baselineskip}
\small 图1.5。从动力学方程到使用基于 NDF 重建的矩方法的宏观模型。
\end{sourceblock}
在时间 t 关闭，可用于计算时间 t + δt 处的矩。然后重复重建过程以找到 n*(t + Δt, v, x)，并将矩更新到下一次。从这个描述中可以看出，重建过程是确保矩闭合一致性的关键步骤。在本书中，我们重点关注基于求积的矩方法（参见第 3 章）。然而，对与基于重构的矩方法相关的技术问题有一个总体概述将是有用的，如下所述。矩反演算法的四个必要属性如下。

（一）积极性。根据定义，NDF 是非负的。因此，重建的 NDF 对于其自变量的所有值都应该是非负的（即 n*≥0）。

(ii) 可实现性。介观变量可以在相空间的有限子集中呈现物理上可实现的值。例如，颗粒尺寸始终为正，颗粒表面积有下限。仅对于物理上可实现的介观尺度变量值，重构的 NDF 才应为非零。

(iii) 可扩展到任意顺序。为了提高矩收敛的精度，常常需要增加用于重建的矩集的阶数。因此，矩反演算法应该能够系统地增加矩的阶数，使得当矩的阶数趋向无穷大时，n*收敛到n。

{\small (iv) 非迭代。因为它要在每个时间步应用于每个网格单元，\par}
{\small 矩反演算法必须具有计算效率。在实践中，这\par}
意味着首选非迭代方法（即，如果可能，使用简单的分析公式）。然而，如果迭代是不可避免的（例如非线性求解器），则该方法应保证收敛。

实际上，在单个算法中获得所有四个属性可能很困难。例如，保证正性的一种方法是选择 n* 的函数形式，该函数形式取决于一小组未知参数。如果函数形式非常简单，则可以通过分析将矩与参数联系起来（例如高斯分布）。然而，很难保证可实现性并用简单的函数形式捕获复杂形状的 NDF，因此无法系统地将其扩展到高阶矩。使用其他方法，例如熵最大化（Massot et al., 2010；Mead \& Papanicolaou, 1984；Tagliani, 1999），可以实现积极性和可实现性，但代价是使用昂贵的迭代求解器来确定参数。另一种选择是使用涉及截断函数展开的 n* 形式（例如 Grad 方法），它可以系统地扩展到更高阶。该选项的优点是可以通过求解根据矩的定义构造的线性系统来找到系数（Favero \& Lage, 2012；Grad, 1949b；Lage, 2011；Strumendo \& Arastoopour, 2008）。然而，由于不能保证系数为非负，因此重建的 NDF 在相空间中的某些点很可能为负（即不可实现）。此外，当相空间的非零子集具有非正交边界时，可实现性条件也难以保证。总之，理想矩反演算法的必要属性很难获得。扩展到任意阶并保证为正的方法几乎总是需要迭代求解器（例如熵最大化方法）。这一规则的一个例外是基于求积的矩方法（Desjardins et al., 2008; Fox et al., 2008; Fox \& Vedula, 2010; Kah et al., 2010; Marchisio \& Fox, 2005; McGraw, 1997; Passalacqua et al., 2011; Vikas et al., 2011b; Yuan \& Fox, 2011; Zou et al., 2010) 其中参数由特征值问题确定（详细信息请参见第 3 章）。此外，尽管不能保证满足一般多元分布函数的可实现性，但基于求积的方法在与矩输运方程的有限体积求解器结合使用时具有多种数值优势（Vikas 等人，2011a）。从基于求积的矩方法导出的矩闭包的构造将在后面的章节中详细描述。

\section*{1.5 第 2--8 章的路线图}\phantomsection\addcontentsline{toc}{section}{1.5 第 2--8 章的路线图}
本书其余章节的结构如下。第二章简要介绍了多分散系统的介观描述。其中描述了根据内坐标的不同选择制定的数密度函数（NDF）的数学定义，然后介绍了各种形式的群体平衡方程（PBE）。第 2 章最后简短讨论了与 PBE 相关的矩输运方程与湍流理论中的系综平均所产生的矩输运方程之间的差异。这种差异非常重要，读者应该记住，在介观尺度水平上，微尺度湍流以流体阻力、质量传递等相关性的形式出现，因此即使微尺度流动是湍流（即单个粒子后面的湍流尾流），介观尺度模型也可以具有非湍流解。因此，在处理介观尺度流动的湍流模型时，必须将单独的系综平均程序应用于 PBE 的矩输运方程（或 PBE 本身）。在本书中，我们主要关注介观尺度模型的公式化和直接求解，因此我们不关注多相湍流模型。对多相湍流模型感兴趣的读者可以在文献中查找评论（Balachandar \& Eaton, 2010; Crowe et al., 1996; Fox, 2012; Jiang et al., 2010; Mashayek \& Pandya, 2003; Salazar \& Collins, 2009）。

第 3 章介绍了基于求积矩方法 (QBMM) 中使用的高斯求积和矩反演算法。为了制定后面章节中讨论的 PBE 的矩闭包，需要该材料的基础知识。由于内部变量的数量通常很大，因此空间非齐次问题的 PBE 直接解决方案通常会很棘手。因此，在本书中，我们重点关注 QBMM 以降低问题的维度。 QBMM 的关键步骤是从 NDF 的时刻开始重建 NDF。因此，在第 3 章中，我们回顾了单变量和多变量 NDF 的基于求积的矩反演。已经熟悉 QBMM 的读者可能仍然会发现本章很有用，因为它介绍了后面章节中使用的许多符号和术语。附录 A 中提供了示例矩反演算法。广义群体平衡方程 (GPBE) 的详细推导是第 4 章的主题。该推导从流体（即连续相）中相互作用的粒子（即分散相）的集合开始，并得出单粒子 GPBE。在此过程中，引入了多粒子联合分布函数级别的闭包，并对应于从微观尺度到介观尺度模型的转变（如1.3节中讨论的）。第 4 章显然是本书中理论性最强的一章，然而，在尝试制定正确捕获相互作用粒子系统的微观物理现象的介观模型之前，必须清楚地掌握多粒子闭合问题。一旦以封闭形式编写，GPBE 就可以提供分散型多相流的介观（或动力学）模型。此外，GPBE 是推导 NDF 矩输运方程的起点。因此，在第 4 章中，我们简要概述了通过相空间变量积分导出矩输运方程所需的数学步骤。正如 1.4 节中所讨论的，这个过程会导致第二个闭包问题，因为矩输运方程通常不是闭包的，而第二个闭包才是 QBMM 的焦点。第五章重点介绍从文献中选择的关键物理和化学过程的介观模型。正如读者可能想象的那样，文献中提出的介观尺度模型的数量是巨大的。因此，第五章的目标是提供一般现象的模型示例，例如物理空间和相空间中的平流、物理空间和相空间中的扩散，以及分散相成核等点过程。在实际应用中，所采用的介观模型的准确性对于模拟的整体准确性至关重要。然而，由于本书的重点主要是如何在通用环境中制定和求解矩输运方程，因此我们的目的并不是对可用的介观尺度模型进行详尽的分类。尽管如此，读者应该会发现第 5 章是对介观尺度模型的有用介绍，也是开发新模型的良好起点。在后续章节中，我们使用第 5 章中的介观尺度模型来说明求解 GPBE 的数值方法。许多分散相系统涉及粒子之间的碰撞，典型的例子是硬球碰撞。 因此，第 6 章专门讨论 QBMM 背景下的硬球碰撞模型主题。特别是，由于具有硬球碰撞的 GBPE 的矩输运方程包含碰撞期间 NDF 变化率的源项，因此有必要推导这些源项的解析表达式 (Fox \& Vedula, 2010)。在第 6 章中，导出了两个具有不同直径/质量的粒子之间的非弹性二元碰撞的整数矩的确切源项。尽管推导的技术细节相当复杂，但产生的源项具有相当简单的形式，在第 6 章中列出以供将来使用。讨论了更简单的碰撞动力学模型，并提供了选定的矩输运方程作为示例。最后，本章最后简要讨论了如何使用 QBMM 和导出的解析源项来闭合矩输运方程中的碰撞通量和碰撞源项。

在前面的章节中介绍了基本概念和模型之后，第 7 章和第 8 章讨论了模拟多分散颗粒和多相系统所需的关键计算方法。正如在模拟代码中经常使用算子分裂所做的那样，我们首先在第 7 章中考虑空间齐次系统的数值方法，然后在第 8 章中讨论空间非齐次系统的处理。在齐次系统的背景下，我们回顾通过离散相空间变量直接求解 GPBE 的求解方法，以及使用矩闭包求解 NDF 矩的间接方法。作为同质系统的标准实践，还提出了通过比较直接和间接解决方案来验证矩闭合的方法。第 8 章专门讨论求解空间非齐次 GPBE 的矩方法。空间非均匀系统的一个关键问题是数值平流（扩散）算子的矩可实现性，因此我们引入基于动理学的有限体积方法（KBFVM）（Pertame，2003）来解决这个问题。附录 B 中提供了应用于矩输运方程的 KBFVM 的更多详细信息。其他关键问题是相空间中存在粒子轨迹穿越 (PTC) 时的矩闭合（这种情况发生在远离碰撞平衡的系统中），以及某些 QBMM 的矩守恒。这两个问题分别在附录 C 和附录 D 中讨论。如附录 C 中的分析示例所示，PTC 对应于双曲 GPBE 中特征的交叉。因此，在第 8 章中，我们证明 QBMM 非常适合在矩输运方程级别捕获 PTC。作为本书的总结，在第 8 章中，使用 QBMM 结合用于空间传输的 KBFVM（Vikas 等人，2011a）解决了具有不同类型 GPBE 的空间一维数值示例。我们提供了数值实现中使用的确切公式，因为作者希望读者尝试重现其中一些示例，从而获得模拟多分散多相系统的宝贵经验，这些经验可应用于读者自己的应用。

\chapter{多分散系统的介观描述}
在本章中，提出了描述多分散型多相流所需的控制方程，但没有从微观模型进行严格推导。 （完整推导请参见第 4 章。）为了清楚起见，本章中控制方程的讨论将仅限于颗粒系统（例如结晶器、流化床和气溶胶过程）。然而，熟悉分散型多相流建模的读者将重新认识到我们的评论适用于更普遍的背景。事实上，将本章中开发的建模概念扩展到许多其他多相系统是简单的，并将在后面的章节中讨论。本章的主要目的是介绍开发多分散型多相流模型所需的关键概念和符号。因此，我们首先对各种形式的数密度函数 (NDF) 进行一般讨论，然后是具有已知 (PBE) 和计算 (GPBE) 粒子速度的 NDF 的传输方程示例。这些输运方程是用“平均”量写的，其精确定义将在第 4 章中介绍。然后，我们考虑通过相空间积分从 NDF 输运方程导出的矩输运方程。最后，我们简要描述如何从矩输运方程开始进行湍流建模。

\section*{2.1 数密度函数（NDF）}\phantomsection\addcontentsline{toc}{section}{2.1 数密度函数（NDF）}
分散相由离散元素构成。我们分析的主要假设之一是元素的特征长度尺度小于感兴趣的属性变化的特征长度尺度（即化学物质浓度、温度、连续相速度）。如果这个假设成立，则颗粒系统可以用连续介质或平均场理论来描述。分散相的每个元素通常由许多称为坐标的属性来识别。如果两个元素具有相同的坐标值，则它们是相同的，否则元素是无法区分的。通常坐标分为内部坐标和外部坐标。外部坐标是空间坐标；事实上，元素在物理空间中的位置并不是元素的内部属性。内坐标指的是元素的更亲密的属性，例如它们的动量（或速度）、它们的焓
（或温度），或其体积、表面积、大小或年龄。当使用欧拉方法时，分散相通常由 NDF 描述。正如第 1 章中提到的，NDF 包含有关无穷小控制体积内的粒子群如何分布在感兴趣的属性上的信息。让我们考虑位于物理点 x eq(x1, x2, x3) 且测量值为 dx eq dx1 dx2 dx3 的无穷小控制体积内的一群分散实体，例如固体颗粒或液滴。令 ϋ ≡(ϋ1, ϋ2, . . . , ϋM) 为内坐标向量，即包含感兴趣的 M 个内坐标的向量。 NDF nΣ(t, x, Σ) 定义为无穷小物理体积 dx 和无穷小相空间体积 dΣ 中实体的预期数量 1：nΣ dx dΣ。请注意，nψ 是时间 t、空间 x 和内坐标向量 ψ 的函数。今后，只有在需要清楚说明时，我们才会包含对这些变量的显式依赖。如第 4 章详细所示，NDF 是分散相的平均数量。事实上，它具有平均函数的数学特征：它相对于时间、物理空间和内坐标（或相位）空间是平滑且可微的。很明显，NDF 的基本定义是基于微尺度系统的随机描述。事实上，尽管总体中任何单个粒子的行为都被视为随机的，但如果总体中粒子的数量很大，则粒子总体的行为可以被视为确定性的。正如第 4 章中所讨论的，NDF 识别计算域中任何时刻和任何给定点的整个粒子群，并考虑与每个粒子的状态相关的概率。我们建议读者参考专门的文本来推测这个平均值的性质，以及对非常稀的系统的处理，在这些系统中，由于总体中的粒子数量太少，确定性描述失败；例如，参见 Ramkrishna (2000) 第 7 章。目前，我们只是在这里强调，这种随机行为可能反映了分散相非常亲密的性质的一些混沌特性，或者反映了例如连续流体相中的微尺度湍流引起的一些其他混沌行为。然而，正如第 2.4 节中将变得更加清楚的那样，在处理 NDF 输运方程解中的介观尺度湍流时，NDF 可以像湍流中的任何其他标量场（尽管是无限维）一样处理。很明显，数量 n xi ( xi ) d xi 代表每单位物理体积以 ψ 为中心的相空间体积 d ψ 中包含的分散实体的数量密度。如果我们对内坐标向量 Ωψ 所有可能值的 NDF 进行积分，我们就可以得到总数浓度 N(t, x)：

\begin{sourceblock}
\includegraphics[page=71,pagebox=mediabox,trim={195.79bp 198.26bp 67.56bp 467.56bp},clip,max width=0.94\textwidth,max height=0.38\textheight]{Marchisio2013.pdf}
\end{sourceblock}
这是位于时间 t 和点 x 的每单位体积的实体总数，对应于密度函数的零阶矩。同样，NDF 的任意矩可以定义为

\begin{sourceblock}
\includegraphics[page=71,pagebox=mediabox,trim={186.94bp 120.50bp 67.56bp 545.31bp},clip,max width=0.94\textwidth,max height=0.38\textheight]{Marchisio2013.pdf}
\end{sourceblock}
1 我们将在第 7 章中看到，随机算法可用于模拟由 NDF 表示的粒子群的演化。然而，NDF 并不是一个随机量；相反，它是随机算法无限数量实现的整体平均值（假设后者是无偏的）。

\footnotetext[2]{当 NDF 表示粒子速度时，输运方程是双曲的，因此解不需要是光滑的，也不需要在空间中的每个点都是可微的。事实上，NDF 中的休克治疗是一个重大挑战。}
其中 k ≡(k1, ..., kM) 是一个向量，其中包含与 xi 的每个分量相关的矩的阶数。

\subsection*{2.1.1 基于长度的NDF}\phantomsection\addcontentsline{toc}{subsection}{2.1.1 基于长度的NDF}
为了引入一些对讨论的发展非常有用的特定概念和属性，我们考虑仅使用一个内坐标即可描述的粒子群。正如已经提到的，除非通过所选的内坐标，否则粒子是无法区分的。让我们想象一下，粒子可以通过它们的大小或长度 L 来表征，在这种情况下，它构成了内坐标向量的唯一分量。使用这种表示法，nL(t, x, L) 是基于长度的 NDF，nL dL 表示长度在 L 和 L+dL 之间的粒子的预期数密度。正如上面方程式中所报告的那样。 (2.1) 总数量浓度可计算为

\begin{sourceblock}
\includegraphics[page=72,pagebox=mediabox,trim={195.68bp 446.58bp 65.07bp 219.62bp},clip,max width=0.94\textwidth,max height=0.38\textheight]{Marchisio2013.pdf}
\end{sourceblock}
而平均粒径可以定义为

\begin{sourceblock}
\includegraphics[page=72,pagebox=mediabox,trim={182.91bp 396.82bp 65.07bp 270.73bp},clip,max width=0.94\textwidth,max height=0.38\textheight]{Marchisio2013.pdf}
\end{sourceblock}
分布的平均特性是通过数密度本身定义的，它代表相对于系统中颗粒数量的平均颗粒尺寸。当然，平均粒度的其他定义也是可能的，这将在下面变得清楚。如果我们将基于长度的 NDF 的 kth 矩定义为

\begin{sourceblock}
\includegraphics[page=72,pagebox=mediabox,trim={200.47bp 304.80bp 65.07bp 361.41bp},clip,max width=0.94\textwidth,max height=0.38\textheight]{Marchisio2013.pdf}
\end{sourceblock}
很容易看出，对于任意 k 值，平均粒径可定义为 mL,k+1/mL,k 之比。例如，Sauter 平均直径通过设置 k = 2 来定义，给出 L32 ≠ mL,3/mL,2，并且对于对表面积是关键参数的质量传递等过程进行建模非常有用。如果体积与长度的三次方成比例，我们可以通过体积形状因子 kV:3 V = kVL3 将颗粒体积 V 和颗粒长度 L 联系起来。生成的体积密度函数 VL(t, x, L) 表示长度在 L 和 L + dL 之间的每单位空间体积的颗粒体积，并由 VL eqkVL3nL 定义。将 VL 标准化为 1 后，体积分数密度函数 αV(t, x, L) 表示具有特定长度的颗粒占总颗粒体积的体积分数：

\begin{sourceblock}
\includegraphics[page=72,pagebox=mediabox,trim={204.01bp 148.74bp 65.07bp 518.75bp},clip,max width=0.94\textwidth,max height=0.38\textheight]{Marchisio2013.pdf}
\end{sourceblock}
根据该体积分数密度函数计算出的平均颗粒长度为

\begin{sourceblock}
\includegraphics[page=72,pagebox=mediabox,trim={187.00bp 94.52bp 65.07bp 571.68bp},clip,max width=0.94\textwidth,max height=0.38\textheight]{Marchisio2013.pdf}
\end{sourceblock}
\footnotetext[3]{对于分形颗粒，该体积不能用于计算颗粒的质量。相反，质量通过分形维数与长度相关。}
颗粒长度的这种定义广泛用于颗粒系统。如上所述，L10是相对于颗粒数密度平均的平均颗粒长度，而L43是相对于颗粒体积分数密度函数平均的平均颗粒长度。当然，所有这些粒子长度的定义都可以从基于长度的 NDF 开始找到，因此没有必要为每种新型 NDF 开发单独的理论。

\subsection*{2.1.2 基于体积的NDF}\phantomsection\addcontentsline{toc}{subsection}{2.1.2 基于体积的NDF}
颗粒系统的另一个非常常见的内部坐标是颗粒体积 V。使用与 NDF 非常相似的定义，我们可以将基于体积的 NDF 定义为体积介于 V 和 V + dV 之间的颗粒的预期数量：nV dx dV。如果体积和长度的三次方为 VL eqkVL3nL，则基于长度的 NDF 和基于体积的 NDF 之间的关系很简单。事实上，nL dL 量化了长度在 L 和 L+dL 之间的粒子的预期数量，它必须等于体积在 V 和 V + dV 之间的粒子的预期数量。而且，容易看出dV = 3kVL2 dL。因此，期望数相等需要 nL(L)dL = nV(V)dV 或 nL(L)dL = nV(kVL3)3kVL2 dL，由此我们发现 nL(L) = 3kVL2nV(kVL3)，这本质上告诉我们如何用 nV 来定义 nL。可以使用相同的过程根据任何明确定义的 L 函数求 NDF。基于体积的数 NDF 的矩定义为

\begin{sourceblock}
\includegraphics[page=73,pagebox=mediabox,trim={196.48bp 357.95bp 67.56bp 308.25bp},clip,max width=0.94\textwidth,max height=0.38\textheight]{Marchisio2013.pdf}
\end{sourceblock}
并且可以很容易地与基于长度的 NDF 的矩相关联：

\begin{sourceblock}
\includegraphics[page=73,pagebox=mediabox,trim={167.03bp 309.86bp 67.56bp 362.29bp},clip,max width=0.94\textwidth,max height=0.38\textheight]{Marchisio2013.pdf}
\end{sourceblock}
在实际应用中，NDF 的选择通常取决于可以通过实验测量哪些矩。

\subsection*{2.1.3 基于质量的 NDF}\phantomsection\addcontentsline{toc}{subsection}{2.1.3 基于质量的 NDF}
在有关颗粒聚集和液滴聚结的文献中，通常首选基于质量的 NDF nM(t, x, M)，因为质量是守恒量。对于这些系统，可以使用 M = kMLDf 将粒子的质量与其特征长度联系起来，其中 1 < Df ≤ 3 是分形维数。4 使用上述过程，我们可以证明基于长度和基于质量的 NDF 之间的关系为 nL(L) = DfkMLDf−1nM(kMLDf)。同样，它们的矩通过 mM,k = kk MmL,Dfk 相关。正如我们将在接下来的章节中看到的，通常情况下，在封闭系统中，总粒子质量是守恒的，即 mM,1 是恒定的。在基于长度的分析中，这对应于分数矩 mL,Df 守恒。理想情况下，无论用于表达问题的 NDF 类型如何，求解 PBE 的数值算法都应保持质量守恒。然而，由于模型中使用的是数值近似，情况并非总是如此。我们将回到第 5 章中选择合适的 NDF 的主题。

\footnotetext[4]{对于球形颗粒，Df = 3；对于扩散限制聚集体，Df ≈1.8。}
\subsection*{2.1.4 基于速度的NDF}\phantomsection\addcontentsline{toc}{subsection}{2.1.4 基于速度的NDF}
当内坐标向量是粒子速度向量时，会出现一个令人感兴趣的特殊情况，我们将其用相空间变量 v.5 表示。事实上，粒子速度是一个特殊的内坐标，因为它通过牛顿定律与粒子位置（即外坐标）相关，因此需要进行特殊处理。我们稍后会回到这个方面，但暂时让我们想象一下其他方面相同的6个粒子正在以不同的速度移动（并且与周围的流体速度不同）。因此，可以定义基于速度的 NDF nU(t, x, v)，其由速度分量 v = (v1, v2, v3) 参数化。为了获得总数浓度（即每单位体积的颗粒数），对所有可能的颗粒速度 Ωv 值进行积分就足够了：

\begin{sourceblock}
\includegraphics[page=74,pagebox=mediabox,trim={196.86bp 464.24bp 65.07bp 201.43bp},clip,max width=0.94\textwidth,max height=0.38\textheight]{Marchisio2013.pdf}
\end{sourceblock}
这只是基于速度的 NDF 的零阶矩。因为我们假设粒子可能具有不同的速度，所以我们可以使用基于速度的 NDF 的前两个矩计算每个空间位置和每个时刻的平均粒子速度 Up：

\begin{sourceblock}
\includegraphics[page=74,pagebox=mediabox,trim={185.30bp 382.30bp 65.07bp 283.38bp},clip,max width=0.94\textwidth,max height=0.38\textheight]{Marchisio2013.pdf}
\end{sourceblock}
这类似于方程式中报告的平均颗粒长度的定义。 （2.4）。读者应注意，平均粒子速度的定义仅在所有粒子具有相同质量时才适用。如果粒子质量是内坐标，则平均粒子速度使用平均分散相动量与平均粒子质量的比率来定义。显然，NDF nψ 和基于速度的 NDF nU 具有相同的数学意义，并且可以在类似的框架中处理。然而，正如已经强调的那样，粒子速度是一个特定的内坐标，因此这两类问题的理论和计算模型通常被分开处理。接下来，我们将推导这些 NDF 的输运方程，将被动内坐标（例如颗粒尺寸、体积、面积或温度）的处理与主动内坐标（例如粒子速度）的处理分开。我们将第一个输运方程称为群体平衡方程 (PBE)7，将第二个方程称为广义群体平衡方程 (GPBE)。8 在接下来的章节中，我们将开发一个通用框架来以数值方式描述和解决这两个问题。然而，读者应该记住，这两个模型之间的关键区别在于，对于 PBE，分散相的速度被假定为已知（例如，与连续相中的流体速度相同），因此 NDF 可以被视为被动标量场，而 GPBE 中出现的标量场通常是主动标量场（即动量）

5 对于多相流，有必要区分各相的速度。例如，对于气固流，我们将用 Uf 表示流体速度矢量。然而，对于当前讨论的大部分内容，我们可以假设粒子在真空中传播。 6 稍后我们将通过引入属性（长度、体积、质量等）也不同的粒子来放宽这一要求。

\footnotetext[7]{PBE 在气溶胶文献中也称为一般动力学方程。 8在文献中，根据研究领域的不同，GPBE 有许多不同的名称。例如，在统计物理学中，当写出速度 NDF 时，它被称为玻尔兹曼动力学方程；在燃烧文献中，当为液滴编写时，它被称为喷雾或威廉姆斯方程。}
粒子之间的交换取决于内部坐标（例如长度），因此分散相的速度必须根据一致的动量平衡来计算。

\section*{2.2 NDF输运方程}\phantomsection\addcontentsline{toc}{section}{2.2 NDF输运方程}
我们现在考虑 NDF9 在时间、相空间和物理空间上的演化。我们将看到 PBE 和 GPBE 的底层数学结构非常相似。

\subsection*{2.2.1 群体平衡方程（PBE）}\phantomsection\addcontentsline{toc}{subsection}{2.2.1 群体平衡方程（PBE）}
PBE 是根据 NDF 编写的简单连续性声明。它可以作为相和物理空间的某些固定子区域中粒子的平衡而导出（Ramkrishna，2000）。让我们考虑物理空间 Ωx 和相空间 Ωψ 中的有限控制体积，其边界分别定义为 ∂Ωx 和 ∂Ωψ。在 PBE 中，假设平流速度 v 是已知的（例如等于连续相中的局部流体速度或直接从该变量导出）。粒子数平衡方程可以写为

\begin{sourceblock}
\includegraphics[page=75,pagebox=mediabox,trim={98.99bp 349.72bp 67.56bp 286.51bp},clip,max width=0.94\textwidth,max height=0.38\textheight]{Marchisio2013.pdf}
\end{sourceblock}
其中 v 是颗粒系统的速度矢量， ̇ 是相空间的连续变化率， h 是表示离散事件的不连续跳跃函数。除了平流速度之外，细颗粒（尺寸小于约一微米）和具有湍流尾流的低惯性颗粒（例如气泡）将分别由于布朗波动和湍流波动而具有扩散通量项。方程中的第一项。 (2.12) 是控制量上的积分，表示累加。第二项是控制体积10边界上的表面积分，表示物理空间中对流引起的数密度净通量。第三项代表相空间中的对流。事实上，粒子在物理空间中移动是因为粒子速度 v，而在相空间中移动则是因为内坐标速度 ̇ xi = d xi/dt。该向量在每个分量中包含连续过程的第 i 个内坐标的变化率。我们通常指的是内坐标以比 PBE 特征解小得多的时间尺度连续变化的过程。因此，尽管所有过程（至少在分子水平上）都是不连续的，但在方程（1）中这些过程可以被视为连续的。 （2.12）。例如，如果由于分子或分子团连续沉积在颗粒表面上而使颗粒尺寸增大，并且这种情况发生的时间和长度尺度远小于与颗粒尺寸相关的时间和长度尺度，则这些过程可以视为连续的，并且它们的介观尺度描述11是第三项
9 同样的分析一般适用于密度函数，包括概率密度函数。 10 控制体积边界上的积分通过无穷小表面单位向量 dAx 定义，其大小等于无穷小表面的测量值，方向由垂直于表面本身的向量定义。

\footnotetext[11]{即使在微观层面，分子尺度事件也是通过连续过程来描述的。}
方程的（2.12）。方程中的项。方程（2.12）指的是物理或相空间中的对流，通常称为漂移项。另一方面，如果颗粒因碰撞和随后的聚集而生长，由于该过程的长度尺度与颗粒尺寸处于同一数量级，因此不能将其视为相空间的连续变化。该过程属于相空间中不连续跳跃的类别，并由等式右侧的术语表示。 （2.12）。这些概念将在第五章中变得更加清晰，我们将在其中研究连续和离散过程的典型模型。如果将雷诺-高斯定理应用于方程 12。 (2.12) 并且控制体积边界上的积分可以用控制体积本身的积分来写，很容易得到

\begin{sourceblock}
\includegraphics[page=76,pagebox=mediabox,trim={98.74bp 452.11bp 65.07bp 184.12bp},clip,max width=0.94\textwidth,max height=0.38\textheight]{Marchisio2013.pdf}
\end{sourceblock}
其中 ∂/∂x = ∇x = (∂/∂x1, ∂/∂x2, ∂/∂x3) 是物理空间中的梯度算子，∂/∂Σ = ∇Σ = (∂/∂Σ1, . . . , ∂/∂Σm) 是相空间中的梯度算子。为了方程。对于任何任意控制量 Ωx 和 Ωψ，要满足 (2.13) 式，被积函数必须满足以下关系

\begin{sourceblock}
\includegraphics[page=76,pagebox=mediabox,trim={178.05bp 367.37bp 65.07bp 304.48bp},clip,max width=0.94\textwidth,max height=0.38\textheight]{Marchisio2013.pdf}
\end{sourceblock}
或者，使用爱因斯坦符号（即重复的罗马索引意味着求和），

\begin{sourceblock}
\includegraphics[page=76,pagebox=mediabox,trim={178.18bp 311.89bp 65.07bp 359.14bp},clip,max width=0.94\textwidth,max height=0.38\textheight]{Marchisio2013.pdf}
\end{sourceblock}
请注意，虽然我们以相同的方式对待 x 和 xi，但它们实际上是不同类型的向量。矢量 x 和 v 是连续介质力学中使用的位置和速度的标准矢量。另一方面，内坐标向量 ψ 是线性代数意义上的长度为 N 的广义向量。方程（2.14）必须与为启动时间给出的初始条件以及物理空间 Ωx 和相空间 ΩΨ 中的边界条件相结合。方程的解析解。 （2.14）适用于一些特殊情况，并且仅在一些非常简单的假设指定的条件下可用。然而，可以使用数值方法来求解该方程，并将在第 7 章和第 8 章中介绍。 (2.14) 提供了每个时刻、计算域中每个物理点以及相空间中每个点的 NDF 知识。正如已经提到的，有时粒子群仅由一个内部坐标来描述，例如粒子长度（即 xi = L），并且 PBE 被称为单变量。当需要两个内坐标时，例如颗粒体积和表面积（即 xi = (v, a)），PBE 被称为双变量。更一般地，高维情况被称为多元 PBE。当内坐标向量的一部分等于粒子速度向量时（即，当粒子的特征不是由唯一的速度场而是由它们自己的速度分布来表征时），会发生另一种重要的情况。在这种情况下，PBE 变为 GPBE，如下所述。

\footnotetext[12]{有关雷诺-高斯输运定理的严格推导，请参见 Aris (1962)。}
\subsection*{2.2.2 广义群体平衡方程（GPBE）}\phantomsection\addcontentsline{toc}{subsection}{2.2.2 广义群体平衡方程（GPBE）}
通过使用与上述 PBE 非常相似的方法，可以导出 NDF 的 GPBE，其中包括粒子速度作为内部变量。我们将这个一般 NDF 表示为 n(t, x, v, xi)（即 n 上没有下标）。最简单的 GPBE（即没有其他内坐标的速度）被称为玻尔兹曼动力学方程，最初是在气体理论的背景下推导出来的（Chapman \& Cowling，1961）。 GPBE的最终形式是

\begin{sourceblock}
\includegraphics[page=77,pagebox=mediabox,trim={155.89bp 516.04bp 67.56bp 155.80bp},clip,max width=0.94\textwidth,max height=0.38\textheight]{Marchisio2013.pdf}
\end{sourceblock}
其中 Ap 是粒子速度（即加速度）的连续变化率，或作用在粒子上的每单位质量的力（例如重力、流体阻力等）。右侧 h 又是不连续跳跃项，但现在包括粒子动量的不连续变化（例如碰撞）。在 GPBE 中，物理空间中从来不存在扩散项。然而，由于随机（例如布朗）力或连续相中的湍流波动（例如惯性粒子后面的湍流尾迹），速度相空间中经常存在扩散术语。等式左边的第一项。 （2.16）是累积项，而第二项是漂移项，它解释了由于粒子速度而导致的物理空间中的对流。第三项是代表粒子速度相空间中对流的漂移项。事实上，作用在粒子上的力会产生均匀的加速度，或者换句话说，粒子速度的连续变化。右侧的术语表示由于离散事件（主要是粒子之间的碰撞）导致的相空间中的不连续跳跃，因此也称为碰撞积分。 13 当两个粒子碰撞时，它们交换动量，因此在同一时刻它们将动量（以及速度）从一组特定值更改为另一组特定值。在这方面，碰撞过程是离散的，并在粒子速度相空间中产生有限且不连续的变化。 14此外，如果碰撞导致内坐标的变化（例如气泡合并、粒子聚集等），则与气体动力学中使用的硬球碰撞相比，碰撞项可能相当复杂。

\subsection*{2.2.3 闭包问题}\phantomsection\addcontentsline{toc}{subsection}{2.2.3 闭包问题}
现在我们已经介绍了 GPBE，现在是时候讨论闭包问题了。一般来说，我们说等式。如果方程中的所有项都可以根据 n(t, x, v, xi) 及其矩的知识来计算，则 (2.16) 是闭合的。 15 在本书中，我们实际上将假设感兴趣的 GPBE 是闭合的，并重点关注计算 n(t, x, v, xi) 的方法。正如第 4 章所讨论的，寻找多粒子相关性闭包的问题，例如，由于碰撞期间的粒子动量交换或通过
13 正如统计物理学中众所周知的那样，碰撞积分不是封闭的，因为它涉及两个粒子之间的速度相关性。因此，必须引入封闭假设来根据 n 定义 h。详细信息请参见第 6 章。 14 仅当碰撞发生的时间和长度尺度比 NDF 中表征变化的时间和长度尺度短得多时，才可以将碰撞描述为不连续。这确实是在硬球势下发生粒子与粒子碰撞的情况，但在不同的势下，例如，如果粒子开始在很远的距离处感知到彼此的存在，情况可能不再是这样。

\footnotetext[15]{由于后者是根据 n 计算的，因此提及矩是多余的。尽管如此，我们无论如何还是要提到它们，以提醒读者，GPBE 级别的闭包包括矩级别的特殊情况闭包。}
由连续流体相介导的相互作用是非常重要的。例如，为了闭合玻尔兹曼动力学方程，人们通常求助于分子混沌假设，从而通过两个单粒子 NDF 的乘积来近似两粒子 NDF。 （详细信息请参见第 6 章。）正如第 1 章中提到的，本书的重点不是可用于关闭 GPBE 的方法。相反，我们将建议读者参考专门的技术文献，了解为封闭 GPBE 而导出的封闭方法。最后，我们想提醒读者，即使 GPBE 是封闭的（即适合本书中讨论的计算方法），这并不意味着矩输运方程将是封闭的。事实上，在大多数实际问题中情况并非如此。尽管如此，了解（未封闭的）矩输运方程的形式通常很有用。在下一节中，我们将简要概述如何从 PBE 和 GPBE 出发找到矩输运方程。

\section*{2.3 矩输运方程}\phantomsection\addcontentsline{toc}{section}{2.3 矩输运方程}
前面已经提到，PBE（方程（2.14））和GPBE（方程（2.16））具有非常相似的结构。它们都量化 NDF 的变化率，代表经历一些连续过程（例如分子生长或由于物理力而加速）和一些离散过程（例如碰撞和聚集）的离散元素群体。现在有趣的是从 NDF 输运方程导出一些感兴趣的积分量（即矩），并定义它们的输运方程。

\subsection*{2.3.1 PBE 的矩输运方程}\phantomsection\addcontentsline{toc}{subsection}{2.3.1 PBE 的矩输运方程}
首先，为了简单起见，我们考虑一个仅以粒子长度作为内坐标的 PBE，其 NDF 为 nL(t, x, L)。 PBE 在方程式中报告。 (2.14) 则变为

\begin{sourceblock}
\includegraphics[page=78,pagebox=mediabox,trim={174.58bp 293.45bp 65.07bp 380.80bp},clip,max width=0.94\textwidth,max height=0.38\textheight]{Marchisio2013.pdf}
\end{sourceblock}
其中粒子长度连续变化率 ̇L 用 GL 表示，离散事件用 hL 表示。与所有 PBE 一样，我们假设速度 v 已知。还请回想一下，对于细颗粒，由于周围流体相中颗粒的布朗运动，物理空间中的额外扩散通量会出现在右侧。 16 然而，由于扩散系数与颗粒质量成反比，因此在许多系统中，扩散通量贡献将非常小，甚至可以忽略不计。将矩变换（第 32 页的方程（2.5））应用于方程。 (2.17) 我们得到17

\begin{sourceblock}
\includegraphics[page=78,pagebox=mediabox,trim={159.35bp 178.23bp 65.07bp 495.88bp},clip,max width=0.94\textwidth,max height=0.38\textheight]{Marchisio2013.pdf}
\end{sourceblock}
其中平均粒子生长率定义为

\begin{sourceblock}
\includegraphics[page=78,pagebox=mediabox,trim={194.35bp 126.19bp 65.07bp 531.50bp},clip,max width=0.94\textwidth,max height=0.38\textheight]{Marchisio2013.pdf}
\end{sourceblock}
16 在推导 PBE 时，隐含地假设颗粒之间的平均距离足够大，使得颗粒运动由流体相驱动。如果情况并非如此，那么描述粒子的正确起点将是 GPBE。因此，对于 PBE，扩散通量项可以通过连续流体相中大小为 L 的颗粒的稀释二元扩散系数来近似。

\footnotetext[17]{如果 L = 0 时 GL(L) 为负，则该方程中可能需要额外的汇项，以考虑由于原点消失而导致的粒子损失。例如，当蒸发速率与表面积成正比时，这种情况会发生在由蒸发液滴组成的系统中。}
hL 的矩变换为

\begin{sourceblock}
\includegraphics[page=79,pagebox=mediabox,trim={204.49bp 587.50bp 67.56bp 78.70bp},clip,max width=0.94\textwidth,max height=0.38\textheight]{Marchisio2013.pdf}
\end{sourceblock}
在下面的方程中，为简单起见，我们假设粒子增长率 GL 与尺寸无关，因此 GL,k = GL。然而，一般来说，建模者必须提供 GL 对尺寸的函数依赖性，以便应用方程（1）。 （2.19）。零阶矩 mL,0 相当于总数浓度 N，其传输方程由式（1）得出。 （2.18）：

\begin{sourceblock}
\includegraphics[page=79,pagebox=mediabox,trim={198.75bp 479.40bp 67.56bp 194.95bp},clip,max width=0.94\textwidth,max height=0.38\textheight]{Marchisio2013.pdf}
\end{sourceblock}
可以看出，由于颗粒尺寸的连续增长不会改变总数浓度（如果 GL > 0），因此漂移项消失了。然而，N 受到颗粒形成速率（例如成核）以及聚集和破碎速率的影响，从而导致颗粒的出现和消失。这些过程都包含在源项 hL,0 中。三阶矩 mL,3 与颗粒相对于悬浮流体所占的体积分数有关，可以从式（1）中轻松找到。 （2.18）：

\begin{sourceblock}
\includegraphics[page=79,pagebox=mediabox,trim={176.88bp 354.16bp 67.56bp 312.04bp},clip,max width=0.94\textwidth,max height=0.38\textheight]{Marchisio2013.pdf}
\end{sourceblock}
其中 αp 表示颗粒相的体积分数。流体相用 αf 表示，αp +αf = 1。由此得出的分散相体积分数的传输方程为

\begin{sourceblock}
\includegraphics[page=79,pagebox=mediabox,trim={163.73bp 279.02bp 67.56bp 392.83bp},clip,max width=0.94\textwidth,max height=0.38\textheight]{Marchisio2013.pdf}
\end{sourceblock}
有趣的是，如果漂移项为空（即粒子尺寸没有增长）并且没有向系统中引入新粒子（即成核速率为空18），则方程的右侧为： (2.23) 为零，导致粒子相的连续性方程为

\begin{sourceblock}
\includegraphics[page=79,pagebox=mediabox,trim={201.12bp 192.48bp 67.56bp 479.38bp},clip,max width=0.94\textwidth,max height=0.38\textheight]{Marchisio2013.pdf}
\end{sourceblock}
由于 v 已知，因此分散相体积分数的闭合双曲平流方程是保守形式。从 PBE 开始寻找矩输运方程的过程可以继续到任意阶。我们应该注意到，在大多数应用中，所得的矩输运方程不会是封闭的。换句话说，k 阶矩输运方程将涉及高于 k 阶的矩。我们将在第 7 章和第 8 章中在基于求积的矩方法的背景下讨论矩闭合方法。

\footnotetext[18]{当成核产生尺寸为零的核时，会出现 hL,3 为零但成核速率非零的非常常见的情况。对于这种情况，事实上，只有 hL,0 是非零的。}
\subsection*{2.3.2 GPBE 的矩输运方程}\phantomsection\addcontentsline{toc}{subsection}{2.3.2 GPBE 的矩输运方程}
与上述 PBE 相同的方法也可以应用于 GPBE（第 37 页的方程（2.16））。将矩变换应用于方程。 (2.16) 对于零阶矩，得到总数浓度的输运方程：

\begin{sourceblock}
\includegraphics[page=80,pagebox=mediabox,trim={200.92bp 548.97bp 65.07bp 122.88bp},clip,max width=0.94\textwidth,max height=0.38\textheight]{Marchisio2013.pdf}
\end{sourceblock}
其中碰撞项的矩变换为

\begin{sourceblock}
\includegraphics[page=80,pagebox=mediabox,trim={195.98bp 495.31bp 65.07bp 170.37bp},clip,max width=0.94\textwidth,max height=0.38\textheight]{Marchisio2013.pdf}
\end{sourceblock}
（参见第 34 页的方程（2.11））Up 是平均粒子速度。该传输方程与方程非常相似。 （2.21）。我们可以注意到，对于颗粒尺寸恒定（h0 = 0）的情况，关系式 αp = kVL3N 将允许我们重写方程： （2.25）作为分散相体积分数的平衡（或者，通过乘以颗粒密度ρp，颗粒质量分数）：

\begin{sourceblock}
\includegraphics[page=80,pagebox=mediabox,trim={191.47bp 400.16bp 65.07bp 271.70bp},clip,max width=0.94\textwidth,max height=0.38\textheight]{Marchisio2013.pdf}
\end{sourceblock}
该表达式是分散相质量守恒的基本表述。如果应用一阶粒子速度矩变换，我们得到

\begin{sourceblock}
\includegraphics[page=80,pagebox=mediabox,trim={178.85bp 336.90bp 65.07bp 334.96bp},clip,max width=0.94\textwidth,max height=0.38\textheight]{Marchisio2013.pdf}
\end{sourceblock}
其中二阶速度矩张量定义为

\begin{sourceblock}
\includegraphics[page=80,pagebox=mediabox,trim={172.39bp 281.20bp 65.07bp 384.48bp},clip,max width=0.94\textwidth,max height=0.38\textheight]{Marchisio2013.pdf}
\end{sourceblock}
平均加速度项为

\begin{sourceblock}
\includegraphics[page=80,pagebox=mediabox,trim={185.43bp 223.51bp 65.07bp 442.17bp},clip,max width=0.94\textwidth,max height=0.38\textheight]{Marchisio2013.pdf}
\end{sourceblock}
和碰撞加速度项

\begin{sourceblock}
\includegraphics[page=80,pagebox=mediabox,trim={191.97bp 165.81bp 65.07bp 499.86bp},clip,max width=0.94\textwidth,max height=0.38\textheight]{Marchisio2013.pdf}
\end{sourceblock}
方程（2.28）是平均粒子速度 Up 的输运方程，其结构与之前报道的方程非常相似。19 左侧第一项代表累积，而第二项代表物理空间中的速度对流。右侧第一项表示作用在粒子上的力产生的加速度，而第二项表示由于离散事件（例如粒子碰撞）引起的速度变化。

\footnotetext[19]{请注意，术语 ∂/∂x·(Up,2N) 是空间梯度向量和由速度向量的二进积产生的二阶张量的标量积。}
我们顺便指出，方程中的几乎所有项。 (2.28) 在时刻级别未关闭。我们还可以注意到，对于颗粒密度和尺寸恒定的情况，关系 ρpαp = ρpkVL3N 将允许我们重写方程： (2.28) 作为分散相动量的平衡 (ρpαpUp):20

\begin{sourceblock}
\includegraphics[page=81,pagebox=mediabox,trim={169.47bp 557.25bp 67.56bp 114.60bp},clip,max width=0.94\textwidth,max height=0.38\textheight]{Marchisio2013.pdf}
\end{sourceblock}
其中 mp 是（恒定）粒子质量。该表达式是分散相动量守恒的基本表述。右侧的术语表示与连续流体相和体积力（例如重力）的动量交换。二阶张量 Up,2 在矩输运方程层面上是未封闭的。最简单的可能闭包是 Up,2 = Up ⊗Up，它忽略了分散相中的速度波动。使用这个闭包，等式。 （2.27）和（2.32）描述了以保守形式编写的分散相的质量和动量守恒。21由于分散相质量和动量是基本守恒量，因此从这些量的一般守恒方程出发开发数值格式通常很有用。为了从 GPBE 中导出它们，我们必须知道粒子质量对内坐标的函数依赖性，我们用 m(xi) 表示。根据定义，颗粒材料密度和分散相体积分数（即分散相质量）与质量的关系为

\begin{sourceblock}
\includegraphics[page=81,pagebox=mediabox,trim={177.83bp 363.37bp 67.56bp 302.31bp},clip,max width=0.94\textwidth,max height=0.38\textheight]{Marchisio2013.pdf}
\end{sourceblock}
同样，分散相动量与粒子质量的关系为

\begin{sourceblock}
\includegraphics[page=81,pagebox=mediabox,trim={169.75bp 310.97bp 67.56bp 354.72bp},clip,max width=0.94\textwidth,max height=0.38\textheight]{Marchisio2013.pdf}
\end{sourceblock}
因此，为了导出分散相质量守恒方程和动量守恒方程，只需将 GPBE 分别乘以 m(xi) 和 vm(xi)，并对相空间变量进行积分即可。在开发多分散型多相流模型时，采用颗粒尺寸调节通常很有用。例如，在气固流中，气相和固体颗粒之间的动量交换项将取决于颗粒尺寸。因此，假设粒子具有内坐标向量 ψ = ζ，则条件粒子速度将取决于 ζ 的值。从数学上讲，我们可以通过 n(v, ze) = fU| ze(v|ze)n ze(ze) 来定义条件概率密度函数22 (PDF) fU | ze，其中，使用 NDF 的定义，

\begin{sourceblock}
\includegraphics[page=81,pagebox=mediabox,trim={200.52bp 162.13bp 67.56bp 503.55bp},clip,max width=0.94\textwidth,max height=0.38\textheight]{Marchisio2013.pdf}
\end{sourceblock}
请注意，fU|ψ 确实是一个条件 PDF，因为

\begin{sourceblock}
\includegraphics[page=81,pagebox=mediabox,trim={206.69bp 109.91bp 67.56bp 555.77bp},clip,max width=0.94\textwidth,max height=0.38\textheight]{Marchisio2013.pdf}
\end{sourceblock}
20 质量恒定的粒子在碰撞过程中动量守恒要求 h1 = 0.21 Ap,0 项的闭包将分散相与连续相耦合，并且需要单独的模型（例如斯托克斯阻力）。

\footnotetext[22]{fU|ψ(v, ze) 也可以解释为速度为 v 的粒子在所有内坐标为 ze 的粒子中所占的比例。}
对于所有 z 值。使用条件 PDF，我们可以定义任意阶的条件速度矩。例如，前两个时刻是

\begin{sourceblock}
\includegraphics[page=82,pagebox=mediabox,trim={187.14bp 581.53bp 65.07bp 84.15bp},clip,max width=0.94\textwidth,max height=0.38\textheight]{Marchisio2013.pdf}
\end{sourceblock}
\begin{sourceblock}
\includegraphics[page=82,pagebox=mediabox,trim={164.56bp 541.19bp 65.07bp 124.48bp},clip,max width=0.94\textwidth,max height=0.38\textheight]{Marchisio2013.pdf}
\end{sourceblock}
以及

从 GPBE（方程（2.16））开始，我们可以通过用 fU|Σ(v|Σ)nΣ(Σ) 代替 n 并在速度空间上积分来找到 nΣ(Σ) 的输运方程：

\begin{sourceblock}
\includegraphics[page=82,pagebox=mediabox,trim={121.89bp 480.48bp 65.07bp 185.21bp},clip,max width=0.94\textwidth,max height=0.38\textheight]{Marchisio2013.pdf}
\end{sourceblock}
\begin{sourceblock}
\includegraphics[page=82,pagebox=mediabox,trim={178.49bp 440.12bp 65.07bp 225.57bp},clip,max width=0.94\textwidth,max height=0.38\textheight]{Marchisio2013.pdf}
\end{sourceblock}
其中

比较方程。将式（2.39）与 PBE（式（2.14））相结合，我们发现，如果速度与内坐标无关，则它们是等价的：23 〈Up|zel = ze〉= Up，如果增长项不依赖于 v，则 〈 ̇zel|zel = ze〉= ̇zel(ze)，并且如果

\begin{sourceblock}
\includegraphics[page=82,pagebox=mediabox,trim={208.89bp 369.02bp 65.07bp 296.67bp},clip,max width=0.94\textwidth,max height=0.38\textheight]{Marchisio2013.pdf}
\end{sourceblock}
一般来说，这些条件不成立，因此等式：对于粒子速度因粒子不同而不同的系统，应使用方程 (2.39) 求 nψ。条件粒子速度的输运方程也可以从 GPBE 出发找到：

\begin{sourceblock}
\includegraphics[page=82,pagebox=mediabox,trim={78.67bp 256.32bp 65.07bp 389.05bp},clip,max width=0.94\textwidth,max height=0.38\textheight]{Marchisio2013.pdf}
\end{sourceblock}
\begin{sourceblock}
\includegraphics[page=82,pagebox=mediabox,trim={170.41bp 215.95bp 65.07bp 449.73bp},clip,max width=0.94\textwidth,max height=0.38\textheight]{Marchisio2013.pdf}
\end{sourceblock}
其中

虽然该方程包含几个未封闭项，但它确实提供了在多相流模型背景下解释颗粒尺寸效应的可能性，而无需求助于完整的 GPBE。相反，等式。 (2.39) 和(2.42) 可以通过适当的闭包一起求解。在典型的应用中，条件粒子速度是通过忽略方程左侧来找到的。 （2.42），仅留下 〈Up|zel=z〉 的代数表达式。这种近似适用于加速度项（例如流体阻力24）比空间传输的影响强得多的情况。

23 回想一下，在 PBE 中，我们用 v 表示已知速度，这在方程 2 中是等效的。 (2.39) 至 Up = v. 24流体阻力项的形式为 (v −Uf)/τ(xi)，其中 τ(xi) 是内坐标为 xi 的粒子的阻力时间尺度。则条件平均值为 (〈Up|ψ = ze〉−Uf)/τ(ze)。对于较小的 τ(ζ)（即较小的斯托克斯数），可以组合右侧的项来找到 〈Up|ψ = ze〉= Uf + τ(ze) · RHS。然后可以将该结果用于等式中。 (2.39) 求解 n xi，其中 RHS 是粒子速度与从方程 (2.39) 中找到的流体速度的偏差项。 （2.42）。因此，在斯托克斯数的一阶范围内，在没有体积力和碰撞的情况下，我们发现 〈Up|ψ = ze〉nψ = Ufnψ

\begin{sourceblock}
\includegraphics[page=82,pagebox=mediabox,trim={196.37bp 52.91bp 196.77bp 629.09bp},clip,max width=0.94\textwidth,max height=0.38\textheight]{Marchisio2013.pdf}
\end{sourceblock}
\section*{2.4 PBE 的流态}\phantomsection\addcontentsline{toc}{section}{2.4 PBE 的流态}
方程式中的 PBE方程（2.14）包含物理空间中的传输术语，并且如所写的，假设平流速度 v(t, x) 是准确已知的。然而，在许多实际应用中，平流速度将是湍流，并且只有低阶湍流统计数据是已知的。如果速度为层流，则 NDF 由于平流而变化，对于细颗粒，由于布朗运动而发生扩散，而当速度为湍流时，由于湍流混合，n 会出现湍流波动，这一点也必须考虑在内。我们通过下面两个简单的例子来展示它是如何工作的。 GPBE 中的湍流处理比 PBE 复杂得多，因为 GPBE 中的湍流通常是由分散相和连续相之间的耦合产生的。换句话说，PBE 中的 NDF 是被动标量，它从连续相的湍流继承了湍流属性，而 GPBE 中的 NDF 是动态主动标量，可以生成自己的湍流。在本书中，我们不考虑在 GPBE 背景下处理多相湍流模型的问题。对这个主题感兴趣的读者可以查阅 Fox (2012) 的评论。一般来说，设备规模应用的多相流模型的当前技术水平受到单相湍流模型的强烈影响，例如 Pope (2000) 中描述的模型。例如，大多数模型从分散相质量和动量方程（即方程（2.27）和（2.32））的临时闭包开始，并引入分散相湍流动力学方程和能量耗散的输运方程。正如 Fox (2012) 中所描述的，一种更基本的方法包括使用介观尺度闭包直接与 GPBE 合作，这些闭包之前已经针对 GPBE 的完全解析解决方案进行了验证。因此，本书的重点是通过仿真方法获得 GPBE 的完全解析解。

\subsection*{2.4.1 层流PBE}\phantomsection\addcontentsline{toc}{subsection}{2.4.1 层流PBE}
在层流的情况下，NDF 以流体速度 v 进行平流。然而，有趣的是，由于作用在每个颗粒上的阻力不同，不同尺寸的颗粒群可能具有不同的速度。因此，如上所述，对于方程的小斯托克斯数限制。 （2.42），有效平流速度 〈Up|ψ = ze〉 可以取决于颗粒尺寸。同样，当粒子足够小时，它们将受到布朗运动的影响，这可以被认为是粒子速度的随机、混沌分量，必须添加到平均速度中。在拉格朗日框架中描述布朗运动时，条件粒子速度 〈Up|ψ = ze〉 可以写为平均分量加上波动分量之和，而在欧拉框架中，它可以用粒子通量来表示，该通量表示为尺寸相关的扩散系数乘以平均 NDF 梯度，使用与动力学理论中定义分子属性（例如流体粘度）相同的方法。在此限制下，方程。 (2.39) 可以写成封闭形式

\begin{sourceblock}
\includegraphics[page=83,pagebox=mediabox,trim={149.82bp 119.04bp 67.56bp 552.82bp},clip,max width=0.94\textwidth,max height=0.38\textheight]{Marchisio2013.pdf}
\end{sourceblock}
其中 Γp(xi) 是有效扩散系数，可以通过 Nernst-Einstein 理论等进行估计（Bird 等人，1960）。请注意，在编写此方程时，我们隐含地假设 h 可以用 n 来表示，因此右侧是封闭的。总之，我们应该注意到方程的数值解。 （2.44）非常具有挑战性，因为有效扩散率通常比平流速度小许多数量级。换句话说，与方程相关的施密特数。 (2.44)（取决于 Ψ）在大多数实际应用中都非常大。

\subsection*{2.4.2 湍流PBE}\phantomsection\addcontentsline{toc}{subsection}{2.4.2 湍流PBE}
在湍流平流速度的情况下，PBE（即 NDF）中的传输量在其平均值附近波动。这些波动是由于连续相动量方程中的非线性对流项造成的。在湍流中，通常引入雷诺平均值（Pope，2000）。25 它包括计算感兴趣的系综平均量（通常是低阶矩）。给定湍流 φ(t, x) 的脉动特性，其在时间和空间固定点的雷诺平均值可以写为

\begin{sourceblock}
\includegraphics[page=84,pagebox=mediabox,trim={202.29bp 473.02bp 65.07bp 193.43bp},clip,max width=0.94\textwidth,max height=0.38\textheight]{Marchisio2013.pdf}
\end{sourceblock}
其中 fφ(t, x, ψ) 是 φ(t, x) 的单点 PDF，由以下概率语句定义：fφ(ψ)dψ = P[ψ ≤φ(t, x) < ψ + dψ]。该表达式右侧的概率是根据湍流中可观察到的所有随机场 φ(t, x) 的集合计算的（Fox，2003；Pope，2000）。 PBE 中主要关注的随机场是由 PBE 定义的流体相速度 Uf 和 NDF nψ（方程（2.44））。我们可以将 Uf 和 nψ 的单点联合 PDF 定义为 fUf,nψ(v, n*

\begin{sourceblock}
\includegraphics[page=84,pagebox=mediabox,trim={67.52bp 385.27bp 65.06bp 286.45bp},clip,max width=0.94\textwidth,max height=0.38\textheight]{Marchisio2013.pdf}
\end{sourceblock}
ψ、v 是流体速度状态空间，其中包括流体速度 Uf 的所有可能值。请注意，从概念上讲，fUf,nxi 与湍流反应流中使用的联合速度-成分 PDF 具有相同的含义（Fox，2003）。事实上，函数 n* Ψ 代表无限个由 Ψ 参数化的标量场。换句话说，每个值 ψ 对应于由湍流速度场 Uf 传输的具有“浓度”n ψ 的被动标量。我们现在可以将雷诺平均流体速度定义为

\begin{sourceblock}
\includegraphics[page=84,pagebox=mediabox,trim={178.24bp 271.87bp 65.07bp 394.59bp},clip,max width=0.94\textwidth,max height=0.38\textheight]{Marchisio2013.pdf}
\end{sourceblock}
同样，雷诺平均 NDF 定义为

\begin{sourceblock}
\includegraphics[page=84,pagebox=mediabox,trim={176.90bp 223.24bp 65.07bp 443.22bp},clip,max width=0.94\textwidth,max height=0.38\textheight]{Marchisio2013.pdf}
\end{sourceblock}
还可以将每个量的波动分量定义为瞬时值和雷诺平均值之间的差：uf = Uf −〈Uf〉和 n′ ψ = nψ −〈nψ〉。这些量用于高阶项的雷诺分解： 〈Ufnψ〉= 〈Uf〉〈nψ〉+ 〈ufn′

\begin{sourceblock}
\includegraphics[page=84,pagebox=mediabox,trim={171.28bp 169.06bp 259.46bp 515.22bp},clip,max width=0.94\textwidth,max height=0.38\textheight]{Marchisio2013.pdf}
\end{sourceblock}
ψ〉称为 NDF 的湍流。物理上，该术语表示给定 n′ xi 的流体元素通过流体速度波动而产生的平流，因此作为空间通量进入输运方程。雷诺平均 NDF 可以从方程式中找到。 （2.44）：26

\begin{sourceblock}
\includegraphics[page=84,pagebox=mediabox,trim={131.54bp 103.77bp 65.07bp 568.08bp},clip,max width=0.94\textwidth,max height=0.38\textheight]{Marchisio2013.pdf}
\end{sourceblock}
\footnotetext[25]{可以在大涡模拟 (LES) 的背景下发展类似的理论（Fox，2012）。 26 进行雷诺平均时，内坐标矢量 ψ 不受影响。因此，雷诺平均值随时间、空间和相空间梯度交换。有关此主题的讨论，请参见 Fox (2003)。}
为简单起见，我们假设 ̇ 只取决于 xi，而不是例如由流动传输的其他标量场。27 使用雷诺分解和湍流通量的梯度扩散模型（Fox，2003），我们可以写 〈Ufnψ〉= 〈Uf〉〈nψ〉−ΓT ∇x〈nψ〉，封闭式（1）中的对流项。 （2.48）。通过求解流体相中的湍流动能 (kf) 和湍流耗散率 (εf) 的单独模型，可以找到湍流扩散率 ΓT。因此，使用这种近似，方程。 (2.48) 减少为

\begin{sourceblock}
\includegraphics[page=85,pagebox=mediabox,trim={85.44bp 516.44bp 67.56bp 150.59bp},clip,max width=0.94\textwidth,max height=0.38\textheight]{Marchisio2013.pdf}
\end{sourceblock}
在高雷诺数湍流中，“分子”扩散项 Гp(xi) 与 ГT 相比可以忽略不计（Fox，2003）。除了最后一项，方程。 (2.49) 是封闭的。在许多实际应用中（例如粒子聚集），hψ 将是 nψ 的二次方。因此，与 〈hxi〉 相关的封闭问题与封闭湍流反应流的化学源项时所面临的问题非常相似（Fox，2003）。例如，如果 hψ 中所考虑的现象的时间尺度比湍流时间尺度慢得多，那么我们可以用 〈nψ〉 来封闭 〈hψ〉。28 然而，一般来说，必须诉诸于使用微混合模型来捕获湍流速度场和 NDF 之间的相互作用。其他常见的非封闭现象是流体剪切引起的颗粒聚集和破碎，这取决于 v 的瞬时空间梯度的大小。在湍流流场中，|∇xv| ≫|∇x〈Uf〉|而且，|∇xv|在空间和时间上变化很大。因此，雷诺平均聚集率和破损率必须使用基于高雷诺数湍流缩放定律的 kf 和 εf 值来估计（Fox，2003）。最后，读者应该清楚地记住这样一个事实，即使当等式： (2.49) 是封闭的，NDF 〈nxi〉仍然取决于许多自变量（即时间、空间和内坐标）。在实践中，直接求解方程。通过离散化所有自变量来求解方程（2.49）是很困难的，并且需要额外的近似。在本书中，我们简化 PBE 的主要方法是矩收敛。

\section*{2.5 矩闭合问题}\phantomsection\addcontentsline{toc}{section}{2.5 矩闭合问题}
作为本章的总结，我们使用方程 1 中的 NDF 输运方程简要介绍一下矩闭合问题。 （2.49）。为了清楚起见，我们将删除尖括号并假设 PBE 中的源项已关闭。然后 PBE 读取

\begin{sourceblock}
\includegraphics[page=85,pagebox=mediabox,trim={119.17bp 149.01bp 67.56bp 518.03bp},clip,max width=0.94\textwidth,max height=0.38\textheight]{Marchisio2013.pdf}
\end{sourceblock}
为了减少自变量的数量，我们使用式（2）中矩的定义。 （2.2）找到对应于方程（2.2）的矩输运方程。 （2.50）：

\footnotetext[27]{例如，如果生长速率取决于流体相中化学物质的浓度，则该假设将不成立。}
\footnotetext[28]{在湍流反应流中，这被称为慢化学极限（Fox，2003）。}
\begin{sourceblock}
\includegraphics[page=86,pagebox=mediabox,trim={88.51bp 555.27bp 65.07bp 56.44bp},clip,max width=0.94\textwidth,max height=0.38\textheight]{Marchisio2013.pdf}
\end{sourceblock}
在此表达式中，闭合项出现在等号的左侧，非闭合项出现在等号的右侧。然后，矩闭合问题包括寻找右侧项的函数逼近，该逼近仅依赖于一组有限值 k 的矩 mxi,k。实际上，在方程的数值模拟过程中，矩闭合会“即时”重新计算。 （2.51）。由于矩取决于时间 t 和空间 x，因此在单次模拟过程中必须多次计算矩闭合。因此，计算算法必须快速且高效（即基于多元迭代解决方案的算法不太可能易于处理）。在本书中，我们重点关注基于求积的矩方法 (QBMM)，该方法从一组有限的矩 mψ,k 重建 nψ。在第3章中，我们详细介绍了QBMM中使用的矩反演算法。在后续章节中，我们将深入讨论如何使用 QBMM 来闭合矩输运方程，以及用于求解方程 1 的鲁棒数值算法。 (2.51) 以及从 GPBE 导出的相关矩输运方程。

\chapter{基于求积的矩方法}
在本章中，我们讨论高斯求积的基本理论，它是基于求积的矩方法（QBMM）的核心。大多数结果的证明均未包括在内，对于需要更广泛分析处理的读者，可以参考文献。除了相关理论的总结之外，还提出了根据一组已知矩计算节点（或节点）和正交近似权重的不同算法，并批判性地讨论了它们的优缺点。需要提醒读者的是，求积公式的大部分理论都是针对一维积分（即单个自变量的积分）发展起来的。因此，下面的讨论从高斯求积理论完全适用的单变量分布开始，然后转向双变量和多元分布。尽管在后一种情况下，求积不再是严格的高斯分布，但其大多数有趣的属性仍然有效。在单变量情况下，权重与节点用于矩求积法 (QMOM) 来求解矩输运方程。因此，我们将使用完整矩集的矩反演算法称为 QMOM，而使用简化集的其他方法将被不同地称为。

\section*{3.1 单变量分布}\phantomsection\addcontentsline{toc}{section}{3.1 单变量分布}
如前面章节所述，对于单变量 QMOM，闭合问题出现在 NDF 矩的输运方程的不同部分中。然而，它通常可以简化为以下积分：

\begin{sourceblock}
\includegraphics[page=87,pagebox=mediabox,trim={206.04bp 122.05bp 67.56bp 543.76bp},clip,max width=0.94\textwidth,max height=0.38\textheight]{Marchisio2013.pdf}
\end{sourceblock}
其中 n(xi) 是未知的单变量 NDF，g(xi) 是将积分中出现的除 NDF 本身之外的所有项组合在一起的函数，Ωψ 是积分区间，也称为支持区间。正如已经报道的那样，在以粒子质量（或粒子尺寸或表面积）作为内坐标的情况下，支持区间变为 Ω xi = [0, ∞)，而在粒子速度的情况下，则支持区间为 Ω xi = (− ∞, + ∞)。封闭问题可以通过任何能够准确计算方程式中报告的积分的数值方案来克服。 (3.1)。这里需要强调的是，由于 NDF 未知，并且只有其矩被传递（因此已知），因此必须从 NDF 的传递矩开始构建该数值方案。高斯求积就是这样一种方案，但是，在讨论细节之前，让我们先介绍一些初步概念。

在高斯求积理论中，NDF 称为权重函数或测度。权重函数在积分区间及其所有时刻必须是非负且非空的，

\begin{sourceblock}
\includegraphics[page=88,pagebox=mediabox,trim={172.77bp 506.73bp 65.07bp 159.08bp},clip,max width=0.94\textwidth,max height=0.38\textheight]{Marchisio2013.pdf}
\end{sourceblock}
必须存在（Gautschi，2004；Press 等，1992）。代表真实多相颗粒系统的任何 NDF 显然都满足这些要求。在本章中，k 阶矩也将表示为 m(k) 或 〈ψk〉，这三个符号将互换使用。另一个重要的定义是正交多项式的定义（Gautschi，2004）。一组多项式\{P0(Ψ), P1(Ψ), .。。 , Pα(xi), .。。 \}，其中 Pα(xi) = kα,0xα+ kα,1xα−1 + · · · + kα,α，相对于权重函数，在积分区间 Ωϋ 中被认为是正交的，如果

\begin{sourceblock}
\includegraphics[page=88,pagebox=mediabox,trim={163.61bp 383.14bp 65.07bp 279.71bp},clip,max width=0.94\textwidth,max height=0.38\textheight]{Marchisio2013.pdf}
\end{sourceblock}
并且是正交的，如果

\begin{sourceblock}
\includegraphics[page=88,pagebox=mediabox,trim={163.05bp 330.51bp 65.07bp 332.34bp},clip,max width=0.94\textwidth,max height=0.38\textheight]{Marchisio2013.pdf}
\end{sourceblock}
积分域 Ω 和权重函数 n(xi) 唯一地定义了多项式族 \{Pα(xi)\}。当多项式的首项系数（即 kα,0）等于 1 时，多项式被定义为单项式。下面报告了两个重要的定理（没有证明）。

定理 3.1 任意一组正交多项式 \{Pα(xi)\} 都具有按以下顺序关联任意三个连续多项式的递推公式：

\begin{sourceblock}
\includegraphics[page=88,pagebox=mediabox,trim={138.81bp 234.49bp 65.07bp 452.55bp},clip,max width=0.94\textwidth,max height=0.38\textheight]{Marchisio2013.pdf}
\end{sourceblock}
\begin{sourceblock}
\includegraphics[page=88,pagebox=mediabox,trim={67.52bp 130.99bp 65.07bp 473.85bp},clip,max width=0.94\textwidth,max height=0.38\textheight]{Marchisio2013.pdf}
\end{sourceblock}
递归关系是正交多项式的构造和计算使用中最重要的属性。事实上，如下所示，了解递归系数可以计算正交多项式的零点，并利用它们计算求积规则。因此，计算这种三项递推关系的系数至关重要。方程中的递归关系(3.5) 生成与所考虑的积分区间中的权重函数正交的一元多项式序列。可以根据式（1）计算a0。 (3.6)，然后根据方程 P1(xi) （3.5）。知道 P0(xi) 和 P1(xi)，就可以从等式 1 计算 a1 和 b1。 （3.6）和（3.7），等等。还很容易证明，如果显式地表示多项式并且逐项进行乘法，最终可以将系数 aα 和 bα 写成 NDF 的矩。此外，构建 N 阶多项式所需的系数可以根据 NDF 的前 2N -1 个矩计算，即 m0、m1、...。。。 , m2N−2, m2N−1.例如，使用 m0、m1、m2 和 m3，可以计算以下系数：

\begin{sourceblock}
\includegraphics[page=89,pagebox=mediabox,trim={184.71bp 518.66bp 269.64bp 161.57bp},clip,max width=0.94\textwidth,max height=0.38\textheight]{Marchisio2013.pdf}
\end{sourceblock}
\begin{sourceblock}
\includegraphics[page=89,pagebox=mediabox,trim={184.71bp 449.94bp 67.56bp 187.77bp},clip,max width=0.94\textwidth,max height=0.38\textheight]{Marchisio2013.pdf}
\end{sourceblock}
反过来，这足以计算与积分区间中的权重函数正交的多项式 P2(xi)（二阶）。下面是另一个重要定理，它确保一旦指定了积分区间，正交多项式的零点必须限制在该区间内。这是正交多项式的一个重要特征，因为正如下面将变得更清楚的那样，它保证内部坐标始终在 NDF 的支持内进行采样。

定理 3.2 正交序列中的每个多项式的所有根都是实数、不同的，并且严格位于积分区间（也称为正交区间）内。每个多项式的根严格位于序列中下一个高阶多项式的根之间（交错根）。

为什么我们对正交多项式感兴趣？因为，正如下一节将解释的，它们的根是高斯求积近似的节点。

\subsection*{3.1.1 高斯求积}\phantomsection\addcontentsline{toc}{subsection}{3.1.1 高斯求积}
在单变量 QMOM 中，封闭问题可以通过任何能够计算式（1）中积分的数值方案来克服。 (3.1)。一个非常简单的方法是使用插值公式：

\begin{sourceblock}
\includegraphics[page=89,pagebox=mediabox,trim={185.30bp 162.91bp 67.56bp 500.67bp},clip,max width=0.94\textwidth,max height=0.38\textheight]{Marchisio2013.pdf}
\end{sourceblock}
其中 wα 和 xiα 分别是插值公式的权重与节点，N 是用于插值的节点数（或公式的阶数）。插值公式的准确性是通过其精确度来量化的。1 当被积函数是小于或等于 d 阶的多项式且至少存在一个 d + 1 阶多项式使得插值公式不精确时，如果插值公式是精确的，则精确度等于 d。插值公式的主要要求之一是具有不同的节点，并且根据它们的放置方式

\footnotetext[1]{准确度也称为精确度。}
积分区间，可以达到不同程度的准确度。例如，具有 N 个等距节点的插值公式（即 Newton-Cotes）的精度为 N −1，而高斯求积的精度为 2N −1。因此，很清楚为什么正交近似如此有趣：在相同数量的节点（以及因此相同的函数评估）下，它们的准确度基本上是其他公式的两倍。然而，重要的是要提醒读者，准确度的定义是基于求积在处理多项式作为被积函数时的性能，因此只有当被积函数足够平滑以至于可以用多项式逼近时，高阶才能转化为高精度。高斯求积是一个插值公式，其 N 个节点是与权重函数 n(xi) 正交的多项式 PN(xi) 的根，如下一个定理中所述 (Gautschi, 1997, 2004)。

定理3.3 公式的充要条件

\begin{sourceblock}
\includegraphics[page=90,pagebox=mediabox,trim={171.93bp 453.00bp 65.07bp 210.58bp},clip,max width=0.94\textwidth,max height=0.38\textheight]{Marchisio2013.pdf}
\end{sourceblock}
要成为高斯求积近似，或者等价地，其精度为 2N −1，其节点 \{xiα\} 与 Ω 中关于权重函数 n(xi) 正交的 N 阶多项式 PN(xi) 的 N 个根重合。定理 3.3 意味着只要被积函数 g(xi) 是阶数小于或等于 2N −1 的多项式，2 N 阶的求积近似就是精确的（即剩余函数 RN(g) 为空）。著名的高斯求积包括与经典正交多项式相关的求积，例如高斯-勒让德多项式、高斯-雅可比多项式、高斯-拉盖尔多项式、高斯-埃尔米特多项式和高斯-切比雪夫多项式。然而，给定一个特定的非经典权函数，可以使用专门的算法来计算相应的正交。执行此计算的一种直接方法是利用其高精度。由于 N 阶高斯求积的精度为 2N -1，因此它将是唯一能够捕获权重函数的前 2N -1 矩的近似值。因此，N个权重和N个节点可以通过求解以下非线性系统来确定：

\begin{sourceblock}
\includegraphics[page=90,pagebox=mediabox,trim={204.36bp 108.47bp 65.07bp 433.31bp},clip,max width=0.94\textwidth,max height=0.38\textheight]{Marchisio2013.pdf}
\end{sourceblock}
该非线性系统可以使用牛顿-拉夫逊法或任何其他非线性方程求解器来求解。然而，这不是最好的方法，因为需要非常好的初始猜测来确保收敛。

\footnotetext[2]{也称为 N 点求积。}
一种更聪明的求根方法是利用等式（1）中的递归关系。 (3.5) 对于正交多项式（Dette \& Studden，1997；Golub \& Welsch，1969；Wilf，1962）。 α = 0, 1, 2, ... 的递归关系序列。。。 , N 可以用矩阵形式解释为

\begin{sourceblock}
\includegraphics[page=91,pagebox=mediabox,trim={84.75bp 480.31bp 67.56bp 108.92bp},clip,max width=0.94\textwidth,max height=0.38\textheight]{Marchisio2013.pdf}
\end{sourceblock}
方程(3.12)清楚地表明PN(xi)的根，即正交近似\{xiα\}的节点，是方程中出现的三对角矩阵的特征值。可以通过对角相似变换使该矩阵对称（保留特征值）以给出雅可比矩阵：

\begin{sourceblock}
\includegraphics[page=91,pagebox=mediabox,trim={132.38bp 310.70bp 67.56bp 276.01bp},clip,max width=0.94\textwidth,max height=0.38\textheight]{Marchisio2013.pdf}
\end{sourceblock}
该过程将求多项式根的病态问题转化为求三对角对称矩阵的特征值和特征向量的良条件问题。如 Wilf (1962) 所示，N 个权重可以计算为 wα = m0φ2 α1，其中 φα1 是雅可比矩阵的第 α 个特征向量 φα 的第一个分量。当然，上述方法需要了解方程1中报告的递归关系中出现的系数aα和bα。 （3.5）。反过来，可以通过利用正交条件并使用权重函数的矩来简单地计算系数。下面介绍并讨论了两种有效执行这些计算的算法。读者可以注意到，如果 NDF 对应于已知正交多项式族的权重函数（例如，高斯 NDF 是 Hermite 多项式的权重函数（Gautschi，2004）），则递归系数是预先已知的，因此等式：可以直接应用方程（3.13）来求权重与节点，而不需要使用矩。在 3.3 节中，我们利用这一事实来开发扩展的正交。相反，对于任意 NDF，可以根据矩计算递归系数。

\subsection*{3.1.2 乘积差（PD）算法}\phantomsection\addcontentsline{toc}{subsection}{3.1.2 乘积差（PD）算法}
{\small 乘积差 (PD) 算法由 Gordon (1968) 开发，基于 Stieltjes 的连分数理论。第一步是构建矩阵 P\par}
从权重函数的矩开始的分量 Pα,β。 P 第一列中的分量为 Pα,1 = δα1（对于 α ε1,...）。。。 , 2N + 1，其中 δα1 是克罗内克增量。 P 的第二列中的分量为 Pα,2 = (−1)α−1mα−1（对于 α ε1,...）。。。，2N。应该强调的是，由于最终权重可以通过乘以真实的 m0 来校正，因此也可以假设标准化分布（即 m0 = 1）来进行计算。其余分量通过 PD 算法找到：Pα,β = P1,β−1Pα+1,β−2 −P1,β−2Pα+1,β−1 对于 β ∈3,...。。。 , 2N + 1 且 α ε1, .。。 , 2N + 2 −j。以N=2为例，P矩阵为

\begin{sourceblock}
\includegraphics[page=92,pagebox=mediabox,trim={119.10bp 473.04bp 65.07bp 156.63bp},clip,max width=0.94\textwidth,max height=0.38\textheight]{Marchisio2013.pdf}
\end{sourceblock}
确定三项递归关系的系数所需的连分数\{zα\}的系数是通过将第一个元素设置为零（z1 = 0）来生成的，然后根据以下递归关系计算其他元素：

\begin{sourceblock}
\includegraphics[page=92,pagebox=mediabox,trim={179.61bp 397.07bp 65.07bp 276.42bp},clip,max width=0.94\textwidth,max height=0.38\textheight]{Marchisio2013.pdf}
\end{sourceblock}
然后从 δα 的和与乘积获得对称三对角雅可比矩阵的系数：对于 α ε1,...，aα = zeta2α + zeta2α−1。。。，N 且 bα = −√ze2α+1ze2α 对于 α ε1,。。。，N-1。实现 PD 算法的 Matlab 脚本可以在附录 A 的 A.1.1 节中找到。

练习 3.1 考虑以下正态（或高斯）分布：

\begin{sourceblock}
\includegraphics[page=92,pagebox=mediabox,trim={182.07bp 287.86bp 65.07bp 376.64bp},clip,max width=0.94\textwidth,max height=0.38\textheight]{Marchisio2013.pdf}
\end{sourceblock}
众所周知，μ 是平均值，σ2 是分布的方差（Walpole \& Myers，1985）。该分布的矩可以通过以下方程轻松计算（仅报告前九个矩）：

\begin{sourceblock}
\includegraphics[page=92,pagebox=mediabox,trim={145.84bp 95.31bp 65.07bp 461.61bp},clip,max width=0.94\textwidth,max height=0.38\textheight]{Marchisio2013.pdf}
\end{sourceblock}
{\small 使用 PD 算法计算 μ = 5 和 σ = 1 时的四阶求积近似（即 N = 4）。对于此计算，需要分布的前八个矩：\par}
\begin{sourceblock}
\includegraphics[page=93,pagebox=mediabox,trim={216.64bp 522.65bp 217.53bp 59.44bp},clip,max width=0.94\textwidth,max height=0.38\textheight]{Marchisio2013.pdf}
\end{sourceblock}
应用PD算法后，得到以下雅可比矩阵：

\begin{sourceblock}
\includegraphics[page=93,pagebox=mediabox,trim={179.66bp 442.78bp 67.56bp 199.62bp},clip,max width=0.94\textwidth,max height=0.38\textheight]{Marchisio2013.pdf}
\end{sourceblock}
得出权重 w1 = 0.0459、w2 = 0.4541、w3 = 0.4541 和 w4 = 0.0459，以及节点 Ψ1 = 2.6656、Ψ2 = 4.2580、Ψ3 = 5.7420 和 Ψ4 = 7.3344。留给读者验证的是，通过这种求积近似，可以准确计算前八个矩，而 m8（以及高阶矩）则会出现错误。

PD算法在很多实际情况下都相当高效；然而，随着 N 的增加，它通常会变得不太稳定。很难先验地预测何时会发生这种情况，因为它取决于矩的绝对值，但通常当 N > 10 时就会出现问题。另一个重要问题与以下事实有关：对于均值 3 为零（或换句话说，m1 = 0）的分布，由于在计算连分数 \{zα\} 的系数期间除以零，算法会崩溃。下一节将介绍的惠勒算法不会遇到这些问题。

\subsection*{3.1.3 惠勒算法}\phantomsection\addcontentsline{toc}{subsection}{3.1.3 惠勒算法}
Sack \& Donovan (1971) 提出了另一种方法来计算方程 1 中报告的递归公式的系数。 （3.5）（也出现在雅可比矩阵中）导致更高的稳定性。该方法基于使用一组不同的基函数 πα(Ψ) 来表示正交多项式的思想，而不是通常的 Ψ 幂。稳定性的提高源于新多项式基能够更好地对积分区间 Ω 进行采样。系数根据修改矩计算，定义如下：

\begin{sourceblock}
\includegraphics[page=93,pagebox=mediabox,trim={154.87bp 122.95bp 67.56bp 542.87bp},clip,max width=0.94\textwidth,max height=0.38\textheight]{Marchisio2013.pdf}
\end{sourceblock}
3 但当内坐标 xi 表示粒子质量、表面积、尺寸等时，m1 通常不为零。在这些情况下，可以安全地使用 PD 算法。当内坐标是粒子速度时，m1 为空的情况更常见，其范围从负实值到正实值，可以导致平均速度为零的分布。平均值为空的另一种常见情况是使用中心矩（相对于分布平均值平移的矩）来构建正交近似。稍后在描述构建多元求积的算法时将讨论这些情况。

类似于标准时刻所做的事情。此外，假设πk(xi)满足以下递归关系：

\begin{sourceblock}
\includegraphics[page=94,pagebox=mediabox,trim={174.02bp 557.63bp 65.07bp 98.82bp},clip,max width=0.94\textwidth,max height=0.38\textheight]{Marchisio2013.pdf}
\end{sourceblock}
当然，其中系数 a' α 和 b' α 必须明确已知。根据这些发现，Wheeler（1974）开发了一种有效的算法，通过一些中间量来计算雅可比矩阵的系数：

\begin{sourceblock}
\includegraphics[page=94,pagebox=mediabox,trim={163.06bp 470.34bp 65.07bp 195.47bp},clip,max width=0.94\textwidth,max height=0.38\textheight]{Marchisio2013.pdf}
\end{sourceblock}
这些数量可以通过初始化来计算

\begin{sourceblock}
\includegraphics[page=94,pagebox=mediabox,trim={177.65bp 365.33bp 65.07bp 266.94bp},clip,max width=0.94\textwidth,max height=0.38\textheight]{Marchisio2013.pdf}
\end{sourceblock}
然后计算 α = 1, 2,...。。。 , N −1 （其中 N 又是求积近似和正交多项式的阶数），得到以下方程： σα,β = σα−1,β+1 −

\begin{sourceblock}
\includegraphics[page=94,pagebox=mediabox,trim={153.34bp 314.08bp 196.89bp 362.30bp},clip,max width=0.94\textwidth,max height=0.38\textheight]{Marchisio2013.pdf}
\end{sourceblock}
βσα−1,β−1 对于 β = α, α + 1,...。。 , 2N -α -1。最后，雅可比矩阵的系数计算如下：

\begin{sourceblock}
\includegraphics[page=94,pagebox=mediabox,trim={187.73bp 243.59bp 65.07bp 404.94bp},clip,max width=0.94\textwidth,max height=0.38\textheight]{Marchisio2013.pdf}
\end{sourceblock}
有关该算法推导的详细信息可以在 Gautschi (1997, 2004) 中找到。 Wheeler算法具有较高的稳定性，通常可用于高阶正交近似的计算（与PD算法相比）。尽管修改矩和基函数 πα(xi) 的选择在算法的性能中起着非常重要的作用，但有趣的是，即使修改矩 νk 与标准矩 μk 相同（或者换句话说，a′

\begin{sourceblock}
\includegraphics[page=94,pagebox=mediabox,trim={72.49bp 146.41bp 389.66bp 538.10bp},clip,max width=0.94\textwidth,max height=0.38\textheight]{Marchisio2013.pdf}
\end{sourceblock}
α=0)Wheeler算法更加稳定。 Wheeler 算法的另一个优点是能够在均值为零（即 m1 = 0）的分布情况下计算雅可比矩阵的系数。实现 Wheeler 算法的 Matlab 脚本可以在附录 A 的 A.1.2 节中找到。

练习 3.2 再次考虑方程 1 中报告的分布。练习 3.1 的 (3.16)。现在让我们使用 Wheeler 算法计算 μ = 0 和 σ = 1 时的四阶正交近似（即 N = 4）。对于此计算，需要分布的前八个矩：

\begin{sourceblock}
\includegraphics[page=95,pagebox=mediabox,trim={227.43bp 522.65bp 228.31bp 59.44bp},clip,max width=0.94\textwidth,max height=0.38\textheight]{Marchisio2013.pdf}
\end{sourceblock}
请注意，由于分布的均值为零并且完全对称（零偏度），因此奇矩全部为零。应用Wheeler算法后，得到以下雅可比矩阵：

\begin{sourceblock}
\includegraphics[page=95,pagebox=mediabox,trim={179.66bp 419.36bp 67.56bp 223.03bp},clip,max width=0.94\textwidth,max height=0.38\textheight]{Marchisio2013.pdf}
\end{sourceblock}
得出权重 w1 = 0.0459、w2 = 0.4541、w3 = 0.4541 和 w4 = 0.0459，以及节点 Ψ1 = -2.3344、Ψ2 = -0.7420、Ψ3 = 0.7420 和 Ψ4 = 2.3344。留给读者验证的是，通过这种求积近似，可以准确计算前八个矩，而 m8（和高阶矩）会产生错误。也很容易验证PD算法在这种情况下无法计算雅可比矩阵的系数。

在实践中，惠勒算法返回可实现矩集的完整权重与节点集。但是，对于无法实现的矩集，它将失败。 Yuan \& Fox (2011) 提出了 Wheeler 算法的自适应版本，该算法检查可实现性并返回可能的最大权重与节点集。为了完成此任务，使用可实现矩集 0 < σα,α 的事实来检查，从 α = 1 开始，对于 1 ≤α ≤N 的每个值，是否 σα,α ≤0。如果对于 α = m，σm,m ≤0，则只能计算 m -1 个权重与节点。自适应惠勒算法对于退化情况特别有用，在这种情况下，精确的 NDF 由少于 N 个加权狄拉克 delta 函数组成。在这种情况下，自适应惠勒算法将返回精确的 NDF。有关自适应正交的更多详细信息，请参阅 Yuan \& Fox (2011)。

\subsection*{3.1.4 矩集的一致性}\phantomsection\addcontentsline{toc}{subsection}{3.1.4 矩集的一致性}
PD 和 Wheeler 算法在计算单变量正交近似值时非常有效。生成的节点始终位于支持范围内（因此代表内部坐标的可实现值），并且权重始终为正，确保当使用求积来克服闭合问题时，获得准确的结果。然而，只有当矩集可实现时，这些属性才会受到尊重，这意味着存在导致该特定矩集的 NDF。如果 PD 或 Wheeler 算法输入无效的矩集，则会计算无法实现的节点，从而危及模拟的稳定性。因此很明显，算法

\footnotetext[4]{矩集的可实现性或一致性的概念将在第 3.1.4 节中介绍。}
能够确定矩集的可实现性和一致性，以及纠正无效矩集，在这种情况下特别有用。因此，让我们分析将一组矩连接在一起的数学约束。 NDF 的矩代表了构成所研究的多相系统的基础元素群的一些重要物理特性。因此，他们必须满足一些简单的规则。例如，密度函数相对于其支持的正值意味着零阶矩必须为正。此外，还有其他简单、直观的规则。例如，如果内部坐标仅假定正值（或者换句话说，矩定义在正支撑上），则一阶矩（以及所有其他矩）必须为正。分布的另一个重要属性是它的方差（即 σ2 = m2 −m2 1/m0），对于 delta 函数分布它必须为零，而对于多分散分布它必须为正。因此，必须是m2≥m2 1/m0。对于高阶矩，数学约束不太直观，并且不能与多相系统的特定全局属性直接相关。幸运的是，矩理论提供了一些有趣的定理，这些定理对于确定一组矩是否无效非常有用。

定理 3.4 如果 Hankel-Hadamard 行列式 (Gautschi, 2004; Shohat \& Tamarkin, 1943) 都是非负的，则称矩集是可实现的：

\begin{sourceblock}
\includegraphics[page=96,pagebox=mediabox,trim={67.52bp 335.95bp 65.07bp 283.70bp},clip,max width=0.94\textwidth,max height=0.38\textheight]{Marchisio2013.pdf}
\end{sourceblock}
请注意，对于 k = 0 且 l = 1，Hankel-Hadamard 行列式为 m0m2 −m2 1；因此，有效性条件简化为 m2 ≥m2 1/m0。以下定理定义了不太严格的有效性条件（Shohat \& Tamarkin，1943）。

定理 3.5 有效性的必要（但非充分）条件是函数 ln(mk) 关于 k 的凸性：

\begin{sourceblock}
\includegraphics[page=96,pagebox=mediabox,trim={157.63bp 211.79bp 65.07bp 442.25bp},clip,max width=0.94\textwidth,max height=0.38\textheight]{Marchisio2013.pdf}
\end{sourceblock}
或等价地， mkmk−2 −m2
证明凸性条件相当于 k = 0, 1 和 l = 0, 1, 2 或前四个矩的 Hankel-Hadamard 行列式的正性是微不足道的。对于高阶矩，等价性会丢失，因为定理 3.4 需要更严格的条件。然而，事实证明凸性条件是有用的，其原因将在下面讨论。通过构建 ln(mk) 的差分表，可以轻松验证函数 ln(mk) 关于 k 的凸性。表的第一列（练习 3.3 的表 3.1 和表 3.2 中的 d0）报告了所有矩的自然对数。后续列中的元素是通过从同一行中的元素减去其下方行中的元素来从前一列中计算的。 ln(mk) 的凸性由该列的元素相对于二阶差分（即 d2）的正值来保证。同样有趣的是，对于对数正态分布的矩，当 ln(mk) 是 k 的二次函数时，三阶差分消失。因此，对数正态分布是应用凸性条件时产生可实现矩（即正二阶差分与最小三阶差分）的最平滑分布。以下示例说明了这两个有效性标准。

练习 3.3 让我们验证练习 3.1 的矩集是可实现的。利用前八个时刻，

\begin{sourceblock}
\includegraphics[page=97,pagebox=mediabox,trim={216.64bp 454.26bp 217.53bp 127.83bp},clip,max width=0.94\textwidth,max height=0.38\textheight]{Marchisio2013.pdf}
\end{sourceblock}
可计算的最大 Hankel-Hadamard 矩阵为 Δ0,3 和 Δ1,3，结果为

\begin{sourceblock}
\includegraphics[page=97,pagebox=mediabox,trim={173.85bp 379.08bp 67.56bp 267.25bp},clip,max width=0.94\textwidth,max height=0.38\textheight]{Marchisio2013.pdf}
\end{sourceblock}
以及

\begin{sourceblock}
\includegraphics[page=97,pagebox=mediabox,trim={163.47bp 316.67bp 67.56bp 329.67bp},clip,max width=0.94\textwidth,max height=0.38\textheight]{Marchisio2013.pdf}
\end{sourceblock}
其行列式分别为 12 和 5736。因为这两个值为正，所以矩集是可以实现的。事实上，正如练习 3.1 中所示，当将此矩集输入 PD（或 Wheeler）算法时，它会产生可实现的节点。正如预期的那样，前七个矩通过使用四阶求积近似完美地再现。计算ln(mk)的差值表如表3.1所示，也可以得出同样的结论。这里，矩集的自然对数在表的第一列（即 d0）中报告，其他元素被计算为后续行中的元素与前一列的差值。第二列 (d2) 包含二阶差分，其元素均为正，证实矩集是可实现的。现在将二阶矩修改为 25（而不是 26，对应的差异仅为 4\%）就足以使矩集无法实现。 Hankel-Hadamard 行列式现在分别等于 -179 352 和 -12 362 344，而差值表（见表 3.2）在包含二阶差值（即 d2）的列中呈现负元素。如果将该矩集输入 PD 算法，则所得的正交将无法重现该矩集。

一般来说，从可实现的矩集开始进行的模拟应该产生可实现的矩集。然而，矩输运方程是与一些有限的离散误差进行数值积分的。正如 Wright (2007) 清楚地报告的那样，大多数问题是由对流项的近似引起的，特别是高阶离散化方案，它可以将一组可实现的矩变成

\begin{center}\small 表3.1。可实现矩集 ln(mk) 的差分表\end{center}
\begin{sourceblock}
\includegraphics[page=98,pagebox=mediabox,trim={58.00bp 483.38bp 58.23bp 73.92bp},clip,width=\textwidth,max height=0.72\textheight]{Marchisio2013.pdf}
\par\vspace{0.45\baselineskip}
\small 表3.2。不可实现矩集 ln(mk) 的差分表
\end{sourceblock}
\begin{sourceblock}
\includegraphics[page=98,pagebox=mediabox,trim={58.00bp 330.92bp 58.23bp 232.82bp},clip,width=\textwidth,max height=0.72\textheight]{Marchisio2013.pdf}
\end{sourceblock}
一个无法实现的。这种现象称为矩损坏，对基于求积的矩方法构成严重威胁，因为当将一组不可实现的矩输入 PD 或惠勒算法时，可能会产生不可实现的节点并导致数值不稳定。一般来说，解决矩损坏的最佳方法是使用保证生成可实现矩的数值算法（详细信息请参阅附录 B）。然而，如果使用“黑匣子”传输代码（作者认为应该避免这种情况！），则必须采用一致的策略来纠正损坏的时刻。可能克服矩腐败问题的策略是在不可实现的矩集出现后立即替换它们。 McGraw (2006) 开发了这种算法。该算法首先验证矩集是否可实现（通过查看二阶差分向量或通过查看 Hankel-Hadamard 行列式）。如果设定的矩无法实现，则继续进行修正。矩校正算法可以基于将负二阶差分转化为正差分的矩的变化。然而，该策略存在两个问题。第一，二阶差分的正性只是有效性的必要条件，而不是充分条件。因此，将校正建立在必要条件的基础上是有风险的，因为校正后可能会出现仍然无法实现的矩设置。第二个问题与这一事实有关，即该校正策略引入了任意决定，即校正后二阶差分必须为正值，或者换句话说，人们希望设定的时刻距不可实现的条件有多远。

{\small 另一种方法是执行迭代最小化，识别并校正调整后使函数 ln(mk) 的平滑度最大化的时刻。这个\par}
是通过最小化 d3 来完成的，d3 是差异表中包含 ln(mk) 三阶差异的列，正如已经报道的那样，它在对数正态分布的时刻消失。最大化 ln(mk) 的平滑度（或者等效地最小化三阶差分）代表比确保二阶差分的正性（这只是必要条件）更严格的约束，并且通常会在几次迭代后成功收敛。对于由单个时刻的值导致无法实现的时刻组来说尤其如此。如果多个时刻的估计错误，算法可能会失败，并且必须使用另一种方法。最简单和最有效的方法是将矩替换为对数正态分布的矩，该分布与原始分布仅共享一些正确的矩，正如 Wright (2007) 所建议的那样。由于损坏通常会影响高阶矩，因此通常可以从低阶矩 m0、m1 和 m2 中找到对数正态分布的参数，正如所强调的，其物理意义尤为重要。然后，使用对数正态分布，可以计算其他矩。我们现在将详细讨论这些矩校正算法。

\paragraph*{麦格劳校正算法}
矩集的可实现性可以通过计算函数ln(mk)的差分表来验证。在第一列 (d0) 中，报告了按升序排列的矩序列。在接下来的列中，报告位于同一前列和两个后续行中的两个相邻元素之间的差异。继续该过程直到第一行出现零。不同列中报告的值代表函数 ln(mk) 的二阶、三阶和高阶差分向量。如果将某个时刻 mk 改为 (mk)1 = ck(mk)0，则其自然对数发生如下变化： ln(mk)1 = ln(ck) + ln(mk)0，其中 (mk)0 和 ln(mk)0 为初始值及其对数，(mk)1 和 ln(mk)1 为最终值及其对数。三阶差分向量的后续变化导致 (d3)1 −(d3)0 = ln(ck)bk，其中 (d3)0 和 (d3)1 分别是三阶差分的初始向量和最终向量，而 bk 是描述三阶差分向量对第 k 阶矩的自然对数单位增量的变化的响应向量。因此最终向量是 (d3)1 = (d3)0 + ln(ck)bk。正如已经讨论过的，校正矩集的有效方法是最小化三阶差分向量。通过简单的几何参数，可以立即看出，修改后的三阶差分向量 (d3)1 在与向量 bk 正交时呈现最小长度：\%(d3)0 + ln(ck)bk \& · bk = 0。因此，可以从以下关系获得最小化 (d3)1 的 ck 值：

\begin{sourceblock}
\includegraphics[page=99,pagebox=mediabox,trim={255.96bp 166.07bp 67.56bp 507.36bp},clip,max width=0.94\textwidth,max height=0.38\textheight]{Marchisio2013.pdf}
\end{sourceblock}
ln(ck) = −cos((d3)0, bk)|(d3)0|
其中 cos((d3)0, bk) 是向量 (d3)0 和向量 bk 形成的角度。所得的最小平方振幅结果为 |(d3)1|2 = |(d3)0 + ln(ck)bk|2 = |(d3)0|2 × cos2((d3)0, bk)。在算法中，选择校正时刻的索引是为了使函数ln(mk)达到最大平滑度。选择用于校正的时刻是导致三阶差分向量的最小可能幅度的时刻。对于索引 k 的任何矩，必须校正的矩的索引 k* 被选择为给出最大 cos2((d3)0, bk) 的索引。校正的实体由正交关系给出，即式（1）。 (3.29)，计算 k = k*。实现该校正算法的 Matlab 脚本可以在附录 A 的 A.2.1 节中找到。

练习3.4 使用矩校正算法修改以下矩集：

\begin{sourceblock}
\includegraphics[page=100,pagebox=mediabox,trim={219.13bp 500.73bp 215.04bp 81.36bp},clip,max width=0.94\textwidth,max height=0.38\textheight]{Marchisio2013.pdf}
\end{sourceblock}
将这组矩与练习 3.1 的矩进行比较，很明显，应该修改的矩是 m3，其原始值为 140。校正算法识别出该组无效，并将 k*= 3 确定为必须更改的矩的阶数，以最大化函数 ln(mk) 的平滑度。校正结果为 m3 = 140.001，该值与原始值惊人地接近。留给读者验证的是，在改变任何其他矩时，该算法不仅能够使矩集可实现，而且能够将矩恢复到其原始值。

练习3.5 使用矩校正算法修改以下矩集：

\begin{sourceblock}
\includegraphics[page=100,pagebox=mediabox,trim={219.13bp 252.71bp 215.04bp 329.39bp},clip,max width=0.94\textwidth,max height=0.38\textheight]{Marchisio2013.pdf}
\end{sourceblock}
在这种情况下，导致集合无效的已更改的矩是 m2 和 m3。当将该矩集输入校正算法时，得到以下答案：

\begin{sourceblock}
\includegraphics[page=100,pagebox=mediabox,trim={219.13bp 92.47bp 215.04bp 489.62bp},clip,max width=0.94\textwidth,max height=0.38\textheight]{Marchisio2013.pdf}
\end{sourceblock}
{\small 可以看出，所有的时刻都被改变了，虽然最后的时刻集是可以实现的，但它不是原来的时刻。\par}
总之，McGraw 的校正算法在大多数研究案例中成功地校正了矩集，或者换句话说，将不可实现的矩集转换为可实现的矩集。此外，当该算法应用于集合中大于四个的矩集合时，如果仅改变一个矩，该算法通常能够恢复原始矩集合。如果改变更多矩，该算法能够将集合转换为有效的集合，但通常不会将其恢复为原始值。当矩集仅包含四个矩时，三阶差分向量仅由一个分量构成，正交性条件（式（3.29））变得毫无意义。在这种极限情况下，构成该集合的四个矩中的任何一个的单个变化都将导致三阶差分为零；然而，只有其中一些会产生可实现的集合（具有正的二阶差异）。在这些情况下，以及在算法失败的所有其他情况下（即使对于由四个以上矩构成的集合），使用下面描述的过程会更方便。

\paragraph*{赖特校正算法}
第二个校正算法包括从对应于对数正态分布的矩集计算新的矩集。对数正态分布具有以下函数形式：

\begin{sourceblock}
\includegraphics[page=101,pagebox=mediabox,trim={174.38bp 375.78bp 67.56bp 289.53bp},clip,max width=0.94\textwidth,max height=0.38\textheight]{Marchisio2013.pdf}
\end{sourceblock}
对应于以下 k 阶矩：

\begin{sourceblock}
\includegraphics[page=101,pagebox=mediabox,trim={245.78bp 331.58bp 67.56bp 348.54bp},clip,max width=0.94\textwidth,max height=0.38\textheight]{Marchisio2013.pdf}
\end{sourceblock}
λk = Nt exp
其中 Υ 是内坐标，Nt 是总粒子数密度（与分布 λ0 的零阶矩相同），μ 和 σ 是与分布均值和方差相关的两个参数。显然，以三个参数（即 Nt、μ 和 σ）为特征的对数正态分布最多可以容纳属于特定集合的三个矩。此外，其中一个必须是零阶矩，而另外两个可以是任何其他矩。该校正过程背后的主要思想如下：选择原始集合的三个矩并用于计算对数正态分布的三个参数（即 Nt、μ 和 σ）；然后使用等式重新计算整个矩集。 （3.31）。正如已经强调的，由于三个矩中的一个必须是零阶矩，因此 Nt 简单地取为等于 m0。如果另外两个选择的矩是i阶和j阶的矩，则可以使用以下关系式来确定

\begin{sourceblock}
\includegraphics[page=101,pagebox=mediabox,trim={168.35bp 137.58bp 67.56bp 529.36bp},clip,max width=0.94\textwidth,max height=0.38\textheight]{Marchisio2013.pdf}
\end{sourceblock}
以及

\begin{sourceblock}
\includegraphics[page=101,pagebox=mediabox,trim={166.15bp 93.78bp 67.56bp 573.15bp},clip,max width=0.94\textwidth,max height=0.38\textheight]{Marchisio2013.pdf}
\end{sourceblock}
{\small 对数正态分布。根据所研究的问题和内坐标的性质，可以选择不同的矩来进行计算。\par}
通常选择低阶矩 m1、m2 和 m3 中的两个。然而，为了容纳原始集合的更多矩，有时将最终矩计算为两个对数正态分布的矩的算术平均值。典型的选择可能是用于第一个对数正态分布 m0、m1 和 m3，以及第二个 m0、m2 和 m3。这样，校正后只有m0和m3与原来相同，但对m1和m2都实现了一定程度的控制。实现该校正算法的 Matlab 脚本可以在附录 A 的 A.2.2 节中找到。

\section*{3.2 多元分布}\phantomsection\addcontentsline{toc}{section}{3.2 多元分布}
高斯求积理论仅适用于单变量分布。然而，在实际情况中，往往需要研究具有多个内坐标的分布。在这些情况下，封闭问题通常采用以下形式：

\begin{sourceblock}
\includegraphics[page=102,pagebox=mediabox,trim={207.84bp 426.42bp 65.07bp 239.40bp},clip,max width=0.94\textwidth,max height=0.38\textheight]{Marchisio2013.pdf}
\end{sourceblock}
其中 n(xi) 是（通常未知的）NDF，g(xi) 是一个将积分中出现的所有非 NDF 项组合在一起的函数， xi = (xi1, xi2, . . . , xiM) 是（维度为 M 的）内坐标向量，方程 1 的简洁表示法为： (3.34) 表示以下多维积分：

\begin{sourceblock}
\includegraphics[page=102,pagebox=mediabox,trim={160.23bp 332.55bp 65.07bp 331.50bp},clip,max width=0.94\textwidth,max height=0.38\textheight]{Marchisio2013.pdf}
\end{sourceblock}
在积分域 Ωψ = Ωψl ⊗Ωψ2 ⊗· · · ⊗ΩψM 上。 M 个内坐标的 N 个节点的一般求积需要了解 N 个权重、wα 和长度为 M 的 N 个向量节点（或节点）， Σα = Σ1,α, Σ2,α,...。。。 , ΣM, α。 N 点、M 维求积（称为多元 QMOM）通过逼近方程 1 中报告的积分来克服封闭问题。 (3.34) 求和如下：

\begin{sourceblock}
\includegraphics[page=102,pagebox=mediabox,trim={186.89bp 223.22bp 65.07bp 440.36bp},clip,max width=0.94\textwidth,max height=0.38\textheight]{Marchisio2013.pdf}
\end{sourceblock}
正如已经提到的，等式。 (3.36) 不是高斯求积近似；它的准确度是未知的（并且很大程度上取决于构造公式的矩的选择），并且从 NDF 的矩推导它的算法也不是众所周知的（与单变量情况不同）。在详细描述三种不同的算法之前，有趣的是在这里回顾一下方程： (3.36) 对应于 NDF 的以下假设：

\begin{sourceblock}
\includegraphics[page=102,pagebox=mediabox,trim={199.84bp 101.39bp 65.07bp 565.55bp},clip,max width=0.94\textwidth,max height=0.38\textheight]{Marchisio2013.pdf}
\end{sourceblock}
{\small 其中\par}
\begin{sourceblock}
\includegraphics[page=102,pagebox=mediabox,trim={192.54bp 51.10bp 65.07bp 615.84bp},clip,max width=0.94\textwidth,max height=0.38\textheight]{Marchisio2013.pdf}
\end{sourceblock}
第一个索引β标识内部坐标，而第二索引α标识特定节点或节点。在本章中，我们将使用等式 1 中报告的常用公式来引用正交近似。 （3.36），或方程中 NDF 的相应假设。 （3.37）。提醒读者也很有用，对于这个多元 NDF，混合矩定义如下：

\begin{sourceblock}
\includegraphics[page=103,pagebox=mediabox,trim={119.75bp 543.82bp 67.56bp 121.99bp},clip,max width=0.94\textwidth,max height=0.38\textheight]{Marchisio2013.pdf}
\end{sourceblock}
其中，三种不同的符号 (mk1,k2,...,kM, m(k) 和 〈Σk1 1 Σk2 2 . . . ΣkM M 〉）将在本书中互换使用，其中 k = (k1, k2, . . . , kM) 称为指数向量。建立多元求积近似，使得对于一组特定的矩，以下等价是精确的：

\begin{sourceblock}
\includegraphics[page=103,pagebox=mediabox,trim={189.31bp 439.68bp 67.56bp 227.25bp},clip,max width=0.94\textwidth,max height=0.38\textheight]{Marchisio2013.pdf}
\end{sourceblock}
或者，换句话说，多元求积近似能够精确地再现用于构造它的矩集。在介绍三种不同的多元矩反演算法之前，需要强调的是，尽管多元求积近似背后的数学理论不像单变量求积近似那么清晰，但理想的矩反演算法应该保留 PD 和 Wheeler 算法的一些特征。首先，理想的算法应该是非迭代的。事实上，这些计算通常以有限体积（或有限元）代码实现，并在大量计算单元上重复。因此，在整个计算域中进行迭代计算是负担不起的。其次，很明显，由于 NDF 在积分区间（或支持度）中为正，因此矩反演算法必须产生正权重。第三，节点（或节点）必须保持在积分区间内。这在许多应用中至关重要，例如在内坐标代表物种浓度或颗粒体积的情况下，仅举两个常见的例子。事实上很清楚，如果内部坐标只能取正值，则必须避免代表无法实现的解的负节点，因为它们会导致不稳定性，从而确实会危及模拟代码的收敛。矩反演算法的最后一个非常重要的期望属性（PD 和算法以及 Wheeler 算法都满足）是能够精确表示 N 点分布。5 只需计算与 N 点分布相对应的矩并检查所获得的权重与节点是否准确地再现了初始分布，即可轻松验证这一点。

\subsection*{3.2.1 穷举式 QMOM}\phantomsection\addcontentsline{toc}{subsection}{3.2.1 穷举式 QMOM}
Wright 等人提出了穷举式 QMOM，即直接反演方法。 (2001b) 在粒子聚结和烧结的背景下，被称为多元 QMOM 近似。该方法基于非常简单的想法

\footnotetext[5]{N 点分布由恰好 N 个具有非零权重的不同 delta 函数组成。}
选择N(M+1)个混合矩并求解相应的N(M+1)个方程组构成的非线性系统

\begin{sourceblock}
\includegraphics[page=104,pagebox=mediabox,trim={131.79bp 575.19bp 65.07bp 91.75bp},clip,max width=0.94\textwidth,max height=0.38\textheight]{Marchisio2013.pdf}
\end{sourceblock}
用 N(M + 1) 个不同的指数向量 ki = (ki1, ki2, ..., kiM) 值获得。紧凑地指示所选矩的一种便捷方法是将不同的指数向量（按全局矩顺序 γi = |ki| = ki1 + ki2 + · · · + kiM 排序）分组到所谓的指数矩阵中。 6 例如，对于具有两个内坐标 (M = 2) 的 NDF，可以从六个混合矩构造两点求积 (N = 2)。如果选择 m0,0、m1,0、m0,1、m1,1、m2,0 和 m0,2，则矩阶矩阵变为

\begin{sourceblock}
\includegraphics[page=104,pagebox=mediabox,trim={189.13bp 415.64bp 65.07bp 205.08bp},clip,max width=0.94\textwidth,max height=0.38\textheight]{Marchisio2013.pdf}
\end{sourceblock}
然后，使用标准非线性求解器（例如 Newton–Raphson 方法或共轭梯度最小化算法）求解 N(M + 1) 方程组，这两种算法均在 Press 等人的文章中进行了描述。 （1992）。表述该问题的一种便捷方法是将其写为权重和加权节点（而不是节点）。这些量通常在权重向量中组合在一起，wT = [w1, w2, .。。 , wN−1, wN]，维度为 N，在加权节点矩阵中，YT = [YT

\begin{sourceblock}
\includegraphics[page=104,pagebox=mediabox,trim={172.70bp 323.02bp 240.46bp 361.27bp},clip,max width=0.94\textwidth,max height=0.38\textheight]{Marchisio2013.pdf}
\end{sourceblock}
N]，维度为 NM，其中 YT α = wαΣα = [wαΣ1,α, wαΣ2,α, .。。 , wαψM,α]。 N(M+1)个未知数中的N(M+1)个方程组（对应于不同的矩）的最终非线性系统可以写成如下：

\begin{sourceblock}
\includegraphics[page=104,pagebox=mediabox,trim={227.48bp 281.68bp 65.07bp 406.15bp},clip,max width=0.94\textwidth,max height=0.38\textheight]{Marchisio2013.pdf}
\end{sourceblock}
其中ZT = [wT YT]表示未知权重和加权节点的向量。通过采用牛顿-拉夫森方法，可以使用以下迭代方案来求解方程： (3.43): Zn+1 = Zn −A−1(K, Xn)F(Zn)。矩阵 A 对应于方程（1）中非线性系统的雅可比行列式（Jacobian）。 (3.43) 是指数矩阵 K 和节点 XT = [ Ψ1, Ψ2, . ] 的函数。。。 , ΣN−1, ΣN] 而已。 A 的各分量可计算为

\begin{sourceblock}
\includegraphics[page=104,pagebox=mediabox,trim={106.17bp 141.78bp 65.07bp 486.78bp},clip,max width=0.94\textwidth,max height=0.38\textheight]{Marchisio2013.pdf}
\end{sourceblock}
aij =
由于矩阵必须求逆（使用牛顿-拉夫森方法时），因此在整个计算过程中它必须是非奇异的。7
6 在整本书中，全局矩阶有时会表示为矩的上标，而指数向量的分量通常出现在下标级别，即 mγ
{\small k1,k2,...,kM。 7 有趣的是，这里要强调的是，雅可比矩阵 A 正是 DQMOM 中用于计算源项的矩阵，例如参见方程 7。 （7.118）对于一般的单变量情况和方程。 （7.128）对于特定的二元情况。\par}
对于 M = 1（单变量情况），可以轻松证明当且仅当节点不同时 A 具有满秩。如果相应的 NDF 是非简并的，则始终是后者。此外，对于单变量情况，可以选择独立的非整数矩而不影响 A 的秩。对于多变量情况，具有不同的节点并不能保证 A 对于每个不同的矩选择都是满秩的。事实上，可以证明，对于固定的 N 和 M，当所有可能的节点集 M ≥ 1 时，某些不同矩是线性相关的。因此，有必要确定一个矩集，对于给定的 M 和 N 值，对于相空间中的所有非简并点，A 始终是非奇异的。因此，在下文中，我们定义最优和可实现的矩集，并讨论确定它们的不同策略。

\paragraph*{最佳矩集的定义}
我们为给定数量的内坐标 M 定义最佳矩集，使其具有以下属性（Fox，2009b）。

1. 最佳矩集由 N(M + 1) 个不同矩组成。

2. 对于 N 个不同的非简并节点的所有可能集合，最佳矩集将产生满秩方阵 A。

3.在添加高阶矩之前，最佳矩集包括特定阶γi的所有线性独立矩，以便实现内坐标的完美对称处理。

最后一个属性确保所有低阶混合矩都包含在集合中。事实证明这非常重要，因为通常忽略混合矩会导致节点位于 M 维相空间的低维子空间上。每当所研究问题的动力学不将节点-sas 限制在低维支持上时，就应该选择一个矩集来定义不限于生成此类行为的 A。因此，建议在许多应用中应使用平等对待所有 M 内坐标的矩集。如果我们放宽第三个条件，则会获得有效（但不是最佳）的矩集。为寻找给定 M 和 N 的最佳矩集而提出的方法如下。

1. 特定阶数的不同矩取决于 M。因此，我们首先定义 A 的所有可能行，最大阶数为 2N。请注意，在此步骤中构造的矩阵 A 的行数将多于列数。

2. 第一步中生成的某些行将线性依赖于任何节点的选择。因此，我们生成一组 N 个最佳节点（定义如下），并从最低顺序开始，一次从 A 中删除一个行（如果它们是线性无关的）。当 A 是满秩且平方时，该过程终止。

3. 上一步中找到的矩选择导致 A 对于特定的节点选择具有满秩。为了使时刻选择是最优的，必须证明 A 对于所有非简并节点的选择都是满秩的。这可以通过随机生成节点并检查 A 的条件数来完成。

4. 必须进行测试模拟以确定权重是否会随着建议的矩设置而变为负值。如果观察到负权重，则矩集因非最佳而被拒绝。为了避免出现负权重，可以采用另一种策略，即根据 w2 α 而不是 wα 来定义权重向量 w。

上述过程适用于给定值 N ≥2。如果失败，则无法找到该 N 值的最佳矩集，因此必须使用下一个较大的 N 重复该过程。让我们说明双变量情况（即 M = 2）的过程。当使用具有两个节点的正交近似（即 N = 2）时，需要六个矩来求解非线性系统并计算节点和权重。在这种情况下，全局零阶矩 m0,0 和全局一阶矩 m1,0 和 m0,1 确实可以选择作为最优矩集的一部分。在全局二阶矩中，只能选择两个（因为第三个是线性相关的）。为了同等对待两个内坐标，必须选择 m2,0 和 m0,2。必须选择最后一个时刻来使自由度饱和。通过选择全局三阶矩（即 m3,0、m2,1、m1,2 和 m3,0）中的任何可用矩，对称性将丢失，这意味着两个内坐标将不会被同等对待。遵循保持对称性的选择的唯一方法是选择四阶矩 m2,2。这种特定的选择将导致有效但非最佳的矩集，因为在进入下一个全局顺序之前，必须选择前一个全局顺序的所有独立矩。该示例如表 3.3 所示。因此，很明显，对于 M = 2，N = 2 没有最佳矩集。尽管表 3.3 中报告的矩集不是最优的，但它仍然有效，这意味着它可以成功地用于构建两点二元求积近似。如果节点数量增加到三个，则需要九个矩才能求解非线性系统。在这种情况下，可以选择全局阶数小于或等于 2 的所有六个矩（即 m0,0、m1,0、m0,1、m2,0、m1,1 和 m0,2）。剩下的三个矩可以从全局三阶的四个矩中选择。无论做出什么选择，最终的矩集都不会相对于两个内坐标对称，如表 3.4 所示。同样，对于这种情况，该选择不是最佳矩集，但它是有效的。最终，当使用四个节点时

\begin{center}\small 表3.3。用于构建 N = 2 的二元求积近似 (M = 2) 的矩。在这种情况下，选择 m0,3 作为三阶矩以使自由度饱和。\end{center}
\begin{sourceblock}
\includegraphics[page=106,pagebox=mediabox,trim={58.00bp 136.29bp 58.23bp 499.24bp},clip,width=\textwidth,max height=0.72\textheight]{Marchisio2013.pdf}
\par\vspace{0.45\baselineskip}
\small 表3.4。用于构建 N = 3 的二元求积近似 (M = 2) 的矩。在这种情况下，在三阶矩中选择 m2,1、m1,2 和 m0,3 以使自由度饱和。
\end{sourceblock}
\begin{sourceblock}
\includegraphics[page=106,pagebox=mediabox,trim={58.00bp 53.16bp 58.23bp 603.11bp},clip,width=\textwidth,max height=0.72\textheight]{Marchisio2013.pdf}
\par\vspace{0.45\baselineskip}
\small 表3.5。用于构建 N = 4 的二元求积近似 (M = 2) 的最佳矩集
\end{sourceblock}
\begin{sourceblock}
\includegraphics[page=107,pagebox=mediabox,trim={58.00bp 536.82bp 58.23bp 85.16bp},clip,width=\textwidth,max height=0.72\textheight]{Marchisio2013.pdf}
\par\vspace{0.45\baselineskip}
\small 表3.6。用于构建 N = 9 的二元求积近似 (M = 2) 的最佳矩集
\end{sourceblock}
\begin{sourceblock}
\includegraphics[page=107,pagebox=mediabox,trim={58.00bp 414.98bp 58.23bp 190.26bp},clip,width=\textwidth,max height=0.72\textheight]{Marchisio2013.pdf}
\end{sourceblock}
(N = 4) 用于构建二元求积近似 (M = 2)，获得由 12 个矩构成的最佳矩集，如表 3.5 所示。可以看出，矩集合关于两个内坐标对称，并且所有低阶混合矩都包含在集合中，从而得到非奇异雅可比矩阵。当使用 N = 9 时，确定了 M = 2 时另一个由 27 个矩组成的最佳矩集，如表 3.6 中所示。对于 M = 3，也可以定义最佳且有效的矩集。当使用 8 个节点时，必须选择总共 32 个矩，其中指数矩阵的分量采用集合 \{0, 1, 2, 3\} 中的整数值。在这种情况下，最佳矩集中的前 20 个矩是全局三阶或更小的不同矩。其余 12 个矩属于全局四阶或五阶矩。表 3.7 给出了完整的最佳矩集（以指数矩阵形式表示）。有趣的是，高阶矩出现在对称组合中。例如，因为 m3,2,0 不是线性早期独立的，所以 m0,2,3（或相同指数的任何其他排列）也不是。这一事实极大地简化了寻找较大 M 值的线性独立矩的任务。对于 M = 3 和 N = 27，总共 108 个矩属于最优集。指数在集合 \{0, 1, 2, 3, 4, 5\} 中取整数值的时刻。最佳矩集中的前 56 个矩是五阶或更小的不同矩。其余 52 个矩的顺序范围从 6 到 9 (Fox, 2009b)。如前所述，矩指数的所有排列都出现在集合中，因此只需列出一个这样的矩。一旦定义了矩集（最优的或简单有效的），就可以开始搜索最终的正交。对于单变量分布，对于多变量情况，收敛确实需要良好的初始猜测。这通常是通过从混合矩（如果可用）中选择相对于每个内部坐标的前 2N -1 个纯矩的子集来计算的（例如，对于 xi1，使用 m0,0,...,0、m1,0,...,0、m2,0,...,0、...、m2N−2,0,...,0 和 m2N−1,0,...,0）。然而，有效计算初始猜测的策略仅针对二元分布（即 M = 2）而开发，并且不能轻易推广到多元分布。对细节感兴趣的读者可参考专业文献（McGraw \& Wright，2003；Rosner et al.，2003；Wright et al.，2001b；Yoon \& McGraw，2004a，2004b）。实现双变量求积算法的 Matlab 脚本可以在附录 A 的 A.3.1 节中找到。

\clearpage
\begin{center}\small 表3.7。用于构建 N = 8 的三变量求积近似 (M = 3) 的最佳矩集\end{center}
\begin{sourceblock}
\includegraphics[page=108,pagebox=mediabox,trim={58.00bp 333.51bp 58.23bp 85.16bp},clip,width=\textwidth,max height=0.72\textheight]{Marchisio2013.pdf}
\end{sourceblock}
读者可以验证本节中描述的强力 QMOM 是否能够从初始猜测足够接近最终精确解的相应时刻开始检索 N 点分布。随着最初猜测的恶化，欠松弛因子的减少可能有助于收敛，但超过某个点该方法就会失败。本节末尾报告了比较不同方法的数值示例。

\subsection*{3.2.2 张量积QMOM}\phantomsection\addcontentsline{toc}{subsection}{3.2.2 张量积QMOM}
张量积 QMOM 是由 Yoon 和 McGraw (2004a, 2004b) 在气溶胶建模背景下以及 Fox (2008, 2009a) 在稀气固流背景下引入的。所提出的算法根据用于反演的混合矩的选择和内坐标的类型（例如颗粒尺寸、体积或速度）而有所不同。在本节中，讨论张量积算法的通用版本（对于以通用内坐标向量 ε 形式编写的任何类型的分布都有效）。该方法构造 N = N1N2 阶的正交近似。。。 NM，其中M是内坐标的数量，N1，N2，.。。 , NM 是在每个 M 方向上首先构建的单变量正交近似的阶数。事实上，单变量求积的构造（从第一个 2Nβ，相对于不同的 M 内坐标，具有 1 ≤β ≤M 纯矩）是该过程的第一步。这是通过对第一个内坐标应用 PD 或 Wheeler 算法来完成的，(m0,0,...,0, m1,0,...,0, ..., m2N1−2,0,...,0, m2N1−1,0,...,0) →(w1

\begin{sourceblock}
\includegraphics[page=108,pagebox=mediabox,trim={67.52bp 53.31bp 65.06bp 621.70bp},clip,max width=0.94\textwidth,max height=0.38\textheight]{Marchisio2013.pdf}
\end{sourceblock}
\begin{sourceblock}
\includegraphics[page=109,pagebox=mediabox,trim={65.02bp 598.06bp 67.55bp 59.82bp},clip,max width=0.94\textwidth,max height=0.38\textheight]{Marchisio2013.pdf}
\end{sourceblock}
\begin{sourceblock}
\includegraphics[page=109,pagebox=mediabox,trim={97.21bp 493.71bp 98.08bp 122.81bp},clip,max width=0.94\textwidth,max height=0.38\textheight]{Marchisio2013.pdf}
\end{sourceblock}
有M个不同的向量，由N1，N2，...组成。。。，NM 个分量，狄拉克 δ 函数以单变量节点（在每个 M 方向）为中心，对应于用于近似多元 NDF 函数形式的基函数集。最终的 N 个多元节点集是通过使用以下张量积获得的： Δ =

\begin{sourceblock}
\includegraphics[page=109,pagebox=mediabox,trim={101.11bp 426.68bp 306.88bp 252.27bp},clip,max width=0.94\textwidth,max height=0.38\textheight]{Marchisio2013.pdf}
\end{sourceblock}
。这会产生一个 N = N1N2 的向量。。。 NM 个分量对应于 M 内坐标的所有可能的节点组合，即 Δ = [δ( xi1 − xi )

\begin{sourceblock}
\includegraphics[page=109,pagebox=mediabox,trim={65.03bp 378.49bp 67.55bp 284.20bp},clip,max width=0.94\textwidth,max height=0.38\textheight]{Marchisio2013.pdf}
\end{sourceblock}
\begin{sourceblock}
\includegraphics[page=109,pagebox=mediabox,trim={195.34bp 285.74bp 196.21bp 327.84bp},clip,max width=0.94\textwidth,max height=0.38\textheight]{Marchisio2013.pdf}
\end{sourceblock}
其中 N = N1N2。。。新墨西哥。如果 N 点求积的最终（未知）权重向量被定义为 w ==(w1, w2, ..., wN−1, wN)，则最终的求积近似（根据公式 (3.37) 中 NDF 函数形式的假设编写）由这两个向量的标量积（即 w·Δ）得出。通过求解N个线性方程来确定未知权重向量，确保求积近似正确地表示N个混合矩。第一个 N1 + (N2 -1) + · · · + (NM -1) 方程是通过强制与 M 个单变量求积一致来确定的，8 而其余方程可以通过固定一些混合矩来确定。不保证所产生的权重的符号，并且根据选择，某些权重可能结果为负。一般来说，其余方程是通过选择最低阶混合矩来确定的，如下面的示例所示。为了更清楚地说明该算法的潜力和局限性，我们讨论了一些例子。

练习 3.6 考虑具有两个内坐标 Ψ1 和 Ψ2 的二元分布 (M = 2)，让我们构造一个由阶数 N1 = N2 = 2 的单变量求积得出的四点求积近似。了解第一个 2N1 = 2N2 = 4（相对于第一个 Ψ1 的纯矩）和第二个 Ψ2（内坐标）就足够了

\footnotetext[8]{不言而喻，N1 + N2 + · · · + NM 方程可以通过强制与单变量求积一致来写成，但只有 N1 + (N2 −1) + · · · + (NM −1) 是线性无关的。}
初始单变量求积的计算。事实上，使用PD（或Wheeler）算法，可以很容易地进行以下计算：

\begin{sourceblock}
\includegraphics[page=110,pagebox=mediabox,trim={161.45bp 570.22bp 157.34bp 87.82bp},clip,max width=0.94\textwidth,max height=0.38\textheight]{Marchisio2013.pdf}
\end{sourceblock}
使用张量积 QMOM，最终的四点求积近似以以下双变量节点为中心：

\begin{sourceblock}
\includegraphics[page=110,pagebox=mediabox,trim={216.23bp 455.07bp 65.07bp 166.03bp},clip,max width=0.94\textwidth,max height=0.38\textheight]{Marchisio2013.pdf}
\end{sourceblock}
对应于 NDF 的以下假设：

\begin{sourceblock}
\includegraphics[page=110,pagebox=mediabox,trim={198.18bp 390.61bp 65.07bp 286.41bp},clip,max width=0.94\textwidth,max height=0.38\textheight]{Marchisio2013.pdf}
\end{sourceblock}
其中 xiα 且 α = 1，.。。 , 4 在等式中定义。 (3.45)，而 w = [w1, w2, w3, w4] 可以通过求解线性系统来计算。该线性系统是通过强制与单变量分布一致而获得的。因为两个单变量求积是在八个纯矩上构造的（四个关于 Ψ1，另外四个关于 Ψ2），并且因为 m0,0 对于两个矩集是相同的，理论上可以写出以下七个方程：

\begin{sourceblock}
\includegraphics[page=110,pagebox=mediabox,trim={88.33bp 185.86bp 65.07bp 399.43bp},clip,max width=0.94\textwidth,max height=0.38\textheight]{Marchisio2013.pdf}
\end{sourceblock}
然而，在这七个中，只有三个是线性独立的（即 N1 + N2 −1 = 3）。值得注意的是，编写这三个方程的另一种方法是强制最终权重的总和等于相应方向上的原始单变量权重。例如，权重 w1 与 w2 在同一直线上（这意味着它们都共享相同的节点 xi* 1,1），

\begin{sourceblock}
\includegraphics[page=110,pagebox=mediabox,trim={218.33bp 89.98bp 65.07bp 594.31bp},clip,max width=0.94\textwidth,max height=0.38\textheight]{Marchisio2013.pdf}
\end{sourceblock}
{\small 相应地对于 xi*\par}
\begin{sourceblock}
\includegraphics[page=110,pagebox=mediabox,trim={170.85bp 50.87bp 65.07bp 619.23bp},clip,max width=0.94\textwidth,max height=0.38\textheight]{Marchisio2013.pdf}
\par\vspace{0.45\baselineskip}
\small 表3.8。用于使用张量积 QMOM 构建 N = 4 的双变量求积近似 (M = 2) 的矩集
\end{sourceblock}
\begin{sourceblock}
\includegraphics[page=111,pagebox=mediabox,trim={58.00bp 383.71bp 58.23bp 93.12bp},clip,width=\textwidth,max height=0.72\textheight]{Marchisio2013.pdf}
\end{sourceblock}
\begin{sourceblock}
\includegraphics[page=111,pagebox=mediabox,trim={58.00bp 383.71bp 58.23bp 94.12bp},clip,width=\textwidth,max height=0.72\textheight]{Marchisio2013.pdf}
\par\vspace{0.45\baselineskip}
\small 图3.1。通过两点单变量求积 (N1 = N2 = 2) 获得的双变量张量积 QMOM (M = 2) 的四个节点在内坐标平面中的位置。
\end{sourceblock}
对于第一个内坐标，可以写出两个方程（即 N1 = 2），而对于第二个内坐标，只能写出一个方程（即 N2 −1 = 1），从而得到

\begin{sourceblock}
\includegraphics[page=111,pagebox=mediabox,trim={215.85bp 288.51bp 67.56bp 395.77bp},clip,max width=0.94\textwidth,max height=0.38\textheight]{Marchisio2013.pdf}
\end{sourceblock}
当然，另一个是前三个的线性组合。方程的推导。通过观察相空间中四个双变量节点的位置可以直观地理解（3.48）–（3.50）（见图3.1）。最后一个缺失的方程，用于求解四阶线性系统，来自于一个额外矩的执行。通常的选择是强制执行一个混合矩，产生一个附加方程，以便最终的线性系统是非奇异的。例如，如果选择 m1,1，则得到以下等式：

\begin{sourceblock}
\includegraphics[page=111,pagebox=mediabox,trim={144.26bp 186.80bp 67.56bp 499.90bp},clip,max width=0.94\textwidth,max height=0.38\textheight]{Marchisio2013.pdf}
\end{sourceblock}
由方程组成的系统。可以求解（3.48）–（3.51）来找到权重。可以看出，四点二元求积近似容纳了八个矩，即相对于内坐标的零阶矩、六个纯矩和一个混合矩。表 3.8 以矩阵形式报告了这八个矩，与表 3.5 的比较清楚地表明这是最佳矩集的子集。

由于速度分布（在三个空间维度上）需要三个内坐标 (M = 3)，因此我们在下面的示例中讨论三变量张量积 QMOM 的正交近似的构造。

{\small 练习 3.7 考虑三个内坐标 Ψ1、Ψ2 和 Ψ3 的三变量分布 (M = 3)。让我们构造一个由单变量产生的八点张量积 QMOM\par}
N1 = N2 = N3 = 2 阶的求积。与前面的示例一样，三个初始单变量求积是通过使用 PD（或 Wheeler）算法从关于三个内坐标的第一个 2N1 = 2N2 = 2N3 = 4 个纯矩构建的：

\begin{sourceblock}
\includegraphics[page=112,pagebox=mediabox,trim={150.99bp 542.22bp 146.88bp 97.25bp},clip,max width=0.94\textwidth,max height=0.38\textheight]{Marchisio2013.pdf}
\end{sourceblock}
使用张量积 QMOM，最终的八点求积近似以以下三变量节点为中心：

\begin{sourceblock}
\includegraphics[page=112,pagebox=mediabox,trim={207.37bp 358.32bp 65.07bp 188.90bp},clip,max width=0.94\textwidth,max height=0.38\textheight]{Marchisio2013.pdf}
\end{sourceblock}
权重 w 可以通过求解线性系统来计算，其中第一个 N1 + (N2 −1) + (N3 −1) = 4 方程是通过强制与单变量求积一致而获得的：

\begin{sourceblock}
\includegraphics[page=112,pagebox=mediabox,trim={197.76bp 255.46bp 65.07bp 379.68bp},clip,max width=0.94\textwidth,max height=0.38\textheight]{Marchisio2013.pdf}
\end{sourceblock}
正如练习3.6，通过观察三个内坐标相空间中八个节点的位置，可以很容易地找到这四个方程（见图3.2）。另外四个方程是通过固定一些混合矩而获得的（即，下面我们显示了固定可能的最低阶矩的最自然选择的方程：m1,1,0、m1,0,1、m0,1,1 和 m1,1,1），结果为

\begin{sourceblock}
\includegraphics[page=112,pagebox=mediabox,trim={77.48bp 149.99bp 65.07bp 521.77bp},clip,max width=0.94\textwidth,max height=0.38\textheight]{Marchisio2013.pdf}
\end{sourceblock}
\begin{sourceblock}
\includegraphics[page=112,pagebox=mediabox,trim={77.48bp 102.70bp 65.07bp 569.07bp},clip,max width=0.94\textwidth,max height=0.38\textheight]{Marchisio2013.pdf}
\end{sourceblock}
\begin{sourceblock}
\includegraphics[page=112,pagebox=mediabox,trim={77.48bp 55.40bp 65.07bp 616.36bp},clip,max width=0.94\textwidth,max height=0.38\textheight]{Marchisio2013.pdf}
\end{sourceblock}
\begin{sourceblock}
\includegraphics[page=113,pagebox=mediabox,trim={58.00bp 494.80bp 58.23bp 59.00bp},clip,width=\textwidth,max height=0.72\textheight]{Marchisio2013.pdf}
\par\vspace{0.45\baselineskip}
\small 图3.2。通过两点单变量求积 (N1 = N2 = N3 = 2) 获得的三变量张量积 QMOM（即 M = 3）的八个节点在相空间中的位置。
\end{sourceblock}
\begin{sourceblock}
\includegraphics[page=113,pagebox=mediabox,trim={74.99bp 426.55bp 67.56bp 245.20bp},clip,max width=0.94\textwidth,max height=0.38\textheight]{Marchisio2013.pdf}
\end{sourceblock}
这个八点张量积 QMOM 固定了所有内坐标的零阶矩，即 9 个纯矩加上 4 个混合矩，总共 14 个矩。此特定选择所适应的最高阶矩是全局阶 γ = 3 的 m1,1,1。将此矩集与表 3.7 中报告的 M = 3 和 N = 8 的最佳矩集进行比较，清楚地表明它是 32 个最佳矩的子集。

总之，张量积 QMOM 的矩反演算法允许我们从 M 个阶 N1、N2、... 的单变量求积开始构建（无需像强力 QMOM 那样的迭代过程）M 维、N 点求积近似。。。 , NM 在每个 M 方向上。正如已经报道的，正交节点的最终数量等于 N1N2。。。新墨西哥。矩反演算法控制（并因此采用）的矩总数等于 N1N2。。。 NM+N1+N2+····+NM。 N 个节点是通过对每个内坐标应用 PD（或 Wheeler）算法 M 次来确定的，从第一个 2Nβ 开始，相对于不同的内坐标有 1 ≤β ≤M 纯矩），然后应用张量积 QMOM。 N个权重是通过求解线性系统来计算的。 N 个线性方程是通过强制与原始单变量求积一致而获得的，结果为 N1 + (N2 -1)。。。 (NM −1) 线性独立方程。这些方程可以通过直接强制一致或采用以下程序来找到。针对第一个内坐标编写 N1 方程，然后编写 N2 -1, N3 -1,...。。。 、NM -1 写入其他剩余方向。正如前面练习中所报告的，这是通过强制沿一个内部坐标的最终权重之和等于相应方向上的原始权重（从单变量求积获得）来完成的。其他剩余方程通过固定混合矩来确定。正如已经提到的，混合矩的选择不能任意完成，因为如示例中所示，对于特定的混合矩，所得方程必须是线性独立的。一般来说，选择索引小于或等于 Nβ -1 且 1 ≤β ≤M 的混合矩可确保线性系统是非奇异的并且看起来是唯一的。理论上，实现大于本节讨论的 Nβ 值是可能的，预计不会出现任何具体问题。

为了结束本节关于张量积 QMOM 的内容，需要强调的是，尽管没有必要，但用平移（即以均值为中心）和旋转（即使用对角协方差矩阵）内坐标来表述问题可能是有利的。事实上，如果改变变量9，使得分布相对于其主坐标被重写，则求积近似的推导计算就会被简化。这些概念将在练习 3.8 中进行说明。实现张量积 QMOM 的 Matlab 脚本可以在附录 A 的 A.3.2 节中找到。

\subsection*{3.2.3 条件QMOM}\phantomsection\addcontentsline{toc}{subsection}{3.2.3 条件QMOM}
条件矩求积法 (CQMOM) 基于条件密度函数的概念（Yuan \& Fox，2011）。条件密度函数又表示当一个或多个其他内坐标固定且等于特定值时，一个内坐标处于无穷小极限内的概率。例如，在通用 NDF 的情况下，表达式

\begin{sourceblock}
\includegraphics[page=114,pagebox=mediabox,trim={138.05bp 412.00bp 65.07bp 261.47bp},clip,max width=0.94\textwidth,max height=0.38\textheight]{Marchisio2013.pdf}
\end{sourceblock}
定义了当所有其他内坐标（β = 1，...，M −1 的 ββ）均采用固定常数值时，第 M 个内坐标在 ΣM 到 ΣM + dΣM 范围内的概率。根据定义，条件密度函数积分为 1 和方程的分母。 (3.58) 通过积分第 M 个内坐标来定义：

\begin{sourceblock}
\includegraphics[page=114,pagebox=mediabox,trim={165.69bp 319.29bp 65.07bp 344.89bp},clip,max width=0.94\textwidth,max height=0.38\textheight]{Marchisio2013.pdf}
\end{sourceblock}
类似地，我们可以定义以下条件密度函数：

\begin{sourceblock}
\includegraphics[page=114,pagebox=mediabox,trim={132.85bp 265.05bp 65.07bp 408.42bp},clip,max width=0.94\textwidth,max height=0.38\textheight]{Marchisio2013.pdf}
\end{sourceblock}
定义在无穷小区间内 ΨM−1 的概率，而 Ψβ（其中 β = 1, . . ., M -2）假设为常数和特定值。通过构造基于越来越小的内坐标子集的条件密度函数，可以获得以下恒等式：

\begin{sourceblock}
\includegraphics[page=114,pagebox=mediabox,trim={118.83bp 171.46bp 65.07bp 500.64bp},clip,max width=0.94\textwidth,max height=0.38\textheight]{Marchisio2013.pdf}
\end{sourceblock}
那么很明显以下等式成立：

\begin{sourceblock}
\includegraphics[page=114,pagebox=mediabox,trim={192.03bp 112.39bp 65.07bp 551.66bp},clip,max width=0.94\textwidth,max height=0.38\textheight]{Marchisio2013.pdf}
\end{sourceblock}
这些概念可用于构造多元求积。

\footnotetext[9]{一般来说，变量的变化包括相对于平均值的平移和旋转到无应力（对角）坐标系。这种转换确保求积保持伽利略不变且可实现（即非负权重）。}
\begin{sourceblock}
\includegraphics[page=115,pagebox=mediabox,trim={58.00bp 589.80bp 58.23bp 59.00bp},clip,width=\textwidth,max height=0.72\textheight]{Marchisio2013.pdf}
\par\vspace{0.45\baselineskip}
\small 图3.3。在 CQMOM 中构建多元求积的第一步。确定第一个内部坐标 Ψ1 的 N1 个节点和 N1 个权重。
\end{sourceblock}
多元求积的构造首先通过使用带有前 2N1 -1 个矩的 Wheeler 算法计算第一个内坐标的 N1 阶单变量求积：

\begin{sourceblock}
\includegraphics[page=115,pagebox=mediabox,trim={162.96bp 437.00bp 165.50bp 189.59bp},clip,max width=0.94\textwidth,max height=0.38\textheight]{Marchisio2013.pdf}
\end{sourceblock}
N1 权重 wα1 由单个索引标识，而 N1 节点 xi1;α1 由两个索引标识，其中第一个表示内部坐标，α1 = 1,...。。。，N1。第一步如图 3.3 所示，回顾使用正交近似和重建 NDF 的等价性，可以获得以下表达式：

\begin{sourceblock}
\includegraphics[page=115,pagebox=mediabox,trim={187.18bp 339.69bp 67.56bp 326.69bp},clip,max width=0.94\textwidth,max height=0.38\textheight]{Marchisio2013.pdf}
\end{sourceblock}
通过使用方程式中报告的定义。 (3.61) 并利用第一个内坐标上的单变量求积，得到 k = [k1, k2, ... 阶的通用混合矩。。。 , kM−1, kM] 可计算为

\begin{sourceblock}
\includegraphics[page=115,pagebox=mediabox,trim={76.69bp 213.17bp 67.56bp 405.56bp},clip,max width=0.94\textwidth,max height=0.38\textheight]{Marchisio2013.pdf}
\end{sourceblock}
当 k = 0 时，求积精确地再现了所有内坐标的零阶矩：

\begin{sourceblock}
\includegraphics[page=115,pagebox=mediabox,trim={189.42bp 153.06bp 67.56bp 513.31bp},clip,max width=0.94\textwidth,max height=0.38\textheight]{Marchisio2013.pdf}
\end{sourceblock}
同样，当取相对于第一内坐标的 k1 阶混合矩和相对于第二内坐标的 k2 阶混合矩时，得到以下等式：

\begin{sourceblock}
\includegraphics[page=115,pagebox=mediabox,trim={147.18bp 84.08bp 67.56bp 578.93bp},clip,max width=0.94\textwidth,max height=0.38\textheight]{Marchisio2013.pdf}
\end{sourceblock}
{\small 因为所有其他条件密度函数都积分为一。方程（3.66）包含关于第二个内部的 k2 阶条件矩的定义\par}
坐标以第一个内坐标的值为条件（即 Σ1 = Σ1;α1），

\begin{sourceblock}
\includegraphics[page=116,pagebox=mediabox,trim={175.03bp 593.60bp 65.07bp 70.46bp},clip,max width=0.94\textwidth,max height=0.38\textheight]{Marchisio2013.pdf}
\end{sourceblock}
为简洁起见，将其表示为 〈Σk2 2 〉α1。对于每个 N1 节点 xi1;α1，前 2N2 个条件矩的知识足以通过应用 Wheeler 算法计算 N1 个不同的求积：

\begin{sourceblock}
\includegraphics[page=116,pagebox=mediabox,trim={146.28bp 470.29bp 143.82bp 146.77bp},clip,max width=0.94\textwidth,max height=0.38\textheight]{Marchisio2013.pdf}
\end{sourceblock}
其中 α1 = 1, .。。，N1。必要的条件矩可以使用等式计算： (3.66) 当 k1 = 0 时，.。。 , N1 且 k2 = 0, .。。，2N2 -1，产生以下范德蒙德形式的 2N2 -1 线性系统（等级为 N1）：

\begin{sourceblock}
\includegraphics[page=116,pagebox=mediabox,trim={175.31bp 357.02bp 65.07bp 265.92bp},clip,max width=0.94\textwidth,max height=0.38\textheight]{Marchisio2013.pdf}
\end{sourceblock}
其中系数矩阵定义为

\begin{sourceblock}
\includegraphics[page=116,pagebox=mediabox,trim={172.57bp 280.94bp 65.07bp 357.32bp},clip,max width=0.94\textwidth,max height=0.38\textheight]{Marchisio2013.pdf}
\end{sourceblock}
以及

\begin{sourceblock}
\includegraphics[page=116,pagebox=mediabox,trim={201.47bp 223.71bp 65.07bp 426.24bp},clip,max width=0.94\textwidth,max height=0.38\textheight]{Marchisio2013.pdf}
\end{sourceblock}
且 k2 = 1，.。。，2N2-1。方程的范德蒙线性系统。只要节点 xi1;α1 不同，方程（3.68）就是非奇异的，并且可以使用 Rybicki 提出的有效算法来求解（Press et al., 1992）。最终的求积（对于前两个内坐标）如图 3.4 所示，对应于 NDF 的以下假设：

\begin{sourceblock}
\includegraphics[page=116,pagebox=mediabox,trim={89.50bp 124.13bp 65.07bp 542.24bp},clip,max width=0.94\textwidth,max height=0.38\textheight]{Marchisio2013.pdf}
\end{sourceblock}
如果现在阶数为 k = [k1, k2, k3, 0, ... 的混合矩。。。 , 0] 使用等式计算。 （3.61）和等式。 (3.71) 得到以下表达式：

\begin{sourceblock}
\includegraphics[page=116,pagebox=mediabox,trim={77.48bp 55.12bp 65.07bp 607.89bp},clip,max width=0.94\textwidth,max height=0.38\textheight]{Marchisio2013.pdf}
\end{sourceblock}
\begin{sourceblock}
\includegraphics[page=117,pagebox=mediabox,trim={58.00bp 484.80bp 58.23bp 59.00bp},clip,width=\textwidth,max height=0.72\textheight]{Marchisio2013.pdf}
\par\vspace{0.45\baselineskip}
\small 图3.4。使用 CQMOM 构建多元求积的第二步。确定第二个内部坐标 ψ2 的 N1N2 节点和 N1N2 权重。
\end{sourceblock}
其中，关于第三个内坐标的 k3 阶条件矩的定义，以前两个内坐标的值为条件（即 Ψ1 = Ψ1;α1 和 Ψ2 = Ψ2;α1,α2），为

\begin{sourceblock}
\includegraphics[page=117,pagebox=mediabox,trim={91.85bp 321.25bp 67.56bp 313.83bp},clip,max width=0.94\textwidth,max height=0.38\textheight]{Marchisio2013.pdf}
\end{sourceblock}
其中

\begin{sourceblock}
\includegraphics[page=117,pagebox=mediabox,trim={168.22bp 273.45bp 67.56bp 392.93bp},clip,max width=0.94\textwidth,max height=0.38\textheight]{Marchisio2013.pdf}
\end{sourceblock}
等式。 (3.72) 变为

\begin{sourceblock}
\includegraphics[page=117,pagebox=mediabox,trim={173.33bp 210.48bp 67.56bp 455.89bp},clip,max width=0.94\textwidth,max height=0.38\textheight]{Marchisio2013.pdf}
\end{sourceblock}
对于 k1 = 0，.。。 , N1 -1 生成以下线性系统：

\begin{sourceblock}
\includegraphics[page=117,pagebox=mediabox,trim={180.70bp 100.07bp 67.56bp 519.32bp},clip,max width=0.94\textwidth,max height=0.38\textheight]{Marchisio2013.pdf}
\end{sourceblock}
{\small 对于每对 k2 = 0, .。。 , N2 −1 且 k3 = 1, .。。 , 2N3 -1，其中 N3 是所采用的正交节点的数量。请注意，这些线性系统具有与方程相同的系数矩阵。 （3.68），因此我们保证找到 ζk2,k3 α1 的解。\par}
对于 α1 的每个值，最后一步是求解由式（1）生成的线性系统。 （3.74），再次采用范德蒙德形式：

\begin{sourceblock}
\includegraphics[page=118,pagebox=mediabox,trim={175.84bp 514.11bp 65.07bp 84.58bp},clip,max width=0.94\textwidth,max height=0.38\textheight]{Marchisio2013.pdf}
\end{sourceblock}
其中系数矩阵定义为

\begin{sourceblock}
\includegraphics[page=118,pagebox=mediabox,trim={164.24bp 425.70bp 65.07bp 202.20bp},clip,max width=0.94\textwidth,max height=0.38\textheight]{Marchisio2013.pdf}
\end{sourceblock}
\begin{sourceblock}
\includegraphics[page=118,pagebox=mediabox,trim={188.17bp 368.12bp 65.07bp 281.83bp},clip,max width=0.94\textwidth,max height=0.38\textheight]{Marchisio2013.pdf}
\end{sourceblock}
以及

总之，对于每对 k2 = 0，.。。 , N2 −1 且 k3 = 1, .。。 , 2N3 −1，通过求解（使用 Rybicki 算法）方程 1 中报告的线性系统来计算量 zek2,k3 α1。 （3.76）。随后，这些量用于求解 α1 = 1 的每个值，...。。。，N1 方程中报告的线性系统。 （3.77）。最后，通过将 Wheeler 算法应用于相对于第三内坐标的前 2N3 个条件矩来计算第三内坐标的正交节点：

\begin{sourceblock}
\includegraphics[page=118,pagebox=mediabox,trim={147.09bp 203.37bp 144.63bp 402.78bp},clip,max width=0.94\textwidth,max height=0.38\textheight]{Marchisio2013.pdf}
\end{sourceblock}
请注意，对于第三个内部坐标，求积节点是根据图 3.4 中的每个 N1N2 节点沿着垂直于相空间中由 Ψ1 和 Ψ2 生成的平面的垂直线计算的，从而得到 N1N2N3 节点和 N1N2N3 权重。这对应于 NDF 的以下功能假设：

\begin{sourceblock}
\includegraphics[page=118,pagebox=mediabox,trim={69.69bp 84.30bp 65.07bp 555.11bp},clip,max width=0.94\textwidth,max height=0.38\textheight]{Marchisio2013.pdf}
\end{sourceblock}
{\small 按照相同的步骤，可以计算其他内坐标的正交节点。\par}
\begin{center}\small 表3.9。用于构建 N1 = N2 = 3 的双变量 CQMOM 正交近似 (M = 2) 的矩，其中 ψ2 以 ψ1 为条件\end{center}
\begin{sourceblock}
\includegraphics[page=119,pagebox=mediabox,trim={58.00bp 497.33bp 58.23bp 94.61bp},clip,width=\textwidth,max height=0.72\textheight]{Marchisio2013.pdf}
\par\vspace{0.45\baselineskip}
\small 表3.10。用于构建 N1 = N2 = 3 的双变量 CQMOM 求积近似 (M = 2) 的矩，其中 xi1 以 ψ2 为条件
\end{sourceblock}
\begin{sourceblock}
\includegraphics[page=119,pagebox=mediabox,trim={58.00bp 367.15bp 58.23bp 231.60bp},clip,width=\textwidth,max height=0.72\textheight]{Marchisio2013.pdf}
\end{sourceblock}
作为一个例子，考虑一个简单的双变量情况 (M = 2)，让我们确定用于构建正交近似的矩。如果 N1 取等于 3，则将使用相对于 xi1 的前 6 个矩来构建初始单变量求积近似，即 m(0, 0)、m(1, 0)、m(2, 0)、m(3, 0)、m(4, 0) 和 m(5, 0)。如果我们还取 N2 等于 3，则必须为 Ψ1 上的三个节点中的每一个建立 Ψ2 上的正交近似。用于计算条件矩并确定三个单变量求积近似的矩为 m(0, 1)、m(1, 1)、m(2, 1)、m(0, 2)、m(1, 2)、m(2, 2)、m(0, 3)、m(1, 3)、m(2, 3)、m(0, 4)、m(1, 4)、 m(2, 4)、m(0, 5)、m(1, 5)，最后是 m(2, 5)。所使用的矩总数（对于这种双变量情况）等于 21，如表 3.9 所示。表 3.9 与表 3.6 的比较表明，这 21 个时刻是 27 个最佳时刻的子集。很明显，如果该过程被反转并且第一个内部坐标以第二个内部坐标为条件，则必须使用表 3.10 中报告的 21 个矩。通过比较表 3.9 和表 3.10 可以清楚地看出，两个集合并集得到的矩集对应于 N = 9 时的最优矩集。通过考虑两种可能的方案，控制的最高阶矩为全局阶 7（即 m(2, 5) 或 m(5, 2)），而控制的最高纯矩为全局阶 5 的 m(5, 0) 和 m(5, 0)。双变量情况下，N1和N2现在取等于4，可以控制更多的矩。为了构造关于 xi1 的第一个单变量求积，使用前 2N1 个纯矩，即 m(0, 0)、m(1, 0)、m(2, 0)、m(3, 0)、m(4, 0)、m(5, 0)、m(6, 0) 和 m(7, 0)。然后，使用以下矩计算条件矩和所得的 N1 个单变量求积（N2 阶）：m(0, 1)、m(0, 2)、m(0, 3)、m(0, 4)、m(0, 5)、m(0, 6)、m(0, 7) 和 m(1, 1)、m(1, 2)、m(1, 3)、m(1, 4)、m(1, 5)、m(1, 6)、m(1, 7) 和 m(2, 1)、m(2, 2)、m(2, 3)、m(2, 4)、m(2, 5)、m(2, 6)、m(2, 7)，最后是 m(3, 1)、m(3, 2)、m(3, 3)、m(3, 4)、m(3, 5)、m(3, 6)、m(3, 7)。本例中使用的矩如表 3.11 所示。在本例中，控制了 36 个矩，包括

\begin{center}\small 表3.11。用于构建 N1 = N2 = 4 的双变量 CQMOM 正交近似 (M = 2) 的矩，其中 ψ2 以 ψ1 为条件\end{center}
\begin{sourceblock}
\includegraphics[page=120,pagebox=mediabox,trim={58.00bp 491.40bp 58.23bp 86.65bp},clip,width=\textwidth,max height=0.72\textheight]{Marchisio2013.pdf}
\end{sourceblock}
十阶以下的混合矩和七阶以下的纯矩。很容易证明，该矩集与通过使用 CQMOM（通过将 Ψ1 调节到 Ψ2）获得的矩集相结合，产生了 N = 16 的 48 个最佳矩集。由读者来验证这对于三变量分布是否也有效。例如，当 N1 = N2 = N3 = 4 时，CQMOM（具有一种排列）控制 256 个最佳矩中的 168 个，最高可达 13 阶，而 64 节点张量积 QMOM 仅控制 76 个矩。将 CQMOM 与之前介绍的方法进行比较很有趣。首先，CQMOM 的一个重要属性是它永远不会有负权重，因为权重是使用单变量求积计算的。因此，只要单变量求积中使用的矩集有效，10 CQMOM 将为多元矩集提供可实现的求积。正如我们之前看到的，张量积 QMOM 无法保证此属性。其次，CQMOM 比张量积 QMOM 控制更大的矩集，几乎与暴力 QMOM 一样多。第三，与暴力 QMOM 不同，CQMOM 提供直接（非迭代）解决方法。实现 CQMOM 算法的 Matlab 脚本可以在附录 A 的 A.3.3 节中找到。

练习 3.8 确定以下二元高斯分布的 4 阶和 9 阶的两个正交近似：

\begin{sourceblock}
\includegraphics[page=120,pagebox=mediabox,trim={70.26bp 203.04bp 66.32bp 451.82bp},clip,max width=0.94\textwidth,max height=0.38\textheight]{Marchisio2013.pdf}
\end{sourceblock}
μ1 = 10、μ2 = 20、σ1 = σ2 = 2 和不同的 ρ 值，通过使用附录 A 中报告的算法 BF、TP 和 CQMOM。随着 ρ 从 0 增加到 1，分布被拉伸，导致相同的纯矩，但不同的混合矩。当 ρ = 0 和 ρ = 0.5 时，评估三种算法的性能和结果很有趣。图 3.5 中报告了 N = 4（左）和 N = 9（右）以及 ρ = 0（上）和 ρ = 0.5（下）的结果。正如预期的那样，当 ρ = 0 时，这三种方法返回完全相同的值。这是因为这三种方法使用相同的纯矩，仅在所使用的混合矩的数量和类型上有所不同。此外，应该强调的是，当使用强力 QMOM 时，迭代过程会收敛

\footnotetext[10]{如前所述，有必要确定一组特定的单变量矩是否可实现，以便应用求积算法。如果矩集对应于某个非负分布函数，则它是可实现的。}
\begin{sourceblock}
\includegraphics[page=121,pagebox=mediabox,trim={58.00bp 367.17bp 58.23bp 59.00bp},clip,width=\textwidth,max height=0.72\textheight]{Marchisio2013.pdf}
\par\vspace{0.45\baselineskip}
\small 图3.5。对于 N = 4（左）和 N = 9（右），μ1 = 10、μ2 = 20、σ1 = σ2 = 2、ρ = 0（上）和 ρ = 0.5（下）的双变量高斯分布的求积近似，即强力 QMOM（菱形）、张量积 QMOM（圆形）和 CQMOM（方形）。
\end{sourceblock}
仅当初始猜测足够接近最终解决方案时。根据初始猜测的距离和欠松弛因子的值，需要多达 1000 次迭代才能计算最终的正交。相反，正如已经报道的那样，张量积 QMOM 和 CQMOM 通过非迭代计算提供最终解决方案。当 ρ = 0.5 时，不同算法的结果不同。这是因为，在改变 ρ 时，混合矩会发生变化，并且由于不同的算法使用相同的纯矩但不同的混合矩，因此每种方法都会产生不同的最终求积。可以看出，ρ = 0.5 时张量积 QMOM 计算的求积近似的节点在相空间中的位置与 ρ = 0 时的位置相同。这显然是因为该方法仅根据纯矩计算节点，当 ρ 增加时纯矩不会变化。当然，对于 ρ = 0 和 ρ = 0.5，最终的权重是不同的，并且当 ρ 增大时，左上节点和右下节点的权重会减小。很容易验证，当 ρ = 0.6 时，这些权重变为负值！正如预期的那样，基于中心旋转矩的求积计算在某些情况下确实有帮助。当 CQMOM 用于 ρ = 0.5 时，最终节点会自动在内坐标空间中移动，而无需旋转参考系。值得强调的是，由于此处的计算是通过以第一个内部坐标为条件进行的，因此当 ρ 变化时，节点仅针对第二个内部坐标进行移动。如果改为使用以第一内部坐标为条件的条件矩，则将检测到另一个方向的移位。这里还需要记住的是，对于 CQMOM（与张量积 QMOM 一样），仅使用最佳矩集的子集。使用穷举式 QMOM 计算的求积可以更有效地适应新的分布形状，因为只有使用穷举式 QMOM 才能采用最佳矩集的所有矩。然而，所需的迭代次数显著增加，并且欠松弛因子必须限制为远小于单位的值以避免发散。在此分布和其他（双变量和三变量）分布上进行的测试反演表明，CQMOM 是三种算法中更稳健、更高效的。

我们上面对 CQMOM 的讨论假设多元矩集对应于连续分布。此外，我们隐含地假设每一步生成的条件矩集都是可实现的。在实践中，这些假设不成立的情况很常见，因此有必要修改CQMOM算法来处理这种情况。正如 Yuan \& Fox (2011) 所示，对 CQMOM 算法的修改实现起来相当简单，并且依赖于第 3.1.3 节末尾提到的自适应 Wheeler 算法。基本思想是，对于每组条件矩，自适应惠勒算法确定可以使用的最大节点数。由于 CQMOM 求积是由多个一维求积构建的，因此对后者使用自适应 Wheeler 算法是确保 CQMOM 始终返回明确定义的多元求积的一种优雅方法。更多细节和例子可以在 Yuan \& Fox (2011) 中找到。

\section*{3.3 扩展求积矩方法（EQMOM）}\phantomsection\addcontentsline{toc}{section}{3.3 扩展求积矩方法（EQMOM）}
3.1.1 节中介绍的高斯求积算法相当于通过狄拉克 delta 函数之和来近似单变量 NDF：

\begin{sourceblock}
\includegraphics[page=122,pagebox=mediabox,trim={200.64bp 245.48bp 65.07bp 421.46bp},clip,max width=0.94\textwidth,max height=0.38\textheight]{Marchisio2013.pdf}
\end{sourceblock}
该正交近似的阶数由 N 决定，将值从 N 增加到 N + 1 需要两个额外的矩（即 QMOM 始终使用偶数个矩）。对于许多应用（Yuan et al., 2012），方程中的形式。 (3.81) 可以通过引入具有由参数 σ 确定的有限（或无限）支持的核密度函数来推广：

\begin{sourceblock}
\includegraphics[page=122,pagebox=mediabox,trim={201.33bp 141.67bp 65.07bp 525.27bp},clip,max width=0.94\textwidth,max height=0.38\textheight]{Marchisio2013.pdf}
\end{sourceblock}
σ 的值是通过固定一个附加矩（总共 2N + 1 个矩，即奇数个矩）来确定的。为了区分使用方程的矩方法。 （3.81）和那些使用方程。式（3.82）中，我们将前者称为矩的求积矩（QMOM），后者称为矩的扩展求积法（EQMOM）（Yuan等，2012）。使用 EQMOM 代替 QMOM 的主要优点是，通过一个额外的矩，可以重建一个平滑的非负 NDF，精确地再现前 2N +1 个矩。然而，使用 EQMOM 和 Eq 时有几个实际考虑因素。 （3.82），我们将在下面讨论。函数 δσ(x, y) 最重要的属性是，它相对于 x 的整数矩导致表达式很容易反转，以找到构建涉及 NDF 的积分的求积近似所需的一组 N 个权重 wα、N 个节点 xiα 和 σ。在极限 σ = 0 的情况下，方程： （3.82）应该平滑地减少到方程。 （3.81）（基于熵最小化的重建方法不是这种情况！）。对于 σ 的所有值，关于方程 Σ 的整数矩。 （3.82）应该很容易评估，以促进 EQMOM 中使用的矩反演算法的应用（Yuan et al., 2012）。此外，δσ(x, y) 的支持度不应超过真实 NDF 的非零值范围（例如，如果 x 始终为非负，则当 x < 0 时，δσ(x, y) 应为空）。当适用时，通常可以从与正交多项式族相对应的权重函数中找到 δσ 的绝佳选择。

\subsection*{3.3.1 与正交多项式的关系}\phantomsection\addcontentsline{toc}{subsection}{3.3.1 与正交多项式的关系}
单变量 EQMOM 的一个方便选择是根据已知正交多项式族的权重函数 w(t) 来定义 δσ (x, y) (Gautschi, 2004)。例如，在区间 [0, ∞) 上，关联的拉盖尔多项式具有权重函数

\begin{sourceblock}
\includegraphics[page=123,pagebox=mediabox,trim={219.56bp 406.57bp 67.56bp 279.53bp},clip,max width=0.94\textwidth,max height=0.38\textheight]{Marchisio2013.pdf}
\end{sourceblock}
其中参数 α > −1 定义了拉盖尔多项式的特定族。11 然后，可以通过使用

\begin{sourceblock}
\includegraphics[page=123,pagebox=mediabox,trim={172.98bp 345.82bp 67.56bp 323.52bp},clip,max width=0.94\textwidth,max height=0.38\textheight]{Marchisio2013.pdf}
\end{sourceblock}
其中右侧是等式的缩放形式。 (3.83) 具有以下性质

\begin{sourceblock}
\includegraphics[page=123,pagebox=mediabox,trim={88.47bp 298.14bp 67.56bp 368.06bp},clip,max width=0.94\textwidth,max height=0.38\textheight]{Marchisio2013.pdf}
\end{sourceblock}
使用方程式找到的 NDF 示例(3.84) 和 N = 2 的四个 σ 值如图 3.6 所示。使用已知的正交多项式族的主要优点是方程中的递归系数 aα 和 bα。 (3.5) 是预先已知的（Gautschi，2004）（即我们不需要从时刻计算它们）。因此，方程中的雅可比矩阵。对于较大的 N 值，可以非常精确地计算式（3.13）（当根据矩计算递归系数时，通常情况并非如此）。雅可比矩阵已知这一事实可用于通过应用高斯求积来非常准确地评估涉及 δσ (x, y) 的积分。例如，等式。 (3.84) 导致

\begin{sourceblock}
\includegraphics[page=123,pagebox=mediabox,trim={83.98bp 161.32bp 67.56bp 502.26bp},clip,max width=0.94\textwidth,max height=0.38\textheight]{Marchisio2013.pdf}
\end{sourceblock}
其中 kα = ψα/σ 和 \{wαβ, tαβ\} 是 N′ 权重和 N′ 节点，根据给定 kα 值的拉盖尔多项式递归系数计算。12 这里，重要的是要认识到方程 1 中的 N′。 (3.86) 的选择可以独立于 N 的值
11 相关拉盖尔多项式 P(α) n (t) 的零点是区间 [0, ∞) 中的 n 个节点 tα (Gautschi, 2004)。通过定义 t = x −a，可以将区间移动到 [a, ∞)，因此后者的节点将为 xα = tα + a。

\footnotetext[12]{请注意，每个 α 值都使用一组不同的拉盖尔多项式，参数化为 kα > 0。因此，必须针对每个 α 值单独计算权重与节点。与任何高斯求积一样，方程。当 g(xi) 是小于 2N′ 阶的多项式时，(3.86) 是精确的（Gautschi，2004）。}
\begin{sourceblock}
\includegraphics[page=124,pagebox=mediabox,trim={58.00bp 473.17bp 58.23bp 60.77bp},clip,width=\textwidth,max height=0.72\textheight]{Marchisio2013.pdf}
\par\vspace{0.45\baselineskip}
\small 图3.6。用方程找到的 n(xi) 的例子。 (3.84) 对于四个 σ 值。其他参数为 w1 = w2 = 0.5、 Σ1 = 0.5 和 Σ2 = 1.5。
\end{sourceblock}
用在等式中。 （3.82）。换句话说，用于评估方程的高斯求积的阶数。 （3.86）（即2N′）与EQMOM中使用的传输矩的数量（即2N + 1）无关，我们可以选择N′≫N来提高方程（3.86）的精度。 （3.86）。最后，使用等式中的结果。 (3.86)，我们可以看到由式(3.86)定义的EQMOM闭包。 (3.84) 使用对偶求积表示法对 NDF 进行近似积分：

\begin{sourceblock}
\includegraphics[page=124,pagebox=mediabox,trim={159.56bp 325.02bp 65.07bp 338.56bp},clip,max width=0.94\textwidth,max height=0.38\textheight]{Marchisio2013.pdf}
\end{sourceblock}
其中 N 个权重 wα、N 个节点 xiα 和 σ 是从 n(xi) 的前 2N + 1 个整数矩中找到的。正如我们将在下面看到的，使用 EQMOM 时的主要计算挑战是寻找 σ 的算法的定义。一般方法是找到前 2N + 1 个整数矩的解析表达式，并使用它们定义迭代过程来查找 σ。请注意，稳健的 EQMOM 算法还必须能够处理退化情况，例如 σ = 0 和/或当 2N + 1 个整数矩可以由少于 N 个正交节点捕获时。

\subsection*{3.3.2 单变量EQMOM}\phantomsection\addcontentsline{toc}{subsection}{3.3.2 单变量EQMOM}
在本节中，我们讨论三种类型的单变量 EQMOM，它们基于具有无限、半无限和有限支持的权重函数（Chalons 等，2010；Yuan 等，2012）。附录 A 的 A.4 节提供了示例数值算法。

\paragraph*{高斯分布}
作为单变量 EQMOM 的示例，我们将考虑具有 xi ε(−∞, +∞) 的 NDF 并使用高斯分布定义核密度函数（Chalons 等人，2010）：

\begin{sourceblock}
\includegraphics[page=124,pagebox=mediabox,trim={176.65bp 61.00bp 65.07bp 610.29bp},clip,max width=0.94\textwidth,max height=0.38\textheight]{Marchisio2013.pdf}
\end{sourceblock}
\begin{sourceblock}
\includegraphics[page=125,pagebox=mediabox,trim={58.00bp 484.17bp 58.23bp 59.00bp},clip,width=\textwidth,max height=0.72\textheight]{Marchisio2013.pdf}
\par\vspace{0.45\baselineskip}
\small 图3.7。用方程找到的 n(xi) 的例子。 (3.88) 对于四个 σ 值。其他参数为 w1 = w2 = 0.5、 Σ1 = 0.5 和 Σ2 = 1.5。
\end{sourceblock}
使用方程式找到的 NDF 示例(3.88) 和 N = 2 的四个 σ 值如图 3.7 所示。请注意，当 y = 0 时，方程的右侧为(3.88) 对应于 Hermite 多项式的（缩放和移位）权重函数。13 通过选择 δσ (x, y)，可以直接评估等式 (3.88) 中的积分。 （3.87）使用 Hermite 多项式的已知递归系数（Gautschi，2004）。整数矩以方程 y 为中心。 （3.88）很容易使用它们的定义来计算：

\begin{sourceblock}
\includegraphics[page=125,pagebox=mediabox,trim={129.20bp 300.23bp 67.56bp 331.33bp},clip,max width=0.94\textwidth,max height=0.38\textheight]{Marchisio2013.pdf}
\end{sourceblock}
其中，如果 k 为偶数，则 [k/2] = k/2；如果 k 为奇数，则 [k/2] = (k −1)/2。因此，以 y = xiα 为中心的前五个整数矩为

\begin{sourceblock}
\includegraphics[page=125,pagebox=mediabox,trim={163.74bp 124.77bp 67.56bp 432.91bp},clip,max width=0.94\textwidth,max height=0.38\textheight]{Marchisio2013.pdf}
\end{sourceblock}
这组五个矩用于 N = 2 的双高斯 EQMOM（Chalons 等人，2010）。

\footnotetext[13]{为 t ε(−∞, ∞) 定义的 Hermite 多项式 P(α) n (t) 具有权重函数 w(t) = |t|2αe−t2 (Gautschi, 2004)。通过平移和缩放，它们可以写成 x = √}
{\small 2σt + y。因此，如果 tα 是 Hermite 多项式 P(0) 的 n 个根\par}
\begin{sourceblock}
\includegraphics[page=125,pagebox=mediabox,trim={162.03bp 60.42bp 323.59bp 629.77bp},clip,max width=0.94\textwidth,max height=0.38\textheight]{Marchisio2013.pdf}
\end{sourceblock}
{\small 2σtα + y 是方程式的 n 节点。 （3.88）。\par}
{\small n (t), then xα =\par}
从方程。 (3.89)，方程的整数矩。 (3.82) 可以表示为

\begin{sourceblock}
\includegraphics[page=126,pagebox=mediabox,trim={170.89bp 584.95bp 65.07bp 78.13bp},clip,max width=0.94\textwidth,max height=0.38\textheight]{Marchisio2013.pdf}
\end{sourceblock}
其中

\begin{sourceblock}
\includegraphics[page=126,pagebox=mediabox,trim={217.86bp 535.97bp 65.07bp 130.96bp},clip,max width=0.94\textwidth,max height=0.38\textheight]{Marchisio2013.pdf}
\end{sourceblock}
是 QMOM 矩（即当 σ = 0 时）。请注意等式。 (3.91) 形成一个下三角线性系统，可以使用前向替换进行反转：

\begin{sourceblock}
\includegraphics[page=126,pagebox=mediabox,trim={140.40bp 307.41bp 65.07bp 198.97bp},clip,max width=0.94\textwidth,max height=0.38\textheight]{Marchisio2013.pdf}
\end{sourceblock}
以及

\begin{sourceblock}
\includegraphics[page=126,pagebox=mediabox,trim={159.33bp 258.47bp 65.07bp 404.61bp},clip,max width=0.94\textwidth,max height=0.38\textheight]{Marchisio2013.pdf}
\end{sourceblock}
为了求解这个 2N + 1 方程组以找到 N 个权重、N 个节点和 σ2，可以使用以下迭代方法。

1. 从初始猜测 σ2 = 0 开始，使用 Wheeler 算法和 2N 个矩 \{m* 0, m* 1,...。。。 , m* 2N−1\} 由方程计算得出。 (3.93) 求 N 个权重 wα 和 N 个节点 xiα。

2. 使用权重与节点（对于 σ2 的当前值），计算矩 m* 2 N

\begin{sourceblock}
\includegraphics[page=126,pagebox=mediabox,trim={223.91bp 122.63bp 196.55bp 544.31bp},clip,max width=0.94\textwidth,max height=0.38\textheight]{Marchisio2013.pdf}
\end{sourceblock}
3. 使用非线性方程求解器（例如有界割线法）找到最小（实数）σ2 ≥0，使得目标函数

\begin{sourceblock}
\includegraphics[page=126,pagebox=mediabox,trim={172.10bp 50.16bp 143.75bp 613.32bp},clip,max width=0.94\textwidth,max height=0.38\textheight]{Marchisio2013.pdf}
\end{sourceblock}
为空。请注意，每次改变 σ2 时，2N 个矩 \{m* 0, m* 1,...。。 , m* 2N−1\} 必须根据方程重新计算。 （3.93），必须重新应用Wheeler算法来找到权重与节点（即步骤1），并且必须重新评估矩m* 2 N（即步骤2）。

应用该算法时，矩集\{m0, m1,...。。 , m2N\} 必须已知（且可实现）。根据目标函数的定义和高斯求积的性质，14 我们有 f(0) ≥0。因此，作为有界搜索算法的第一步，可以确定上限 σ+，使得在矩集 \{m* 0, m* 1, … 的条件下 f(σ+) < 0。。。 , m* 2N−1\} 从方程中找到。 (3.93) 使用 σ2 = σ+ 是可以实现的。如果不存在这样的 σ+，则可以选择 σ2 使其最小化 |f(σ2)|和矩集 \{m* 0, m* 1, .。。 , m* 2N−1\} 是可实现的。对于 N = 2 的情况（即双高斯 NDF (Chalons et al., 2010)），可以从方程 (2) 中找到三阶多项式。 (3.94) 具有确定 σ2 的唯一实根。方程式中设定的时刻。 N = 2 时的 (3.93) 和 (3.94) 如下：

\begin{sourceblock}
\includegraphics[page=127,pagebox=mediabox,trim={189.15bp 374.41bp 67.56bp 245.28bp},clip,max width=0.94\textwidth,max height=0.38\textheight]{Marchisio2013.pdf}
\end{sourceblock}
假设五个时刻\{m0,.。。 , m4\} 是可实现的，这个由五个方程组组成的系统可以针对五个未知数 \{w1, w2, Ψ1, Ψ2, σ2\} 进行解析求解。对于二节点求积，权重与节点取决于 \{m*

\begin{sourceblock}
\includegraphics[page=127,pagebox=mediabox,trim={136.29bp 263.78bp 67.56bp 355.60bp},clip,max width=0.94\textwidth,max height=0.38\textheight]{Marchisio2013.pdf}
\end{sourceblock}
\begin{sourceblock}
\includegraphics[page=127,pagebox=mediabox,trim={65.02bp 193.64bp 67.56bp 444.48bp},clip,max width=0.94\textwidth,max height=0.38\textheight]{Marchisio2013.pdf}
\end{sourceblock}
其中

\begin{sourceblock}
\includegraphics[page=127,pagebox=mediabox,trim={163.55bp 151.25bp 67.56bp 520.61bp},clip,max width=0.94\textwidth,max height=0.38\textheight]{Marchisio2013.pdf}
\end{sourceblock}
请注意，权重与节点是 σ2 的隐函数，因为它们依赖于 m* 2 和 m* 3。然后可以使用 (3.96) 来计算 m* 4：

\begin{sourceblock}
\includegraphics[page=127,pagebox=mediabox,trim={206.94bp 81.71bp 67.56bp 601.98bp},clip,max width=0.94\textwidth,max height=0.38\textheight]{Marchisio2013.pdf}
\end{sourceblock}
\footnotetext[14]{根据前 2N 个矩计算的高斯求积会产生 m* 2N 的最小可实现值（Dette \& Studden，1997）。}
这导致必须求解 σ2 的以下表达式：

\begin{sourceblock}
\includegraphics[page=128,pagebox=mediabox,trim={163.89bp 601.36bp 65.08bp 79.93bp},clip,max width=0.94\textwidth,max height=0.38\textheight]{Marchisio2013.pdf}
\end{sourceblock}
经过一些繁琐的代数（Chalons et al., 2010），方程。 (3.100) 可以重写为三阶多项式：

\begin{sourceblock}
\includegraphics[page=128,pagebox=mediabox,trim={67.52bp 478.02bp 65.08bp 124.38bp},clip,max width=0.94\textwidth,max height=0.38\textheight]{Marchisio2013.pdf}
\end{sourceblock}
其中

\begin{sourceblock}
\includegraphics[page=128,pagebox=mediabox,trim={176.15bp 424.67bp 65.08bp 246.05bp},clip,max width=0.94\textwidth,max height=0.38\textheight]{Marchisio2013.pdf}
\end{sourceblock}
在 η† > 1+(q†)2 的条件下，若 q† = 0，且 η† ≤3，则式 (1) 中的多项式为： (3.101) 在区间 0 ≤ s ≤ 1 内有唯一的实根 15：

\begin{sourceblock}
\includegraphics[page=128,pagebox=mediabox,trim={222.98bp 368.33bp 65.08bp 311.89bp},clip,max width=0.94\textwidth,max height=0.38\textheight]{Marchisio2013.pdf}
\end{sourceblock}
其中

\begin{sourceblock}
\includegraphics[page=128,pagebox=mediabox,trim={124.73bp 322.67bp 122.28bp 351.00bp},clip,max width=0.94\textwidth,max height=0.38\textheight]{Marchisio2013.pdf}
\end{sourceblock}
σ2 的期望值与该根的关系为

\begin{sourceblock}
\includegraphics[page=128,pagebox=mediabox,trim={196.82bp 264.72bp 65.08bp 409.01bp},clip,max width=0.94\textwidth,max height=0.38\textheight]{Marchisio2013.pdf}
\end{sourceblock}
对于选定的 q† 值，σ2 对 η† 的依赖性示例如图 3.8 所示。请注意，多项式的系数仅取决于已知矩 mk，而不取决于权重与节点。16 因此，可以首先计算 σ2，然后使用方程 1 计算矩 m* k。 （3.95），使用惠勒算法从中找到权重与节点。在 s = 0 的极限条件下（即 η† = 3 且 q† = 0），所有矩 mk 均为高斯矩，并且只需要一个正交节点（即 Ψ1 = Ψ2 = 〈Ψ〉）。在极限 s = 1（即 η† = 1 + (q†)2）下，矩集恰好对应于两个加权狄拉克 δ 函数。

15 条件 η† ≥1 + (q†)2 将由任何可实现的矩集满足，对于两个加权狄拉克 δ 函数的情况，其等式成立。当 q† = 0 时，NDF 关于其平均值对称，并且 η† = 3 是使用 δσ 的高斯分布可以实现的 η† 的最大可能值。如果 η† 的真实值大于该值，则目标函数 f(σ2) 永远不会为零（即多项式 P(s) 不成立），但 f(σ2) 的最小值出现在 σ2 = 〈ψ′2〉 处。可以在式（1）中找到 P(s) 的三个根的解析表达式。 (3.101)。在高斯极限中，在 s = 0 处存在三重根，否则其中两个根是复数，并且所需的根是实数。

\footnotetext[16]{据作者所知，目前尚不清楚是否可以找到 N > 2 的类似多项式，从而可以直接求解 σ2。尽管如此，上述用于确定 σ2 的迭代算法已被发现是稳健的。}
\begin{sourceblock}
\includegraphics[page=129,pagebox=mediabox,trim={58.00bp 465.41bp 58.23bp 60.42bp},clip,width=\textwidth,max height=0.72\textheight]{Marchisio2013.pdf}
\par\vspace{0.45\baselineskip}
\small 图3.8。对于 (q†)2 的选定值，参数 σ2/〈ψ′2〉 出现在双高斯分布中，与 η† 相对。
\end{sourceblock}
\paragraph*{伽玛分布}
上述针对高斯分布的矩反演算法可应用于 δσ 的其他选择（例如，在半无限或有限区间上定义的 NDF）。第一步是解析计算 δσ(x, y) 相对于 x 的整数矩。例如，使用等式中的定义。 (3.84) 对于 x ∈[0, ∞) (Yuan et al., 2012)，这些 y = xiα 的整数矩是

\begin{sourceblock}
\includegraphics[page=129,pagebox=mediabox,trim={196.40bp 298.22bp 67.57bp 368.19bp},clip,max width=0.94\textwidth,max height=0.38\textheight]{Marchisio2013.pdf}
\end{sourceblock}
因此

\begin{sourceblock}
\includegraphics[page=129,pagebox=mediabox,trim={150.67bp 241.60bp 67.57bp 425.34bp},clip,max width=0.94\textwidth,max height=0.38\textheight]{Marchisio2013.pdf}
\end{sourceblock}
其中

\begin{sourceblock}
\includegraphics[page=129,pagebox=mediabox,trim={166.37bp 188.92bp 67.57bp 472.12bp},clip,max width=0.94\textwidth,max height=0.38\textheight]{Marchisio2013.pdf}
\end{sourceblock}
\begin{sourceblock}
\includegraphics[page=129,pagebox=mediabox,trim={65.02bp 166.13bp 221.14bp 519.16bp},clip,max width=0.94\textwidth,max height=0.38\textheight]{Marchisio2013.pdf}
\end{sourceblock}
α + Hk−1(xiα, σ)，其中 Hk−1(x, y) 是 x 和 y 中的 k −1 阶齐次多项式，我们可以将整数矩重写为

\begin{sourceblock}
\includegraphics[page=129,pagebox=mediabox,trim={184.98bp 111.24bp 67.57bp 555.71bp},clip,max width=0.94\textwidth,max height=0.38\textheight]{Marchisio2013.pdf}
\end{sourceblock}
其中左侧第二项可以用 σ 和矩 m* 表示

\begin{sourceblock}
\includegraphics[page=129,pagebox=mediabox,trim={110.78bp 74.24bp 337.58bp 609.75bp},clip,max width=0.94\textwidth,max height=0.38\textheight]{Marchisio2013.pdf}
\end{sourceblock}
{\small 0.换句话说，方程。 (3.109) 的形式类似于等式。 (3.91) 其中，高阶矩以低阶矩和 σ 形式已知。例如，直到 k = 4 我们有\par}
\begin{sourceblock}
\includegraphics[page=130,pagebox=mediabox,trim={172.30bp 563.38bp 65.08bp 57.72bp},clip,max width=0.94\textwidth,max height=0.38\textheight]{Marchisio2013.pdf}
\end{sourceblock}
正如上面针对高斯分布所讨论的，在给定可实现矩集 mk, k = 0,... 的情况下，可以迭代求解该方程组以确定 α = 1, 2 时的 σ 和 \{wα, xiα\}。。。 , 4. 扩展到更大的 N 使用相同的迭代过程。总之，我们应该注意方程的形式。 (3.109)（和方程（3.110））对于 δσ 的所有可能选择，不能保证右侧根据已知矩是闭合的。事实上，方程的选择。 (3.84) 使用 t = Ψ/σ 是专门制作的，以便所得矩系统将是封闭的。在定义 δσ 的其他选择（即使用其他类型的正交多项式的权重函数）时，读者应牢记这一重要的技术细节。接下来讨论在有限区间上具有此属性的 δσ 的示例。

\paragraph*{贝塔分布}
单变量 EQMOM 可以应用于在有限区间定义的 NDF（Yuan 等，2012）。例如，如果 NDF 仅在区间 xi ∈[0, 1] 上非零，那么我们可以使用 beta 分布定义核密度函数：

\begin{sourceblock}
\includegraphics[page=130,pagebox=mediabox,trim={182.12bp 320.82bp 65.08bp 352.20bp},clip,max width=0.94\textwidth,max height=0.38\textheight]{Marchisio2013.pdf}
\end{sourceblock}
其中 σ > 0 且 B 是 beta 函数。使用方程式找到的 NDF 示例(3.111) 和 N = 2 的四个 σ 值如图 3.9 所示。请注意，σ 控制峰的宽度，因此较小的 σ 值会产生更接近 delta 函数的峰。等式的右侧。 (3.111) 对应于雅可比的（缩放和移位）权重函数

\begin{sourceblock}
\includegraphics[page=130,pagebox=mediabox,trim={58.00bp 81.40bp 58.23bp 454.88bp},clip,width=\textwidth,max height=0.72\textheight]{Marchisio2013.pdf}
\par\vspace{0.45\baselineskip}
\small 图3.9。用方程找到的 n(xi) 的例子。 (3.111) 对于 σ 的四个值。其他参数为 w1 = w2 = 0.5、 Σ1 = 1/3 和 Σ2 = 2/3。
\end{sourceblock}
多项式，17 和积分，如方程 17 所示。 （3.87）可以使用相应的递归系数来评估（Gautschi，2004）。对于方程中的贝塔分布。 （3.111），整数矩可由递推公式求得：

\begin{sourceblock}
\includegraphics[page=131,pagebox=mediabox,trim={169.36bp 567.04bp 67.57bp 99.89bp},clip,max width=0.94\textwidth,max height=0.38\textheight]{Marchisio2013.pdf}
\end{sourceblock}
m(α) 0 = 1。因此，我们可以将 NDF 的整数矩表示为

\begin{sourceblock}
\includegraphics[page=131,pagebox=mediabox,trim={199.93bp 498.62bp 67.57bp 168.31bp},clip,max width=0.94\textwidth,max height=0.38\textheight]{Marchisio2013.pdf}
\end{sourceblock}
其中

\begin{sourceblock}
\includegraphics[page=131,pagebox=mediabox,trim={148.21bp 447.44bp 67.57bp 213.59bp},clip,max width=0.94\textwidth,max height=0.38\textheight]{Marchisio2013.pdf}
\end{sourceblock}
这与等式非常相似。 （3.108）。因为 Ψα 只出现在式（1）中的分子中。 (3.114)，乘积将产生 k 阶 Σα 的多项式。因此，我们可以将 NDF 的整数矩表示为 mk = γkm*

\begin{sourceblock}
\includegraphics[page=131,pagebox=mediabox,trim={231.07bp 403.24bp 159.81bp 280.41bp},clip,max width=0.94\textwidth,max height=0.38\textheight]{Marchisio2013.pdf}
\end{sourceblock}
1，其中系数 γk 仅取决于参数 σ。例如，直到 k = 4 我们有

\begin{sourceblock}
\includegraphics[page=131,pagebox=mediabox,trim={100.19bp 278.38bp 67.57bp 314.61bp},clip,max width=0.94\textwidth,max height=0.38\textheight]{Marchisio2013.pdf}
\end{sourceblock}
这与等式非常相似。 (3.110)。正如上面针对高斯分布所讨论的，在给定可实现矩集 mk, k = 0,... 的情况下，可以迭代求解该方程组以确定 α = 1, 2 时的 σ 和 \{wα, xiα\}。。。 , 4.

\subsection*{3.3.3 使用 EQMOM 评估积分}\phantomsection\addcontentsline{toc}{subsection}{3.3.3 使用 EQMOM 评估积分}
虽然 EQMOM 允许我们捕获额外的时刻，但在评估积分（例如方程 1）时，使用核密度函数可能会导致闭包问题。 （3.9）：

\begin{sourceblock}
\includegraphics[page=131,pagebox=mediabox,trim={156.42bp 128.32bp 67.57bp 535.26bp},clip,max width=0.94\textwidth,max height=0.38\textheight]{Marchisio2013.pdf}
\end{sourceblock}
如前所述，如果 δσ 是使用已知正交多项式族的权重函数来定义的，则方程： (3.116) 可以如式(3.116)中所述闭合。 (3.86) 使用递归

\footnotetext[17]{为 t ε[−1, 1] 定义的雅可比多项式 P(α,β) n (t) 具有权重函数 w(t) = (1−t)α(1+t)β (Gautschi, 2004)。通过平移和缩放，它们可以用于任何有限区间 ψ ∈[a, b]：t = (a + b −2ψ)/(a −b)。因此，如果 tα 是区间 [−1, 1] 中的节点，则 xiα = [a + b −(a −b)tα]/2 就是 [a, b] 中的节点。}
系数。否则，假设可以解析计算 δσ(Ψ, Ψα) 的整数矩：18

\begin{sourceblock}
\includegraphics[page=132,pagebox=mediabox,trim={196.20bp 589.66bp 65.08bp 76.15bp},clip,max width=0.94\textwidth,max height=0.38\textheight]{Marchisio2013.pdf}
\end{sourceblock}
闭合问题可以使用高斯求积来处理，如下所述。然而，我们应该注意到，使用递归系数可以进行更准确的积分评估，因为矩反演问题通常是病态的（Gautschi，2004）。考虑 2N′（可实现）矩的集合 \{m(α) 0 , m(α) 1 ,...。。。 , m(α) 2N′−1\} 对于 α ε (1, 2, . . . , N) 由方程定义。 （3.117）。使用自适应惠勒算法，对于 α 的每个值，这组矩导致 N′ 权重和 N′ 节点：19 \{w(α) 1 , w(α) 2 ,...。。。 , w(α) N′ \} 和 \{ xi(α) 1 , xi(α) 2 , .。。 , xi(α) N′ \}。方程中的积分。然后可以使用 NDF 的双正交表示来编写 (3.116)：

\begin{sourceblock}
\includegraphics[page=132,pagebox=mediabox,trim={168.05bp 448.54bp 65.08bp 215.04bp},clip,max width=0.94\textwidth,max height=0.38\textheight]{Marchisio2013.pdf}
\end{sourceblock}
请注意，N' 可以独立于 N 进行选择，以便可以在保持矩数固定（即等于 2N + 1）的同时增加求积的阶数。通过选择 δσ 使其具有接近真实限制 NDF 的形式，这一事实可以有利地用于限制情况（例如，如果限制 NDF 是加权高斯，则方程（3.88）是 δσ(x, y) 的不错选择）。在这些情况下，N 可以很小（例如 N = 1），同时通过调整等式中的 N' 仍保留高斯求积近似的精度。 （3.118）。作为一个例子，考虑 δσ 是方程 1 中的高斯分布的情况。 （3.88）。在这种情况下，σ 充当缩放参数，因此方程中的积分也起作用。 (3.118) 可以写成 ∫

\begin{sourceblock}
\includegraphics[page=132,pagebox=mediabox,trim={174.03bp 303.55bp 65.08bp 362.86bp},clip,max width=0.94\textwidth,max height=0.38\textheight]{Marchisio2013.pdf}
\end{sourceblock}
其中 N′ 权重 w* β 和 N′ 节点 φβ 从前 2N′ 个标准法向矩中找到：20

\begin{sourceblock}
\includegraphics[page=132,pagebox=mediabox,trim={176.15bp 247.91bp 65.08bp 415.92bp},clip,max width=0.94\textwidth,max height=0.38\textheight]{Marchisio2013.pdf}
\end{sourceblock}
（有关 N′ = 4 时的权重与节点的示例，请参见练习 3.2。）请注意，由于矩 μk 是固定的，因此对于给定的 N′ 值，权重 w* β 和节点 φβ 只需计算一次。因此，与 QMOM 相比，使用 EQMOM 的额外计算成本主要与 σ 的计算相关。在极限情况下 σ = 0，方程： (3.119) 简化为

\begin{sourceblock}
\includegraphics[page=132,pagebox=mediabox,trim={147.75bp 142.32bp 65.08bp 521.26bp},clip,max width=0.94\textwidth,max height=0.38\textheight]{Marchisio2013.pdf}
\end{sourceblock}
18 这一步是选择 δσ(x, y) 以获得易于计算矩的主要原因。 19对于某些情况，权重不依赖于 α，因此可以在 xiα = 0 时进行计算。同样，如果 σ 输入 δσ 的定义作为缩放参数（即 σδσ(xi, ϋα) = δσ=1((xi −ϋα)/σ, 0)），则可以根据 δσ=1(x, 0) 的整数矩找到权重与节点：如果 φβ 是所得节点，则 xi(α) β = σφβ + xiα。请注意方程中的分布。 (3.84) 没有这种形式，这解释了为什么必须针对每个 α 值单独计算节点。

\footnotetext[20]{更好的替代方案是使用 Hermite 多项式的递归系数直接形成雅可比矩阵（Gautschi，2004）。}
这是 QMOM 近似。在极限情况 σ2 = 〈Σ′2〉 下，我们有 Σα = 〈Σ〉，因此 ∫

\begin{sourceblock}
\includegraphics[page=133,pagebox=mediabox,trim={176.73bp 584.90bp 67.57bp 81.51bp},clip,max width=0.94\textwidth,max height=0.38\textheight]{Marchisio2013.pdf}
\end{sourceblock}
这是基于高斯 NDF 的 N′ 节点高斯求积近似（Gautschi，2004）。总而言之，我们可以注意到方程。 (3.119) 等价于将单变量 NDF 近似为对偶求积表示：

\begin{sourceblock}
\includegraphics[page=133,pagebox=mediabox,trim={133.69bp 501.95bp 67.57bp 164.45bp},clip,max width=0.94\textwidth,max height=0.38\textheight]{Marchisio2013.pdf}
\end{sourceblock}
其中我们明确显示取决于时间 t 和位置 x 的数量。方程的扩展。 (3.123) 到除高斯函数之外的其他核密度函数（例如方程（3.87））是简单的。从 NDF 的双求积表示中，σ 在确定矩输运方程中未闭项的 EQMOM 近似中的作用通过时间/空间无关的权重 w* β 和节点 φβ 与底层权重函数明确相关。

\subsection*{3.3.4 多元EQMOM}\phantomsection\addcontentsline{toc}{subsection}{3.3.4 多元EQMOM}
原则上，3.3.2节中介绍的EQMOM可以推广到包括多个内坐标。然而，根据核密度函数的假设形式，可能需要使用多元非线性方程求解器来查找参数（即类似于第 3.2.1 节中讨论的强力 QMOM）。一个有趣的替代方案是扩展第 3.2.3 节中描述的 CQMOM 算法。在这里我们考虑使用这两种方法的示例。

\paragraph*{穷举式 QMOM}
在文献中（Chalons 等人，2010），仅考虑了具有四个节点的二元 EQMOM，由具有对角协方差矩阵的加权高斯分布表示。然而，很可能可以为其他分布函数开发强力 QMOM 算法。以多高斯表示为例，近似 NDF 可以写为

\begin{sourceblock}
\includegraphics[page=133,pagebox=mediabox,trim={198.13bp 179.32bp 67.57bp 487.62bp},clip,max width=0.94\textwidth,max height=0.38\textheight]{Marchisio2013.pdf}
\end{sourceblock}
其中，对于 M 内坐标，核密度函数定义为

\begin{sourceblock}
\includegraphics[page=133,pagebox=mediabox,trim={144.93bp 123.69bp 67.57bp 541.85bp},clip,max width=0.94\textwidth,max height=0.38\textheight]{Marchisio2013.pdf}
\end{sourceblock}
M × M 对称协方差矩阵 Σ 包含缩放分量 σij。未知缩放分量的总数为 M(M+1)/2，必须确定这些

\footnotetext[21]{在对应于方程的双正交表示中。 （3.87）中，权重与节点既不是时间独立的，也不是空间独立的，因为 kα 取决于传递的矩。}
\footnotetext[22]{对于对角协方差矩阵，Σ = diag(σ11, σ22)。}
通过固定相同数量的多元矩。回想一下，CQMOM 中使用的最佳矩集包含 (M + 1)nM 个矩，其中 n 是单变量求积中的节点数23，因此 N = nM。由于确定 Σ 所需的附加矩必须独立于最佳矩集，因此多元 EQMOM 将使用总共 (M +1)(nM + M/2) 个多元矩。下面，我们将这个矩集合称为扩展最优矩集合。为了与 3.3.2 节中讨论的单变量 EQMOM 保持一致，Σ 的对角分量由 2n 阶纯矩确定（例如，在 2D 相空间中，m(2n, 0) 和 m(0, 2n) 分别用于求 σ11 和 σ22）。然后，可以通过使用总阶数为 2n、但 i 和 j 分量中的阶数为 n 的交叉矩来确定 Σ 的非对角分量。例如，在 2D 相空间中，使用矩 m(n, n) 求 σ12，在 3D 相空间中，使用矩 m(n, 0, n) 求 σ13。读者可以验证这些矩都不包含在最佳矩集中，因此它们可用于扩展最佳矩集中。多元矩 m(k) 由方程计算得出。 (3.124) 可以写成

\begin{sourceblock}
\includegraphics[page=134,pagebox=mediabox,trim={126.93bp 410.11bp 65.08bp 245.51bp},clip,max width=0.94\textwidth,max height=0.38\textheight]{Marchisio2013.pdf}
\end{sourceblock}
其中关于 xi 的整数矩定义为

\begin{sourceblock}
\includegraphics[page=134,pagebox=mediabox,trim={198.39bp 363.02bp 65.08bp 310.73bp},clip,max width=0.94\textwidth,max height=0.38\textheight]{Marchisio2013.pdf}
\end{sourceblock}
QMOM 矩定义为

\begin{sourceblock}
\includegraphics[page=134,pagebox=mediabox,trim={210.32bp 295.84bp 65.08bp 371.10bp},clip,max width=0.94\textwidth,max height=0.38\textheight]{Marchisio2013.pdf}
\end{sourceblock}
在这些表达式中，我们使用了索引符号 i = (i1,...,iM)，使得 xii = xii1 1。。。 xiM M。请注意，在极限条件下 |Σ| →0，只有零阶矩 mΣ(0) = 1 是非零的，式（1） (3.126) 恢复到 QMOM 矩。中心高斯矩 mΣ(i) 是协方差矩阵 Σ 的已知函数。对于给定的矩阶 γ = k1 + · · · + kM，方程： (3.126) 有一个下三角形式，可以使用前向替换进行反转：

\begin{sourceblock}
\includegraphics[page=134,pagebox=mediabox,trim={112.58bp 114.28bp 65.08bp 490.58bp},clip,max width=0.94\textwidth,max height=0.38\textheight]{Marchisio2013.pdf}
\end{sourceblock}
其中 ei 是第 i 个单位向量。因此，给定一个已知的扩展最优矩集 m(k) 和一个已知的协方差矩阵 Σ，方程： (3.129) 可用于找到最优矩集

\footnotetext[23]{这里，为了简化符号，与 3.2 节中报告的相比，我们假设 N1 = N2 = · · · = NM = n。}
m*(k) 用于多元求积（例如 CQMOM）。换句话说，一旦 m*(k) 已知，我们就可以计算一组 N 个权重 wα 和 N 个节点 xiα。因此，仍然需要解释如何从扩展的最佳时刻确定 Σ。计算 Σ 的算法类似于单变量情况中使用的算法。使用求根（或最小化）方案，我们必须改变 Σ 的组成部分，直到目标函数最小化。从扩展最优矩找到的目标函数的组成部分是

\begin{sourceblock}
\includegraphics[page=135,pagebox=mediabox,trim={81.71bp 355.83bp 67.57bp 150.24bp},clip,max width=0.94\textwidth,max height=0.38\textheight]{Marchisio2013.pdf}
\end{sourceblock}
以 f1(Σ) 和 f1,2(Σ) 为例，方程右侧的矩 m* 为： (3.130) 的阶数为 γ < 2n，可从方程 (3.130) 中找到。 (3.129) 对于给定的 Σ 值。根据这些矩，计算权重与节点，然后用于计算 γ = 2n 阶矩：

\begin{sourceblock}
\includegraphics[page=135,pagebox=mediabox,trim={84.62bp 261.03bp 67.57bp 405.91bp},clip,max width=0.94\textwidth,max height=0.38\textheight]{Marchisio2013.pdf}
\end{sourceblock}
按照此过程，从方程计算目标函数所需的所有时刻。 (3.130) 现已推出。通过迭代 Σ，在矩 m* 仍然可实现的约束下，可以最小化目标函数分量的平方和。需要注意的是，虽然在mΣ(i, 0, ..., 0)和mΣ(i, j, 0, ..., 0)等时刻中只出现了Σ的少数分量，但权重与节点涉及所有分量。因此，一般来说，不可能将单独组件的求解算法解耦。多元 EQMOM 计算积分的应用与单变量 EQMOM 非常相似。对于多高斯 EQMOM，我们有

\begin{sourceblock}
\includegraphics[page=135,pagebox=mediabox,trim={150.03bp 103.25bp 67.57bp 560.33bp},clip,max width=0.94\textwidth,max height=0.38\textheight]{Marchisio2013.pdf}
\end{sourceblock}
{\small 其中 N′ = (n′)M，n′ 是用于表示每个单变量高斯分布的正交点数。平方根矩阵 Σ1/2，定义为 Σ =\par}
{\small Σ1/2 2 是变换张量积节点 φβ 的便捷（但不是唯一）选择。的\par}
后者（以及权重 w* β）是使用具有单位协方差矩阵（或 Hermite 多项式的递归系数）的 M 维高斯分布找到的。 （有关 n = 2、3 的相关示例，请参见图 3.5。）如练习 3.8 所示，对于单位协方差矩阵，所有具有高斯矩的多元求积方法都会产生相同的 w* β 和 φβ 值（与 Hermite 多项式的零点相关（Gautschi，2004 年））。与单变量情况一样，w*β 和 φβ 是常数，并且可以针对给定的 M 和 n′ 值进行预先计算。因此，多元 NDF 的双正交表示为

\begin{sourceblock}
\includegraphics[page=136,pagebox=mediabox,trim={130.05bp 506.49bp 65.08bp 159.91bp},clip,max width=0.94\textwidth,max height=0.38\textheight]{Marchisio2013.pdf}
\end{sourceblock}
我们可以再次清楚地看到 Σ 在正交近似中所起的作用。总之，虽然多元 EQMOM 在矩方面得到了很好的定义，但关键的步骤是计算给定（可实现的）扩展最优矩集的 Σ。据作者所知，本节中描述的强力 QMOM 已应用于仅使用 Σ = diag(σ11, σ22) 的 CQMOM 求解双变量情况 (M = n = n′ = 2) 的矩输运方程 (Chalons et al., 2010)。

\paragraph*{扩展CQMOM}
为了避免需要迭代求解器，我们可以扩展第 3.2.3 节中描述的 CQMOM 算法，以包含计算参数 σαβ 所需的附加矩。然而，为了应用该算法，我们必须假设条件 PDF 对于条件变量的所有值都具有相同的形式。作为一个例子，考虑双变量的情况（Ψ1，Ψ2）。以 Ψ1 为条件，二元 NDF 可写为 f(Ψ1, Ψ2) = f(Ψ1) f(Ψ2|Ψ1)。通过扩展 CQMOM (ECQMOM)，对于 xi2 以 xi1 为条件，我们可以假设 24 NDF 近似为

\begin{sourceblock}
\includegraphics[page=136,pagebox=mediabox,trim={152.24bp 274.36bp 65.08bp 392.58bp},clip,max width=0.94\textwidth,max height=0.38\textheight]{Marchisio2013.pdf}
\end{sourceblock}
请注意，虽然扩展参数 σ1 对于求和中的所有项都相同，但条件参数 σ2,α 可以取决于 α。此外，用于 δσ1 的函数形式不必与用于 δσ2,α 的函数形式相同（例如 Ψ1 可以使用 beta EQMOM，而 Ψ2 使用高斯 EQMOM）。虽然方程中的形式。 （3.134）并不像等式（3.134）中那样普遍。 （3.124），我们将看到它导致了一种与 CQMOM 中使用的方法非常相似的直接矩反演方法。从方程中找到的双变量矩。 （3.134）有形式

\begin{sourceblock}
\includegraphics[page=136,pagebox=mediabox,trim={117.99bp 123.04bp 65.08bp 504.48bp},clip,max width=0.94\textwidth,max height=0.38\textheight]{Marchisio2013.pdf}
\end{sourceblock}
1,i 和 m(αβ) 2, j 由右侧的积分定义。因此，对于 j = 0，相对于 xi1 的单变量矩减少到 m(α)

\footnotetext[24]{δσ2,α(Σ2, Σ2,αβ) 可以有更一般的形式。例如，它可能取决于 Ψ1，因为在求解 δσ2,α(Ψ2, Ψ2,αβ) 中的参数时将已知 δσ1(Ψ1, Ψ1,α)。然而，如果尝试对 Ψ1 进行过于复杂的依赖，数值反演算法很快就会变得棘手。}
\begin{sourceblock}
\includegraphics[page=137,pagebox=mediabox,trim={65.02bp 589.42bp 67.57bp 58.62bp},clip,max width=0.94\textwidth,max height=0.38\textheight]{Marchisio2013.pdf}
\end{sourceblock}
β=1wαβ。请注意等式。 (3.136) 与单变量 EQMOM 具有相同的形式，因此可以使用相同的算法来查找参数 wα、Σ1、α 和 σ1。此外，通过构造，每个 α 值的 m(α) 1,i 是这些参数的已知函数。遵循与 CQMOM 中使用的逻辑相同的逻辑，我们接下来考虑具有固定 j 和 i ≥0 的时刻：

\begin{sourceblock}
\includegraphics[page=137,pagebox=mediabox,trim={194.68bp 388.07bp 67.57bp 153.71bp},clip,max width=0.94\textwidth,max height=0.38\textheight]{Marchisio2013.pdf}
\end{sourceblock}
我们引入了符号

\begin{sourceblock}
\includegraphics[page=137,pagebox=mediabox,trim={208.59bp 337.69bp 67.57bp 329.25bp},clip,max width=0.94\textwidth,max height=0.38\textheight]{Marchisio2013.pdf}
\end{sourceblock}
且 w* αβ = wαβ/wα。在等式中。 (3.137) 中，权重 wα 和节点 xi1,α （以及矩 m(α) 1,i ）已知。同样，方程左侧的双变量矩。 (3.137) 也是已知的（即它们是通过求解矩输运方程找到的）。新变量 μj α 对应于给定 Ψ1 的 Ψ2 的条件矩。此外，方程的形式。固定 α 的式（3.138）与 EQMOM 中使用的矩系统完全相同。因此，假设 μj α 可以从等式计算。 (3.137) 对于 j = 0, 1, .。。，2N，等式。 (3.138) 可以反演得到 w* αβ、 ϋ2,αβ 和 σ2,α。与 CQMOM 类似，条件矩是通过求解线性系统找到的：

\begin{sourceblock}
\includegraphics[page=137,pagebox=mediabox,trim={65.03bp 116.54bp 67.57bp 467.83bp},clip,max width=0.94\textwidth,max height=0.38\textheight]{Marchisio2013.pdf}
\end{sourceblock}
1、i替换节点。在 σ1 →0 的极限下，这些矩变为 m(α) 1,i = xii 1,α，因此系数矩阵与此极限下的 CQMOM 中的相同。一般来说，只要矩向量 φα = [m(α)

{\small 1,0, m(α) 1,1, .。。 , m(α) 1,N−1]T 是线性无关的，方程组中的系统可以求解（3.139）来找到条件矩。\par}
作为一个例子，考虑一个四节点双高斯分布：

\begin{sourceblock}
\includegraphics[page=138,pagebox=mediabox,trim={144.10bp 587.63bp 65.08bp 89.50bp},clip,max width=0.94\textwidth,max height=0.38\textheight]{Marchisio2013.pdf}
\end{sourceblock}
其中高斯函数是

\begin{sourceblock}
\includegraphics[page=138,pagebox=mediabox,trim={173.42bp 537.19bp 65.08bp 134.47bp},clip,max width=0.94\textwidth,max height=0.38\textheight]{Marchisio2013.pdf}
\end{sourceblock}
在等式中。 (3.140) 中，有 13 个参数，(wαβ, Ψ1,α, Ψ2,αβ; α, β = 1, 2) 和 (σ1, σ2,α; α = 1, 2)。这些参数必须通过求解对应于以下 13 个二元矩（直至五阶）的非线性系统来找到：

\begin{sourceblock}
\includegraphics[page=138,pagebox=mediabox,trim={180.21bp 429.00bp 177.75bp 204.87bp},clip,max width=0.94\textwidth,max height=0.38\textheight]{Marchisio2013.pdf}
\end{sourceblock}
与 1D EQMOM 中一样，m4,0 用于查找 σ1，而需要 m0,4 和 m1,4 才能查找 α = 1, 2 时的 σ2,α。对于 j = 0，从等式中找到的五个矩： (3.136) 当 i = 0 时，.。。 , 4 是

\begin{sourceblock}
\includegraphics[page=138,pagebox=mediabox,trim={123.43bp 302.06bp 119.32bp 319.00bp},clip,max width=0.94\textwidth,max height=0.38\textheight]{Marchisio2013.pdf}
\end{sourceblock}
该系统具有与第 3.3.2 节中的一维双高斯求积完全相同的形式，因此可以直接求逆（即无需迭代）以找到五个参数（w1、w2、Σ1,1、Σ1,2 和 σ2 1）。从方程中找到条件矩的线性系统。 (3.139) 是

\begin{sourceblock}
\includegraphics[page=138,pagebox=mediabox,trim={108.45bp 212.01bp 110.12bp 450.75bp},clip,max width=0.94\textwidth,max height=0.38\textheight]{Marchisio2013.pdf}
\end{sourceblock}
其中μ0 1 = μ0 2 = 1。只要 Ψ1,1 Ψ1,2 且 w1,w2 > 0，该系统就可以反演得到 μj α。根据条件矩的定义，我们有

\begin{sourceblock}
\includegraphics[page=138,pagebox=mediabox,trim={124.76bp 93.45bp 120.64bp 524.48bp},clip,max width=0.94\textwidth,max height=0.38\textheight]{Marchisio2013.pdf}
\end{sourceblock}
{\small 对于每个 α，该系统再次具有与一维双高斯求积中完全相同的形式，因此可以反转以找到五个参数（wα1、wα2、Σ2,α1、Σ2,α2 和 σ2）\par}
{\small 2,α)。完成此操作后，方程式中将出现 13 个参数。 (3.140) 是已知的。\par}
除退化情况外，上述 2D 矩反演算法能够恢复 13 个二元矩。实际上，存在几种可能的简并情况，例如，当真实分布 f 可以由少于四个高斯分布的总和精确表示时。另一种可能的简并性发生在 Σ1 的方差为零，但 Σ2 的方差是有限的（即 f 的支撑是垂直于 Σ1 轴的线）时。通过使用自适应惠勒算法来处理所有可能的简并情况，只要输入的二元矩集是可实现的，就可以使用 ECQMOM 构建鲁棒的二维矩反演算法来查找未知参数。由于 ECQMOM 仅需要多个 1D EQMOM 求解，因此计算速度比强力 QMOM 算法快得多。

\section*{3.4 直接求积矩方法（DQMOM）}\phantomsection\addcontentsline{toc}{section}{3.4 直接求积矩方法（DQMOM）}
从技术上讲，矩的直接求积法 (DQMOM) 并不是矩反演算法。相反，通过 DQMOM，使用矩反演算法来查找权重 wα(0, x) 和节点 xiα(0, x) 的初始值，然后求解 wα(t, x) 和 xiα(t, x) 的传输方程。显然，DQMOM 与 EQMOM 不兼容，因为它依赖于近似 NDF 的双正交表示。 Marchisio \& Fox (2005) 和 Fox (2003) 从 NDF 传输方程开始，描述了这些传输方程的形式推导。另一种替代方法是从整数矩的输运方程出发，通过变量的非线性变化来推导 DQMOM 输运方程

\begin{sourceblock}
\includegraphics[page=139,pagebox=mediabox,trim={158.34bp 311.67bp 160.86bp 355.28bp},clip,max width=0.94\textwidth,max height=0.38\textheight]{Marchisio2013.pdf}
\end{sourceblock}
开发 DQMOM 的最初动机是缺乏 M > 1 的鲁棒多元矩反演算法（即现在通过 CQMOM 的开发克服了这一问题）。原则上，DQMOM 仅需要多元矩反演算法来找到初始条件，然后通过求解权重与节点的传输方程来找到权重与节点的时间演化。事实上，对于 M = 1，求解矩的输运方程的单变量 QMOM 算法相当于抛物线系统（即具有扩散项的系统）的 DQMOM（达到数值误差）。对于纯双曲系统，由于特征交叉，DQMOM 方程的解可能会有所不同。25（详细信息请参阅附录 C。）因此，对于双曲系统，当节点 N 的数量固定时，求解矩输运方程是唯一物理上有效的选择。为了理解 DQMOM，我们在这里仅考虑与 NDF 输运方程相关的单变量方程：

\begin{sourceblock}
\includegraphics[page=139,pagebox=mediabox,trim={190.65bp 119.58bp 67.57bp 551.31bp},clip,max width=0.94\textwidth,max height=0.38\textheight]{Marchisio2013.pdf}
\end{sourceblock}
其中 S (xi) 是通用源项，且 Γ > 0 是（常数）空间扩散系数。假设速度 u(t, x) 已知，初始 NDF f(0, x, xi) 也是已知的。的

\footnotetext[25]{当特性交叉时，非零权重的数量 N 必须局部增加，以便使用 DQMOM 获得物理上有效的解决方案。然而，由于交叉的数量可能趋于无穷大，因此数值复杂性很快就会变得棘手。}
对应于方程的矩输运方程。 (3.142) 是

\begin{sourceblock}
\includegraphics[page=140,pagebox=mediabox,trim={165.95bp 593.59bp 65.08bp 73.85bp},clip,max width=0.94\textwidth,max height=0.38\textheight]{Marchisio2013.pdf}
\end{sourceblock}
其中源项积分必须通过求积来闭合：

\begin{sourceblock}
\includegraphics[page=140,pagebox=mediabox,trim={181.21bp 536.24bp 178.75bp 127.34bp},clip,max width=0.94\textwidth,max height=0.38\textheight]{Marchisio2013.pdf}
\end{sourceblock}
在该表达式中，使用矩反演算法从每个时间步长和每个空间位置的矩中找到权重与节点。按照 Marchisio \& Fox (2005) 中描述的方法，DQMOM 方程对应于 Eq. (3.143) 是

\begin{sourceblock}
\includegraphics[page=140,pagebox=mediabox,trim={154.29bp 423.30bp 65.08bp 223.86bp},clip,max width=0.94\textwidth,max height=0.38\textheight]{Marchisio2013.pdf}
\end{sourceblock}
其中 aα 和 bα 是通过强制方程找到的源项。 (3.144) 与 k = 0, 1, . 的矩输运方程一致。。。 , 2N -1。为了清楚起见，我们将考虑 N = 2 的情况。矩约束的线性系统得出

\begin{sourceblock}
\includegraphics[page=140,pagebox=mediabox,trim={87.10bp 324.92bp 65.08bp 318.76bp},clip,max width=0.94\textwidth,max height=0.38\textheight]{Marchisio2013.pdf}
\end{sourceblock}
等式左侧的系数矩阵26如果 Ψ1 Ψ2 则 (3.145) 将是非奇异的，这对应于非零方差。右侧的矢量是由于扩散项而产生的，并且涉及节点的梯度平方。在实践中，这些梯度可能变得非常大，使得源项非常僵硬。显然，如果初始 NDF 的方差为零（参见 Marchisio \& Fox (2005) 的示例），或者如果节点几乎相等，导致系数矩阵病态，则必须特别小心。尽管如此，只要 Γ > 0 并且数值算法完全解决了空间梯度问题，方程的解就可以得到。与方程（3.144）的解相比，方程（3.144）仍然相当准确。 （3.143）。另一方面，当 Γ 非常小（或为零）时，DQMOM 解的精度可能会因数值误差而严重降低。设计方程的数值方法时要认识到的重要一点。 （3.144）的特点是权重与节点不是守恒量，而矩在没有源项的情况下是守恒的。因此，保守的有限体积方案非常适合求解方程（1）。 (3.143)，但由于非零源项 aα 和 bα，不保证 DQMOM 方程的守恒性。相对于 DQMOM 求解矩输运方程的明显缺点是必须在每个时间步和每个网格单元中应用矩反演算法来查找权重与节点。然而，由于用矩输运方程进行高阶通量重构将不可避免地导致无法实现的矩

\footnotetext[26]{正如已经提到的，该系数矩阵（也表示为 A）对应于方程 2 中报告的非线性系统的雅可比行列式。 (3.43) 并在强力 QMOM 中用于计算正交近似。}
即使矩输运方程闭合，如果使用高阶通量，也无法避免矩反演算法（详细信息请参见附录 B）。这一事实对于多变量情况尤其如此，对于多变量情况，确保矩可实现性比单变量情况更加困难。最后，对于双曲和近双曲系统，EQMOM 矩反演算法提供了空间通量的更高保真度表示。由于 DQMOM 与 EQMOM 不兼容，因此应通过使用 EQMOM 求解矩输运方程来处理空间扩散非常小的矩系统。

\chapter{广义群体平衡方程}
本章的目标是从分散多相系统的微观描述出发，提供广义群体平衡方程（GPBE）的相对严格的推导。我们首先使用概率密度函数 (PDF) 方法定义离散粒子系统的数密度函数 (NDF)。一旦定义了 NDF，我们就通过引入条件期望值的概念来推导 GPBE。后者包含微尺度物理的介观尺度表示，因此包含关闭 GPBE 所需的介观尺度模型。接下来，我们将在后面的章节中详细解释如何从 GPBE 中找到选定的感兴趣的多元矩的输运方程。本章最后对矩输运方程中的矩闭合进行了简短描述。

\section*{4.1 基于粒子的 NDF 定义}\phantomsection\addcontentsline{toc}{section}{4.1 基于粒子的 NDF 定义}
在 NDF 的基于粒子的定义中，我们从微观尺度开始，写出介观尺度上分散相粒子特性变化率的动态方程。我们首先考虑的最简单的系统是真空中相互作用的粒子的集合，其中粒子通过碰撞和短程力相互作用。这样的系统称为粒度系统。然后我们考虑分散两相系统，其中颗粒分散在流体中。

\subsection*{4.1.1 粒状系统 NDF 的定义}\phantomsection\addcontentsline{toc}{subsection}{4.1.1 粒状系统 NDF 的定义}
首先考虑周围没有流体相的 Np 颗粒颗粒系统。第 n 个粒子的动力学可以通过其质心位置 X(n),1 来描述

\footnotetext[1]{由于粒子的尺寸有限，因此粒子所占的体积将不为零。然而，就粒子动力学而言，我们可以通过粒子的质心来标记粒子，质心是一个没有体积的点。}
它的速度 U(n),2 和 M 内坐标 ϋ(n).3 为简单起见，我们现在假设粒子具有相同的质量。不失一般性，粒子的动力学可以通过微分方程组来描述：

\begin{sourceblock}
\includegraphics[page=143,pagebox=mediabox,trim={205.28bp 541.69bp 67.56bp 105.57bp},clip,max width=0.94\textwidth,max height=0.38\textheight]{Marchisio2013.pdf}
\end{sourceblock}
以及

\begin{sourceblock}
\includegraphics[page=143,pagebox=mediabox,trim={207.44bp 484.42bp 67.56bp 188.79bp},clip,max width=0.94\textwidth,max height=0.38\textheight]{Marchisio2013.pdf}
\end{sourceblock}
右侧的算子分别表示由于体积力引起的连续加速度 A(n)、由于不连续碰撞（或粒子与粒子相互作用）引起的变化 C(n) xi 以及由于“增长”G(n) （即内坐标的连续变化）引起的变化。一般来说，这些算子取决于所有粒子的完整变量集，并且是包含粒子的系统体积的固有函数。换句话说，系统体积必须足够大以包含所有非零体积的粒子。同样，我们可以说算子是粒子数密度的本征函数。为了简化符号，我们将所有粒子变量的集合表示如下：U和C(n)

\begin{sourceblock}
\includegraphics[page=143,pagebox=mediabox,trim={193.21bp 327.55bp 67.56bp 332.57bp},clip,max width=0.94\textwidth,max height=0.38\textheight]{Marchisio2013.pdf}
\end{sourceblock}
以及

\begin{sourceblock}
\includegraphics[page=143,pagebox=mediabox,trim={195.36bp 277.23bp 67.56bp 401.36bp},clip,max width=0.94\textwidth,max height=0.38\textheight]{Marchisio2013.pdf}
\end{sourceblock}
请注意，该系统共有 Nd = Np(3 + 3 + M) 个自由度。对于每组初始条件，方程。求解 (4.1)–(4.3) 即可找到 \{X(n)\}、\{U(n)\} 和 \{xi(n)\}。初始条件是从已知的分布函数中随机选择的，我们可以假设有无限多种可能的组合。每个组合称为粒度流的实现，所有可能实现的集合形成一个整体。请注意，由于粒子的大小有限，因此它们不能位于同一点：因此 X(n) X(m) 表示 n m。此外，碰撞算子 C(n) U 将生成混沌轨迹，因此在相对少量的碰撞后，粒子位置将变得不相关。相比之下，对于悬浮在流体中的粒子来说，碰撞受到抑制，并且相关性可以是长期存在的并且是长范围的。当我们稍后介绍流体-粒子系统时，我们将使这些概念更加精确。虽然粒子相关性的确切性质不是下面介绍的多粒子联合 PDF 定义的一个因素，但重要的是要记住，它们将具有
2 我们可以使用粒子动量而不是速度。然而，由于粒子质量可以是内坐标之一，因此使用速度并不失一般性。

\footnotetext[3]{内坐标的概念非常笼统，可以包括粒子质量、半径、几何形状、角动量等。正如我们稍后所述，对内坐标如何因微观物理而变化进行建模通常并不简单。}
对粒子动力学有很大的影响。假设我们有无限的实现集合，我们可以在时间 t by4 定义一个多粒子联合 PDF

\begin{sourceblock}
\includegraphics[page=144,pagebox=mediabox,trim={87.43bp 540.51bp 65.07bp 87.08bp},clip,max width=0.94\textwidth,max height=0.38\textheight]{Marchisio2013.pdf}
\end{sourceblock}
fNp

其中右侧的概率是使用系综中的所有实现来计算的，Nd 个状态空间变量是 \{x(n)\}、\{V(n)\} 和 \{η(n)\}。尽管它的自由度太多，无法用于建模实际系统，但原则上，这个联合 PDF 完整地描述了粒状系统的动力学。通过积分其他粒子的所有状态变量，可以将多粒子联合 PDF 简化为单粒子联合 PDF。正式地，对于第 n 个粒子，这产生

\begin{sourceblock}
\includegraphics[page=144,pagebox=mediabox,trim={85.53bp 424.69bp 65.07bp 241.59bp},clip,max width=0.94\textwidth,max height=0.38\textheight]{Marchisio2013.pdf}
\end{sourceblock}
其中积分是对所有 m n 的积分。在所有粒子分布相同的特殊情况下，单粒子联合 PDF 对于所有 n 都是相同的，我们可以简单地用 f1.5 表示它。请注意，与 fNp 不同，f (n) 1 不提供粒子动力学的完整描述，因为通常不能从单粒子 PDF 中找到两个粒子（或更多）相关性。而且，f(n) 1 中相空间变量的数量为3+3+M，比Nd小1/Np倍。第 2 章中介绍的通用 NDF 是根据 fNp:6 的期望值定义的

\begin{sourceblock}
\includegraphics[page=144,pagebox=mediabox,trim={89.25bp 183.50bp 65.07bp 376.52bp},clip,max width=0.94\textwidth,max height=0.38\textheight]{Marchisio2013.pdf}
\end{sourceblock}
因此，如果所有粒子均等分布，我们发现

\begin{sourceblock}
\includegraphics[page=144,pagebox=mediabox,trim={189.74bp 140.65bp 65.07bp 545.67bp},clip,max width=0.94\textwidth,max height=0.38\textheight]{Marchisio2013.pdf}
\end{sourceblock}
4 概率陈述 P[A∩B∩C] 被解释为事件 A、B 和 C 同时发生的概率。概率计算为集合中语句 A ∩B ∩C 为真的所有实现的分数。 5 有人可能会争辩说，由于粒子编号系统是任意的，因此粒子必须是相同分布的。参见 Subramaniam (2000) 对不成立的例子的讨论。无论如何，在推导实际系统中的 NDF 模型时，可以方便地假设 f (n)

\footnotetext[1]{= f1.6 回想一下，该方程中的随机变量是粒子位置、速度和内坐标的完整集合。状态空间变量 x、v 和 xi 是固定的。}
其中状态空间变量 x、v 和 xi 是固定的。方程的第一行。 (4.9)是NDF的定义，第二行是多粒子联合PDF的期望值的定义。在第三行中，我们整合了除第n个粒子之外的所有粒子信息，只留下单粒子联合PDF。在最后一行，我们应用多元 delta 函数的定义。第 2 章中定义的另一个 NDF（即 nxi 和 nU）可以从方程（1）中找到。 （4.9）通过积分适当的内坐标。现在从等式可以明显看出。 (4.9) 一般 NDF 是根据以下概率陈述定义的：7

\begin{sourceblock}
\includegraphics[page=145,pagebox=mediabox,trim={89.71bp 504.58bp 67.56bp 155.53bp},clip,max width=0.94\textwidth,max height=0.38\textheight]{Marchisio2013.pdf}
\end{sourceblock}
右侧的概率再次相对于粒状流的所有实现的集合来定义。 NDF 的动态行为将由等式的右侧决定。 (4.1)-(4.3)，稍后我们将回到这方面。然而，首先，我们将研究基于颗粒流的单一实现来估计 NDF 的方法。

\subsection*{4.1.2 NDF估计方法}\phantomsection\addcontentsline{toc}{subsection}{4.1.2 NDF估计方法}
NDF 的正式定义如式（1）所示。 （4.11）在数学上是一致的，但在实践中很难实现。因此，定义从颗粒流的单一实现中估计 NDF 的方法是有用的。请注意，从数学上讲，统计估计是随机变量，因此不应与确定性的 NDF 混淆。为了区分估计的 NDF 和 n，我们将用 \textasciicircum{}n 表示估计值。因此，例如，如果估计是无偏的，则 〈ˆn〉= n，其中预期值是相对于等式 1 中定义的多粒子联合 PDF fNp 获取的。 （4.7）。

\paragraph*{体积平均值}
由于粒子是离散的，因此粒子位于给定点 x 的概率为零。因此，为了用有限的粒子样本来估计 NDF，我们需要引入一个以 x 为中心、带宽为 W 的核密度函数 hW(x)。例如，一个常数核密度函数定义为

\begin{sourceblock}
\includegraphics[page=145,pagebox=mediabox,trim={178.30bp 204.87bp 67.56bp 458.96bp},clip,max width=0.94\textwidth,max height=0.38\textheight]{Marchisio2013.pdf}
\end{sourceblock}
可以用于此目的。请注意，该内核是各向同性的、非负的，并且定义为使体积积分标准化：

\begin{sourceblock}
\includegraphics[page=145,pagebox=mediabox,trim={212.58bp 149.83bp 67.56bp 523.73bp},clip,max width=0.94\textwidth,max height=0.38\textheight]{Marchisio2013.pdf}
\end{sourceblock}
NDF 的估计量现在可以定义为

\begin{sourceblock}
\includegraphics[page=145,pagebox=mediabox,trim={121.51bp 90.17bp 67.56bp 576.17bp},clip,max width=0.94\textwidth,max height=0.38\textheight]{Marchisio2013.pdf}
\end{sourceblock}
\footnotetext[7]{概率语句 P[A ∪B ∪C] 被解释为事件 A、事件 B 或事件 C 发生的概率。概率计算为集合中语句 A ∪B ∪C 为真的所有实现的分数。}
请注意，n 本质上是 x 周围半径为 W 的邻域内所有粒子的体积平均值。 NDF 估计器的统计特性可以根据多粒子联合 PDF 找到。例如，NDF 估计量的平均值为

\begin{sourceblock}
\includegraphics[page=146,pagebox=mediabox,trim={136.81bp 527.30bp 65.07bp 103.57bp},clip,max width=0.94\textwidth,max height=0.38\textheight]{Marchisio2013.pdf}
\end{sourceblock}
因此，当且仅当，估计器才是无偏的

\begin{sourceblock}
\includegraphics[page=146,pagebox=mediabox,trim={158.90bp 483.70bp 65.07bp 187.91bp},clip,max width=0.94\textwidth,max height=0.38\textheight]{Marchisio2013.pdf}
\end{sourceblock}
例如，当所有颗粒均匀分布时，就会发生这种情况。在此限制下，估计器的归一化标准差将像

\begin{sourceblock}
\includegraphics[page=146,pagebox=mediabox,trim={200.04bp 418.94bp 65.07bp 248.00bp},clip,max width=0.94\textwidth,max height=0.38\textheight]{Marchisio2013.pdf}
\end{sourceblock}
其中 V 是域的体积，VW 是样本体积。由于我们必须使 σ <1 才能准确估计，因此系统中的粒子总数 Np 必须大于 V/VW。8 此约束对可用于估计 NDF 的样本量设置了下限。最后，请注意，如果粒子均匀分布，则方程： (4.15) 变为

\begin{sourceblock}
\includegraphics[page=146,pagebox=mediabox,trim={173.83bp 335.48bp 65.07bp 338.01bp},clip,max width=0.94\textwidth,max height=0.38\textheight]{Marchisio2013.pdf}
\end{sourceblock}
它将单粒子联合 PDF 的体积平均值与 NDF 估计器的期望值联系起来。该结果是等式的体积平均模拟。 （4.10）。在实践中，等式。 (4.18) 没有实际用途，因为 f1 未知。然而，它确实表明，至少在统计意义上，体积平均与 NDF 的正式定义是一致的。

\paragraph*{时间平均值}
对于体积平均值不足以估计 NDF 的系统（由于从方程（4.17）中找到的 σ 值太大），我们可以尝试通过定义时间平均值来改进统计数据。理论上，时间平均值应该仅应用于统计上平稳的系统。对于 NDF，这将时间平均限制为 f (n) 1 出现在方程 1 中的系统。 (4.10) 不依赖于 t。然而，实际上，时间平均也可以用于与时间平均间隔Δt相比f(n) 1 变化缓慢的系统。方程中估计器的时间平均值。 (4.14) 的定义为

\begin{sourceblock}
\includegraphics[page=146,pagebox=mediabox,trim={84.59bp 109.36bp 65.07bp 554.23bp},clip,max width=0.94\textwidth,max height=0.38\textheight]{Marchisio2013.pdf}
\end{sourceblock}
8 读者将认识到 Np/V 是系统中粒子的平均数密度。在具有实际意义的颗粒系统中，平均数密度的范围从悬浮在流体中的纳米颗粒系统的极大到大颗粒（例如矿物加工中的矿石颗粒）的非常小。对于前者，体积平均对于估计 NDF 很有用，因为相对于流体流动中的最小尺度，可以选择非常小的 VW。

请注意，就像 \textasciicircum{}n 一样，时间平均值 \textasciicircum{}n 是一个随机变量。因此，它与 NDF 并不相同。相反，它的期望值是

\begin{sourceblock}
\includegraphics[page=147,pagebox=mediabox,trim={123.65bp 574.45bp 67.56bp 89.13bp},clip,max width=0.94\textwidth,max height=0.38\textheight]{Marchisio2013.pdf}
\end{sourceblock}
如果系统统计平稳，则 〈ˆn〉= 〈ˆn〉。此外，如果估计量是无偏的，则 〈ˆn〉= n。正是在这种统计意义上，方程中的时间平均值。 (4.19) 表示 NDF，因此不应假设 \textasciicircum{}n 与 n 是等价的，从而将它们混淆。当应用时间平均值来估计 n 时，时间间隔 Δt 的长度必须选择得足够大，以将估计器的方差降低到可接受的水平。 \textasciicircum{}n 的方差可以从多粒子联合 PDF 中找到。一般来说，其标准化标准差将像

\begin{sourceblock}
\includegraphics[page=147,pagebox=mediabox,trim={193.50bp 448.80bp 67.56bp 216.12bp},clip,max width=0.94\textwidth,max height=0.38\textheight]{Marchisio2013.pdf}
\end{sourceblock}
其中 τW 是随机变量 hW(x −X(n)(t)) 的自相关时间。从物理上讲，τW 是 x 周围半径为 W 的体积中的粒子被新粒子替换之前必须经过多少时间的度量。因此，比率 NΔt = Δt/τW 是估计器时间平均部分中使用的有效样本大小，NW = VWNp/V 是体积平均中使用的有效样本大小。方程中估计器的总体有效样本量。 (4.19) 因此

\begin{sourceblock}
\includegraphics[page=147,pagebox=mediabox,trim={179.50bp 338.27bp 67.56bp 335.27bp},clip,max width=0.94\textwidth,max height=0.38\textheight]{Marchisio2013.pdf}
\end{sourceblock}
并且参数W和Δt的选择必须使得Neff≫1以控制估计误差。两种限制情况，其中方程中定义的估计量。 (4.19) 效果良好如下。

(i) 当由于平均颗粒数密度较大（例如细颗粒和纳米颗粒系统）而 NW ≫1 时，我们可以取 Δt = 0 并仅使用体积平均。

(ii) 当由于（例如）频繁的粒子与粒子碰撞而导致 τW <1 时，我们可以设置 W = 0 并仅使用时间平均。

另一方面，对于具有低粒子数密度和低碰撞频率的系统，估计器 \textasciicircum{}n 将产生较差的 NDF 表示。尽管如此，这并不意味着不能为此类系统定义 NDF。事实上，它仍然由方程式精确定义。 （4.9）。相反，它只是简单地指出，使用颗粒流的单一实现来估计 NDF 将极其困难。这种说法的实际结果是，对于估计器的标准偏差较大的系统，很难验证 GPBE 中项的闭合模型（使用微尺度系统的 DNS 或实验测量）。

\subsection*{4.1.3 流体-粒子系统的 NDF 定义}\phantomsection\addcontentsline{toc}{subsection}{4.1.3 流体-粒子系统的 NDF 定义}
上面导出的 NDF 的定义是在颗粒被真空包围的极限下进行的。在更一般的情况下，当它们被流体包围时，每个粒子将通过作用在其表面上的流体应力而加速。对由此产生的表面力的精确描述非常复杂，并且涉及流体-颗粒耦合系统的完整微观解决方案。9然而，为了定义 NDF，我们不需要知道表面力的确切形式。知道它们存在就足够了，并且通常它们不仅取决于周围的流体，还取决于描述粒子系统的完整随机变量集：粒子位置 \{X(n)\}、粒子速度 \{U(n)

p \}，粒子内坐标 \{ xi(n) p \}。形式上，流体的存在将在方程式中产生一个新的连续加速度项。 (4.2),10

\begin{sourceblock}
\includegraphics[page=148,pagebox=mediabox,trim={194.33bp 531.77bp 65.07bp 139.13bp},clip,max width=0.94\textwidth,max height=0.38\textheight]{Marchisio2013.pdf}
\end{sourceblock}
和粒子内坐标

\begin{sourceblock}
\includegraphics[page=148,pagebox=mediabox,trim={196.48bp 474.74bp 65.07bp 196.16bp},clip,max width=0.94\textwidth,max height=0.38\textheight]{Marchisio2013.pdf}
\end{sourceblock}
可包括相间的传质。在大多数流体-颗粒流的描述中，介观尺度流体-颗粒加速度 A(n) fp 包括以下形式的阻力项

\begin{sourceblock}
\includegraphics[page=148,pagebox=mediabox,trim={163.66bp 420.16bp 65.07bp 255.32bp},clip,max width=0.94\textwidth,max height=0.38\textheight]{Marchisio2013.pdf}
\end{sourceblock}
其中 C(n) D 是阻力系数，U(n) f 是 X(n)(t) 附近的代表性流体速度。11 U(n) f 的正确选择将取决于系统。然而，合适的选择可能是粒子周围动量边界层边缘的流体速度的表面平均值。 12 读者应该清楚，开发 A(n) fp 的介观模型的主要挑战是充分考虑其对其他粒子的依赖性，而不必求解完整的多粒子联合 PDF。对于当前的讨论，我们不需要知道流体速度和特性（例如成分和焓）的介观模型的确切形式。相反，我们将简单地假设 U(n) f 和 xi(n) f 服从 13

\begin{sourceblock}
\includegraphics[page=148,pagebox=mediabox,trim={207.77bp 253.80bp 65.07bp 417.03bp},clip,max width=0.94\textwidth,max height=0.38\textheight]{Marchisio2013.pdf}
\end{sourceblock}
以及

\begin{sourceblock}
\includegraphics[page=148,pagebox=mediabox,trim={209.92bp 198.76bp 65.07bp 472.08bp},clip,max width=0.94\textwidth,max height=0.38\textheight]{Marchisio2013.pdf}
\end{sourceblock}
9 例如，在未被粒子占据的域的体积中，流体可以遵循纳维-斯托克斯方程。 10 流体还可能改变内坐标，因此需要内坐标的附加源项。对于目前的讨论，我们将忽略这种可能性。 11 在气固流的阻力模型中，U(n) f 被称为“第 n 个固体粒子看到的流体速度”。在所有情况下，重要的是要认识到 U(n) f 的统计特性通常与拉格朗日流体粒子的统计特性不同。只有在斯托克斯数为零的情况下，所有固体颗粒都完全跟随流体，统计特性才会相同。例如，对于浸没在粘性流体中的纳米颗粒，这种情况可能会发生。 12 显然，我们不能使用粒子位置 x = X(n)(t) 处的流体速度，因为根据定义，该点没有流体。使用动量边界层定义 U(n)

{\small f 的优点是位置取决于粒子雷诺数。\par}
\footnotetext[13]{与 Peirano \& Minier (2002) 不同，这里的时间导数位于第 n 个粒子之后的拉格朗日参考系中。}
其中，介观尺度粒子流体加速度 A(n) pf 通过动量守恒定律取决于 A(n) fp，而 A(n) f 取决于流体相中的压力、应力和体力。同样，G(n) p 和 C(n) f 和 C(n) fψ 分别通过质量和化学物质守恒定律取决于 G(n) pψ。 14 然后，描述具有 Np 粒子的流体粒子系统所需的随机变量集将被增强，以包括流体速度和流体的附加特性（或流体内坐标）：\{X(n)\}、\{U(n)

\begin{sourceblock}
\includegraphics[page=149,pagebox=mediabox,trim={156.66bp 544.33bp 233.19bp 138.22bp},clip,max width=0.94\textwidth,max height=0.38\textheight]{Marchisio2013.pdf}
\end{sourceblock}
f\}。然而，系综的定义基本保持不变。系综中的单个实现仍然由一组可能的初始条件定义，但现在扩展为包括 \{U(n) f (0)\} 和 \{xi(n) f (0)\}。在实践中，初始微尺度流体速度场U(0,x)可以以与颗粒表面上的零滑移边界条件一致的方式指定。然后可以计算表面项 \{A(n) fp \}，并与等式 1 一起使用。 (4.25) 定义\{U(n) f (0)\}。考虑到无限的实现集合，我们现在可以通过以下方式定义多粒子流体联合 PDF：

\begin{sourceblock}
\includegraphics[page=149,pagebox=mediabox,trim={105.64bp 394.43bp 67.56bp 224.86bp},clip,max width=0.94\textwidth,max height=0.38\textheight]{Marchisio2013.pdf}
\end{sourceblock}
fNp

其中流体速度的状态空间变量为 \{V(n) f \}，流体特性的状态空间变量为 \{η(n) f \}。然后，按照与获得方程式相同的步骤进行操作。 （4.9），我们可以定义流体-粒子NDF：

\begin{sourceblock}
\includegraphics[page=149,pagebox=mediabox,trim={150.99bp 306.10bp 67.56bp 360.18bp},clip,max width=0.94\textwidth,max height=0.38\textheight]{Marchisio2013.pdf}
\end{sourceblock}
其中 vf 是流体速度的相空间变量，Σf 是流体成分。在相同分布粒子的极限下，方程： (4.29) 变为 n(t, x, vp, Ψp, vf, Ψf) = Np f1(t, x, vp, Ψp, vf, Ψf)。因此，我们现在根据明确定义的单粒子联合 PDF，对流体粒子 NDF 有了数学上一致的定义。15 对于实际应用，剩下的主要挑战是为 f (n) 1 的动力学找到适当的介观模型。我们将在第 5 章中讨论这个问题。为了理解流体粒子 NDF 中包含的物理范围，考虑两个极限情况是有启发性的。

(i) 在粒子与流体完全解耦的极限下，流体-粒子 NDF 是可分离的： n(t, x, vp, Ψp, vf, Ψf) = n(t, x, vp, Ψp) ff(t, x, vf, Ψf)，其中 ff(t, x, vf, Ψf) 是由 ff(t, x, vf, ψf) = 〈δ(Uf(t, x) −vf)δ(φf(t, x) − ψf)〉，其中 Uf(t, x) 和 φf(t, x) 是流体速度和特征属性场。用于计算期望值的系综
14 G(n) f 模型将因粒子和流体速度之间的差异、相间的传质以及流体相中的化学反应而产生贡献。我们将在这里重点关注前两个贡献。 C(n) pxi 项主要表示粒子数的不连续变化，这是由于
{\small 颗粒成核，其中流体相中相应的质量变化为 C(n)\par}
{\small f ψ。 15将该结果扩展到 Np 由于聚集和破坏等原因而随时间变化的情况，这一点很简单。人们必须简单地以 Np 的值为条件，并包括改变系统中粒子数量的过程的动态描述。\par}
右侧是流体流动所有可能实现的集合16（在没有与颗粒相互作用的情况下）。当方程式中的阻力系数时，就会出现这种限制。 (4.25) 为空，在求解流体速度时可以忽略粒子的存在（即单向耦合）。

(ii) 在颗粒均具有与流体相同的速度且流体流动不受颗粒存在影响的情况下，17 流体-颗粒 NDF 变为

\begin{sourceblock}
\includegraphics[page=150,pagebox=mediabox,trim={122.13bp 475.10bp 65.07bp 155.76bp},clip,max width=0.94\textwidth,max height=0.38\textheight]{Marchisio2013.pdf}
\end{sourceblock}
其中 f (n) 1 (t, x, ψp|vf) 是单粒子条件 PDF，假设 U(n) p = vf。因此，一般的NDF可以表示为

\begin{sourceblock}
\includegraphics[page=150,pagebox=mediabox,trim={168.18bp 404.16bp 65.07bp 262.12bp},clip,max width=0.94\textwidth,max height=0.38\textheight]{Marchisio2013.pdf}
\end{sourceblock}
或者，对于相同分布的粒子，n(t, x, vf, Ψp) = Np f1(t, x, Ψp|vf) ff(t, x, vf, Ψf)。换句话说，假设所有粒子都具有与流体相同的速度，则 Np f1(t, x, ϋp|vf) 项是时间 t 和点 x 处内坐标的 NDF。当方程中的阻力系数为非常稀的系统时，就会出现这种限制。 (4.25) 是无限的并且不存在质量传递。

这两个极限之间的流体-粒子 NDF 的形式将取决于表面力模型和间隙流体动力学的确切形式。因此，从理论上讲，精确计算 NDF 的唯一方法是使用直接数值模拟 (DNS) 来模拟完整的微尺度流体-颗粒系统。然而，通过为 NDF 输运方程（即 GPBE）中出现的多粒子统计提供介观尺度闭包，可以获得 NDF 的有用近似值。总之，我们已经证明，通过采用明确定义的多粒子联合 PDF，可以直接定义各种形式的 NDF。同样，我们还展示了在某些条件下如何使用体积和时间平均值来根据流量的单个实现来估计 NDF。现在我们转向更困难的任务，即从定义粒子位置、速度、内坐标和表面力的微分方程出发，推导描述 NDF 演化的 GPBE。这个 GPBE 将有未封闭的术语，因此我们将在第 5 章中回顾各种介观尺度模型。

\section*{4.2 从多粒子流体联合PDF到GPBE}\phantomsection\addcontentsline{toc}{section}{4.2 从多粒子流体联合PDF到GPBE}
在本节中，我们从多粒子流体 PDF 的输运方程出发，对流体粒子 NDF 的 GPBE 进行简化推导。这需要三个步骤。

\footnotetext[16]{在层流中，ff 简化为δ 函数。在湍流中，可以使用 PDF 方法对 ff 进行建模。 17因此，根据定义，相之间不存在质量传递。}
1. 从粒子特性的介观尺度模型（方程（4.1）、（4.3）、（4.23）和（4.26））开始，我们首先推导出 fNp 的广义输运方程。

2. 给定 fNp 的输运方程，我们将通过积分除了与第 n 个粒子相关的自由度之外的所有自由度来找到 f (n) 1 的输运方程。由于信息丢失，此步骤将生成几个必须建模的未封闭项。

3. 最后，使用方程中给出的流体-粒子 NDF 的定义。 （4.29），我们将推导出GPBE。

熟悉气体动力学动力学理论的读者应该会发现与玻尔兹曼动力学方程的推导有许多相似之处。

\subsection*{4.2.1 多粒子联合PDF的输运方程}\phantomsection\addcontentsline{toc}{subsection}{4.2.1 多粒子联合PDF的输运方程}
多粒子联合 PDF 输运方程的一般形式（刘维尔方程的推广）为

\begin{sourceblock}
\includegraphics[page=151,pagebox=mediabox,trim={76.18bp 371.37bp 67.56bp 265.79bp},clip,max width=0.94\textwidth,max height=0.38\textheight]{Marchisio2013.pdf}
\end{sourceblock}
请注意，该表达式的左侧实际上是 fNp 的连续性方程，它表明多粒子联合 PDF 沿着相空间中的轨迹是恒定的。等式右侧的项。 (4.32) 的贡献在于 Np 粒子碰撞算子，它会产生粒子速度 U(n) p 和内坐标 ψ(n) p 的不连续变化，以及粒子成核或蒸发。左侧第一项是 fNp 的累积。左侧的其余项表示粒子位置空间 x(n)、粒子速度空间 V(n) p 、粒子内坐标空间 η(n) p 、流体速度空间 V(n) f 和流体内坐标空间 η(n) f 的连续变化。为了简化符号，我们以可扩展的紧凑形式编写了导数。例如，

\begin{sourceblock}
\includegraphics[page=151,pagebox=mediabox,trim={159.06bp 200.41bp 67.56bp 465.88bp},clip,max width=0.94\textwidth,max height=0.38\textheight]{Marchisio2013.pdf}
\end{sourceblock}
方程中出现的自变量总数。因此，方程（4.32）相当大，实际上对于实际应用来说太大了。然而，如前所述，通过耦合方程。 （4.32）用纳维-斯托克斯方程求出流体作用在粒子上的力，Np-粒子系统就完全确定了。尽管没有明确写出，但读者应该记住，相空间通量和碰撞项的介观模型取决于完整的自变量集。例如，表面项取决于所有状态变量：A(n)

\begin{sourceblock}
\includegraphics[page=151,pagebox=mediabox,trim={265.98bp 109.24bp 94.52bp 572.94bp},clip,max width=0.94\textwidth,max height=0.38\textheight]{Marchisio2013.pdf}
\end{sourceblock}
f\}）。确定这些函数的唯一已知方法是使用所有可能的初始条件集对微尺度流体-粒子系统进行直接数值模拟。显然，这种方法是棘手的。因此，我们减少了自变量的数量，并引入了试图捕捉多粒子相互作用的“平均”效应的介观模型。

\subsection*{4.2.2 单粒子联合PDF的输运方程}\phantomsection\addcontentsline{toc}{subsection}{4.2.2 单粒子联合PDF的输运方程}
流体-粒子系统最简单有用的表示是单粒子联合 PDF f (n) 1，通过积分除与第 n 个粒子相关的变量之外的所有自变量来定义：

\begin{sourceblock}
\includegraphics[page=152,pagebox=mediabox,trim={146.82bp 550.28bp 65.07bp 115.99bp},clip,max width=0.94\textwidth,max height=0.38\textheight]{Marchisio2013.pdf}
\end{sourceblock}
从等式开始(4.32) 并积分除与第 n 个粒子产量相关的变量之外的所有变量

\begin{sourceblock}
\includegraphics[page=152,pagebox=mediabox,trim={78.67bp 446.12bp 65.07bp 190.22bp},clip,max width=0.94\textwidth,max height=0.38\textheight]{Marchisio2013.pdf}
\end{sourceblock}
我们引入了单粒子碰撞算子，定义为

\begin{sourceblock}
\includegraphics[page=152,pagebox=mediabox,trim={145.11bp 394.50bp 65.07bp 271.78bp},clip,max width=0.94\textwidth,max height=0.38\textheight]{Marchisio2013.pdf}
\end{sourceblock}
以及单粒子条件期望值，例如，定义为

\begin{sourceblock}
\includegraphics[page=152,pagebox=mediabox,trim={121.87bp 341.74bp 65.07bp 323.82bp},clip,max width=0.94\textwidth,max height=0.38\textheight]{Marchisio2013.pdf}
\end{sourceblock}
与 A(n) fp 相比，该表达式的左侧取决于一个小得多的集合

\begin{sourceblock}
\includegraphics[page=152,pagebox=mediabox,trim={169.53bp 300.13bp 212.77bp 382.42bp},clip,max width=0.94\textwidth,max height=0.38\textheight]{Marchisio2013.pdf}
\end{sourceblock}
自变量： 〈A(n) f )。因此，原则上，它应该比多粒子对应物更容易处理。然而，由于粒子间相互作用的直接信息丢失，单粒子条件期望值必须用介观尺度闭包表示。开发闭包，通过求解方程式来产生 f (n) 1 的准确预测。 （4.35）是一项重要的任务，介观尺度闭合的形式将高度依赖于系统的物理特性。读者还应该清楚地区分单粒子条件期望值，例如方程： （4.37）和等式的右侧。 (4.1)-(4.3)。后者可能涉及所有粒子位置、速度和内坐标，而前者仅取决于第 n 个粒子的介观变量。在文献中，离散元方法 (DEM) 经常声称可以解决方程 1 定义的微观问题。 (4.1)–(4.3)，但这只有当粒子在真空中运动时才成立。添加流体后，改变一个粒子的位置/速度可能会影响系统中的所有其他粒子。因此，将粒子上的流体力建模为仅取决于粒子速度的阻力项的 DEM 模拟与方程 1 的微尺度 DNS 不同。 （4.1）-（4.3）和流体速度场。

\subsection*{4.2.3 NDF 的输运方程}\phantomsection\addcontentsline{toc}{subsection}{4.2.3 NDF 的输运方程}
{\small NDF 由等式定义。 （4.29），并根据 GPBE 演化。从等式开始（4.35），我们首先将自变量重写为（见方程（4.29）的右侧）\par}
\begin{sourceblock}
\includegraphics[page=153,pagebox=mediabox,trim={76.18bp 584.35bp 67.56bp 58.32bp},clip,max width=0.94\textwidth,max height=0.38\textheight]{Marchisio2013.pdf}
\end{sourceblock}
其中方程中的所有项。现在在相空间位置 (x, vp, Ψp, vf, Ψf) 评估 (4.38)。对所有粒子求和即可得出 GPBE：

\begin{sourceblock}
\includegraphics[page=153,pagebox=mediabox,trim={76.18bp 481.07bp 67.56bp 162.43bp},clip,max width=0.94\textwidth,max height=0.38\textheight]{Marchisio2013.pdf}
\end{sourceblock}
为了考虑可变的粒子数，我们将碰撞项 S1 概括为包括由于成核、聚集和破碎而导致的 Np 变化。这些过程还需要模型才能结束方程。 （4.39）。该方程可以与方程进行比较。 （2.16）（第37页），可以看出它们具有相同的一般形式。然而，现在很明显，只有为条件相空间速度（〈Afp〉1、〈Ap〉1、〈Gp〉1、〈Apf〉1、〈Af〉1 和〈Gf〉1）和不连续“源”项 S1 提供介观尺度闭包，才能求解 GPBE。请注意，我们在方程 1 中去掉了条件相空间速度上的上标 (n)。 （4.39）。形式上，这意味着（例如）〈Afp〉1 的定义是

\begin{sourceblock}
\includegraphics[page=153,pagebox=mediabox,trim={103.91bp 313.61bp 106.44bp 345.67bp},clip,max width=0.94\textwidth,max height=0.38\textheight]{Marchisio2013.pdf}
\end{sourceblock}
或者，如果有人认为所有粒子在统计上都是相同的，那么 〈A(n) fp 〉1 和 f (n) 1 不应该依赖于 n。尽管这个假设对于推导方程不是必需的。 （4.39），通常调用它来简化物理建模。换句话说，如果所有粒子在统计上都是相同的，则可以通过考虑来自 DNS 的随机选择粒子的单粒子统计来验证介观尺度模型。

\subsection*{4.2.4 闭包问题}\phantomsection\addcontentsline{toc}{subsection}{4.2.4 闭包问题}
正如第一章所述，在本书中，我们将考虑以下情况的计算模型： (4.39) 以封闭形式出现。特别是，这意味着 GPBE 中的所有相空间通量和源项（〈Afp〉1、〈Ap〉1、〈Gp〉1、〈Apf〉1、〈Af〉1、〈Gf〉1 和 S1）将仅取决于自变量（t、x、vp、ψp、vf 和Ψf) 和 NDF n(t, x, vp, Ψp, vf, Ψf)。然而，我们应该注意到，为这些项找到准确的介观尺度闭包是一项非常重要的任务。一般来说，这样的闭包将是非线性的，并且可能涉及卷积积分（例如，描述聚合和碰撞过程）。因此，方程的精确数值解。 （4.39）也将是不平凡的，并且将取决于介观尺度闭合的特定形式。求解方程时的另一个可怕的复杂性。 (4.39) 的事实是 NDF 通常取决于大量自变量。为了降低问题的维数，有时可以通过求解输运方程来处理 NDF 的选定矩。一般来说，矩输运方程不会闭合，需要进一步的建模近似。在下一节中，我们从方程开始推导最常用的矩输运方程。 （4.39）。

\section*{4.3 矩输运方程}\phantomsection\addcontentsline{toc}{section}{4.3 矩输运方程}
在本节中，我们将从方程 1 开始推导出 NDF 选定矩的几个示例传输方程。 （4.39）。因为条件相空间速度出现在方程中。由于方程（4.39）尚未被介观模型所取代，我们不会详细讨论如何使用矩闭包在宏观尺度上逼近这些项。然而，熟悉气体动力学理论的读者（Chapman 和 Cowling，1961）会认识到，开发当前的宏观尺度闭合是一个重大挑战。这里的目标比较温和。我们主要感兴趣的是从数学角度理解从 GPBE 出发寻找矩输运方程的过程。然而，正如后面章节中详细讨论的那样，我们的首要目标是说明封闭的 GPBE 是开发多分散型多相流模型的正确起点。在这种情况下，读者必须了解 GPBE 中的矩输运方程是如何产生的，这一点很重要。当我们继续推导矩输运方程时，我们还将讨论流体相的宏观模型和粒子所看到的流体的介观模型（即〈Apf〉1、〈Af〉1和〈Gf〉1）之间的一致性问题。这一讨论将产生一组一致性约束，这些约束采用相空间变量上的流体相介观尺度模型积分的形式。作为示例，我们提出了一些满足一致性约束的最小一致性模型。然而，由于这些模型的设计只是为了满足一致性约束（例如质量、物种、动量、能量等守恒），因此它们并不排除更复杂的基于物理的模型。作为一般规则，由于感兴趣的矩是 GPBE 的线性（积分）变换，因此生成的输运方程处于保守形式（即，对于封闭系统，具有累积和空间通量项，其空间体积积分为零）。正如后面章节中将讨论的，保守形式的输运方程非常适合有限体积方法（Leveque，2002），特别是对于可以生成不连续解或冲击的矩系统。

\subsection*{4.3.1 关于相空间积分的一些说明}\phantomsection\addcontentsline{toc}{subsection}{4.3.1 关于相空间积分的一些说明}
从方程计算矩。 （4.39）基本上涉及对 NDF 中的一些自变量（即 vp、 Ψp、vf 和 Ψf）进行积分，同时保持其他变量固定（即 t 和 x）。为了促进这项任务，回顾一些基本的集成规则是很有用的。

\paragraph*{(i) 时间和空间导数的处理。设 g(vp, Ψp, vf, Ψf) 为任意值}
其变量的函数。然后

\begin{sourceblock}
\includegraphics[page=154,pagebox=mediabox,trim={114.93bp 55.70bp 65.07bp 555.14bp},clip,max width=0.94\textwidth,max height=0.38\textheight]{Marchisio2013.pdf}
\end{sourceblock}
(ii) 速度导数的处理。对于速度导数，有必要使用分部积分。例如，

\begin{sourceblock}
\includegraphics[page=155,pagebox=mediabox,trim={151.92bp 524.24bp 67.56bp 87.48bp},clip,max width=0.94\textwidth,max height=0.38\textheight]{Marchisio2013.pdf}
\end{sourceblock}
该方程右侧第一项是全导数，因此可以对其进行形式积分以找到

\begin{sourceblock}
\includegraphics[page=155,pagebox=mediabox,trim={100.43bp 455.37bp 76.40bp 211.28bp},clip,max width=0.94\textwidth,max height=0.38\textheight]{Marchisio2013.pdf}
\end{sourceblock}
(4.42) 其中第二项对应于垂直于限制速度相空间的表面的通量。由于速度相空间延伸至无穷大，因此该通量必须为零。因此，部分速度变量（即 vp 和 vf）在积分过程中由于通量而产生的贡献被设置为零。

(iii) 内坐标导数的处理。对于内坐标导数，再次需要使用分部积分：18

\begin{sourceblock}
\includegraphics[page=155,pagebox=mediabox,trim={101.56bp 297.13bp 67.56bp 342.06bp},clip,max width=0.94\textwidth,max height=0.38\textheight]{Marchisio2013.pdf}
\end{sourceblock}
该方程右侧第一项是全导数，因此可以形式积分。然而，对于内坐标，相空间通常不会延伸到无穷大。为了清楚地看到会发生什么，请考虑具有以零和无穷大为界的单个内部坐标的情况：

\begin{sourceblock}
\includegraphics[page=155,pagebox=mediabox,trim={113.34bp 204.35bp 67.56bp 462.30bp},clip,max width=0.94\textwidth,max height=0.38\textheight]{Marchisio2013.pdf}
\end{sourceblock}
右侧的两项是无穷远处的通量（我们可以安全地将其设置为零）和原点处的通量。根据 〈Gp〉1 和 g 的形式，原点处的通量不必为零。例如，如果当 ψp = 0 时 g 不为零，则在 ψp = 0 时计算的乘积 〈Gp〉1n 必须为零，以便抵消通量项。不幸的是，在一些重要的应用中，通量项不为零，因此必须注意在推导矩输运方程时如何处理通量项。例如，如果 ψp 代表蒸发液滴的表面，并且 〈Gp〉1 是常数（即蒸发速率与表面积成正比），那么当 ψp = 0 时，n 将不为零。物理上，非零通量是由于蒸发导致液滴消失，因此不能忽略。

\footnotetext[18]{Σf 也有类似的表达式。}
总之，从方程开始计算矩输运方程。 (4.39) 涉及使用上述规则对 g 的特定选择在相空间上进行积分。下面，我们假设相空间边界处的通量项可以忽略不计。然而，读者应该记住，必须针对特定情况验证这一假设。

\subsection*{4.3.2 分散相数传输}\phantomsection\addcontentsline{toc}{subsection}{4.3.2 分散相数传输}
总数浓度 N(t, x) 对应于 NDF 的零阶矩（即 g = 1），定义为

\begin{sourceblock}
\includegraphics[page=156,pagebox=mediabox,trim={197.26bp 493.73bp 65.07bp 179.05bp},clip,max width=0.94\textwidth,max height=0.38\textheight]{Marchisio2013.pdf}
\end{sourceblock}
其输运方程可由式（1）求出： （4.39）：

\begin{sourceblock}
\includegraphics[page=156,pagebox=mediabox,trim={198.71bp 439.08bp 65.07bp 235.27bp},clip,max width=0.94\textwidth,max height=0.38\textheight]{Marchisio2013.pdf}
\end{sourceblock}
其中我们明确假设相空间边界处的通量为零。数均分散相速度定义为

\begin{sourceblock}
\includegraphics[page=156,pagebox=mediabox,trim={184.14bp 379.23bp 65.07bp 288.31bp},clip,max width=0.94\textwidth,max height=0.38\textheight]{Marchisio2013.pdf}
\end{sourceblock}
和粒子数源项

\begin{sourceblock}
\includegraphics[page=156,pagebox=mediabox,trim={191.72bp 335.14bp 65.07bp 337.64bp},clip,max width=0.94\textwidth,max height=0.38\textheight]{Marchisio2013.pdf}
\end{sourceblock}
请注意，源项的符号取决于粒子是在系统中创建还是销毁。还要注意方程中的空间传输项。 (4.46) 通常不会闭合，除非所有粒子都具有相同的速度。方程中的传输方程。 （4.46）主要用于具有颗粒聚集和破碎的系统（即当N（t，x）不恒定时）。在这种情况下，它通常会耦合到涉及高阶矩的矩输运方程组。

\subsection*{4.3.3 分散相体积传输}\phantomsection\addcontentsline{toc}{subsection}{4.3.3 分散相体积传输}
为了导出分散相体积分数 αp 的输运方程，我们让第一个内坐标 ψp1 等于颗粒体积 (Vp)。然后分散相体积分数定义为

\begin{sourceblock}
\includegraphics[page=156,pagebox=mediabox,trim={189.89bp 155.24bp 65.07bp 517.54bp},clip,max width=0.94\textwidth,max height=0.38\textheight]{Marchisio2013.pdf}
\end{sourceblock}
从等式开始(4.39) 当 g = ϋp1 时，αp 的输运方程为

\begin{sourceblock}
\includegraphics[page=156,pagebox=mediabox,trim={178.33bp 98.97bp 65.07bp 572.88bp},clip,max width=0.94\textwidth,max height=0.38\textheight]{Marchisio2013.pdf}
\end{sourceblock}
{\small 其中体积平均分散相速度定义为\par}
\begin{sourceblock}
\includegraphics[page=156,pagebox=mediabox,trim={177.02bp 50.19bp 65.07bp 616.47bp},clip,max width=0.94\textwidth,max height=0.38\textheight]{Marchisio2013.pdf}
\end{sourceblock}
对于所有粒子具有相同体积的情况，UV = UN。请注意，与 UN 一样，体积平均分散相速度通常不会采用封闭形式。粒子体积源项定义为

\begin{sourceblock}
\includegraphics[page=157,pagebox=mediabox,trim={175.03bp 575.60bp 67.56bp 97.18bp},clip,max width=0.94\textwidth,max height=0.38\textheight]{Marchisio2013.pdf}
\end{sourceblock}
\begin{sourceblock}
\includegraphics[page=157,pagebox=mediabox,trim={183.32bp 535.64bp 67.56bp 137.14bp},clip,max width=0.94\textwidth,max height=0.38\textheight]{Marchisio2013.pdf}
\end{sourceblock}
以及

Gp V 的连续贡献可能是由于从流体到固相的传质。另一方面，不连续项 SV 可能是由于流体相中体积非零的颗粒成核而出现的。对于粒子体积守恒的系统，方程的右侧： (4.50) 将为空。最后，我们应该注意，根据定义，分散相体积分数不能大于 1 (0 ≤αp ≤1)。19 这一约束意味着当 αp = 1 时，UV 的发散度必须为非负。在二流体模型中，通常通过在 UV 传输方程中添加压力项来强制执行此条件。方程中的分散相体积分数方程。当所有颗粒的分散相材料密度相同时，使用方程（4.49）。换句话说，密度 ρp 不是内坐标，而是颗粒体积变化，例如由于聚结和破碎而变化。对于这种情况，体积平均速度与质量平均速度相同。否则，使用粒子质量找到平均分散相速度的更自然选择。

\subsection*{4.3.4 流体相体积输运}\phantomsection\addcontentsline{toc}{subsection}{4.3.4 流体相体积输运}
为了推导流体相体积分数 αf 的输运方程，我们让内坐标 Σf1 等于粒子看到的流体体积。20 然后流体相体积分数定义为

\begin{sourceblock}
\includegraphics[page=157,pagebox=mediabox,trim={189.64bp 273.10bp 67.56bp 399.68bp},clip,max width=0.94\textwidth,max height=0.38\textheight]{Marchisio2013.pdf}
\end{sourceblock}
和体积平均流体速度

\begin{sourceblock}
\includegraphics[page=157,pagebox=mediabox,trim={174.77bp 215.90bp 67.56bp 450.76bp},clip,max width=0.94\textwidth,max height=0.38\textheight]{Marchisio2013.pdf}
\end{sourceblock}
从等式开始(4.39) 当 g = Σf1 时，αf 的输运方程为

\begin{sourceblock}
\includegraphics[page=157,pagebox=mediabox,trim={174.16bp 165.61bp 67.56bp 506.24bp},clip,max width=0.94\textwidth,max height=0.38\textheight]{Marchisio2013.pdf}
\end{sourceblock}
19 事实上，由于颗粒的形状通常不允许系统完全被分散相占据，因此 αp 的上限通常会小于 1.20 我们隐含地假设每个粒子看到的流体体积是不同的，并且所有这些体积的并集等于总流体体积。由于 n 代表每单位体积的颗粒数，所以 ϋf1n 是每单位体积的流体体积。原则上，如果所有粒子都看到相同的流体体积，则 Σf1 可以是常数。然而，我们将看到，如果系统在空间上不均匀，则 xif1 一般不能恒定。直观上，该事实与当分散相体积分数αp增加（减少）时颗粒周围的流体体积减少（增加）的图像是一致的。最简单的情况是 ψf1 = g(t, x, ψp1)，其中 g 是满足式（1）的函数。 （4.54）。在这种情况下，NDF 可写为 n(t, x, vp, Ψp1, vf, Ψf1) = n(t, x, vp, Ψp1, vf)δ(Ψf1 − g(t, x, Ψp1))。如果我们选择线性函数 g(t, x, Σp1) = αfΣp1/αp，则每个粒子都有与其体积 Σp1 成正比的流体体积。

其中

\begin{sourceblock}
\includegraphics[page=158,pagebox=mediabox,trim={176.27bp 602.05bp 65.07bp 64.60bp},clip,max width=0.94\textwidth,max height=0.38\textheight]{Marchisio2013.pdf}
\end{sourceblock}
流体相体积分数源项定义为

\begin{sourceblock}
\includegraphics[page=158,pagebox=mediabox,trim={178.51bp 555.25bp 65.07bp 117.54bp},clip,max width=0.94\textwidth,max height=0.38\textheight]{Marchisio2013.pdf}
\end{sourceblock}
以及

\begin{sourceblock}
\includegraphics[page=158,pagebox=mediabox,trim={185.23bp 510.34bp 65.07bp 162.44bp},clip,max width=0.94\textwidth,max height=0.38\textheight]{Marchisio2013.pdf}
\end{sourceblock}
根据定义，αp + αf = 1，并且在相间不存在传质的情况下，方程的右侧为： (4.50) 为零。在此限制下，方程。 (4.56) 应减少为

\begin{sourceblock}
\includegraphics[page=158,pagebox=mediabox,trim={199.62bp 440.13bp 65.07bp 234.22bp},clip,max width=0.94\textwidth,max height=0.38\textheight]{Marchisio2013.pdf}
\end{sourceblock}
它简单地表明流体相体积分数是一个守恒量（分散相体积分数也是如此）。

\paragraph*{颗粒所见流体体积变化的一致模型}
为了方程。为了在没有传质的情况下等价（4.56）和（4.60），我们必须有

\begin{sourceblock}
\includegraphics[page=158,pagebox=mediabox,trim={183.82bp 323.06bp 65.07bp 348.79bp},clip,max width=0.94\textwidth,max height=0.38\textheight]{Marchisio2013.pdf}
\end{sourceblock}
因此，方程 1 中的 〈Gf1〉1 必须非零。 （4.58）；或者，换句话说，粒子看到的流体体积必须沿着粒子轨迹变化。从物理上讲，这种变化与其他粒子的存在相关（即不能使用单粒子物理学对其进行建模）。 〈Gf1〉1 的简单一致模型可以写为

\begin{sourceblock}
\includegraphics[page=158,pagebox=mediabox,trim={149.77bp 233.72bp 65.07bp 432.44bp},clip,max width=0.94\textwidth,max height=0.38\textheight]{Marchisio2013.pdf}
\end{sourceblock}
其中 τpf 是 ψf1 松弛回到其条件平均值 αfψp1/αp 的时间尺度。等式右侧的第一项。 （4.62）指出，粒子和流体之间相对速度方向上的 NDF 空间梯度会导致粒子看到的流体体积的变化。使用 NDF 梯度来解释有限体积粒子并不罕见。 （对于点粒子，αf = 1，因此 Gf V = 0。）例如，NDF 的空间梯度出现在第 6 章中硬球碰撞的 Boltzmann-Enskog 方程中（例如，参见第 235 页的方程（6.71））。最后，请注意，使用等式。 （4.62），即使包括传质，也没有必要知道 UpV。

\subsection*{4.3.5 分散相质量输运}\phantomsection\addcontentsline{toc}{subsection}{4.3.5 分散相质量输运}
{\small 如果分散相具有不同体积和不同质量（或材料密度）的颗粒，则至少需要两个内坐标来描述颗粒状态。为了推导分散相质量密度 ρp 的输运方程，我们让\par}
内坐标 ψp2 等于颗粒质量 (Mp)。21 分散相质量密度定义为

\begin{sourceblock}
\includegraphics[page=159,pagebox=mediabox,trim={187.88bp 596.63bp 67.56bp 76.15bp},clip,max width=0.94\textwidth,max height=0.38\textheight]{Marchisio2013.pdf}
\end{sourceblock}
从等式开始(4.39) 当 g = ϋp2 时，ρp 的输运方程为

\begin{sourceblock}
\includegraphics[page=159,pagebox=mediabox,trim={178.38bp 542.73bp 67.56bp 129.13bp},clip,max width=0.94\textwidth,max height=0.38\textheight]{Marchisio2013.pdf}
\end{sourceblock}
其中质量平均分散相速度定义为

\begin{sourceblock}
\includegraphics[page=159,pagebox=mediabox,trim={174.42bp 492.82bp 67.56bp 173.83bp},clip,max width=0.94\textwidth,max height=0.38\textheight]{Marchisio2013.pdf}
\end{sourceblock}
对于所有粒子质量相同的情况，UM = UN；对于所有粒子具有相同材料密度的情况，UM = UV。请注意，与 UN 和 UV 一样，质量平均分散相速度通常不会采用封闭形式。在大多数应用中，平均粒子速度 Up 将设置为质量平均分散相速度。因此，除非另有说明，我们将在本书的其余部分中设置 Up = UM。粒子质量源项定义为

\begin{sourceblock}
\includegraphics[page=159,pagebox=mediabox,trim={177.79bp 387.31bp 67.56bp 285.48bp},clip,max width=0.94\textwidth,max height=0.38\textheight]{Marchisio2013.pdf}
\end{sourceblock}
\begin{sourceblock}
\includegraphics[page=159,pagebox=mediabox,trim={182.74bp 347.67bp 67.56bp 325.11bp},clip,max width=0.94\textwidth,max height=0.38\textheight]{Marchisio2013.pdf}
\end{sourceblock}
以及

由于在传输过程中质量守恒，因此 Gp 的连续贡献是由于从流体到固相的质量传递所致。同样，由于流体相中质量非零的颗粒成核，可能会出现不连续项 SM。对于分散相和连续相之间没有传质的系统，方程的右侧为： (4.64) 将为空。

\paragraph*{流体相连续性方程}
由于流体-颗粒系统的总质量守恒，我们可以定义流体相质量密度 ρf

\begin{sourceblock}
\includegraphics[page=159,pagebox=mediabox,trim={178.02bp 198.71bp 67.56bp 475.64bp},clip,max width=0.94\textwidth,max height=0.38\textheight]{Marchisio2013.pdf}
\end{sourceblock}
其中质量平均流体速度 Uf 由总动量守恒方程定义，如下所述。请注意，ρp 和 ρf 分别不等于实际材料密度 ρp 和 ρf。相反，它们等于流体-颗粒系统每单位体积的固体（流体）质量。为了区分这两种密度，当材料密度恒定时，我们可以将它们与体积分数 αp 和 αf 联系起来：ρpαp = ρp 和 ρfαf = ρf。这些关系一般不成立。如果材料密度因粒子而异，则 ρf 取决于内坐标，但 αp 和 ρp 则不然。

\footnotetext[21]{颗粒质量可写为 Mp = ρpVp，其中 ρp 是颗粒的材料密度，Vp 是颗粒的体积。请注意，除了质量之外，材料密度或体积也可以包含在内坐标向量中。 （一般来说，我们会使用质量和体积作为内坐标。）因此，例如，将粒子质量固定为相等并不意味着所有粒子都具有相同的体积。}
\paragraph*{混合物连续性方程}
一般情况下，混合物质量密度可以由 ρmix ≡ρp + ρf 定义，并且满足混合物连续性方程：

\begin{sourceblock}
\includegraphics[page=160,pagebox=mediabox,trim={190.63bp 565.54bp 65.07bp 108.81bp},clip,max width=0.94\textwidth,max height=0.38\textheight]{Marchisio2013.pdf}
\end{sourceblock}
其中质量平均混合物速度定义为

\begin{sourceblock}
\includegraphics[page=160,pagebox=mediabox,trim={191.12bp 519.14bp 65.07bp 151.83bp},clip,max width=0.94\textwidth,max height=0.38\textheight]{Marchisio2013.pdf}
\end{sourceblock}
ρp 和 ρf 的输运方程用于多相流的二流体模型。通过将具有相同内部坐标的所有粒子视为“流体”，可以开发 k 流体模型。因此，例如，如果所有粒子都是相同的，除了一些粒子的质量为 M1，另一些粒子的质量为 M2（这意味着它们具有不同的固体密度），那么我们可以将粒子相视为两种“流体”。从数学上讲，这直接来自本例的 NDF 形式。

\paragraph*{体积和质量恒定}
因为在第 5 章中开发数值方法时它会很有用，所以接下来我们将考虑 Ψp1 = Vp 和 Ψp2 = Mp 恒定的情况。这种情况对应于所有具有相同体积和质量的粒子，因此唯一的相空间变量是速度 vp 和 vf。从等式开始（4.63），我们发现对于这种情况，ρp = MpN = ρpαp。同样，平均粒子速度都相等，因此我们可以将它们设置为等于平均粒子速度：Up ≡UM = UV = UN。在这些条件下，方程。 (4.64) 变为

\begin{sourceblock}
\includegraphics[page=160,pagebox=mediabox,trim={200.51bp 302.48bp 65.07bp 369.36bp},clip,max width=0.94\textwidth,max height=0.38\textheight]{Marchisio2013.pdf}
\end{sourceblock}
请注意，为了与所有粒子具有相同体积和质量的假设保持一致，该表达式的右侧为空。如果流体的材料密度也恒定（即流体不可压缩），则

\begin{sourceblock}
\includegraphics[page=160,pagebox=mediabox,trim={201.97bp 235.66bp 65.07bp 438.70bp},clip,max width=0.94\textwidth,max height=0.38\textheight]{Marchisio2013.pdf}
\end{sourceblock}
将等式加在一起。 (4.71) 和(4.72) 给出了速度场的可实现性约束，即∇x·Uvol = 0，其中Uvol equationαpUp + αfUf。如前所述，必须将此约束纳入分散相动量输运方程中的条件源项中。请注意，一般来说，Uvol Umix 除非流体和颗粒具有相同的材料密度。

\paragraph*{多组分分散相}
在许多应用中，颗粒将由多种化学物质组成。在这种情况下，有必要引入一个内坐标向量 ψp，其分量是每种化学物质的质量。显然，这些内坐标之和等于粒子质量。根据定义，如果 ψpα 是组分 α 的质量，则相空间积分可得出组分分散相质量密度：

\begin{sourceblock}
\includegraphics[page=160,pagebox=mediabox,trim={187.97bp 54.90bp 65.07bp 617.88bp},clip,max width=0.94\textwidth,max height=0.38\textheight]{Marchisio2013.pdf}
\end{sourceblock}
同样，由于所有粒子分量都随粒子移动，因此 ρpα 的输运方程具有以下形式

\begin{sourceblock}
\includegraphics[page=161,pagebox=mediabox,trim={174.31bp 589.67bp 67.56bp 82.17bp},clip,max width=0.94\textwidth,max height=0.38\textheight]{Marchisio2013.pdf}
\end{sourceblock}
其中右侧的传质项和成核项对于每个组件可能不同。求和方程(4.74) 对于所有分量 α 都可得出方程 (4.74)。 （4.64）。

\subsection*{4.3.6 流体相质量输运}\phantomsection\addcontentsline{toc}{subsection}{4.3.6 流体相质量输运}
在有关湍流两相流的文献中（Minier \& Peirano，2001；Peirano \& Minier，2002；Simonin 等人，1993），通常使用单独的分布函数来处理流体相，该函数在相空间上的积分可得出流体相质量密度。在这里，我们使用从 n(t, x, vp, Ψp, vf, Ψf) 开始的不同方法。在这种方法中，我们让内坐标 xif2 等于粒子看到的流体质量。22 流体相质量密度由下式给出

\begin{sourceblock}
\includegraphics[page=161,pagebox=mediabox,trim={188.86bp 440.55bp 67.56bp 232.23bp},clip,max width=0.94\textwidth,max height=0.38\textheight]{Marchisio2013.pdf}
\end{sourceblock}
我们还定义了相应的流体质量平均分散相速度

\begin{sourceblock}
\includegraphics[page=161,pagebox=mediabox,trim={173.66bp 389.76bp 67.56bp 276.90bp},clip,max width=0.94\textwidth,max height=0.38\textheight]{Marchisio2013.pdf}
\end{sourceblock}
使用这个定义并从等式开始。 (4.39) 当 g = ϋf2 时，ρf 的输运方程为

\begin{sourceblock}
\includegraphics[page=161,pagebox=mediabox,trim={177.86bp 309.36bp 67.56bp 333.86bp},clip,max width=0.94\textwidth,max height=0.38\textheight]{Marchisio2013.pdf}
\end{sourceblock}
其中

SM 表示由于颗粒成核（即非恒定 N）而导致的流体质量损失。通过强制方程。 (4.77) 等于式(4.77) (4.68)，我们可以推导出式(4.68)中〈Gf2〉1的一致模型。 （4.78）。

\paragraph*{粒子所看到的流体质量变化的一致模型}
正式地，为了方程。为了使（4.77）和（4.68）一致，我们必须有

\begin{sourceblock}
\includegraphics[page=161,pagebox=mediabox,trim={171.35bp 190.50bp 67.56bp 481.36bp},clip,max width=0.94\textwidth,max height=0.38\textheight]{Marchisio2013.pdf}
\end{sourceblock}
这与等式非常相似。 （4.61）。 〈Gf2〉1 的简单一致模型是23

\begin{sourceblock}
\includegraphics[page=161,pagebox=mediabox,trim={126.26bp 140.90bp 67.56bp 525.25bp},clip,max width=0.94\textwidth,max height=0.38\textheight]{Marchisio2013.pdf}
\end{sourceblock}
22 与粒子看到的流体体积一样，我们隐含地假设每个粒子看到的流体质量是不同的。在这种情况下，等式。 （4.75）和（4.81）是对 Ψf2 定义的一致性约束。最简单的情况是 ψf2 = g(t, x, ψp2)，其中 g 是满足式（1）的函数。 （4.75）。在这种情况下，NDF可以写十为n(t, x, vp, Ψp2, vf, Ψf2) = n(t, x, vp, Ψp2, vf)δ(Ψf2 -g(t, x, Ψp2))。如果我们选择线性函数 g(t, x, Ψp2) = ρfΨp2/ρp，则每个粒子都有与其质量 Ψp2 成正比的流体质量。

\footnotetext[23]{一个更简单的模型是}
\begin{sourceblock}
\includegraphics[page=161,pagebox=mediabox,trim={142.42bp 50.20bp 144.94bp 622.18bp},clip,max width=0.94\textwidth,max height=0.38\textheight]{Marchisio2013.pdf}
\end{sourceblock}
其中右侧第一项代表相之间的质量传递，最后一项是将 Ψf2 松弛回其条件平均值 ρfΨp2/ρp 的模型。方程中引入的质量平均流体速度。 (4.68)，以 NDF 形式给出：

\begin{sourceblock}
\includegraphics[page=162,pagebox=mediabox,trim={180.07bp 568.55bp 65.07bp 98.10bp},clip,max width=0.94\textwidth,max height=0.38\textheight]{Marchisio2013.pdf}
\end{sourceblock}
如式中所示。 （4.62），值得注意的是等式的右侧。对于没有传质的系统，(4.80) 不为零。事实上，只要粒子和流体具有不同的速度（即非零粒子斯托克斯数），就需要包括方程（1）中给出的模型。 (4.80) 在 GPBE 中。请注意，仅当存在颗粒时才定义颗粒所见的流体。在实践中，Uf 是通过求解流体相动量传输方程来找到的，即使在不存在粒子的情况下，该方程也能得到很好的定义，在这种情况下，方程： (4.81) 被解释为 GPBE 的一致性要求（方程（4.39））。尽管如此，从方程中知道〈Gf2〉1。 (4.80) 对于导出涉及 vf 的高阶矩非常有用，例如粒子所看到的波动流体速度。

\paragraph*{多组分流体相}
如果颗粒由多种化学物质组成，那么通常流体相也是由多种化学物质组成。在这种情况下，有必要引入一个内坐标向量 ψf，其分量是粒子所看到的每种化学物质的质量。显然，这些内坐标之和等于粒子所看到的流体质量。根据定义，如果 ψfα 是组分 α 的质量，则相空间积分可得出组分流体相质量密度：

\begin{sourceblock}
\includegraphics[page=162,pagebox=mediabox,trim={188.78bp 330.56bp 65.07bp 342.22bp},clip,max width=0.94\textwidth,max height=0.38\textheight]{Marchisio2013.pdf}
\end{sourceblock}
假设所有流体成分以相同的流体速度移动，24 ρfα 的传输方程具有以下形式

\begin{sourceblock}
\includegraphics[page=162,pagebox=mediabox,trim={172.47bp 272.92bp 65.07bp 401.43bp},clip,max width=0.94\textwidth,max height=0.38\textheight]{Marchisio2013.pdf}
\end{sourceblock}
其中右侧的传质项和成核项对于每个组件可能不同。求和方程(4.83) 在所有分量 α 上可得出方程 (4.83)。 （4.77）。请注意等式。 (4.83) 的右侧可能包含一个附加项，对应于流体相中物质之间的均相化学反应。对于每个分量，〈Gfα〉1 的模型可以具有方程式的形式。 (4.80) 用 ρfα 代替 ρf。总而言之，
它将“校正项”均匀地分布在所有存在的粒子中，无论其质量或速度如何。我们排除这种类型的模型，因为它没有考虑 NDF 中的空间梯度（仅在统计中）。在没有流体速度波动的限制下，方程中的模型(4.80) 变为

\begin{sourceblock}
\includegraphics[page=162,pagebox=mediabox,trim={144.70bp 146.22bp 145.79bp 526.16bp},clip,max width=0.94\textwidth,max height=0.38\textheight]{Marchisio2013.pdf}
\end{sourceblock}
其中 n 不再包含 vf 并且 Uf(t, x) 是通过求解方程找到的。 (4.92) 第 124 页。 24 还可以假设每个分量以相同的平均速度加上扩散分量移动。例如，等式。 (4.83) 可以替换为

\begin{sourceblock}
\includegraphics[page=162,pagebox=mediabox,trim={159.83bp 87.01bp 156.18bp 592.11bp},clip,max width=0.94\textwidth,max height=0.38\textheight]{Marchisio2013.pdf}
\end{sourceblock}
{\small 其中 Yfα 是流体相中的物质质量分数，Гf 是有效扩散系数。由于颗粒的存在，有效扩散系数可能明显大于纯流体中的分子扩散系数。\par}
由于物理系统模型中质量守恒的重要性，方程。 （4.64）和（4.68）应始终包含在用于模拟多分散型多相流的矩输运方程组中。一般来说，为了充分描述系统的物理特性，需要几个额外的矩输运方程。最重要的是，需要描述颗粒相和流体相动量，如下所述。

\subsection*{4.3.7 分散相动量输运}\phantomsection\addcontentsline{toc}{subsection}{4.3.7 分散相动量输运}
分散相动量密度定义为

\begin{sourceblock}
\includegraphics[page=163,pagebox=mediabox,trim={177.80bp 492.88bp 67.56bp 179.90bp},clip,max width=0.94\textwidth,max height=0.38\textheight]{Marchisio2013.pdf}
\end{sourceblock}
其中内坐标 ψp2 再次对应于粒子质量。从等式开始(4.39) 当 g = ϋp2vp 时，分散相动量密度的输运方程为

\begin{sourceblock}
\includegraphics[page=163,pagebox=mediabox,trim={121.83bp 412.11bp 67.56bp 259.74bp},clip,max width=0.94\textwidth,max height=0.38\textheight]{Marchisio2013.pdf}
\end{sourceblock}
其中（未封闭的）动量对流项（二阶张量）定义为

\begin{sourceblock}
\includegraphics[page=163,pagebox=mediabox,trim={165.17bp 368.03bp 67.56bp 304.75bp},clip,max width=0.94\textwidth,max height=0.38\textheight]{Marchisio2013.pdf}
\end{sourceblock}
请注意，一般来说，方程中的传输和源项。 （4.85）不会是封闭形式。根据方程中的定义。 （4.84），我们可以识别三种限制情况。

1. 如果粒子质量 (Mp) 恒定，则 ρpUp = MpNUN。

2. 如果颗粒材料密度（ρp）恒定，则ρpUp = ρpαpUV。

3. 如果 Mp 和 ρp 恒定，则 ρpUp = ρpαpUN。

这些情况表明，使用的最佳平均速度取决于保持恒定的内部坐标。然而，由于动量守恒具有根本重要性，因此最好使用方程（1）中的质量平均速度。 （4.85）。由于这个原因，等式。 (4.85) 应始终包含在用于模拟分散相的矩输运方程组中。分散相动量源项出现在方程的右侧。 (4.85) 定义如下。第一个学期，

\begin{sourceblock}
\includegraphics[page=163,pagebox=mediabox,trim={161.01bp 168.37bp 67.56bp 504.41bp},clip,max width=0.94\textwidth,max height=0.38\textheight]{Marchisio2013.pdf}
\end{sourceblock}
描述由于周围流体而作用在颗粒上的表面力（例如阻力、升力和压力）。第二个学期，

\begin{sourceblock}
\includegraphics[page=163,pagebox=mediabox,trim={163.34bp 105.66bp 67.56bp 567.12bp},clip,max width=0.94\textwidth,max height=0.38\textheight]{Marchisio2013.pdf}
\end{sourceblock}
{\small 这是由于作用在单个粒子上的体积力（例如重力）造成的。第三学期，\par}
\begin{sourceblock}
\includegraphics[page=163,pagebox=mediabox,trim={168.97bp 54.90bp 67.56bp 617.88bp},clip,max width=0.94\textwidth,max height=0.38\textheight]{Marchisio2013.pdf}
\end{sourceblock}
描述由于液相的传质而添加到颗粒相的动量。第四学期，

\begin{sourceblock}
\includegraphics[page=164,pagebox=mediabox,trim={174.22bp 596.47bp 65.07bp 76.32bp},clip,max width=0.94\textwidth,max height=0.38\textheight]{Marchisio2013.pdf}
\end{sourceblock}
描述了由于碰撞（参见第 6 章）和粒子成核而导致的粒子动量的不连续变化。一般来说，这些项都不会以封闭形式出现（即使 GPBE 是封闭的），并且必须提供宏观模型才能封闭式（1）。 （4.85）。

\paragraph*{体积和质量恒定}
再次考虑一个系统，其中所有粒子具有相同的体积和质量，并且分散相动量密度为 ρp = ρpαp。这种情况下分散相动量密度的输运方程为25

\begin{sourceblock}
\includegraphics[page=164,pagebox=mediabox,trim={73.69bp 445.27bp 65.07bp 226.58bp},clip,max width=0.94\textwidth,max height=0.38\textheight]{Marchisio2013.pdf}
\end{sourceblock}
其中，分散相应力张量由 upup N = FN −Up ⊗Up 定义。分散相应力张量是二阶矩，因此不是封闭的。在 4.4 节中，我们回顾了 upup N 的一些经典闭包，它们对于 NDF 的特定限制行为有效。然而，一般来说，分散相应力张量的闭合并不平凡，但对于确定方程解的性质极其重要。 （4.91）。事实上，正是由于未能为此项提供准确的闭包，才导致二流体模型不适定（即系统的保守形式不是双曲线）。正如我们将在第 8 章中看到的，离散相应力张量的精确闭合总是会产生（弱）双曲矩方程。最后，请注意，碰撞项可以生成单独的空间通量项（称为碰撞通量），该空间通量项将添加到分散相应力张量中。有关该术语的更多详细信息可以在第 6 章中找到。

\paragraph*{多组分分散相}
对于多组分粒子，粒子中的所有组分都具有相同的速度，因此可以从方程（1）中找到Up。 （4.85）。与单组分情况的唯一区别是，每个组分的传质和成核项可能不同（例如，只有一个组分可能会改变相），在这种情况下，方程的右侧为： (4.85) 将包括每个组成部分的捐款总额。

\subsection*{4.3.8 流体相动量传递}\phantomsection\addcontentsline{toc}{subsection}{4.3.8 流体相动量传递}
对于流体相动量密度，我们可以假设以下形式的动量平衡

\begin{sourceblock}
\includegraphics[page=164,pagebox=mediabox,trim={84.73bp 128.76bp 65.07bp 545.49bp},clip,max width=0.94\textwidth,max height=0.38\textheight]{Marchisio2013.pdf}
\end{sourceblock}
其中Sf表示流体相中的粘性和压应力张量，Af表示流体相上的体力。对流项 Uf ⊗Uf f 需要一个单独的模型来描述由于微尺度湍流等原因可能产生的流体速度波动。在流体速度不波动的极限条件下，

\footnotetext[25]{对于这种情况，质量平均值由下标 N 表示的数量平均值代替。}
Uf ⊗Uf f = Uf ⊗Uf。然而，一般来说，颗粒所见的流体速度可能会由于例如与颗粒相的动量耦合而波动。如式中所示。 (4.85)，术语 Gp p 和 S p 由相间的传质产生。再次令 ψf2 为粒子所见的流体质量，对流项可写为

\begin{sourceblock}
\includegraphics[page=165,pagebox=mediabox,trim={151.97bp 558.17bp 67.56bp 108.49bp},clip,max width=0.94\textwidth,max height=0.38\textheight]{Marchisio2013.pdf}
\end{sourceblock}
这可以被视为对 Ψf2 定义的一致性约束。请注意等式。 (4.92) 应与直接从 GPBE 中找到的矩输运方程一致。 （4.39），它提供了有关粒子相和粒子相所看到的流体相的信息。换句话说，等式。 (4.92) 应该与从方程 (4.92) 推导 g = ψf2vf 的矩输运方程得到的结果一致。 （4.39）。事实上，如果我们忽略传质和粒子成核，则乘以方程： (4.39) 由 ψf2vf 得到，然后在相空间上积分，得到

\begin{sourceblock}
\includegraphics[page=165,pagebox=mediabox,trim={127.67bp 436.99bp 67.56bp 234.86bp},clip,max width=0.94\textwidth,max height=0.38\textheight]{Marchisio2013.pdf}
\end{sourceblock}
其中 Uf 由方程式给出。 （4.81）。方程中的其余项。 (4.94) 的定义为

\begin{sourceblock}
\includegraphics[page=165,pagebox=mediabox,trim={147.83bp 334.90bp 67.56bp 278.91bp},clip,max width=0.94\textwidth,max height=0.38\textheight]{Marchisio2013.pdf}
\end{sourceblock}
\begin{sourceblock}
\includegraphics[page=165,pagebox=mediabox,trim={170.77bp 301.27bp 67.56bp 371.52bp},clip,max width=0.94\textwidth,max height=0.38\textheight]{Marchisio2013.pdf}
\end{sourceblock}
以及

Apf f 项是颗粒看到的流体的质量平均加速度，部分归因于相间的动量传递，Af f 是归因于流体相中的力。 Gf f 项可以使用等式进行评估。 （4.80）：

\begin{sourceblock}
\includegraphics[page=165,pagebox=mediabox,trim={84.81bp 223.73bp 67.56bp 447.24bp},clip,max width=0.94\textwidth,max height=0.38\textheight]{Marchisio2013.pdf}
\end{sourceblock}
其中右侧第一项是由质量传递导致的相间动量交换产生的，第二项取代了等式左侧的对流项。 (4.94) 与等式中的相应项。 (4.92)，最后一项包含粒子质量平均流体速度：

\begin{sourceblock}
\includegraphics[page=165,pagebox=mediabox,trim={170.59bp 136.65bp 67.57bp 530.00bp},clip,max width=0.94\textwidth,max height=0.38\textheight]{Marchisio2013.pdf}
\end{sourceblock}
值得注意的是，不需要额外的模型来获得一致的流体相动量密度对流。为了保持一致性，我们可以要求 Uf p = Uf，以便方程式右侧的最后一项。 (4.99) 为空；26 然而，由于动量守恒，由于 〈Apf2〉1 的贡献，最后一项将被取消，如下所示。

\footnotetext[26]{在极限 τpf →0 下，该一致性条件将是精确的，因为在此极限下 ψf2 →ρfψp2/ρp。}
\paragraph*{粒子所看到的流体加速度的一致模型}
通过比较等式。 (4.94) 与等式。 （4.92），插入方程式后（4.99），我们看到

\begin{sourceblock}
\includegraphics[page=166,pagebox=mediabox,trim={198.31bp 558.73bp 65.08bp 100.26bp},clip,max width=0.94\textwidth,max height=0.38\textheight]{Marchisio2013.pdf}
\end{sourceblock}
以及

\begin{sourceblock}
\includegraphics[page=166,pagebox=mediabox,trim={205.54bp 520.44bp 65.08bp 152.25bp},clip,max width=0.94\textwidth,max height=0.38\textheight]{Marchisio2013.pdf}
\end{sourceblock}
方程中的约束。 （4.101）指出，在相间的质量传递过程中，动量必须守恒，并且它可以与方程（4.101）中的约束相结合。 (4.103)。方程中流体加速度项的最小一致模型。 (4.39) 则可以写成

\begin{sourceblock}
\includegraphics[page=166,pagebox=mediabox,trim={169.36bp 449.30bp 65.08bp 224.16bp},clip,max width=0.94\textwidth,max height=0.38\textheight]{Marchisio2013.pdf}
\end{sourceblock}
and27

\begin{sourceblock}
\includegraphics[page=166,pagebox=mediabox,trim={146.24bp 398.25bp 65.08bp 272.71bp},clip,max width=0.94\textwidth,max height=0.38\textheight]{Marchisio2013.pdf}
\end{sourceblock}
〈Apf2〉1 = 1
等式右侧的前两项。 (4.105) 表示相之间的动量传递，值得注意的是，为了一致性，不需要涉及速度差或 NDF 空间梯度的附加项。 28 第三项是由时间尺度 τpf 的弛豫过程中动量守恒得出的，它对平均流体相动量平衡的贡献恰好抵消了方程 (4.105) 中的最后一项。 （4.99）。但请注意，方程右侧的第一项。 (4.105) 生成方程式中的项 Gp p。 （4.92），第二项对应于没有质量传递的情况下的动量传递。方程中的最后一项。 (4.104) 是一个简单的流体速度波动耗散模型，其中 τf 是流体速度返回其质量平均值的特征弛豫时间。 29 解释宏观流体湍流的更复杂的模型可以在文献中找到（Minier \& Peirano，2001；Simonin 等，1993）。然而，读者应该记住，本书中讨论的介观尺度模型不需要（或不需要）宏观尺度湍流模型。介观尺度模型最多将包括微尺度湍流的影响
27 粒子成核产生的贡献类似于方程 2 中涉及 〈Gp2〉1 的传质项。 (4.105)。同样重要的是要认识到 〈Afp2〉1 的模型将具有以下形式的传质贡献

\begin{sourceblock}
\includegraphics[page=166,pagebox=mediabox,trim={175.60bp 171.16bp 171.94bp 505.15bp},clip,max width=0.94\textwidth,max height=0.38\textheight]{Marchisio2013.pdf}
\end{sourceblock}
因此，仅当 〈Gp2〉1 < 0 时，〈Apf2〉1 才会因传质而产生贡献，即从颗粒到流体的传质。 28 对于湍流（Minier \& Peirano，2001；Simonin 等，1993），将取决于速度差 Uf −Up 的修正项添加到 〈Apf2〉1。然而，这些模型不包括 〈Gf2〉1 模型中的速度差项。从方程。 （4.80），我们看到，即使n非常小（即αp<1），它的空间梯度也可以足够大，以至于〈Gf2〉1不能被忽略。 29 在介观尺度模型中，设置 τf = 0 会强制粒子看到的流体速度等于质量平均流体速度。例如，这对于单向耦合是合适的，其中颗粒不干扰流体。一般来说，由于其他粒子或微尺度湍流的存在而产生的流体波动可以通过添加 vf 的相空间扩散项来建模，类似于用于宏观湍流的那些（Minier \& Peirano, 2001）。时间尺度 τf 将对应于微尺度湍流的耗散时间尺度。

（例如粒子尾迹）对介观尺度变量的影响。在单动流体极限中，我们取τf = 0，使得vf = Uf，这在颗粒相非常稀且动量耦合不重要的情况下是合适的。

\paragraph*{流体相的有效粘度闭合}
方程中 Uf ⊗Uf f 的简单闭包。包含微尺度流体速度波动影响的方程（4.92）是将其替换为 Uf ⊗Uf，然后将 Sf 建模为流体速度 Uf 的牛顿流体，其有效粘度等于流体的分子运动粘度 (νf) 和微尺度有效运动粘度 (νeff) 之和。这导致 Sf = −pfI + τf，其中 pf 是流体相中的压力（例如，从状态方程中找到的压力），并且 τf = 2ρf(νf + νeff)[Df −1 3 trace(Df)I]，其中 Df = 1 2[∇xUf + (∇xUf)T]。对于等体积颗粒，有效粘度可以用混合长度模型来描述：νeff= Cναpdp|Up -Uf|，其中 dp 是颗粒直径，Cν 是有序单位模型常数。通过这种闭包，可以通过求解方程找到质量平均流体速度。 （4.92）耦合到分散相的矩输运方程。在矩输运方程的其他地方，脉动流体速度 vf 被 Uf 代替，这与单动流体极限类似。

\paragraph*{混合动量方程}
方程中给出的混合连续性方程。 （4.69）可以与通过将方程相加得到的混合动量方程相结合。 （4.85）和（4.92）：

\begin{sourceblock}
\includegraphics[page=167,pagebox=mediabox,trim={93.59bp 338.91bp 67.57bp 332.94bp},clip,max width=0.94\textwidth,max height=0.38\textheight]{Marchisio2013.pdf}
\end{sourceblock}
虽然这个表达式不包含相间动量传递项（因此看起来更简单），但它几乎没有实用价值，因为左侧的空间通量项不能用 Umix 来写。只有在 Umix 非常接近 Uf（即非常小的斯托克斯数）的极限情况下，才有可能使用方程 1 的封闭形式准确地模拟分散型多相流。 (4.106)。对于有限斯托克斯数，最好求解分散相和流体相的单独动量方程。

\subsection*{4.3.9 高阶矩输运}\phantomsection\addcontentsline{toc}{subsection}{4.3.9 高阶矩输运}
到目前为止，我们推导出的矩输运方程在速度变量中是一阶的。为了描述一阶矩的涨落，有必要推导二阶矩（有时甚至是高阶矩）的输运方程。和以前一样，这是从等式开始完成的。 (4.39) 对 g 有一个特定的选择。为了说明这是如何完成的，我们将考虑函数 g = ψp2v2 pi，它会产生粒子质量平均矩

\begin{sourceblock}
\includegraphics[page=167,pagebox=mediabox,trim={170.01bp 120.86bp 67.57bp 551.45bp},clip,max width=0.94\textwidth,max height=0.38\textheight]{Marchisio2013.pdf}
\end{sourceblock}
其中 ψp2 再次等于粒子质量。请注意，该矩是总分散相能量的一个组成部分：

\begin{sourceblock}
\includegraphics[page=167,pagebox=mediabox,trim={208.34bp 50.16bp 67.57bp 621.04bp},clip,max width=0.94\textwidth,max height=0.38\textheight]{Marchisio2013.pdf}
\end{sourceblock}
它用于定义颗粒温度，这是分散相中速度波动幅度的度量。一般来说，二阶分散相速度矩不会等于一阶分散相速度矩（U2 pi p U2 pi）的平方，二阶流体速度矩（U2 fi f U2 fi）也是如此。例如，如第 5.1 节所述，介观速度波动可能通过流体-粒子阻力介观模型中的随机分量引入，或者通过惯性粒子的粒子轨迹穿越 (PTC) 引入。在这种情况下，可能需要求解高阶矩的输运方程。

\paragraph*{分散相速度二阶矩}
从等式开始（4.39），对应于方程（4.39）的矩输运方程。 (4.107) 是（不暗示对 i 求和）

\begin{sourceblock}
\includegraphics[page=168,pagebox=mediabox,trim={71.19bp 439.85bp 65.08bp 218.46bp},clip,max width=0.94\textwidth,max height=0.38\textheight]{Marchisio2013.pdf}
\end{sourceblock}
其中传质项定义为

\begin{sourceblock}
\includegraphics[page=168,pagebox=mediabox,trim={167.60bp 399.25bp 65.08bp 273.06bp},clip,max width=0.94\textwidth,max height=0.38\textheight]{Marchisio2013.pdf}
\end{sourceblock}
和对流项

\begin{sourceblock}
\includegraphics[page=168,pagebox=mediabox,trim={162.43bp 356.33bp 65.08bp 315.98bp},clip,max width=0.94\textwidth,max height=0.38\textheight]{Marchisio2013.pdf}
\end{sourceblock}
涉及·p 的其他术语的定义类似。请注意，与矩方法中通常的情况一样，对流项是不封闭的，甚至涉及更高阶的矩。方程右侧的其余项。 （4.109）的定义方式与等式（4.109）中的对应方式非常相似。 （4.85）。因此，我们不会详细讨论它们，只是说它们通常不会以封闭的形式出现。只有在 4.4 节中讨论的特殊情况下，速度 NDF 接近高斯分布，我们才能为方程 1 提供准确的闭包。 (4.111) 就二阶速度矩而言。

\paragraph*{流体相速度的二阶矩}
接下来考虑流体相速度的流体质量平均二阶矩，定义为

\begin{sourceblock}
\includegraphics[page=168,pagebox=mediabox,trim={172.69bp 184.08bp 65.08bp 482.57bp},clip,max width=0.94\textwidth,max height=0.38\textheight]{Marchisio2013.pdf}
\end{sourceblock}
其中 ψf2 是粒子看到的流体质量。从等式开始(4.39) 其中 g = ϋf2v2 fi，对应于方程 (4.39) 的矩输运方程。 (4.112) 具有以下形式（不暗示对 i 求和）

\begin{sourceblock}
\includegraphics[page=168,pagebox=mediabox,trim={80.02bp 102.92bp 65.08bp 556.30bp},clip,max width=0.94\textwidth,max height=0.38\textheight]{Marchisio2013.pdf}
\end{sourceblock}
其中传质项定义为

\begin{sourceblock}
\includegraphics[page=168,pagebox=mediabox,trim={169.41bp 62.12bp 65.08bp 609.99bp},clip,max width=0.94\textwidth,max height=0.38\textheight]{Marchisio2013.pdf}
\end{sourceblock}
和对流项

\begin{sourceblock}
\includegraphics[page=169,pagebox=mediabox,trim={160.13bp 591.68bp 67.57bp 74.97bp},clip,max width=0.94\textwidth,max height=0.38\textheight]{Marchisio2013.pdf}
\end{sourceblock}
涉及·f 的其他术语的定义类似。正如预期的那样，对流项取决于粒子速度，因为它是针对粒子所看到的流体而写的。使用〈Gf2〉1（方程（4.80））、〈Af2〉1（方程（4.104））和〈Apf2〉1（方程（4.105））的一致模型，方程右侧的项(4.113) 可以用更熟悉的形式重写（对 j 求和，但不对 i 求和）：

\begin{sourceblock}
\includegraphics[page=169,pagebox=mediabox,trim={77.30bp 442.18bp 67.57bp 185.28bp},clip,max width=0.94\textwidth,max height=0.38\textheight]{Marchisio2013.pdf}
\end{sourceblock}
值得注意的是，左侧的对流项现在用流体速度来表示。在等式中。 (4.116) 右侧的术语有以下起源。

(i) 前两个是由流体相中的应力和体积力产生的。

(ii) 第三项是由于相间的动量传递（包括当 〈Gp2〉1 > 0 时对 〈Afp2〉1 的传质贡献）。

(iii) 第四项、第五项和第六项是由相间的质量传递产生的。

(iv) 涉及 τf 的第七项是流体速度方差耗散项。请注意，当 τf →0（即单动流体极限）时，我们有 U2 fi f →U2 fi。

(v) 在极限 τpf →0 中，最后一项将为空。因此，为了保持一致性，我们可以要求

\begin{sourceblock}
\includegraphics[page=169,pagebox=mediabox,trim={89.93bp 265.16bp 188.27bp 421.18bp},clip,max width=0.94\textwidth,max height=0.38\textheight]{Marchisio2013.pdf}
\end{sourceblock}
在等式中。 (4.116) 似乎没有明确的齐次流体相速度方差源项。尽管如此，质量和动量传递项都是流体相速度方差的潜在来源。例如，流体阻力项可以是流体相速度方差的源或汇，具体取决于混合矩 UpiUfi p 和 UpiUfi f 的大小，我们接下来将对此进行描述。

\paragraph*{分散相和流体相速度的二阶矩}
介观尺度相间动量交换项通常用速度差 (vp −vf) 表示。因此，在最简单的情况下，这些项在速度差中是线性的，方程中的二阶矩输运方程。 （4.109）和（4.116）将包含涉及以下混合矩的源项：

\begin{sourceblock}
\includegraphics[page=169,pagebox=mediabox,trim={133.28bp 59.89bp 135.81bp 586.93bp},clip,max width=0.94\textwidth,max height=0.38\textheight]{Marchisio2013.pdf}
\end{sourceblock}
当在组件中展开时，会产生两个新的二阶混合矩：

\begin{sourceblock}
\includegraphics[page=170,pagebox=mediabox,trim={160.41bp 571.71bp 157.96bp 75.10bp},clip,max width=0.94\textwidth,max height=0.38\textheight]{Marchisio2013.pdf}
\end{sourceblock}
请注意，这些质量平均矩是使用不同的质量定义的（即粒子质量 Ψp2 与粒子所见的流体质量 Ψf2）。然而，可以在 Σf2 的定义上引入一致性约束 30，使得 Uf ⊗Up f = Uf ⊗Up p （在极限 τpf →0 下是精确的），在这种情况下只需要一个新的矩输运方程。该传输方程可以从方程（2）开始找到。 (4.39) 其中 g = ϋp2vpivf j，并且将具有类似于方程 1 的形式。 （4.109）。最后，值得注意的是，在单动流体极限（即 τf →0）下，速度混合矩应简化为 Up ⊗Uf p →Up ⊗Uf，在这种情况下，仅对粒子速度的高阶矩感兴趣。作为本节的总结，我们应该提醒读者，作为一般规则，由于空间通量或介观模型中的非线性，矩输运方程永远不会闭合。在这里，我们导出了一些具有特定物理意义（例如密度、速度等）的低阶矩的输运方程。然而，由于没有一个传输方程是闭合的，因此有必要提供矩闭合。在下一节中，我们将简要回顾一些对由 GPBE 中的源项确定的特定情况有效的矩闭包。

\section*{4.4 GPBE 的矩闭合}\phantomsection\addcontentsline{toc}{section}{4.4 GPBE 的矩闭合}
在本章中，我们介绍了从多分散型多相流的微观尺度描述开始推导介观模型所需的基本步骤。第一步是识别一组微尺度变量，此处用 (X(n), U(n) p , Ψ(n) p , U(n) f , Ψ(n) f ) 表示，这些变量足以表征流动的介观特性。原则上，可以使用微尺度模型的直接数值模拟来研究这些变量（及其源项）的时间演化。下一步（无疑是最困难的）步骤是在给定一组固定的介观尺度变量 (X(n) = x, U(n) 的情况下，假设微尺度源项的条件期望的介观尺度模型

\begin{sourceblock}
\includegraphics[page=170,pagebox=mediabox,trim={67.52bp 220.01bp 65.07bp 451.23bp},clip,max width=0.94\textwidth,max height=0.38\textheight]{Marchisio2013.pdf}
\end{sourceblock}
f = Σf)。换句话说，必须通过假设（或者在极少数情况下直接推导）出现在 GPBE 中的介观尺度模型，将全耦合系统的完整多粒子描述简化为单粒子描述。一般来说，介观尺度模型包含确定性和概率性成分（例如代表微尺度速度波动的平均阻力和布朗运动）。第五章详细描述了多相系统的介观模型。然而，读者应该认识到介观模型在确定多分散型多相流的最终宏观描述的预测能力方面的根本重要性。简而言之，如果在介观尺度上缺少重要的物理现象，就不可能在宏观尺度上获得令人满意的描述。事实上，正如第 1 章所讨论的，使用基于动力学的建模方法的一个主要优点是物理模型的分离。

\footnotetext[30]{回想一下等式。 (4.75)、(4.81) 和(4.93) 也表示对 Σf2 定义的一致性约束。一般来说，矩输运方程组在极限 τpf→0 下呈现最简单的形式，在这种情况下，相对于粒子质量或流体质量的加权是等效的，因为两者在此极限下线性相关（参见方程（4.80）中的最后一项）。}
纯数学建模步骤（即中观尺度→宏观尺度）的建模步骤（即微观尺度→中观尺度）。对于后者，我们当然指的是闭合矩输运方程所需的矩闭合。推导给定 GPBE 的矩闭包的数学步骤绝不是显而易见的，甚至不能保证有效。如果我们考虑由下式给出的 GPBE

\begin{sourceblock}
\includegraphics[page=171,pagebox=mediabox,trim={98.58bp 501.11bp 119.70bp 126.57bp},clip,max width=0.94\textwidth,max height=0.38\textheight]{Marchisio2013.pdf}
\end{sourceblock}
我们一方面可以区分积累和自由传输术语，另一方面可以区分相空间传输和源项。一般来说，自由传输项往往会使 NDF 的数学形式变得复杂，因为它以 NDF 中出现的速度在整个空间域中传输信息。事实上，如果不存在相空间传输和源项，则几乎没有希望找到适用于任何任意初始 NDF 的矩闭包。 （幸运的是，对于特定的初始 NDF，情况并非那么绝望！）事实上，最成功的矩闭包是针对相空间输运和/或源项主导自由输运项的系统开发的。下面，我们简要回顾一些与多相流更相关的示例。

\paragraph*{单动流体相}
第一个例子是 vf = Uf 且 ϋf = φf 的情况，因此流体相中不存在微观波动。这种情况下的 NDF 为 n(vp, Ψp, vf, Ψf) = n(vp, Ψp)δ(vf - Uf)δ(Ψf -φf)，其中 Uf 通过求解方程 2 求出。 (4.92) 和 φf 是流体相的平均流体成分矢量31。 GPBE 简化为

\begin{sourceblock}
\includegraphics[page=171,pagebox=mediabox,trim={92.70bp 242.29bp 67.57bp 428.67bp},clip,max width=0.94\textwidth,max height=0.38\textheight]{Marchisio2013.pdf}
\end{sourceblock}
其中所有介观模型均在 vf = Uf 和 ψf = φf 下进行评估。这里，术语单动力学适用于在给定时间和位置具有唯一值的介观尺度变量。 GBPE 的单动流体极限是通过在方程 (1) 的流体加速度模型中取 τf →0 得到的。 (4.104)。

\paragraph*{单动力分散相}
第二个例子是分散相速度等于给定其他介观尺度变量的条件预期分散相速度的情况：vp = 〈Up|ψp, vf, ψf〉。在这种情况下，分散相的波动“受制于”流体相波动，并且可能取决于颗粒内坐标（例如颗粒尺寸）。这种情况下的 NDF 为 n(vp, ψp, vf, ψf) = n(ψp, vf, ψf)δ(vp −〈Up|ψp, vf, ψf〉) 并且 GPBE 简化为

\footnotetext[31]{唯一的例外是粒子所看到的流体质量，它在本章前面表示为 ϋf2。在单动流体极限中，我们令 ψf2 = ρfψp2/ρp，其中 ψp2 是粒子质量。}
\begin{sourceblock}
\includegraphics[page=172,pagebox=mediabox,trim={100.08bp 588.02bp 96.44bp 55.47bp},clip,max width=0.94\textwidth,max height=0.38\textheight]{Marchisio2013.pdf}
\end{sourceblock}
其中，所有介观模型均在 vp = 〈Up|Σp, vf, Σf〉 处进行评估。对于斯托克斯数非常小的颗粒，发现了单动力学分散相。通常，对于小斯托克斯数，流体相也被假定为单动力，导致 vp = 〈Up|ψp, ψf〉 以及 GPBE 的额外简化。请注意，对于单动分散相，必须使用从 GPBE 导出的一致模型方程来找到条件速度。对于上述两种单动力学情况，GPBE 都得到了简化，但仍必须求解才能找到 NDF。因此，通常需要额外的关闭才能进一步进行。为此目的，有时可以应用基于接下来描述的时刻的闭合方法之一。

\paragraph*{水动力极限}
正如第 1 章所讨论的，流体动力极限是指碰撞项 S1 足够强以确定速度分布函数的形式的情况。在其最初的公式中（Chapman \& Cowling，1961），流体动力学极限利用了弹性碰撞过程中分散相质量 ρp、动量 ρpUp 和总能量 ρpEp 守恒的事实。反过来，这意味着在相应的矩输运方程中，由 S1 得出的项为零。在无限快碰撞的极限（即零克努森数）下，五个矩方程（ρp、ρpUp、ρpEp）是无粘流体的欧拉方程，速度分布函数是麦克斯韦分布。对于非常小但有限的克努森数，矩输运方程组可导出气体动力学的纳维-斯托克斯方程。有关用于开发流体动力闭合件的努森数展开程序的更多详细信息，请参见 Struchtrup (2005)。该过程的一个重要限制是，所得输运方程仅适用于较小的努森数，并且在许多多分散型多相流中，分散相努森数并不小。因此，需要不依赖于截断克努森数（或类似）展开式的方法来闭合矩输运方程。

\paragraph*{参数化NDF}
封闭矩输运方程组的一种流行方法是根据介观尺度变量假设 NDF 的函数形式。优选地，函数形式的参数可以用几个低阶矩的封闭形式来写。然后可以仅求解低阶矩的输运方程，以确定假定的 NDF 中的参数。然后，NDF 的函数形式已知，并且可用于评估矩输运方程中出现的积分。作为示例，考虑速度 NDF 假设为高斯分布的情况：32

\begin{sourceblock}
\includegraphics[page=172,pagebox=mediabox,trim={130.12bp 83.32bp 127.67bp 582.69bp},clip,max width=0.94\textwidth,max height=0.38\textheight]{Marchisio2013.pdf}
\end{sourceblock}
\footnotetext[32]{回想一下，对于麦克斯韦分布，Σp 是各向同性的。因此，与麦克斯韦分布相比，高斯 NDF 允许增加五个自由度。}
Σp 是粒子速度协方差矩阵（即对称二阶张量）。该假定 NDF 的前十个矩直到二阶为 ρp、ρpUM 和 ρp(Σp + Up ⊗Up)；因此，NDF 完全是通过求解这些时刻的 10 个输运方程来确定的。使用假定 NDF 的一个明显优点是，保证 NDF 对于相空间变量的所有值都是可实现的（非负）。另一方面，一个明显的缺点是，一旦NDF的形式被固定，就没有空间来描述更复杂的功能形状。例如，高斯 NDF 是单峰的（即，对于 Σp 的所有选择，n 在 v = Up 处都有一个最大值）。因此，它无法描述双峰物理，例如粒子轨迹穿越 (PTC) 或墙壁上的镜面反射，为此它将简单地用单峰高斯表示双峰 NDF。一般来说，由于我们事先不知道 GPBE 的解是否接近特定的已知 NDF，因此通常很难为给定问题指定参数化 NDF。此规则的唯一例外是当 GPBE 中的相空间源项之一支配其他项时，使得 NDF 保持接近已知形式。例如，如果微尺度流体速度波动迫使粒子速度波动在小斯托克斯数的限制下服从高斯形式，则可能会发生这种情况。这种情况很少见，因此，为了在 NDF 的形状上提供更大的灵活性，使用截断的无穷级数的想法似乎很有吸引力。

\paragraph*{NDF 的功能扩展}
为了以系统的方式增加自由度的数量，可以使用函数展开来表示 NDF (Grad, 1949b)。以速度分布为例，形式展开为

\begin{sourceblock}
\includegraphics[page=173,pagebox=mediabox,trim={166.11bp 316.20bp 168.64bp 350.73bp},clip,max width=0.94\textwidth,max height=0.38\textheight]{Marchisio2013.pdf}
\end{sourceblock}
其中 Cα 是展开系数，pα 是已知的 α 阶正交多项式族，w 是该族相应的权重函数。例如，如果 w 是麦克斯韦分布（其参数取决于 NDF 的五个流体动矩），则张量积 Hermite 多项式用于 pα。理论上，几乎任何 NDF 都可以通过一组独特的展开系数以这种方式展开。实际上，仅使用有限展开式来表示 NDF：

\begin{sourceblock}
\includegraphics[page=173,pagebox=mediabox,trim={163.48bp 198.51bp 166.01bp 468.43bp},clip,max width=0.94\textwidth,max height=0.38\textheight]{Marchisio2013.pdf}
\end{sourceblock}
有限的展开系数集是由有限的低阶矩集确定的。使用矩的定义，重构的 NDF n* 可用于评估矩：

\begin{sourceblock}
\includegraphics[page=173,pagebox=mediabox,trim={157.69bp 118.72bp 160.22bp 544.36bp},clip,max width=0.94\textwidth,max height=0.38\textheight]{Marchisio2013.pdf}
\end{sourceblock}
括号中的积分项是正交多项式族中已知的。在有限 N 和已知矩的情况下，该线性方程组的形式为 M = AC，可以通过求解来找到展开系数 Cα(t, x)。因此，假定的 NDF n*(t, x, vp) 是有限矩集的唯一函数，后者是通过使用 n* 求解矩输运方程来封闭未封闭项而找到的。事实是
Cα(t, x) 是从（小）线性系统中找到的，使得该方法在计算上非常容易处理，特别是因为系数矩阵 A 可以提前计算和存储（使用变量的变化，A 可以根据流体动力学参数编写）。虽然在数学上非常有吸引力，但有限函数展开式有一个主要缺点：对于有限 N，NDF n* 对于 vp 的某些值来说可能无法实现（负）。换句话说，给定一个可实现的矩集 mk，线性函数展开将产生一个不可实现的 NDF n*。每当通过求解线性系统从矩确定系数时，就会出现此缺点。如第一章所述，当使用不可实现的 NDF 来封闭矩输运方程时，矩输运方程的解很有可能变得不可实现。当这种情况发生时，数值模拟将以不可恢复的方式失败。实际上，生成不可实现的 NDF 的可能性很高，除非 GBPE 的真实解接近权重函数 w(t, x, vp)。因此，我们不认为使用 NDF 的函数展开作为闭合大多数系统的矩输运方程的可行替代方案。

\paragraph*{离散NDF}
闭合矩输运方程的一种越来越流行的方法是假设相空间变量的离散形式。以速度 NDF 为例，速度相空间可以在以 vp = 0 为中心的均匀对称格子上离散化。出于说明目的，我们假设使用 N = 16 个格点，并将相应的速度表示为 Uα。离散 NDF 的正式定义是

\begin{sourceblock}
\includegraphics[page=174,pagebox=mediabox,trim={174.44bp 324.36bp 171.98bp 343.53bp},clip,max width=0.94\textwidth,max height=0.38\textheight]{Marchisio2013.pdf}
\end{sourceblock}
值得注意的是，16 个速度 Uα 是固定的，因此不随空间和时间变化。33 因此，离散 NDF 的矩只会因权重 ρα 的变化而变化。给定一组 16 个矩（或矩的线性组合），我们可以使用 n* 来评估它们，这会产生 M = Aρ 形式的线性系统，其中 A 取决于速度，M 是已知矩。该线性系统可以求解得到 ρα，因此离散 NDF 是由其矩唯一确定的。实际上，16 个矩输运方程可以重写为 16 个权重的方程组。然而，从概念上讲，离散 NDF 只是具有特定形式的重构 NDF，用于封闭矩输运方程。与上述函数展开方法一样，不能保证离散 NDF 中的权重为非负，因此重构的 NDF 可能无法实现。此外，当节点数量增加时，获得负权重的可能性会增加（就像为了找到更好的重建而可能做的那样）。在文献中，离散 NDF 与格子玻尔兹曼方法 (Aidun \& Clausen, 2010) 一起使用来模拟流体动力学极限下的流体流动（即小努森数），其中权重保持为正（因为真正的 NDF 接近麦克斯韦）。不幸的是，负权重的可能性使得具有固定节点的离散 NDF 不适合作为矩输运方程的通用闭包方法。

\footnotetext[33]{虽然我们不需要在这里使用它，但实际上 n* 仅在 xβ 表示的物理空间中的格节点上定义。在一个时间步长中，GPBE 中的自由传输项在给定离散速度 Uα 的情况下在离散空间晶格节点之间移动信息。因此，如果 n 最初是空间晶格上的点分布，则它将始终保持如此。}
\paragraph*{基于求积的矩方法}
第3章详细介绍了使用基于求积的矩方法（QBMM）重建NDF。与上面介绍的矩闭包相比，QBMM具有以下相似之处和优点。

(i) 重建的 NDF 总是可实现的。

(ii)矩反演算法是直接的。

(iii) 重建的 NDF 准确地再现了一组特定的时刻。

(iv) 通过使用更多矩可以提高重建 NDF 的准确性。

(v) 重建的 NDF 可以是离散的 (QMOM) 或连续的 (EQMOM)。

(vi) 连续重建的 NDF 具有任意精度的离散正交表示。

作者认为，QBMM 是闭合由不接近已知极限状态的 GPBE 产生的矩输运方程的最佳方法。尽管如此，读者应该清楚的是，随着内坐标数量的增加，找到 NDF 的准确重建的难度也随之增加。因此，在尝试解决宏观矩问题之前，首先尝试减少 NDF 中的介观自由度的数量总是有利的。例如，如果微尺度物理由几个不同的相互作用的粒子群组成，那么最好编写多个低维 GBPE（即每个群一个）而不是单个高维 GBPE 来描述系统。在接下来的章节中，我们将详细讨论如何使用 QBMM 求解矩输运方程，并提供许多示例。为了理解这些例子，读者应该首先掌握从第 4.3 节中介绍的 GBPE 导出矩输运方程的方法。此外，需要很好地理解第 3 章中介绍的矩反演算法，以便遵循第 7 章和第 8 章中介绍的求解方法的讨论。

\chapter{物理与化学过程的介观模型}
本章阐述并讨论了一些最广泛使用的用于描述颗粒过程的介观尺度模型。读者应该记住，本章并不是对所有可能的介观模型的全面讨论，而只是一些示例模型的集合，在这种情况下用于强调多相系统模拟中通常遇到的主要问题。尽管下一节中说明的一些模型具有普遍有效性，但当应用于固液系统（例如结晶和沉淀过程）、固气系统（例如流化悬浮液、气溶胶反应器）、液-液系统（例如乳液）、气-液系统（例如鼓泡塔和气-液搅拌罐）和液-气系统的模拟时，它们通常采用略有不同的形式（例如蒸发和非蒸发喷雾）。在下文中，我们经常将分散相的元素称为“颗粒”，以表示固体颗粒（例如晶体、固体无定形颗粒和固体聚集体）和流体颗粒（例如液滴和气泡）。本章的其余部分组织如下。在第 5.1 节概述了介观尺度模型开发背后的原理之后，第 5.2-5.5 节讨论了相空间和平流过程（以及随后的扩散过程）的具体示例。随后，第 5.6-5.8 节介绍了相空间点（不连续）过程。对于每个过程，都会导出并讨论最终 GPBE 中出现的介观模型的相应函数形式。讨论还包括量化每个过程的速率的最常见方法，以及不同流态的适用范围。

\section*{5.1 介观尺度建模概述}\phantomsection\addcontentsline{toc}{section}{5.1 介观尺度建模概述}
在本节中，我们将简要概述介观尺度建模方法及其在关闭 GPBE 中的意义。我们首先描述第 4 章导出的 GPBE 中出现的各种术语的物理意义，然后讨论如何从微观物理中导出介观模型。

\subsection*{5.1.1 GPBE 中的介观尺度模型}\phantomsection\addcontentsline{toc}{subsection}{5.1.1 GPBE 中的介观尺度模型}
正如第 4 章中已经讨论的，分散相的状态完全由内坐标向量 ψp 和速度 vp 确定，而粒子位置由外坐标向量 x 定义。此外，颗粒被认为是分散相中不可区分的元素。通常，假设外坐标和内坐标方向的连续变化率不仅是 ψp 和 vp 的函数，而且也是颗粒所见的悬浮流体状态（即主要连续相）的函数。颗粒所见的悬浮流体的状态由其速度 vf 和将颗粒所见的所有其他相关流体相变量组合在一起的矢量 ψf 来识别。假设粒子是不可区分的，并且其状态变化率仅是 vp、ψp、vf 和 ψf 的函数，这意味着具有相同状态和相同速度的粒子，被相同状态和速度的流体包围，将不可避免地以相同的速率沿相同的方向演化。然而，我们将在下面展示代表微尺度上非常快速的力或通量的扩散过程，这些力或通量在介观尺度上无法解析，可以将其建模为粒子状态的随机变化。因此，当存在扩散过程时，具有相同初始状态的两个粒子不需要以完全相同的方式演化。

\paragraph*{多分散型多相流的介观描述}
根据第 4 章中的微尺度模型推导，多分散型多相流的 GPBE 具有以下形式：1

\begin{sourceblock}
\includegraphics[page=177,pagebox=mediabox,trim={70.24bp 255.70bp 67.56bp 327.04bp},clip,max width=0.94\textwidth,max height=0.38\textheight]{Marchisio2013.pdf}
\end{sourceblock}
其中 ψp1 是粒子质量，ψf1 是粒子所看到的流体质量。2 相空间向量 ψf = (ψf1, . . . , ψfMf) 中的其余分量表示粒子所看到的流体属性（例如温度、化学浓度等），需要这些属性来描述相和有限尺寸粒子之间的质量和能量传递（参见第 4.3.4 节）。 GPBE 解释了微观物理和化学过程对数密度函数 (NDF) n(t, x, vp, Ψp, vf, Ψf) 演化的影响。等式左侧的项。 (5.1) 表示由于粒子速度 vp = (vp1, vp2, vp3) 导致的实际空间中的累积、自由传输，以及考虑粒子所见流体特性的空间梯度所需的通量模型。右侧第一项表示由流体-粒子引起的连续相空间平流过程
1 在这里，我们通过将 xif1 指定为粒子所见的流体质量，明确包含第 4.3.6 节中讨论的一致性项。因此，〈Gf〉1 和〈Apf〉1 仅包含相间质量/能量/动量传递产生的贡献。

\footnotetext[2]{指定粒子所见流体质量的最简单方法是 Σf1 = ρfΣp1/ρp，其中 Σp1 是粒子质量。如果所有粒子具有相同的质量，则 ψp1 是常数。然而，在推导矩输运方程时，最简单的方法是在推导过程中将 Ψf1 和 Ψp1 作为自变量，然后在最终结果中代入 Ψf1 = ρfψp1/ρp，这相当于取极限 τpf →0。}
动量交换和体积力，分别产生相空间速度 〈Afp〉1 和 〈Ap〉1。右侧第二项代表由于相空间速度 〈Gp〉1（例如相间的质量交换）而导致的内坐标向量 Ψp = (Ψp1, . . . , ΨpMp) 连续变化对 NDF 的影响。第三项量化了 NDF 的变化，该变化与粒子 vf 由于相空间速度 〈Apf〉1 和 〈Af〉1 所看到的流体速度的连续变化相关。第四项表示由于相空间速度 〈Gf〉1 导致粒子 ψf 看到的流体性质连续变化的影响。第五项表示由不连续或点过程（例如粒子成核和粒子-粒子碰撞）引起的 NDF 变化。最后，方程的最后两项。 (5.1) 对于任意松弛时间 τpf 值，强制粒子所见流体的质量和动量的介观尺度守恒。正如第 4 章所讨论的，我们主要对 GPBE 在极限 τpf →0 下的行为感兴趣。

\paragraph*{介观尺度平流模型}
一般来说，相空间平流项被建模为纯平流和相空间扩散的贡献之和（Gardiner，2004）：3

\begin{sourceblock}
\includegraphics[page=178,pagebox=mediabox,trim={82.85bp 311.00bp 65.07bp 277.56bp},clip,max width=0.94\textwidth,max height=0.38\textheight]{Marchisio2013.pdf}
\end{sourceblock}
纯平流速度 (Afp, Gp, Apf, Gf) 是矢量，其长度分别取决于相空间矢量 vp、ψp、vf 和 ψf 的长度。一般来说，这些速度可以取决于 NDF 中的所有自变量，即 (t, x, vp, Ψp, vf, Ψf)。开发介观尺度模型时的主要任务之一是确定与微尺度模型一致的纯平流速度的正确函数依赖性。同样，相空间扩散系数出现在方程右侧的偏导数中。 (5.2) 是自变量的矩阵函数（即二阶十元）。每个矩阵的大小由其上下文决定，并且矩阵通常不是方形的。4 我们将在第 5.5 节中看到，扩散项作为“随机”波动微尺度过程的模型而出现，这些过程发生在相对于纯平流速度的短时间尺度上。作为一般规则，扩散矩阵不会彼此独立地定义，因此有必要从一致的微尺度模型中导出它们。此外，由于扩散项对 NDF 零阶矩和一阶矩的输运方程没有贡献，因此在不考虑介观描述的多分散多相系统宏观模型的直接推导中，它们经常被忽略。然而，这可能会导致宏观模型不一致，例如，当纯平流速度相对可忽略不计时

\footnotetext[3]{〈Ap〉1 的模型可能包含随机分量，例如由于外力场的波动而产生的随机分量。然而，我们在本章中假设〈Ap〉1 = Ap。}
\footnotetext[4]{例如，Bpv 是正方形，其大小由 vp 的长度决定，而 BpΨ 的行数等于 Ψp 的长度，列数等于 vp 的长度。}
到扩散项。从数学上讲，相空间扩散项倾向于“平滑”NDF 中的不连续性（即 GPBE 在相空间中具有抛物线方向）。因此，只要底层的微观物理产生扩散项，就应该包括它们。正如我们将在第 5.5 节中看到的，方程中的许多扩散矩阵。由于某些状态空间变量的独立性，方程（5.2）将为零。例如，如果相之间的动量交换不依赖于内坐标 ψp，则 Bpψ 将为空。

\paragraph*{纯平流过程}
对于纯平流过程，我们指的是导致外坐标和内坐标连续变化的连续现象。粒子在物理空间中位置的连续变化由物理空间平流（或自由传输）项量化：

\begin{sourceblock}
\includegraphics[page=179,pagebox=mediabox,trim={143.51bp 455.99bp 67.56bp 214.97bp},clip,max width=0.94\textwidth,max height=0.38\textheight]{Marchisio2013.pdf}
\end{sourceblock}
其中粒子速度矢量 vp = (vp1, vp2, vp3) 可以使用不同的策略进行计算，如以下各节所述。粒子速度的连续变化也出现在纯平流项中：

\begin{sourceblock}
\includegraphics[page=179,pagebox=mediabox,trim={66.22bp 386.65bp 67.55bp 284.32bp},clip,max width=0.94\textwidth,max height=0.38\textheight]{Marchisio2013.pdf}
\end{sourceblock}
(5.4) 其中 Afp = (Afp1, Afp2, Afp3) 是由于流体-颗粒动量交换而产生的纯平流速度，Ap = (Ap1, Ap2, Ap3) 是由于颗粒上的体力而产生的纯平流速度。内坐标的连续变化（例如粒子体积、表面积、长度等）也出现在相空间平流项中：

\begin{sourceblock}
\includegraphics[page=179,pagebox=mediabox,trim={116.63bp 297.60bp 67.56bp 372.79bp},clip,max width=0.94\textwidth,max height=0.38\textheight]{Marchisio2013.pdf}
\end{sourceblock}
其中，Gp = (Gp1,Gp2,...,GpM) 是表示内坐标向量的连续变化率的向量。 GPBE 中的类似项源自粒子所见流体的纯平流速度：Apf 和 Gf。我们将回到第 5.2-5.4 节中对由于流体-颗粒相互作用而产生的纯平流过程进行建模的主题。然而，如等式所示。根据方程（4.80），当分散相体积分数不可忽略时，矢量 ψf 必须至少有一个与粒子所见的流体质量相对应的分量。

\paragraph*{相空间扩散过程}
方程中的相空间扩散项。 (5.2) 在 GPBE 中生成大量项（其中许多项为零）。例如，在粒子速度波动是由于微尺度流体湍流5（即Bpψ=0，Bpψf=0）的极限情况下，仅考虑流体-粒子相互作用项，产生速度相空间中的扩散项
5 根据定义，仅当存在连续流体相时才会出现微尺度流体湍流。系数 Bpv 描述颗粒相与连续相的相互作用。相比之下，Bpvf 模拟了粒子所看到的流体速度的快速波动，这些波动不包含在介观尺度阻力项 Ap 中。在介观尺度粒子动量平衡中，生成 Bpvf 的项将取决于流体相质量密度 ρf，因此当流体材料密度 (ρf) 为零时，该项将为空。无论如何，Bpvf 模拟流体-粒子系统中粒子相的动量传递，其中总动量守恒（参见方程（5.17）的讨论）。

由于流体速度相空间方向上的梯度，

\begin{sourceblock}
\includegraphics[page=180,pagebox=mediabox,trim={122.60bp 533.97bp 65.07bp 80.10bp},clip,max width=0.94\textwidth,max height=0.38\textheight]{Marchisio2013.pdf}
\end{sourceblock}
和 (Bfψ = 0, Bfψf = 0) 由于流体速度相空间中的梯度而导致的流体速度相空间中的扩散项，

\begin{sourceblock}
\includegraphics[page=180,pagebox=mediabox,trim={125.35bp 409.71bp 65.07bp 206.59bp},clip,max width=0.94\textwidth,max height=0.38\textheight]{Marchisio2013.pdf}
\end{sourceblock}
扩散矩阵 Bpvf 和 Bfvf 是对称的，最重要的是，微尺度动量守恒要求它们是相关的，因此，最多只需要从微尺度模型中确定六个扩散系数（参见方程（5.17）的讨论）。最简单的情况发生在扩散矩阵各向同性时：

\begin{sourceblock}
\includegraphics[page=180,pagebox=mediabox,trim={140.36bp 285.46bp 65.07bp 353.61bp},clip,max width=0.94\textwidth,max height=0.38\textheight]{Marchisio2013.pdf}
\end{sourceblock}
因此只需要一个扩散系数就可以关闭方程（1）。 （5.6）和（5.7）。当我们回到 5.5 节中由于流体-粒子相互作用而引起的扩散过程建模主题时，这些不同类型的介观尺度扩散模型将被清楚地表述出来。

\paragraph*{点流程}
方程中的点过程项 S1。 （5.1）可以代表大范围的不同物理过程。然而，术语点过程旨在将其与 GPBE 中可以用导数编写的其他术语区分开来。一般来说，我们可以根据 NDF 在数学表达式中的表现方式，将点过程进一步分为三类。

(i) 零阶。此类别中的术语不依赖于 NDF。一个常见的例子是由于连续相的质量交换而导致新颗粒的成核。

(ii) 第一顺序。此类别中的术语与 NDF 线性相关。一个常见的例子是由于连续相产生的表面应力而导致颗粒破裂。

{\small (iii) 第二顺序。此类别中的术语与 NDF 呈二次方关系。一个常见的\par}
{\small 示例是两个不同尺寸的粒子的聚集。\par}
一般来说，零阶点过程由相空间子集中的粒子源项建模，该子集对应于添加到系统中的粒子的属性。另一方面，一阶点过程通常会产生涉及 NDF 的相空间子集上的积分项。例如，对于颗粒破碎，相空间中所有较大尺寸都存在线性积分源项和线性汇项。最后，二阶点过程涉及积分源项和汇项，涉及在相空间变量的不同值处评估的两个 NDF。从数学角度来看，GPBE 中处理二阶点过程比其他类型更复杂，因为它会导致对 NDF 的非线性依赖。我们将在第 5.6-5.8 节中回到对由流体-粒子和粒子-粒子相互作用引起的点过程进行建模的主题，以及在第 6 章中对由粒子-粒子碰撞引起的点过程进行建模的主题。

\subsection*{5.1.2 介观尺度模型的制定}\phantomsection\addcontentsline{toc}{subsection}{5.1.2 介观尺度模型的制定}
正如第 4 章所讨论的，条件相空间速度（〈Afp〉1、〈Ap〉1、〈Gp〉1、〈Apf〉1、〈Af〉1 和 〈Gf〉1）和不连续“源”项 S1 需要通过使用适当的“平均”模型来实现介观尺度闭包微观化学和物理过程。导出这些闭合模型的最简单过程是假设所有粒子在统计上都是相同的。在这种假设下，相空间速度可以从与平均场量相互作用的单粒子统计数据（即单个粒子的粒子特性的变化率）导出，平均场量模拟所有其他粒子的影响。多粒子统计对相空间速度的影响通常以校正的形式引入，例如基于分散相体积分数，这将减少到稀极限下的孤立粒子统计。一个有用的起点是考虑拉格朗日参考系中粒子属性随粒子速度移动的变化。

\paragraph*{拉格朗日单粒子模型}
单个粒子之后，介观尺度变量的变化率可以写成拉格朗日形式：6

\begin{sourceblock}
\includegraphics[page=181,pagebox=mediabox,trim={185.99bp 143.43bp 67.56bp 424.20bp},clip,max width=0.94\textwidth,max height=0.38\textheight]{Marchisio2013.pdf}
\end{sourceblock}
其中每种情况下的左侧都是粒子之后的时间导数，右侧的第一项对应于粒子变量的连续变化，
6 当使用欧拉-拉格朗日模拟以有限颗粒样本来近似分散型多相流时，该公式特别方便。正如第 5.2 节和第 5.3 节中所讨论的，虽然一些介观变量是密集的（即独立于粒子质量），但在推导单粒子模型时通常最好从保守的广泛变量（例如粒子质量或粒子动量）开始。例如，在第 4 章中，我们发现 ψf 必须至少有一个分量，对应于粒子所见的流体质量，以便描述分散相体积分数不恒定的情况。

右侧的最后一项对应于不连续的变化（例如碰撞）。介观尺度建模器的目标是为方程右侧的项提供函数形式。 (5.9) 取决于
(i) 描述粒子的介观尺度变量，即 X* p(t)、U* p(t)、 Σ* p(t)、U* f(t) 和 Σ* f (t)；和
(ii) 在粒子位置评估的 NDF 的（与时间相关的）矩：x = X* p(t)。

请注意，介观模型对 NDF 矩的依赖性用于通过“平均场”变量（例如分散相体积分数）引入多粒子效应。因此，方程的右侧。 （5.9）与第 4 章中介绍的精确微尺度模型不同（即第 103 页的方程（4.1）–（4.3））。形式上，考虑到这些依赖性，方程右侧的介观尺度模型。 (5.9) 可以明确地表示为（例如）A*

\begin{sourceblock}
\includegraphics[page=182,pagebox=mediabox,trim={308.62bp 447.04bp 99.00bp 236.91bp},clip,max width=0.94\textwidth,max height=0.38\textheight]{Marchisio2013.pdf}
\end{sourceblock}
f )，其中函数参数表示确定单粒子属性所需的变量。正如第 4 章中详细描述的，这些函数通常与条件相空间速度相关，如下所示： 〈Afp〉1 = A* fp(t, x, vp, Ψp, vf, Ψf)，即通过评估 GPBE 中出现的相空间变量处的函数。然而，如果 A* fp 包含一个随机分量来模拟微观尺度上快速变化的过程，则需要一个额外的步骤来导出条件相空间速度（Gar-diner，2004）。最后，方程中的不连续项。 （5.9）通常使用随机点过程建模。

\paragraph*{随机加速度模型}
如上所述，介观模型可以包含随机过程，这导致方程中出现扩散项。 （5.2）。以 A* fp 为例，我们可以写

\begin{sourceblock}
\includegraphics[page=182,pagebox=mediabox,trim={102.07bp 262.31bp 65.07bp 421.28bp},clip,max width=0.94\textwidth,max height=0.38\textheight]{Marchisio2013.pdf}
\end{sourceblock}
我们引入了 delta 相关维纳过程 W(t) 来模拟快速时间尺度（Gardiner，2004）。例如，随机向量 Wpvp(t) = [Wpvp1(t), Wpvp2(t), Wpvp3(t)]T 具有独立分量 Wpvpi(t)（均值为零 〈Wpvpi(t)〉= 0），其不同时间的期望值由 〈Wpvpi(t)Wpvpi(s)〉= δ(t −s) 给出。等式右侧随机向量的长度。 (5.10) 与下标中的向量相同（即 Wpψp(t) 与 ψp 具有相同的长度）。为了完整起见，我们包含了查找方程中的扩散项所需的所有可能的随机向量。 （5.2）；然而，读者应该意识到，如果底层微尺度模型没有在介观尺度模型中引入随机分量，则系数（例如 bpvp）可能为零。此外，所有系数都出现在方程的右侧。 (5.10) 可以依赖于与 A* fp 相同的变量，如上所述。7 使用维纳过程（Gardiner，2004），生成方程式中的介观尺度模型的完整单粒子模型集。 (5.2) is8

\footnotetext[7]{实际上，通常假设扩散系数仅取决于 t 和 x 就足够了。然后，潜在的随机过程被近似为高斯分布。}
\footnotetext[8]{我们再次假设 A* p = Ap，因此体积力没有波动分量。}
\begin{sourceblock}
\includegraphics[page=183,pagebox=mediabox,trim={131.44bp 576.05bp 67.56bp 57.72bp},clip,max width=0.94\textwidth,max height=0.38\textheight]{Marchisio2013.pdf}
\end{sourceblock}
为了建立单粒子模型和介观尺度扩散系数之间的联系，我们可以重写方程： (5.11) 使用向量和矩阵：Z*= Z + BW，其中 Z*= [A* f]T 和 Z = [Afp, Gp, Apf, Gf]T 是长度为 N 的列向量，W = [Wpvp, WpΨp, Wpvf, WpΨf]T 是长度为 N 的列向量，B 是大小为 N × N 的矩阵，定义为

\begin{sourceblock}
\includegraphics[page=183,pagebox=mediabox,trim={126.80bp 531.88bp 309.05bp 152.06bp},clip,max width=0.94\textwidth,max height=0.38\textheight]{Marchisio2013.pdf}
\end{sourceblock}
\begin{sourceblock}
\includegraphics[page=183,pagebox=mediabox,trim={183.20bp 457.45bp 185.73bp 185.03bp},clip,max width=0.94\textwidth,max height=0.38\textheight]{Marchisio2013.pdf}
\end{sourceblock}
读者应注意，微尺度模型用于确定 B 中的非零项，因此对于以下讨论，可以假设 B 是已知的。使用矩阵表示法和维纳过程的属性 (Gardiner, 2004)，对称 N × N 扩散矩阵 D 可以定义为

\begin{sourceblock}
\includegraphics[page=183,pagebox=mediabox,trim={220.86bp 365.66bp 67.56bp 308.96bp},clip,max width=0.94\textwidth,max height=0.38\textheight]{Marchisio2013.pdf}
\end{sourceblock}
如果我们在方程左侧表示介观尺度平流项： (5.2) 由

\begin{sourceblock}
\includegraphics[page=183,pagebox=mediabox,trim={168.88bp 327.43bp 171.41bp 357.81bp},clip,max width=0.94\textwidth,max height=0.38\textheight]{Marchisio2013.pdf}
\end{sourceblock}
然后我们有 〈Z〉1n = Zn −∇η · (Dn)，其中长度为 N 的相空间梯度向量定义为

\begin{sourceblock}
\includegraphics[page=183,pagebox=mediabox,trim={186.53bp 268.58bp 189.06bp 397.56bp},clip,max width=0.94\textwidth,max height=0.38\textheight]{Marchisio2013.pdf}
\end{sourceblock}
将 〈Z〉1 的定义与方程进行比较。根据方程（5.2），读者可以观察到 D 与 GPBE 中的相空间扩散系数直接相关。此外，这些扩散系数是由单粒子介观模型通过方程式确定的。 （5.12）。换句话说，GPBE 中平流项的形式可以通过单个粒子的介观尺度（拉格朗日）模型精确确定，其形式如下： （5.9）。因此，推导 GPBE 封闭形式的首选方法是从拉格朗日粒子模型开始。这种选择提供了直接连接到 GPBE 的欧拉-拉格朗日公式的额外优势，可用于模拟分散型多相流。此规则的例外是粒子所看到的流体特性的介观模型。在第 4 章中，我们发现，为了确保与流体相宏观模型的一致性，〈Gf〉1 必须包含空间梯度项，如方程式 1 所示。 （4.62）和（4.80）。仅观察单个分散相粒子的拉格朗日行为时，对此类非零项的需求并不明显。

\paragraph*{微尺度动量守恒}
{\small 作为扩散系数的具体示例，让我们考虑唯一的介观尺度变量是 U* 的情况\par}
{\small p and U∗\par}
{\small f.此外，我们假设拉格朗日\par}
这些变量的单粒子模型是

\begin{sourceblock}
\includegraphics[page=184,pagebox=mediabox,trim={169.47bp 606.56bp 306.31bp 80.04bp},clip,max width=0.94\textwidth,max height=0.38\textheight]{Marchisio2013.pdf}
\end{sourceblock}
p dt = Afp + bpvpWpvp(t) + bpvfWpvf(t),

\begin{sourceblock}
\includegraphics[page=184,pagebox=mediabox,trim={169.68bp 580.36bp 65.07bp 101.31bp},clip,max width=0.94\textwidth,max height=0.38\textheight]{Marchisio2013.pdf}
\end{sourceblock}
f dt = Apf + bfvpWpvp(t) + bfvfWpvf(t)。

如果我们分别用 mp 和 mf 表示（恒定）粒子和流体质量 9，则微尺度动量守恒可以表示为

\begin{sourceblock}
\includegraphics[page=184,pagebox=mediabox,trim={202.57bp 498.48bp 65.07bp 172.79bp},clip,max width=0.94\textwidth,max height=0.38\textheight]{Marchisio2013.pdf}
\end{sourceblock}
因此，就方程中的模型系数而言。 （5.13），我们必须有
mpAfp + mfApf = 0，mpbpvp + mfbfvp = 0，

\begin{sourceblock}
\includegraphics[page=184,pagebox=mediabox,trim={402.10bp 442.90bp 65.07bp 245.30bp},clip,max width=0.94\textwidth,max height=0.38\textheight]{Marchisio2013.pdf}
\end{sourceblock}
mpbpvf + mfbfvf = 0。

第一个约束是相间动量传递的通常确定性条件。其余两个约束影响 GPBE 中出现的扩散矩阵。从方程。 （5.13）和等式。 （5.12），我们看到本例中 GPBE 中出现的扩散矩阵（在应用等式（5.15）中的约束之后）由 10 定义

\begin{sourceblock}
\includegraphics[page=184,pagebox=mediabox,trim={125.41bp 334.18bp 65.07bp 332.99bp},clip,max width=0.94\textwidth,max height=0.38\textheight]{Marchisio2013.pdf}
\end{sourceblock}
其中对称 3 × 3 扩散矩阵由 Dpf = 1 2(bpvpbT pvp + bpvfbT pvf) 定义。为了进一步进行，有必要引入 bpvp 和 bpvf 的介观尺度模型。最简单的可能模型是各向同性的，例如，导致 bpvp = 0 和 bpvf ∝(ρf/ρp)I，其中 I 是 3 × 3 单位矩阵。因此，GPBE 中保持微尺度动量的扩散矩阵的简单一致模型为 11

\begin{sourceblock}
\includegraphics[page=184,pagebox=mediabox,trim={150.95bp 227.62bp 65.07bp 439.55bp},clip,max width=0.94\textwidth,max height=0.38\textheight]{Marchisio2013.pdf}
\end{sourceblock}
其中 Dpf 是标量系数。请注意，虽然 D 是 6×6 矩阵，但其最大秩为 3。因此，为了在微观尺度上守恒动量，相空间中的扩散仅存在于六个维度中的三个维度中。类似的结果可以在
9 在 GPBE 中，质量将被相密度 ρp 和 ρf 取代，如果唯一的介观变量是 U* p 和 U* f，则它们必须恒定。 10矩阵乘积 A ⊗B 定义为

\begin{sourceblock}
\includegraphics[page=184,pagebox=mediabox,trim={153.77bp 103.42bp 151.32bp 554.31bp},clip,max width=0.94\textwidth,max height=0.38\textheight]{Marchisio2013.pdf}
\end{sourceblock}
\footnotetext[11]{这里，作为一个例子，我们假设速度波动是由连续相中的微尺度湍流引起的，因此 Dpf 将与粒子 uf·uf f 所看到的流体的速度方差成正比。}
等式的更一般背景。 （5.2）由于微观尺度上质量、动量、能量等的守恒。我们将在 5.5 节中更详细地讨论这一点。

\paragraph*{确定性加速度模型}
由于构建包含确定性和随机成分的成熟拉格朗日模型非常具有挑战性，因此第一步通常是集中精力对确定性部分进行建模。这将是我们在第 5.2-5.4 节中开发流体-颗粒相互作用模型时的方法。在这里，作为一个例子，我们考虑体积力加速度项 Ap 的确定性模型。这些力是由于力场的存在而产生的，典型的例子是电场和引力场。众所周知，当连续相不存在时（即粒子在真空中运动），体积力也会起作用。在重力的情况下，作用在单个粒子上的每单位质量的最终力可以写为 Ap = g。在电场的情况下，力可以计算为电场强度 E 和与颗粒相关的电荷 qp 除以颗粒质量的乘积：Ap = qpE/Mp。请注意（与始终具有相同符号的重力相反）电场力取决于与粒子相关的电荷的符号。我们还应该注意到，体力也出现在粒子 Af 所看到的流体速度的介观尺度模型中。一般来说，由于这些力作用在每个相上，与其他相是否存在无关，因此通常很容易找到合适的介观模型。此外，同一系统中可以存在多个体力，在这种情况下 Ap 将是所有这些力的总和。为了避免繁琐的符号，在本章中，除非另有必要，否则我们将单独讨论每个力（隐含地假设总力是通过求和得出的）。12在本章的大部分内容中，我们将忽略粒子所看到的流体相特性的波动，这相当于假设流体的相空间变量“受制于”它们的平均值，即单动流体极限。 （更多细节请参见第 4.4 节。）在实践中，这意味着 vf = Uf 和 ϋf = φf，其中流体相的平均速度和平均属性是从方程 1 形式的单独守恒方程中找到的。 （4.92）。然而，当颗粒材料密度与流体材料密度（例如气泡流）相比较小时，即使在层流中，忽略流体相波动也是不合理的（Cartellier et al., 2009）。对于这种情况，应将随机分量添加到相空间速度模型中。

\subsection*{5.1.3 与宏观模型的关系}\phantomsection\addcontentsline{toc}{subsection}{5.1.3 与宏观模型的关系}
在 4.3 节中，针对 NDF 的选定矩导出了示例宏观输运方程。介绍了方程中介观尺度平流模型的精确形式。 （5.2），明确地推导出这些模型产生的一些示例矩源项是有意义的。为此，我们将使用第 4.3.1 节中提出的规则进行相空间积分。为简单起见，我们仅考虑涉及 〈Afp〉1 的平流项，并假设唯一感兴趣的相空间变量是 v 和 vf，并且方程 1 中的模型。 (5.2) 简化为

\begin{sourceblock}
\includegraphics[page=185,pagebox=mediabox,trim={149.63bp 94.23bp 152.16bp 576.73bp},clip,max width=0.94\textwidth,max height=0.38\textheight]{Marchisio2013.pdf}
\end{sourceblock}
\footnotetext[12]{虚拟质量力是该规则的例外，因为它涉及时间导数 ̇U*。对于这种情况，Afp 必须在拉格朗日意义上定义，并考虑方程 1 中的虚拟质量。 （5.9）。}
令 g(v, vf) 为我们的任意函数，即方程中的积分。 (4.41) 可以重写为

\begin{sourceblock}
\includegraphics[page=186,pagebox=mediabox,trim={72.20bp 550.98bp 65.07bp 75.14bp},clip,max width=0.94\textwidth,max height=0.38\textheight]{Marchisio2013.pdf}
\end{sourceblock}
通过应用分部积分，可以进一步减少扩散项以找到

\begin{sourceblock}
\includegraphics[page=186,pagebox=mediabox,trim={78.61bp 464.25bp 65.07bp 172.71bp},clip,max width=0.94\textwidth,max height=0.38\textheight]{Marchisio2013.pdf}
\end{sourceblock}
观察该方程的右侧，很明显，〈Afp〉1 的介观模型将在宏观尺度矩输运方程中生成源项，这些项取决于矩的阶数，而矩阶又取决于 g 中出现的相空间变量的幂。

\paragraph*{零阶矩}
对于零阶矩，任意函数设置为 g = 1。 （5.19），这个选择产生 ∫ ∂ ∂vp

\begin{sourceblock}
\includegraphics[page=186,pagebox=mediabox,trim={217.21bp 331.84bp 232.52bp 345.96bp},clip,max width=0.94\textwidth,max height=0.38\textheight]{Marchisio2013.pdf}
\end{sourceblock}
dvp dvf = 0
正如预期的那样，因为粒子加速度不会改变粒子数。

\paragraph*{一阶矩}
对于一阶速度矩，任意函数设置为 g = vpi。从方程。 （5.19），这个选择产生

\begin{sourceblock}
\includegraphics[page=186,pagebox=mediabox,trim={120.26bp 215.30bp 229.67bp 451.35bp},clip,max width=0.94\textwidth,max height=0.38\textheight]{Marchisio2013.pdf}
\end{sourceblock}
Afpin dvp dvf = − Afpi NN,
其中 N 是由方程定义的总数浓度。 （4.45）。该表达式是由于流体-粒子相互作用而在 xi 方向上产生的（负）宏观粒子加速度。请注意，如果我们设置 g = vfi （即粒子看到的相空间流体速度），则该积分为零；然而，Apf 的相应积分将不为零。如前所述，在这里我们可以观察到方程中扩散项的介观模型。 (5.2) 不影响一阶矩，但 Afp 的介观模型决定了分散相和连续相之间的宏观动量交换。

\paragraph*{二阶矩}
{\small 有许多涉及 vp 和 vf 分量的二阶矩。尽管如此，导致非零加速度项的三种类型是 g1 = v2\par}
{\small pi，g2=vpivpj，g3=vpivfj。将这些函数插入方程。 (5.19) 导致\par}
\begin{sourceblock}
\includegraphics[page=187,pagebox=mediabox,trim={65.03bp 549.99bp 106.83bp 55.47bp},clip,max width=0.94\textwidth,max height=0.38\textheight]{Marchisio2013.pdf}
\end{sourceblock}
读者可以观察到相空间扩散系数在二阶矩的输运方程中显示为源 (sink13) 项。在许多加速项的介观尺度模型中，扩散系数被简单地忽略。例如，当分散相材料密度远大于连续相材料密度 (ρf ≪ρp)（例如气体颗粒流）时，这将是合适的。然而，当情况并非如此时（例如气泡流），连续相中的微尺度湍流波动可能会导致分散相中的显著速度波动。在这种情况下，当介观模型中包含扩散项时，颗粒相的空间分散将明显更大。当我们在 5.5 节中介绍扩散动量传输的介观尺度模型时，我们将更详细地讨论这一点。与二阶矩一样，三阶和高阶矩的输运方程将包含纯平流项和扩散项的贡献。通过遵循从 5.1.2 节中介绍的介观模型开始的推导，这些贡献将被明确地知道（即不需要额外的物理模型来闭合输运方程）。

\section*{5.2 相空间平流：传质和传热}\phantomsection\addcontentsline{toc}{section}{5.2 相空间平流：传质和传热}
当属于分散多相系统的颗粒与周围流体相交换质量和热量时，会引起识别颗粒状态的内坐标的一些变化。在质量传递的情况下，驱动力通常与颗粒和周围流体之间的化学势差有关。例如，在过饱和溶液中，溶质分子将从溶液扩散到颗粒，并在被纳入颗粒晶格后，其质量将增加，从而其体积、表面积和“尺寸”也会增加。另一个流行的例子来自气泡。当包含不同气体分子的气泡与周围的流体相交换其中一些分子时，其质量根据传质速率而变化。在传热的情况下，驱动力通常与颗粒和周围流体之间的温差有关。在5.1.2节中，我们介绍了方程（1）中介观尺度模型的一般公式。 （5.9）。在这里，我们对 G* p.14 的介观模型感兴趣。特别是，我们对方程中的确定性分量 Gp 感兴趣。 （5.11）。在下文中，我们将内坐标向量表示为 ψ* p，描述 ψ* p 的变化，但为了清楚起见，通常会用特定符号来表示 ψ* p 的特定分量。例如，如果 ψ*p1 是粒子质量，我们用符号 Mp 来表示。然而，读者应该记住，在所有情况下，我们都在对粒子后面的内坐标变化进行建模，如第 5.1.2 节所述。

13 对称方阵 Bpvp 是非负的，导致源项。非对称方阵 Bpvf 对应于粒子速度和流体速度之间的互相关，因此不需要具有特定的符号。

\footnotetext[14]{在第 5.3 节中，我们将讨论质量传递引起的动量变化的介观尺度模型，其中我们将看到 Afp 包含与 Gp 成比例的贡献。}
\paragraph*{颗粒质量平衡}
一般来说，连续性陈述是用粒子质量 (xi* p1 = Mp) 来写的，所有其他表达式都从中导出。由于分子与流体相的交换而引起的颗粒质量的变化率通常写为

\begin{sourceblock}
\includegraphics[page=188,pagebox=mediabox,trim={110.11bp 526.07bp 65.07bp 133.04bp},clip,max width=0.94\textwidth,max height=0.38\textheight]{Marchisio2013.pdf}
\end{sourceblock}
其中 J(U* f ) 是摩尔通量，量化每单位时间和每单位颗粒表面积（可以是所有介观变量的函数）接近（或离开）颗粒的分子数量，Mw 是分子的分子量，Ap 是颗粒表面积。进一步详细说明方程。 (5.20) 有必要将颗粒质量 Mp 与颗粒表面积 Ap 联系起来。第一步通常是将颗粒质量写为颗粒材料密度 ρp 和颗粒体积 Vp 的乘积。当分散相元素的物质密度一定时，颗粒质量的所有变化都反映在颗粒体积上，两个变量的使用是完全等价的。当分散相元素的材料密度不恒定时，即使颗粒质量恒定，颗粒体积也可能发生一些变化。一个非常常见的例子是气泡在具有大压力梯度的容器中移动（例如由于流体静力学），其中，在不改变单个气泡的质量的情况下，压力差会导致气泡体积的显著变化。在下一节中，我们将研究用于描述粒子“大小”的内坐标的各种选择。另请注意，在单粒子模型中，如果 Ap 和 Mp 因粒子而异，则它们都可以包含在 xi* p 的定义中。在一般情况下，可以引入多个内坐标来解释粒子体积和密度的变化（即 xi*

\begin{sourceblock}
\includegraphics[page=188,pagebox=mediabox,trim={136.09bp 322.59bp 298.96bp 361.11bp},clip,max width=0.94\textwidth,max height=0.38\textheight]{Marchisio2013.pdf}
\end{sourceblock}
p2 = Mp)，当体积和密度由于化学变化等原因而独立变化时，这将是合适的。同样，如果颗粒的化学成分很重要，则可能需要引入颗粒物质质量分数 Yp 和流体相物质质量分数 Yf 的内坐标，它们可以通过质量交换关系（如方程式 1）耦合。 （5.20）或化学反应。从概念上讲，重要的一点是建模者可以自由定义内坐标 xi* p 和 xi* f，以包含描述单粒子水平的化学和物理所需的任何变量。在本章的其余部分中，我们将尝试通过不在每一步中引入 xi* p 的新组件来保持符号的易处理性。然而，读者应该记住，任何颗粒特征（例如直径、表面积、化学成分等）都可能是多分散的，因此可以包含在 xi* p 的定义中。

\paragraph*{粒子-焓平衡}
同样，在相间热交换的情况下，会引起内坐标的连续变化。例如，如果内坐标之一是颗粒的温度 Tp，则颗粒焓的变化率可以计算为
d dt(Mpcp

\begin{sourceblock}
\includegraphics[page=188,pagebox=mediabox,trim={239.64bp 91.18bp 65.07bp 595.47bp},clip,max width=0.94\textwidth,max height=0.38\textheight]{Marchisio2013.pdf}
\end{sourceblock}
{\small 哪里 cp\par}
\begin{sourceblock}
\includegraphics[page=188,pagebox=mediabox,trim={280.46bp 62.23bp 166.21bp 621.70bp},clip,max width=0.94\textwidth,max height=0.38\textheight]{Marchisio2013.pdf}
\end{sourceblock}
{\small p 是颗粒的比热，q(U*\par}
{\small f ) 是从连续相到颗粒的热通量的介观模型。显然，在以下情况下\par}
质量和比热恒定的颗粒，15 得到以下颗粒温度的介观模型： dTp dt = Ap Mpcp

\begin{sourceblock}
\includegraphics[page=189,pagebox=mediabox,trim={262.64bp 587.76bp 67.56bp 91.29bp},clip,max width=0.94\textwidth,max height=0.38\textheight]{Marchisio2013.pdf}
\end{sourceblock}
然而，一般来说，除了Tp之外，粒子还可以用其他介观变量来描述。同样，热通量可能是由于对流、辐射和内部加热造成的，因此可能取决于多个介观变量。我们在 5.2.3 节中简要考虑了对流传输问题，并建议读者参考 Crowe 等人的书。 (1998) 对其他模型的详细讨论。此外，如上所述，在单粒子模型中 Ap 和 Mp 都可以包含在 xi* p 的定义中。在更一般的情况下，方程中的单粒子焓平衡。由于质量传递、化学反应、辐射等，方程（5.21）可以在右侧包括源/汇项，这会将单粒子能量平衡与质量平衡耦合起来。因此，与单粒子质量平衡一样，多组分和化学反应系统通常需要额外的内坐标。最后，重要的是要记住，方程 1 中出现的 G* p 和 G* f 的单粒子模型中有两种类型的项。 （5.9）。第一种类型表示相之间的质量/能量/物质转移，第二种类型表示一个相内内坐标的变化（例如均相化学反应）。对于第一种类型，微尺度质量/物种/能量守恒将耦合 G* p 和 G* f 的单粒子模型，如第 5.1.2 节中针对动量所述。在以下部分中，我们不会明确考虑颗粒所见流体相的单颗粒模型。尽管如此，读者应该记住，这些模型必须与第 4 章中所示的粒子属性的单粒子模型一致。

\subsection*{5.2.1 颗粒尺寸的介观变量}\phantomsection\addcontentsline{toc}{subsection}{5.2.1 颗粒尺寸的介观变量}
作为实践问题，限制描述粒子“大小”所需的内坐标数量是有用的。因此，在本节中，我们将回顾使用一个或两个内坐标的颗粒尺寸的可用描述。

\paragraph*{等维粒子}
当分散相的元素可以归类为等维时，即它们具有几乎相同的尺寸或在多个方向上分布，并且具有恒定的材料密度，通常使用单个内坐标来识别元素的“尺寸”。这可以是颗粒质量（或体积）、颗粒表面积或颗粒长度。事实上，在等维粒子的情况下，这些量都是彼此相关的。例如，在球形或立方体颗粒的普通情况下，颗粒体积和颗粒表面积可以很容易地写为 Vp = kVd3 p 和 Ap = kAd2 p，或者换句话说，作为特征长度 dp（即球体的直径和立方体的边缘）、体积形状因子 kV 和表面积形状因子 kA 的函数。对于等维物体，特征长度的选择很简单，kA 和 kV 之间的比率始终等于 6。然而，该方法也可以扩展到非等维对象。在这种情况下，只有在以下情况下，扩展才非常有用：

\footnotetext[15]{例如，如果颗粒质量随其温度变化，则热通量必须包括与颗粒和流体之间的温差相对应的分量，以及由于质量传递导致的总热能变化的分量。}
颗粒每个尺寸的变化率使得当颗粒收缩或增大时其形状和形态不会改变。尽管在这种情况下粒子并不是严格的等维物体，但它的大小仍然可以通过单个内坐标来量化。事实上，对于一个通用物体（具有特定形状），总是可以在多种可能的特征长度中自由定义一个特征长度 d* p，16 并将真实颗粒体积和颗粒表面积表示为该特征长度的函数：Vp = k* p)3 和 Ap = k* p)2。一般来说，实际形状因子 k* V(d* A(d* A 和 k* V) 之间的比率不会是 6，但是，在引入等效特征长度 dp = φcd* p 后，形状校正因子 φc = 6k*

\begin{sourceblock}
\includegraphics[page=190,pagebox=mediabox,trim={208.61bp 528.58bp 264.81bp 155.46bp},clip,max width=0.94\textwidth,max height=0.38\textheight]{Marchisio2013.pdf}
\end{sourceblock}
A，获得等效或平均形状因子（其比率因此被迫等于六）：

\begin{sourceblock}
\includegraphics[page=190,pagebox=mediabox,trim={169.07bp 475.16bp 166.62bp 183.16bp},clip,max width=0.94\textwidth,max height=0.38\textheight]{Marchisio2013.pdf}
\end{sourceblock}
在这些假设下，如果颗粒尺寸是（第一）内坐标（即 xi* p1 = dp）并且假设 ρp 为常数，则方程 11 为： (5.20) 变为

\begin{sourceblock}
\includegraphics[page=190,pagebox=mediabox,trim={216.01bp 407.87bp 65.07bp 265.86bp},clip,max width=0.94\textwidth,max height=0.38\textheight]{Marchisio2013.pdf}
\end{sourceblock}
颗粒尺寸的变化率也用 Gp 表示，在颗粒增长的情况下可以为正值，在颗粒收缩的情况下可以为负值。这可能是指示由于质量交换而引起的相空间平流速率的最流行的方法，也许是因为很容易测量不同时刻颗粒尺寸的变化，例如通过简单的成像技术。如果内坐标是粒子体积（即 xi* p1 = Vp），则定义变为

\begin{sourceblock}
\includegraphics[page=190,pagebox=mediabox,trim={200.65bp 292.17bp 198.21bp 380.21bp},clip,max width=0.94\textwidth,max height=0.38\textheight]{Marchisio2013.pdf}
\end{sourceblock}
而如果内坐标是颗粒表面积（即 xi* p1 = Ap），

\begin{sourceblock}
\includegraphics[page=190,pagebox=mediabox,trim={197.72bp 236.89bp 195.03bp 434.87bp},clip,max width=0.94\textwidth,max height=0.38\textheight]{Marchisio2013.pdf}
\end{sourceblock}
由于 Mw、ρp、kA 和 kV 通常是常数，因此唯一相关的介观尺度依赖性是与 xi* 相关的介观依赖性

\begin{sourceblock}
\includegraphics[page=190,pagebox=mediabox,trim={340.22bp 203.42bp 106.45bp 480.52bp},clip,max width=0.94\textwidth,max height=0.38\textheight]{Marchisio2013.pdf}
\end{sourceblock}
p1 以及从连续相到颗粒的摩尔通量 J(U* f )。上述定义颗粒尺寸的方法是模拟相空间中由于传质引起的平流的最简单方法，对细节感兴趣的读者可以参考专门文献（Gavi et al., 2007, 2008; Marchisio et al., 2001a, 2001b, 2001c, 2002）。这种方法已成功应用于流体-颗粒系统（即流化悬浮液、结晶过程）的模拟，其中固体颗粒可以合理地假设为等维的，以及流体-流体系统（即喷雾、气泡、液滴）的模拟，其中颗粒实际上是球形或略椭圆形。

16 通常选择特征长度是为了表示易于测量的量。例如，如果使用筛分来测量颗粒尺寸，则该特征长度可能是第二大的。事实上，当筛分或筛分颗粒时，第二大尺寸（而不是最大尺寸）决定颗粒是否通过筛子。当改为使用静态光散射来确定颗粒尺寸时，可以选择最大和第二大之间的中间尺寸。

\paragraph*{非等维粒子}
对于非等维颗粒，通常使用多个内坐标来描述颗粒的形状、形态和尺寸。为了说明与传质引起的相空间平流相关的问题，让我们分析一个简单的例子。我们考虑将针状颗粒描述为长方体，其长度为 xi* p1，宽度和深度相等为 xi* p2。显然，粒子体积可以很容易地从这两个量计算出来（即 Vp = xi*

\begin{sourceblock}
\includegraphics[page=191,pagebox=mediabox,trim={119.06bp 529.22bp 216.43bp 140.73bp},clip,max width=0.94\textwidth,max height=0.38\textheight]{Marchisio2013.pdf}
\end{sourceblock}
p2）。在引入两个特征尺寸的变化率之后，问题可以用这两个内坐标来描述，即 Gp1 = ̇ xi*

\begin{sourceblock}
\includegraphics[page=191,pagebox=mediabox,trim={145.44bp 504.61bp 276.46bp 178.28bp},clip,max width=0.94\textwidth,max height=0.38\textheight]{Marchisio2013.pdf}
\end{sourceblock}
分别为p2。这可以通过量化接近颗粒不同面的分子摩尔通量来轻松完成。如果我们用 J12(U*

\begin{sourceblock}
\includegraphics[page=191,pagebox=mediabox,trim={87.93bp 468.75bp 101.02bp 202.63bp},clip,max width=0.94\textwidth,max height=0.38\textheight]{Marchisio2013.pdf}
\end{sourceblock}
p2，获得以下速率：Gp1 = 2MwJ22/ρp 和 Gp2 = 4MwJ12/ρp。例如，另一种描述可以基于颗粒体积和表面积，在这种情况下，这两个新内坐标的变化率将变为 ̇Vp = (xi*

\begin{sourceblock}
\includegraphics[page=191,pagebox=mediabox,trim={65.02bp 419.15bp 68.06bp 250.44bp},clip,max width=0.94\textwidth,max height=0.38\textheight]{Marchisio2013.pdf}
\end{sourceblock}
p2Gp1。对更多细节感兴趣的读者可参阅专业文献（Heineken et al., 2011; Immanuel \& Doyle, 2005; Ma et al., 2007; Pinto et al., 2007; Qamar et al., 2011b; Qamar \& Seidel-Morgenstern, 2009; Wang et al., 2008）。

\paragraph*{分形粒子}
在许多实际情况下（例如在聚集期间），构成分散相的颗粒不能被认为是等维或非等维填充对象。相反，由于它们内部包含一些空隙，它们更方便地被描述为分形物体。在这种情况下，颗粒被认为是由初级颗粒组成，通常其结构由初级颗粒的尺寸（或体积）来确定， xi* p1 = dp 和每个颗粒的初级颗粒数量 xi* p2 = Np。在这种情况下，区分初级粒子（即聚集在一起形成分形对象的那些粒子）和次级粒子（对应于分形对象本身）是有用的。根据颗粒的结构，由特定 dp 和 Np 值表征的颗粒仍然可以具有非常不同的表观尺寸。这种结构通过分形维数进行量化，一个可能的定义是

\begin{sourceblock}
\includegraphics[page=191,pagebox=mediabox,trim={207.02bp 181.27bp 209.54bp 484.88bp},clip,max width=0.94\textwidth,max height=0.38\textheight]{Marchisio2013.pdf}
\end{sourceblock}
其中 Rg 是颗粒的回转半径，通常定义为初级颗粒中心与次级颗粒质心的均方根距离，kg 是统一数量级的比例因子。很明显，对于给定的 dp 和 Np，较小的分形维数对应于较大的回转半径，反之亦然。分形维数通常为一到三，代表非常不同的结构，从松散的线性聚集体到紧凑的完全填充的物体。当质量交换引起的相空间平流涉及分形物体时，可能会出现不同的情况。通常，传质会导致初级颗粒的尺寸以及整个次级颗粒的形态和结构发生改变。初级颗粒尺寸的变化率可由式（1）计算。 (5.23)

假设初级粒子表现为孤立粒子。可以引入校正“拥挤”因子 kc（范围从零到一）来解释其他初级粒子的存在，17 正如 Marchisio (2009) 所示。一种可能的假设是，在传质过程中，整个次级粒子生长，使初级粒子的数量保持恒定。在第一种情况下，传质引起初级粒子尺寸的变化，但不影响每个次级粒子的初级粒子数量，从而得出以下内坐标连续变化率的表达式：

\begin{sourceblock}
\includegraphics[page=192,pagebox=mediabox,trim={159.22bp 520.61bp 156.78bp 153.12bp},clip,max width=0.94\textwidth,max height=0.38\textheight]{Marchisio2013.pdf}
\end{sourceblock}
式中，kc为上述拥挤因子，kV和kA为初级颗粒的体积和表面积形状因子。另一种可能的情况是，由于质量交换，初级粒子结合在一起，失去了它们的特性。在正摩尔通量 J > 0 的情况下，或者换句话说，在颗粒生长的情况下，这种情况似乎更有可能发生。在第二种情况（通常称为湮灭）中，传质引起初级颗粒尺寸和每个聚集体的初级颗粒数量的变化。如果在初级粒子生长过程中，假设次级粒子的回转半径及其分形维数几乎保持恒定，则可以得到以下两个内坐标连续变化率的表达式：

\begin{sourceblock}
\includegraphics[page=192,pagebox=mediabox,trim={102.47bp 373.69bp 100.03bp 299.26bp},clip,max width=0.94\textwidth,max height=0.38\textheight]{Marchisio2013.pdf}
\end{sourceblock}
另一种可能的情况是，由于质量传递，仅初级粒子的数量发生变化，而它们的尺寸或多或少保持不变。在负摩尔通量 J < 0 的情况下，或者换句话说，在收缩粒子的情况下，该假设似乎是现实的。事实上，在这种情况下，外部颗粒更有可能先于内部颗粒被消耗。因此，两个内坐标连续变化率的结果表达式为

\begin{sourceblock}
\includegraphics[page=192,pagebox=mediabox,trim={153.13bp 258.75bp 150.70bp 414.21bp},clip,max width=0.94\textwidth,max height=0.38\textheight]{Marchisio2013.pdf}
\end{sourceblock}
不同的颗粒过程的行为会有很大不同，并且在每种情况下，都必须假设微观物理场景才能制定介观尺度模型。有关这些和其他示例的更多详细信息，请读者参阅 Lee 等人的著作。 (1962)、Marchisio 和 Barresi (2009) 以及 Zucca 等人。 （2006）。

\subsection*{5.2.2 结晶和非晶颗粒的尺寸变化}\phantomsection\addcontentsline{toc}{subsection}{5.2.2 结晶和非晶颗粒的尺寸变化}
晶体颗粒的生长过程涉及几个步骤，例如溶质分子从溶液主体扩散到晶体表面、晶体表面上的吸附、表面上的扩散、附着到台阶、沿台阶扩散以及整合到晶体扭结位点。这些不同过程的速率取决于材料的类型和操作条件。溶质分子的最终总摩尔通量，之前用 J(U*

\begin{sourceblock}
\includegraphics[page=192,pagebox=mediabox,trim={302.52bp 92.51bp 144.15bp 591.43bp},clip,max width=0.94\textwidth,max height=0.38\textheight]{Marchisio2013.pdf}
\end{sourceblock}
f ), 当然是

\footnotetext[17]{在许多颗粒系统中，事实证明，拥挤因素与可用表面积的减少有关。从简单的几何考虑，可以证明该因子的尺度为 kc ∝N−1/3}
{\small p，而通过实验拟合发现了不同的缩放指数。\par}
受每个单个过程的速率影响，通常由一个或两个过程控制整个现象。对于等维粒子，尺寸增长率 Gp 与 J 的关系为 Gp = MwJ/ρp。一般来说，最终摩尔通量是颗粒 Tf 所见的液相温度、过饱和度 18 S 以及一些颗粒特性（例如尺寸 dp 或表面积 Ap）的函数。为了清楚起见，我们将根据这些变量编写介观模型，但读者应该记住，在单粒子模型中，这些变量分别对应于流体相 xi* f 和粒子 xi* p 的内坐标。

\paragraph*{动力学状态}
人们已经发展了不同的理论来解释晶体生长的过程并推导出相对速率。一种流行的理论是二维增长理论。根据这一理论，当分子被吸收到晶体表面上时，它们会形成一个新的二维圆核，就像经典的三维成核一样，当它的尺寸大于临界尺寸时，它会进一步生长，覆盖整个晶体表面。根据新核形成和现有核生长的相对速率，可能出现不同的情况。如果核扩散形成完整层的速率远大于新核形成的速率，则我们处于单核生长极限，导致以下速率：

\begin{sourceblock}
\includegraphics[page=193,pagebox=mediabox,trim={269.60bp 363.30bp 67.56bp 298.88bp},clip,max width=0.94\textwidth,max height=0.38\textheight]{Marchisio2013.pdf}
\end{sourceblock}
Gp = C1hAp[ln(S )]1/2 exp
其中h为台阶高度，Ap为晶体表面积，C1和C2具有从理论推导出来的精确表达式，但由于是鲜为人知且难以测量的参数函数，通常用作通过拟合实验数据获得的经验常数。更多细节可以在 Ohara \& Reid (1973) 和 Myerson (2002) 的著作中找到。在相反的情况下，当二维成核速度远快于面的扩散和覆盖时，我们处于多核生长极限，导致以下速率：

\begin{sourceblock}
\includegraphics[page=193,pagebox=mediabox,trim={166.30bp 242.99bp 67.56bp 417.42bp},clip,max width=0.94\textwidth,max height=0.38\textheight]{Marchisio2013.pdf}
\end{sourceblock}
其中 C3 也可以从理论中推导出来，但通常是通过拟合实验数据获得的。替代生长模型基于诞生和扩散机制，假设核扩散速率有限，并且核可以在任何位置形成，包括不完整的层，并且核之间不存在相互生长，也已被制定，得出以下表达式：

\begin{sourceblock}
\includegraphics[page=193,pagebox=mediabox,trim={281.70bp 142.58bp 67.56bp 519.61bp},clip,max width=0.94\textwidth,max height=0.38\textheight]{Marchisio2013.pdf}
\end{sourceblock}
Gp = C4(S −1)2/3[ln(S )]1/6 exp
其中 C4 和 C5 通常是通过拟合实验数据获得的。对细节感兴趣的读者可以参考文献（Garside et al., 1975; Janssen Van Rosmalen et al., 1975）。上面给出的模型对于制定和制定非常有用

\footnotetext[18]{过饱和度 S 通常定义为局部瞬时溶质浓度与平衡浓度或液相溶解度的比率。在沉淀的情况下，过饱和度被定义为构成沉淀的离子浓度与流体相中的溶解度乘积的乘积之比。}
测试不同的机制，但所有这些都无法定量预测实验观察到的正确行为（并且在相当大的操作条件范围内）。一种有趣且更可靠的方法是所谓的 Burton–Cabrera–Frank 模型（Nyvlt 等人，1985），该模型假设晶面上沿着螺旋位错进行连续生长过程，从而得出以下表达式：

\begin{sourceblock}
\includegraphics[page=194,pagebox=mediabox,trim={279.65bp 547.08bp 65.07bp 118.35bp},clip,max width=0.94\textwidth,max height=0.38\textheight]{Marchisio2013.pdf}
\end{sourceblock}
Gp = C6Tf(S −1)ln(S )tanh
其中 C6 和 C7 也来自理论或实验。方程（5.27）很好地描述了增长率与过饱和度的相关性从近抛物线到线性的变化，因此在实际应用中得到了广泛的应用。

\paragraph*{经验增长模型}
另一种（非常流行）的方法基于非常简单的想法，即使用以下经验表达式计算粒子的增长率：Gp = kg(S −1)g，其中 g 和 kg 均来自实验，而 kg 又经常用阿伦尼乌斯型定律进行描述：

\begin{sourceblock}
\includegraphics[page=194,pagebox=mediabox,trim={205.78bp 397.26bp 65.07bp 282.19bp},clip,max width=0.94\textwidth,max height=0.38\textheight]{Marchisio2013.pdf}
\end{sourceblock}
这些经验定律的有效性有限，不应在其衍生的过饱和值范围之外使用。

\paragraph*{传质限制制度}
如果溶质分子在晶格中的整合速率和面的形成是速率控制步骤，则先前提出的生长模型是有效的。例如，对于悬浮在蒸气中的结晶颗粒来说，几乎总是如此。然而，当晶体悬浮在液体中时，或者当生长的颗粒是无定形的时，因此溶质分子的结合没有障碍（或非常低的障碍），那么溶质分子从溶液主体扩散到颗粒表面的过程可能成为控制步骤。众所周知，粒子周围形成了可以用简单模型描述的浓度（以及动量和热）边界层。这些模型假设边界层中存在线性浓度分布，并引入传质系数 kd。由此产生的增长率（在等维物体的扩散控制增长的情况下）变为

\begin{sourceblock}
\includegraphics[page=194,pagebox=mediabox,trim={192.36bp 165.46bp 65.07bp 508.27bp},clip,max width=0.94\textwidth,max height=0.38\textheight]{Marchisio2013.pdf}
\end{sourceblock}
其中 ceq 是液相中的平衡溶质浓度（或溶解度），传质系数 kd 根据以舍伍德数表示的相关性计算：Sh = kddp

\begin{sourceblock}
\includegraphics[page=194,pagebox=mediabox,trim={253.32bp 105.95bp 65.07bp 575.15bp},clip,max width=0.94\textwidth,max height=0.38\textheight]{Marchisio2013.pdf}
\end{sourceblock}
其中 D 是本体溶液中溶质分子的分子扩散系数。在蠕变（或斯托克斯）流动条件下，以及在刚性球体的情况下，可以从理论中轻松导出舍伍德数的以下相关性：

\begin{sourceblock}
\includegraphics[page=194,pagebox=mediabox,trim={206.33bp 52.94bp 65.07bp 633.16bp},clip,max width=0.94\textwidth,max height=0.38\textheight]{Marchisio2013.pdf}
\end{sourceblock}
其中 P ́eclet 数定义为

\begin{sourceblock}
\includegraphics[page=195,pagebox=mediabox,trim={219.68bp 602.65bp 67.56bp 84.39bp},clip,max width=0.94\textwidth,max height=0.38\textheight]{Marchisio2013.pdf}
\end{sourceblock}
其中流体相的施密特数定义为

\begin{sourceblock}
\includegraphics[page=195,pagebox=mediabox,trim={227.11bp 547.65bp 229.65bp 126.70bp},clip,max width=0.94\textwidth,max height=0.38\textheight]{Marchisio2013.pdf}
\end{sourceblock}
粒子雷诺数为

\begin{sourceblock}
\includegraphics[page=195,pagebox=mediabox,trim={233.63bp 490.87bp 67.56bp 179.62bp},clip,max width=0.94\textwidth,max height=0.38\textheight]{Marchisio2013.pdf}
\end{sourceblock}
Rep =
需要提醒读者的是，U*f 是粒子看到的流体相的速度，|U*

\begin{sourceblock}
\includegraphics[page=195,pagebox=mediabox,trim={109.68bp 456.92bp 355.41bp 227.58bp},clip,max width=0.94\textwidth,max height=0.38\textheight]{Marchisio2013.pdf}
\end{sourceblock}
f|是滑移速度，dp 是颗粒直径，νf 是流体相的运动粘度。请注意等式。 (5.33) 取决于粒子速度 U* p，并且在零斯托克斯数限制下有效，其中 U* p = U* f，以便粒子跟随流体。方程中的相关性。 (5.31) 仅对于 Rep < 1 和 Sc > 200 有效。对于较大的粒子雷诺数，可以使用以下相关式：Sh = 2 + 0.724Re0.48 p Sc1/3，对于 100 < Rep ≤2000 有效，Sh = 2 + 0.425Re0.55 p Sc1/3，对于 2000 < Rep 有效≤105。在其他可用的相关性中，重要的是要引用 Ranz \& Marshall (1952) 针对宏观粒子提出的相关性：Sh = 2.0 + 0.6Re1/2 p Sc1/3。这些表达式假设流体速度 U* f 已知。对于在湍流流体中移动的微米级（或更小）粒子，仅知道系综平均流体速度 〈Uf〉，最好采用由 Armenante \& 推导的介观模型

\begin{sourceblock}
\includegraphics[page=195,pagebox=mediabox,trim={206.74bp 317.12bp 134.21bp 365.70bp},clip,max width=0.94\textwidth,max height=0.38\textheight]{Marchisio2013.pdf}
\end{sourceblock}
Kirwan (1989)：Sh = 2.0 + 0.52 p /νf 是修正粒子雷诺数，εf 是流体湍流耗散率。然而，我们应该提醒读者，在介观尺度模型中，仅包括由流体中存在的颗粒产生的微尺度湍流，在这种情况下，宏观尺度湍流相关性不适用。事实上，方程。 （5.29）对于扩散控制增长是大小依赖的（通过舍伍德数）具有重要的影响，这将在接下来的章节中强调（特别是在第 7.3 节中讨论矩方法时）。

\paragraph*{溶解模型}
溶解的逆过程，即溶质分子离开结晶颗粒或无定形颗粒的晶格，因为颗粒与欠饱和溶液（即 S < 1）接触，通常由扩散控制。因此，方程中报告的介观模型。 (5.29) 可用于计算由于溶解而引起的尺寸变化率。

\subsection*{5.2.3 非等温系统}\phantomsection\addcontentsline{toc}{subsection}{5.2.3 非等温系统}
与传质类似，非等温粒子的对流热通量（在没有传质的情况下）建模为

\begin{sourceblock}
\includegraphics[page=195,pagebox=mediabox,trim={212.25bp 51.81bp 67.56bp 634.89bp},clip,max width=0.94\textwidth,max height=0.38\textheight]{Marchisio2013.pdf}
\end{sourceblock}
其中 kh 是传热系数，Tf 是颗粒看到的流体温度，Tp 是颗粒温度。系数 kh 一般是根据努塞尔数的相关性来计算的：

Nu = khdp

\begin{sourceblock}
\includegraphics[page=196,pagebox=mediabox,trim={252.67bp 570.96bp 221.24bp 108.85bp},clip,max width=0.94\textwidth,max height=0.38\textheight]{Marchisio2013.pdf}
\end{sourceblock}
其中 αc p 是分散相的导热率。在写方程。 （5.34），我们假设粒子内部不存在热梯度。努塞尔数是通过相关性计算的，例如 Ranz \& Marshall (1952) 提出的：Nu = 2 + 0.6Re1/2 p Pr1/3，其中普朗特数定义为

\begin{sourceblock}
\includegraphics[page=196,pagebox=mediabox,trim={230.53bp 482.81bp 229.78bp 188.66bp},clip,max width=0.94\textwidth,max height=0.38\textheight]{Marchisio2013.pdf}
\end{sourceblock}
αc f 是流体相的热导率。上面报告的相关性与 Rep < 5 × 104 的实验具有良好的一致性。可对非均匀场效应 (Michaelides \& Feng, 1994)、吹气效应 (Renksizbulut \& Yuen, 1983) 和湍流效应 (Yearling \& Gould, 1993) 进行校正。在存在传质的情况下，方程中的模型。必须修改方程（5.34）以包括由于质量的增加/损失而导致的粒子热函的增加/损失。

\subsection*{5.2.4 气泡的传质}\phantomsection\addcontentsline{toc}{subsection}{5.2.4 气泡的传质}
另一个有趣的例子是分散在连续液相中的气泡，与之进行质量交换。同样对于这种情况，由于质量传递引起的内部坐标的变化率是从单个气泡的简单质量平衡开始写入的。遵循气液系统的标准符号，单粒子质量平衡变为 dMp

\begin{sourceblock}
\includegraphics[page=196,pagebox=mediabox,trim={213.75bp 288.96bp 65.07bp 392.04bp},clip,max width=0.94\textwidth,max height=0.38\textheight]{Marchisio2013.pdf}
\end{sourceblock}
式中，kd 为传质系数，Ap 为气泡表面积，Δc 为传质驱动力（即气泡中活性成分的浓度与颗粒看到的液相平衡浓度之差）。此外，在这种情况下，基本假设是气泡是均匀的，并且气泡内部没有检测到显著的浓度梯度。因此，总体传质率是通过 kd 和 Ap 的估计来量化的。事实上，气泡的一个重要的复杂问题是它们的形状不是固定的，但根据主相和分散相的性质以及它们的尺寸，它们可以呈现球形、椭球体、凹坑椭球体帽、裙边和球帽形状（Tomiyama 等人，2002a）。与不同形状相对应的状态由方程（1）中粒子雷诺数假设的值来识别。 （5.33），即 E ̈otv ̈os （或 Bond）数，

\begin{sourceblock}
\includegraphics[page=196,pagebox=mediabox,trim={218.56bp 120.87bp 65.07bp 552.97bp},clip,max width=0.94\textwidth,max height=0.38\textheight]{Marchisio2013.pdf}
\end{sourceblock}
和莫顿数，

\begin{sourceblock}
\includegraphics[page=196,pagebox=mediabox,trim={217.84bp 78.19bp 65.07bp 590.10bp},clip,max width=0.94\textwidth,max height=0.38\textheight]{Marchisio2013.pdf}
\end{sourceblock}
{\small 其中 μf 和 ρf 是液相的粘度和材料密度，g 是重力加速度，σ 是液相和气相之间的界面张力，并且\par}
Δρ 是液相和气相之间的材料密度差。这种情况通常通过用单个内坐标（即等效球直径）描述气泡并引入纵横比（定义为气泡的短轴和长轴之间的比率）来处理。该长宽比 E 可以通过使用 Moore (1965) 提出的经验方程作为莫顿数的函数来计算：E = 1/(1 + 0.043Re5/3 p Mo1/3)。另一种方法是使用 Wellek 等人提出的相关性。 (1966) 对于液-液滴：E = 1/(1 + 0.1613Eo0.757)，这对于 Eo < 40 和 Mo ≤10−6 有效，而对于 Eo > 40 和 Rep > 1.2 流体颗粒通常为球形。一旦已知特征E值，气泡的实际面积Ap与等体积球体面积Aeq的比值可计算如下：

\begin{sourceblock}
\includegraphics[page=197,pagebox=mediabox,trim={164.84bp 471.68bp 67.56bp 191.08bp},clip,max width=0.94\textwidth,max height=0.38\textheight]{Marchisio2013.pdf}
\end{sourceblock}
空气质量当量

就清洁气泡的传质系数而言（例如，参见 Clift 等人 (1978) 的评论），在蠕动（或斯托克斯）流动条件下移动的球形气泡的情况下，提出了以下关系式：Sh = 1 + (1 + 0.564Pe2/3)3/4。对于 Rep ≥70 的球形颗粒，舍伍德数可以用以下关系表示（Lochiel \& Calderbank，1964）：

\begin{sourceblock}
\includegraphics[page=197,pagebox=mediabox,trim={182.72bp 368.36bp 67.56bp 293.82bp},clip,max width=0.94\textwidth,max height=0.38\textheight]{Marchisio2013.pdf}
\end{sourceblock}
其中，如前所述，D 是分子从一相转移到另一相的分子扩散率。如果雷诺数足够大，则括号中的项与单位和方程式相比较小。 (5.39) 简化为 Higbie (1935) 提出的经典结果：Sh = 2Pe1/2/√π。在雷诺数较高时，气泡呈现近似扁球体的形状。在这种情况下，舍伍德数（基于真实的气泡表面）变为

\begin{sourceblock}
\includegraphics[page=197,pagebox=mediabox,trim={146.96bp 225.23bp 67.56bp 413.61bp},clip,max width=0.94\textwidth,max height=0.38\textheight]{Marchisio2013.pdf}
\end{sourceblock}
仔细观察方程。 （5.40）表明偏心率对传质系数的影响很小（E 非常低时只有几个百分点）。然而，净传质速率可能会大大增加，高于相同体积球体的值，因为气泡的肥大导致传质面积增加，如方程 1 所示。 （5.38）。对于湍流液相中的气泡，Danckwerts (1951) 提供了另一种描述，他建议对渗透模型进行改进，假设 kd 与气泡界面暴露于具有可变接触时间的湍流涡流所产生的平均表面更新率相关，从而得出 kd =

\begin{sourceblock}
\includegraphics[page=197,pagebox=mediabox,trim={139.52bp 121.31bp 344.71bp 566.89bp},clip,max width=0.94\textwidth,max height=0.38\textheight]{Marchisio2013.pdf}
\end{sourceblock}
Ds，其中 s 是表面元素的替换率。正如 Lamont 和 Scott (1970) 所建议的，可以假设后一个量与液相中的微尺度湍流运动有关，范围从最小的粘性尺度到惯性尺度。通过使用柯尔莫哥洛夫的各向同性湍流理论，最终的湍流传质系数为 kd = KD0.5(εf/νf)0.25，其中 εf 是液相中的湍流耗散率，K 是单位量级的模型常数，通常取等于 1.13。然而，我们必须再次提醒读者，介观模型并不能描述宏观湍流效应，因此涉及εf的相关性不足以描述介观物理。

\subsection*{5.2.5 液滴的热/质量传递}\phantomsection\addcontentsline{toc}{subsection}{5.2.5 液滴的热/质量传递}
讨论当分散在气体中的液滴分别由于冷凝或蒸发而增大或收缩时发生的同时传热和传质是很有趣的。在这种情况下，颗粒尺寸（表面积或体积）的变化率与蒸发和冷凝潜热引起的能量交换密切相关。为了说明主要问题，让我们考虑由单组分液体组成的液滴的情况，该液滴蒸发后与周围气相发生质量交换。 （相应的冷凝情况可以通过质量和热通量的符号反转来进行类似处理。）根据 Marchisio \& Fox (2007) 第 5 章的讨论，我们让 Y1 表示从液滴中蒸发的化学成分的气相摩尔分数，Y2 表示其他（停滞）成分的摩尔分数，Tf 表示液滴看到的气相温度。蒸发的液滴被材料和热边界层包围。我们分别用 Yf1 和 Ys1 表示气相主体（距离液滴表面较远）和液滴表面上蒸发组分的摩尔分数。类似地，我们将 Tf 和 Ts 分别定义为距液滴表面较远距离处和液滴表面上的气相温度。在这四个量中，Yf1 和 Tf 通常被假定为已知（并且通常在内坐标向量 xi* f 中分组在一起，以表示液滴所见的气相特性）。通过在液滴表面施加热力学平衡（使用克劳修斯-克拉佩龙定律，其有效性远离临界点），质量和焓守恒（考虑到液滴温度的变化必须与蒸发潜热相关）产生以下表达式：

\begin{sourceblock}
\includegraphics[page=198,pagebox=mediabox,trim={163.34bp 211.72bp 65.07bp 400.67bp},clip,max width=0.94\textwidth,max height=0.38\textheight]{Marchisio2013.pdf}
\end{sourceblock}
式中，cp f 为气相比热，Lv 为蒸发组分蒸发潜热，pf 为本体气相压力，ps1 为气相蒸发组分分压，R 为理想气体常数。参数 pref 和 Tref 是液体沸腾的参考压力和温度。常数 Mw1 和 Mw2 是两种化学物质的分子量。一旦定义了所有物理常数并指定了 Tp、pf、Tf 和 Yf1（通过求解液滴所见气相的输运方程），就可以用三个未知数 Ys1、Ts 和 ps1 求解这三个方程。了解这三个量足以计算从液滴到气相的质量和焓传递。液滴质量的变化率（假设为球形液滴）由下式给出
{\small dMp\par}
\begin{sourceblock}
\includegraphics[page=198,pagebox=mediabox,trim={191.14bp 50.16bp 65.07bp 632.93bp},clip,max width=0.94\textwidth,max height=0.38\textheight]{Marchisio2013.pdf}
\end{sourceblock}
其中，ρf 为气相物质密度，D 为蒸发物质分子在气相中的扩散系数，系数 Bm 定义为

\begin{sourceblock}
\includegraphics[page=199,pagebox=mediabox,trim={212.77bp 580.50bp 67.56bp 92.97bp},clip,max width=0.94\textwidth,max height=0.38\textheight]{Marchisio2013.pdf}
\end{sourceblock}
这一特定系统的舍伍德数可以通过使用经验相关性来计算，例如 Faeth 和 Fendell 提出的相关性（Kuo，1986）：

Sh = 2 + 0.55RepSc

\begin{sourceblock}
\includegraphics[page=199,pagebox=mediabox,trim={222.52bp 509.88bp 67.56bp 165.95bp},clip,max width=0.94\textwidth,max height=0.38\textheight]{Marchisio2013.pdf}
\end{sourceblock}
为了消除对方程右侧的显式液滴尺寸依赖性。 （5.42），由于质量交换引起的连续变化率通常以液滴表面积来表示（总是假设为球形）：

\begin{sourceblock}
\includegraphics[page=199,pagebox=mediabox,trim={189.45bp 428.58bp 67.56bp 244.89bp},clip,max width=0.94\textwidth,max height=0.38\textheight]{Marchisio2013.pdf}
\end{sourceblock}
确定方程时的隐含假设。 （5.42）和（5.45）中液滴内的传质和传热阻力为零。这是一个常见的假设，基于液滴内准稳态条件的想法，通常被称为快速混合模型。在 Bm 近似恒定的特殊情况下，方程的右侧： (5.45) 是常数，因此表面积的变化率与液滴直径无关（即 d2 蒸发定律）。或者，可以使用薄皮模型，该模型假设无限扩散阻力，仅液滴表面上的薄层被加热和蒸发。第二种假设更适合液滴在蒸发过程开始后立即达到沸腾温度的情况。在多组分液滴的情况下推荐的所有其他方法（例如扩散限制模型）需要液滴内部某种形式的空间离散化，从而导致更高的计算成本。为了闭合方程组，有必要描述第 5.2.3 节中讨论的液滴温度 Tp 的演变。例如，在快速混合模型中考虑蒸发对热通量的影响，颗粒-焓平衡得出

\begin{sourceblock}
\includegraphics[page=199,pagebox=mediabox,trim={197.93bp 188.91bp 67.56bp 473.96bp},clip,max width=0.94\textwidth,max height=0.38\textheight]{Marchisio2013.pdf}
\end{sourceblock}
dTp dt = 6LeNu
其中右侧的最后一项取决于等式： (5.45) cp p 是颗粒的比热。路易斯数定义为

\begin{sourceblock}
\includegraphics[page=199,pagebox=mediabox,trim={213.86bp 121.69bp 216.38bp 550.38bp},clip,max width=0.94\textwidth,max height=0.38\textheight]{Marchisio2013.pdf}
\end{sourceblock}
这个特定系统的努塞尔数可以通过使用 Faeth 和 Fendell 提出的经验相关性来计算（Kuo，1986）：

{\small Nu = 2 + 0.55RepPr\par}
\begin{sourceblock}
\includegraphics[page=199,pagebox=mediabox,trim={223.91bp 51.85bp 67.56bp 623.97bp},clip,max width=0.94\textwidth,max height=0.38\textheight]{Marchisio2013.pdf}
\end{sourceblock}
给定 Tp，ps1 的表达式是封闭的，从而固定传质速率。上述讨论适用于单组分液滴。在许多应用中，液/气相将包含多种化学物质，因此需要额外的内坐标来描述蒸发物理（Sazhin，2006）。在多组分液滴的单颗粒模型中，最简单的介观模型必须包括颗粒质量 Mp、组分质量分数 Yp 和 Yf 以及温度 Tp 和 Tf。总之，我们提醒读者，在质量和能量传递的单粒子模型中，任何一个系数和无量纲数对于不同的粒子都可能不同。因此，一般来说，所有系数和无量纲数都是介观变量的函数（U*

\begin{sourceblock}
\includegraphics[page=200,pagebox=mediabox,trim={297.47bp 504.57bp 127.89bp 179.36bp},clip,max width=0.94\textwidth,max height=0.38\textheight]{Marchisio2013.pdf}
\end{sourceblock}
f）。然而，在实践中，本节中描述的许多相关性是在仅使用状态空间变量的平均值的宏观模型的背景下开发的。因此，除非使用单分散数据拟合相关性（实验数据通常不是这种情况），否则当将相关性应用于多分散系统时，可能需要对相关性中出现的数值常数进行一些调整。另一方面，使用微尺度模型的直接数值模拟开发的相关性通常专门适用于单分散颗粒（例如等直径球体）。理想情况下，可以通过对相互作用的多分散颗粒的模拟，为每个内坐标值（即颗粒直径）开发单独的相关性。无论如何，在将可用相关性应用于单粒子模型之前，读者应该敏锐地意识到可用相关性的适用范围。最后，读者应该记得，流体相平流 G* f 的介观模型必须与 G* p 的粒子相模型一致，以便整个系统在介观尺度上守恒质量和能量。

\subsection*{5.2.6 质量传递引起的动量变化}\phantomsection\addcontentsline{toc}{subsection}{5.2.6 质量传递引起的动量变化}
当颗粒相的速度与流体相的速度不同时，相间的质量传递也会导致动量的传递。例如，如果我们令 xi* p1 为粒子的质量，并 xi* f1 为粒子所见的流体质量，则介观尺度的质量守恒导致

\begin{sourceblock}
\includegraphics[page=200,pagebox=mediabox,trim={181.31bp 220.69bp 65.07bp 449.78bp},clip,max width=0.94\textwidth,max height=0.38\textheight]{Marchisio2013.pdf}
\end{sourceblock}
其中 G* p1 是从流体到颗粒的传质速率（即负号表示从颗粒到流体的传质）。19 每个相相应的动量变化可以写为

\begin{sourceblock}
\includegraphics[page=200,pagebox=mediabox,trim={159.13bp 113.96bp 65.07bp 530.89bp},clip,max width=0.94\textwidth,max height=0.38\textheight]{Marchisio2013.pdf}
\end{sourceblock}
其中 G* p1 的符号决定了由于转移质量而产生的动量贡献，并且通过构造，动量在介观尺度上守恒。通过操纵

\footnotetext[19]{如果粒子成核/溶解也存在，则其对质量/动量传递的贡献将简单地添加到 G*}
{\small 本节中出现的所有公式中的 p1。\par}
等式。 （5.48）和（5.49），对于质量传递引起的加速度可得出以下表达式：

\begin{sourceblock}
\includegraphics[page=201,pagebox=mediabox,trim={155.69bp 551.85bp 67.56bp 84.88bp},clip,max width=0.94\textwidth,max height=0.38\textheight]{Marchisio2013.pdf}
\end{sourceblock}
因此，如上所述，动量传递的介观模型 A* fp 和 A* pf 将包含由于质量传递而产生的贡献，只要 U*

\begin{sourceblock}
\includegraphics[page=201,pagebox=mediabox,trim={285.22bp 515.82bp 176.69bp 167.78bp},clip,max width=0.94\textwidth,max height=0.38\textheight]{Marchisio2013.pdf}
\end{sourceblock}
\section*{5.3 相空间平流：动量传递}\phantomsection\addcontentsline{toc}{section}{5.3 相空间平流：动量传递}
分散相中颗粒动量的变化通常是由于体积力和流体-颗粒相互作用力造成的。在身体力中，重力可能是最重要的。然而，由于体积力单独作用于每个相，因此它们不会导致相之间的动量传递。相反，流体-粒子力导致连续相和分散相之间的动量传递。其中最重要的是浮力和阻力，出于下面将变得更清楚的原因，必须以一致的方式定义它们。然而，正如 Maxey 和 Riley (1983) 的工作中所详述的，附加力会影响分散相中粒子的运动，例如附加质量或虚拟质量力 (Auton 等人，1988)、萨夫曼升力 (Saffman，1965)、Basset 历史术语以及布朗力和热泳力。 20 所有这些力及其方程都将在以下部分中讨论。量化将被提出和讨论。根据所研究的多相系统的不同，相间动量传递的介观尺度模型有很大差异，并且针对不同的系统开发了不同的半经验关系。由于分散相的性质特别重要，可用的介观模型通常分为适用于流体-流体系统的模型和适用于流体-固体系统的模型。主要区别在于，在流体-流体系统中，分散相的元素是可变形颗粒（即气泡或液滴），而在流体-固体系统中，分散相由恒定形状的颗粒构成。下面报道的介观模型适用的典型流体-流体系统是气-液、液-液和液-气系统。所报道的流固系统的介观模型对于气固系统和液固系统都有效。作为一般规则，Afp 的介观模型应从单粒子动量平衡开始推导：

4 作用在粒子上的力 5 + 4 由于传质引起的动量变化 5 , dMpU* p dt =
20 值得一提的是，布朗力和热泳力经常被忽视为实际力，仅量化它们对粒子运动的影响。例如，由于布朗力是引起粒子随机运动的随机力，因此它被视为扩散过程。因此，这种现象通常被视为物理空间扩散过程，详细信息将在下一节中给出。类似地，由于热泳力导致物理空间中由温度梯度驱动的下一个粒子通量，因此它通常被表示为物理空间的平流过程。此外，对于该术语，其他详细信息将在下一节中报告。

其中 Mp 是粒子质量。然而，在下文中，我们将假设粒子质量是恒定的，并且不存在由于质量传递而引起的净动量交换。21 即使在这种情况下，我们可以写出 Mp ̇U* p = - 4 作用于粒子 5 的力，在定义 Afp 的单粒子模型时也必须小心。这是因为右侧的力（例如虚拟质量和历史力）可能涉及时间导数 ̇U* p，因此动量平衡可能是 ̇U* p 的隐函数。然而，通常可以将右侧的力最多表示为 ̇U* p 的线性函数，在这种情况下，可以将单粒子动量平衡写为显式形式 ̇U* p = U* p(t) 和其他介观尺度变量的函数，这是方程 1 隐含的形式。 （5.9）。下面，我们分别讨论单粒子动量平衡中出现的各个力。因为大多数这些力都是 U* p(t) 的显式函数，所以它们直接（在没有其他力的情况下）得出 Afp 模型。尽管如此，读者应该记住，当包含虚拟质量力（取决于 U* p）时，有必要修改 Afp 以包含附加质量因子。

\subsection*{5.3.1 浮力和阻力}\phantomsection\addcontentsline{toc}{subsection}{5.3.1 浮力和阻力}
在分散多相系统中必须考虑的许多相关的流体-颗粒相互作用力中，最重要的是颗粒表面上连续相的应力产生的，从而产生阻力和浮力。对于球形孤立颗粒在以非常小的滑移速度为特征的层流条件下流动的情况，22 该问题有一个解析解，得出作用在直径为 dp 的球体上的浮力和阻力的表达式如下：

\begin{sourceblock}
\includegraphics[page=202,pagebox=mediabox,trim={164.04bp 320.64bp 65.07bp 346.60bp},clip,max width=0.94\textwidth,max height=0.38\textheight]{Marchisio2013.pdf}
\end{sourceblock}
Ffp = −
其中第一项和第二项是由作用在颗粒表面上的应力张量的法向分量引起的，而最后一项是由切向分量引起的。第一项是浮力；当流体和颗粒具有相同的速度（即零滑移速度，或 U* f = U* p）时，它等于被颗粒排开的流体的重量。在浮力中，流体压力梯度在粒子位置 X*(t) 处进行评估。第二项和第三项分别称为形式阻力和摩擦阻力。这两项通常组合在一起形成动阻力，得出以下表达式：

\begin{sourceblock}
\includegraphics[page=202,pagebox=mediabox,trim={154.51bp 177.96bp 65.07bp 486.76bp},clip,max width=0.94\textwidth,max height=0.38\textheight]{Marchisio2013.pdf}
\end{sourceblock}
Ffp = −
正如已经提到的，这仅对于滑移速度非常小的孤立球体有效。无量纲项 24νf/(|U*

\begin{sourceblock}
\includegraphics[page=202,pagebox=mediabox,trim={276.63bp 144.56bp 190.09bp 539.04bp},clip,max width=0.94\textwidth,max height=0.38\textheight]{Marchisio2013.pdf}
\end{sourceblock}
p|dp) 在等式中(5.52)代表摩擦系数，通常称为阻力系数，用CD表示。请注意，如果暴露于流 AD 的颗粒体积 Vp 和横截面表面积与球体不同，我们可以将阻力概括为
21 恒定的粒子质量并不意味着质量传递导致的净动量变化为零。例如，粒子可能以相同的速率失去和获得质量，因此 dMp/dt = 0，但由于速度不同，失去和获得的质量的动量不必相同。

\footnotetext[22]{这种惯性项不重要的非常慢的流动条件也称为蠕动流或斯托克斯流。}
\begin{sourceblock}
\includegraphics[page=203,pagebox=mediabox,trim={187.20bp 616.28bp 67.56bp 55.47bp},clip,max width=0.94\textwidth,max height=0.38\textheight]{Marchisio2013.pdf}
\end{sourceblock}
Ffp = −Vp
然而，有必要将 Vp 和 AD 与用于描述颗粒尺寸的内坐标联系起来（参见第 5.2.1 节）。通过除以方程。通过粒子质量 (ρpVp)（5.53），可以很容易地计算出由于浮力和阻力而产生的粒子加速度，并且通过假设所有粒子在统计上是相同的，可以针对孤立球体 (αf = 1) 获得以下纯平流速度表达式：

\begin{sourceblock}
\includegraphics[page=203,pagebox=mediabox,trim={177.75bp 511.12bp 67.56bp 155.92bp},clip,max width=0.94\textwidth,max height=0.38\textheight]{Marchisio2013.pdf}
\end{sourceblock}
Afp = −1
其中球体的 AD/Vp = 3/(2dp)。很明显，在蠕动（或斯托克斯）流动的情况下，当方程： (5.52) 成立，则阻力系数的以下定义成立：

\begin{sourceblock}
\includegraphics[page=203,pagebox=mediabox,trim={192.04bp 438.87bp 67.56bp 232.27bp},clip,max width=0.94\textwidth,max height=0.38\textheight]{Marchisio2013.pdf}
\end{sourceblock}
值得注意的是，CD 取决于介观速度 U*p 和 U*f，以及多分散颗粒的内坐标 dp。为了避免麻烦的符号，我们不表示 Afp 对所有内部坐标的显式依赖。尽管如此，读者应该记住，方程中出现的任何粒子/流体参数。对于不同的粒子，方程（5.54）可能会有所不同。当应用方程式中的单粒子模型时。对于多分散系统（5.54），主要挑战是找到考虑颗粒特性（即形状、尺寸、密度等）多分散性的阻力系数的准确表达式。方程中的表达式。当然，方程（5.55）仅对于滑移速度非常小的球体有效，并且通常假设其对于Rep ≤0.1有效。对于较高粒子雷诺数的球形粒子，可以使用以下修正：

\begin{sourceblock}
\includegraphics[page=203,pagebox=mediabox,trim={131.43bp 258.56bp 67.56bp 404.15bp},clip,max width=0.94\textwidth,max height=0.38\textheight]{Marchisio2013.pdf}
\end{sourceblock}
这些修正考虑到以下事实：颗粒周围的主要连续相的流动不够慢，无法满足用于推导斯托克斯阻力定律的假设，并且在某一点，流动开始分离，并在球体后面形成涡流。涡流的存在减少了尾流并增加了形状阻力。在 CD 与 Rep 曲线的平坦部分中，阻力几乎完全由形状阻力造成。在这个区域，称为牛顿阻力定律的适用区域，阻力系数可以近似为常数值 0.445。当粒子雷诺数进一步增加时，边界层变得湍流，分离点向后移动，阻力系数急剧降低。这种转变发生在临界雷诺数处，通常假设该雷诺数等于 3.5×105。以较高Rep为特征的流动通常被称为超临界流动，阻力系数从0.445急剧下降至0.07。另一个非常简单的相关性类似于等式中报告的。 (5.56) 是 Schiller \& Naumann (1935) 的阻力模型：

\begin{sourceblock}
\includegraphics[page=203,pagebox=mediabox,trim={139.40bp 51.14bp 67.56bp 611.37bp},clip,max width=0.94\textwidth,max height=0.38\textheight]{Marchisio2013.pdf}
\end{sourceblock}
其中粒子雷诺数照常定义。 Morsi \& Alexander (1972) 提出了另一种更详细的阻力系数随粒子雷诺数变化的描述：

\begin{sourceblock}
\includegraphics[page=204,pagebox=mediabox,trim={200.63bp 572.22bp 65.07bp 99.35bp},clip,max width=0.94\textwidth,max height=0.38\textheight]{Marchisio2013.pdf}
\end{sourceblock}
其中

\begin{sourceblock}
\includegraphics[page=204,pagebox=mediabox,trim={91.56bp 441.63bp 65.07bp 142.26bp},clip,max width=0.94\textwidth,max height=0.38\textheight]{Marchisio2013.pdf}
\end{sourceblock}
当周围的连续相可以被视为连续体（而不是稀薄气体）时，这些方程对于孤立的球形颗粒是有效的。事实上，量化构成主相的分子的平均自由程（即两次后续碰撞之间的平均时间间隔）与粒径之间的比率非常重要：

\begin{sourceblock}
\includegraphics[page=204,pagebox=mediabox,trim={227.90bp 349.00bp 65.07bp 324.46bp},clip,max width=0.94\textwidth,max height=0.38\textheight]{Marchisio2013.pdf}
\end{sourceblock}
当这个量远小于1时，用于描述主相和次级相元素之间相互作用的连续体方法是合适的。当这个比率大于1时，主相和次要相元素之间的相互作用必须用冲击颗粒表面和从颗粒表面反弹的单个分子来描述；这就是所谓的自由分子体系。对于中间值，主相和副相之间的相互作用处于滑移状态，其中仍然可以使用主相的连续介质方法，尽管需要进行一些修正以考虑与颗粒表面相邻的流体之间的速度跳跃。此修正得出阻力系数的以下表达式：

\begin{sourceblock}
\includegraphics[page=204,pagebox=mediabox,trim={159.77bp 195.57bp 65.07bp 475.62bp},clip,max width=0.94\textwidth,max height=0.38\textheight]{Marchisio2013.pdf}
\end{sourceblock}
通常称为 Cunningham 校正因子（Cunningham，1910；Knudsen \& Weber，1911；Millikan，1910）。我们参考 Crowe (2006) 的最新工作来回顾最新的修正。23 还得出了其他修正
23 重要的是要提醒读者，非连续效应与连续相的可压缩性有关。事实上，对于理想气体，动力学理论预测：

\begin{sourceblock}
\includegraphics[page=204,pagebox=mediabox,trim={219.30bp 108.93bp 65.05bp 568.23bp},clip,max width=0.94\textwidth,max height=0.38\textheight]{Marchisio2013.pdf}
\end{sourceblock}
其中 Ma 是连续相的马赫数（即滑移速度与声速之间的比率），k 是连续相的比热之间的比率。修正阻力系数C*
{\small 因此，D 是这三个无量纲数中的两个的函数，例如Kn 和 Rep。出于实际目的，任何 Ma < 0.2 的等温流都被认为是不可压缩的，如果 Kn < 10−3，则连续相被认为是连续相。\par}
考虑到颗粒的形状、其他颗粒的存在 (αf < 1) 和壁效应，并且根据多相系统的类型（即流体-固体、流体-流体等）而有很大不同。因此，接下来我们讨论一些选定的示例系统的阻力系数的相关性。

\paragraph*{流体-流体系统的阻力系数}
流体-流体系统中阻力系数 CD 的计算因多种因素而变得复杂。第一个是，在流体-流体系统中，由于流体元素在一个区域中聚集而发生聚结，导致分散相的体积分数非常高（接近于1），这是很常见的。 24另一个问题与这样一个事实有关，即根据其尺寸的不同，流体颗粒可以变形，可以呈现不同的形状，从而导致不同类型的相互作用。接下来的讨论将集中在气-液系统上，但需要谨慎的是，这些结论也适用于液-液和液-气系统。对于分散在连续液相中的气泡，当气泡尺寸非常小或气体和液体之间的表面张力非常高时，呈现球形。在这种情况下，气泡运动或多或少是直线运动，阻力由粘性力控制，因此取决于气泡雷诺数。另一方面，当气泡尺寸较大或表面张力较低时，气泡呈非球形。在这种情况下，气泡运动不再是直线运动，并且阻力由表面张力控制；它不再取决于 Rep，而是取决于 E ̈otv ̈os 数量。另一个重要因素与气液界面上积聚的杂质的存在有关。事实上，实验证据（根据最终速度收集）表明，其中一些杂质不会显著改变宏观特性，例如主相或分散相的密度或粘度或两相之间的界面能，但它们的存在会导致阻力系数发生重要变化。后一个因素对于球形气泡尤其重要并且很难考虑。根据杂质，气液系统大致可分为三类：（i）纯净系统，（ii）轻微污染系统，（iii）完全污染系统。尽管这种区别可能看起来含糊不清，但它在实践中被广泛使用。例如，纯净的空气-水系统可以对应于仔细蒸馏两次或多次的水（即双蒸馏水），受污染的空气-水系统可以对应于自来水中移动的气泡，部分污染的空气-水系统可以对应于中等程度的污染
24 需要提醒读者的是，本书中讨论的多相系统是分散多相系统，这意味着除了主要连续相之外，还存在由颗粒、液滴或气泡构成的次要相。此处描述的情况（分散相的体积分数值接近1）不包含在此类别中；事实上，如果分散相的所有颗粒聚结在一起，则分散相不再分散。这种情况通常由矩值反映，并且通常标记为胶凝或相转化。在凝胶化过程中，NDF 的零阶矩变为零，而其他矩则发散为非物理值。克服这一限制的一种实用方法是在分散相体积分数变得非常高的区域中将矩的演变冻结为其预胶凝值。此外，为了将这些区域与分散相按预期演变的其他区域连接起来，可以方便地采用对阻力系数的修正。正如已经提到的，这些情况在流体-流体系统中非常常见，例如在低压区域（搅拌器叶片或挡板后面）的气-液搅拌罐中或每当发生相转化时（例如在气体和液体之间的界面处）都会遇到这些情况。当然，在非聚结系统的情况下，例如在固体颗粒不聚集在一起的情况下，通过引入使颗粒分开的力可以避免这种情况。这通常是通过引入分散相体积分数的最大填充极限以及与实际分散相体积分数和该极限之间的差值成比例的力来完成的。

（富山等人，1998）。完全污染的气泡表现为固体颗粒；事实上，杂质在气液界面的积累使得界面变得坚硬。相反，在纯系统中，由于不存在杂质，因此可以在气泡内形成再循环模式，从而降低粘性阻力。对于纯系统和孤立气泡，Hadamard (1911) 提出了以下关系：

\begin{sourceblock}
\includegraphics[page=206,pagebox=mediabox,trim={224.59bp 553.15bp 65.07bp 127.18bp},clip,max width=0.94\textwidth,max height=0.38\textheight]{Marchisio2013.pdf}
\end{sourceblock}
这仅对非常小的 Rep 有效。 （5.63）与式（5.63）进行比较。 （5.55）是针对固体颗粒开发的，很明显它会导致阻力值更小。方程中包含的表达式。 (5.63) 对于小 Rep 有效，而对于较高 Rep，必须应用通过线性化纳维-斯托克斯方程中的对流项而得出的校正，从而得到以下表达式：

\begin{sourceblock}
\includegraphics[page=206,pagebox=mediabox,trim={224.60bp 455.14bp 65.07bp 225.19bp},clip,max width=0.94\textwidth,max height=0.38\textheight]{Marchisio2013.pdf}
\end{sourceblock}
并且，使用与等式非常相似的方法。 （5.56），可以使用以下具有更普遍有效性的表达式：

\begin{sourceblock}
\includegraphics[page=206,pagebox=mediabox,trim={206.81bp 391.15bp 65.07bp 280.87bp},clip,max width=0.94\textwidth,max height=0.38\textheight]{Marchisio2013.pdf}
\end{sourceblock}
CD=最小
对于受污染系统中的球形孤立气泡，也开发了类似的表达式：25

\begin{sourceblock}
\includegraphics[page=206,pagebox=mediabox,trim={206.81bp 335.79bp 65.07bp 336.23bp},clip,max width=0.94\textwidth,max height=0.38\textheight]{Marchisio2013.pdf}
\end{sourceblock}
CD=最小
对于较大的非球形气泡，阻力系数不再取决于污染程度和 Rep，而仅取决于 E ̈otv ̈os 数。一种常用的依赖是

\begin{sourceblock}
\includegraphics[page=206,pagebox=mediabox,trim={209.56bp 266.58bp 65.07bp 400.35bp},clip,max width=0.94\textwidth,max height=0.38\textheight]{Marchisio2013.pdf}
\end{sourceblock}
当这些不同的阻力系数组合在一起时，它们通常被称为富山阻力定律（Tomiyama，2004）。总之，对于清洁系统，可以得到以下表达式：

\begin{sourceblock}
\includegraphics[page=206,pagebox=mediabox,trim={169.51bp 194.82bp 65.07bp 472.12bp},clip,max width=0.94\textwidth,max height=0.38\textheight]{Marchisio2013.pdf}
\end{sourceblock}
CD = 最大
用于轻度污染系统的相关性是

\begin{sourceblock}
\includegraphics[page=206,pagebox=mediabox,trim={169.51bp 143.42bp 65.07bp 523.52bp},clip,max width=0.94\textwidth,max height=0.38\textheight]{Marchisio2013.pdf}
\end{sourceblock}
CD = 最大
而对于完全污染的系统，采用以下相关性：

\begin{sourceblock}
\includegraphics[page=206,pagebox=mediabox,trim={195.33bp 92.02bp 65.07bp 574.92bp},clip,max width=0.94\textwidth,max height=0.38\textheight]{Marchisio2013.pdf}
\end{sourceblock}
CD = 最大

\footnotetext[25]{正如刚才提到的，杂质和表面活性剂往往会积聚在界面处，使其变得坚硬，从而抑制内部运动。因此，受污染系统对流体颗粒的阻力系数与固体颗粒的阻力系数相同。}
这些关系对于在层流条件下移动的孤立气泡是有效的。在湍流的情况下，湍流涡流撞击气泡表面的作用是增加阻力。这通常是通过引入有效流体粘度（而不是连续相的分子粘度 μf）来解释的，有效流体粘度定义为 μf,eff= μf + Clρfε1/3 f d4/3 p，其中 εf 是流体相中的湍流耗散率，Cl 是通常取等于 0.02 的常数。这种有效粘度用于计算气泡/颗粒雷诺数（Bakker \& van den Akker，1994），它解释了由于气泡周围动量传输增加而导致的湍流减少，这又与气泡尺寸和湍流长度尺度的比率有关。然而，提醒读者，介观尺度模型不包括宏观尺度湍流，因此，使用基于宏观尺度湍流的有效粘度是不合适的。当每个气泡可以被认为是孤立的时，前面提到的所有相关性在稀释系统中都非常有效。然而，当气泡开始感知彼此的存在26时，所呈现的任何相关性都不会产生可接受的预测。总体性能通常根据阻力定律预测现实容器中分散相总体浓度（例如总体含气量）或预测从一种状态过渡到另一种状态的能力进行评估27，如 Montante 等人报道。 (2007) 和佩蒂蒂等人。 （2009）。定义能够重现实验测量行为的阻力系数的纯经验方法是诉诸终端速度的概念。让我们想象一个只有重力、浮力和阻力作用的单个粒子。粒子移动（仅由重力和浮力驱动），连续相处于静止状态，并在一定时间后达到稳定状态。达到稳定状态时粒子移动的恒定速度称为终端速度，此处用 U∞ 表示。众所周知，在等式中设置 Afp = 0。 (5.54) 并考虑连续相处于静止状态（即 U* f = 0），分散相和连续相之间的速度差就是终端速度，即 U∞= |U*

\begin{sourceblock}
\includegraphics[page=207,pagebox=mediabox,trim={276.63bp 301.90bp 188.86bp 382.60bp},clip,max width=0.94\textwidth,max height=0.38\textheight]{Marchisio2013.pdf}
\end{sourceblock}
f|。这导致阻力系数的近似值如下：

\begin{sourceblock}
\includegraphics[page=207,pagebox=mediabox,trim={202.83bp 258.06bp 67.56bp 413.81bp},clip,max width=0.94\textwidth,max height=0.38\textheight]{Marchisio2013.pdf}
\end{sourceblock}
方程（5.71）允许我们计算特定终端速度的稳态阻力系数。这个公式特别有用，因为它提出了一种计算 CD 的策略，可以轻松应用于非常不同的流体-流体系统。事实上，考虑到正在研究的特定系统，我们可以确定终端速度和
26 这种情况在流体-流体系统中尤其重要，因为分散相的元素在进入物理接触之前就开始相互作用。当总粒子数密度（和分散相体积分数）足够大以允许这种情况发生时，该状态被标记为三向耦合。根据这一定义，当分散相稀释到不影响主要连续相时，系统处于单向耦合条件下：连续相影响分散相，但反之则不然。相反，当第二相浓度增加并且检测到对连续相的一些影响时，系统被认为处于双向耦合条件下。当分散相的元素达到足够高的浓度以允许通过主要连续相相互作用（例如粒子/气泡穿过彼此的尾流）时，系统被定义为处于三向耦合条件下。最终，当分散相的浓度使得分散相元素之间的相互作用相关时，系统被称为处于四向耦合条件下。

\footnotetext[27]{典型例子是气液搅拌罐，根据容器内的气体分布，检测到三种状态，即完全分散、装载和淹没。气液搅拌罐中的其他重要状态转变是指搅拌器叶片后面的气腔的数量和形状，即涡流粘附腔、大腔和参差不齐的腔。}
然后使用该信息计算阻力系数，确保对多相系统中预测的气泡滞留量进行更严格的控制。例如，很容易量化由于其他气泡的存在或由于任何其他影响（例如细颗粒或超细颗粒的存在）而导致的终端速度的变化。实验证据表明，例如，如果一个孤立的气泡以特定的终端速度移动，那么其他气泡的存在会降低终端速度，突出显示与较大阻力系数相关的某种障碍或拥挤效应。使用经验阻力系数的一个重要缺点是它仅对分散相体积分数αp在空间上均匀的均匀系统有效。如果真实阻力系数强烈依赖于局部分散相体积分数，并且系统在空间上不均匀，则方程： （5.71）不会捕获对分散相体积分数的正确依赖性。理想情况下，可以在不同的分散相体积分数下测量 CD，并且可以建立相关性以正确捕获 CD 对 αp 的依赖性。尽管经验方法实施起来很简单并且不是很复杂，但它在实践中非常有效，因为气泡的终端速度可以通过多种不同的方式确定。对于非常流行的系统，例如经典的清洁空气-水系统，气泡终端速度由 Mendelson (1967) 提出的相关性及其扩展给出：

\begin{sourceblock}
\includegraphics[page=208,pagebox=mediabox,trim={96.32bp 358.28bp 65.07bp 275.69bp},clip,max width=0.94\textwidth,max height=0.38\textheight]{Marchisio2013.pdf}
\end{sourceblock}
当然，这仅适用于在清水中移动的孤立气泡，并且应用终端速度的修正（即减小）来考虑其他气泡或部分污染的存在。或者，终端速度可以通过实验或通过详细的模拟来确定。用于确定单个气泡或气泡群（群）的终端速度的模拟可以基于流体体积（VOF）方法或类似方法（Bertola 等人，2004 年；Petitti 等人，2007 年）。

\paragraph*{流固系统的阻力系数}
正如已经报道的那样，固体颗粒的情况可以用比流体-流体系统更简单的方法来处理。 Wen \& Yu (1966) 提出了计算悬浮在流体中的固体颗粒的阻力系数的最流行的相关式之一：

\begin{sourceblock}
\includegraphics[page=208,pagebox=mediabox,trim={174.25bp 168.53bp 65.07bp 506.95bp},clip,max width=0.94\textwidth,max height=0.38\textheight]{Marchisio2013.pdf}
\end{sourceblock}
其中 β = 3.65 是考虑其他粒子存在的校正因子。在这种情况下，当颗粒系统非常稀（即αf ≈1）时，方程： (5.73) 简化为等式。 （5.57）。 Di Felice (1994) 提出了类似的表达式，其中指数现在是粒子雷诺数的函数：

\begin{sourceblock}
\includegraphics[page=208,pagebox=mediabox,trim={165.09bp 88.71bp 65.07bp 585.06bp},clip,max width=0.94\textwidth,max height=0.38\textheight]{Marchisio2013.pdf}
\end{sourceblock}
{\small 当颗粒浓度增加超过某个临界值时（5.73）可能会变得不准确。在这些情况下（这在密集的几乎拥挤的环境中很常见）\par}
流体-固体系统）阻力是通过使用额尔古纳方程计算的，该方程估计流体流经球体填充床时所经历的压降：

\begin{sourceblock}
\includegraphics[page=209,pagebox=mediabox,trim={188.42bp 575.90bp 67.56bp 90.26bp},clip,max width=0.94\textwidth,max height=0.38\textheight]{Marchisio2013.pdf}
\end{sourceblock}
在实践中，等式的组合。 (5.73)和(5.75)用于覆盖αf的所有值。另一种广泛使用的阻力模型是 Syamlal 和 O’Brien (1989) 提出的模型：

\begin{sourceblock}
\includegraphics[page=209,pagebox=mediabox,trim={178.61bp 504.81bp 181.39bp 158.18bp},clip,max width=0.94\textwidth,max height=0.38\textheight]{Marchisio2013.pdf}
\end{sourceblock}
其中粒子终端速度是根据 Garside \& Al-Dibouni (1977) 提出的相关性计算的，U∞= 0.5\{A −0.06Rep + [(0.06Rep)2 + 0.12Rep(2B −A) + A2]1/2\}，其中 A = α4.14 f 且

\begin{sourceblock}
\includegraphics[page=209,pagebox=mediabox,trim={184.25bp 416.68bp 187.64bp 244.39bp},clip,max width=0.94\textwidth,max height=0.38\textheight]{Marchisio2013.pdf}
\end{sourceblock}
尽管这两种表达式都很常用，但它们无法预测固液悬浮液的一些重要宏观特性，例如膨胀和沉降剖面。为了克服这一限制，Mazzei 和 Lettieri (2007) 基于 Richardson 和 Zaki (1954) 描述均匀流体-固体悬浮液的膨胀曲线的经验相关性，建立了阻力系数的关系。其主要特征在于，该表达式在整个流体动力学范围内以及对于悬浮液空隙率的任何值都与 Richardson 和 Zaki 相关性一致。它有以下公式：

\begin{sourceblock}
\includegraphics[page=209,pagebox=mediabox,trim={180.65bp 270.94bp 67.56bp 392.08bp},clip,max width=0.94\textwidth,max height=0.38\textheight]{Marchisio2013.pdf}
\end{sourceblock}
其中

\begin{sourceblock}
\includegraphics[page=209,pagebox=mediabox,trim={220.49bp 225.61bp 67.56bp 447.85bp},clip,max width=0.94\textwidth,max height=0.38\textheight]{Marchisio2013.pdf}
\end{sourceblock}
以及

\begin{sourceblock}
\includegraphics[page=209,pagebox=mediabox,trim={171.89bp 175.25bp 67.56bp 482.09bp},clip,max width=0.94\textwidth,max height=0.38\textheight]{Marchisio2013.pdf}
\end{sourceblock}
其中

\begin{sourceblock}
\includegraphics[page=209,pagebox=mediabox,trim={211.80bp 130.30bp 67.56bp 549.76bp},clip,max width=0.94\textwidth,max height=0.38\textheight]{Marchisio2013.pdf}
\end{sourceblock}
值得注意的是，上面讨论的流体-固体阻力系数仅对单分散颗粒（即具有相同直径和材料密度的颗粒）有效。使用流体-颗粒流微观方程的直接数值模拟（DNS），一些作者（Beetstra 等人，2007；Buhrer-Skinner 等人，2009；Holloway 等人，2010；Tenneti 等人，2010、2012；Yin 和 Sundaresan，2009）提出了改进的阻力系数来解释多分散颗粒。

\paragraph*{由于湍流而进行的阻力系数修正}
正如已经提到的，通常会进行修正以考虑湍流的影响。详细介绍过去几十年中开发的不同治疗方法和矫正方法超出了本书的范围（Balachandar \& Eaton，2010）。在不同的方法中，值得一提的是完全解析的模拟、28 拉格朗日点粒子模拟、29 以及不同的欧拉方法（Prosperetti \& Tryggvason，2009）。然而，读者应该记住，本章讨论的介观尺度模型没有考虑宏观尺度湍流。因此，在介观尺度建模的背景下，应谨慎使用下面讨论的与经典湍流模型一起使用的相关性。一般来说，校正阻力系数的主要因素有两个，即湍流强度和湍流长度尺度与颗粒尺寸之间的长度尺度比。当颗粒小于柯尔莫哥洛夫长度尺度（即 dp ≪ηk）时，通常采用上面报告的标准系数，因为环境流动对颗粒来说几乎是层流。对于尺寸与流体相流动尺度（即 dp ≥ηk）相当的较大颗粒，颗粒看到的环境流动可能更复杂。事实上，由于在湍流中存在级联的长度尺度，因此通常存在多种效应。直观上，小于颗粒尺寸的湍流结构会影响颗粒周围的边界层，而远大于颗粒尺寸的湍流结构会导致颗粒堆积，也称为惯性聚类。根据所采用的方法，在阻力系数的修正中必须仅考虑第一个效应或两种效应。湍流的存在通常会降低临界粒子雷诺数（Crowe 等人，1998）。临界雷诺数 30 Rec p 通常被认为是阻力系数达到 0.3 时的数值。 Rep 的进一步增加首先导致阻力系数减小，直到达到最小值，然后阻力系数开始再次增加。 CD 再次达到 0.3 时的 Rep 值称为元临界雷诺数，通常写为 ReM p。湍流作用下的临界粒子雷诺数可按下式计算(Crowe et al., 1998)：

\begin{sourceblock}
\includegraphics[page=210,pagebox=mediabox,trim={189.49bp 245.76bp 65.07bp 416.95bp},clip,max width=0.94\textwidth,max height=0.38\textheight]{Marchisio2013.pdf}
\end{sourceblock}
记录记录

其中Ir是相对湍流强度，定义为湍流动能与滑移速度大小的比值。类似的表达式对于超临界粒子雷诺数有效：

\begin{sourceblock}
\includegraphics[page=210,pagebox=mediabox,trim={188.13bp 160.89bp 65.07bp 501.81bp},clip,max width=0.94\textwidth,max height=0.38\textheight]{Marchisio2013.pdf}
\end{sourceblock}
对数ReM

28 对于完全解析的模拟，我们通常指的是流固界面的详细描述，其中环境湍流的所有尺度以及粒子引入的流动尺度都通过对流体相采用 DNS 和对粒子采用拉格朗日跟踪来完全解析。 29 在这种方法中，流体相通常用 DNS 处理，分散相用点粒子描述，其演化通常用拉格朗日跟踪处理。在这种情况下，由于流体-固体界面未解析，因此需要某种形式的校正。

\footnotetext[30]{在处理湍流时，所有相关的无量纲数均使用可用数量进行评估。例如，在 DNS 中，将采用流体和粒子瞬时速度，而在大涡模拟 (LES) 或雷诺平均纳维斯托克斯方程 (RANS) 模拟中，将使用过滤值或雷诺平均值。}
r Rep−1（对于 Rep < 50 和 0.05 < Ir < 0.5），并且 CD = 0.133(1 + 150/Rep)1.565 + 4Ir（对于 50 < Rep < 700 和 0.05 < Ir < 0.5）。对于 Rec 之间具有中间值的粒子雷诺数，所得阻力系数可以计算为 CD = 162I1/3 p 和 ReM p，可以使用以下表达式： CD = 0.3(Rep/Rec p)−3 for 0.9Rec p < Rep < Rem p，其中 Rem p )(3.45+20Ir) = (Rec p)3(ReM p 可以根据表达式 (Rem

\begin{sourceblock}
\includegraphics[page=211,pagebox=mediabox,trim={107.85bp 546.60bp 70.55bp 124.11bp},clip,max width=0.94\textwidth,max height=0.38\textheight]{Marchisio2013.pdf}
\end{sourceblock}
p。关于阻力系数对湍流的依赖性的不同数据之间存在相当大的差异，但 Torobin \& Gauvin (1960) 和 Clamen \& Gauvin (1969) 开发的相关性似乎是最合适的。在开发的不同方法中，值得一提的是，由于惯性聚类（由介观尺度模型精确捕获），必须在宏观模型中考虑虚拟颗粒尺寸（Xu＆Kato，1999）。为了结束我们对阻力系数的讨论，重要的是要提醒读者，几乎所有现有的相关性都是在相间平均动量传递的宏观模型的背景下发展起来的。因此，当在单粒子模型的背景下使用时，我们不能指望它们完全捕获对介观变量的所有依赖性。尽管如此，最近使用微尺度模型 DNS 的工作（Beetstra 等人，2007 年；Buhrer-Skinner 等人，2009 年；Holloway 等人，2010 年；Yin 和 Sundare-san，2009 年）为改善可用阻力相关性带来了巨大希望。 Tenneti 等人的工作。 （2010）关于将随机分量纳入阻力模型的观点尤其值得注意。

\subsection*{5.3.2 虚拟质量和升力}\phantomsection\addcontentsline{toc}{subsection}{5.3.2 虚拟质量和升力}
虚拟质量（或附加质量）力对应于添加到粒子上的惯性，因为当粒子在流体中移动时，周围流体的体积会随粒子加速或减速。为了简单起见，虚拟质量力通常使用随粒子移动的流体质量来建模；然而，实际上大部分流体都被加速，尽管程度不同。从历史上看，虚拟质量力是在研究流体中的摆运动后首次报道的（Stokes，1851）。由此产生的粒子加速度以及粒子速度的变化率（假设所有粒子在统计上是相等的）通过以下力平衡进行量化（Auton，1987）：

\begin{sourceblock}
\includegraphics[page=211,pagebox=mediabox,trim={170.57bp 192.11bp 67.56bp 472.00bp},clip,max width=0.94\textwidth,max height=0.38\textheight]{Marchisio2013.pdf}
\end{sourceblock}
其中 CvmVp 是体积为 Vp 的粒子加速的流体体积，DU* f/Dt 是粒子沿流体粒子轨迹看到的流体速度的变化率，我们可以近似为

\begin{sourceblock}
\includegraphics[page=211,pagebox=mediabox,trim={181.13bp 115.86bp 67.56bp 555.89bp},clip,max width=0.94\textwidth,max height=0.38\textheight]{Marchisio2013.pdf}
\end{sourceblock}
在介观尺度上调用动量守恒，可以得出粒子所看到的流体速度的以下表达式：

\begin{sourceblock}
\includegraphics[page=211,pagebox=mediabox,trim={168.33bp 53.44bp 67.56bp 610.67bp},clip,max width=0.94\textwidth,max height=0.38\textheight]{Marchisio2013.pdf}
\end{sourceblock}
其中 Vf 是粒子看到的流体体积（与 CvmVp 不同）。正如第 4.3.4 节中所讨论的（其中我们用 ψf1 表示 Vf），流体相体积分数与 Vf 之间的关系如下： （4.54）在第 117 页，平均流体速度 Uf 是通过方程（4.54）得出的。 (4.81) 第 122 页。在单动流体极限中，我们可以通过沿流体粒子轨迹的平均流体速度 Uf(t, x) 的变化率来近似 DU* f/Dt：

\begin{sourceblock}
\includegraphics[page=212,pagebox=mediabox,trim={186.57bp 541.66bp 65.07bp 122.97bp},clip,max width=0.94\textwidth,max height=0.38\textheight]{Marchisio2013.pdf}
\end{sourceblock}
收集方程中的项。 （5.82）涉及 U* p 的时间导数，得出单动流体极限中虚拟质量力（在没有其他力的情况下31）的介观模型：

\begin{sourceblock}
\includegraphics[page=212,pagebox=mediabox,trim={220.43bp 471.99bp 65.07bp 193.82bp},clip,max width=0.94\textwidth,max height=0.38\textheight]{Marchisio2013.pdf}
\end{sourceblock}
Afp =
其中 Cvm 是虚拟质量系数，通常与形状相关（即，它取决于描述颗粒形状的内坐标），并且通常包括对流体相体积分数 αf 的依赖性（Auton，1987）。最常见的值为 0.5，由 Drew 等人提出。 (1979) 对于刚性球形颗粒，而 0.25 用于水中的气泡 (Cook \& Harlow, 1986)。还发现该系数随分散相元素的局部浓度而变化（Biesheuvel \& Spoel-stra，1989；Homsy 等，1980）。例如，Sankaranarayanan 等人。 (2002) 开发了虚拟质量系数模型，并观察到它随着分散相体积分数的增加而几乎线性增加。从等式中可以明显看出。根据（5.86），当连续相的密度高于分散相的密度（即ρf≫ρp）时，虚拟质量力尤其重要，气液系统就是典型的例子。在气液系统中，虚拟质量力被发现对于正确预测状态转变和流体动力学行为非常重要，特别是在气泡塔中（例如，参见 Monahan \& Fox (2007) 和 Monahan et al. (2005) 的工作）。在气液搅拌罐的情况下，阻力似乎是主导力。在具有平均剪切力的流体中移动的颗粒会受到垂直于流体流动方向的升力。剪切升力源自颗粒周围粘性流的惯性效应，并且取决于在颗粒位置 x = X*(t) 处评估的流体相32的平均涡度。对于球形粒子，升力（也称为萨夫曼升力）引起的粒子加速度等于（Auton，1987；Drew \& Lahey，1993；Drew \& Passman，1999；Saffman，1965）

\begin{sourceblock}
\includegraphics[page=212,pagebox=mediabox,trim={215.23bp 184.66bp 65.07bp 486.30bp},clip,max width=0.94\textwidth,max height=0.38\textheight]{Marchisio2013.pdf}
\end{sourceblock}
法新社 = CL
其中 CL 是可以从理论推导出来的升力系数（Maxey \& Riley，1983）。在文献中可以找到关于升力的重要性和意义的相互矛盾的结论。例如，Delnoij 等人。 （1997）报告说，电梯
31 这是非常重要的一点，因为当包含其他力时，附加质量项将修改 Afp 的模型。事实上，对于一般公式，应该从方程（1）中的单粒子动量平衡开始。 (5.82) 并在右侧添加其他力。 Afp 的最终介观尺度模型将右侧的所有项乘以附加质量因子 Cvmρf/(ρp + Cvmρf)。换句话说，由于增加了质量，无法通过简单地将各个力的模型相加来找到 Afp。更多详细信息，请参见第 5.3.4 节。

\footnotetext[32]{流体相的平均涡度是平均流体速度场的旋度，ωc = ∇× Uf，是连续相平均旋转的量度。}
力似乎对于准确表示气液系统中的流动行为和 Monahan 等人的工作是必要的。 (2005)证明了使用欧拉-欧拉双流体模型时升力对均匀流动预测稳定性的重要性。 Auton (1987) 表明 CL 应取等于 0.5 以确保与实验的良好一致性。其他地方也报道了相同的值（Deen 等人，2001；Lain 等人，2002），而 Drew 和 Passman（1999）建议使用 CL = 0.25。富山等人。 (2002b) 确定了对气液系统有效的净横向升力系数的经验相关性。对于 dp < 4.4 mm，发现 CL 是 Rep 的函数，而对于 dp > 4.4 mm，发现 CL 是 E ̈otv ̈os 数的函数。此外，当 dp = 5.8 mm 时，CL 的符号从正变为负。已经提出了基于毛细管数的其他类似相关性（Sankara-narayanan \& Sundaresan，2002），表明对于大多数应用来说 CL 应为正。总结上面介绍的力，在不同的多相系统中进行的测试表明，所涉及的不同力的重要性顺序通常将浮力和阻力排在第一位33，然后是液-固系统的升力和虚拟质量力，以及液-液系统的虚拟质量和升力（例如，参见 Diaz 等人（2008）和 Barton（1995）对非阻力的研究），而最常见的力则为浮力和阻力。对于液-液和液-固系统，相应常数的值为 CL = 0.25 和 Cvm = 0.5。当然，由于在计算代码中实现所有力很简单（Vikas 等人，2011b），因此为了通用性，最好将它们全部包括在内。

\subsection*{5.3.3 Boussinesq——巴塞特力、布朗力和热泳力}\phantomsection\addcontentsline{toc}{subsection}{5.3.3 Boussinesq——巴塞特力、布朗力和热泳力}
正如已经提到的，在流体中运动的物体的加速度产生称为虚拟质量力的力。虚拟质量力是由于施加到粒子周围的流体的瞬时加速度而产生的。然而，粒子加速会产生与粒子及其尾迹周围边界层的瞬态发展相关的附加效应。这些附加效应首先由 Boussinesq (1885) 和 Basset (1888) 在加速斯托克斯流中发现，通常被量化为单一力，有时被称为 Boussinesq-Basset 或历史力 (Crowe et al., 1998)。该术语描述了由于颗粒在流体中移动加速而导致边界层发展滞后而产生的力。换句话说，它解释了由于颗粒边界层和尾流的时间发展而产生的不稳定粘性阻力效应，因此它解决了当相对流体-颗粒速度随时间变化时边界层发展的时间延迟。 Boussinesq-Basset 力（以相应的加速度表示）可以写为

\begin{sourceblock}
\includegraphics[page=213,pagebox=mediabox,trim={146.05bp 145.68bp 67.56bp 517.13bp},clip,max width=0.94\textwidth,max height=0.38\textheight]{Marchisio2013.pdf}
\end{sourceblock}
该表达式经常被重写，以突出历史力内核的存在：

\begin{sourceblock}
\includegraphics[page=213,pagebox=mediabox,trim={148.48bp 84.76bp 67.56bp 578.38bp},clip,max width=0.94\textwidth,max height=0.38\textheight]{Marchisio2013.pdf}
\end{sourceblock}
\footnotetext[33]{对于 ρf ≪ρp 的气体-粒子系统，只有阻力很重要，因为所有其他力都与 ρf/ρp 比率成正比。}
其中历史力内核为 KB(s) = dp/(4πνfs)1/2，表示一种对粒子边界层和尾流的历史进行加权的滤波器。 Boussinesq-Basset 力很难在单粒子模型中实现，因为它在 ̇U* p(t) 中并不明确。因此，尽管已经开发了一些有效的计算策略（Dorgan \& Loth，2007；van Hinsberg et al.，2011），但由于多相流代码中的实际原因，它通常被忽略。尽管如此，应该指出的是，当粒子以极高的速率加速时，布辛涅斯克-巴塞特力可能非常大。布朗力是众所周知的力，对于悬浮在连续相中的非常小的颗粒来说，它变得很重要。布朗力可以定义为由于连续相分子与悬浮颗粒碰撞而产生的瞬时动量交换。当粒子小到足以感知连续相的分子性质（和运动）时（即粒子克努森数足够大时），它就会表现出随机运动，这一点早在2000年前就被罗马卢克莱修观察到了。布朗力通常被描述为随机过程（Gardiner，2004），并且可以将其建模为维纳过程：34

\begin{sourceblock}
\includegraphics[page=214,pagebox=mediabox,trim={182.01bp 424.04bp 65.07bp 246.93bp},clip,max width=0.94\textwidth,max height=0.38\textheight]{Marchisio2013.pdf}
\end{sourceblock}
其中Mp是粒子质量，kB是玻尔兹曼常数，Tf是连续相的绝对温度。维纳过程 W(t) 在第 5.1.2 节中介绍。比较方程。 (5.90) 与等式。 （5.10），我们可以观察到等式： (5.90) 表示速度相空间中的各向同性扩散系数。在实践中，由于布朗运动仅对非常小的粒子有意义，因此通常将其作为零粒子斯托克斯数极限内粒子浓度的物理空间扩散项。由于这个原因，等式。 （5.90）在实际计算中很少使用。热泳（或索雷效应）是当颗粒在流体相中经受大温度梯度时观察到的现象。物理解释非常简单。众所周知并且已经讨论过，由于主相分子与颗粒之间的碰撞，颗粒与周围流体交换动量。这些碰撞的强度和频率当然与流体相的温度有关。与高流体相温度区域相邻的低流体相温度区域中的颗粒将与热区域中的分子交换比与冷区域中的分子更多的动量。这种动量交换的不平衡转化为将粒子从高温区域推向低温区域的力。这种力对于气溶胶以及一般情况下在温度梯度较大的气体中移动的小颗粒尤其重要。典型的例子是在火焰中形成和/或传输的颗粒。热泳的量化取决于颗粒尺寸与（流体相分子的）平均自由程之比。该比率是方程式中引入的液相克努森数的倒数。 （5.60）。当比值小于1（努森数大于1）时，热泳加速度的单粒子模型一般表示为35

\begin{sourceblock}
\includegraphics[page=214,pagebox=mediabox,trim={236.97bp 104.21bp 65.07bp 569.25bp},clip,max width=0.94\textwidth,max height=0.38\textheight]{Marchisio2013.pdf}
\end{sourceblock}
\footnotetext[34]{从技术上讲，布朗力不是一个纯粹的平流术语，可以包含在第 5.5 节中扩散过程的讨论中。}
\footnotetext[35]{从技术上讲，热泳并不是一个真正的纯平流术语，可以包含在第 6 章关于碰撞过程的讨论中。}
其中 pf 和 Tf 是在颗粒位置处计算的流体相压力和温度，λf 是流体相分子的平均自由程。该方程对于球形颗粒有效，但可以扩展到不同的形状。当颗粒尺寸与平均自由程之比大于1（或者相反，努森数小于1）时，由于颗粒内建立了温度梯度，该机制更加复杂。对于这种情况，通常采用以下表达式：

\begin{sourceblock}
\includegraphics[page=215,pagebox=mediabox,trim={133.64bp 528.11bp 67.56bp 143.25bp},clip,max width=0.94\textwidth,max height=0.38\textheight]{Marchisio2013.pdf}
\end{sourceblock}
其中kf和kp分别是流体和颗粒的热导率，Cs = 1.17是热滑移系数，Ct = 2.18是热交换系数，Cm = 1.14是动量交换系数。关于这个主题的更多细节可以在 Talbot (1981) 中找到。与布朗力一样，热泳力对于非常小的颗粒很重要，并且通常作为物理空间扩散项包含在颗粒浓度输运方程中。总之，由于不同的原因，Boussinesq-Basset、布朗和热泳力很少用于分散型多相流模拟。布辛涅斯克-巴塞特力被忽略，因为只有快速加速的粒子才需要它，而且它的形式使其模拟难以实现。布朗力和热泳力对于非常小的颗粒很重要，这通常意味着颗粒斯托克斯数接近于零。对于此类粒子，无需求解分散相动量密度的输运方程。相反，布朗力和热泳力在粒子浓度传输方程（与流体相动量方程耦合）中生成物理空间扩散项。

\subsection*{5.3.4 介观尺度加速度模型的最终表达式}\phantomsection\addcontentsline{toc}{subsection}{5.3.4 介观尺度加速度模型的最终表达式}
为了完成我们对动量传递的讨论，我们必须考虑在所有流体-粒子力存在的情况下介观加速度模型的最终形式。当包含虚拟质量力时，必须从单个粒子上的力平衡开始推导介观加速度模型：36

\begin{sourceblock}
\includegraphics[page=215,pagebox=mediabox,trim={106.49bp 222.76bp 109.02bp 437.85bp},clip,max width=0.94\textwidth,max height=0.38\textheight]{Marchisio2013.pdf}
\end{sourceblock}
我们隐含地假设粒子密度和速度不随时间变化。对于粒子看到的流体也有类似的表达式：

\begin{sourceblock}
\includegraphics[page=215,pagebox=mediabox,trim={104.24bp 152.07bp 106.77bp 508.09bp},clip,max width=0.94\textwidth,max height=0.38\textheight]{Marchisio2013.pdf}
\end{sourceblock}
在这些力平衡的右侧，我们包括前面部分讨论的身体力和流体-粒子力。通过将时间导数从虚拟质量力移到右侧，我们可以将力的平衡组合成矩阵形式：

\begin{sourceblock}
\includegraphics[page=215,pagebox=mediabox,trim={117.15bp 50.73bp 119.68bp 588.53bp},clip,max width=0.94\textwidth,max height=0.38\textheight]{Marchisio2013.pdf}
\end{sourceblock}
其中 I 是 3 × 3 单位矩阵，力由相间力（其中 Fpf + Ffp = 0）和相内力（Fp 和 Ff）表示。左边的矩阵可以求逆

\begin{sourceblock}
\includegraphics[page=216,pagebox=mediabox,trim={117.41bp 567.05bp 114.96bp 95.62bp},clip,max width=0.94\textwidth,max height=0.38\textheight]{Marchisio2013.pdf}
\end{sourceblock}
其中 κ = ρfVp[ρpVf + (ρpVp + ρfVf)Cvm]。因此，介观尺度加速度模型的最终表达式为

\begin{sourceblock}
\includegraphics[page=216,pagebox=mediabox,trim={118.67bp 502.07bp 65.07bp 164.08bp},clip,max width=0.94\textwidth,max height=0.38\textheight]{Marchisio2013.pdf}
\end{sourceblock}
和 dU*

\begin{sourceblock}
\includegraphics[page=216,pagebox=mediabox,trim={121.35bp 459.65bp 65.07bp 206.49bp},clip,max width=0.94\textwidth,max height=0.38\textheight]{Marchisio2013.pdf}
\end{sourceblock}
其中有效体积系数由下式给出

\begin{sourceblock}
\includegraphics[page=216,pagebox=mediabox,trim={156.27bp 409.87bp 65.07bp 256.29bp},clip,max width=0.94\textwidth,max height=0.38\textheight]{Marchisio2013.pdf}
\end{sourceblock}
从方程式。 (5.93) 和 (5.94) 中，值得注意的是，虽然流体-粒子力与没有虚拟质量力时相同，但只要 Cvm 0，粒子和流体的体力就会耦合。然而，与质量成线性比例的体力（例如重力）（即 Fp = ρpVpg 和 Ff = ρfVfg）将保持不耦合。在非常稀的极限下，我们有 Vp ≪Vf，因此方程： （5.93）简化为经典粒子模型，其中虚拟质量力仅影响粒子质量。对于在 ρf ≪ρp 极限内粒子所看到的流体也是如此。本章讨论的主要流体-粒子力得出以下流体-粒子力表达式：37

\begin{sourceblock}
\includegraphics[page=216,pagebox=mediabox,trim={95.66bp 246.54bp 65.07bp 408.92bp},clip,max width=0.94\textwidth,max height=0.38\textheight]{Marchisio2013.pdf}
\end{sourceblock}
其中右侧的项分别是浮力、阻力、虚拟质量和升力。关于比较方程式。 （5.93）和（5.94）与等式。 （5.9），我们可以通过等式观察到这一点。 (5.96) 和重力，介观尺度加速度模型由下式给出38

\begin{sourceblock}
\includegraphics[page=216,pagebox=mediabox,trim={105.72bp 101.28bp 65.07bp 504.57bp},clip,max width=0.94\textwidth,max height=0.38\textheight]{Marchisio2013.pdf}
\end{sourceblock}
\footnotetext[37]{我们用虚拟体积 Vpγ 来表示浮力，以便考虑颗粒所排出的流体体积。}
\footnotetext[38]{介观流体压力必须与 Cvm 无关，以便静水压力仅取决于混合物密度。}
以及

\begin{sourceblock}
\includegraphics[page=217,pagebox=mediabox,trim={129.79bp 567.15bp 67.56bp 76.91bp},clip,max width=0.94\textwidth,max height=0.38\textheight]{Marchisio2013.pdf}
\end{sourceblock}
Apf = Vp
有趣的是，虚拟质量力主要通过有效体积系数γ起作用，并且往往会减少动量交换项。最后一步是使用第 4 章中介绍的质量平均矩的定义来推导宏观动量平衡中的力项。对于分散相动量平衡（参见第 123 页的方程（4.85）），该过程导致

\begin{sourceblock}
\includegraphics[page=217,pagebox=mediabox,trim={81.36bp 429.10bp 67.57bp 205.08bp},clip,max width=0.94\textwidth,max height=0.38\textheight]{Marchisio2013.pdf}
\end{sourceblock}
对于流体相动量平衡（参见第 125 页的方程（4.94）），宏观加速度项为

\begin{sourceblock}
\includegraphics[page=217,pagebox=mediabox,trim={83.22bp 326.67bp 67.57bp 307.52bp},clip,max width=0.94\textwidth,max height=0.38\textheight]{Marchisio2013.pdf}
\end{sourceblock}
从等式可以看出。 (5.100) 中，虚拟质量力将阻力和升力减小了 1/γ。浮力未修改，因为我们选择根据有效体积 Vpγ 来定义它。我们提醒读者，粒子 A* pf 所看到的流体的介观加速度模型必须与粒子相 A* fp 的介观尺度模型一致，以确保整个系统在介观尺度上保持动量守恒。 （有关更多详细信息，请参阅第 4.3.8 节。）如等式附近所讨论的。 （5.14）第 144 页，这是通过在给定 Afp 模型（源自本节中介绍的力项）的情况下约束 Apf 模型来在单粒子模型中完成的。因此，如方程式中所示。 （5.98）和（5.99）中，没有必要为 Apf 中出现的动量传递项导出单独的模型。

\section*{5.4 物理空间平流}\phantomsection\addcontentsline{toc}{section}{5.4 物理空间平流}
物理空间平流代表物理空间中粒子的运动。正如方程式中已经报道的那样。 (5.3) 式中，它由 GPBE 左侧第二项表示（见式（5.1））：

\begin{sourceblock}
\includegraphics[page=217,pagebox=mediabox,trim={143.51bp 83.47bp 144.84bp 587.49bp},clip,max width=0.94\textwidth,max height=0.38\textheight]{Marchisio2013.pdf}
\end{sourceblock}
{\small 尽管粒子速度经常被作为内坐标矢量的变量之一，但它显然起着重要且不同的作用，这与\par}
被动和主动内部变量的概念。粒子速度是主动变量，因为它产生空间传输，而其他内坐标（例如粒子大小）是被动变量。被动变量简单地随特定速度一起携带，因此可以求解物理空间平流项以获得被动变量的固定值。活动变量的情况并非如此，如第 8.1.3 节中详细描述的。接下来将讨论评估物理空间平流项的不同策略。重要的是要提醒读者，所有这些策略都需要了解粒子所看到的流体速度。在目前的情况下（即单动流体极限），通常通过求解流体相动量平衡方程（例如方程式）来找到此信息。 （4.92）：

\begin{sourceblock}
\includegraphics[page=218,pagebox=mediabox,trim={102.45bp 485.96bp 98.82bp 188.40bp},clip,max width=0.94\textwidth,max height=0.38\textheight]{Marchisio2013.pdf}
\end{sourceblock}
其中 Sf 表示流体相中的粘性和压力应力张量，Af 表示流体相中的体力。右侧第二项表示固体颗粒的存在对连续相的影响。单向耦合时可以忽略，但两向或三向耦合时必须包含。如果 NDF 是粒子速度的函数，那么 GPBE 的解将为建模者提供计算物理空间平流项的基本信息。只要粒子斯托克斯数不小，就会使用这种方法，并且将导致粒子速度分布的发展。有关此主题的更多详细信息，请参阅第 8 章。另一种方法是对 NDF 与粒子速度进行积分。例如，让我们考虑一个通用的 NDF n(t, x, vp, Ψp)，它是时间 t、空间 x、粒子速度 vp 和内坐标 Ψp 的函数。通过对粒子速度进行积分，可以获得以下 NDF：

\begin{sourceblock}
\includegraphics[page=218,pagebox=mediabox,trim={181.40bp 297.52bp 65.08bp 375.97bp},clip,max width=0.94\textwidth,max height=0.38\textheight]{Marchisio2013.pdf}
\end{sourceblock}
同样，如果粒子速度是从方程中积分出来的。根据式（5.3），得到物理空间平流项的表达式如下：

\begin{sourceblock}
\includegraphics[page=218,pagebox=mediabox,trim={82.35bp 226.28bp 65.08bp 444.68bp},clip,max width=0.94\textwidth,max height=0.38\textheight]{Marchisio2013.pdf}
\end{sourceblock}
根据定义，条件速度为

\begin{sourceblock}
\includegraphics[page=218,pagebox=mediabox,trim={101.89bp 171.31bp 65.08bp 492.71bp},clip,max width=0.94\textwidth,max height=0.38\textheight]{Marchisio2013.pdf}
\end{sourceblock}
它表示内坐标等于 ψp 的粒子的预期平均速度，或者等效地，在内坐标条件下的平均粒子速度。物理空间平流项的计算需要了解条件粒子速度 〈Up|ψp〉。这可以通过不同的方法（在欧拉框架中）确定，即伪均质或尘埃气体模型、平衡或代数欧拉模型以及欧拉双流体模型（Balachandar，2009；Balachandar＆Eaton，2010；Ferry＆Balachandar，2001）。这三种方法在不同的流态中是足够的，这些流态取决于主相和分散相之间的相耦合类型，并且由主相和分散相的状态决定。流态通常由方程（1）中定义的粒子斯托克斯数来表征。第 10 页上的方程（1.20）表示特征粒子弛豫时间 (τp) 与流体特征时间 (τf) 之间的比率，以及粒子与粒子碰撞的相对重要性。这里需要指出的是，对于许多用于计算 GPBE 矩的求解方法，了解少量 ψp 值的条件粒子速度〈Up|ψp〉通常足以计算空间通量。事实上，正如在第 7 章中将变得更加清楚的那样，粒子群通常可以用具有分散相体积分数 αpk 和内坐标 ψpk 值（其中 k = 1）的 N 个粒子类（或组或节点）来表示。。。，N，其中N相对较小。

\subsection*{5.4.1 伪均质或含尘气体模型}\phantomsection\addcontentsline{toc}{subsection}{5.4.1 伪均质或含尘气体模型}
对于含尘气体模型，假设粒子速度等于流体速度：39

\begin{sourceblock}
\includegraphics[page=219,pagebox=mediabox,trim={65.02bp 472.04bp 415.73bp 214.66bp},clip,max width=0.94\textwidth,max height=0.38\textheight]{Marchisio2013.pdf}
\end{sourceblock}
p = Uf。当粒子非常快地松弛到周围流体的速度时，这种情况是现实的。众所周知，这是通过评估粒子弛豫时间来确定的，粒子弛豫时间通常是通过考虑经历粘性阻力和重力的单个粒子的拉格朗日动量平衡方程来定义的：

\begin{sourceblock}
\includegraphics[page=219,pagebox=mediabox,trim={188.66bp 399.63bp 67.57bp 270.47bp},clip,max width=0.94\textwidth,max height=0.38\textheight]{Marchisio2013.pdf}
\end{sourceblock}
\begin{sourceblock}
\includegraphics[page=219,pagebox=mediabox,trim={65.03bp 376.40bp 352.76bp 308.46bp},clip,max width=0.94\textwidth,max height=0.38\textheight]{Marchisio2013.pdf}
\end{sourceblock}
p/(18μf)。方程中的动量平衡（5.105）没有考虑由于粒子与粒子碰撞引起的瞬时变化。粒子弛豫时间的定义导致粒子斯托克斯数的定义如下：

\begin{sourceblock}
\includegraphics[page=219,pagebox=mediabox,trim={238.27bp 316.70bp 67.57bp 354.05bp},clip,max width=0.94\textwidth,max height=0.38\textheight]{Marchisio2013.pdf}
\end{sourceblock}
Stp =
对于层流，流体相的特征时间τf可以定义为特征速度Uf与特征尺寸L之间的比值。例如，在通道流限制在两个平行板内的情况下，L可以取等于板之间的距离，而Uf可以是摩擦速度。另一种常见的选择是根据粘性尺度进行计算，即用流体相的运动粘度除以摩擦速度的平方。对于湍流，通常假定 τf 为流体相中的柯尔莫哥洛夫时间尺度。仅当粒子弛豫时间趋于零（即 Stp <1）时才能应用尘埃气体模型。在这些条件下，方程。 (5.105) 得出 〈Up|ψp〉= Uf + O(τp(ψp))。正如已经提到的，只有当颗粒立即适应局部流体流动时，该模型显然才合适。当颗粒非常小和/或连续相高度粘稠和/或分散相与主相密度比非常小时，通常会发生这种情况。含尘气体模型假设只有一个粒子速度场，与流体的速度场相同。使用这种方法，显然无法预测优先积累和偏析效应，因为粒子在连续相中作为标量传输。如果系统非常稀（单向耦合），则假定连续相的特性（即密度和粘度）与流体的特性相同。如果固体颗粒浓度开始对流体相产生影响（双向耦合），则通常在方程（1）中引入连续相的修正密度和粘度。 （4.92）。

\footnotetext[39]{尘埃气体模型忽略了颗粒的布朗运动，因此不适用于描述亚微米颗粒的扩散传输。}
\subsection*{5.4.2 平衡或代数欧拉模型}\phantomsection\addcontentsline{toc}{subsection}{5.4.2 平衡或代数欧拉模型}
对于非常小但有限的粒子斯托克斯数 (0 < Stp ≪1)，粒子以与流体不同的速度移动，但同一位置的所有粒子都具有唯一的速度 40，该速度取决于内坐标的值。事实上，在这些流动条件下，两个最初不同的粒子速度之间的差异减小，并且当瞬态消失时，具有相同内坐标的所有粒子将松弛到唯一的速度。一个非常直观的例子是单分散、非相互作用的颗粒落入静止的流体中。无论它们的初始速度如何，一段时间后，所有颗粒都会松弛到颗粒最终沉降速度，这与流体速度不同。在这种情况下，最终粒子的终端速度表示平衡速度，其取决于内坐标，即粒子尺寸和形状。因此，对于非常小的粒子斯托克斯数，当粒子不受干扰地放松并很快达到平衡速度时，可以预期会出现独特的粒子速度场。发生这种情况的特征时间尺度与粒子速度弛豫的时间尺度成正比（参见方程（5.105））。因此，对于以比流体的特征时间尺度小得多的有限粒子弛豫时间为特征的颗粒系统，以及对于颗粒与颗粒碰撞可以忽略不计的颗粒系统，可以实现这种独特的颗粒速度。事实上，对于这些情况，任何最初不同的粒子速度都会以指数方式快速收敛（或衰减）到平衡速度。 41 除了粒子斯托克斯数的条件外，第二个重要条件是粒子与粒子的碰撞不是很频繁，并且不会影响粒子速度。由于粒子斯托克斯数与粒子尺寸和密度有关，并且碰撞频率与粒子浓度有关，因此通常根据这些量对不同的可能情况进行参数化，即小粒子与大粒子以及稀系统与稠密系统。使用平衡模型计算粒子速度基于对小粒子弛豫时间有效的简单展开式。事实上，我们知道当 τp →0 时，粒子速度 U*p 趋向于流体速度 Uf。如果粒子速度在方程的左侧。将（5.105）式替换为流体速度，将导数替换为全（或实质）导数（因为Uf是欧拉流体速度），并将所得方程相对于右侧出现的粒子速度进行求解，得到以下方程（Maxey，1987）：

\begin{sourceblock}
\includegraphics[page=220,pagebox=mediabox,trim={161.47bp 216.87bp 65.08bp 450.06bp},clip,max width=0.94\textwidth,max height=0.38\textheight]{Marchisio2013.pdf}
\end{sourceblock}
其中 D/Dt 是总流体导数。该代数方程的主要优点是直接将粒子速度与流体速度及其局部和时间导数联系起来。因此，在任何计算过程中，一旦评估了欧拉流体速度，就可以相对容易地获得相应的粒子速度，而无需求助于额外的微分方程的解。此外，同样容易获得具有不同弛豫时间的粒子的条件速度，或者在其他情况下
40 只有当忽略布朗运动时，这种说法才成立，这是平衡模型推导中的隐含假设。由于布朗运动对于悬浮在流体中的亚微米颗粒很重要，因此平衡模型不适用于描述纳米颗粒传输等。 41 粒子速度的收敛不应与粒子路径的收敛相混淆。事实上，即使具有平衡速度，最初以不同速度放置在相邻点的两个粒子也会很快收敛到单一速度，但处于不同的空间位置！例如，这在粒子路径总是发散的湍流中非常重要。然而，众所周知，非常简单的流场也可以产生混沌粒子路径。

具有不同大小和密度的单词，通常显示为内坐标 ψp。方程的扩展方法。 (5.107) 的推导还可以包括高阶项，从而得到以下表达式 (Druzhinin, 1987)：

\begin{sourceblock}
\includegraphics[page=221,pagebox=mediabox,trim={124.61bp 568.31bp 67.57bp 98.63bp},clip,max width=0.94\textwidth,max height=0.38\textheight]{Marchisio2013.pdf}
\end{sourceblock}
其中 af = DUf/Dt −g。方程（5.108）仅适用于在阻力和重力作用下运动的粒子。如果包括其他力（例如第 5.3 节中报告的力），则方程： （5.108）采用类似的形式。例如，当考虑重力、浮力、阻力、虚拟质量、升力和巴塞特力时，可以使用以下表达式：

\begin{sourceblock}
\includegraphics[page=221,pagebox=mediabox,trim={65.47bp 424.70bp 67.57bp 208.73bp},clip,max width=0.94\textwidth,max height=0.38\textheight]{Marchisio2013.pdf}
\end{sourceblock}
其中流体加速度现在定义为

\begin{sourceblock}
\includegraphics[page=221,pagebox=mediabox,trim={187.37bp 367.68bp 67.57bp 298.14bp},clip,max width=0.94\textwidth,max height=0.38\textheight]{Marchisio2013.pdf}
\end{sourceblock}
L 是一个线性算子，定义为

\begin{sourceblock}
\includegraphics[page=221,pagebox=mediabox,trim={129.13bp 310.94bp 67.57bp 354.05bp},clip,max width=0.94\textwidth,max height=0.38\textheight]{Marchisio2013.pdf}
\end{sourceblock}
Cvm 是虚拟质量系数，CL 是升力系数。总之，平衡或代数欧拉模型可用于非常小的 Stp，并且可以轻松处理单向、二向和三向耦合的情况。在单向耦合（稀系统）的情况下，求解连续相的动量平衡方程（参见方程（4.92）），然后求解方程（4.92）。 (5.109)（或其任何简化形式）用于计算不同粒子类别的条件粒子速度：Upk ≡〈Up|ψp = ψpk〉。对于较稠密的颗粒系统（两向或三向耦合），求解混合物的动量平衡方程（其密度和速度分别在方程（4.69）和（4.70）中定义），然后，对于单向耦合情况，方程（1）被求解。 （5.109）（或其任何简化形式）用于计算条件粒子速度。重要的是要提醒读者，虽然平衡模型能够解释双向和三向耦合，但由于颗粒与颗粒碰撞系统的重要性，即使是中等密度的颗粒系统通常也无法得到很好的描述。

\subsection*{5.4.3 欧拉二流体模型}\phantomsection\addcontentsline{toc}{subsection}{5.4.3 欧拉二流体模型}
当粒子斯托克斯数不小时，粒子速度〈Up|ψp〉的截断展开式不再有效。在这些条件下，平均粒子速度必须根据第 4.3.7 节中描述的分散相动量方程计算。让我们暂时考虑相同粒子的非常稀的群体。这些粒子的平均速度可以通过求解方程式找到。 （4.91）。对于小颗粒
斯托克斯数、粒子轨迹穿越点 (PTC) 并不重要，在这种情况下可以忽略速度波动（即 upup ≪1）。在此限制下，分散相动量平衡减少为

\begin{sourceblock}
\includegraphics[page=222,pagebox=mediabox,trim={150.61bp 574.79bp 146.97bp 97.06bp},clip,max width=0.94\textwidth,max height=0.38\textheight]{Marchisio2013.pdf}
\end{sourceblock}
如果考虑浮力、重力、阻力、虚拟质量和升力，则分散相动量平衡变为（参见第 5.3.4 节）

\begin{sourceblock}
\includegraphics[page=222,pagebox=mediabox,trim={73.14bp 479.39bp 65.08bp 154.80bp},clip,max width=0.94\textwidth,max height=0.38\textheight]{Marchisio2013.pdf}
\end{sourceblock}
同样，在单动流体极限中，流体相动量平衡由下式给出

\begin{sourceblock}
\includegraphics[page=222,pagebox=mediabox,trim={75.19bp 393.02bp 65.08bp 241.17bp},clip,max width=0.94\textwidth,max height=0.38\textheight]{Marchisio2013.pdf}
\end{sourceblock}
综上所述，欧拉二流体模型由方程（2）表示。除了流体应力张量 τf 的本构模型之外，还有（5.112）和（5.113）。正如已经提到的，等式。方程（5.112）是在粒子速度分布非常窄（即粒子斯托克斯数较小）的假设下推导出来的，并且粒子必须具有相同的内坐标。如果这些简化不成立，例如在密集条件下，当粒子与粒子的碰撞变得重要时，则必须考虑粒子速度波动，如第 4 章末尾所讨论的。

\subsection*{5.4.4 物理空间平流指南}\phantomsection\addcontentsline{toc}{subsection}{5.4.4 物理空间平流指南}
可以使用三个参数来确定计算物理空间平流项的最佳方法：粒子浓度、粒子斯托克斯数和多分散性。当由 ρp/ρf 定义的颗粒载荷远小于 1 时，系统可以被认为是单向动量耦合稀释的，而当 ρp/ρf ≥0.1 时，则需要双向动量耦合。我们认为小颗粒斯托克斯数极限由 Stp ≪1 确定，而中等和大颗粒斯托克斯数分别由 Stp ≈1 和 Stp ≫1 确定。重要的是要记住，一般来说，系统中不同物理点的粒子斯托克斯数可能会有所不同，具体取决于连续相特征时间尺度的定义。多分散性（就颗粒平流而言）可以通过多分散指数 (PDI) 来识别，多分散指数 (PDI) 定义为最大/最重颗粒与最小/最轻颗粒计算的颗粒斯托克斯数之间的比率。 42 给定多相系统的预期多分散性和多相系统经历的过程类型，可以使用以下模型之一计算相应的颗粒速度。

\footnotetext[42]{有多种计算 PDI 的方法。如果不存在密度不均匀性并且尺寸是唯一的变量，则多分散性完全由粒度分布定义。如果已知这一点，则可以根据 d10\% 和 d90\% 计算较小和最大颗粒的颗粒斯托克斯数。这些}
1. 扩散方程：亚微米粒子进行布朗运动。

2. 拟均质或含尘气体模型：非常小的颗粒斯托克斯数和有限的多分散性（如果系统是稀的，则动量平衡方程仅适用于连续相；如果系统是稠密的，则动量平衡方程仅适用于连续相和分散相的混合物）。

3.基于平均颗粒尺寸的单一条件速度的平衡或代数欧拉模型：小颗粒斯托克斯数和有限的多分散性（动量平衡方程仅适用于连续相（如果系统是稀的）或对于连续相和分散相的混合物（如果系统是稠密的）。

4. 具有条件速度的平衡或代数欧拉模型：小颗粒斯托克斯数和不可忽略的多分散性（如果系统是稀的，则动量平衡方程仅适用于连续相；如果系统是稠密的，则动量平衡方程仅适用于连续相和分散相的混合物）。

5. 具有基于平均颗粒尺寸的颗粒相速度的欧拉双流体模型：小颗粒斯托克斯数和有限的多分散性（在稀和稠密系统中）。

对于所有其他情况，有必要求解从 GPBE 导出的矩输运方程，如第 4 章所述。第 8 章详细介绍了使用基于求积的矩方法寻找 GPBE 近似解的数值算法。

\section*{5.5 扩散过程}\phantomsection\addcontentsline{toc}{section}{5.5 扩散过程}
正如第 5.1 节中所讨论的，GPBE 中出现的扩散过程是未解决的微尺度物理的介观尺度近似的结果。这些现象在物理空间和相空间中产生平流，其特征是时间和长度尺度上的随机混沌性质，比前面章节中讨论的纯平流项更小。例子包括布朗力、晶体中的随机晶格缺陷以及单个颗粒周围的不稳定流体流动。就介观尺度的单粒子模型而言，扩散过程表现为随机过程（详细信息请参见第 5.1.2 节）。方程中的一般数学公式。 （5.11）引入了大量的模型系数，这些系数可能必须通过物理推理来指定。然而，正如第 5.1.2 节所述，最基本的模型仅使用可能扩散过程的一小部分，因此，我们在这里仅考虑最重要的示例。方程中的单粒子介观模型。第 143 页上的（5.11）直接得出方程（5.11）中的扩散系数。 (5.2) 通过等式给出的关系（5.12）。作为一般规则，最好从方程形式的介观模型开始。 （5.11），为此我们已经引入了纯平流项的模型（即 Gp、Afp、
可以根据以下标准从累积粒度分布中识别两个数字：d10\%是10\%的颗粒小于d10\%的粒度，对于d90\%也类似。如果粒子群分布在粒子尺寸和密度上，我们可以类似地定义 d10\% 和 ρp,10\% 以及 d90\% 和 ρp,90\% 作为在累积分布中分离 10\% 和 90\% 粒子的尺寸和密度。因此，可以使用这些数字来计算较小和较轻的颗粒以及较大和较重的颗粒的颗粒斯托克斯数。

等）见第 5.2 和 5.3 节。这一规则的一个例外是方程 1 中出现的粒子所看到的流体质量的扩散项。 （4.80）（第 121 页），这是由于正在建模的单个粒子的邻域中存在多个粒子而导致的。因此，对于这一项，我们可以尝试向后确定等式中的系数。 （5.11）需要生成方程中的扩散系数。 （4.80）。同样，使用 5.1.2 节中介绍的介观尺度质量和动量守恒的概念，我们可以指定方程 1 中系数的形式。 （5.11）对于粒子（即 A* pf 和 G* f）所看到的流体，使用为 A* fp 和 G* p 开发的随机模型。

\subsection*{5.5.1 相空间扩散}\phantomsection\addcontentsline{toc}{subsection}{5.5.1 相空间扩散}
在本节中，我们考虑使用相空间扩散建模的两个示例。第一个示例涉及生长粒子的粒度分布，第二个示例涉及粒子速度波动。

\paragraph*{增长率离散度}
相空间扩散的一个例子是结晶过程中的生长速率分散。许多实验表明，当晶体在过饱和作用下生长时，所得的 NDF 如预期的那样移动到更高的粒径，但同时在粒径空间中扩散。这一观察结果表明，颗粒尺寸的变化率 G*p 必须包含纯平流分量 Gp（参见第 5.2.2 节）和自扩散分量 cpψp，其中 ψp 是颗粒质量。该机制最早由 Larson 等人提出。 (1985) 和 Zumstein \& Rousseau (1987)，尽管在一定时期内，增长率离散度与规模依赖性增长相混淆，因为后者对 NDF 的演化具有类似的影响（例如参见 Girolami \& Rousseau (1985) 和 Mydlarz \& Briedis (1992) 的工作）。在大多数已发表的著作中，增长率离散系数是通过拟合实验数据来计算的。事实上，关于该量如何受操作条件（最明显的是过饱和、滑移速度、固体浓度和温度）影响的信息严重缺乏。据认为，生长速率分散是由晶格缺陷的随机形成引起的，这使得每个晶体以略有不同的生长速率生长。就粒子 ψf 所看到的流体质量而言，介观尺度上的质量守恒导致在零粒子斯托克斯数（即 U*

\begin{sourceblock}
\includegraphics[page=224,pagebox=mediabox,trim={74.71bp 173.80bp 65.08bp 473.86bp},clip,max width=0.94\textwidth,max height=0.38\textheight]{Marchisio2013.pdf}
\end{sourceblock}
其中 B 是方程式中扩散系数定义中出现的矩阵。 （5.12）。得到的 GPBE 由下式给出43

\begin{sourceblock}
\includegraphics[page=224,pagebox=mediabox,trim={104.62bp 110.87bp 100.97bp 557.37bp},clip,max width=0.94\textwidth,max height=0.38\textheight]{Marchisio2013.pdf}
\end{sourceblock}
其中 n(t, Ψp, Ψf) 是 NDF。由于 B 为一阶，因此该 GPBE 中的扩散系数仅允许在相空间的一个方向上进行色散。计算式（1）中 GPBE 的低阶矩是有指导意义的。 (5.114) 使用所描述的方法

\footnotetext[43]{为了清楚起见，在本节中，我们在编写 GPBE 时将省略空间平流项。}
在第 4.3 节中。为了清楚起见，我们假设 Gp 和 cpψp 是常数。对应于粒子数密度 Np 的零阶矩是恒定的。与 Ψp 和 Ψf 相关的一阶矩产生

\begin{sourceblock}
\includegraphics[page=225,pagebox=mediabox,trim={174.13bp 567.52bp 220.27bp 105.95bp},clip,max width=0.94\textwidth,max height=0.38\textheight]{Marchisio2013.pdf}
\end{sourceblock}
dt = −GpNp,
反映了系统中的质量守恒。与 Ψ2 p 和 Ψ2 f 相关的二阶矩产生

\begin{sourceblock}
\includegraphics[page=225,pagebox=mediabox,trim={65.02bp 478.57bp 128.61bp 172.83bp},clip,max width=0.94\textwidth,max height=0.38\textheight]{Marchisio2013.pdf}
\end{sourceblock}
\begin{sourceblock}
\includegraphics[page=225,pagebox=mediabox,trim={181.85bp 441.27bp 183.19bp 232.20bp},clip,max width=0.94\textwidth,max height=0.38\textheight]{Marchisio2013.pdf}
\end{sourceblock}
通过将质量方差和协方差定义为

\begin{sourceblock}
\includegraphics[page=225,pagebox=mediabox,trim={79.94bp 388.17bp 84.17bp 281.46bp},clip,max width=0.94\textwidth,max height=0.38\textheight]{Marchisio2013.pdf}
\end{sourceblock}
很容易证明

\begin{sourceblock}
\includegraphics[page=225,pagebox=mediabox,trim={132.25bp 330.98bp 133.59bp 340.29bp},clip,max width=0.94\textwidth,max height=0.38\textheight]{Marchisio2013.pdf}
\end{sourceblock}
因此，对于较大的 t，颗粒和流体质量之间的相关系数变为 -1（或者如果初始方差为零，则对于所有 t 都等于 -1）。该结果直接作为介观尺度质量守恒的结果得出（即方程（5.114）），并且对于 Gp 和 cpψp 的所有选择都是正确的。

\paragraph*{流体和分散相速度波动}
分散相速度的波动可能由许多不同的物理过程引起，但在这里我们重点关注由于与流体相的相互作用而产生的波动，例如，由于不稳定的边界层和粒子周围的尾流（Ten-neti et al., 2012）。当表 1.1 中定义的粒子雷诺数大于 1 时，就会出现这种波动。最简单的单粒子介观模型可以写为 A* fp = Afp + bpvfWpvf(t)，其中纯平流项 Afp 如 5.3 节所述。令 Ψp 和 Ψf 分别为粒子和粒子所见流体的质量，介观尺度动量守恒导致 A* pf = -ΨpAfp/Ψf -ΨpbpvfWpvf(t)/Ψf，因此在方程 1 中 Apf = -ΨpAfp/Ψf 和 bfvf = -Ψpbpvf/Ψf。 （5.11）。在这里，为了清楚起见，我们假设 G*p 和 G*f 没有由于随机过程而产生的波动分量。此外，我们暂时忽略式（1）中的“扩散”贡献。 （4.80）（参见下面第 5.5.3 节），因此从方程中找到扩散矩阵。 (5.12) 将仅取决于 bpvf 和 bfvf：

\begin{sourceblock}
\includegraphics[page=225,pagebox=mediabox,trim={84.47bp 50.16bp 87.00bp 615.58bp},clip,max width=0.94\textwidth,max height=0.38\textheight]{Marchisio2013.pdf}
\end{sourceblock}
其中 3 × 3 扩散矩阵由 Bfp = bpvfbT pvf 定义。当用 n(t, vp, vf) 表示（同质）NDF 时，得到的 GPBE 为

\begin{sourceblock}
\includegraphics[page=226,pagebox=mediabox,trim={102.46bp 545.11bp 98.82bp 94.25bp},clip,max width=0.94\textwidth,max height=0.38\textheight]{Marchisio2013.pdf}
\end{sourceblock}
正定矩阵 Bfp 将由颗粒周围流体流动的结构决定，并且很可能具有显著的各向异性（Tenneti 等人，2012）。请注意，我们没有在方程中包含流体速度波动耗散模型。 (4.104) 编写 GPBE 时。在这里，为了清楚起见，我们将专门关注流体和粒子速度的相互依赖性，为此考虑 vp 和 vf 的一维速度相空间就足够了。这种情况下的 GPBE 由下式给出

\begin{sourceblock}
\includegraphics[page=226,pagebox=mediabox,trim={101.10bp 424.24bp 65.08bp 243.79bp},clip,max width=0.94\textwidth,max height=0.38\textheight]{Marchisio2013.pdf}
\end{sourceblock}
其中我们假设 Afp = (1/τp)(vf −vp) 并且 Bfp 在相空间中是均匀的。方程的一阶矩。 (5.115) 分别关于 Ψp 和 Ψf，证实相密度 ρp 和 ρf 是常数。分别关于 Ψpvp 和 Ψfvf 产量的二阶矩，

\begin{sourceblock}
\includegraphics[page=226,pagebox=mediabox,trim={148.76bp 333.29bp 136.39bp 339.39bp},clip,max width=0.94\textwidth,max height=0.38\textheight]{Marchisio2013.pdf}
\end{sourceblock}
并表达宏观尺度的动量守恒。方程的三阶矩。 (5.115) 分别关于 Ψpv2 p、Ψfv2 f 和 Ψpvpvf 产量（详见第 4.3.9 节），

\begin{sourceblock}
\includegraphics[page=226,pagebox=mediabox,trim={165.50bp 217.01bp 161.85bp 420.42bp},clip,max width=0.94\textwidth,max height=0.38\textheight]{Marchisio2013.pdf}
\end{sourceblock}
以及

\begin{sourceblock}
\includegraphics[page=226,pagebox=mediabox,trim={138.79bp 172.89bp 121.59bp 493.27bp},clip,max width=0.94\textwidth,max height=0.38\textheight]{Marchisio2013.pdf}
\end{sourceblock}
我们除以恒定相密度。很容易证明方差和协方差的方程（例如 U2

\begin{sourceblock}
\includegraphics[page=226,pagebox=mediabox,trim={345.92bp 142.02bp 68.88bp 542.84bp},clip,max width=0.94\textwidth,max height=0.38\textheight]{Marchisio2013.pdf}
\end{sourceblock}
p) 具有完全相同的形式，因此 Bfp 可以被视为方差源项。然而，也可以证明波动使以下数量保持不变：44

\begin{sourceblock}
\includegraphics[page=226,pagebox=mediabox,trim={180.40bp 84.03bp 177.95bp 588.65bp},clip,max width=0.94\textwidth,max height=0.38\textheight]{Marchisio2013.pdf}
\end{sourceblock}
\footnotetext[44]{当方程中的流体速度脉动耗散模型时。 (4.104) 被包括在内，将存在一个额外的方差减少机制。}
因此，如果初始方差为零，则 UpUf * 完全由 U2 p * 和 U2 f * 决定。在稳定状态下，速度波动与 Bfp 的关系为

\begin{sourceblock}
\includegraphics[page=227,pagebox=mediabox,trim={90.46bp 560.35bp 92.98bp 91.55bp},clip,max width=0.94\textwidth,max height=0.38\textheight]{Marchisio2013.pdf}
\end{sourceblock}
f *)1/2（即流体-粒子速度相关系数为-1，这意味着由于介观尺度动量守恒，波动是反相关的）。剩下的介观尺度建模任务是确定无量纲介观尺度参数 C(Rep, Stp) = τpBfp/|Uf −Up|2（更一般地说，Bfp）如何依赖于粒子雷诺数和斯托克斯数，例如使用微尺度方程的 DNS（Tenneti 等人，2012）。总之，上面讨论的两个例子说明了未解决的微尺度现象如何在介观尺度上产生相空间扩散项。一般来说，开发单个粒子的介观模型的起点将是使用无量纲参数（例如粒子雷诺数和斯托克斯数）的 DNS 派生相关性。如上面两个例子所示，粒子所看到的流体的介观模型通常可以通过调用介观尺度的质量、化学物质、动量等守恒来推导。

\subsection*{5.5.2 物理空间扩散}\phantomsection\addcontentsline{toc}{subsection}{5.5.2 物理空间扩散}
物理空间扩散最常见的例子是布朗运动，这已经在 5.3.3 节中讨论过。让我们考虑悬浮在流体中的一群颗粒。众所周知，流体分子相互作用，导致速度分布在平衡条件下为经典麦克斯韦分布。流体分子还与悬浮颗粒相互作用（并交换动量）。当粒子惯性很大时（通过有限粒子斯托克斯数测量），流体-粒子力（即 A* fp 和 A* pf）是这些相互作用的相关介观模型。然而，当粒子惯性非常小时，粒子对与周围流体分子的微小的、几乎瞬时的动量交换也很敏感。由此产生的力产生随机粒子速度波动 (Gardiner, 2004)，当 1 ≤ Stp 时，可以将其建模为粒子速度相空间中的随机过程，或者当 Stp = 0 时，将其建模为物理空间中的随机游走。对于后者，布朗运动用正定各向同性扩散张量来描述： Г = ГI =⇒∇· (Г ∇n)，其中扩散系数使用以下公式计算斯托克斯-爱因斯坦方程：

\begin{sourceblock}
\includegraphics[page=227,pagebox=mediabox,trim={219.37bp 198.95bp 67.57bp 474.79bp},clip,max width=0.94\textwidth,max height=0.38\textheight]{Marchisio2013.pdf}
\end{sourceblock}
其中 kB 是玻尔兹曼常数，Tf 是流体温度，μf 是流体粘度，dp 是颗粒扩散直径。值得注意的是，在极限 Stp = 0 的情况下，GPBE 可以简化为 NDF n(t, x, ψp) 的 PBE，其中平流速度为 Uf，扩散系数如式（1）所示。 (5.116)。流体速度 Uf 和流体特性 φf（包括 Tf）是通过求解单独的传输方程得出的。事实上，作为一般规则，当使用 GPBE 求 NDF n(t, x, vp, Ψp, vf, Ψf) 时，布朗力项将作为速度相空间扩散项出现在 A* fp 和 A* pf 的模型中。 （有关布朗运动 Langevin 模型的详细信息，请参阅 Gardiner (2004) 第 1 章。）

\footnotetext[45]{该结果是在由颗粒存在引起的流体相速度波动的极限条件下得出的。在由多个源产生流体相速度波动的系统中，流体-粒子速度相关系数将趋向于零，并且取决于源项的相对大小。}
\subsection*{5.5.3 混合相和物理空间扩散}\phantomsection\addcontentsline{toc}{subsection}{5.5.3 混合相和物理空间扩散}
GPBE 方程中未封闭项的一些模型（Reeks，2005a，2005b）使用涉及 x 和一些相空间变量的混合二阶导数。作为一个例子，让我们考虑以下 n(t, x, vp, vf, ϋf) 的 GPBE：

\begin{sourceblock}
\includegraphics[page=228,pagebox=mediabox,trim={83.77bp 516.23bp 65.08bp 116.60bp},clip,max width=0.94\textwidth,max height=0.38\textheight]{Marchisio2013.pdf}
\end{sourceblock}
这是方程的空间非均匀版本。 (5.115)。这里， ψf 是粒子所看到的流体质量，其变化率由方程 1 给出。 （4.62）。后者在等式左侧引入了 x 和 Σf 之间的混合导数。 （5.117）（见方程（5.1））。现在假设我们想简化方程式中的 GPBE。 （5.117）通过对流体变量 vf 和 ϋf 进行积分（从而消除混合导数）。将新的 NDF 定义为

\begin{sourceblock}
\includegraphics[page=228,pagebox=mediabox,trim={196.72bp 423.40bp 194.27bp 249.39bp},clip,max width=0.94\textwidth,max height=0.38\textheight]{Marchisio2013.pdf}
\end{sourceblock}
方程的正式积分(5.117) 导致

\begin{sourceblock}
\includegraphics[page=228,pagebox=mediabox,trim={139.47bp 368.27bp 65.08bp 296.74bp},clip,max width=0.94\textwidth,max height=0.38\textheight]{Marchisio2013.pdf}
\end{sourceblock}
其中未知条件流体速度定义为

\begin{sourceblock}
\includegraphics[page=228,pagebox=mediabox,trim={194.43bp 323.94bp 241.55bp 349.55bp},clip,max width=0.94\textwidth,max height=0.38\textheight]{Marchisio2013.pdf}
\end{sourceblock}
vfn dvf dψf。

请注意，尽管我们已经消除了方程式中的混合导数。 (5.118)，新的GPBE不再关闭。为了关闭方程。 (5.118)，我们遵循 Reeks (2005b) 和 let46 的例子

\begin{sourceblock}
\includegraphics[page=228,pagebox=mediabox,trim={195.29bp 254.71bp 192.84bp 418.42bp},clip,max width=0.94\textwidth,max height=0.38\textheight]{Marchisio2013.pdf}
\end{sourceblock}
其中 μ 是一个非负系数，取决于粒子位置与粒子看到的流体速度之间的相关性。该模型导致封闭的 GPBE：

\begin{sourceblock}
\includegraphics[page=228,pagebox=mediabox,trim={104.79bp 194.30bp 65.07bp 470.71bp},clip,max width=0.94\textwidth,max height=0.38\textheight]{Marchisio2013.pdf}
\end{sourceblock}
然而，闭包在 x 和 vp 之间引入了新的混合导数。从数学角度来看，我们可以问混合导数项是否可以通过介观尺度随机过程建模，例如方程（1）中给出的过程。 （5.11）。通过隔离方程中混合导数中涉及变量的项。 (5.117) （即 x 和 ϋf）并且在方程式中（5.119）（即 x 和 vp），可以证明这两个方程都是双曲线的（不是扩散模型所需的抛物线）。因此，不幸的是，似乎没有介观模型需要描述 xif（即，对于非稀释系统，其中流体相体积分数不均匀）和 vp（即，对于具有未解析流体的系统）

\footnotetext[46]{除了空间梯度之外，标准闭包还包括一个关于 vp 的导数项。类似的项已作为方程式中的最后一项包含在内。 （5.118）。事实上，由于 Bfp 模拟了流体速度波动引起的强迫，因此可以认为 μ 应与 Bfp 成正比。}
未包含在 NDF 和 GPBE 中的速度波动可以很容易地用拉格朗日单粒子模型来表示。 47 然而幸运的是，这种影响可以从方程 1 中的 GPBE 开始，在欧拉矩方法中以直接的方式解释。 （5.1）。

\section*{5.6 零阶点过程}\phantomsection\addcontentsline{toc}{section}{5.6 零阶点过程}
点过程通过不连续事件影响粒子的状态。连续事件和不连续事件之间的主要区别在于，对于前者来说，在无穷小的时间间隔内，内坐标会发生无穷小的变化，而对于后者来说，变化是有限的。一个经典的例子是阻力和碰撞事件对粒子速度的影响之间的差异。拖曳力连续作用，导致粒子速度平滑变化，而碰撞会瞬间将粒子速度从碰撞前值更改为碰撞后值。48 请注意，在方程式 1 给出的 GPBE 中， (5.1) 点过程都集中到 S1 中，但这里我们将它们分开，为每种类型引入新的符号。

\subsection*{5.6.1 分散相的形成}\phantomsection\addcontentsline{toc}{subsection}{5.6.1 分散相的形成}
在许多多相系统中，分散相并不是一开始就存在的，而是在多相系统的演化过程中产生的。然而，也可能出现截然不同的情况。在由单一化学成分构成的系统的情况下，分散相可以自发形成，因为局部热力学条件允许状态转变。一个这样的例子是沸水，例如，水在热表面上流动，并且由于局部温度值，可以形成一些水蒸气气泡。在这种情况下，驱动力是当地水温与液-汽转变平衡温度之间的距离（通常表示为比率）。另一个重要的例子是熔体结晶，其中固体结晶颗粒在液相内形成。在此类系统中，新的分散相元素的形成通常通过成核过程发生。重要的实际例子是结晶和沉淀过程，其中，例如，由于连续相相对于一种（或多种）组分的过饱和，固体颗粒（晶体）成核。另一个有趣的例子是火焰中粒子的形成。在这种情况下，气相中发生的一些化学反应导致过饱和的积累，并随后形成固体产物（例如二氧化钛颗粒、烟灰等）。在专业文献中，后一个过程经常被称为起始。流体-流体系统还存在许多其他例子，特别是不混溶液体的分层和沸腾液体系统中的蒸气形成。

47 在欧拉-拉格朗日模拟中，拉格朗日粒子根据 n*(t, x, vp) 的封闭 GPBE 进行模拟，其中 〈Uf|vp〉= Uf （即它们忽略了 μ 的贡献）。由于 μ 与 1/Stp 成正比（Reeks，2005a，2005b），因此此类代码不适合近似具有小甚至中等颗粒斯托克斯数的分散相流。 48 这在形式上仅适用于涉及硬球势的粒子之间的相互作用。事实上，在这种情况下，相互作用的粒子仅在接触点处感知到彼此的存在。当相互作用由平滑势控制时，理论上碰撞不应被视为点过程。然而，即使在这些情况下，碰撞也常常被描述为在相互作用的粒子之间最短距离处发生的不连续过程，如麦克斯韦分子的情况，正如麦克斯韦 (1867) 的原始作品中所解释的那样。欲了解更多详细信息，请读者参阅第 6 章。

在由两种或多种化学成分构成的系统的情况下，局部条件可以允许自发形成由部分化学成分构成的新的（或多个）热力学相。在这种情况下，驱动力由参与新相形成的化学成分的浓度与相应的平衡浓度的比率来确定。该比率通常称为过饱和度。对于多种组分，新相的形成可以通过成核或旋节线分解过程发生。事实上，根据当地的热力学条件，不同的情况都是可能的。新热力学相的形成通常与小分子波动有关，这些波动可以将第二相的许多分子聚集在一起。尽管新相是热力学稳定的，但可能需要克服能量势垒，这通常由自由能相对于组成参数的二阶导数的符号表示。当存在这种能量势垒时，颗粒通过成核形成，这意味着只有大于临界尺寸（称为核）的颗粒才是局部稳定的并且可以进一步生长，形成新的分散相元素。因此，只有当随机波动将临界数量的分子聚集在一起（形成核）时，才会发生成核事件。在这种情况下，颗粒形成局限于特定位置，其速率由能垒和新形成相分子的热能控制。当没有需要克服的能量势垒时，由于新相是局部和全局稳定的，因此即使分散相的两个分子在一起也比分开时更稳定。在这种情况下，分散相元素是通过旋节线分解形成的，其速率由扩散控制。旋节线分解通常不如成核局部化，并且通常导致完全相分离而不是形成稳定的分散体系。下面，我们将从粒子相 NDF 的角度讨论零级点过程的两个重要示例。然而，一些零级点过程，例如从流体相形成分散相，伴随着流体相的状态变化（即两相的总质量和动量守恒）。因此，从颗粒和流体相的角度来看，整个过程需要流体相质量的减少等于所形成颗粒的质量，这由式（1）中的SM项表示。 （4.68）第119页。换句话说，方程中的术语SM。 （4.77）表示质量从流体相添加到颗粒相的速率，必须遵循形成分散相的零级点过程的源项。利用 delta 函数的性质，我们可以将源项形式化地写成等式： (5.1) 对于零阶点过程为

\begin{sourceblock}
\includegraphics[page=230,pagebox=mediabox,trim={141.18bp 194.68bp 65.08bp 492.03bp},clip,max width=0.94\textwidth,max height=0.38\textheight]{Marchisio2013.pdf}
\end{sourceblock}
其中 J(φf) ≥0 是具有性质 ψpc 的颗粒的形成速率，这取决于流体相 φf 的平均组成。这里我们假设新粒子的速度和这些粒子看到的流体的速度都等于平均流体速度 Uf。必须指定新粒子 ψpc 和这些粒子 ψfc 所看到的流体的属性，并且后者必须与方程（1）中的质量平衡一致。 （4.77）。例如，如果 Σp1 和 Σf1 是质量，则 Σpc1 = mc 是新粒子的质量，并且 Σfc1 = −mc 与整体质量平衡保持一致。49 事实上，使用第 4 章中的符号，相应的数字变化，

\footnotetext[49]{在许多情况下，mc 非常小，因此对质量和动量平衡的影响可以忽略不计。然而，为了保持一致性，最好在从 GPBE 计算矩输运方程时将其包括在内。}
\begin{sourceblock}
\includegraphics[page=231,pagebox=mediabox,trim={181.15bp 621.70bp 183.69bp 64.63bp},clip,max width=0.94\textwidth,max height=0.38\textheight]{Marchisio2013.pdf}
\end{sourceblock}
mass,

\begin{sourceblock}
\includegraphics[page=231,pagebox=mediabox,trim={166.92bp 554.45bp 167.80bp 92.35bp},clip,max width=0.94\textwidth,max height=0.38\textheight]{Marchisio2013.pdf}
\end{sourceblock}
和动力，

\begin{sourceblock}
\includegraphics[page=231,pagebox=mediabox,trim={150.74bp 478.95bp 151.62bp 167.85bp},clip,max width=0.94\textwidth,max height=0.38\textheight]{Marchisio2013.pdf}
\end{sourceblock}
验证各相之间质量和动量的总体守恒。在接下来的两节中，我们将举例说明如何找到两个不同系统的形成速率 J。

\subsection*{5.6.2 溶液中晶体的成核}\phantomsection\addcontentsline{toc}{subsection}{5.6.2 溶液中晶体的成核}
就 GPBE 的解而言，定义分散相新元素形成速率 J(φf) 的函数形式就足够了。例如，溶液中颗粒的形成通常分为均质或非均匀。对于前者，颗粒在本体溶液中形成，没有任何表面的支持，而对于后者，涉及外表面（例如，由已经形成的颗粒或由外来颗粒提供）。根据经典理论（Myerson，2002），成核率规定了尺寸为 dc 的单分散核的爆发。所得的成核速率为

\begin{sourceblock}
\includegraphics[page=231,pagebox=mediabox,trim={184.52bp 257.40bp 67.57bp 404.69bp},clip,max width=0.94\textwidth,max height=0.38\textheight]{Marchisio2013.pdf}
\end{sourceblock}
其中D是溶质的分子扩散系数（明确的表达式可以在Bird等人，1960中找到），d是其分子直径，kB是玻尔兹曼常数，Tf是绝对流体温度，σ是颗粒与溶液（在均相成核的情况下）或异物表面（在异相成核的情况下）之间的界面张力，S是流体中的过饱和度，̃v是溶质分子体积。核的尺寸通过成核理论的应用来估计：

\begin{sourceblock}
\includegraphics[page=231,pagebox=mediabox,trim={167.70bp 145.38bp 67.57bp 528.87bp},clip,max width=0.94\textwidth,max height=0.38\textheight]{Marchisio2013.pdf}
\end{sourceblock}
有关方程的理论和相关应用的更多详细信息。 (5.121) 可以在 Myerson (2002) 编辑的手册中找到。

\subsection*{5.6.3 沸腾液体中蒸气泡的成核}\phantomsection\addcontentsline{toc}{subsection}{5.6.3 沸腾液体中蒸气泡的成核}
{\small 通常认为，当液体温度高于平衡温度（即饱和温度）时，蒸汽（在大量液体内）开始形成。\par}
温度，Tsat）。在此假设下，蒸发速率（Tf > Tsat）可写为 J = 6 ̇mcp

\begin{sourceblock}
\includegraphics[page=232,pagebox=mediabox,trim={242.48bp 590.53bp 65.08bp 88.45bp},clip,max width=0.94\textwidth,max height=0.38\textheight]{Marchisio2013.pdf}
\end{sourceblock}
其中 ̇mcp（以 kg m−3 s−1 表示）表示单位时间内单位体积液汽混合物蒸发的液体质量，dc 是有核气泡的大小（即 mc = ρpkVd3 c）。根据蒸发速率，不同的状态是可能的。对于核沸腾，可以使用以下方程：

\begin{sourceblock}
\includegraphics[page=232,pagebox=mediabox,trim={205.69bp 511.70bp 65.08bp 168.25bp},clip,max width=0.94\textwidth,max height=0.38\textheight]{Marchisio2013.pdf}
\end{sourceblock}
其中 β 是蒸发频率，ρf 是液相的相密度。蒸发频率 β (s−1) 可以从理论考虑和强简化假设中得出：

\begin{sourceblock}
\includegraphics[page=232,pagebox=mediabox,trim={185.89bp 427.96bp 65.08bp 237.85bp},clip,max width=0.94\textwidth,max height=0.38\textheight]{Marchisio2013.pdf}
\end{sourceblock}
其中Mw是水分子的分子量，Lv是蒸发潜热。该方程是在平坦界面和形成具有恒定直径的分散气泡的假设下推导的。然而，通过将速率乘以范围从零到无穷大的校正因子，可以在非常不同的操作条件下使用相同的方程。 Lee (1979) 和 Tanasawa (1991) 的工作报告了有关气泡成核的更多细节。

\section*{5.7 一阶点过程}\phantomsection\addcontentsline{toc}{section}{5.7 一阶点过程}
一阶点过程将涉及单个粒子的所有不连续过程组合在一起。由于只有一个粒子参与不连续事件，因此可以假设该速率与 NDF 成正比。一个常见的例子是颗粒破碎：当颗粒（固体、液滴或气泡）受到某种压力时，它会破碎成碎片。在这种情况下，必须包括正源项和负源项，因为一个粒子将消失，而其他粒子（碎片）将同时进入相空间。另一个例子是墙壁和边界上的粒子沉积，其中仅存在负源项（Bensaid 等，2009，2010）。在这里，我们将 GPBE 中的源项写为

\begin{sourceblock}
\includegraphics[page=232,pagebox=mediabox,trim={218.99bp 185.36bp 216.54bp 497.23bp},clip,max width=0.94\textwidth,max height=0.38\textheight]{Marchisio2013.pdf}
\end{sourceblock}
其中 h+ 是生产率，h− 是由于一阶点过程导致的粒子损失率。一阶过程通常用 b(vp, Ψp, vf, Ψf)dt 来量化，它是速度为 vp 和内坐标 Ψp 的粒子在速度为 vf 和内坐标 Ψf 的流体中在无穷小的时间间隔 dt 内经历所研究过程的概率。量b是过程的频率并且具有反比时间的单位，并且其对颗粒和流体速度以及流体性质的依赖性用于例如剪切引起的破裂。根据这个定义，由于一阶过程而消失的每单位体积和单位时间的粒子数量可以很容易地写为（为了清楚起见，抑制 t 和 x 依赖性）

\begin{sourceblock}
\includegraphics[page=232,pagebox=mediabox,trim={150.42bp 52.15bp 65.08bp 633.16bp},clip,max width=0.94\textwidth,max height=0.38\textheight]{Marchisio2013.pdf}
\end{sourceblock}
为了量化由一阶过程引起的粒子形成速率，必须定义附加函数。这通常以条件概率密度函数 (PDF) 的形式引入，该函数表示从具有 vp 和 Ψp 的母粒子形成具有 v′ p 和 Ψ′ p 的子粒子的概率。这个PDF必须满足一些约束，例如归一化条件：50

\begin{sourceblock}
\includegraphics[page=233,pagebox=mediabox,trim={185.26bp 550.48bp 67.57bp 121.32bp},clip,max width=0.94\textwidth,max height=0.38\textheight]{Marchisio2013.pdf}
\end{sourceblock}
它表明产生子粒子的概率是统一的。 （请注意，如果子粒子始终具有与母粒子相同的速度，则 P 或以下公式中都不需要对速度的依赖。）在许多情况下，而不是等式中的 PDF， （5.127），根据相应的条件 NDF N(v′

\begin{sourceblock}
\includegraphics[page=233,pagebox=mediabox,trim={211.37bp 477.18bp 264.25bp 207.68bp},clip,max width=0.94\textwidth,max height=0.38\textheight]{Marchisio2013.pdf}
\end{sourceblock}
使用 p|vp, ψp)。相应的数学对象使得相空间积分现在产生

\begin{sourceblock}
\includegraphics[page=233,pagebox=mediabox,trim={170.63bp 437.82bp 67.57bp 233.99bp},clip,max width=0.94\textwidth,max height=0.38\textheight]{Marchisio2013.pdf}
\end{sourceblock}
其中ν是一阶过程形成的新粒子的数量。显然，两个函数之间的关系是一个简单的重整化。函数 N 通常称为子分布函数。由于一阶不连续过程而形成新粒子的速率现在可以表示为

\begin{sourceblock}
\includegraphics[page=233,pagebox=mediabox,trim={83.83bp 349.07bp 67.57bp 322.74bp},clip,max width=0.94\textwidth,max height=0.38\textheight]{Marchisio2013.pdf}
\end{sourceblock}
或者，就 P 而言，为

\begin{sourceblock}
\includegraphics[page=233,pagebox=mediabox,trim={93.05bp 296.95bp 67.57bp 374.87bp},clip,max width=0.94\textwidth,max height=0.38\textheight]{Marchisio2013.pdf}
\end{sourceblock}
请注意，这些表达式是在假设流体速度和成分在一阶过程中不发生变化的情况下编写的。因此，今后当不需要时，我们将在符号中抑制 vf 和 ψf。在以下部分中，我们将回顾 b 和 N 的一些最常见的函数形式。

\subsection*{5.7.1 颗粒过滤和沉积}\phantomsection\addcontentsline{toc}{subsection}{5.7.1 颗粒过滤和沉积}
当颗粒在固体壁附近流动时，由于与壁的碰撞，可能会发生沉积。颗粒沉积是一种表面过程，因为它发生在壁的表面上，并且取决于表面的范围等。然而，在许多实际应用中，例如颗粒过滤，颗粒被迫通过复杂的孔隙和通道结构（例如，在陶瓷或纤维过滤器中），必须将其描述为多孔介质。对于这种情况，沉积过程被建模为体积一阶过程。让我们考虑一个由填充物体组成的系统，其中充满颗粒的流体流过物体之间的空隙。这是催化过程、地下水流和过滤过程中的常见情况。现在让我们将这组物体描述为多孔介质，其特征在于其孔隙率 ε 和平均物体尺寸，

\footnotetext[50]{如果子粒子取决于流体特性，则对 vf 和 εf 的依赖可以包含在 P 和 N 的定义中。}
总干事。孔隙率很容易定义为未被物体占据且可自由填充流体的空间量，而平均物体尺寸通常被解释为等效球体直径 DG = 6/s，这又由多孔介质的比表面积 s 确定。如果系统被描述为表面粒子速度为 51 vp 的连续体，则粒子沉积速率的系数，或者换句话说，一阶过程的系数为（Elimelech 等人，1998）

\begin{sourceblock}
\includegraphics[page=234,pagebox=mediabox,trim={214.49bp 530.33bp 65.08bp 142.36bp},clip,max width=0.94\textwidth,max height=0.38\textheight]{Marchisio2013.pdf}
\end{sourceblock}
这一表达式是通过考虑一个尺寸为 DG 的孤立球形物体而得出的，在颗粒和物体（通常称为收集器）之间不存在排斥力的情况下，颗粒在其周围以非常低的速度 vp 流动。 52 在颗粒沉积的情况下，子 NDF N 完全相同，因为当颗粒撞击收集器时，它们会从多相系统中消失，而不会生成新颗粒。相反，如果粒子在附着到收集器后分离（或释放）回到多相系统中，则会生成粒子。由于在实际系统中，方程的条件。 (5.131) 的推导不被尊重，引入传输效率 η 是为了考虑与收集器周围颗粒流态和其他因素相关的非理想情况。总传输效率表示为 η = ηD + ηI - ηDηI，其中 ηD 是与布朗运动相关的效率， ηI 与拦截相关。由于颗粒受到随机布朗力的作用，因此它们可以靠近收集器并粘附在其上。该过程的效率可以量化为

\begin{sourceblock}
\includegraphics[page=234,pagebox=mediabox,trim={162.35bp 328.90bp 65.08bp 338.40bp},clip,max width=0.94\textwidth,max height=0.38\textheight]{Marchisio2013.pdf}
\end{sourceblock}
这里使用 P ́eclet 数的不同定义，Γ 是方程 1 中定义的斯托克斯-爱因斯坦扩散系数。 (5.116)，桑原函数为

\begin{sourceblock}
\includegraphics[page=234,pagebox=mediabox,trim={164.72bp 252.79bp 65.08bp 408.59bp},clip,max width=0.94\textwidth,max height=0.38\textheight]{Marchisio2013.pdf}
\end{sourceblock}
方程（5.132）表明，当颗粒速度增加（或颗粒尺寸增加）时，导出沉积速率系数的理论条件不再有效，实际沉积速率小于理论沉积速率（即 ηD <1）。此外，由于颗粒的尺寸 dp 有限，因此必须校正沉积速率系数（假设点颗粒得出）。这种修正是通过拦截效率引入的：

\begin{sourceblock}
\includegraphics[page=234,pagebox=mediabox,trim={208.22bp 143.53bp 65.08bp 527.83bp},clip,max width=0.94\textwidth,max height=0.38\textheight]{Marchisio2013.pdf}
\end{sourceblock}
其中 MG = dp/DG 且 s = (3 −2ε)/ε。这些效率（包括考虑粒子-收集器排斥力的效率）的更准确表达式已经
51 在没有形成多孔介质的物体的情况下，粒子占据整个空间时的速度。

\footnotetext[52]{这一假设意味着，当粒子撞击物体表面时，它们会立即附着在物体上。如果出现排斥力，则必须引入效率，考虑到颗粒与收集器之间并非所有碰撞都成功的事实。}
公式化（Ma et al., 2009; Wero ́nski \& Elimelech, 2008）。结束对粒子沉积 h− 的讨论，总速率出现在方程式中。 (5.126) 是

\begin{sourceblock}
\includegraphics[page=235,pagebox=mediabox,trim={186.54bp 579.87bp 67.57bp 92.82bp},clip,max width=0.94\textwidth,max height=0.38\textheight]{Marchisio2013.pdf}
\end{sourceblock}
当粒子仍然粘附在固体表面时，h+ 为零。

\subsection*{5.7.2 颗粒破碎}\phantomsection\addcontentsline{toc}{subsection}{5.7.2 颗粒破碎}
在一阶过程中，颗粒破碎是最常见的过程之一。当某种形式的应力干扰保持颗粒完整的力时，就会发生颗粒破碎。固体颗粒（即晶体、聚集体等）的破碎动力学与流体颗粒（即气泡或液滴）的破碎动力学非常不同。一般来说，有两种类型的应力会导致破损：(i) 由碰撞引起的应力和 (ii) 由流体中的速度梯度引起的应力。在碰撞的情况下，相关的力通常是由粒子与粒子的碰撞或粒子与壁的碰撞引起的。后者又可以分为静止壁和移动壁的冲击（Ramachandran et al., 2009）。晶体撞击结晶器中的叶轮叶片（也称为二次成核）或颗粒撞击弯管中的管壁是非常常见的例子。这些情况下涉及的能量通常足以破坏化学键（即离子键、共价键或氢键）。对于这种情况，由于很难从第一原理推导出破坏核，因此使用纯经验或半经验表达式（Myerson，2002；Randolph＆Larson，1971）。在由流体中的速度梯度引起的应力的情况下，所涉及的能量要小得多，因此只有较弱的键被破坏，即由范德华力或静电起源的其他吸引力引起的键，它们通常将单独的初级粒子聚集在一起。对于这种情况，相对容易推导出破碎内核的理论表达式。一般来说，断裂核被写成断裂力和内聚力之间比率的非线性函数。文献中采用了不同的函数形式（例如幂律和指数），导致了不同的表达式。通过强调对粒径和破碎驱动力的函数依赖性，获得了不同的表达式。一般来说，破碎驱动力可以用撞击、碰撞或与表面碰撞破碎的粒子的相对速度Ur来确定；或者在颗粒因流体引起的应力而破碎的情况下使用流体剪切速率 Gf。例如，幂律破坏核是

\begin{sourceblock}
\includegraphics[page=235,pagebox=mediabox,trim={215.00bp 177.37bp 67.57bp 505.10bp},clip,max width=0.94\textwidth,max height=0.38\textheight]{Marchisio2013.pdf}
\end{sourceblock}
而指数核是

\begin{sourceblock}
\includegraphics[page=235,pagebox=mediabox,trim={248.64bp 118.36bp 67.57bp 542.79bp},clip,max width=0.94\textwidth,max height=0.38\textheight]{Marchisio2013.pdf}
\end{sourceblock}
b = C4Gf exp
在这些表达式中，C1–C6 是常数，其值通常通过实验数据拟合来确定（Marchisio et al., 2006）。对于颗粒因撞击静止壁或运动表面而破碎的情况，用冲击速度Ur代替剪切速率Gf，相应的常数具有不同的单位和不同的含义。幂律核在很大程度上纯粹是在经验基础上使用的，但也得到了 Ghadiri 和Zhang (2002) 的机械模型的支持，该模型预测破损程度与颗粒尺寸 dp 和 U2 r 成正比，并且受到 Moreno-Atanasio 和 Ghadiri (2006) 的模型的支持，该模型预测了损坏率，定义为破损接触与初始键数的比率，与 d5/3 p 和到 U2 r。有关推导的更多示例和更多详细信息，请参阅 Ramachandran 等人的著作。 (2009) 和 Rajniak 等人。 （2008）。指数核的函数形式具有理论基础，下面将针对湍流中粒子破碎的特殊情况进行讨论。

\paragraph*{湍流中的颗粒破碎}
在稀颗粒湍流中，作用在颗粒上的流体剪切速率近似为 Gf = √εf/νf，其中 εf 是连续相中的湍流耗散速率，必须从湍流模型中找到。后者引入了很大的不确定性，因为对于许多分散的多相流（例如气液流），不存在精确的湍流模型。然而，对于在湍流流体中经历破碎的颗粒，文献特别丰富，但几乎所有可用的模型都是高度经验性的。最流行的工程模型之一是（Coulaloglou \& Tavlarides，1977）

\begin{sourceblock}
\includegraphics[page=236,pagebox=mediabox,trim={176.02bp 381.41bp 65.08bp 279.12bp},clip,max width=0.94\textwidth,max height=0.38\textheight]{Marchisio2013.pdf}
\end{sourceblock}
其中 C1 和 C2 是通过拟合实验数据确定的两个常数。该内核最初是在局部各向同性假设下针对流体液滴导出的，并假设液滴尺寸在惯性子范围内。如果涡流传递到颗粒的动能大于液滴表面能，则会发生破裂事件。破碎频率是通过考虑具有足够动能导致系统中存在的破碎的涡流的分数而得出的。这可以通过估计颗粒破碎的分数并通过流体速度 DLL(r) 的二阶结构函数估计破碎所需的特征时间 (tb) 来完成。该函数表示给定距离 r 处两点之间速度差的协方差，根据湍流理论，该函数与 (εfr)2/3 成正比。假设子粒子的运动类似于湍流中两团流体的相对运动，并且距离r是粒子尺寸，则可以得到以下近似值：(dp/tb)2/3 ≈DLL(dp) ∝(εfL)2/3。这可用于求解 tb，从而得出方程： （5.138）。 Narsimhan 等人提出了另一种流行的工程内核，源自分散在另一种液体中的液滴。 （1979）：

\begin{sourceblock}
\includegraphics[page=236,pagebox=mediabox,trim={239.86bp 160.43bp 65.08bp 503.64bp},clip,max width=0.94\textwidth,max height=0.38\textheight]{Marchisio2013.pdf}
\end{sourceblock}
b = C3 erfc
它是通过将湍流涡旋到达具有最大稳定直径 d*p 的液滴表面视为泊松过程而得出的。原始表达被 Alopaeus 等人修改。 (2002) 通过添加对涡流碰撞频率中的 ε1/3 f 和能量平衡中的粘性应力的依赖性，考虑了碰撞转变为破裂的概率。拉科宁等人。 (2006)考虑到液滴和气泡的破碎机制密切相关，使用这个核来描述气泡的破碎。他们进一步修改了表达式，用连续相粘度代替分散相粘度，考虑到反对气泡破裂的力的一个组成部分是由周围液体的粘性应力给出的。破坏核的最终表达式为

\begin{sourceblock}
\includegraphics[page=237,pagebox=mediabox,trim={135.54bp 570.69bp 67.57bp 87.57bp},clip,max width=0.94\textwidth,max height=0.38\textheight]{Marchisio2013.pdf}
\end{sourceblock}
在这种情况下，C4、C5 和 C6 也是可以通过理论估计的常数，但随后通过拟合实验数据进行调整。看来常数 C4 的长度为负三分之二，根据理论，应该与整数长度尺度相关联。这意味着C4不是真正的动力学常数，而是包含一些标度参数；因此，在使用方程式时应格外小心。 （5.140）。没有可靠的方法来先验计算 C4；因此其值是通过拟合实验数据来确定的。对于 C5 和 C6，Laakkonen 等人采用的值。通常使用(2006)，即分别为0.04和0.01。在上述模型的相同唯象简化下，考虑到惯性子范围内涡流群体的不同贡献，Luo \& Svendsen (1996) 推导出了分散在湍流液体中的气泡的以下工程关系式：

\begin{sourceblock}
\includegraphics[page=237,pagebox=mediabox,trim={104.89bp 373.71bp 67.57bp 286.87bp},clip,max width=0.94\textwidth,max height=0.38\textheight]{Marchisio2013.pdf}
\end{sourceblock}
p, β = 3/2, cf = f 2/3 + (1 −f)2/3 −1 为表面积的增加系数，f = (dp/d′

\begin{sourceblock}
\includegraphics[page=237,pagebox=mediabox,trim={65.03bp 352.05bp 358.04bp 333.77bp},clip,max width=0.94\textwidth,max height=0.38\textheight]{Marchisio2013.pdf}
\end{sourceblock}
p)3 是两个碎片之一的体积分数，dp 是子颗粒的尺寸，d' p 是母颗粒的尺寸。方程（5.141）不是纯粹的破坏核，而是破坏核和仅对二元破坏（即形成两个片段）有效的子分布函数的乘积。值得一提的是，对于该内核，唯一必须任意固定的参数是 zemin。该量代表有助于颗粒破碎的湍流涡流的最小长度尺度（相对于母体颗粒尺寸标准化）。该数量通常为零（Sanyal 等，2005）。文献中提出的其他内核是 Martinez-Baz ́an 等人的内核。 (1999) 以及 Sathyagal 和 Ramkrishna (1996)。对于不同内核（尤其是流体-流体系统）的回顾和比较，读者可以参考 Patruno 等人的工作。 （2009）。与液滴的破碎一样，固体聚集体的破碎在很大程度上取决于颗粒尺寸与最小湍流涡流之间的比率。当颗粒的尺寸大于最小涡流的尺寸时，由于作用于颗粒表面的压力波动而产生瞬时法向应力，可能会发生破裂。对于小于湍流微尺度的颗粒尺寸，破碎可能是由作用在颗粒相对两侧的湍流动态速度差产生的剪切应力引起的。在这些条件下，破碎核可以写成聚集强度 (τs) 的函数，可以估计为 τs = 9kcorφpF/(8πdo)，其中 do 是初级粒子的直径，kco 是配位数，F 是两个初级粒子之间的粒子间力（Ayazi Shamlou 等，1994）。产生的破损内核是

\begin{sourceblock}
\includegraphics[page=237,pagebox=mediabox,trim={168.42bp 50.16bp 67.57bp 615.90bp},clip,max width=0.94\textwidth,max height=0.38\textheight]{Marchisio2013.pdf}
\end{sourceblock}
粒子间力 F 可以计算为 F = AHdo/(12H2 o)，其中 AH 是液体粒子系统的 Hamaker 参数，Ho 是两个初级粒子之间的距离。配位数基于实验观察，可计算为 kc ≈15φ1.2 p，其中 φp 是聚集体中固体的体积分数。在致密（或固体）颗粒的情况下， φp 接近于 1，而在分形聚集体的情况下，一旦知道了聚集体的分形维数 Df，就可以确定 φp： φp = (0.414Df −0.211)(dp/do)Df−3，其中 dp 是颗粒的尺寸，do 是初级颗粒的尺寸 (Vanni, 2000b)。从相同的基本考虑出发，但使用高度简化的现象学方法，Kramer \& Clark (1999) 开发了固体骨料的水动力破碎模型。他们通过定义失效的表现、诱导和位置对失效模式进行分类，并指出当应力张量的最大特征值大于聚合强度 (τp) 时，聚合体会破裂。假设强度等于或低于破坏强度的聚集体键的数量不是应变率的线性函数，并且在破裂模型中还应考虑尺寸依赖性，破裂核的最终形式为 b = C7να

\begin{sourceblock}
\includegraphics[page=238,pagebox=mediabox,trim={372.80bp 444.09bp 96.08bp 237.81bp},clip,max width=0.94\textwidth,max height=0.38\textheight]{Marchisio2013.pdf}
\end{sourceblock}
p，其中 C7、α、β 和 γ 是模型常数。通过半理论的考虑和对实验数据的拟合，提出了指数 γ 的几个值，即 γ = 0、1/3、2/3 和 1 (Luo \& Svendsen, 1996)。 Pandya \& Spielman (1982) 发现 γ = 1，而 Peng \& Williams (1994) 通过将模型与实验数据进行拟合，发现 γ 可以取 1 到 3 之间的值，并且 β = 1/2。 Serra \& Casamitjana (1998) 拟合实验数据，发现破碎核与湍流耗散率之间的关系取决于固体分数。当固体分数较低时，β = 0.9，而对于较高的固体分数值，β 则要高得多。这可以通过以下事实来解释：固体分数越高，粒子与粒子的碰撞就越有效。他们的结果证实了 Spicer 和 Pratsinis (1996) 的发现，他们使用了 Oles (1991) 的数据并发现对于低固体浓度值，β = 0.8。在对草酸钙的成核、生长、团聚和破碎动力学的研究中，Zauner \& Jones (2000) 发现破碎率和湍流耗散率 εf 之间存在线性相关性（即 β = 1），并且破碎核与尺寸无关（γ = 0）。 Wojcik \& Jones (1998) 使用尺寸相关的破碎内核，发现内核与湍流耗散率之间存在线性关系 (β = 1)。

\paragraph*{子分配函数}
对于子分布函数，根据分散相的性质（即固体与流体颗粒），可能出现非常不同的情况，如 Zac-cone 等人所解释的。 (2007) 和帕特鲁诺等人。 （2009）。一般来说，最简单的函数形式是 delta 函数的求和：

\begin{sourceblock}
\includegraphics[page=238,pagebox=mediabox,trim={188.10bp 133.27bp 65.08bp 533.68bp},clip,max width=0.94\textwidth,max height=0.38\textheight]{Marchisio2013.pdf}
\end{sourceblock}
其中 ν 是子代的总数，xii 是一个函数，表示第 i 个子代的内坐标与父代 ψ′ p 的内坐标之间的关系。方程（5.143）表明，当粒子破碎时，会形成 ν 个碎片，并且表征每个碎片的内坐标等于 xii。这里需要注意的是，在制定子分布函数的形状时，必须考虑一些额外的约束。例如，如果考虑的内坐标是粒子质量（即 Mp），则在破碎过程中，所有子碎片的质量总和必须等于母粒子的质量。在二元断裂（ν = 2）的情况下，两个碎片相同的情况被标记为对称断裂，而不是质量不等时的不对称断裂。如果 f 是两个子颗粒之一的质量分数，则当 f 与 0.5 相差很大时，破碎过程称为侵蚀。所有上述功能均指形成具有已知特征的碎片，这意味着，采用这种描述，每当发生破损事件时，都会形成相同的碎片。然而，更合理的是，由于构成分散相的许多元素的多样性，有时会发生一种类型的破碎（例如对称），而其他时候则遵循另一种机制（即侵蚀）。为了表示这种情况，可以采用连续函数，而不是增量函数的求和。最简单的连续函数是均匀分布，它表明所有可能的片段都是等概率的。只有指定了内部坐标，才能讨论有关此分布的函数形式的更多详细信息。例如，在单个内坐标 m = ψp 的情况下，仅假设正值并在破碎过程中表现出附加属性（即碎片内坐标之和等于母粒子的内坐标之和），与粒子质量的情况一样，可得到以下表达式：

\begin{sourceblock}
\includegraphics[page=239,pagebox=mediabox,trim={179.81bp 371.77bp 67.57bp 278.04bp},clip,max width=0.94\textwidth,max height=0.38\textheight]{Marchisio2013.pdf}
\end{sourceblock}
当然，形成内坐标 m 大于母粒子 m' 的碎片的概率为零。通过将形成的碎片的质量加起来等于母粒子的质量，

\begin{sourceblock}
\includegraphics[page=239,pagebox=mediabox,trim={192.92bp 289.02bp 67.57bp 377.18bp},clip,max width=0.94\textwidth,max height=0.38\textheight]{Marchisio2013.pdf}
\end{sourceblock}
很容易证明这样的函数仅对于二元破坏（即 ν = 2）是一致的。另一个非常流行的子分布函数是使用三个系数定义的抛物线分布。通过在感兴趣的范围内施加分布的正性（即 0 < m < m′）并通过施加碎片总质量守恒，可以获得以下函数：

\begin{sourceblock}
\includegraphics[page=239,pagebox=mediabox,trim={87.41bp 170.57bp 67.57bp 479.06bp},clip,max width=0.94\textwidth,max height=0.38\textheight]{Marchisio2013.pdf}
\end{sourceblock}
其中系数C代表第三自由度（0≤C≤3）。方程（5.146）总是导致形成的碎片总数等于二，并且通过改变系数C，可以获得不同的形状。对于 0 ≤ C < 2，很可能形成不同大小的碎片（即侵蚀），而对于 2 < C ≤ 3，对称断裂的可能性更大。 C = 2 恢复等式中的均匀分布。 （5.144）。当内坐标是粒子质量时，可以很容易地使用上面报告的子分布函数。在这种情况下，子分布函数具有反质量单位（即 kg−1）。一旦颗粒质量和颗粒尺寸之间的关系得到满足，就可以用颗粒尺寸 dp 来写出相同的分布

\begin{sourceblock}
\includegraphics[page=240,pagebox=mediabox,trim={130.30bp 625.04bp 158.14bp 59.82bp},clip,max width=0.94\textwidth,max height=0.38\textheight]{Marchisio2013.pdf}
\end{sourceblock}
导出： N(dp|d′ p 在等维常密度粒子的情况下。在这种情况下，子分布函数具有粒子尺寸的倒数单位。对于这种分布，代表形成的碎片数量和质量守恒的约束现在用相对于分布的子粒子尺寸的零阶和三阶矩来表示。这个版本的子分布函数经常在实际问题中使用。例如，Coulaloglou \& Tavlarides (1977) 从 Valentas 等人 (1966) 的早期工作出发，制定了以下函数：

\begin{sourceblock}
\includegraphics[page=240,pagebox=mediabox,trim={161.60bp 501.31bp 65.08bp 159.64bp},clip,max width=0.94\textwidth,max height=0.38\textheight]{Marchisio2013.pdf}
\end{sourceblock}
根据 Konno 等人的说法，导致形成三个碎片。 (1983) 是搅拌釜中液-液分散体最有可能发生的事件。该分布的最可能事件是碎片的尺寸相当于母体颗粒尺寸的69\%，或者换句话说，碎片的体积（或质量）为母体颗粒的三分之一。另一种经常使用的函数形式是 beta 分布（Hsia \& Tavlarides，1983；Konno 等人，1983；Lee 等人，1987），主要是因为它通过仅改变几个参数就可以灵活地表示非常不同的机制（即对称破坏、均匀分布、侵蚀）。 Diemer \& Olson (2002b) 对 Beta 分布进行了推广，他们通过根据两个自由参数制定模型来调整 Hill \& Ng (1996) 的先前模型。更改这些参数时，分布的形状会发生变化，但质量（或体积）守恒得到保证。最终的表达式为53

\begin{sourceblock}
\includegraphics[page=240,pagebox=mediabox,trim={128.64bp 297.85bp 65.08bp 366.22bp},clip,max width=0.94\textwidth,max height=0.38\textheight]{Marchisio2013.pdf}
\end{sourceblock}
其中 ν 是形成的碎片数量，q 是形状参数。当 q 大于 1 时，最可能的事件是形成 ν 个同等大小的碎片，而当 q 小于 1 时，倾向于形成具有 ν -1 个卫星碎片的一个大碎片。方程中的分布（5.148）与 Laakkonen 等人提出的非常相似。

53 正标准化变量 z 的 β 分布为

\begin{sourceblock}
\includegraphics[page=240,pagebox=mediabox,trim={196.91bp 196.12bp 194.46bp 482.91bp},clip,max width=0.94\textwidth,max height=0.38\textheight]{Marchisio2013.pdf}
\end{sourceblock}
如果 z 是子母质量比，则相应的子分布函数为

\begin{sourceblock}
\includegraphics[page=240,pagebox=mediabox,trim={67.52bp 116.38bp 152.15bp 518.87bp},clip,max width=0.94\textwidth,max height=0.38\textheight]{Marchisio2013.pdf}
\end{sourceblock}
对于致密粒子，我们有 m = ρpkVd3 p，并且

\begin{sourceblock}
\includegraphics[page=240,pagebox=mediabox,trim={140.59bp 67.71bp 137.94bp 601.97bp},clip,max width=0.94\textwidth,max height=0.38\textheight]{Marchisio2013.pdf}
\end{sourceblock}
{\small 可以很容易地简化为等式： （5.148）。\par}
（2007）。在从物理推理和结合物理参数导出的湍流破坏模型中，最流行的模型之一是由 Luo 和 Svendsen (1996) 制定的。子分布函数包含在破损内核中，可以通过进行标准归一化来提取和隔离。该模型包含涡流尺度积分下限作为唯一的自由参数，即涡流的最小长度尺度，相对于颗粒尺寸进行归一化，能够引起有效的颗粒破碎。这个量，在方程式中表示为 zemin。 (5.141) 对最终预测有很大影响，通常取为零。该分布仅对二元破碎有效，并且在低湍流耗散率下显示出一种指示不对称碎片形成的形式，最小值对应于对称破碎 f = 0.5，对于 f = 0 和 f = 1，概率接近无穷大。正如预期的那样，分布的形状是局部湍流条件通过湍流耗散率的函数，并且曲线随着湍流强度的增加而变平。对称破坏的概率增加得非常快，但是对这个参数的依赖似乎是高度非单调的，一些作者（Lasheras et al., 2002）认为这种行为并不能准确地描述物理现实。 Wang 等人开发的一种最新模型源自 Luo 和 Svendsen (1996) 引入的方法，但克服了其一些局限性。 （2003）。子分布函数的最终形状是有上限的，并且当子与父尺寸比率等于零或为一时预测零概率以及以相同尺寸为特征的子颗粒的最小概率。然而，同样对于这种方法，由于其公式的数学复杂性，需要求解三维积分，因此在模拟代码中的实现很麻烦。 Martinez-Baz ́an 等人提出的分布要简单得多。 (1999)，它是在纯机械原理的基础上推导出来的。该分布是通过建立一对粒子形成的概率与作用于两个子粒子的湍流动能和表面能之间的差值之间的比例关系而得出的。经过一些修改，假设二元破坏的原始想法也扩展到三元破坏。除了 Coulaloglou 和 Tavlarides (1977)、Diemer 和 Olson (2002b) 以及 Luo 和 Svendsen (1996) 之外，上述许多模型都存在一个缺点，即质量守恒定律并不总是得到遵守。不预测质量守恒的模型一旦在封闭系统的模拟代码中实现，就会导致预测体积分数的演变，从物理角度来看，这种演变是没有意义的。大多数子分布函数可以轻松扩展到二元问题。让我们考虑两个例子。在第一个示例中，描述了具有两种组分A和B的颗粒。颗粒系统根据这些颗粒的尺寸dp和颗粒的组成来定义，例如表示为颗粒中组分A的质量分数。 例如，当粒子破裂时，我们可以假设组分 A 的数量与碎片的质量成比例地分配给子体。在这些假设下，以及遵循 beta 分布的二元破坏的附加假设，所得的二元分布为

\begin{sourceblock}
\includegraphics[page=241,pagebox=mediabox,trim={128.55bp 96.17bp 67.57bp 566.57bp},clip,max width=0.94\textwidth,max height=0.38\textheight]{Marchisio2013.pdf}
\end{sourceblock}
{\small p 和 φ' 代表母粒子的属性，而 dp 和 φ 代表子粒子的属性 (Buffo et al., 2011, 2012)。第二个例子是分形聚合的碎片化。在这种情况下，聚合体由主要组成\par}
{\small where d′\par}
具有无穷小接触点的粒子。当聚集体分裂成两个相等的碎片时，聚集体的质量和表面积都减半。如果用颗粒体积 Vp 和表面积 Ap 来描述颗粒，则子分布函数为 (Zucca et al., 2007)

\begin{sourceblock}
\includegraphics[page=242,pagebox=mediabox,trim={156.29bp 556.69bp 65.08bp 114.59bp},clip,max width=0.94\textwidth,max height=0.38\textheight]{Marchisio2013.pdf}
\end{sourceblock}
不假设子粒子相等的其他形式也是可能的。总之，很难先验地定义哪种子分布函数最能描述破损。对于固体颗粒来说尤其如此，在这种情况下，根据颗粒的性质（例如致密生长的晶体与松散的分形聚集体），控制整个过程的机制非常不同。在这些情况下，始终建议通过与实验数据比较来进行广泛的验证。对于流体颗粒，情况稍微简单一些。事实上，对于这种情况，β 分布似乎是最合适的。然而，在这种情况下，两个参数 ν 和 q 的识别也必须通过与实验数据进行比较来进行。涉及不同特性流体的实验（例如，参见 Andersson \& Andersson (2006) 和 Zaccone 等人 (2007)）表明，在湍流条件下破裂期间，液滴会变形，其表面会延伸，直到形成一种活性状态。在流体颗粒破裂之前不久，内部流体重新分配机制决定了形成的子体的大小和数量。在气泡的情况下，构成分散相的流体提供的阻力非常小，因此形成两个相同尺寸的气泡（即ν = 2且q > 1）。对于液滴，分散相的较高粘度阻碍了这一过程，导致产生更多具有 U 形分布 (q < 1) 的液滴 (ν = 3 或 4)。所有这些对于清洁系统都是有效的；一旦杂质和其他物质能够在流体-流体界面积聚并改变气泡或液滴表面特性，就会发生完全不同的情况。不幸的是，目前还没有可靠的模型来描述此类系统。

\section*{5.8 二阶点过程}\phantomsection\addcontentsline{toc}{section}{5.8 二阶点过程}
二阶点过程是涉及两个粒子的不连续事件。正如已经提到的，根据定义，这些过程会导致粒子内坐标的瞬时变化。粒子之间的碰撞和聚集通常被视为二阶点过程，这意味着在碰撞事件中仅涉及两个粒子54，并且假设粒子动量的变化是瞬时的并且在碰撞时发生。聚集、附聚和合并是其他常见的例子。如果碰撞后两个粒子粘在一起形成新的单个粒子，则不仅粒子速度而且其他属性（例如粒子质量）都会瞬间改变。因此，我们将首先根据通用二阶点过程推导 GPBE 的源项。为简单起见，我们将考虑颗粒系统（即没有流体相）并表示颗粒速度和
54 这一假设具有广泛的适用性，因为在大多数实际情况下，碰撞仅是两个粒子之间的实际相互作用。这对于稀系统来说当然是正确的，在稀系统中发现三个粒子非常接近的概率非常低，并且对于较稠密的系统也可以被认为是正确的，因为涉及两个以上粒子的任何相互作用都可以合理地描述为两个粒子之间的许多后续碰撞。

属性由单个相空间向量 η 决定。和以前一样，我们将 GPBE 中的源项写为

\begin{sourceblock}
\includegraphics[page=243,pagebox=mediabox,trim={216.49bp 601.76bp 219.03bp 80.83bp},clip,max width=0.94\textwidth,max height=0.38\textheight]{Marchisio2013.pdf}
\end{sourceblock}
其中 h+ 是生产率，h− 是二阶点过程造成的颗粒损失率。55

\subsection*{5.8.1 源项的推导}\phantomsection\addcontentsline{toc}{subsection}{5.8.1 源项的推导}
让我们想象一个粒子的质心位于位置 x′，并由相空间向量 η′ 表征。此外，让我们想象另一个粒子的质心位于物理点 ̃x 处，并具有相空间向量 ̃η。我们定义该二阶点过程的频率 56 a(x′, η′; ̃x, ̃η)，因此量 a(x′, η′; ̃x, ̃η)dt 表示在时间间隔 dt 内经历点过程的粒子的分数。57 该频率关于粒子排序的排列是对称的，从而得到 a(x′, η′; ̃x, ̃η)dt。 ̃x, ̃η) = a( ̃x, ̃η; x′, η′)。我们现在定义一个对数密度函数，使得 n(2)(t, x′, η′, ̃x, ̃η)dx′ dη′ d ̃x d ̃η 表示具有 (x′, η′) 和 ( ̃x, ̃η) 的粒子对的预期数量。正如第 4 章所讨论的，由于了解该函数需要解涉及 n(3) 的方程，因此经常使用以下闭包：n(2)(t, x′, η′, ̃x, ̃η) ≈n(t, x′, η′)n(t, ̃x, ̃η)。给定这些定义，我们可以将涉及粒子对 (x′, η′) 和 ( ̃x, ̃η) 的每单位时间事件数计算为 a(x′, η′; ̃x, ̃η)n(t, x′, η′)n(t, ̃x, ̃η)dx′ dη′ d ̃x d ̃η。如果我们现在想要计算每单位时间和单位体积发生的涉及位于 ( ̃x, ̃η) 附近的粒子的事件数量，我们必须考虑测试粒子和场粒子之间所有可能的相互作用。在二元粒子相遇的命名法中，位于 (̃x, ̃η) 的粒子称为测试粒子，而位于 (x', η') 的粒子称为场粒子。最终碰撞频率结果为

\begin{sourceblock}
\includegraphics[page=243,pagebox=mediabox,trim={136.29bp 301.85bp 67.57bp 371.64bp},clip,max width=0.94\textwidth,max height=0.38\textheight]{Marchisio2013.pdf}
\end{sourceblock}
其中积分是在场粒子的所有位置 x′ 和相空间变量 η′ 上的。在大多数实际情况下，发生相遇时测试颗粒和现场颗粒之间的物理空间距离（即 |x′ − ̃x|）是颗粒直径的数量级58，并且在该距离内，NDF 不会发生显著变化。因此，经常调用以下闭包：n(t, x′, η′)n(t, ̃x, ̃η) ≈n(t, ̃x, η′)n(t, ̃x, ̃η)。如果将这个简化的表达式代入方程： (5.151) 并且积分的顺序反转，我们得到

\begin{sourceblock}
\includegraphics[page=243,pagebox=mediabox,trim={127.71bp 191.89bp 67.57bp 481.09bp},clip,max width=0.94\textwidth,max height=0.38\textheight]{Marchisio2013.pdf}
\end{sourceblock}
方程式中括号之间出现的频率的积分。 （5.152）称为内核：

\begin{sourceblock}
\includegraphics[page=243,pagebox=mediabox,trim={179.56bp 145.08bp 67.57bp 529.20bp},clip,max width=0.94\textwidth,max height=0.38\textheight]{Marchisio2013.pdf}
\end{sourceblock}
55 对于硬球碰撞等二阶点过程，粒子总数保持不变。然而，具有特定相空间矢量的粒子数量总是会增加或减少。 56 作为标准频率，该量具有反时间单位。 57这个量可以看作是二阶点过程在无穷小时间步dt内发生的概率。

\footnotetext[58]{在粒子与硬球势相互作用的情况下，该距离恰好等于相互作用粒子的半径之和。当粒子与其他电势相互作用时，这个距离更大，但仍然保持在这个数量级。}
特定二阶点过程的核具有单位时间空间体积的维度，它粗略地表示位于物理点̃x且内坐标靠近η′和̃η的经历二阶点过程的粒子的体积通量。该表达式通常用于导出经历不同物理过程（即硬球碰撞、布朗运动、差异沉降等）的颗粒系统的内核。由于粒子相遇主要由大于粒子尺寸的距离处的局域传输现象和较短距离处的相互作用势主导，因此内核自然地分为两项。第一个解释了相互作用区域中的粒子传输，而第二个通常称为效率，解释了短程相互作用。下面将讨论一些例子，对更多细节感兴趣的读者可以参考 Ramkrishna (2000)。现在让我们推导由于二阶点过程而出现在 GPBE 中的源项（对应于第 137 页方程（5.1）右侧的最后一项）。 GPBE 在 NDF n(t, x, vp, Ψp, vf, Ψf) 上运行；因此，源项必须表示粒子改变其数量密度的净速率。与之前对颗粒系统的说明类似，推导是通过考虑内坐标约为 ̃η 的测试粒子与位于约 η' 的场粒子的相互作用来完成的，它们的相遇会产生以 η 为特征的粒子。然而，现在 η 的定义包括所有相空间变量（vp、Σp、vf 和 Σf）。该源项由正项和负项构成。积极的术语是（Ramkrishna，2000）

\begin{sourceblock}
\includegraphics[page=244,pagebox=mediabox,trim={140.27bp 367.37bp 65.08bp 300.44bp},clip,max width=0.94\textwidth,max height=0.38\textheight]{Marchisio2013.pdf}
\end{sourceblock}
当用核来写时，其中 δ 是避免多重计数的对称因子，J( ̃η, η) 是将碰撞前的相空间变量与碰撞后的值相关联的变量变换的雅可比行列式。对于产生单个粒子（例如聚集、聚结等）的二阶点过程，对称因子对于相同粒子为 2，对于不同粒子为 1，而对于粒子数量不变的二阶点过程，对称因子始终等于 1。变量变换的雅可比行列式是包含测试粒子 ̃η 的内坐标向量相对于内坐标 η 最终值的偏导数的矩阵的行列式，而场粒子 η′ 的内坐标保持不变：

\begin{sourceblock}
\includegraphics[page=244,pagebox=mediabox,trim={156.65bp 188.99bp 65.08bp 459.39bp},clip,max width=0.94\textwidth,max height=0.38\textheight]{Marchisio2013.pdf}
\end{sourceblock}
̃η、η′ 和 η 之间的关系可以通过根据内部坐标的事件前值和事件后值关联连续性语句写十来导出，并且取决于所研究的二阶点过程。正源项也可以写成频率（而不是核）：

\begin{sourceblock}
\includegraphics[page=244,pagebox=mediabox,trim={107.63bp 101.06bp 65.08bp 566.75bp},clip,max width=0.94\textwidth,max height=0.38\textheight]{Marchisio2013.pdf}
\end{sourceblock}
{\small 负源项更容易找到：\par}
\begin{sourceblock}
\includegraphics[page=244,pagebox=mediabox,trim={158.25bp 56.39bp 65.08bp 617.88bp},clip,max width=0.94\textwidth,max height=0.38\textheight]{Marchisio2013.pdf}
\end{sourceblock}
当按照内核编写时，或者

\begin{sourceblock}
\includegraphics[page=245,pagebox=mediabox,trim={137.96bp 603.76bp 67.57bp 70.51bp},clip,max width=0.94\textwidth,max height=0.38\textheight]{Marchisio2013.pdf}
\end{sourceblock}
当用频率来写时。下面将通过几个例子进行分析和讨论。硬球碰撞的源项在第 6 章中介绍。

\subsection*{5.8.2 聚合和合并的源项}\phantomsection\addcontentsline{toc}{subsection}{5.8.2 聚合和合并的源项}
现在让我们考虑两个粒子碰撞并粘在一起形成单个粒子的情况。一般来说，质量和动量将守恒，因此碰撞后的相空间变量将通过代数约束与碰撞前两个粒子的值相关。根据守恒量的数量，Jaco-bian J(̃η; η) 的计算可能会很复杂。为了简单起见，我们将从一个粒子非常小的例子开始，因此它们的速度与流体的速度相同。然后我们将考虑更复杂的情况，其中碰撞前的粒子速度不同。

\paragraph*{细颗粒聚集}
考虑以 NDF 为特征的粒子群，NDF 是时间、物理空间和粒子质量的函数，即 n(t, x, ψp)。我们将 ψ′ p 和 ̃ ψp 定义为粒子的质量，这些粒子粘在一起后将产生质量为 ψp 的粒子。在这种情况下， 是单个守恒语句，即 ψp = ψ′ p + ̃ψp，导致雅可比行列式等于 1。最终的源项（此处根据内核编写）是

\begin{sourceblock}
\includegraphics[page=245,pagebox=mediabox,trim={95.53bp 312.31bp 67.57bp 353.89bp},clip,max width=0.94\textwidth,max height=0.38\textheight]{Marchisio2013.pdf}
\end{sourceblock}
\begin{sourceblock}
\includegraphics[page=245,pagebox=mediabox,trim={76.82bp 249.14bp 67.57bp 403.51bp},clip,max width=0.94\textwidth,max height=0.38\textheight]{Marchisio2013.pdf}
\end{sourceblock}
如果颗粒密度恒定，则这种描述等效于基于颗粒体积的描述。另一方面，考虑一个具有两个内坐标（例如粒子质量 ψp1 和粒子表面积 ψp2）的系统，这两个内坐标在聚合事件期间是相加的： ψp = ψ′

\begin{sourceblock}
\includegraphics[page=245,pagebox=mediabox,trim={146.41bp 195.04bp 177.87bp 488.40bp},clip,max width=0.94\textwidth,max height=0.38\textheight]{Marchisio2013.pdf}
\end{sourceblock}
p2)，和 ̃ Ψp = ( ̃ Ψp1, ̃ Ψp2)。该线性变换产生单位雅可比行列式，对应的源项是

\begin{sourceblock}
\includegraphics[page=245,pagebox=mediabox,trim={89.27bp 151.81bp 67.57bp 514.39bp},clip,max width=0.94\textwidth,max height=0.38\textheight]{Marchisio2013.pdf}
\end{sourceblock}
这是方程的多元版本。 （5.160）。如果相关的内坐标是颗粒尺寸， ψp = dp，则必须对颗粒的形状做出一些假设。如果假设聚集后所得颗粒的形状与原始颗粒的形状相似，59并且如果颗粒质量（即聚合过程中的守恒量）

\footnotetext[59]{例如，在气泡或液滴的聚结过程中，这是合理的，因为在两个球形流体颗粒聚集后，将形成另一个球形流体颗粒。该假设在附聚过程中也被认为是有效的（即晶体在过饱和溶液中的聚集），因为由于分子生长，聚集晶体的接触点之间留下的空隙被新的结晶材料填充。}
\begin{sourceblock}
\includegraphics[page=246,pagebox=mediabox,trim={67.52bp 588.24bp 65.06bp 69.38bp},clip,max width=0.94\textwidth,max height=0.38\textheight]{Marchisio2013.pdf}
\end{sourceblock}
\begin{sourceblock}
\includegraphics[page=246,pagebox=mediabox,trim={91.20bp 521.56bp 65.08bp 117.81bp},clip,max width=0.94\textwidth,max height=0.38\textheight]{Marchisio2013.pdf}
\end{sourceblock}
这完全等价于等式。 （5.160）。上面给出的例子明确假设粒子具有与流体相同的速度（即粒子斯托克斯数为空）。在下一个例子中，我们考虑具有有限斯托克斯数的粒子的更一般情况。

\paragraph*{液滴的聚结}
\begin{sourceblock}
\includegraphics[page=246,pagebox=mediabox,trim={67.52bp 382.03bp 65.56bp 288.40bp},clip,max width=0.94\textwidth,max height=0.38\textheight]{Marchisio2013.pdf}
\end{sourceblock}
\begin{sourceblock}
\includegraphics[page=246,pagebox=mediabox,trim={194.26bp 336.80bp 65.08bp 328.95bp},clip,max width=0.94\textwidth,max height=0.38\textheight]{Marchisio2013.pdf}
\end{sourceblock}
结果源项是

\begin{sourceblock}
\includegraphics[page=246,pagebox=mediabox,trim={98.42bp 255.62bp 65.08bp 385.30bp},clip,max width=0.94\textwidth,max height=0.38\textheight]{Marchisio2013.pdf}
\end{sourceblock}
正如我们接下来将讨论的，方程中使用的合并内核。 (5.164) 通常取决于速度差： β(xi′

\begin{sourceblock}
\includegraphics[page=246,pagebox=mediabox,trim={178.15bp 218.11bp 207.73bp 466.74bp},clip,max width=0.94\textwidth,max height=0.38\textheight]{Marchisio2013.pdf}
\end{sourceblock}
p−vp|)。由于不同的内核主要是粒径 dp 和一些附加变量的函数，因此仅突出和讨论这些函数依赖性。

\subsection*{5.8.3 细颗粒聚集核}\phantomsection\addcontentsline{toc}{subsection}{5.8.3 细颗粒聚集核}
细颗粒的聚集频率（即 Stp = 0）很大程度上取决于使颗粒紧密靠近的传输现象。一旦定义了频率，就可以使用等式 1 中的定义来执行内核的计算。 （5.153）。我们将讨论仅限于最终结果，对程序细节感兴趣的读者可以查阅 Ramkrishna (2000)。然而，回想一下，对于细颗粒，NDF 中仅出现内坐标 ψp。布朗运动是一种重要现象，它使细颗粒紧密接触并引发碰撞（可能还有聚集）事件。正如已经解释的那样，由于布朗运动而移动的粒子表现出方程（1）中的扩散系数。 (5.116)。通过计算穿过以尺寸为 dp 的粒子为中心、半径等于 (dp + d' p)/2 的球体的尺寸为 d' p 的粒子的净体积通量，可以轻松导出所得的聚集核 60（尺寸为 dp 和 d' p 的粒子），从而得出以下众所周知的方程（Elimelech 等人，1998）：

\begin{sourceblock}
\includegraphics[page=247,pagebox=mediabox,trim={127.10bp 553.88bp 67.57bp 107.48bp},clip,max width=0.94\textwidth,max height=0.38\textheight]{Marchisio2013.pdf}
\end{sourceblock}
其中 μf 是流体粘度。如果聚集颗粒的尺寸大致相同，则将得到尺寸无关的核：β(dp, d′ p) = 8kBTf/(3μf)。这些内核是针对连续体状态导出的，换句话说，假设连续相被粒子感知为连续体。正如已经讨论过的，只有当粒子努森数（相对于连续相）定义为

\begin{sourceblock}
\includegraphics[page=247,pagebox=mediabox,trim={221.28bp 461.28bp 67.57bp 212.18bp},clip,max width=0.94\textwidth,max height=0.38\textheight]{Marchisio2013.pdf}
\end{sourceblock}
小于 1，其中 λf 是连续相分子的平均自由程。当克努森数远小于单位数时，我们就处于连续统（或斯托克斯-爱因斯坦）状态。对于较高的 Kn* 值，我们进入自由分子（或 Epstein）状态，其中使用以下聚合内核：

\begin{sourceblock}
\includegraphics[page=247,pagebox=mediabox,trim={130.76bp 362.05bp 67.57bp 300.69bp},clip,max width=0.94\textwidth,max height=0.38\textheight]{Marchisio2013.pdf}
\end{sourceblock}
为了使用涵盖两种状态以及它们之间的过渡的单个表达式，通常使用 Fuchs (1964) 方程：

\begin{sourceblock}
\includegraphics[page=247,pagebox=mediabox,trim={85.14bp 295.91bp 67.57bp 364.02bp},clip,max width=0.94\textwidth,max height=0.38\textheight]{Marchisio2013.pdf}
\end{sourceblock}
其中修正后的扩散系数61为

\begin{sourceblock}
\includegraphics[page=247,pagebox=mediabox,trim={158.45bp 248.60bp 67.57bp 416.72bp},clip,max width=0.94\textwidth,max height=0.38\textheight]{Marchisio2013.pdf}
\end{sourceblock}
p + l2)3/2]/(6dpl) −2dp，l = 8Γ/(πc)，c = 8kBTf/(πMp)，其中 Mp 是粒子质量。其中 g = [(2dp + l)3 −(4d2 使细颗粒紧密靠近并引起聚集的另一种机制是流体速度中空间梯度的存在。让我们考虑平行运动的两组颗粒，它们遵循两条距离小于其尺寸的线。如果它们以相同的速度移动，则它们永远不会发生碰撞，而速度差异将导致碰撞。从这个简单的论证和计算来看，与用于近运动内核的那些非常相似，获得以下聚合 kernel62（von Smoluchowski，1917）：

\begin{sourceblock}
\includegraphics[page=247,pagebox=mediabox,trim={189.86bp 107.69bp 67.57bp 564.06bp},clip,max width=0.94\textwidth,max height=0.38\textheight]{Marchisio2013.pdf}
\end{sourceblock}
60 这种聚集核以及一般情况下粒子因布朗运动而聚集的现象，被标记为近运动。

\footnotetext[61]{比较方程 6 中的标准斯托克斯-爱因斯坦方程很有用。 （5.116）与方程的校正表达式。 （5.169）。}
\footnotetext[62]{这通常称为正交运动聚合内核。}
其中 Gf 是流体剪切速率。对于非常简单的流动，假设剪切速率等于主要的局部流体速度梯度。一般来说，对于层流，流体剪切率计算为应变率张量的第二个不变量。 63 在湍流系统中，流体剪切率通常通过湍流耗散率和连续相的运动粘度的值计算，结果为 Gf = √εf/νf，如 Saffman \& Turner (1956) 中所述。仅当粒子小于柯尔莫哥洛夫长度尺度 ηk = (ν3 f /εf)1/4 时，这才有效，从而得出以下表达式：

\begin{sourceblock}
\includegraphics[page=248,pagebox=mediabox,trim={105.78bp 517.02bp 65.08bp 149.13bp},clip,max width=0.94\textwidth,max height=0.38\textheight]{Marchisio2013.pdf}
\end{sourceblock}
因为在这种情况下，粒子在被困在同一个涡流中时会发生碰撞。当粒子大于柯尔莫哥洛夫长度尺度时，通常使用以下表达式：

\begin{sourceblock}
\includegraphics[page=248,pagebox=mediabox,trim={174.07bp 454.17bp 65.08bp 216.83bp},clip,max width=0.94\textwidth,max height=0.38\textheight]{Marchisio2013.pdf}
\end{sourceblock}
这些内核对于无惯性粒子（即 Stp = 0）有效，并且只能通过采用临时校正来扩展到有限斯托克斯数。例如，Ammar \& Reeks (2009) 针对 Saffman \& Turner (1956) 提出的内核导出了基于局部斯托克斯数的修正。近运动聚集与正运动聚集的相对重要性通过 P ́eclet 数来量化：

\begin{sourceblock}
\includegraphics[page=248,pagebox=mediabox,trim={214.33bp 353.24bp 65.08bp 334.59bp},clip,max width=0.94\textwidth,max height=0.38\textheight]{Marchisio2013.pdf}
\end{sourceblock}
事实上，当 Pe ≤ 0.001 时，只需考虑近运动聚集，而当 Pe ≥ 10 时，近运动贡献可以忽略不计。总聚合内核通常被写为上述内核以及其他冲突机制产生的其他可能贡献的总和。然而，以这种方式获得的数量仅占由于碰撞而产生的总粒子体积通量。为了考虑到并非所有碰撞都能成功产生新聚合的事实，引入了效率因子 α。该效率因子是一个介于 0 和 1 之间的数字（尽管在某些情况下，它也可以采用大于 1 的值），它量化了导致聚合事件的成功冲突数量与使用上面报告的内核计算的冲突总数的比较。有许多因素阻碍两个粒子成功聚集。事实上，即使两个粒子非常接近，只有排斥力和吸引力之间的平衡才能决定粒子是否会粘在一起。事实上，为了让颗粒粘在一起，需要吸引力，并且通常必须克服潜在的排斥力。颗粒系统中典型的吸引力（范德华力）是由永久粒子之间的吸引力产生的。

63 众所周知，应变率张量的分量定义为

\begin{sourceblock}
\includegraphics[page=248,pagebox=mediabox,trim={208.20bp 103.06bp 205.34bp 569.45bp},clip,max width=0.94\textwidth,max height=0.38\textheight]{Marchisio2013.pdf}
\end{sourceblock}
剪切率通常计算为应变率张量的第二个不变量（第一个不变量是迹）：

\begin{sourceblock}
\includegraphics[page=248,pagebox=mediabox,trim={211.99bp 68.79bp 209.74bp 611.71bp},clip,max width=0.94\textwidth,max height=0.38\textheight]{Marchisio2013.pdf}
\end{sourceblock}
{\small 其中重复的索引意味着求和（即爱因斯坦符号）。\par}
偶极子（基索姆力）、永久偶极子和相应的感应偶极子之间（德拜力）以及两个瞬时感应偶极子之间（伦敦色散力）。其他吸引力是指颗粒表面上存在稳定电荷（静电力）或在颗粒表面上形成化学键的可能性，如氢键的情况。其他值得注意的情况是悬浮在过饱和溶液中的生长晶体以及由于粘合剂的存在而在颗粒接触点上形成的稳定颈。正如 Marshall 等人所讨论的，后一个因素在结晶和造粒中非常重要。 （2011）。相反，排斥力是由连续相的存在引起的。对于悬浮在气态主相中的颗粒，会产生非常弱（如果不是完全可以忽略不计）的排斥力。对于悬浮在液体中的颗粒，排斥力通常是由颗粒周围的离子双电层引起的。事实上，当带电粒子悬浮在液相（水相）中时，它立即被交替的离子层和反离子层包围。当两个粒子彼此靠近时，两个离子云必须重叠，并且由于渗透压，会产生排斥力。另一个原因是两个碰撞颗粒之间（在它们即将发生物理接触之前）液膜的粘性排出，从而产生阻力。在这种情况下，流体可以（在分形多孔聚集体中）或不能（在固体聚集体中）渗透碰撞颗粒的事实起着非常重要的作用。根据作用在颗粒上的相关力，可以得出不同的聚集效率表达式。这些表达式通常是通过将聚集效率写为 Fuchs 稳定性比（即 α = 1/W）的倒数而导出的，对于尺寸为 dp 和 d' 的两个颗粒

\begin{sourceblock}
\includegraphics[page=249,pagebox=mediabox,trim={137.09bp 324.29bp 67.57bp 324.42bp},clip,max width=0.94\textwidth,max height=0.38\textheight]{Marchisio2013.pdf}
\end{sourceblock}
其中 h 是粒子间距离，V(h) 是粒子间电势。粒子间势被计算为涉及吸引力 Va(h) 和排斥力 Vr(h) 的势总和，对于保守系统来说，它表示作用在粒子上的力（因为粒子间力为 F = −∂V/∂h）。范德华吸引力通常用以下势函数表示：

\begin{sourceblock}
\includegraphics[page=249,pagebox=mediabox,trim={69.97bp 212.04bp 67.57bp 441.09bp},clip,max width=0.94\textwidth,max height=0.38\textheight]{Marchisio2013.pdf}
\end{sourceblock}
其中 AH 是 Hamaker 参数，dp 和 d′ p 是两个碰撞粒子的尺寸。通常，AH 的范围为 10−21–10−19 J，并且可以使用不同的方法进行估计。根据 Lifshitz 理论，可以从界面张力确定 Hamaker 参数（Israelachvili，1992）：AH ≈24πγD0，其中 D0 是截止距离，通常可以假设为 1-2 Å 左右。非导电粒子系统的 Hamaker 参数值可以根据以下方程计算（Venables \& Schmuttenmaer, 2000）：AH = AH(ω = 0) + AH(ω > 0) f(h)，其中 AH (ω = 0) 和 AH (ω > 0) 分别是总结相互作用的零频和非零频分量影响的两个常数，h 是原子中心。延迟的影响包含在函数 f(h) 中（Vanni \& Baldi，2002）。这三个项取决于溶液的组成，这些计算的详细信息可以在 Lince 等人中找到。 （2008）。与排斥力（双层性质）相关的势通常用粒子的表面势 Ψs 来表示。在不同的简化假设下可以得出不同的表达式（读者可以参考 Elimelech 等人 (1998) 的详细评论）。最流行的方法之一是 Gregory（Elimelech 等人，1998）针对尺寸为 dp 和 d′ 的粒子开发的方法

\begin{sourceblock}
\includegraphics[page=250,pagebox=mediabox,trim={140.43bp 550.38bp 65.08bp 97.28bp},clip,max width=0.94\textwidth,max height=0.38\textheight]{Marchisio2013.pdf}
\end{sourceblock}
其中 ε0 和 εr 是绝对和相对介电常数，κ 是 Debye-H ̈uckel 参数，z 是离子的化合价（本例中假设为对称 z-z 电解质），e 是基本电荷。表面电势又可以写成溶液中某些电势决定离子 (PDI) 浓度的函数，如下所示：

\begin{sourceblock}
\includegraphics[page=250,pagebox=mediabox,trim={214.76bp 460.22bp 65.08bp 211.63bp},clip,max width=0.94\textwidth,max height=0.38\textheight]{Marchisio2013.pdf}
\end{sourceblock}
Ψs = kBTf
其中 Ci 是 PDI 的浓度，Cpzc i 是 PDI 在零电荷点（当颗粒表面电荷为零时）的浓度。64 不同化学系统的典型计算示例可以在 Gavi 等人中找到。 （2007、2010）。方程（5.174）通常用于通过布朗运动而靠近的粒子，并忽略任何类型的短程流体动力相互作用。事实上，唯一考虑的相互作用是与吸引的范德华力和排斥双层力相关的相互作用。通过修改方程来考虑流体动力学效应。 （5.174）如下：

\begin{sourceblock}
\includegraphics[page=250,pagebox=mediabox,trim={130.58bp 318.07bp 65.08bp 342.56bp},clip,max width=0.94\textwidth,max height=0.38\textheight]{Marchisio2013.pdf}
\end{sourceblock}
其中 Go f 是表达粒子之间流体动力学相互作用的函数。例如，在多孔颗粒65的情况下，可以假设函数 Go f (h) 等于 1，或者可以近似为

\begin{sourceblock}
\includegraphics[page=250,pagebox=mediabox,trim={175.38bp 239.85bp 65.08bp 429.12bp},clip,max width=0.94\textwidth,max height=0.38\textheight]{Marchisio2013.pdf}
\end{sourceblock}
在两个完全不可穿透的粒子（斯托克斯蠕动流）的情况下，dp 大小相等。对于不同尺寸 dp 和 d′ p 的粒子，只能使用近似表达式。例如，对于趋于零的 h，我们得到 Gf(h) = 2h(dp + d′ p)/(dpd′ p)。一旦粒子相互作用势 V(h) 被定义，方程： （5.174）可以积分。通常针对输入参数的不同值（碰撞粒子的尺寸，即 dp 和 d' p、表面电势 Ψs 等）以数值方式进行积分。然后将结果收集在查找表中，该查找表定义了碰撞效率对这些参数的依赖性：α = α(dp, d′ p, Ψs)。然而，对于一些简化的情况，积分方程。 （5.174）进行分析，得到效率的函数形式。可能的场景数量非常多，因此不可能进行全面、完整的讨论。一般来说，每个实际案例都必须单独分析，不同的情况

\footnotetext[64]{kBTf/(ze) 一词通常显示为 RT/(zF)，其中 F 是法拉第常数，等于阿伏伽德罗数乘以基本电荷。}
\footnotetext[65]{实际上，在这种情况下，流体阻力不起任何作用，可以完全忽略不计，从而得出式（1）。 （5.174）。}
作用在粒子上的力和驱动它们碰撞的机制必须被识别和量化。下面将仅给出几个导致聚合效率的分析表达式的示例。第一个例子涉及非常小的（亚微米或布朗）粒子通过一些明确的吸引和排斥势相互作用。方程中出现的积分。 (5.174) 可以通过使用 Reerink \& Overbeek (1954) 近似来解析求解，得到以下聚合效率表达式：

\begin{sourceblock}
\includegraphics[page=251,pagebox=mediabox,trim={191.55bp 526.42bp 67.57bp 151.70bp},clip,max width=0.94\textwidth,max height=0.38\textheight]{Marchisio2013.pdf}
\end{sourceblock}
kBTf

其中 κ 是 Debye-H ̈uckel 参数，φmax 是粒子聚集的能垒（即势能与粒子间距离关系图中的最大值）。该最大值是 PDI 浓度、溶液总离子强度和其他操作参数的函数。第二个例子是指颗粒在范德华吸引力和粘性排斥力的作用下经历剪切诱导的聚集。在这种情况下，聚集效率得出 (Wang et al., 2005a) α = kFl−0.18，其中 k 是单位阶的前置因子，通常通过实验数据拟合来确定，Fl 是流量数，66 对于不渗透颗粒，定义为 Fl = 6πμf(dp + d′ p)Gf/AH，其中 μf 是流体粘度，Gf 是流体剪切力。在可渗透粒子的情况下（如分形物体的情况），采用类似（但更复杂）的方程，更多细节可以在 Wang 等人的文章中找到。 （2005b）。另一个值得一提的有趣案例是大颗粒在过饱和溶液中聚集的情况。在大颗粒的情况下，范德华吸引力不是很重要，聚集效率通常由生长颗粒形成稳定连接颈的能力决定。仅当生长该连接颈所需的时间 τn 小于两个颗粒保持紧密接触的时间 τf 时，这才有可能。对于颗粒在剪切作用下聚集的情况，该时间与剪切速率的倒数相关，τf = 1/Gf。生长连接颈的时间计算为 τn = Db/( f(δ) ̇dp)，其中 ̇dp 是粒子生长速率（例如参见方程（5.23）），Db 是稳定桥的直径，f(δ) 是取决于碰撞粒子的尺寸比 δ 的形状函数。大多数这些参数可以通过独立实验或拟合来确定。然后，聚集（或更准确地说，团聚）效率计算为 α = e−τn/τc。从该表达式可以明显看出，当粒子在时间间隔 τc（远长于形成稳定颈部所需的时间 τn）内彼此保持靠近时，它们的比率趋于零，效率趋于一致。相反，当 τc <τn 时，效率趋于零。有关特定示例的更多详细信息，请参阅 Marchi-sio 等人。 （2003a）。在许多情况下，从第一原理推导聚集效率是不可能的，而是采用经验效率。67由于在这些情况下效率取决于粒子的组成，因此必须将其作为内坐标的一部分。一个例子是二元聚集，即 A 型粒子与 B 型粒子聚集（Lee 等人，2008 年；Mat-soukas 等人，2006 年、2009 年）。用于表达这种效率的高度经验形式，提供与实验研究的二元系统的定性类比（Rajniak

\footnotetext[66]{流数可以解释为剪切力和范德华力之间的比率。 67这在造粒等领域非常常见。在造粒过程中，通过添加称为粘合剂的粘性物质来诱导颗粒聚集。}
et al., 2007) 为 α(c, c′) = exp[−aAAcc′ −aBB(1 −c)(1 −c′) −aAB(c + c′ −2cc′)]，其中 c 和 c′ 是两个聚集颗粒内组分 A 的浓度（例如质量分数）。三个参数 aAA、aBB 和 aAB 控制自聚集（即涉及相同类型颗粒的 A-A 和 B-B 聚集）与交叉聚集（即颗粒 A 和颗粒 B 之间的聚集）的强度。显然，当 aAA = aBB = aAB = 0 时，组合依赖性被消除，聚合效率降低至统一。例如，如果 aAA = -3、aBB = +2、aAB = -2，如 Rajniak 等人讨论的情况。 (2007)，组分 A 的自聚集优于组分 B 的自聚集，而交叉聚集的强度介于两者之间。这些系数的值通常通过拟合实验数据来确定。

\subsection*{5.8.4 液滴和气泡的聚结核}\phantomsection\addcontentsline{toc}{subsection}{5.8.4 液滴和气泡的聚结核}
由于液滴和气泡是可变形的，聚结流体颗粒的处理变得复杂。尽管如此，对于这些情况，可用的内核通常被编写为冲突频率和效率的乘积。碰撞频率一般是通过类比气体动力学理论推导出来的（详细见第6章），其计算基于在直径等于两个碰撞粒子半径之和的碰撞圆柱体中运动的一个粒子所扫过的粒子数。对于尺寸为 dp 和 d′ p 的两个粒子，碰撞频率可以近似为

\begin{sourceblock}
\includegraphics[page=252,pagebox=mediabox,trim={173.85bp 373.84bp 65.08bp 302.99bp},clip,max width=0.94\textwidth,max height=0.38\textheight]{Marchisio2013.pdf}
\end{sourceblock}
其中第一项代表碰撞圆柱体的截面，第二项，即涉及 vr = vp −v′ p 的项，量化了两个碰撞粒子的相对速度。函数η≤1是下面将描述的聚结效率。对于具有有限斯托克斯数的粒子（例如液体喷雾），vr 由 GPBE 可知。否则， vr 必须通过其他方式进行近似。一旦两个粒子非常靠近并引发碰撞事件，它们可能会或可能不会合并。可能的情况有很大不同，具体取决于主要阶段的性质。如果主相是气体，则碰撞颗粒成功聚结的概率主要由韦伯数和冲击参数决定。后一个量通常被假定为由碰撞粒子中心的矢量与粒子速度差 vr 所形成的角度的正弦。根据这个定义，当速度差平行于连接液滴中心的矢量时，冲击参数为零；相反，当两者正交时，冲击参数等于 1。韦伯数定义为

\begin{sourceblock}
\includegraphics[page=252,pagebox=mediabox,trim={217.04bp 165.10bp 65.08bp 508.74bp},clip,max width=0.94\textwidth,max height=0.38\textheight]{Marchisio2013.pdf}
\end{sourceblock}
其中ρp是分散相的材料密度，σ是连续相和分散相之间的界面张力。韦伯数是流体惯性相对于界面张力的相对重要性的量度，对于非常高的界面张力来说韦伯数非常小。当韦伯数小于10时，聚结的可能性很大，效率η非常接近于1。相比之下，随着 We 的增加，只有当影响因子接近 0.2 时，聚结才成为可能，对应于约 10-15° 的影响角。在这些条件下，不可能发生合并（即 η ≈0），并且碰撞后的液滴据说会经历横向或正面弹跳以及反射或拉伸分离。

如果连续相是液体，聚结的主要障碍是两个颗粒之间的小空间中液体膜的排出。在这些情况下，效率通常被量化为液滴接触和薄膜排水的特征时间之比的函数（通常是负指数函数）。例如，在小气泡由于湍流速度波动而聚结的情况下，聚结核呈现以下形式（Buffo 等人，2012；Laakkonen 等人，2006；Petitti 等人，2010）

\begin{sourceblock}
\includegraphics[page=253,pagebox=mediabox,trim={135.94bp 542.33bp 67.57bp 140.42bp},clip,max width=0.94\textwidth,max height=0.38\textheight]{Marchisio2013.pdf}
\end{sourceblock}
而聚结效率则写为：

\begin{sourceblock}
\includegraphics[page=253,pagebox=mediabox,trim={154.30bp 484.67bp 67.57bp 179.40bp},clip,max width=0.94\textwidth,max height=0.38\textheight]{Marchisio2013.pdf}
\end{sourceblock}
对特定系统的更多细节感兴趣的读者应该查阅专业文献。

\chapter{硬球碰撞模型}
在本章中，我们考虑具有相同（单分散）或不同（多分散）密度和直径的光滑（即无摩擦）球形颗粒之间的碰撞模型。为简单起见，我们仅考虑粒子质量和直径守恒的碰撞，并排除可能改变这些属性的其他过程（例如表面凝聚或聚集（Fox 等，2008））。同样，假设光滑的球形粒子意味着粒子角动量在碰撞过程中不会改变，因此在动力学方程中只需要考虑粒子速度。我们将讨论限制在硬球碰撞上，这意味着碰撞后的粒子速度可以写为碰撞前粒子速度的显式函数，但也考虑具有恒定恢复系数和有限尺寸粒子效应的非弹性碰撞。有关硬球碰撞的更多详细信息，请参阅 Cercignani 的书籍（Cercignani，1975、1988、2000；Cercignani 等人，1994）。我们还将简要讨论更简单的碰撞模型，这些模型通常用于近似稀极限下的硬球碰撞模型。其中包括适用于单分散颗粒的麦克斯韦模型 (Maxwell, 1879) 和两个线性化碰撞模型（即 BGK (Bhatnagar et al., 1954) 和 ES-BGK (Holway, 1966)），以及将 ES-BGK 模型扩展到多分散颗粒的模型。关于非硬球类型碰撞的动力学模型的讨论可以在 Struchtrup (2005) 中找到。作为一个例子，在第 6.4 节中，我们将考虑控制速度数密度函数 f(t, x, v) 的动力学方程：

\begin{sourceblock}
\includegraphics[page=254,pagebox=mediabox,trim={210.92bp 201.00bp 65.07bp 485.44bp},clip,max width=0.94\textwidth,max height=0.38\textheight]{Marchisio2013.pdf}
\end{sourceblock}
其中 t 是时间，x 是空间位置，v 是粒子速度，C 表示粒子与粒子的碰撞。为了简化表示，只有当自变量不同于集合(t,x,v)时，我们才会将其包含在公式中。方程中的密度函数。 (6.1) 被归一化使得

\begin{sourceblock}
\includegraphics[page=254,pagebox=mediabox,trim={217.95bp 112.13bp 215.51bp 554.32bp},clip,max width=0.94\textwidth,max height=0.38\textheight]{Marchisio2013.pdf}
\end{sourceblock}
其中 αp 是分散相体积分数。由于碰撞不会改变粒子的数量，因此碰撞模型具有以下属性：

\begin{sourceblock}
\includegraphics[page=254,pagebox=mediabox,trim={219.95bp 50.16bp 217.50bp 617.50bp},clip,max width=0.94\textwidth,max height=0.38\textheight]{Marchisio2013.pdf}
\end{sourceblock}
此外，动量守恒意味着

\begin{sourceblock}
\includegraphics[page=255,pagebox=mediabox,trim={214.96bp 594.16bp 217.50bp 72.30bp},clip,max width=0.94\textwidth,max height=0.38\textheight]{Marchisio2013.pdf}
\end{sourceblock}
最后，碰撞引起的动能变化要求

\begin{sourceblock}
\includegraphics[page=255,pagebox=mediabox,trim={213.25bp 544.65bp 215.79bp 121.82bp},clip,max width=0.94\textwidth,max height=0.38\textheight]{Marchisio2013.pdf}
\end{sourceblock}
其中弹性碰撞等式成立。通过推导每个碰撞模型的矩输运方程，将可以验证这些条件是否适用于所有模型。更一般地，由于碰撞而引起的矩变化可以通过在相空间上对函数 ψ(v) 乘以 C 进行积分来找到。本章的主要目标之一是找出单分散和多分散粒子碰撞引起的整数矩的变化。对于整数矩，泛函可以表示为 ψ(v) = vl1 1 vl2 2 vl3 3，其中 l1、l2 和 l3 是非负整数。 Struchtrup (2005) 中详细介绍了如何使用速度矩推导宏观输运方程。本章其余部分的组织如下。首先，在 6.1 节中，我们考虑弹性和非弹性粒子的单分散硬球碰撞的碰撞项。我们为配对相关函数引入了 Boltzmann (1872) 和 Enksog (1921) 的动力学闭包，然后导出任意阶速度矩和整数矩的精确源项。其次，在 6.2 节中，我们考虑多分散硬球碰撞的精确源项，推导任意和整数阶矩的精确表达式。接下来，在第 6.3 节中，我们考虑从精确的碰撞源项导出的单分散和多分散系统的简化动力学模型，并讨论它们相对于硬球碰撞模型的属性。在 6.4 节中，我们讨论从方程 1 导出的矩输运方程的性质。 （6.1）与硬球碰撞模型。最后，在第 6.5 节中，我们简要描述了如何应用基于求积的矩方法来闭合速度矩的碰撞源项。

\section*{6.1 单分散硬球碰撞}\phantomsection\addcontentsline{toc}{section}{6.1 单分散硬球碰撞}
在本节中，我们将使用 Fox \& Vedula (2010) 中的符号简要回顾二元硬球碰撞的碰撞模型。由于弹性硬球碰撞而导致的数密度函数的变化（Boltzmann，1872；Cercignani，1988；Chapman \& Cowling，1961；Enksog，1921）服从以下形式的（未闭合）积分表达式：

\begin{sourceblock}
\includegraphics[page=255,pagebox=mediabox,trim={68.72bp 152.65bp 71.26bp 513.81bp},clip,max width=0.94\textwidth,max height=0.38\textheight]{Marchisio2013.pdf}
\end{sourceblock}
(6.2) 其中dp为粒径，f(2)为配对相关函数，

\begin{sourceblock}
\includegraphics[page=255,pagebox=mediabox,trim={211.48bp 98.94bp 67.56bp 574.18bp},clip,max width=0.94\textwidth,max height=0.38\textheight]{Marchisio2013.pdf}
\end{sourceblock}
\footnotetext[1]{尽管对于单分散情况来说这并不是严格必要的，但我们使用下标 12 的符号来表示碰撞中涉及的两个粒子。这样做的原因是，在多分散情况下，粒子 1 和 2 具有不同的属性，只需很少改变符号即可直接修改碰撞积分。}
\begin{sourceblock}
\includegraphics[page=256,pagebox=mediabox,trim={58.00bp 478.90bp 58.23bp 59.00bp},clip,width=\textwidth,max height=0.72\textheight]{Marchisio2013.pdf}
\par\vspace{0.45\baselineskip}
\small 图6.1。两个直径为 dp 的硬球之间的直接碰撞。
\end{sourceblock}
是粒子中心之间方向的单位向量（所谓的碰撞向量），v12 = v1 −v2（大小为 v12）是直接碰撞之前的速度差。 （这些符号的几何解释参见图 6.1。）直接碰撞的碰撞前速度（v1 和 v2）与碰撞后速度（v′ 1 和 v′ 2）相关（Cercignani，1988；Chapman \& Cowling，1961）

\begin{sourceblock}
\includegraphics[page=256,pagebox=mediabox,trim={200.68bp 354.68bp 65.07bp 314.07bp},clip,max width=0.94\textwidth,max height=0.38\textheight]{Marchisio2013.pdf}
\end{sourceblock}
12 = v′ 1 -v′ 2 是直接碰撞后的速度差。表面 S+ 是单位半球，其上 v12 · x12 = -v′，其中 v′ 12 · x12 = v12 cos θ > 0（即导致碰撞的速度差），其中 θ 是碰撞角（碰撞发生时其值使得 cos θ 为正）。由于碰撞是弹性的，因此反向碰撞之前的预碰撞速度与直接碰撞之后的值相同（非弹性碰撞的情况并非如此，如下所述）。因此，动量和能量守恒直接从碰撞前速度的定义得出，因为

\begin{sourceblock}
\includegraphics[page=256,pagebox=mediabox,trim={180.03bp 219.00bp 175.91bp 449.74bp},clip,max width=0.94\textwidth,max height=0.38\textheight]{Marchisio2013.pdf}
\end{sourceblock}
归一化对相关函数

f (2)(x, v1; y, v2) αp(x)αp(y)

是在同一时间 t 在位置 x 找到速度为 v1 的粒子和在位置 y 找到速度为 v2 的粒子的概率。有关粒子位置和速度之间相关性的信息不包含在单粒子数密度函数中，因此必须在我们使用方程（1）之前进行建模。 （6.2）。对于非弹性碰撞，硬球碰撞模型由下式给出（Fox \& Vedula, 2010；Jenkins \& Mancini, 1989；Jenkins \& Richman, 1985）

\begin{sourceblock}
\includegraphics[page=256,pagebox=mediabox,trim={69.31bp 52.94bp 65.07bp 603.69bp},clip,max width=0.94\textwidth,max height=0.38\textheight]{Marchisio2013.pdf}
\end{sourceblock}
其中 χ = 1/e2 是直接碰撞 (v1, v2) 与反向碰撞 (v'' 1 , v'' 2 ) 的预碰撞速度相关的因子，0 ≤e ≤1 是恢复系数。这些速度的关系为（Cercignani，1988）

\begin{sourceblock}
\includegraphics[page=257,pagebox=mediabox,trim={194.08bp 563.24bp 67.56bp 105.51bp},clip,max width=0.94\textwidth,max height=0.38\textheight]{Marchisio2013.pdf}
\end{sourceblock}
\begin{sourceblock}
\includegraphics[page=257,pagebox=mediabox,trim={65.02bp 543.41bp 309.00bp 143.19bp},clip,max width=0.94\textwidth,max height=0.38\textheight]{Marchisio2013.pdf}
\end{sourceblock}
12 = v'' 1 -v'' 2 （震级为 v'' 12）是逆碰撞前的速度差。请注意，直接碰撞后的碰撞后速度由下式给出

\begin{sourceblock}
\includegraphics[page=257,pagebox=mediabox,trim={194.75bp 489.43bp 67.56bp 179.33bp},clip,max width=0.94\textwidth,max height=0.38\textheight]{Marchisio2013.pdf}
\end{sourceblock}
因此，每次直接碰撞都会导致能量损失，其量取决于 e：

\begin{sourceblock}
\includegraphics[page=257,pagebox=mediabox,trim={138.15bp 437.51bp 140.68bp 235.77bp},clip,max width=0.94\textwidth,max height=0.38\textheight]{Marchisio2013.pdf}
\end{sourceblock}
与弹性情况一样，有必要为方程式中出现的成对相关函数提供闭包。 （6.5）。

\subsection*{6.1.1 玻尔兹曼碰撞模型}\phantomsection\addcontentsline{toc}{subsection}{6.1.1 玻尔兹曼碰撞模型}
由于配对相关函数的存在，方程中的碰撞模型(6.2) 未封闭。因此，为了闭合动力学方程（方程 6.1），我们必须提供用 f 表示的 f (2) 的闭合。最简单的闭包是 Boltzmann Stoßzahlansatz（Boltzmann，1872）：

\begin{sourceblock}
\includegraphics[page=257,pagebox=mediabox,trim={162.14bp 291.37bp 67.56bp 386.60bp},clip,max width=0.94\textwidth,max height=0.38\textheight]{Marchisio2013.pdf}
\end{sourceblock}
f (2)(x, v1; x + dpx12, v2) ≈f(x, v1) f(x, v2),
这相当于假设碰撞粒子的速度在碰撞之前是不相关的，并且 dp 远小于粒子之间的平均距离。根据这一假设，（弹性）玻尔兹曼碰撞模型可写为

\begin{sourceblock}
\includegraphics[page=257,pagebox=mediabox,trim={122.00bp 201.82bp 67.56bp 464.64bp},clip,max width=0.94\textwidth,max height=0.38\textheight]{Marchisio2013.pdf}
\end{sourceblock}
因为积分中的所有项都是在同一空间位置评估的，所以我们放弃了 x 依赖性。弹性玻尔兹曼碰撞模型具有由 C = 0 定义的唯一平衡解。该麦克斯韦分布由下式给出

\begin{sourceblock}
\includegraphics[page=257,pagebox=mediabox,trim={171.30bp 106.90bp 67.56bp 557.58bp},clip,max width=0.94\textwidth,max height=0.38\textheight]{Marchisio2013.pdf}
\end{sourceblock}
其中 Up 是平均速度，定义为

\begin{sourceblock}
\includegraphics[page=257,pagebox=mediabox,trim={208.45bp 50.19bp 210.98bp 616.47bp},clip,max width=0.94\textwidth,max height=0.38\textheight]{Marchisio2013.pdf}
\end{sourceblock}
θp 是粒子能量2，定义为

\begin{sourceblock}
\includegraphics[page=258,pagebox=mediabox,trim={192.25bp 594.96bp 189.80bp 71.70bp},clip,max width=0.94\textwidth,max height=0.38\textheight]{Marchisio2013.pdf}
\end{sourceblock}
请注意方程左侧的传输项。 (6.1) 的大小可以大于或小于碰撞项。对于碰撞项比传输项重要得多的情况，方程的解为：玻尔兹曼碰撞模型的方程（6.1）是局部麦克斯韦模型，其中 αp、Up 和 θp 取决于空间和时间，但 f 可以通过方程 1 很好地近似。 （6.10）。在此限制下，粒子表现为理想气体，平均速度服从欧拉方程。一般来说，弹性玻尔兹曼碰撞模型可以很好地描述各种浓度的分子气体，包括极低浓度或稀薄气体。然而，对于宏观粒子或稠密气体，dp 远小于粒子之间的平均距离的假设不再成立。此外，宏观粒子之间的碰撞通常是非弹性的。因此，接下来我们将从方程（1）开始考虑由于非弹性、有限尺寸、硬球碰撞而导致的速度矩的变化。 （6.5）。

\subsection*{6.1.2 任意矩的碰撞项}\phantomsection\addcontentsline{toc}{subsection}{6.1.2 任意矩的碰撞项}
Boltzmann-Enskog 碰撞产生的速度矩的源项可以使用硬球碰撞积分的属性以紧凑的形式表示（Cer-cignani，1988）。出于一般性的考虑，我们考虑由于方程（1）中的碰撞项而导致的属性 ψ ψ ψ(v) 的变化。 （6.5）。相应的碰撞源项由下式给出

\begin{sourceblock}
\includegraphics[page=258,pagebox=mediabox,trim={67.52bp 214.77bp 65.07bp 345.28bp},clip,max width=0.94\textwidth,max height=0.38\textheight]{Marchisio2013.pdf}
\end{sourceblock}
在右侧的第一个积分中。现在，我们可以在等式右侧出现的第一个多维积分中使用符号的变化。 (6.12)用相应的不带素变量(指逆碰撞后或直接碰撞前的值)替换双素变量(指逆碰撞前的值)，同时将无素变量(指逆碰撞后或直接碰撞前的值)替换为相应的单素变量(指直接碰撞后的值)。换句话说，符号的一致变化允许我们用直接碰撞之前和之后的速度变量来表示由于逆碰撞而产生的速度矩的源项：

\begin{sourceblock}
\includegraphics[page=258,pagebox=mediabox,trim={72.50bp 80.96bp 65.07bp 585.50bp},clip,max width=0.94\textwidth,max height=0.38\textheight]{Marchisio2013.pdf}
\end{sourceblock}
\footnotetext[2]{对于分子系统，温度的定义包括粒子质量。对于颗粒系统，我们将 θp 称为颗粒温度。}
其中 ψ(v′ 1) = ψ (v1 − ω(v12 · x12)x12)。请注意，直接碰撞的配对相关函数是在第二个粒子的 x + dpx12 处计算的，因此变量的一致变化仅使用直接碰撞的 f (2) 形式。接下来一系列步骤的目标是获得式（1）的形式。 （6.13）其中成对相关函数仅取决于一个空间位置（即碰撞期间粒子接触的点）。我们首先将方程式中的对分布函数 f (2)(x −dpx12, v1; x, v2) 与 f (2)(x, v1; x + dpx12, v2) 相加和相减。 （6.13）：

\begin{sourceblock}
\includegraphics[page=259,pagebox=mediabox,trim={69.96bp 471.86bp 67.56bp 146.26bp},clip,max width=0.94\textwidth,max height=0.38\textheight]{Marchisio2013.pdf}
\end{sourceblock}
（注意此过程中包含的系数 1/2。）使用关于碰撞接触点的泰勒级数展开式 3，

\begin{sourceblock}
\includegraphics[page=259,pagebox=mediabox,trim={171.56bp 410.99bp 174.09bp 269.67bp},clip,max width=0.94\textwidth,max height=0.38\textheight]{Marchisio2013.pdf}
\end{sourceblock}
我们现在可以重写方程。 (6.14) 形式4

\begin{sourceblock}
\includegraphics[page=259,pagebox=mediabox,trim={193.24bp 360.95bp 67.56bp 325.75bp},clip,max width=0.94\textwidth,max height=0.38\textheight]{Marchisio2013.pdf}
\end{sourceblock}
3 我们将重复使用涉及向量 x = (x1, x2, x3) 和梯度算子 ∇y = (∂y1, ∂y2, ∂y3) 的多元泰勒级数展开式。在展开式中，当 n = 0, 1, ... 时，将出现 (x · ∇y)nψ(y) 形式的项。。。并定义如下：

\begin{sourceblock}
\includegraphics[page=259,pagebox=mediabox,trim={153.29bp 249.86bp 185.54bp 379.51bp},clip,max width=0.94\textwidth,max height=0.38\textheight]{Marchisio2013.pdf}
\end{sourceblock}
\begin{sourceblock}
\includegraphics[page=259,pagebox=mediabox,trim={153.29bp 198.97bp 154.49bp 454.59bp},clip,max width=0.94\textwidth,max height=0.38\textheight]{Marchisio2013.pdf}
\end{sourceblock}
请注意，右侧的和可以重写为多项式展开式。我们将在推导过程中的几个点上利用这一事实。 4 一个简单的例子可能有助于阐明此步骤中使用的两个泰勒级数展开式。考虑在 x ± a 处评估的函数 h(x)。泰勒级数展开式为

\begin{sourceblock}
\includegraphics[page=259,pagebox=mediabox,trim={151.11bp 129.06bp 153.65bp 548.89bp},clip,max width=0.94\textwidth,max height=0.38\textheight]{Marchisio2013.pdf}
\end{sourceblock}
两个展开式相加和相减得到

\begin{sourceblock}
\includegraphics[page=259,pagebox=mediabox,trim={111.50bp 67.58bp 112.71bp 581.74bp},clip,max width=0.94\textwidth,max height=0.38\textheight]{Marchisio2013.pdf}
\end{sourceblock}
{\small 其形式与用于查找方程的多元展开式相同。 （6.16）和等式。分别为（6.17）。\par}
其中碰撞源项由下式给出

\begin{sourceblock}
\includegraphics[page=260,pagebox=mediabox,trim={98.38bp 564.07bp 65.07bp 75.47bp},clip,max width=0.94\textwidth,max height=0.38\textheight]{Marchisio2013.pdf}
\end{sourceblock}
碰撞通量项由下式给出

\begin{sourceblock}
\includegraphics[page=260,pagebox=mediabox,trim={98.92bp 475.15bp 65.07bp 164.40bp},clip,max width=0.94\textwidth,max height=0.38\textheight]{Marchisio2013.pdf}
\end{sourceblock}
在大多数应用中，n = 1 的求和项。。。 , ∞ 在 C 和 G 的定义中被忽略，今后我们也会这样做。该近似相当于假设 dp 远小于 f (2) 空间梯度的特征长度尺度（即 |dp ∇f (2)| ≪α2 p），因此只有 dp 阶以下的项才重要。下一步是为碰撞源和碰撞通量项中出现的成对相关函数提供闭包。对于中等密度的流动，已知有限尺寸颗粒的碰撞频率大于使用 Boltzmann Stoßzahlansatz 发现的碰撞频率（Carnahan \& Starling，1969；Enksog，1921）。为了考虑这种影响，配对相关函数可以建模为两个单粒子速度分布函数和一个径向分布函数的乘积：

\begin{sourceblock}
\includegraphics[page=260,pagebox=mediabox,trim={77.83bp 312.35bp 113.43bp 368.32bp},clip,max width=0.94\textwidth,max height=0.38\textheight]{Marchisio2013.pdf}
\end{sourceblock}
2 dpx12, v1

\begin{sourceblock}
\includegraphics[page=260,pagebox=mediabox,trim={231.88bp 285.04bp 65.07bp 396.41bp},clip,max width=0.94\textwidth,max height=0.38\textheight]{Marchisio2013.pdf}
\end{sourceblock}
请注意，这个假设只是将对相关函数建模的问题转化为对 g0 建模的新问题。 g0 的常用模型假设径向分布函数既不明确依赖于碰撞角度（即 x12）也不明确依赖于 x。前者相当于假设速度为 v2 的粒子相对于速度为 v1.5 的粒子没有优先空间方向。然后可以将径向分布函数建模为分散相体积分数的函数。例如，典型的模型是（Carnahan \& Starling，1969）

\begin{sourceblock}
\includegraphics[page=260,pagebox=mediabox,trim={67.52bp 130.77bp 205.38bp 524.70bp},clip,max width=0.94\textwidth,max height=0.38\textheight]{Marchisio2013.pdf}
\end{sourceblock}
对于密排球体，p ≈0.63。忽略方程中 g0 的空间梯度。 (6.18) 并使用单粒子密度的泰勒级数近似
5 对于具有相同属性（例如质量、直径等）的粒子之间的碰撞，该假设是合理的。然而，对于具有不同属性的粒子之间的碰撞，当外力对不同粒子的影响不同时，它不太可能成立。例如，小颗粒的流体阻力比大颗粒更强，导致每种类型具有不同的平均速度。那么碰撞更有可能发生在较快移动粒子的上游侧（即速度差的方向）。在这种情况下，g0 的模型需要包含此类信息，以便准确地对碰撞频率进行建模。

关于空间点 x 的函数产生

\begin{sourceblock}
\includegraphics[page=261,pagebox=mediabox,trim={76.95bp 562.52bp 67.56bp 83.23bp},clip,max width=0.94\textwidth,max height=0.38\textheight]{Marchisio2013.pdf}
\end{sourceblock}
通过这个表达式，由于碰撞（即方程（6.15））引起的任意速度矩的变化现在可以用单粒子密度函数来表示。最后一步是使用定义 v′ 1 = v1 − ω(v12 · x12)x12 将 ψ(v′ 1) 写为 x12 的显式函数。采用关于 v1 的泰勒级数展开，我们发现

\begin{sourceblock}
\includegraphics[page=261,pagebox=mediabox,trim={137.30bp 465.25bp 67.56bp 201.69bp},clip,max width=0.94\textwidth,max height=0.38\textheight]{Marchisio2013.pdf}
\end{sourceblock}
仅保留方程中的一阶空间导数项。 （6.19）（即达到总体 dp 阶），碰撞源项和碰撞通量项为

\begin{sourceblock}
\includegraphics[page=261,pagebox=mediabox,trim={71.14bp 363.98bp 67.56bp 268.97bp},clip,max width=0.94\textwidth,max height=0.38\textheight]{Marchisio2013.pdf}
\end{sourceblock}
以及

\begin{sourceblock}
\includegraphics[page=261,pagebox=mediabox,trim={116.45bp 289.16bp 67.56bp 356.63bp},clip,max width=0.94\textwidth,max height=0.38\textheight]{Marchisio2013.pdf}
\end{sourceblock}
在这些定义中，我们可以观察到涉及 x12 的积分有两种类型：

\begin{sourceblock}
\includegraphics[page=261,pagebox=mediabox,trim={71.67bp 232.60bp 67.56bp 433.86bp},clip,max width=0.94\textwidth,max height=0.38\textheight]{Marchisio2013.pdf}
\end{sourceblock}
对于 n ≥1，其中为 x12 的每个分量定义第二个积分。在接下来的两节中，我们将展示可以分析计算这些积分。方程中密度函数的积分。 （6.21）和等式。然后可以使用例如求积（参见第 6.5 节）来处理（6.22），或者，在麦克斯韦模型的特殊情况下，可以通过分析来完成它们。然后，在第 6.1.5 节中，我们将看一个具体示例，其中 ψ 对应于整数速度矩。对于这种情况，方程中的泰勒级数。 (6.20) 是有限的。

\subsection*{6.1.3 碰撞角度和变换矩阵}\phantomsection\addcontentsline{toc}{subsection}{6.1.3 碰撞角度和变换矩阵}
碰撞截面 S+ 上的积分可以通过定义从实验室参考系到碰撞参考系的线性变换 L(v12) 来分析计算（Fox \& Vedula，2010）。线性变换不依赖于 x12（仅依赖于 v12），因此可以在积分之外进行。我们可以定义实验室坐标系和碰撞坐标系之间的变换矩阵 L。此时，没有必要知道 L 的精确表达式（下面将给出）。使用 L，单位向量 x12 变换为 x†
12 = Lx12，其中

\begin{sourceblock}
\includegraphics[page=262,pagebox=mediabox,trim={210.29bp 557.98bp 65.07bp 98.90bp},clip,max width=0.94\textwidth,max height=0.38\textheight]{Marchisio2013.pdf}
\end{sourceblock}
dx† 12 = sin θ dθ dφ，其中 0 ≤θ ≤π/2 和 0 ≤φ ≤2π 是碰撞角。在碰撞参考系中，v12 · x12 = v12 cos θ，其中 v12 = |v1 −v2|与碰撞角度无关。从实验室坐标系 (x) 到碰撞参考系 (x†) 的变换矩阵仅取决于相对速度矢量 v12 = v1 −v2。按照惯例，v12 与 x† 3 轴对齐。因此，定义为 x† = Lx 的正交变换矩阵可以写为

\begin{sourceblock}
\includegraphics[page=262,pagebox=mediabox,trim={157.74bp 428.57bp 65.07bp 234.44bp},clip,max width=0.94\textwidth,max height=0.38\textheight]{Marchisio2013.pdf}
\end{sourceblock}
其中球面角 0 ≤θ1 ≤π 和 0 ≤φ1 ≤2π 与 v12 的三个分量相关，关系为 v12,1 = v12 sin θ1 cos φ1, v12,2 = v12 sin θ1 sin φ1, v12,3 = v12 cos θ1。

\begin{sourceblock}
\includegraphics[page=262,pagebox=mediabox,trim={402.10bp 364.22bp 65.07bp 323.98bp},clip,max width=0.94\textwidth,max height=0.38\textheight]{Marchisio2013.pdf}
\end{sourceblock}
请注意，L 的第一行和第二行的顺序并不重要，并且 L 仅在 v12 > 0 时定义。L 的分量可以写为 v12 的函数，但它们同样可以用 Lij(θ1, φ1) 表示，即作为球面角 θ1 和 φ1 的函数。以下伪代码说明了如何根据向量 y = v12 计算球面角：

\begin{sourceblock}
\includegraphics[page=262,pagebox=mediabox,trim={58.00bp 167.04bp 58.23bp 418.54bp},clip,width=\textwidth,max height=0.66\textheight]{Marchisio2013.pdf}
\end{sourceblock}
读者不应将球面角（θ1 和 φ1）（参数化 v12）与碰撞角（θ 和 φ）（参数化 x12）混淆。事实上，碰撞角的积分是用固定的球面角值（即 v12 的固定值）完成的。总之，对于给定的 v12 值，我们可以计算球面角（θ1 和 φ1），然后根据等式计算变换矩阵 L。 （6.25）。变换矩阵的分量出现在速度矩的碰撞项的定义中。在下一节中，我们将展示方程中的积分。碰撞项的方程（6.23）可以写成 v12 分量的显式函数，因此不需要计算球面角来计算积分。然而，有时用球面角来表达线性变换以简化操作是有用的（例如，当使用符号数学程序计算矩的精确表达式时）。

\subsection*{6.1.4 碰撞角积分}\phantomsection\addcontentsline{toc}{subsection}{6.1.4 碰撞角积分}
我们现在想要明确地评估方程式中出现的碰撞项中的积分。 （6.23）。我们将对涉及碰撞角度的每个项分别执行此操作，然后在最后重建最终结果。项 x12 · ∇vψ(v) 可以在碰撞参考系中使用 x12 = LTx† 12 写为
x12 · ∇vψ(v) = (L11 sin θ cos φ + L21 sin θ sin φ + L31 cos θ)∂ψ(v)

\begin{sourceblock}
\includegraphics[page=263,pagebox=mediabox,trim={334.74bp 478.07bp 143.49bp 208.98bp},clip,max width=0.94\textwidth,max height=0.38\textheight]{Marchisio2013.pdf}
\end{sourceblock}
+ (L12 sin θ cos φ + L22 sin θ sin φ + L32 cos θ)∂ψ(v)

\begin{sourceblock}
\includegraphics[page=263,pagebox=mediabox,trim={343.59bp 426.16bp 67.56bp 234.92bp},clip,max width=0.94\textwidth,max height=0.38\textheight]{Marchisio2013.pdf}
\end{sourceblock}
+ (L13 sin θ cos φ + L23 sin θ sin φ + L33 cos θ)∂ψ(v)

多次应用多项式定理6后，我们发现

\begin{sourceblock}
\includegraphics[page=263,pagebox=mediabox,trim={116.66bp 364.30bp 211.29bp 298.80bp},clip,max width=0.94\textwidth,max height=0.38\textheight]{Marchisio2013.pdf}
\end{sourceblock}
1 ∂vj 2 ∂vk 3 × (L11 sin θ cos φ + L21 sin θ sin φ + L31 cos θ)i × (L12 sin θ cos φ + L22 sin θ sin φ + L32 cos θ) j

\begin{sourceblock}
\includegraphics[page=263,pagebox=mediabox,trim={186.80bp 314.30bp 67.56bp 370.50bp},clip,max width=0.94\textwidth,max height=0.38\textheight]{Marchisio2013.pdf}
\end{sourceblock}
可以进一步重写为

\begin{sourceblock}
\includegraphics[page=263,pagebox=mediabox,trim={100.74bp 225.73bp 145.53bp 414.22bp},clip,max width=0.94\textwidth,max height=0.38\textheight]{Marchisio2013.pdf}
\end{sourceblock}
× (L11 sin θ cos φ)i1(L21 sin θ sin φ)i2(L31 cos θ)i3 × (L12 sin θ cos φ) j1(L22 sin θ sin φ) j2(L32 cos θ) j3

\begin{sourceblock}
\includegraphics[page=263,pagebox=mediabox,trim={170.87bp 177.07bp 67.56bp 507.73bp},clip,max width=0.94\textwidth,max height=0.38\textheight]{Marchisio2013.pdf}
\end{sourceblock}
6 多项式系数的定义例如为 n i, j, k = n!我！ j！k！ ,
其中 i + j + k = n。多项式定理指出，对于任何正整数 m 和非负整数 n

\begin{sourceblock}
\includegraphics[page=263,pagebox=mediabox,trim={141.98bp 78.45bp 144.52bp 590.25bp},clip,max width=0.94\textwidth,max height=0.38\textheight]{Marchisio2013.pdf}
\end{sourceblock}
{\small 其中对所有非负整数索引 k1 到 km 的序列进行求和，使得所有 ki 的总和为 n。\par}
最后作为

\begin{sourceblock}
\includegraphics[page=264,pagebox=mediabox,trim={113.34bp 557.54bp 65.07bp 73.91bp},clip,max width=0.94\textwidth,max height=0.38\textheight]{Marchisio2013.pdf}
\end{sourceblock}
其中我们定义了 i = i1 + i2 + i3、j = j1 + j2 + j3、k = k1 + k2 + k3、n1 = i1 + j1 + k1、n2 = i2 + j2 + k2 和 n3 = i3 + j3 + k3。方程中碰撞角的第一个积分。 （6.23）仅涉及方程最后一行中的正弦和余弦。 (6.30)，乘以 (v12 · x12)n+1 dx†

\begin{sourceblock}
\includegraphics[page=264,pagebox=mediabox,trim={130.33bp 497.55bp 244.78bp 184.02bp},clip,max width=0.94\textwidth,max height=0.38\textheight]{Marchisio2013.pdf}
\end{sourceblock}
\paragraph*{第一个积分}
考虑方程中的第一个积分。 （6.23）并使用等式中的结果。 (6.30)，将参考更改为碰撞系后，碰撞角上的两个积分为

\begin{sourceblock}
\includegraphics[page=264,pagebox=mediabox,trim={93.23bp 407.63bp 65.07bp 258.07bp},clip,max width=0.94\textwidth,max height=0.38\textheight]{Marchisio2013.pdf}
\end{sourceblock}
\begin{sourceblock}
\includegraphics[page=264,pagebox=mediabox,trim={67.52bp 365.75bp 65.07bp 298.94bp},clip,max width=0.94\textwidth,max height=0.38\textheight]{Marchisio2013.pdf}
\end{sourceblock}
其中 B(x, y) 是 beta 函数，它根据 gamma 函数定义为

\begin{sourceblock}
\includegraphics[page=264,pagebox=mediabox,trim={208.71bp 317.24bp 205.91bp 357.02bp},clip,max width=0.94\textwidth,max height=0.38\textheight]{Marchisio2013.pdf}
\end{sourceblock}
将方程式中的两个结果相乘。 （6.31）和等式。 (6.32) 并重新排列伽马函数得出

\begin{sourceblock}
\includegraphics[page=264,pagebox=mediabox,trim={79.14bp 221.08bp 65.07bp 418.87bp},clip,max width=0.94\textwidth,max height=0.38\textheight]{Marchisio2013.pdf}
\end{sourceblock}
其中我们使用恒等式 Γ(x + 1) = x!。请注意，仅当 n1 和 n2 为偶数时，此表达式才为非零。然后我们就可以使用该身份

\begin{sourceblock}
\includegraphics[page=264,pagebox=mediabox,trim={138.83bp 155.32bp 65.07bp 511.62bp},clip,max width=0.94\textwidth,max height=0.38\textheight]{Marchisio2013.pdf}
\end{sourceblock}
将积分项重写为

\begin{sourceblock}
\includegraphics[page=264,pagebox=mediabox,trim={79.14bp 70.88bp 65.07bp 569.44bp},clip,max width=0.94\textwidth,max height=0.38\textheight]{Marchisio2013.pdf}
\end{sourceblock}
{\small 其中 n1 和 n2 是非负偶整数。\par}
从方程。 （6.30）和等式。 （6.35），我们可以观察到，为了计算方程（6.35）中的第一个积分，需要以下恒等式： （6.23）：

\begin{sourceblock}
\includegraphics[page=265,pagebox=mediabox,trim={74.99bp 544.85bp 67.56bp 93.21bp},clip,max width=0.94\textwidth,max height=0.38\textheight]{Marchisio2013.pdf}
\end{sourceblock}
在收集各种项后，我们现在可以将第一个积分写为

\begin{sourceblock}
\includegraphics[page=265,pagebox=mediabox,trim={70.45bp 430.75bp 67.56bp 184.89bp},clip,max width=0.94\textwidth,max height=0.38\textheight]{Marchisio2013.pdf}
\end{sourceblock}
可以重写为

\begin{sourceblock}
\includegraphics[page=265,pagebox=mediabox,trim={68.44bp 236.69bp 125.10bp 294.50bp},clip,max width=0.94\textwidth,max height=0.38\textheight]{Marchisio2013.pdf}
\end{sourceblock}
\begin{sourceblock}
\includegraphics[page=265,pagebox=mediabox,trim={143.72bp 221.77bp 67.56bp 466.06bp},clip,max width=0.94\textwidth,max height=0.38\textheight]{Marchisio2013.pdf}
\end{sourceblock}
我们使用 i、j 和 k 的定义来分布微分算子。使用多项式定理，该表达式可以简化为

\begin{sourceblock}
\includegraphics[page=265,pagebox=mediabox,trim={79.11bp 54.82bp 67.56bp 522.21bp},clip,max width=0.94\textwidth,max height=0.38\textheight]{Marchisio2013.pdf}
\end{sourceblock}
利用 n1 和 n2 是偶数这一事实，该表达式可以重写为

\begin{sourceblock}
\includegraphics[page=266,pagebox=mediabox,trim={80.82bp 495.54bp 65.07bp 77.97bp},clip,max width=0.94\textwidth,max height=0.38\textheight]{Marchisio2013.pdf}
\end{sourceblock}
其中 m1、m2 和 m3 是非负整数，对于固定 n，方程 1 中的求和。 (6.40) 扩展到 mi 的所有值，使得 n = 2m1 + 2m2 + m3。接下来我们想展示等式的右侧。 (6.40) 不依赖于球面角 φ1 和 θ1。首先考虑运营商

\begin{sourceblock}
\includegraphics[page=266,pagebox=mediabox,trim={89.81bp 359.00bp 65.07bp 273.27bp},clip,max width=0.94\textwidth,max height=0.38\textheight]{Marchisio2013.pdf}
\end{sourceblock}
以及

\begin{sourceblock}
\includegraphics[page=266,pagebox=mediabox,trim={88.57bp 258.69bp 65.07bp 373.58bp},clip,max width=0.94\textwidth,max height=0.38\textheight]{Marchisio2013.pdf}
\end{sourceblock}
由于 m1 和 m2 的对称性，方程： （6.41）和等式。 (6.42) 将单独出现在求和中，乘以相同的因子。根据方程中变换矩阵的定义： （6.25），我们有以下关系：

\begin{sourceblock}
\includegraphics[page=266,pagebox=mediabox,trim={182.62bp 97.15bp 65.07bp 491.56bp},clip,max width=0.94\textwidth,max height=0.38\textheight]{Marchisio2013.pdf}
\end{sourceblock}
{\small 这是求和中出现的唯一涉及 L1j 和 L2j 的项。请注意，这意味着方程中求和中的所有项。 (6.40) 涉及 m1 和 m2 是\par}
二次相关于 v12 的分量，而涉及 m3 的分量则线性相关。为了简化符号，我们定义等式的右侧。 (6.40) 使用求和-微分运算符：

\begin{sourceblock}
\includegraphics[page=267,pagebox=mediabox,trim={66.64bp 419.29bp 67.56bp 114.63bp},clip,max width=0.94\textwidth,max height=0.38\textheight]{Marchisio2013.pdf}
\end{sourceblock}
其中 m1、m2 和 m3 是非负整数。碰撞源项中第一个积分的最终形式为

\begin{sourceblock}
\includegraphics[page=267,pagebox=mediabox,trim={127.10bp 344.85bp 67.56bp 318.54bp},clip,max width=0.94\textwidth,max height=0.38\textheight]{Marchisio2013.pdf}
\end{sourceblock}
请注意，v12 的“额外”因子并未包含在等式 1 中运算符的定义中。 （6.44）是故意的。我们将看到，如果没有这个因子，整数矩碰撞积分将是闭合的，因此，如果我们用一个常数（即麦克斯韦模型）替换该因子，则矩方程在每个阶都是闭合的。

\paragraph*{第二次积分}
方程中的第二个积分。 (6.23) 可以针对 x12 = (x12,1, x12,2, x12,3) 的每个分量进行计算。例如，对于第 m 个分量，我们发现

\begin{sourceblock}
\includegraphics[page=267,pagebox=mediabox,trim={75.06bp 88.31bp 90.04bp 479.63bp},clip,max width=0.94\textwidth,max height=0.38\textheight]{Marchisio2013.pdf}
\end{sourceblock}
\begin{sourceblock}
\includegraphics[page=267,pagebox=mediabox,trim={399.61bp 68.48bp 67.56bp 619.72bp},clip,max width=0.94\textwidth,max height=0.38\textheight]{Marchisio2013.pdf}
\end{sourceblock}
该表达式右侧的项涉及三个新的碰撞角积分：

\begin{sourceblock}
\includegraphics[page=268,pagebox=mediabox,trim={71.14bp 495.46bp 65.07bp 92.66bp},clip,max width=0.94\textwidth,max height=0.38\textheight]{Marchisio2013.pdf}
\end{sourceblock}
\begin{sourceblock}
\includegraphics[page=268,pagebox=mediabox,trim={71.14bp 363.11bp 65.07bp 225.00bp},clip,max width=0.94\textwidth,max height=0.38\textheight]{Marchisio2013.pdf}
\end{sourceblock}
以及

\begin{sourceblock}
\includegraphics[page=268,pagebox=mediabox,trim={71.14bp 222.87bp 65.07bp 365.24bp},clip,max width=0.94\textwidth,max height=0.38\textheight]{Marchisio2013.pdf}
\end{sourceblock}
其中双阶乘由 n!! 定义为奇数正整数！ = n(n −2)(n − 4)。。。 (3)(1) 我们使用了恒等式 (2k)！ = (2k−1)!! k！ 2k。请注意等式。仅当 n1 为奇数且 n2 为偶数时，式 (6.47) 才非零，而式 (6.47) 则为非零。仅当 n1 为偶数且 n2 为奇数时，(6.48) 才非零。方程中的结果。仅当 n1 和 n2 均为偶数时，(6.49) 才非零。正如上面对第一个积分所做的那样，可以重写方程式中的求和。 (6.46) 用三个求和微分运算符表示：

\begin{sourceblock}
\includegraphics[page=268,pagebox=mediabox,trim={71.79bp 56.53bp 65.07bp 565.30bp},clip,max width=0.94\textwidth,max height=0.38\textheight]{Marchisio2013.pdf}
\end{sourceblock}
使用与上述第一积分类似的操作，可以证明第一个算子由下式定义：

\begin{sourceblock}
\includegraphics[page=269,pagebox=mediabox,trim={68.31bp 445.01bp 67.56bp 92.22bp},clip,max width=0.94\textwidth,max height=0.38\textheight]{Marchisio2013.pdf}
\end{sourceblock}
其中 m1、m2 和 m3 是非负整数。同样，第二个运算符定义为

\begin{sourceblock}
\includegraphics[page=269,pagebox=mediabox,trim={68.59bp 272.83bp 67.56bp 297.37bp},clip,max width=0.94\textwidth,max height=0.38\textheight]{Marchisio2013.pdf}
\end{sourceblock}
其中 m1、m2 和 m3 是非负整数。最后，第三个运算符定义为

\begin{sourceblock}
\includegraphics[page=269,pagebox=mediabox,trim={68.31bp 76.32bp 67.56bp 457.60bp},clip,max width=0.94\textwidth,max height=0.38\textheight]{Marchisio2013.pdf}
\end{sourceblock}
{\small 其中 m1、m2 和 m3 是非负整数。\par}
\paragraph*{最终表格}
使用碰撞角度积分的解析表达式，我们现在可以重写方程式中的碰撞源项。 (6.21) 作为两个贡献的总和，

\begin{sourceblock}
\includegraphics[page=270,pagebox=mediabox,trim={195.93bp 572.53bp 65.07bp 113.84bp},clip,max width=0.94\textwidth,max height=0.38\textheight]{Marchisio2013.pdf}
\end{sourceblock}
其中

\begin{sourceblock}
\includegraphics[page=270,pagebox=mediabox,trim={100.14bp 509.86bp 65.07bp 153.53bp},clip,max width=0.94\textwidth,max height=0.38\textheight]{Marchisio2013.pdf}
\end{sourceblock}
以及

\begin{sourceblock}
\includegraphics[page=270,pagebox=mediabox,trim={109.41bp 353.25bp 65.07bp 217.12bp},clip,max width=0.94\textwidth,max height=0.38\textheight]{Marchisio2013.pdf}
\end{sourceblock}
同样，方程中碰撞通量项的分量。 (6.22) 可以写成

\begin{sourceblock}
\includegraphics[page=270,pagebox=mediabox,trim={70.86bp 241.28bp 65.07bp 375.88bp},clip,max width=0.94\textwidth,max height=0.38\textheight]{Marchisio2013.pdf}
\end{sourceblock}
此时，碰撞源项中的剩余积分是相对于速度密度函数的。现在我们将考虑特殊情况，其中 ψ 表示相对于速度分量的整数矩。

\subsection*{6.1.5 整数矩的碰撞项}\phantomsection\addcontentsline{toc}{subsection}{6.1.5 整数矩的碰撞项}
对于基于求积的矩方法，需要整数矩的碰撞项。对于这种特殊情况，Fox \& Vedula (2010) 中开发的方法得出了以有限和表示的积分的封闭表达式。对于整数矩，我们有 ψ(v) = vl1 1 vl2 2 vl3 3 以及方程式中 n 的和中的项。当 n > max(l1, l2, l3) 时，(6.55)–(6.57) 为零。对于这种特殊情况，我们将 m = 0 时定义 I(m) l1l2l3

\begin{sourceblock}
\includegraphics[page=270,pagebox=mediabox,trim={133.47bp 50.16bp 65.07bp 616.77bp},clip,max width=0.94\textwidth,max height=0.38\textheight]{Marchisio2013.pdf}
\end{sourceblock}
并且，对于 m = 1, 2, 3，我们定义

\begin{sourceblock}
\includegraphics[page=271,pagebox=mediabox,trim={67.62bp 544.34bp 67.56bp 82.78bp},clip,max width=0.94\textwidth,max height=0.38\textheight]{Marchisio2013.pdf}
\end{sourceblock}
使用第 6.1.4 节中运算符的定义，我们可以将这些项重写为 v1 分量中的多项式：

\begin{sourceblock}
\includegraphics[page=271,pagebox=mediabox,trim={123.32bp 456.51bp 67.56bp 217.82bp},clip,max width=0.94\textwidth,max height=0.38\textheight]{Marchisio2013.pdf}
\end{sourceblock}
其中系数 S (m) i1i2i3 根据 v12 的分量（即 L 的分量）定义，其中 m = 0

\begin{sourceblock}
\includegraphics[page=271,pagebox=mediabox,trim={99.87bp 355.97bp 67.56bp 308.63bp},clip,max width=0.94\textwidth,max height=0.38\textheight]{Marchisio2013.pdf}
\end{sourceblock}
并且，对于 m = 1, 2, 3，通过

\begin{sourceblock}
\includegraphics[page=271,pagebox=mediabox,trim={76.15bp 272.21bp 67.56bp 393.79bp},clip,max width=0.94\textwidth,max height=0.38\textheight]{Marchisio2013.pdf}
\end{sourceblock}
在这些表达式中，i = i1 + i2 + i3，j = j1 + j2 + i3，k = k1 + k2 + k3，求和运算符定义为7

\begin{sourceblock}
\includegraphics[page=271,pagebox=mediabox,trim={169.21bp 182.66bp 67.56bp 479.80bp},clip,max width=0.94\textwidth,max height=0.38\textheight]{Marchisio2013.pdf}
\end{sourceblock}
对于 i1 = i2 = i3 = 0 的特殊情况，我们定义 S (m) 000 = 0。读者将认识到系数 K(m) ijk 与第 6.1.4 节中讨论的碰撞角度的积分有关。确切的定义是

\begin{sourceblock}
\includegraphics[page=271,pagebox=mediabox,trim={95.20bp 86.59bp 67.56bp 587.52bp},clip,max width=0.94\textwidth,max height=0.38\textheight]{Marchisio2013.pdf}
\end{sourceblock}
\footnotetext[7]{我们利用 L13 = 0 的事实，使得 j3 = i3。}
\begin{sourceblock}
\includegraphics[page=272,pagebox=mediabox,trim={121.70bp 502.32bp 65.07bp 58.94bp},clip,max width=0.94\textwidth,max height=0.38\textheight]{Marchisio2013.pdf}
\end{sourceblock}
这些系数以及方程中的求和。 (6.60)，最容易使用符号数学程序计算。8 表 6.1-6.9 给出了选定矩高达五阶 9 的 I(m) l1l2l3(ω, v, g) 的一些示例。为了清楚起见，在这些表中，我们用 g = (g1, g2, g3) 和 g2 = g2 1 + g2 2 + g2 3 表示速度差向量。 γ = l1 + l2 + l3 阶整数矩的碰撞源项的最终表达式现在可以写成方程的形式： (6.54) 作为

\begin{sourceblock}
\includegraphics[page=272,pagebox=mediabox,trim={78.59bp 350.19bp 65.07bp 281.71bp},clip,max width=0.94\textwidth,max height=0.38\textheight]{Marchisio2013.pdf}
\end{sourceblock}
碰撞通量项为

\begin{sourceblock}
\includegraphics[page=272,pagebox=mediabox,trim={133.62bp 296.61bp 65.07bp 369.86bp},clip,max width=0.94\textwidth,max height=0.38\textheight]{Marchisio2013.pdf}
\end{sourceblock}
读者可以验证零阶矩和一阶矩的源项恰好为零（即质量和平均动量守恒）。有趣的是

\begin{center}\small 表6.1. I(0)\end{center}
l1l2l3(ω, v, g) 项表示二阶矩

\begin{sourceblock}
\includegraphics[page=272,pagebox=mediabox,trim={58.00bp 96.19bp 58.23bp 465.42bp},clip,width=\textwidth,max height=0.72\textheight]{Marchisio2013.pdf}
\end{sourceblock}
8 为了有效地简化结果表达式，用球面角定义 L 的分量可能会有所帮助，如方程 8 所示。 （6.25）。如果正确完成，则 S (m) 的最终表达式
{\small i1i2i3 不应该依赖于球面角，而应该是 v12 分量幂的齐次多项式。\par}
\footnotetext[9]{八节点求积需要这些矩（及其排列）。}
\begin{center}\small 表6.2. I(0)\end{center}
三阶矩的 l1l2l3(ω, v, g) 项

\begin{sourceblock}
\includegraphics[page=273,pagebox=mediabox,trim={58.00bp 235.64bp 58.23bp 75.82bp},clip,width=\textwidth,max height=0.72\textheight]{Marchisio2013.pdf}
\end{sourceblock}
请注意，因为对于 m = 1, 2, 3，I(m) l1l2l3(ω, v1, v1 −v2) 具有 γ + 1 阶 v1 和 v2 分量中的齐次多项式形式，因此方程 1 中的积分(6.70) 可以用高达 γ + 1 阶的整数矩来表示。换句话说，虽然 Gl1l2l3 在 γ 阶矩方面不是封闭的，但其值仅取决于下一个更高阶的整数矩。相反，从方程。 (6.68) 我们可以观察到，由于|v1 −v2|（v1 和v2 的非线性函数）的出现，C(0) l1l2l3 并不闭合。然而，如果该非线性项可以用不依赖于 v1 或 v2 的函数来近似，则剩余项 I(0) l1l2l3(ω, v1, v1 −v2) 将是 γ 阶齐次多项式，因此碰撞源项将是封闭的。这种近似被称为麦克斯韦模型（Maxwell，1879），我们将在稍后更详细地讨论它。最后，我们可以注意到方程。 (6.69) 将根据 γ + 1 阶矩及其空间梯度而闭合。这个结果很有趣，因为使用基于求积的矩方法给定 γ 阶矩，γ + 1 阶矩的闭合通常是相当准确的。使用变量的对称变化，可以直接显示

\clearpage
\begin{center}\small 表6.3. I(0)\end{center}
\begin{sourceblock}
\includegraphics[page=274,pagebox=mediabox,trim={58.00bp 202.03bp 58.23bp 74.82bp},clip,width=\textwidth,max height=0.72\textheight]{Marchisio2013.pdf}
\par\vspace{0.45\baselineskip}
\small 表6.4. I(m)
\end{sourceblock}
l1l2l3(ω, v, g) m = 1, 2, 3 的一阶矩项

\begin{sourceblock}
\includegraphics[page=274,pagebox=mediabox,trim={58.00bp 53.16bp 58.23bp 515.64bp},clip,width=\textwidth,max height=0.72\textheight]{Marchisio2013.pdf}
\par\vspace{0.45\baselineskip}
\small 表6.5. I(1)
\end{sourceblock}
\begin{sourceblock}
\includegraphics[page=275,pagebox=mediabox,trim={58.00bp 420.23bp 58.23bp 74.82bp},clip,width=\textwidth,max height=0.72\textheight]{Marchisio2013.pdf}
\end{sourceblock}
\clearpage
\begin{center}\small 表6.6. I(1)\end{center}
\begin{sourceblock}
\includegraphics[page=275,pagebox=mediabox,trim={58.00bp 53.16bp 58.23bp 297.37bp},clip,width=\textwidth,max height=0.72\textheight]{Marchisio2013.pdf}
\end{sourceblock}
\clearpage
\begin{center}\small 表6.7. I(1)\end{center}
\begin{sourceblock}
\includegraphics[page=276,pagebox=mediabox,trim={58.00bp 248.17bp 58.23bp 74.82bp},clip,width=\textwidth,max height=0.72\textheight]{Marchisio2013.pdf}
\end{sourceblock}
因此，由于表 6.4 中的一阶矩项对于 g 和 -g 是相同的，因此得出 C(1) 100 = C(1) 010 = C(1) 001 = 0 以及方程 1 中的第一个非零贡献。 (6.71) 为 C(1) 200。换句话说，一阶碰撞源项 C(1) 100、C(1) 010 和 C(1) 001 不直接影响质量和平均动量的传递。然而，后者受到G100、G010和G001的影响，这可以解释为“非理想气体”对压力张量的贡献。我们将在 6.4 节中回到这一点。 I(m) l1l2l3 的解析表达式将直接与基于求积的矩方法一起使用，以评估每个整数矩的碰撞源和碰撞通量项。这些术语在正交背景下的数值实现将在第 6.5 节中讨论。

\section*{6.2 多分散硬球碰撞}\phantomsection\addcontentsline{toc}{section}{6.2 多分散硬球碰撞}
{\small 现在，我们将上一节中提出的公式推广到具有不同直径 dα 和 dβ 以及不同质量 mα 的两个粒子之间的二元碰撞的情况\par}
\clearpage
\begin{center}\small 表6.8. I(1)\end{center}
\begin{sourceblock}
\includegraphics[page=277,pagebox=mediabox,trim={58.00bp 324.57bp 58.23bp 74.82bp},clip,width=\textwidth,max height=0.72\textheight]{Marchisio2013.pdf}
\end{sourceblock}
和 mβ 如图 6.2 所示。在矩输运方程中，对于相似粒子之间的碰撞，碰撞项将由 Cαα(ψ) 和 Cββ(ψ) 表示，对于不同粒子之间的碰撞，碰撞项将由 Cαβ(ψ) 和 Cβα(ψ) 表示。例如，对于二元系统，1 型粒子的碰撞项是 C1 = C11 + C12，2 型粒子的碰撞项是 C2 = C21 + C22。按照 Fox \& Vedula (2010) 中介绍的方法，我们将推导出适用于 α 和 β 的任意组合的 Cαβ(ψ) 的通用公式。

\subsection*{6.2.1 任意矩的碰撞项}\phantomsection\addcontentsline{toc}{subsection}{6.2.1 任意矩的碰撞项}
对于任意时刻 ψ，碰撞项定义为 (Fox \& Vedula, 2010)

\begin{sourceblock}
\includegraphics[page=277,pagebox=mediabox,trim={65.03bp 128.55bp 67.56bp 501.12bp},clip,max width=0.94\textwidth,max height=0.38\textheight]{Marchisio2013.pdf}
\end{sourceblock}
\begin{sourceblock}
\includegraphics[page=277,pagebox=mediabox,trim={205.33bp 83.49bp 67.56bp 588.94bp},clip,max width=0.94\textwidth,max height=0.38\textheight]{Marchisio2013.pdf}
\end{sourceblock}
{\small 其中 eαβ = eβα 是 α 和 β 类型粒子之间碰撞的恢复系数。请注意，下标 1 和 2 与之前一样对应于碰撞伙伴，并且\par}
\clearpage
\begin{center}\small 表6.9. I(1)\end{center}
\begin{sourceblock}
\includegraphics[page=278,pagebox=mediabox,trim={58.00bp 53.16bp 58.23bp 74.82bp},clip,width=\textwidth,max height=0.72\textheight]{Marchisio2013.pdf}
\end{sourceblock}
\begin{sourceblock}
\includegraphics[page=279,pagebox=mediabox,trim={58.00bp 478.16bp 58.23bp 59.00bp},clip,width=\textwidth,max height=0.72\textheight]{Marchisio2013.pdf}
\par\vspace{0.45\baselineskip}
\small 图6.2。两个不同直径 d1 和 d2 的硬球之间的直接碰撞。
\end{sourceblock}
下标 α 和 β 区分粒子属性。配对分布函数 f (2)(x, mα, v1; x + dαβx12, mβ, v2) 被编写为使得碰撞伙伴 1 的质量为 mα，碰撞伙伴 2 的质量为 mβ。正如第 6.1.2 节中所做的那样，我们现在可以使用关于碰撞接触点的泰勒级数展开式

\begin{sourceblock}
\includegraphics[page=279,pagebox=mediabox,trim={173.29bp 351.40bp 180.12bp 323.22bp},clip,max width=0.94\textwidth,max height=0.38\textheight]{Marchisio2013.pdf}
\end{sourceblock}
重新表达方程。 （6.73）就碰撞源项和碰撞通量项而言：

\begin{sourceblock}
\includegraphics[page=279,pagebox=mediabox,trim={181.08bp 315.52bp 67.56bp 371.18bp},clip,max width=0.94\textwidth,max height=0.38\textheight]{Marchisio2013.pdf}
\end{sourceblock}
其中碰撞源项由下式给出

\begin{sourceblock}
\includegraphics[page=279,pagebox=mediabox,trim={77.12bp 236.57bp 67.56bp 402.50bp},clip,max width=0.94\textwidth,max height=0.38\textheight]{Marchisio2013.pdf}
\end{sourceblock}
碰撞通量项由下式给出

\begin{sourceblock}
\includegraphics[page=279,pagebox=mediabox,trim={75.58bp 158.01bp 67.56bp 481.06bp},clip,max width=0.94\textwidth,max height=0.38\textheight]{Marchisio2013.pdf}
\end{sourceblock}
通过引入对相关函数 gαβ 并使用围绕点 x 的截断泰勒级数展开到 dαβ 中的一阶，可以进一步简化该表达式（请参阅第 6.1.2 节中给出的详细信息）：

\begin{sourceblock}
\includegraphics[page=279,pagebox=mediabox,trim={76.48bp 55.60bp 67.56bp 590.93bp},clip,max width=0.94\textwidth,max height=0.38\textheight]{Marchisio2013.pdf}
\end{sourceblock}
请注意，我们假设当 α β 时，gαβ 可能依赖于 x12，并将在下面讨论可能的依赖关系。我们还引入了每种粒子类型（fα 和 fβ）的速度分布函数，它们不必相同（除非 α = β）。一般来说，gαβ 取决于每种颗粒类型的体积分数和颗粒直径。然而，它也可以取决于速度分布函数的其他矩。例如，如果平均粒子速度 Uα 和 Uβ 非常不同，则可以预期较慢粒子类型的上游侧的碰撞频率会更高。单位向量 x12 表示碰撞时粒子中心的相对位置。如果我们考虑相对于平均速度差 (Uα −Uβ) · x12 的方向，我们可以将配对相关函数对平均速度差的依赖性建模为 10

\begin{sourceblock}
\includegraphics[page=280,pagebox=mediabox,trim={161.40bp 463.69bp 65.07bp 196.21bp},clip,max width=0.94\textwidth,max height=0.38\textheight]{Marchisio2013.pdf}
\end{sourceblock}
其中 g0,αβ 仅取决于体积分数和颗粒直径，cαβ 是模型常数，必须选择使 gαβ(x12) 为非负值11（即 |cαβ| ≤1）。能量比例因子 Eαβ 可以定义为

\begin{sourceblock}
\includegraphics[page=280,pagebox=mediabox,trim={184.35bp 392.70bp 65.07bp 292.53bp},clip,max width=0.94\textwidth,max height=0.38\textheight]{Marchisio2013.pdf}
\end{sourceblock}
其中 θα 和 θβ 分别是颗粒类型 α 和 β 的颗粒温度。请注意等式。 (6.80) 对应于速度分布为麦克斯韦分布时的预期平方速度差。读者应该注意，方程中给出的形式。 （6.79）可以推广到包括矩的空间梯度。例如，我们可能有

\begin{sourceblock}
\includegraphics[page=280,pagebox=mediabox,trim={159.42bp 308.11bp 65.07bp 369.35bp},clip,max width=0.94\textwidth,max height=0.38\textheight]{Marchisio2013.pdf}
\end{sourceblock}
其中 H 是取决于矩的函数（例如，在流体力学极限中，H 仅取决于分散相体积分数、平均速度和颗粒温度）。为了保证正性，参数καβ必须满足|καβ| < 1/|dαβ ∇xH|。方程中的模型。 (6.81) 可用于描述极限 dα <dβ 下的浮力，在这种情况下，H 将取决于压力 pα。事实上，除了阻力项之外，第 5 章中描述的所有相动量交换项都将由流体-颗粒极限中 gαβ(x12) 的函数形式产生（即 dα ≪dβ）。 12 因此，在对具有不同颗粒尺寸的多相系统进行建模时，应注意开发 gαβ(x12) 的物理精确表达式。在下面的讨论中，我们将使用方程式。 （6.79）推导出碰撞项。尽管如此，修改方程 1 中形式的结果表达式还是很简单的。 (6.81) 因为在这两种情况下对 x12 的依赖性是相同的。

10 该模型应解释为 (Uα −Uβ) · x12 中截断泰勒级数展开式的前两项。 gαβ(x12) 对 (Uα −Uβ)·x12 的完全依赖性可能要复杂得多。 11 条件|cαβ| ≤1 由方程 1 中给出的 Eαβ 的定义得出。 （6.80）。因为x12 是单位向量，所以(Uα -Uβ) · x12 的最大值/最小值为±|Uα -Uβ|。 E1/2 αβ 的最小值为 |Uα −Uβ|当颗粒温度为零时发生。因此，gαβ 的符号由 1 ± cαβ 确定，从而得出 cαβ 的界限。

\footnotetext[12]{此处，dα 为分子直径，dβ 为粒径。在流体-粒子极限下，dα趋于零，而分子数趋于无穷大，使得分子所占的体积分数恒定。请注意，在此极限下，流体相 NDF fα 可以视为流体动力学极限中的近麦克斯韦分布。}
下一步是使用定义 v′ 1 = v1 − ωαβ(v12 · x12)x12 将 ψ(v′ 1) 写为 x12 的显式函数。在采用关于 v1 的泰勒级数展开时，我们发现

\begin{sourceblock}
\includegraphics[page=281,pagebox=mediabox,trim={133.22bp 574.30bp 67.56bp 92.65bp},clip,max width=0.94\textwidth,max height=0.38\textheight]{Marchisio2013.pdf}
\end{sourceblock}
然后我们可以将碰撞源项表示为

\begin{sourceblock}
\includegraphics[page=281,pagebox=mediabox,trim={151.67bp 531.24bp 67.56bp 152.18bp},clip,max width=0.94\textwidth,max height=0.38\textheight]{Marchisio2013.pdf}
\end{sourceblock}
和碰撞通量项为

\begin{sourceblock}
\includegraphics[page=281,pagebox=mediabox,trim={188.27bp 488.08bp 67.56bp 195.33bp},clip,max width=0.94\textwidth,max height=0.38\textheight]{Marchisio2013.pdf}
\end{sourceblock}
其中

\begin{sourceblock}
\includegraphics[page=281,pagebox=mediabox,trim={77.39bp 386.26bp 67.56bp 236.97bp},clip,max width=0.94\textwidth,max height=0.38\textheight]{Marchisio2013.pdf}
\end{sourceblock}
\begin{sourceblock}
\includegraphics[page=281,pagebox=mediabox,trim={77.39bp 299.98bp 67.56bp 321.83bp},clip,max width=0.94\textwidth,max height=0.38\textheight]{Marchisio2013.pdf}
\end{sourceblock}
\begin{sourceblock}
\includegraphics[page=281,pagebox=mediabox,trim={77.39bp 228.57bp 67.56bp 409.21bp},clip,max width=0.94\textwidth,max height=0.38\textheight]{Marchisio2013.pdf}
\end{sourceblock}
\begin{sourceblock}
\includegraphics[page=281,pagebox=mediabox,trim={77.39bp 141.70bp 67.56bp 479.11bp},clip,max width=0.94\textwidth,max height=0.38\textheight]{Marchisio2013.pdf}
\end{sourceblock}
以及

\begin{sourceblock}
\includegraphics[page=281,pagebox=mediabox,trim={75.28bp 54.62bp 67.56bp 584.15bp},clip,max width=0.94\textwidth,max height=0.38\textheight]{Marchisio2013.pdf}
\end{sourceblock}
\begin{sourceblock}
\includegraphics[page=282,pagebox=mediabox,trim={80.12bp 568.44bp 65.07bp 58.29bp},clip,max width=0.94\textwidth,max height=0.38\textheight]{Marchisio2013.pdf}
\end{sourceblock}
从这些定义中，我们可以观察到，除了方程中涉及 x12 的两类积分外， (6.23)，当 α β 时，出现第三种类型：

\begin{sourceblock}
\includegraphics[page=282,pagebox=mediabox,trim={155.67bp 505.41bp 65.07bp 161.05bp},clip,max width=0.94\textwidth,max height=0.38\textheight]{Marchisio2013.pdf}
\end{sourceblock}
其中涉及张量积 x12 ⊗x12。下面，我们将考虑整数矩的积分值。

\subsection*{6.2.2 碰撞角的三次积分}\phantomsection\addcontentsline{toc}{subsection}{6.2.2 碰撞角的三次积分}
第三个积分，如方程式所示。 （6.91），可以逐个评估： ∫

\begin{sourceblock}
\includegraphics[page=282,pagebox=mediabox,trim={73.79bp 280.75bp 65.07bp 280.18bp},clip,max width=0.94\textwidth,max height=0.38\textheight]{Marchisio2013.pdf}
\end{sourceblock}
并涉及九组不同的碰撞角度积分。这些积分得出六个系数 K(mn) ijk = K(nm) ijk，其定义为

\begin{sourceblock}
\includegraphics[page=282,pagebox=mediabox,trim={83.58bp 56.53bp 65.07bp 458.94bp},clip,max width=0.94\textwidth,max height=0.38\textheight]{Marchisio2013.pdf}
\end{sourceblock}
\begin{sourceblock}
\includegraphics[page=283,pagebox=mediabox,trim={81.09bp 555.34bp 67.56bp 58.18bp},clip,max width=0.94\textwidth,max height=0.38\textheight]{Marchisio2013.pdf}
\end{sourceblock}
我们剩下的就是 beta 函数。请注意，在方程式中的积分中。根据方程（6.92），这些系数与 L 的分量相关联，使得 K(mn) ijk LmpLnq 出现在最终结果中。另请注意方程中的 1/π 因子。包含 (6.93)–(6.98) 是为了消除 beta 函数中出现的 π 相关项，从而使六个系数都是有理数。当我们代入下面的定义时，相同的因素将从碰撞术语中删除。

\subsection*{6.2.3 整数矩的碰撞项}\phantomsection\addcontentsline{toc}{subsection}{6.2.3 整数矩的碰撞项}
对于 γ = l1 + l2 + l3 阶的整数矩，方程中的碰撞项： (6.88) 和 (6.90) 可以写成

\begin{sourceblock}
\includegraphics[page=283,pagebox=mediabox,trim={95.39bp 377.37bp 67.56bp 296.61bp},clip,max width=0.94\textwidth,max height=0.38\textheight]{Marchisio2013.pdf}
\end{sourceblock}
其中系数 S (p,q) i1i2i3 = S (q,p) i1i2i3 根据 p、q = 1, 2, 3 的 g 分量定义：

\begin{sourceblock}
\includegraphics[page=283,pagebox=mediabox,trim={65.97bp 311.57bp 69.48bp 369.73bp},clip,max width=0.94\textwidth,max height=0.38\textheight]{Marchisio2013.pdf}
\end{sourceblock}
\begin{sourceblock}
\includegraphics[page=283,pagebox=mediabox,trim={99.90bp 277.34bp 67.57bp 399.22bp},clip,max width=0.94\textwidth,max height=0.38\textheight]{Marchisio2013.pdf}
\end{sourceblock}
对于 i1 = i2 = i3 = 0 的特殊情况，我们定义 S (p,q) 000 = 0。请注意，在 S (p,q) i1i2i3 的定义中，g 的大小为 i + 2 次幂。这样做是为了使 I(p,q) l1l2l3(ω, v, g) 就 v 和 g 的分量而言将是 γ + 2 阶齐次多项式。表 6.10-6.12 给出了 I(p,q) l1l2l3 的示例。我们现在可以将 γ 阶整数矩的碰撞源项和碰撞通量项表示为

\begin{sourceblock}
\includegraphics[page=283,pagebox=mediabox,trim={143.65bp 165.20bp 67.57bp 517.25bp},clip,max width=0.94\textwidth,max height=0.38\textheight]{Marchisio2013.pdf}
\end{sourceblock}
以及

\begin{sourceblock}
\includegraphics[page=283,pagebox=mediabox,trim={183.11bp 133.34bp 67.57bp 549.10bp},clip,max width=0.94\textwidth,max height=0.38\textheight]{Marchisio2013.pdf}
\end{sourceblock}
在等式中。 （6.101），右侧的四个贡献被定义为13

\begin{sourceblock}
\includegraphics[page=283,pagebox=mediabox,trim={75.87bp 81.40bp 67.57bp 585.06bp},clip,max width=0.94\textwidth,max height=0.38\textheight]{Marchisio2013.pdf}
\end{sourceblock}
\footnotetext[13]{由于不再需要跟踪碰撞伙伴，我们现在转而将积分变量表示为 vα 和 vβ。}
\clearpage
\begin{center}\small 表6.10。 I(p,q) 的示例\end{center}
\begin{sourceblock}
\includegraphics[page=284,pagebox=mediabox,trim={58.00bp 53.16bp 58.23bp 75.04bp},clip,width=\textwidth,max height=0.72\textheight]{Marchisio2013.pdf}
\end{sourceblock}
\clearpage
\begin{center}\small 表6.11。 I(p,q) 的示例\end{center}
\begin{sourceblock}
\includegraphics[page=285,pagebox=mediabox,trim={58.00bp 53.16bp 58.23bp 75.04bp},clip,width=\textwidth,max height=0.72\textheight]{Marchisio2013.pdf}
\par\vspace{0.45\baselineskip}
\small 表6.12。 I(p,q) 的示例
\end{sourceblock}
\begin{sourceblock}
\includegraphics[page=286,pagebox=mediabox,trim={58.00bp 412.77bp 58.23bp 75.04bp},clip,width=\textwidth,max height=0.72\textheight]{Marchisio2013.pdf}
\end{sourceblock}
在第 6.4 节中，我们将应用这些结果来找到方程的多分散版本的矩输运方程。 （6.1）。

\section*{6.3 动力学模型}\phantomsection\addcontentsline{toc}{section}{6.3 动力学模型}
由于 Boltzmann-Enskog 碰撞积分的形式，前面几节中介绍的积分相当复杂，这导致碰撞源项的表达式相当复杂（尤其是高阶矩）。在文献中，大量工作致力于寻找更简单的动力学模型来近似玻尔兹曼碰撞积分（即，对于忽略碰撞通量项的稀释情况）。对于单分散颗粒，最广泛使用的动力学模型是所谓的 BGK 模型，该模型由 Bhatnagar 等人独立提出。 （1954）和韦兰德（1954）。由于涉及不同粒子之间碰撞的“交叉碰撞”项，将 BGK 类模型扩展到多分散粒子更加复杂（参见 McCormack (1973) 和 Andries 等人 (2002a) 以及其中引用的参考文献）。事实上，正如 Garz ́o 等人首先讨论的那样。 (1989)，几乎没有任何多分散动力学模型遵循“不可微性原理”(Andries et al., 2002a)，该原理指出，当所有粒子相同时，多分散动力学模型应简化为单分散模型。在本节中，我们将仅介绍最简单的非弹性碰撞动力学模型。读者可以在动力学理论文献中找到对动力学模型更详尽的描述（Cercignani，2000；Struchtrup，2005）。

\subsection*{6.3.1 单分散颗粒}\phantomsection\addcontentsline{toc}{subsection}{6.3.1 单分散颗粒}
对于弹性碰撞，已经提出了几种不同的动力学模型，以封闭玻尔兹曼硬球碰撞项（方程 6.9）。对于非弹性碰撞 (e < 1)，必须正确解释颗粒能量耗散对 e 值的依赖性。完成此任务的一种方法是从方程（1）中的精确（未封闭）碰撞积分开始。 （6.68）。根据 I(0) 的定义
{\small ijk 在方程中给出。 （6.60），可得\par}
表明矩 Mγ ijk 的近似碰撞源项，定义为

\begin{sourceblock}
\includegraphics[page=287,pagebox=mediabox,trim={135.55bp 592.02bp 67.57bp 74.43bp},clip,max width=0.94\textwidth,max height=0.38\textheight]{Marchisio2013.pdf}
\end{sourceblock}
就高达 γ 阶的矩而言是封闭的。该表达式与等式不同。 (6.68) 因为我们替换了 |v1 −v2|平均速度差定义为

\begin{sourceblock}
\includegraphics[page=287,pagebox=mediabox,trim={163.57bp 529.73bp 67.57bp 135.79bp},clip,max width=0.94\textwidth,max height=0.38\textheight]{Marchisio2013.pdf}
\end{sourceblock}
在平衡状态下，速度分布为麦克斯韦分布，很容易证明 〈g〉= 4 θp/π，其中 θp 是颗粒温度。我们应该注意到方程。 （6.109）对应于非弹性麦克斯韦粒子（Maxwell，1879），最重要的是，它仍然包含对 ω = (1 + e)/2 的精确依赖。因此，我们将该动力学模型称为非弹性麦克斯韦碰撞模型。单分散颗粒的类 BGK 动力学模型具有线性化形式

\begin{sourceblock}
\includegraphics[page=287,pagebox=mediabox,trim={209.47bp 420.69bp 67.57bp 253.14bp},clip,max width=0.94\textwidth,max height=0.38\textheight]{Marchisio2013.pdf}
\end{sourceblock}
其中（重复的罗马索引意味着求和）“均衡”分布函数建模为

\begin{sourceblock}
\includegraphics[page=287,pagebox=mediabox,trim={232.90bp 365.15bp 67.57bp 304.37bp},clip,max width=0.94\textwidth,max height=0.38\textheight]{Marchisio2013.pdf}
\end{sourceblock}
f ∗(v) = αp [det(2πλ)]1/2 exp
在这个多元高斯分布中，Up 是平均粒子速度，λ−1 是二阶张量 λ 的倒数，其定义使得给定的速度矩集与方程 1 一致。 （6.109）。由于质量和平均动量守恒，方程中的第一个非零项。 (6.109) 对应于二阶矩。让 κ 表示由 I(0) 构造的对称二阶张量

\begin{sourceblock}
\includegraphics[page=287,pagebox=mediabox,trim={65.03bp 250.65bp 67.55bp 393.08bp},clip,max width=0.94\textwidth,max height=0.38\textheight]{Marchisio2013.pdf}
\end{sourceblock}
从方程中找到的二阶矩的碰撞项。 (6.109) 由下式给出

\begin{sourceblock}
\includegraphics[page=287,pagebox=mediabox,trim={143.96bp 200.60bp 67.57bp 465.86bp},clip,max width=0.94\textwidth,max height=0.38\textheight]{Marchisio2013.pdf}
\end{sourceblock}
该表达式右侧的积分可以显式求值：

\begin{sourceblock}
\includegraphics[page=287,pagebox=mediabox,trim={163.19bp 148.17bp 67.57bp 517.98bp},clip,max width=0.94\textwidth,max height=0.38\textheight]{Marchisio2013.pdf}
\end{sourceblock}
其中 σp 是速度协方差张量，θp = tr(σp)/3。使用 ω(ω −1) = (e2 −1)/4，可以直接证明对于非弹性麦克斯韦碰撞模型，均匀颗粒温度服从

\begin{sourceblock}
\includegraphics[page=287,pagebox=mediabox,trim={65.02bp 52.15bp 67.57bp 599.11bp},clip,max width=0.94\textwidth,max height=0.38\textheight]{Marchisio2013.pdf}
\end{sourceblock}
从方程中找到的二阶矩的碰撞项。 (6.111) 是

\begin{sourceblock}
\includegraphics[page=288,pagebox=mediabox,trim={207.24bp 596.24bp 65.08bp 74.72bp},clip,max width=0.94\textwidth,max height=0.38\textheight]{Marchisio2013.pdf}
\end{sourceblock}
将该表达式等同于等式。 (6.115) 并求解未知模型参数，我们发现

\begin{sourceblock}
\includegraphics[page=288,pagebox=mediabox,trim={177.03bp 517.43bp 65.08bp 127.79bp},clip,max width=0.94\textwidth,max height=0.38\textheight]{Marchisio2013.pdf}
\end{sourceblock}
以及

其中 0 < ze ≤3/2 是模型常数。在弹性极限（ω = 1）下，方程： (6.119) 对应于 Prandtl 数由 Pr = 1/ze 给出的 ES-BGK 模型，以及 ze = 1 时的 BGK 模型。对于非弹性碰撞，我们将非弹性 BGK 和 ES-BGK 模型称为方程 (6.119) 中给出的线性化模型。 (6.111) 与矩阵方程。 （6.119）使用弹性极限中的 z 值。唯一剩下的未封闭项是〈g〉。在弹性极限平衡时，g 将呈高斯分布，均值为零，协方差矩阵为 2θpI。这一事实导致近似值 〈g〉= 4 θp/π，因此我们可以将碰撞时间定义为 τcol = ze √πdp/

\begin{sourceblock}
\includegraphics[page=288,pagebox=mediabox,trim={152.51bp 414.84bp 286.21bp 262.95bp},clip,max width=0.94\textwidth,max height=0.38\textheight]{Marchisio2013.pdf}
\end{sourceblock}
。在评估矩的碰撞源项时，动力学模型的实用性最为明显。如前所述，等式。 (6.109) 是封闭的，但会导致高阶矩的复杂多项式表达式。相反，动力学模型得出

\begin{sourceblock}
\includegraphics[page=288,pagebox=mediabox,trim={195.68bp 348.18bp 65.08bp 327.34bp},clip,max width=0.94\textwidth,max height=0.38\textheight]{Marchisio2013.pdf}
\end{sourceblock}
其中矩 M*γ ijk 可以使用方程式中 f* 的定义进行分析计算。 （6.112）：

\begin{sourceblock}
\includegraphics[page=288,pagebox=mediabox,trim={195.20bp 292.57bp 65.08bp 373.88bp},clip,max width=0.94\textwidth,max height=0.38\textheight]{Marchisio2013.pdf}
\end{sourceblock}
由于 f 是多元高斯分布，因此可以轻松计算其高阶矩（例如使用矩生成函数）。然而，读者应该记住，动力学模型仅确保二阶以下的矩与方程 1 相同。 （6.109）。因此，当真实的速度分布函数远离平衡时，三阶和高阶矩可能很难近似。

\subsection*{6.3.2 多分散颗粒}\phantomsection\addcontentsline{toc}{subsection}{6.3.2 多分散颗粒}
如前所述，已经有许多尝试将类似 BGK 的动力学模型扩展到多组分混合物（例如 Andries 等人，2002a；Garz ́o 等人，1989；Goldman 和 Sirovich，1967；Gross 和 Krook，1956；Holway，1973；Kosuge，2009；McCormack，1973；以及西罗维奇，1962）。然而，只有 Andries 等人的模型。 (2002a) 满足多分量玻尔兹曼碰撞项的关键数学性质（即非负分布函数、不可微性原理、混合平衡、H 定理）。一般来说，没有一个动力学模型在所有情况下都能产生与玻尔兹曼碰撞项相当的结果（Andries 等人，2002b；Kosuge，2009；Struchtrup，2005）。然而，由于它们在实践中使用起来更加简单，因此它们可以为某些应用提供有用的近似值。

{\small 对于非弹性碰撞，恢复系数将明确出现在动力学模型中，如上面单分散情况所示。我们将考虑具有非弹性的二元情况\par}
为了简单起见，碰撞。然而，动力学模型将被编写为可以轻松扩展到两种以上的粒子类型。我们将用碰撞产生的矩源项来表达模型：

\begin{sourceblock}
\includegraphics[page=289,pagebox=mediabox,trim={132.87bp 564.44bp 67.57bp 94.36bp},clip,max width=0.94\textwidth,max height=0.38\textheight]{Marchisio2013.pdf}
\end{sourceblock}
其中 C*γ ijk,1 (C*γ ijk,2) 是对应于矩 Mγ ijk,1 (Mγ ijk,2) 的 1 类（2 类）粒子的碰撞源项。选择模型中的参数是为了给出麦克斯韦分子高达二阶矩的正确行为（如单分散情况所做的那样）。术语 Gijk,αβ 表示具有平均速度 Uαβ 和协方差矩阵 σαβ（定义如下）的高斯矩。碰撞率参数 καβ 是非负的，并确定向高斯分布的松弛率：

\begin{sourceblock}
\includegraphics[page=289,pagebox=mediabox,trim={145.20bp 405.68bp 67.57bp 227.61bp},clip,max width=0.94\textwidth,max height=0.38\textheight]{Marchisio2013.pdf}
\end{sourceblock}
其中 θ12 = E12/3，d12 = (d1 + d2)/2，χαβ = dαβ/dβ。高斯分布中的平均速度定义为

\begin{sourceblock}
\includegraphics[page=289,pagebox=mediabox,trim={146.68bp 323.39bp 67.57bp 327.00bp},clip,max width=0.94\textwidth,max height=0.38\textheight]{Marchisio2013.pdf}
\end{sourceblock}
其中

\begin{sourceblock}
\includegraphics[page=289,pagebox=mediabox,trim={212.31bp 284.63bp 67.57bp 388.05bp},clip,max width=0.94\textwidth,max height=0.38\textheight]{Marchisio2013.pdf}
\end{sourceblock}
协方差矩阵可以表示为以下形式

\begin{sourceblock}
\includegraphics[page=289,pagebox=mediabox,trim={142.80bp 153.73bp 67.57bp 428.93bp},clip,max width=0.94\textwidth,max height=0.38\textheight]{Marchisio2013.pdf}
\end{sourceblock}
其中 σα 是 α 型粒子的协方差矩阵，右侧出现的其他矩阵定义为

\begin{sourceblock}
\includegraphics[page=289,pagebox=mediabox,trim={105.84bp 92.40bp 67.57bp 582.21bp},clip,max width=0.94\textwidth,max height=0.38\textheight]{Marchisio2013.pdf}
\end{sourceblock}
{\small 请注意，为了很好地定义高斯分布，方程中的所有协方差矩阵。 (6.126) 必须是非负的。这导致条件 4 ≥(1 + e12)μ21 ≥0，或者等效地，2 ≥μ21 ≥0，这始终为真。\par}
通过设置 α2 = 0，可以将多分散动力学模型与单分散情况进行比较。很容易证明

\begin{sourceblock}
\includegraphics[page=290,pagebox=mediabox,trim={167.25bp 581.21bp 65.08bp 85.73bp},clip,max width=0.94\textwidth,max height=0.38\textheight]{Marchisio2013.pdf}
\end{sourceblock}
对应于等式中的 z = 1/2。 (6.119)。单分散动力学模型中的弛豫率为

\begin{sourceblock}
\includegraphics[page=290,pagebox=mediabox,trim={224.82bp 523.37bp 203.47bp 140.25bp},clip,max width=0.94\textwidth,max height=0.38\textheight]{Marchisio2013.pdf}
\end{sourceblock}
1 τcol
这与 κ11 一致。因此，当仅存在一种颗粒类型时，多分散动力学模型简化为单分散 ES-BGK 模型（Holway，1966，1973）。总之，我们提醒读者，动力学模型提供了硬球碰撞项的简化描述，精确到二阶矩。目前，动力学模型仅适用于碰撞通量项可忽略不计的稀释状态。据作者所知，还没有尝试构建碰撞通量项的动力学模型，这对于中等密度的流动很重要。

\section*{6.4 矩输运方程}\phantomsection\addcontentsline{toc}{section}{6.4 矩输运方程}
现在我们将研究两个系统的整数矩输运方程。第一个对应于具有由方程给出的动力学方程的单分散情况。 （6.1）。第二个是具有动力学方程的双分散系统

\begin{sourceblock}
\includegraphics[page=290,pagebox=mediabox,trim={191.04bp 319.71bp 65.08bp 351.79bp},clip,max width=0.94\textwidth,max height=0.38\textheight]{Marchisio2013.pdf}
\end{sourceblock}
速度数密度函数 f1(t, x, v1) 和 f2(t, x, v2) 定义为

\begin{sourceblock}
\includegraphics[page=290,pagebox=mediabox,trim={166.64bp 267.17bp 164.19bp 399.29bp},clip,max width=0.94\textwidth,max height=0.38\textheight]{Marchisio2013.pdf}
\end{sourceblock}
其中 α1 和 α2 是分散相体积分数，并且

\begin{sourceblock}
\includegraphics[page=290,pagebox=mediabox,trim={146.50bp 215.82bp 144.05bp 450.63bp},clip,max width=0.94\textwidth,max height=0.38\textheight]{Marchisio2013.pdf}
\end{sourceblock}
其中 U1 和 U2 是平均分散相速度。请注意，碰撞期间的质量守恒意味着

\begin{sourceblock}
\includegraphics[page=290,pagebox=mediabox,trim={183.28bp 152.53bp 180.83bp 513.93bp},clip,max width=0.94\textwidth,max height=0.38\textheight]{Marchisio2013.pdf}
\end{sourceblock}
而动量守恒意味着

\begin{sourceblock}
\includegraphics[page=290,pagebox=mediabox,trim={185.09bp 101.18bp 182.65bp 565.28bp},clip,max width=0.94\textwidth,max height=0.38\textheight]{Marchisio2013.pdf}
\end{sourceblock}
\begin{sourceblock}
\includegraphics[page=290,pagebox=mediabox,trim={166.21bp 60.91bp 163.77bp 605.55bp},clip,max width=0.94\textwidth,max height=0.38\textheight]{Marchisio2013.pdf}
\end{sourceblock}
{\small 以及\par}
其中 ρ1 和 ρ2 是颗粒材料密度（即颗粒质量除以颗粒体积）。速度密度函数的整数矩定义为

\begin{sourceblock}
\includegraphics[page=291,pagebox=mediabox,trim={156.61bp 565.32bp 67.57bp 101.15bp},clip,max width=0.94\textwidth,max height=0.38\textheight]{Marchisio2013.pdf}
\end{sourceblock}
其中 γ = l1 + l2 + l3 是矩的阶数。一般来说，二阶矩与“流体动力学”变量相关。例如，

\begin{sourceblock}
\includegraphics[page=291,pagebox=mediabox,trim={143.50bp 471.00bp 135.58bp 181.26bp},clip,max width=0.94\textwidth,max height=0.38\textheight]{Marchisio2013.pdf}
\end{sourceblock}
其中 θα 是颗粒类型 α 的颗粒温度。因此，我们将最仔细地研究直到 γ = 2 的矩输运方程。对于单分散情况，我们分别将流体动力学变量表示为 αp、Up 和 θp，并将矩表示为 Mγ l1l2l3。

\subsection*{6.4.1 单分散颗粒}\phantomsection\addcontentsline{toc}{subsection}{6.4.1 单分散颗粒}
对于单分散系统，矩输运方程由方程导出。 (6.1) 是

\begin{sourceblock}
\includegraphics[page=291,pagebox=mediabox,trim={76.18bp 295.28bp 67.56bp 355.02bp},clip,max width=0.94\textwidth,max height=0.38\textheight]{Marchisio2013.pdf}
\end{sourceblock}
在该方程的左侧，碰撞通量由方程定义。 (6.70)，出现。在右侧，碰撞源项由方程式定义。出现 (6.68) 和 (6.69)。动通量（或自由传输）对应于矩 Mγ+1。

\paragraph*{零阶矩}
对于零阶矩，方程： (6.131) 减少为

\begin{sourceblock}
\includegraphics[page=291,pagebox=mediabox,trim={137.58bp 156.92bp 67.57bp 514.05bp},clip,max width=0.94\textwidth,max height=0.38\textheight]{Marchisio2013.pdf}
\end{sourceblock}
可以根据流体动力学变量将其重写为连续性方程：

\begin{sourceblock}
\includegraphics[page=291,pagebox=mediabox,trim={200.95bp 97.94bp 67.57bp 573.91bp},clip,max width=0.94\textwidth,max height=0.38\textheight]{Marchisio2013.pdf}
\end{sourceblock}
{\small 请注意，该方程在流体动力学变量方面是封闭的，并且对应于可压缩颗粒气体，其中 αp 扮演密度的角色。\par}
\paragraph*{一阶矩}
对于一阶矩，方程。 (6.131) 减少为

\begin{sourceblock}
\includegraphics[page=292,pagebox=mediabox,trim={76.32bp 514.49bp 65.08bp 99.32bp},clip,max width=0.94\textwidth,max height=0.38\textheight]{Marchisio2013.pdf}
\end{sourceblock}
可以用流体动力学变量重写为平均动量平衡：

\begin{sourceblock}
\includegraphics[page=292,pagebox=mediabox,trim={178.46bp 462.18bp 65.08bp 209.68bp},clip,max width=0.94\textwidth,max height=0.38\textheight]{Marchisio2013.pdf}
\end{sourceblock}
对称二阶压力张量 P 由其分量定义：

\begin{sourceblock}
\includegraphics[page=292,pagebox=mediabox,trim={176.35bp 331.16bp 65.08bp 252.65bp},clip,max width=0.94\textwidth,max height=0.38\textheight]{Marchisio2013.pdf}
\end{sourceblock}
右侧括号中的项可以用速度协方差矩阵 σp 表示，定义为

\begin{sourceblock}
\includegraphics[page=292,pagebox=mediabox,trim={161.41bp 264.75bp 65.08bp 401.71bp},clip,max width=0.94\textwidth,max height=0.38\textheight]{Marchisio2013.pdf}
\end{sourceblock}
使用方程式中碰撞通量的定义。 (6.70) 和表 6.4 中的结果，我们可以将 P 重写为

\begin{sourceblock}
\includegraphics[page=292,pagebox=mediabox,trim={81.41bp 202.48bp 90.00bp 463.97bp},clip,max width=0.94\textwidth,max height=0.38\textheight]{Marchisio2013.pdf}
\end{sourceblock}
\vspace{-1.5em}\hfill (6.138)\par
(6.138)。利用恒等式
\begin{sourceblock}
\includegraphics[page=292,pagebox=mediabox,trim={100.00bp 138.00bp 45.00bp 525.00bp},clip,max width=0.94\textwidth,max height=0.38\textheight]{Marchisio2013.pdf}
\end{sourceblock}

以及

\begin{sourceblock}
\includegraphics[page=292,pagebox=mediabox,trim={80.00bp 94.00bp 45.00bp 587.00bp},clip,max width=0.94\textwidth,max height=0.38\textheight]{Marchisio2013.pdf}
\end{sourceblock}

{\small 事实是\par}
\begin{sourceblock}
\includegraphics[page=292,pagebox=mediabox,trim={150.89bp 50.16bp 148.44bp 617.50bp},clip,max width=0.94\textwidth,max height=0.38\textheight]{Marchisio2013.pdf}
\end{sourceblock}
我们可以计算方程中的积分。 (6.138) 求

\begin{sourceblock}
\includegraphics[page=293,pagebox=mediabox,trim={163.96bp 593.66bp 67.57bp 82.19bp},clip,max width=0.94\textwidth,max height=0.38\textheight]{Marchisio2013.pdf}
\end{sourceblock}
\begin{sourceblock}
\includegraphics[page=293,pagebox=mediabox,trim={295.49bp 562.35bp 112.27bp 115.11bp},clip,max width=0.94\textwidth,max height=0.38\textheight]{Marchisio2013.pdf}
\end{sourceblock}
其中颗粒温度由 θp = 1 3 定义。请注意等式。 （6.139）仅取决于零阶矩和二阶矩。通常将 P 重写为分解式 P = τp + pI，它定义了颗粒压力

\begin{sourceblock}
\includegraphics[page=293,pagebox=mediabox,trim={186.25bp 502.66bp 67.57bp 175.14bp},clip,max width=0.94\textwidth,max height=0.38\textheight]{Marchisio2013.pdf}
\end{sourceblock}
和一个无痕颗粒应力张量：

\begin{sourceblock}
\includegraphics[page=293,pagebox=mediabox,trim={166.19bp 441.18bp 67.57bp 234.68bp},clip,max width=0.94\textwidth,max height=0.38\textheight]{Marchisio2013.pdf}
\end{sourceblock}
从方程。 (6.140) 我们可以注意到，由于 g0 的存在，粒状“状态方程”是 αp 的高度非线性函数。在稠密颗粒流中，这种“非理想气体”行为使得方程中的平均动量方程成立。 （6.135）非常刚性并且与方程中的连续性方程紧密耦合。 (6.133)。我们还可以注意到，在稀释状态（αp<1）中，方程中涉及 g0 的项。 (6.140) 和 (6.141) 可以忽略不计。最后，值得注意的是，p 和 τp 的这些表达式是精确的，因为它们直接来自硬球碰撞模型，没有任何简化。然而，为了评估σp，需要考虑六个二阶矩。在流体动力学模型中，仅使用对角二阶矩之和，并且使用纳维-斯托克斯极限中的梯度扩散模型来封闭出现在 τp 中的非对角矩。14 在欧拉极限中，τp 为零，并且方程中 g0 的贡献。 （6.140）被忽略，因此颗粒压力表现得像理想气体。

14 在小克努森数的限制下，弹性玻尔兹曼方程的 Chapman-Enskog 展开式（Chapman，1916；Enksog，1921）产生 σp 的一阶项，其形式为

\begin{sourceblock}
\includegraphics[page=293,pagebox=mediabox,trim={155.28bp 192.34bp 157.81bp 480.17bp},clip,max width=0.94\textwidth,max height=0.38\textheight]{Marchisio2013.pdf}
\end{sourceblock}
\begin{sourceblock}
\includegraphics[page=293,pagebox=mediabox,trim={65.02bp 167.32bp 321.98bp 514.20bp},clip,max width=0.94\textwidth,max height=0.38\textheight]{Marchisio2013.pdf}
\end{sourceblock}
是碰撞时间。在此限制下，粒状应力张量的形式为

\begin{sourceblock}
\includegraphics[page=293,pagebox=mediabox,trim={183.02bp 132.95bp 185.55bp 540.23bp},clip,max width=0.94\textwidth,max height=0.38\textheight]{Marchisio2013.pdf}
\end{sourceblock}
颗粒粘度由下式给出

\begin{sourceblock}
\includegraphics[page=293,pagebox=mediabox,trim={219.51bp 85.54bp 222.04bp 586.31bp},clip,max width=0.94\textwidth,max height=0.38\textheight]{Marchisio2013.pdf}
\end{sourceblock}
{\small 可以对非弹性 Boltzmann-Enskog 方程（Garz ́o et al., 2007；Jenk-ins \& Richman, 1985）进行类似的展开，以找到 τp 的近似值，当 τcol 远小于系统的所有其他特征时间时，该近似值将有效。\par}
\paragraph*{二阶矩}
对于二阶矩，方程： (6.131) 减少为

\begin{sourceblock}
\includegraphics[page=294,pagebox=mediabox,trim={70.37bp 423.76bp 65.02bp 95.71bp},clip,max width=0.94\textwidth,max height=0.38\textheight]{Marchisio2013.pdf}
\end{sourceblock}
碰撞通量和碰撞源项可以使用第 6.1.5 节中的定义进行评估。术语 G200、G110 等取决于三阶矩（并且根据这些矩是封闭的）。在流体动力学模型中，G200、G110等中出现的三阶矩项使用涉及二阶中心矩的空间梯度的梯度扩散模型来封闭。项 C(1) 200、C(1) 110 等也取决于三阶矩及其空间梯度，因此在三阶处闭合。项 C(0) 200、C(0) 110 等在任何有限矩集方面都不是封闭的。然而，对于对角二阶矩，我们有

\begin{sourceblock}
\includegraphics[page=294,pagebox=mediabox,trim={100.90bp 158.33bp 65.08bp 384.72bp},clip,max width=0.94\textwidth,max height=0.38\textheight]{Marchisio2013.pdf}
\end{sourceblock}
将等式中的三个表达式相加。 (6.143) 除以三给出颗粒温度的碰撞源项：

\begin{sourceblock}
\includegraphics[page=294,pagebox=mediabox,trim={78.25bp 96.17bp 65.08bp 570.30bp},clip,max width=0.94\textwidth,max height=0.38\textheight]{Marchisio2013.pdf}
\end{sourceblock}
{\small 通过交换积分变量 v1 和 v2，我们可以找到等效表达式：\par}
\begin{sourceblock}
\includegraphics[page=294,pagebox=mediabox,trim={78.25bp 50.16bp 65.08bp 616.47bp},clip,max width=0.94\textwidth,max height=0.38\textheight]{Marchisio2013.pdf}
\end{sourceblock}
将等式总结在一起。 (6.144) 和 (6.145) 将结果除以二然后得到

\begin{sourceblock}
\includegraphics[page=295,pagebox=mediabox,trim={132.68bp 588.76bp 67.57bp 77.70bp},clip,max width=0.94\textwidth,max height=0.38\textheight]{Marchisio2013.pdf}
\end{sourceblock}
其中我们使用了 ω(ω−1) = −(1−e2)/4 这一事实。因为方程中的结果。 (6.146) 总是小于或等于 0，我们可以观察到非弹性碰撞总是降低颗粒温度。当 f 是麦克斯韦方程时，方程中的积分(6.146) 与 α2 pθ3/2 p 成比例。总之，我们应该重申，在颗粒流的流体动力学模型中，只有方程中对角元素之和的输运方程。 (6.142) 与连续性方程（方程 6.133）和平均动量方程（方程 6.135）结合使用（即颗粒温度的传递方程）。所得方程中的未封闭项使用梯度扩散模型封闭，该模型在 f 接近麦克斯韦的极限内有效。在基于求积的矩方法中，方程中的所有矩。使用式（6.142），并通过根据 f 的矩值重建 f 来实现闭包。原则上，由于重建过程不依赖于 f 接近麦克斯韦方程，因此基于求积的矩方法将适用于更广泛的流态，包括碰撞可以忽略不计的情况。

\subsection*{6.4.2 多分散颗粒}\phantomsection\addcontentsline{toc}{subsection}{6.4.2 多分散颗粒}
对于双分散 15 系统，矩输运方程源自方程： (6.129) 是

\begin{sourceblock}
\includegraphics[page=295,pagebox=mediabox,trim={100.32bp 265.44bp 67.57bp 351.25bp},clip,max width=0.94\textwidth,max height=0.38\textheight]{Marchisio2013.pdf}
\end{sourceblock}
以及

\begin{sourceblock}
\includegraphics[page=295,pagebox=mediabox,trim={99.07bp 155.08bp 67.57bp 461.61bp},clip,max width=0.94\textwidth,max height=0.38\textheight]{Marchisio2013.pdf}
\end{sourceblock}
这些方程具有与单分散情况相同的形式，期望它们包含由于不同粒子之间的碰撞而产生的附加项。请注意，此类碰撞会影响空间通量（例如 Gl1l2l3,12）和碰撞源项（例如 Cl1l2l3,12）。因为等式。 (6.148) 可由式(6.148)求出。 (6.147) 通过指数的排列，此后我们将仅考虑式(6.147)。 (6.147)。

\footnotetext[15]{如前所述，由于碰撞是二元的，因此扩展到多分散系统非常简单，只需为系统中每个可能的碰撞伙伴添加碰撞项即可。}
\paragraph*{零阶矩}
对于零阶矩，方程： (6.147) 减少为

\begin{sourceblock}
\includegraphics[page=296,pagebox=mediabox,trim={130.77bp 568.15bp 65.08bp 102.22bp},clip,max width=0.94\textwidth,max height=0.38\textheight]{Marchisio2013.pdf}
\end{sourceblock}
可以用流体动力学变量重写为质量连续方程：16

\begin{sourceblock}
\includegraphics[page=296,pagebox=mediabox,trim={194.62bp 510.54bp 65.08bp 163.81bp},clip,max width=0.94\textwidth,max height=0.38\textheight]{Marchisio2013.pdf}
\end{sourceblock}
请注意，该方程对于流体动力学变量 α1 和 U1 而言是封闭的。然而，如果将由 U = (ρ1α1U1 + ρ1α1U1)/ρ 和 ρ = ρ1α1 + ρ2α2 定义的质量平均速度用作参考速度，则需要根据 U 提供 U1 的闭包。这通常是通过引入涉及 α1 空间梯度的梯度扩散项来实现的。

\paragraph*{一阶矩}
对于一阶矩，方程。 (6.147) 减少为

\begin{sourceblock}
\includegraphics[page=296,pagebox=mediabox,trim={86.00bp 221.48bp 80.69bp 312.19bp},clip,max width=0.94\textwidth,max height=0.38\textheight]{Marchisio2013.pdf}
\end{sourceblock}
（6.151）可以用流体动力学变量重写为平均动量平衡：

\begin{sourceblock}
\includegraphics[page=296,pagebox=mediabox,trim={162.80bp 165.37bp 65.08bp 508.98bp},clip,max width=0.94\textwidth,max height=0.38\textheight]{Marchisio2013.pdf}
\end{sourceblock}
方程中 C12 碰撞引起的动量交换。 (6.152) 由其组成部分定义：

\begin{sourceblock}
\includegraphics[page=296,pagebox=mediabox,trim={152.05bp 84.85bp 65.08bp 553.77bp},clip,max width=0.94\textwidth,max height=0.38\textheight]{Marchisio2013.pdf}
\end{sourceblock}
\footnotetext[16]{我们引入粒子密度 ρ1，它只是粒子质量 m1 除以其体积 πd3 1/6，以便明确地守恒质量和总平均动量。}
对称二阶压力张量 P1 由其分量定义：

\begin{sourceblock}
\includegraphics[page=297,pagebox=mediabox,trim={144.97bp 509.66bp 67.57bp 73.56bp},clip,max width=0.94\textwidth,max height=0.38\textheight]{Marchisio2013.pdf}
\end{sourceblock}
其中速度协方差矩阵用σ1表示。为了证明 P1 确实是对称的，有必要使用 6.2 节中的定义来检查 G(0) l1l2l3,12 和 G(1) l1l2l3,12 的一阶矩。根据方程中的定义。 (6.107) 和表 6.4 中的结果，我们可以将贡献 G(0) l1l2l3,12 表示为对称二阶张量：

\begin{sourceblock}
\includegraphics[page=297,pagebox=mediabox,trim={74.99bp 387.35bp 67.57bp 264.19bp},clip,max width=0.94\textwidth,max height=0.38\textheight]{Marchisio2013.pdf}
\end{sourceblock}
其中 μ12 = 2m2/(m1 + m2)。通过分离出每个相的平均速度，
(v1 −v2) ⊗(v1 −v2) = [(v1 −U1) −(v2 −U2) + (U1 −U2)] ⊗[(v1 −U1) −(v2 −U2) + (U1 −U2)],
并利用这一事实

\begin{sourceblock}
\includegraphics[page=297,pagebox=mediabox,trim={145.87bp 275.32bp 148.41bp 391.14bp},clip,max width=0.94\textwidth,max height=0.38\textheight]{Marchisio2013.pdf}
\end{sourceblock}
我们可以计算方程中的积分。 (6.155) 求

\begin{sourceblock}
\includegraphics[page=297,pagebox=mediabox,trim={74.83bp 221.98bp 67.57bp 448.99bp},clip,max width=0.94\textwidth,max height=0.38\textheight]{Marchisio2013.pdf}
\end{sourceblock}
并且，通过指数的排列，

\begin{sourceblock}
\includegraphics[page=297,pagebox=mediabox,trim={72.76bp 168.37bp 67.57bp 502.59bp},clip,max width=0.94\textwidth,max height=0.38\textheight]{Marchisio2013.pdf}
\end{sourceblock}
\begin{sourceblock}
\includegraphics[page=297,pagebox=mediabox,trim={65.03bp 113.93bp 67.56bp 538.81bp},clip,max width=0.94\textwidth,max height=0.38\textheight]{Marchisio2013.pdf}
\end{sourceblock}
因此，在 d1 ≪d2 (χ21 →∞) 的极限下，当 ρ1G(0) 12 趋于零时，碰撞通量项 ρ2G(0) 21 将保持有限。 17 同样，根据方程 1 中 G(1) 12 的定义， （6.108）和

\footnotetext[17]{换句话说，在流体-颗粒极限（其中 d1 是分子直径）中，流体相中的压力将与颗粒相无关。}
表6.10中的结果可以表明G(1) 12 具有与G(0) 12 相同的对称性。然而，不可能用闭合形式的矩来表达G(1) 12。使用方程式中的结果。 (6.156) 并忽略 G(1) 12 的贡献，我们可以将 P1 重写为

\begin{sourceblock}
\includegraphics[page=298,pagebox=mediabox,trim={83.98bp 523.20bp 65.08bp 121.17bp},clip,max width=0.94\textwidth,max height=0.38\textheight]{Marchisio2013.pdf}
\end{sourceblock}
根据分解 P1 = τ1 + p1I，1 型颗粒的颗粒压力为

\begin{sourceblock}
\includegraphics[page=298,pagebox=mediabox,trim={95.71bp 464.68bp 65.08bp 206.68bp},clip,max width=0.94\textwidth,max height=0.38\textheight]{Marchisio2013.pdf}
\end{sourceblock}
无迹颗粒应力张量为

\begin{sourceblock}
\includegraphics[page=298,pagebox=mediabox,trim={84.47bp 379.37bp 65.08bp 256.13bp},clip,max width=0.94\textwidth,max height=0.38\textheight]{Marchisio2013.pdf}
\end{sourceblock}
与单分散情况一样，p1 和 τ1 的这些表达式是精确的，并且在二阶矩方面是封闭的。有趣的是，颗粒压力和颗粒应力张量取决于多分散颗粒流中的平均速度差（通过 E12），并且在颗粒温度为零（即无限颗粒马赫数）的极限下，如果平均相速度不同，颗粒压力张量将不为零。 18 这些事实直接来自使用各个相平均速度（而不是质量平均速度）定义流体动力学变量。仅在 d1 <d2 的流体-粒子极限中，p1 和 τ1 才与平均速度差无关。对于方程中的碰撞源项。 （6.153），我们可以使用式（6.153）中的结果。 (6.155) 将 ρ1[C(2) 100,12,C(2) 010,12,C(2) 001,12] 写为向量源项，

\begin{sourceblock}
\includegraphics[page=298,pagebox=mediabox,trim={76.46bp 186.80bp 65.08bp 471.87bp},clip,max width=0.94\textwidth,max height=0.38\textheight]{Marchisio2013.pdf}
\end{sourceblock}
很容易证明 C(2) 21 = −C(2) 12。该项具有固体-固体阻力的形式（尽管具有二阶张量阻力系数），其最大值出现在颗粒温度为零时。请注意，为了稳定性（即固体-固体阻力应迫使 U1 松弛到 U2），参数 c12 = c21 必须大于或等于 0（即 0 ≤ c12 ≤ 1）。我们应该强调方程式。 (6.161) 是精确的，直接遵循硬球碰撞积分的定义。在流体-粒子极限
18 在流体动力学极限下，与两种颗粒类型的颗粒温度相比，平均速度差必须很小（Goldman 和 Sirovich (1967)）。因此，在流体动力学极限 E12 ≈3(θ1 +θ2) 下，τ1 对平均速度差 (U1 −U2) 的依赖性可以忽略不计。本节中推导的矩方程不限于流体动力学极限，即使在颗粒温度为零时也可以应用。

(d1 ≪d2)，等式。 （6.161）可以解释为非各向同性阻力模型。然而，对于无粘性流体 σ1 = θ1I 且 |σ2| ≪θ1，从而从方程中找到阻力系数。 (6.161) 几乎是各向同性的。使用等式中的定义。 (6.103)，我们也可以将 ρ1[C(0) 100,12,C(0) 010,12,C(0) 001,12] 写为向量源项：

\begin{sourceblock}
\includegraphics[page=299,pagebox=mediabox,trim={86.09bp 544.33bp 67.57bp 122.13bp},clip,max width=0.94\textwidth,max height=0.38\textheight]{Marchisio2013.pdf}
\end{sourceblock}
右侧的积分不是封闭的，但我们可以通过以下方式近似它

\begin{sourceblock}
\includegraphics[page=299,pagebox=mediabox,trim={96.41bp 488.27bp 67.57bp 178.19bp},clip,max width=0.94\textwidth,max height=0.38\textheight]{Marchisio2013.pdf}
\end{sourceblock}
其中 c0 是单位阶常数。 C(0) 12 的近似值为

\begin{sourceblock}
\includegraphics[page=299,pagebox=mediabox,trim={139.51bp 428.57bp 67.57bp 241.50bp},clip,max width=0.94\textwidth,max height=0.38\textheight]{Marchisio2013.pdf}
\end{sourceblock}
该源项采用具有标量系数的固-固阻力模型的形式。形式与等式相同。 （6.161），表明 C(0) 12 和 C(2) 12 可以组合起来产生具有张量阻力系数的总固-固阻力项。使用基于求积的矩方法，方程。 (6.162) 可以使用求积权重与节点进行评估。另一种选择是使用动力学模型，但这本质上相当于获得方程 1 中 c0 的准确估计。 （6.163）。另外两个碰撞源向量 C(1) 12 和 C(3) 12 可以使用等式 2 中的定义进行评估。 （6.104）和（6.106）。如前所述，C(1) 12 将在二阶及更低阶矩及其梯度方面闭合。相反，C(3) 12 不会根据任何有限矩集而闭合。尽管如此，它可以使用基于求积的矩方法来近似，如第 6.5 节所述。在流体-粒子极限（d1 <d2）下，C(1) 12 和 C(3) 12 都不会向流体相动量方程贡献涉及流体特性空间梯度（即浮力、升力等）的项。如前所述，这些项由 gαβ(x12) 模型产生，并且会出现，例如，在使用等式 1 中的表达式时。 （6.81）。对于后者，方程。 (6.161) 变为

\begin{sourceblock}
\includegraphics[page=299,pagebox=mediabox,trim={70.00bp 213.98bp 67.57bp 461.59bp},clip,max width=0.94\textwidth,max height=0.38\textheight]{Marchisio2013.pdf}
\end{sourceblock}
21 = -C(2) 12。在流体-颗粒极限中，19 流体压力为 p1 = ρ1α1θ1，并且假设平均速度差可忽略不计20 且 κ21 = −κ12 确保 C(2) 与 θ1/2 1 相比，方程 1 (6.165) 简化为 C(2) 12 = −κ12α2p1 ∇xH。因此，例如，如果 H = -ln p1 且 κ12 = 1,21，我们就可以求出流体相上的浮力：C(2) 12 = α2 ∇xp1。显然，可以提出 H 的其他函数形式来对第 5 章中描述的相动量交换项进行建模。相反，使用第 1 章中描述的多尺度方法，可以采用微尺度模拟来开发
19 在流体-粒子极限中，χ12 = 1/2，μ12 = 2，e12 = 1，g0,12 = 1。此外，ρ1 = ρ1α1 是流体密度，α1 是分子占据的体积（即 α1 + α2 1）。

\footnotetext[20]{对于浮力，只需考虑平均速度为零的情况即可。 21这种情况下的可实现条件是|κ12| < 2/|d2 ∇xln p1|，因此我们必须令 κ12 = min(1, 2/|d2 ∇x ln p1|)。然而，由于在正常流量条件下，压力不太可能在距离 d2/2 上快速变化，因此该条件将产生 κ12 = 1。}
H 的物理精确表达式随后可以在方程式中使用。 （6.165）。总之，使用动力学理论方法，通过考虑 d1 ≪d2 的极限，可以从第一原理推导出气体-颗粒流的双流体模型。捕获相相互作用项过程中的重要步骤是 gαβ(x12) 的定义。

\paragraph*{二阶矩}
作为二阶矩的一个例子，我们只考虑一个：

\begin{sourceblock}
\includegraphics[page=300,pagebox=mediabox,trim={85.17bp 475.25bp 65.08bp 169.26bp},clip,max width=0.94\textwidth,max height=0.38\textheight]{Marchisio2013.pdf}
\end{sourceblock}
与一阶矩一样，该表达式由于左侧的动力学和碰撞通量以及右侧的碰撞源项而有所贡献。由类粒子碰撞（G1,200,11 和 C200,11）引起的贡献与上述单分散粒子的情况具有相同的形式。因此，我们将简要地了解异粒子碰撞引起的术语。碰撞通量项 Gm,200,12 有两个贡献：

\begin{sourceblock}
\includegraphics[page=300,pagebox=mediabox,trim={97.15bp 368.47bp 65.08bp 297.99bp},clip,max width=0.94\textwidth,max height=0.38\textheight]{Marchisio2013.pdf}
\end{sourceblock}
以及

\begin{sourceblock}
\includegraphics[page=300,pagebox=mediabox,trim={72.53bp 319.00bp 70.33bp 344.58bp},clip,max width=0.94\textwidth,max height=0.38\textheight]{Marchisio2013.pdf}
\end{sourceblock}
(6.168) 方程中的项。 (6.167) 可以使用高达三阶的整数矩重写为闭合形式。另一方面，等式的右侧。 （6.168）不能用任何有限的整数矩集来写，因此必须是封闭的（例如使用求积）。碰撞源项 C200,12 有四个贡献：

\begin{sourceblock}
\includegraphics[page=300,pagebox=mediabox,trim={78.66bp 54.89bp 65.08bp 439.12bp},clip,max width=0.94\textwidth,max height=0.38\textheight]{Marchisio2013.pdf}
\end{sourceblock}
方程中的积分。 (6.170) 和 (6.171) 可以根据矩进行明确评估。相反，方程中的积分。 (6.169) 和 (6.172) 就有限矩集而言不是封闭的，并且必须封闭（例如使用求积）。

\section*{6.5 求积在碰撞项中的应用}\phantomsection\addcontentsline{toc}{section}{6.5 求积在碰撞项中的应用}
随着阶数的增加，上面讨论的矩输运方程变得越来越复杂。此外，这些方程不是封闭的。在基于求积的矩方法中，速度分布函数是根据有限的矩集重建的，从而提供了闭包。在本节中，我们将说明如何应用闭合假设来求解具有硬球碰撞的矩输运方程。为了清楚起见，我们将考虑由方程式控制的单分散情况。 （6.131）。形式上，我们可以用保守形式重新表达这个方程：

\begin{sourceblock}
\includegraphics[page=301,pagebox=mediabox,trim={186.87bp 439.21bp 67.57bp 230.28bp},clip,max width=0.94\textwidth,max height=0.38\textheight]{Marchisio2013.pdf}
\end{sourceblock}
其中 Fγ l1l2l3 是矩 Mγ l1l2l3 的空间通量，S γ l1l2l3 是其源项。如果注意确保矩始终保持可实现，则可以使用有限体积方法有效地求解此类方程（Leveque，2002）（Vikas 等人，2009）。 （有关更多详细信息，请参阅附录 B。）后者可以使用根据重建的速度分布函数定义的基于动力学的通量来优雅地完成。在本节的其余部分中，我们将展示方程式如何。 (6.173) 使用基于求积的矩方法进行处理。

\subsection*{6.5.1 通量项}\phantomsection\addcontentsline{toc}{subsection}{6.5.1 通量项}
形式上，方程中的通量项。 (6.173) 可以表示为

\begin{sourceblock}
\includegraphics[page=301,pagebox=mediabox,trim={163.71bp 248.34bp 67.57bp 400.63bp},clip,max width=0.94\textwidth,max height=0.38\textheight]{Marchisio2013.pdf}
\end{sourceblock}
其中 Gl1l2l3 是方程式中定义的碰撞通量。 （6.70）。在基于动力学的通量算法中，Fγ l1l2l3 的分量被分为两部分，分别对应于通量方向上的正速度分量和负速度分量。例如，考虑第一个分量 Fγ 1,l1l2l3。使用 Mγ+1 l1+1l2l3 和 G1,l1l2l3 的定义，我们可以将通量拆分为

\begin{sourceblock}
\includegraphics[page=301,pagebox=mediabox,trim={191.39bp 167.32bp 67.57bp 513.85bp},clip,max width=0.94\textwidth,max height=0.38\textheight]{Marchisio2013.pdf}
\end{sourceblock}
其中

\begin{sourceblock}
\includegraphics[page=301,pagebox=mediabox,trim={74.73bp 116.58bp 97.84bp 549.37bp},clip,max width=0.94\textwidth,max height=0.38\textheight]{Marchisio2013.pdf}
\end{sourceblock}
(6.176) and

\begin{sourceblock}
\includegraphics[page=301,pagebox=mediabox,trim={75.55bp 52.94bp 67.57bp 603.54bp},clip,max width=0.94\textwidth,max height=0.38\textheight]{Marchisio2013.pdf}
\end{sourceblock}
分别是对应于负 v1 速度和正 v1 速度的通量分量。 Fγ l1l2l3 的第二个和第三个分量也有类似的定义，因此我们将得到 Fγ l1l2l3 + F+ l1l2l3 = F− l1l2l3 形式的分解。读者应该注意，对于每个分量，方程中的外积分。 (6.176) 和 (6.177) 每个方向都会发生变化。在有限体积算法中，方程中的通量分裂。 （6.175）直接用于计算相对于给定表面元素的逆风和顺风通量分量（例如 Desjardins 等，2008），因此很自然地有单独计算它们的公式。等式的基于求积的闭包。 (6.176) 和 (6.177) 则如下进行。令 f ∗(t, x, v) 表示从传输矩 Mγ l1l2l3 (方程 6.173) 重建的速度分布函数，并将方向 x1 的负整数矩和正整数矩定义为22

\begin{sourceblock}
\includegraphics[page=302,pagebox=mediabox,trim={157.06bp 437.88bp 65.08bp 199.98bp},clip,max width=0.94\textwidth,max height=0.38\textheight]{Marchisio2013.pdf}
\end{sourceblock}
可以为方向 x2 和 x3 定义类似的矩。由于 f * 是已知函数，因此可以针对 (l1, l2, l3) 的任何选择计算（通常是解析式）这些矩。因此，使用每组矩，我们可以构造一个任意高阶求积表示，该表示可用于评估方程（2）中的积分。 （6.176）和（6.177）。形式上，我们可以将重建过程表示为

\begin{sourceblock}
\includegraphics[page=302,pagebox=mediabox,trim={123.30bp 339.43bp 121.55bp 320.58bp},clip,max width=0.94\textwidth,max height=0.38\textheight]{Marchisio2013.pdf}
\end{sourceblock}
其中N±是求积节点的数量，n±α是权重，v±α是节点-sas。可以使用例如 CQMOM (Yuan \& Fox, 2011) 来计算每个矩集的权重与节点。读者应该认识到，针对方向 x1、x2 和 x3 计算的矩将生成不同的求积，因此将构造六组求积权重与节点来评估通量项。如果使用EQMOM来重建f*，则可以直接从第二个（高斯）求积中找到正负节点，而无需计算式（1）中的半矩。 （6.178）。然而，重要的是（当使用不带半矩的 EQMOM 时）采用相对大量的正交节点以确保准确地再现半矩。给定六组求积，现在可以评估通量分量。例如，从等式。 (6.176) 和 (6.177) 我们发现

\begin{sourceblock}
\includegraphics[page=302,pagebox=mediabox,trim={82.61bp 154.95bp 65.08bp 505.31bp},clip,max width=0.94\textwidth,max height=0.38\textheight]{Marchisio2013.pdf}
\end{sourceblock}
以及

\begin{sourceblock}
\includegraphics[page=302,pagebox=mediabox,trim={80.53bp 101.93bp 65.08bp 558.32bp},clip,max width=0.94\textwidth,max height=0.38\textheight]{Marchisio2013.pdf}
\end{sourceblock}
\footnotetext[22]{重建的分布函数可以是连续的（EQMOM）或离散的（QMOM），但我们假设它总是可实现的（即非负）。对于 f ∗ 是一组加权 delta 函数的情况，矩 M− 的计算}
{\small i,l1l2l3 and M+\par}
{\small i,l1l2l3 是微不足道的。使用 EQMOM，使用根据核密度函数定义的高斯求积来评估积分。\par}
其中 β 的求和是针对整个 N−+ N+ 权重与节点集进行的。23

\begin{sourceblock}
\includegraphics[page=303,pagebox=mediabox,trim={65.02bp 614.83bp 245.45bp 69.38bp},clip,max width=0.94\textwidth,max height=0.38\textheight]{Marchisio2013.pdf}
\end{sourceblock}
1,l1l2l3 = 0 (F+ 1,l1l2l3 = 0)。） (6.179) 和 (6.180) 中，涉及 I(1) l1l2l3 的碰撞通量通常需要比自由传输项更多的正交节点才能准确计算。因此，EQMOM 特别有用，因为正交节点的数量可以独立于传输矩的数量进行设置。从等式中的表达式可以明显看出重要的一点。 （6.179）和（6.180）的要点是，只要知道 I(m) l1l2l3 的函数形式就足以解释碰撞对空间通量的影响。从表 6.4-6.9 可以看出，这些函数形式可以很容易地编码成数值算法，因此，方程(6.179) 和 (6.180) 可以直接计算。这种情况可以与为矩输运方程提出的“传统”闭包进行比较（例如 Struchtrup，2005），其中随着矩阶数的增加，导出闭包所需的操作变得越来越复杂。同样重要的是要注意方程式。即使在传统矩闭包失败的高度非平衡情况下，（6.179）和（6.180）也是精确的。例如，如果精确的速度分布函数是单动力学的，

\begin{sourceblock}
\includegraphics[page=303,pagebox=mediabox,trim={206.48bp 424.74bp 207.52bp 253.06bp},clip,max width=0.94\textwidth,max height=0.38\textheight]{Marchisio2013.pdf}
\end{sourceblock}
那么 N−+ N+ = 1，因此由于碰撞对空间通量的贡献将为零，并且将精确计算空间通量（即仅以速度 Up 进行自由传输）。计算高度非平衡情况下的精确空间通量的能力是使用基于求积的矩方法的主要吸引力之一。

\subsection*{6.5.2 来源项}\phantomsection\addcontentsline{toc}{subsection}{6.5.2 来源项}
方程中的源项。 (6.173) 由下式给出

\begin{sourceblock}
\includegraphics[page=303,pagebox=mediabox,trim={198.98bp 303.00bp 67.57bp 378.16bp},clip,max width=0.94\textwidth,max height=0.38\textheight]{Marchisio2013.pdf}
\end{sourceblock}
其中右侧的两项在等式中定义。 （6.68）和等式。分别为（6.71）。其中第一个使用求积法进行评估：

\begin{sourceblock}
\includegraphics[page=303,pagebox=mediabox,trim={136.26bp 239.44bp 67.57bp 427.50bp},clip,max width=0.94\textwidth,max height=0.38\textheight]{Marchisio2013.pdf}
\end{sourceblock}
其中 N 个权重和 N 个节点是从重建的速度分布函数 f 的矩（即从其离散正交表示）中找到的。请注意，这些时刻与从等式中找到的时刻相关。 (6.178) 由 M* i,l1l2l3 = M− i,l1l2l3 + M+ i,l1l2l3 得出，因此不需要单独计算。对于空间通量，方程的计算： （6.182）很简单，只需要 I(0) l1l2l3 的函数形式（见表 6.1-6.3）。此外，方程。即使对于传统矩闭包失败的高度非平衡情况，（6.182）也是精确的。 C(1) l1l2l3 的评估比较复杂，因为它涉及空间梯度，

\begin{sourceblock}
\includegraphics[page=303,pagebox=mediabox,trim={107.40bp 101.81bp 67.57bp 564.14bp},clip,max width=0.94\textwidth,max height=0.38\textheight]{Marchisio2013.pdf}
\end{sourceblock}
\footnotetext[23]{或者，可以使用由矩 M* 构造的 N 节点求积来定义 β 的求和}
{\small 我，l1l2l3。如果 N−= N+ = N，则两种方法的形式精度均为 2N −1 (Gautschi, 1997)。从 M* 求积\par}
{\small i,l1l2l3 = M−\par}
{\small i,l1l2l3 + M+\par}
{\small i,l1l2l3 通量方程中的求和较短（N 与 2N），但计算另一个正交的额外计算成本可能更高。\par}
其中向量 Hl1l2l3 的分量由等式定义： （6.72）。如前所述，Hm,l1l2l3 是关于 v1 和 v2 的 γ + 1 阶齐次多项式。因此，形式上，我们可以将 Hl1l2l3 表示为

\begin{sourceblock}
\includegraphics[page=304,pagebox=mediabox,trim={192.48bp 559.74bp 190.03bp 107.04bp},clip,max width=0.94\textwidth,max height=0.38\textheight]{Marchisio2013.pdf}
\end{sourceblock}
这样24

\begin{sourceblock}
\includegraphics[page=304,pagebox=mediabox,trim={168.69bp 507.33bp 65.08bp 159.46bp},clip,max width=0.94\textwidth,max height=0.38\textheight]{Marchisio2013.pdf}
\end{sourceblock}
其中 Mn 表示适当的 n 阶速度矩。原则上，系数 an 可以从 I(m) l1l2l3 的定义中找到，因此方程： （6.184）可以根据速度矩展开。例如，来自 H1,200 对 C(1) 200 的贡献是

\begin{sourceblock}
\includegraphics[page=304,pagebox=mediabox,trim={79.43bp 263.18bp 74.77bp 251.22bp},clip,max width=0.94\textwidth,max height=0.38\textheight]{Marchisio2013.pdf}
\end{sourceblock}
而来自 H2,200 的贡献是

\begin{sourceblock}
\includegraphics[page=304,pagebox=mediabox,trim={88.15bp 83.67bp 83.48bp 460.92bp},clip,max width=0.94\textwidth,max height=0.38\textheight]{Marchisio2013.pdf}
\end{sourceblock}
\footnotetext[24]{为了清楚起见，这里我们滥用了这种符号，但读者应该认识到 (∇x · Mn) 涉及三个不同的矩（每个方向一个）。}
H3,200 的贡献具有类似的形式，因此 C(1) 200 的最终表达式涉及高达三阶的所有 20 个速度矩的空间梯度。25 使用由 f * 构造的 N 节点求积表示，方程中出现的矩： (6.184) 可以正式表示为

\begin{sourceblock}
\includegraphics[page=305,pagebox=mediabox,trim={214.04bp 559.33bp 216.08bp 107.62bp},clip,max width=0.94\textwidth,max height=0.38\textheight]{Marchisio2013.pdf}
\end{sourceblock}
给定矩，可以使用一致的有限体积表示来计算梯度 (∇x · Mn)。据作者所知，这种计算 C(1) l1l2l3 的方法尚未在任何模拟代码中实现。总之，我们在本节中演示了如何使用基于求积的矩方法来评估由碰撞引起的矩输运方程中的项。主要的观察结果是，知道项 I(m) l1l2l3 的函数形式就足够了，这些形式是在 6.1 节中推导和列表的。我们还观察到，与传统的矩闭包不同，本节中开发的闭包适用于高度非平衡的流动。

25 由于 C(1) 000 和 C(1) 100 为零，C(1) 200 中的非零贡献表示由于密度和速度梯度而导致的颗粒温度的源项。在单动力学情况下，这些贡献将导致颗粒温度增加到非零值（从而使分布不再单动力学）。从物理上讲，C(1) 200 中的项之所以出现，是因为具有有限直径 dp 和有限间隔的粒子在具有不同速度（由于平均速度梯度）时可能会发生碰撞，即使颗粒温度为零。

\chapter{均匀系统的求解方法}
在本章中，我们考虑在空间均匀性假设下用一组通用内坐标 ε = (Ψ1, Ψ2, ..., ΨM) 求解广义群体平衡方程 (GPBE)。在此假设下，所有空间梯度均为零，并且 GPBE 仅取决于时间 t 和 ε。如前所述，我们在这里将 GPBE 称为描述数密度函数 (NDF) 在空间（本章中忽略）、时间以及由内坐标生成的相空间中演化的一般方程。第8章讨论了非齐次系统的扩展。还值得一提的是，GPBE在不同领域有不同的名称。它在结晶、降水以及一般颗粒过程文献中被称为群体平衡方程 (PBE)。在气溶胶的模拟中，它被称为一般气溶胶动力学方程（Friedlander，2000），而在喷雾的模拟中，它通常被称为威廉姆斯-玻尔兹曼方程（Williams，1985），特别是在这种情况下，液滴速度被作为内坐标之一。当处理固体粒子时，如果将粒子速度作为唯一的内坐标，则 GPBE 称为玻尔兹曼方程。一般来说，当粒子速度是唯一的内坐标时，GPBE 也称为动力学方程 (KE)。1 许多与 KE 相关的挑战都发生在非齐次系统中，这些问题将在第 8 章中讨论。

\section*{7.1 方法概述}\phantomsection\addcontentsline{toc}{section}{7.1 方法概述}
根据内坐标的数量（单变量与多变量）以及粒子速度是否是内坐标向量的一部分，已经开发了非常不同的求解方法。因此，本章将分别讨论有和没有粒子速度的情况的方法。首先，讨论了为解决单变量 PBE（即一个内坐标）而开发的方法。其次，双变量和多变量 PBE 的求解方法（即

\footnotetext[1]{在应用数学文献中，所有在相空间中演化的具有内坐标的平衡方程通常被称为动力学方程 (KE)。使用这种分类，GPBE 或 PBE 可以简称为 KE。然而，由于物理空间和相空间中的动力学导致了相当明显的差异，因此我们更愿意细化定义并使用 KE 来描述唯一内坐标是速度的情况。}
两个或多个内坐标）被呈现。最后，提出了 KE（即粒子速度是唯一的内坐标）的求解方法。在空间均匀性的假设下（即 NDF 和 GPBE 中的所有项都不依赖于 x），控制 GPBE 变为 2

\begin{sourceblock}
\includegraphics[page=307,pagebox=mediabox,trim={174.20bp 555.81bp 67.56bp 117.48bp},clip,max width=0.94\textwidth,max height=0.38\textheight]{Marchisio2013.pdf}
\end{sourceblock}
其中 n(t, xi) 是仅依赖于时间和内坐标的 NDF，̇xi(xi) 是连续过程的内坐标变化率，3 D(t, xi) 是相空间扩散张量，S 是由于相空间不连续跳跃而产生的源项。如第 5 章所述，涉及 ̇ 的漂移项代表连续现象的影响。4 当研究的内坐标是粒子速度时，漂移项量化了连续力的影响，例如体力和流体-粒子力（例如浮力、阻力和升力）。等式右侧的项。 （7.1）通常被标记为不连续点过程并表示不连续事件的影响。 （示例请参见第 5 章。）现在让我们介绍可用于求解方程 1 的不同方法。 （7.1），可大致分为四类：（i）类和截面方法，（ii）矩方法，（iii）基于求积的矩方法和（iv）蒙特卡罗方法。

\paragraph*{类法和分段法}
类方法（CM）基于内坐标的离散化，使GPBE成为一组宏观天平。这些方法已经由颗粒系统模拟领域的几位作者开发并广泛应用于单变量 PBE，而在双变量和多变量 PBE 的情况下，它们的应用则更加困难（Hounslow 等，1988；Kumar \& Ramkrishna，1996a，1996b；Litster 等，1995；Marhal 等，1988；瓦尼，1999，2000a；永利，2004）。一般来说，多变量系统的应用基于变量数量的减少（Briesen，2004；Puel 等人，2003a，2003b；Zhang \& Doherty，2004）和单变量方法的扩展（Kumar 等人，2006，2007，2008b）的组合。离散化内坐标的相同基本思想已在其他领域用于模拟蒸发喷雾，从而产生了通常所说的喷雾多流体模型（Greenberg et al., 1986, 1993; Laurent \& Massot, 2001）。然而，最终的方程与上述方法非常相似。当用于求解 KE 时，相同的方法被称为离散速度模型 (DVM)，并且一起
2 GPBE 用点积向量表示法表示。在分量形式中，重复的希腊指数意味着在内坐标上求和，GPBE 为

\begin{sourceblock}
\includegraphics[page=307,pagebox=mediabox,trim={183.47bp 144.94bp 184.81bp 532.11bp},clip,max width=0.94\textwidth,max height=0.38\textheight]{Marchisio2013.pdf}
\end{sourceblock}
对于均匀系统，漂移 (̇ϋα) 和扩散 (Dαβ) 分量可以是 ϋ 和 t 的函数。 3 连续变化率也可以是时间的函数。然而，在本章中，我们将仅讨论内坐标依赖性，因为这些是导致最重要的数值问题的原因。 4 有趣的是，由于漂移项的数值解存在问题，特别是当相空间离散化时，有时这种观点是相反的。事实上，在最近的作品中，一些连续过程被描述为实际上是不连续过程。正如我们将看到的，这种策略解决了一些问题（即相空间中的数值扩散），但由于控制过程的时间尺度不同，通常会使问题变得非常棘手。对细节感兴趣的读者可以参考 Kumar 等人的著作。 （2008b）。

PBE 的标准方法，将在 7.2 节中介绍和讨论。正如在下一节中将变得更加清楚的那样，对于 CM 和 DVM，NDF 通常在离散内坐标的每个区间中通过低阶多项式来近似。通常使用零阶多项式（即每个单类中的常数 NDF），从而产生零阶方法。5 值得一提的是，这种选择是任意的，并且对方法的性能有一些重要影响。或者，可以假设属于特定区间的所有粒子具有相同的内坐标值，因此集中在一个点，称为枢轴点。相反，高阶方法基于使用高阶多项式或其他基函数对感兴趣区间内的 NDF 进行近似。后者的数值方法属于加权残差法（MWR）的范畴。使用一组加权函数，实际上可以计算残差的加权和，然后在选定的点将其驱动为零。当局部基函数是正交多项式并且加权函数是狄拉克δ函数时，该方法被称为搭配技术（Crowley et al., 2000; Gelbard \& Sein-feld, 1978; Immanuel et al., 2002; Immanuel \& Doyle, 2003），而当加权函数与基函数相同时，该方法被称为 Galerkin方法6 (Chiu \& Christofides, 1999; Godin et al., 1999; Mahoney \& Ramkrishna, 2002; Wulkow et al., 2001)。由于内部坐标离散化的子域也称为有限元，因此这些方法也被称为有限元方法（FEM）。 FEM确实非常强大，但由于其计算成本较高，不被认为适用于非齐次系统，因此不会详细讨论。另一类不会讨论的方法是基于相空间有限体积离散化的方法。这些方法在求解方程式的漂移和扩散项方面非常有效。 （7.1），但是扩展到涉及不连续跳跃（即由于例如碰撞）的问题并不简单。另一个重要的优点在于，很容易将有限体积方法扩展到多元 PBE（例如 Gunawan 等，2004），但是，即使在单变量情况下，对于一般非齐次 PBE 的求解来说，计算成本也太高。

\paragraph*{矩方法 (MOM)}
MOM 是由 Hulburt \& Katz (1964) 为颗粒系统引入的。在他们的开创性工作中，这些作者展示了如何根据 NDF 的矩求解 PBE 并推导相应的输运方程。类似的方法可用于求解 KE，关于矩输运方程的推导的详细讨论可以在 Struchtrup (2005) 和 Truesdell \& Muncaster (1980) 的著作中找到。该技术的主要问题在于闭包问题，即不可能编写仅涉及低阶矩的 NDF 低阶矩的输运方程。自 Hulburt \& Katz (1964) 的工作以来，已经取得了很大进展（Frenklach，2002；Frenklach \& Harris，1987；Kaza-kov \& Frenklach，1998），并且提出了不同的数值闭包（Alexiadis 等，2004；Kostoglou \& Karabelas，2004；Strumendo \& Arastoopour， 2008）。基础
5 这是一种常见但令人困惑的符号。事实上，根据近似排除的项对这些方法进行分类在概念上更正确。对于后一种方法，在每个区间内用恒定值逼近 NDF 的方法可以表示为一阶方法，因为当将 NDF 逼近为泰勒级数时，排除的第一项是一阶导数项。

\footnotetext[6]{粒度评估软件包 PARSIVAL 基于伽辽金法。}
MOM 的概念和早期为克服封闭问题而开发的策略（通过在 NDF 上采用不同的函数假设）将在第 7.3 节中讨论。

\paragraph*{基于求积的矩方法 (QBMM)}
在矩方法的背景下，McGraw (1997) 提出的一个基于求积公式的特定闭包特别有趣，并且产生了一类全新的方法。这些 QBMM 在 7.4 节中详细讨论。更清楚的是，通过 QBMM，NDF 被重构为狄拉克 δ 函数（和其他基函数）的总和，这样这些基函数的权重和相空间中的位置都可以自由移动和模拟 NDF，正确地表示 NDF 的大量矩。第 3 章提供了高斯求积和矩反演算法的技术细节。因此，第 7.4 节的重点是 QBMM 在 GPBE 中的应用。

\paragraph*{蒙特卡罗方法}
蒙特卡罗 (MC) 方法基于一组控制颗粒系统行为的随机微分方程的解，并且广义上 MC 方法基于系统行为的人工实现。不同的作者已经开发了几种方法（Balthasar \& Kraft, 2003; Lattuada et al., 2004; Lee \& Mat-soukas, 2000; Lin et al., 2002; Maisels et al., 2004a, 2004b; Mishra, 2000; Rosner \& Yu, 2001; Rosner et al., 2003；Song \& Qiu，1999；Van Peborgh Gooch \& Hounslow，1996），MC 方法的简短概述在第 7.5 节中介绍。由于本章的重点是矩方法，因此第 7.5 节中讨论的重点是向读者介绍 MC 方法作为验证矩闭包的潜在工具。

\section*{7.2 分类法和分段法}\phantomsection\addcontentsline{toc}{section}{7.2 分类法和分段法}
正如上一节所介绍的，类方法和截面方法基于内坐标的离散化，使得 GPBE 成为状态空间中的一组宏观平衡。事实上，离散化的精细度将由 GPBE 中出现的积分和导数项的近似所需的精度决定。正如已经预料到的，这些方法根据描述中使用的内坐标的数量而不同，并且取决于内坐标的性质。因此，接下来我们将分别讨论单变量、双变量和多变量PBE，以及这些方法在求解KE时的用途。

\subsection*{7.2.1 单变量 PBE}\phantomsection\addcontentsline{toc}{subsection}{7.2.1 单变量 PBE}
设 Ψ 表示唯一内坐标，n(t, Ψ) 表示 NDF。单变量 PBE 为

\begin{sourceblock}
\includegraphics[page=309,pagebox=mediabox,trim={197.48bp 87.45bp 67.56bp 585.25bp},clip,max width=0.94\textwidth,max height=0.38\textheight]{Marchisio2013.pdf}
\end{sourceblock}
{\small 式中，ψ为内坐标的连续变化率，D为内坐标的扩散变化，S为内坐标的不连续跳跃。\par}
将漂移、扩散和点过程表示为唯一的源项 (Ramkrishna, 2000) 得到 ∂n

\begin{sourceblock}
\includegraphics[page=310,pagebox=mediabox,trim={188.60bp 589.23bp 65.07bp 83.46bp},clip,max width=0.94\textwidth,max height=0.38\textheight]{Marchisio2013.pdf}
\end{sourceblock}
其中 S (t, xi, n) 是将 NDF 映射为时间和内坐标函数的函数。现在，我们为以 xii 和 xii+1 为界的内坐标定义一个有限区间 Ii。通过对 Ii 上的 PBE 积分，可以得到 Ii 区间内元素的宏观平衡：

\begin{sourceblock}
\includegraphics[page=310,pagebox=mediabox,trim={194.24bp 509.47bp 65.07bp 156.21bp},clip,max width=0.94\textwidth,max height=0.38\textheight]{Marchisio2013.pdf}
\end{sourceblock}
如果我们将 Ii 区间中的粒子数表示为

\begin{sourceblock}
\includegraphics[page=310,pagebox=mediabox,trim={215.32bp 455.85bp 65.07bp 209.83bp},clip,max width=0.94\textwidth,max height=0.38\textheight]{Marchisio2013.pdf}
\end{sourceblock}
我们认识到

\begin{sourceblock}
\includegraphics[page=310,pagebox=mediabox,trim={214.98bp 408.20bp 65.07bp 257.49bp},clip,max width=0.94\textwidth,max height=0.38\textheight]{Marchisio2013.pdf}
\end{sourceblock}
其中 i = 1，.。。，M，对应于如下划分PM： Ψ1 < Ψ2 < Ψ3 <··· < ΨM−1 < ΨM < ΨM+1。选择 Ψ1 和 ΨM 的值，以便排除的动态可以忽略不计，而间隔的数量是固定的，以便达到所需的精度。所得的 M 个常微分方程 (ODE) 集合是未封闭的。事实上，方程右边是不可能写出来的。 （7.6）仅就Ni而言，据说这反映了缺乏自主性。离散方法的使用取决于自主性的恢复，这取决于我们能够完全用因变量 Ni 来表达右侧。通常，采用关于单个有限区间内 NDF 函数形式的不同假设，一般来说，这些可以写为 n(t, xi) = Ni(t)ki(xi)，其中 xii < xi < xii+1，其中 i = 1，...。。。 、M 和 ki(xi) 是必须满足属性 ∫ xii+1 的多项式

\begin{sourceblock}
\includegraphics[page=310,pagebox=mediabox,trim={217.88bp 232.81bp 209.06bp 446.42bp},clip,max width=0.94\textwidth,max height=0.38\textheight]{Marchisio2013.pdf}
\end{sourceblock}
在零阶多项式的情况下，每个类中都采用（常数）近似（Greenberg et al., 1986, 1993）：

\begin{sourceblock}
\includegraphics[page=310,pagebox=mediabox,trim={173.67bp 169.87bp 171.22bp 503.52bp},clip,max width=0.94\textwidth,max height=0.38\textheight]{Marchisio2013.pdf}
\end{sourceblock}
其中 i = 1，.。。，M。另一种假设是关键点 ζi，其中区间内的所有粒子都集中：

\begin{sourceblock}
\includegraphics[page=310,pagebox=mediabox,trim={194.45bp 96.80bp 65.07bp 580.33bp},clip,max width=0.94\textwidth,max height=0.38\textheight]{Marchisio2013.pdf}
\end{sourceblock}
{\small 通常，离散化不是根据属于单个类别的粒子数量来编写，而是根据属于特定类别的粒子的体积（或质量）来编写。如果 ρp 是粒子的密度，Vp(xi) 是函数关系\par}
在粒子体积和内坐标之间，7 Ii 区间内的粒子质量为 (Laurent \& Massot, 2001)

\begin{sourceblock}
\includegraphics[page=311,pagebox=mediabox,trim={196.93bp 577.87bp 67.56bp 87.81bp},clip,max width=0.94\textwidth,max height=0.38\textheight]{Marchisio2013.pdf}
\end{sourceblock}
而 Ii 区间内相应的基本假设为 n(t, xi) = Mi(t)ki(xi)，其中 ψi < ψ < ψi+1，并且

\begin{sourceblock}
\includegraphics[page=311,pagebox=mediabox,trim={192.98bp 509.07bp 195.51bp 156.61bp},clip,max width=0.94\textwidth,max height=0.38\textheight]{Marchisio2013.pdf}
\end{sourceblock}
同样，每个区间内 NDF 恒定近似的情况会导致

\begin{sourceblock}
\includegraphics[page=311,pagebox=mediabox,trim={156.82bp 448.73bp 159.35bp 218.64bp},clip,max width=0.94\textwidth,max height=0.38\textheight]{Marchisio2013.pdf}
\end{sourceblock}
下面讨论一些克服和解决自治恢复问题的策略。仅针对不连续、漂移和扩散过程分别介绍和分析最流行的方法。

\paragraph*{二阶和一阶点过程}
让我们首先讨论与选择特定离散化相关的含义以及通过使用两种不同的方法恢复二阶点过程的自主性：固定（和移动）枢轴技术和单元平均技术。如果 ψ 表示粒子的附加属性（例如粒子质量），则二阶点过程的源项为

\begin{sourceblock}
\includegraphics[page=311,pagebox=mediabox,trim={111.22bp 274.76bp 67.56bp 391.44bp},clip,max width=0.94\textwidth,max height=0.38\textheight]{Marchisio2013.pdf}
\end{sourceblock}
其中 β(ψ, ψ′) 是质量分别等于 ψ 和 ψ′ 的两个粒子的聚合核。代入等式。 (7.4) 给出

\begin{sourceblock}
\includegraphics[page=311,pagebox=mediabox,trim={75.67bp 206.56bp 73.82bp 459.12bp},clip,max width=0.94\textwidth,max height=0.38\textheight]{Marchisio2013.pdf}
\end{sourceblock}
(7.10) 将第二个积分表示为 xii 和 xii+1 之间的积分之和，并重新排列总和，得出

\begin{sourceblock}
\includegraphics[page=311,pagebox=mediabox,trim={132.83bp 95.91bp 67.56bp 532.10bp},clip,max width=0.94\textwidth,max height=0.38\textheight]{Marchisio2013.pdf}
\end{sourceblock}
\footnotetext[7]{显然，这种关系取决于内坐标的性质。在颗粒尺寸的情况下（即 xi = Lp）Vp = kVL3}
{\small p，其中 kV 是体积形状因子，而对于颗粒表面积 (Σ = Ap)\par}
\begin{sourceblock}
\includegraphics[page=311,pagebox=mediabox,trim={96.30bp 52.01bp 370.02bp 634.30bp},clip,max width=0.94\textwidth,max height=0.38\textheight]{Marchisio2013.pdf}
\end{sourceblock}
{\small Vp = kVA3/2\par}
{\small A，其中 kA 是表面形状因子。\par}
自主性的恢复依赖于我们写出方程中积分的能力。 (7.11) 作为 Ni 的函数。这可以通过使用中值定理来完成。例如，第二个积分对应于以下等式：

\begin{sourceblock}
\includegraphics[page=312,pagebox=mediabox,trim={147.12bp 565.53bp 65.07bp 100.15bp},clip,max width=0.94\textwidth,max height=0.38\textheight]{Marchisio2013.pdf}
\end{sourceblock}
其中 zi 和 zj 是枢轴点或枢轴。在固定枢轴技术中，枢轴在模拟过程中是固定的（Kumar \& Ramkrishna，1996a），而在移动枢轴技术（Kumar \& Ramkrishna，1996b）中，枢轴可以在区间内自由移动。8 应用中值定理并使用枢轴的概念，方程： (7.11) 变为

\begin{sourceblock}
\includegraphics[page=312,pagebox=mediabox,trim={143.08bp 467.04bp 65.07bp 199.38bp},clip,max width=0.94\textwidth,max height=0.38\textheight]{Marchisio2013.pdf}
\end{sourceblock}
其中右侧第一项代表新粒子的生成，第二个求和仅限于内坐标落入区间 Ii 的粒子对（即 zej 和 zek）。方程中的最后一项。 (7.13) 是粒子的损失。方程中的生成项。 (7.13) 导致了大部分问题。事实上，如果 ψ 是粒子质量，那么，当属于区间 I j 的粒子与属于区间 Ik 的粒子聚集时，就形成了质量等于 zej + zek 的粒子。除非使用线性划分，否则不太可能精确地在 zei = zej + zek 处具有枢轴。这意味着属于 Ii 区间的新形成的质量为 zej + zek 的粒子可能具有不等于 zej + zek 的相应枢轴 zei。因此，新形成的粒子必须在相邻关键点之间共享，以保留分布的一些全局属性。 Ramkrishna (2000) 和 Kumar \& Ramkrishna (1996a) 通过推导 NDF r 阶矩的输运方程广泛讨论了这个问题。对于二阶点过程，我们得到

\begin{sourceblock}
\includegraphics[page=312,pagebox=mediabox,trim={93.02bp 260.07bp 65.07bp 410.17bp},clip,max width=0.94\textwidth,max height=0.38\textheight]{Marchisio2013.pdf}
\end{sourceblock}
与方程式的比较(7.13) 乘以 zir 并对索引 i 的所有可能值求和，

\begin{sourceblock}
\includegraphics[page=312,pagebox=mediabox,trim={106.20bp 185.26bp 65.07bp 484.98bp},clip,max width=0.94\textwidth,max height=0.38\textheight]{Marchisio2013.pdf}
\end{sourceblock}
表明方程。 （7.14）和等式。 (7.15) 是相同的，因此当且仅当所选网格使得属于 Ij 区间的粒子与属于 Ik 区间的粒子聚集并生成属于 Ii 区间的粒子时，该方法是内部一致的，其枢轴符合关系 zei = zej + zek。如果满足此关系，则数值格式对于 NDF 的所有矩而言是内部一致的。如果网格不尊重这种关系（通常是这种情况），则需要将新生成的粒子重新分配给相邻的粒子

\footnotetext[8]{这种移动枢轴技术与 QBMM 不同，稍后将在 7.4 节中介绍 QBMM。事实上，在QBMM中，正交近似的枢轴或节点不仅可以在区间Ii中自由移动，而且可以在相空间中自由移动，从而赋予该方法更大的灵活性。}
关键点。进行重新分配是为了保留分布的一些全局属性。例如，保留 r 阶矩会导致重新分配方程 (ζj + zek)r = γ(i) i (zei)r + γ(i) i+1 (zei+1)r，并且由于有两个自由度可用，因此选择 γ(i) i 和 γ(i) i+1 的值以保留分布的两个不同矩。例如，如果选择正确预测零阶矩 m0 和一阶矩 m1，则属于区间 Ii 的粒子将根据以下公式重新分配到关键点 zei 和 zei+1：

\begin{sourceblock}
\includegraphics[page=313,pagebox=mediabox,trim={153.50bp 515.97bp 67.56bp 156.46bp},clip,max width=0.94\textwidth,max height=0.38\textheight]{Marchisio2013.pdf}
\end{sourceblock}
很明显，这种方法仅对于 NDF 的两个选定时刻表现出内部一致，而对于其他时刻，该方法将表现不一致。 9 当然，一致性的定义可以扩展到 NDF 的任何特定属性，而不仅仅是通过它们的时刻 10（例如，参见 Ramkrishna (2000)）。每当系统中出现新粒子时，就需要执行此重新分配过程。不仅二阶点过程，而且零阶和一阶点过程都会产生新的粒子，并且为了保留分布的一些积分属性，有必要重新分配粒子以适合所选的网格。例如，在破碎的源项中，由属于区间 Ij 的粒子破碎生成的子粒子可能不会落入具有保留所选矩的枢轴的区间。可以证明（Ramkrishna，2000）破坏内部一致性的条件是

\begin{sourceblock}
\includegraphics[page=313,pagebox=mediabox,trim={170.12bp 336.58bp 67.56bp 329.11bp},clip,max width=0.94\textwidth,max height=0.38\textheight]{Marchisio2013.pdf}
\end{sourceblock}
其中 εi > ķj 且其中 b(Ψ|Ψ j) 是子分布函数。仅当 r = 0 时才完全满足此条件，在其他情况下，有必要进行类似于二阶点过程的粒子重新分配过程。让我们分析一下不同网格产生的不同方法。如果使用线性网格，则相空间被划分为多个由下界 xii 和上限 xii+1 界定的区间，并根据以下方案由枢轴点 ζi 表征： xii = 1 2 ih, xii+1 = 3 2ih, ζi = ih，对于 i ∈1, .。。 , M，其中 h 是离散化的大小。很明显，该方案对于聚合始终是一致的。事实上，任何时候，一个属于区间 Ij 且枢轴等于 zej 的粒子与一个属于区间 Ik 且枢轴等于 zk 的粒子聚合，所得的属于区间 Ii 的粒子都有一个枢轴 ζi = zej + zek，它是先前粒子的总和，正如内部一致性条件所要求的（见方程（7.16））。在聚集和破碎的情况下，内坐标在 xii 和 xii+1 之间（或者换句话说，属于区间 Ii）并以等于 zei 的关键点为特征的粒子 Ni 数量的演化方程变为

\begin{sourceblock}
\includegraphics[page=313,pagebox=mediabox,trim={124.37bp 116.07bp 67.56bp 550.87bp},clip,max width=0.94\textwidth,max height=0.38\textheight]{Marchisio2013.pdf}
\end{sourceblock}
9 这一事实意味着根据其定义计算出的矩将不同于通过使用方程 1 导出的输运方程计算出的矩。 （7.7）。

\footnotetext[10]{一般来说，分布 g(xi) 的任何全局属性都可以在重新分配过程中使用。然而，在本书分析的案例中，NDF 最重要的全局属性是矩；因此，我们将讨论限制在 g(xi) = xir 的情况。}
其中 βi, j = β(zei, zej) 是聚合核，ai = a(zei) 是破坏核，bi, j = b(zei|zej) 是子分布函数。 11 Batterham 等人。 (1981) 引入了基于恒定比率等于 2 的几何网格的离散化。在该几何网格中，内部坐标被离散化为由下界 xii 和上限 xii+1 界定的区间 Ii，并根据以下方案由关键点 ζi 表征： xii = 3 4zei, zei+1 = 3 2zei, zei+1 = 2zei，对于 i ε1, .。。，M−1。很明显，这种几何网格相对于线性网格的优势在于其粗糙度。例如，如果使用线性网格需要 100、1000 或 10 000 个类来覆盖某些特定范围，则使用几何网格则分别需要 6、10 和 13 个类来覆盖相同的范围。这种低得多的计算成本是以牺牲方案的内部一致性为代价获得的。因此，为了保留一些全局属性，必须使用上面描述的重新分配过程。如果选择零阶矩和一阶矩，如方程中所述。根据式 (7.16)，得到区间 Ii 发生纯聚集的粒子数变化方程： 12

\begin{sourceblock}
\includegraphics[page=314,pagebox=mediabox,trim={81.79bp 422.60bp 65.07bp 244.34bp},clip,max width=0.94\textwidth,max height=0.38\textheight]{Marchisio2013.pdf}
\end{sourceblock}
通过使用相同的方法，该方法可以扩展到断裂（Ramkrishna，2000）。一般来说，固定枢轴技术可以针对任何网格制定，同时保留阶数 r1 和 r2 分布的两个矩。由此产生的 Ii 区间粒子数演化方程为 (Kumar \& Ramkrishna, 1996a)

\begin{sourceblock}
\includegraphics[page=314,pagebox=mediabox,trim={103.71bp 291.70bp 65.07bp 337.62bp},clip,max width=0.94\textwidth,max height=0.38\textheight]{Marchisio2013.pdf}
\end{sourceblock}
其中

\begin{sourceblock}
\includegraphics[page=314,pagebox=mediabox,trim={150.84bp 185.09bp 65.07bp 419.56bp},clip,max width=0.94\textwidth,max height=0.38\textheight]{Marchisio2013.pdf}
\end{sourceblock}
以及

11 有趣的是，在牢不可破的单体聚集和断裂的情况下，方程： (7.18) 是严格的群体平衡方程。事实上，它描述了所有可能类别的颗粒系统（Vanni，2000a）。在这种情况下，内部坐标不能连续变化，而是只能根据构成单个聚集体的单体的数量呈现离散值。例如，由四个单体构成的聚集体的特征在于无量纲质量为四，并且当该聚集体与另一个颗粒（例如无量纲质量为五）聚集时，将生成无量纲质量等于九的聚集体，并且在破碎的情况下类似。因此，很明显，在具有线性网格和牢不可破的单分散单体的情况下，由于聚集和破碎而导致的源项不能生成不落入特定区间的粒子。

\footnotetext[12]{方程 (7.19) 也由 Hounslow 等人得出。 (1988) 在修正了 Batterham 等人提出的原始方程之后。 （1981）。然而，这个方程只不过是应用于比率等于2的几何网格并叠加零阶矩和一阶矩守恒（以粒子质量作为内坐标）之后的定枢点技术。}
这是固定枢轴技术最通用的表述，其中包含大多数先前开发的类方法的特殊情况。例如，马夏尔等人。 (1988) 开发了一种通用方法，其中聚集（和断裂）被视为化学反应，并确定化学计量系数，以保留 NDF 的一阶矩。利斯特等人。 (1995) 开发了一种基于几何网格的方法，其中比率为 21/q，其中 q 是整数。 13 在其他可用方法中，值得一提的是 Vanni (1999) 为纯破坏的通用网格开发的方案。尽管固定枢点技术是最流行的离散化方法，但它存在数密度过度预测的问题，特别是在粒子聚集的情况下。为了克服这个问题，制定了移动枢轴技术，该技术比固定枢轴技术提供更多时刻的一致性附加属性。我们建议读者参阅 Kumar 和 Ramkrishna (1996b) 的原著以了解更多详细信息。相反，有趣的是，我们要多花几句话来介绍单元平均技术（Kumar 等人，2006、2007、2008b、2009），该技术的开发是为了克服固定枢轴技术遇到的一些问题。使用固定枢轴技术，当形成新粒子时，根据通过施加两个矩守恒而先验计算出的一些重新分配因子 γ(i) i 和 γ(i) i+1，将其重新分配到相邻部分。然而，很明显，这种重新分配完全是任意的，并且纯粹基于网格结构。改进该方法的一个简单方法是根据两个选定时刻的净变化率自适应地计算这些重新分配因子。当选择零阶矩和一阶矩时，该过程基于对零阶矩变化率的评估，由于 Ii 区间内粒子的总出生率（Kumar \& Ramkrishna，1997），

\begin{sourceblock}
\includegraphics[page=315,pagebox=mediabox,trim={174.04bp 314.53bp 170.24bp 351.69bp},clip,max width=0.94\textwidth,max height=0.38\textheight]{Marchisio2013.pdf}
\end{sourceblock}
以及对 Ii 区间内一阶矩的变化率的评估，

\begin{sourceblock}
\includegraphics[page=315,pagebox=mediabox,trim={158.58bp 246.69bp 152.31bp 419.53bp},clip,max width=0.94\textwidth,max height=0.38\textheight]{Marchisio2013.pdf}
\end{sourceblock}
因此，Ii 区间内所有新生粒子的内坐标平均值可计算为 ζi = Magg i /Bagg

\begin{sourceblock}
\includegraphics[page=315,pagebox=mediabox,trim={65.02bp 182.97bp 67.55bp 488.83bp},clip,max width=0.94\textwidth,max height=0.38\textheight]{Marchisio2013.pdf}
\end{sourceblock}
i+1 = (zi - zi)/(zi+1 - zi)。通过总结和收集 Ii 区间内粒子诞生和死亡的所有贡献，并通过引入略有不同的符号，得到以下全局表达式：

\begin{sourceblock}
\includegraphics[page=315,pagebox=mediabox,trim={127.26bp 89.83bp 94.87bp 549.40bp},clip,max width=0.94\textwidth,max height=0.38\textheight]{Marchisio2013.pdf}
\end{sourceblock}
dNi dt = 巴格

\footnotetext[13]{很明显，随着 q 的增加，离散化变得越来越精细，而当 q = 1 时，该方法又变成了 Hounslow 等人描述的方法。 （1988）。}
其中 Heaviside 阶跃函数定义为

\begin{sourceblock}
\includegraphics[page=316,pagebox=mediabox,trim={199.78bp 573.79bp 198.28bp 76.02bp},clip,max width=0.94\textwidth,max height=0.38\textheight]{Marchisio2013.pdf}
\end{sourceblock}
且 λ±(xi) = ( xi − z i ± 1 )/( z i − z i ± 1 )。类似地，由于破损事件导致的区间 Ii 中零阶矩的净变化率可以写成离散形式：

\begin{sourceblock}
\includegraphics[page=316,pagebox=mediabox,trim={187.58bp 498.06bp 178.79bp 165.70bp},clip,max width=0.94\textwidth,max height=0.38\textheight]{Marchisio2013.pdf}
\end{sourceblock}
而破损事件间隔中的一阶时刻的净变化率为

\begin{sourceblock}
\includegraphics[page=316,pagebox=mediabox,trim={185.45bp 435.23bp 174.21bp 228.53bp},clip,max width=0.94\textwidth,max height=0.38\textheight]{Marchisio2013.pdf}
\end{sourceblock}
由此可以计算出所观察的区间内的内坐标平均值为 ζi = Mbreak i /Bbreak i，从而得出非常相似的破损方程。对更多细节感兴趣的读者可以参考 Kumar 等人的工作。 （2008b）。基于相同概念的分段方法已在蒸发喷雾领域独立开发（Greenberg 等人，1993），自然产生非常相似的方程。要了解该领域的概述，读者可以参考 Laurent 等人的著作。 （2004）。

\paragraph*{零阶点过程}
让我们以成核为例，它是一个零级点过程，因为它以独立于 NDF 的速率将新粒子引入系统。零阶点过程通常通过计算对每个类的贡献来以非常简单的方式处理。例如，如果仅发生成核，则 PBE 变为

\begin{sourceblock}
\includegraphics[page=316,pagebox=mediabox,trim={232.48bp 238.00bp 65.07bp 436.36bp},clip,max width=0.94\textwidth,max height=0.38\textheight]{Marchisio2013.pdf}
\end{sourceblock}
其中J是成核率，即每单位体积产生的新粒子的数量，其内坐标在 和 之间 + d 之间。离散化 PBE 中的源项简单地变为

\begin{sourceblock}
\includegraphics[page=316,pagebox=mediabox,trim={205.96bp 166.05bp 65.07bp 499.64bp},clip,max width=0.94\textwidth,max height=0.38\textheight]{Marchisio2013.pdf}
\end{sourceblock}
其中 Ji 表示引入由 xii 和 xii+1 界定的内坐标区间 Ii 中的新粒子的数量。

\paragraph*{相空间漂移项}
漂移项表示内部坐标 Σ 由于变化率 ̇ Σ 的连续变化：

\begin{sourceblock}
\includegraphics[page=316,pagebox=mediabox,trim={213.68bp 50.16bp 65.07bp 622.58bp},clip,max width=0.94\textwidth,max height=0.38\textheight]{Marchisio2013.pdf}
\end{sourceblock}
因此，对分区进行积分后，我们得到

\begin{sourceblock}
\includegraphics[page=317,pagebox=mediabox,trim={174.67bp 589.38bp 67.56bp 76.31bp},clip,max width=0.94\textwidth,max height=0.38\textheight]{Marchisio2013.pdf}
\end{sourceblock}
其中指出属于 Ii 类的粒子的累积等于来自周围区间 Ii−1 和 Ii+1 的粒子的净通量。值得注意的是，方程右侧的项。 (7.26) 涉及在点 Ψi+1 和 Ψi 处评估的 NDF。不幸的是，NDF 是未知的，但它可以通过每个区间内的多项式来近似。当NDF由每个区间内的零阶多项式描述时（如式（7.7）所示），它由Ii区间内的常数值Ni来近似。事实上，这是一个不连续的近似，因为，例如， xi = xii 对应于两个区间之间的边界，其中 NDF 在 Ii−1 区间中由 Ni−1 近似，在 Ii 区间中由 Ni 近似。当评估每个单一区间的下边界和上边界处的粒子通量时，这种区别变得非常重要。最简单的方法是一阶迎风方案，其中，对于正 ̇ xi(xii)，n(t, xii) 的值取等于 Ni−1；对于负 ̇ xi(xii)，取等于 Ni。可以证明，这种方法是一阶的，14 正如 Laurent 等人所证明的那样。 （2004）。在 ̇ xi(xii) 为正的情况下，属于 Ii 区间的粒子数密度方程为

\begin{sourceblock}
\includegraphics[page=317,pagebox=mediabox,trim={175.62bp 377.89bp 67.56bp 287.91bp},clip,max width=0.94\textwidth,max height=0.38\textheight]{Marchisio2013.pdf}
\end{sourceblock}
在几何网格的情况下，如果 s 是几何网格的比率（即 s = Ψi+1/Ψi），我们可以写成

\begin{sourceblock}
\includegraphics[page=317,pagebox=mediabox,trim={169.45bp 323.73bp 67.56bp 347.30bp},clip,max width=0.94\textwidth,max height=0.38\textheight]{Marchisio2013.pdf}
\end{sourceblock}
但这是一阶近似方案，会受到数值扩散的影响。事实上，如果使用此方案处理表现出急剧增加或减少的内坐标依赖性的 NDF，则在增长（或根据 ψ 符号缩小）的同时，它将在相空间中扩散，就好像存在附加扩散项一样。为了克服数值扩散，可以开发高阶方案，其中 Ni 被写为 Ni−1 和 Ni+1 的函数：

\begin{sourceblock}
\includegraphics[page=317,pagebox=mediabox,trim={148.02bp 213.50bp 67.56bp 452.72bp},clip,max width=0.94\textwidth,max height=0.38\textheight]{Marchisio2013.pdf}
\end{sourceblock}
正如 Marchal 等人所描述的。 （1988）和大卫等人。 （1991）。通过计算 Ni 的变化率，可以得出另一种二阶方法：

\begin{sourceblock}
\includegraphics[page=317,pagebox=mediabox,trim={148.63bp 147.06bp 67.56bp 523.98bp},clip,max width=0.94\textwidth,max height=0.38\textheight]{Marchisio2013.pdf}
\end{sourceblock}
并选择常数 a、b 和 c，以便正确预测三个矩。如果选择 m0、m1 和 m3，我们得到

\begin{sourceblock}
\includegraphics[page=317,pagebox=mediabox,trim={129.68bp 82.57bp 67.56bp 591.84bp},clip,max width=0.94\textwidth,max height=0.38\textheight]{Marchisio2013.pdf}
\end{sourceblock}
\footnotetext[14]{在近似 NDF 的泰勒展开式中第一个被排除的项是一阶导数项。}
文献中提供了替代的高阶方法，其中大多数都具有有限体积代码（例如 QUICK 和 MUSCL）中使用的空间离散化方案的显著特征。其中一些基于多项式重建，并进行了一些细微修改，以确保 NDF 的正性，如 Laurent (2006) 中所解释的，而另一些则基于香农熵的最大化（例如，参见 Massot 等人 (2010)）。粒子质量的离散化，如方程中所述。 （7.8），对于漂移项的处理特别受欢迎。在这种情况下，属于 Ii 区间的粒子质量的变化率变为

\begin{sourceblock}
\includegraphics[page=318,pagebox=mediabox,trim={191.34bp 496.09bp 183.54bp 169.59bp},clip,max width=0.94\textwidth,max height=0.38\textheight]{Marchisio2013.pdf}
\end{sourceblock}
如果使用粒子体积的通用表达式 Vp(xi) = αxiγ，并按分部积分规则求解积分，则可得到以下表达式：

\begin{sourceblock}
\includegraphics[page=318,pagebox=mediabox,trim={162.39bp 431.77bp 154.61bp 233.91bp},clip,max width=0.94\textwidth,max height=0.38\textheight]{Marchisio2013.pdf}
\end{sourceblock}
如果 NDF 由 Ii 区间内的零阶多项式近似，我们得到

\begin{sourceblock}
\includegraphics[page=318,pagebox=mediabox,trim={199.33bp 377.97bp 196.88bp 291.18bp},clip,max width=0.94\textwidth,max height=0.38\textheight]{Marchisio2013.pdf}
\end{sourceblock}
那么，如果采用一阶迎风格式，则最终方程为

\begin{sourceblock}
\includegraphics[page=318,pagebox=mediabox,trim={100.99bp 316.65bp 65.07bp 340.11bp},clip,max width=0.94\textwidth,max height=0.38\textheight]{Marchisio2013.pdf}
\end{sourceblock}
其中右侧第一项表示离开该区间的颗粒的质量，因为它们正在生长到下一个区间，第二项表示来自较小尺寸区间的颗粒的质量，最后一项表示由于来自连续相的通量而导致的区间内质量的增加。一旦指定了 ̇ xi ( xi ) 的表达式，方程的最后一项就可以得到。 (7.32) 可以被明确地写成，导致 Mi 的演化方程被写成涉及 Mi 的两项和涉及 Mi−1 的一项的总和。显然，离散化的精细程度决定了解的精度。很容易想象，在非常精细的网格的情况下，不同的方法给出了非常相似的性能，但其中一些方法可能会遇到离散数或质量密度的负（无法实现）解的可能性。当解决不连续性问题时，这种情况尤其严重：观察到偏离负值的不合理振荡，正如 Kumar 和 Ramkrishna (1997) 所解释的那样。此外，类方法和截面方法在数值上等价于受数值扩散影响的有限差分和有限体积方法。后者可以通过使用高阶方案来减少，但不能完全避免。

\paragraph*{相空间扩散项}
扩散项表示内坐标 ψ 随扩散率 D 的连续变化，

\begin{sourceblock}
\includegraphics[page=318,pagebox=mediabox,trim={213.45bp 50.21bp 65.07bp 623.17bp},clip,max width=0.94\textwidth,max height=0.38\textheight]{Marchisio2013.pdf}
\end{sourceblock}
因此，在对分区进行积分后，我们得到

\begin{sourceblock}
\includegraphics[page=319,pagebox=mediabox,trim={168.67bp 584.38bp 67.56bp 73.89bp},clip,max width=0.94\textwidth,max height=0.38\textheight]{Marchisio2013.pdf}
\end{sourceblock}
其中指出属于 Ii 类的粒子的累积等于来自周围区间 Ii−1 和 Ii+1 的粒子的净扩散通量。值得注意的是，方程右侧的项。 (7.34) 涉及在点 xii+1 和 xii 处评估的 NDF 的梯度。由于 NDF 未知，通常其梯度是根据可用信息重建的，例如采用以下近似值：ni = Ni/(xii − xii−1)。在 D 为常数的情况下，梯度结果为

\begin{sourceblock}
\includegraphics[page=319,pagebox=mediabox,trim={144.70bp 463.18bp 67.56bp 195.13bp},clip,max width=0.94\textwidth,max height=0.38\textheight]{Marchisio2013.pdf}
\end{sourceblock}
得出相空间扩散项的简单方案。可以通过与漂移项离散化所提出的参数非常相似的参数来获得改进的版本。

\subsection*{7.2.2 双变量和多变量 PBE}\phantomsection\addcontentsline{toc}{subsection}{7.2.2 双变量和多变量 PBE}
对于两个具有 NDF n(t, Ψ1, Ψ2) 的内坐标的情况，PBE 具有以下形式：

\begin{sourceblock}
\includegraphics[page=319,pagebox=mediabox,trim={131.47bp 346.61bp 67.56bp 322.57bp},clip,max width=0.94\textwidth,max height=0.38\textheight]{Marchisio2013.pdf}
\end{sourceblock}
其中左侧第二项和第三项代表相空间中的平流，而右侧项代表相空间中的扩散和不连续点过程。和以前一样，可以使用通用源项编写 PBE：

\begin{sourceblock}
\includegraphics[page=319,pagebox=mediabox,trim={130.07bp 270.22bp 67.56bp 398.96bp},clip,max width=0.94\textwidth,max height=0.38\textheight]{Marchisio2013.pdf}
\end{sourceblock}
对不同 CM 在该方程中的应用的详细讨论感兴趣的读者可以参考专门文献（Immanuel \& Doyle, 2005; Lau-renzi et al., 2002; Qamar \& Warnecke, 2007; Vale \& McKenna, 2005; Xiong \& Pratsinis, 1991, 1993）。下面仅报告一个简短的总结。15 现在让我们讨论针对单变量 PBE 提出的离散化双变量问题的简单几何扩展。给定由 Ψ(i) 1 和 Ψ(i+1) 1 界定的内坐标 Ψ1 的区间，记为 I(i) 1，以及由 Ψ(j) 2 界定的第二内坐标 Ψ2 的区间，并且
15 类方法可以通过遵循两种替代原理应用于双变量情况。第一个是本质上的几何，其中每个内部坐标都被相同地对待。由此产生的二维笛卡尔网格由平行于内部坐标轴的线的交点定义的矩形形成（Kumar 等人，2008a；Vale 和 McKenna，2005）。很明显，在这种情况下，任何事件形成的粒子都必须分布在由四个不同枢轴定义的四个不同间隔上。对于具有 n 个内坐标的一般多维 PBE，这将导致分布在 2n 个间隔上。相反，第二种方法基于一个新概念：以 n 个内坐标为特征的粒子群的特征为需要保留的最少 n + 1 个数量（与之前方法中的 2n 个数量相比）。事实上，众所周知，在二维空间中，三角形是可用于表示相空间的离散形式的最简单的几何对象。类似地，四面体可以用于三维空间，等等。正如 Chakraborty \& Kumar (2007) 和 Nandanwar \& Kumar (2008) 所强调的，这些推测提出了一种离散化相空间的最低有效方法。

xi(j+1) 2，记为 I(j) 2，可以将属于 I(i) 1 和 I(j) 2 区间的粒子数定义为

\begin{sourceblock}
\includegraphics[page=320,pagebox=mediabox,trim={197.95bp 579.67bp 65.07bp 84.57bp},clip,max width=0.94\textwidth,max height=0.38\textheight]{Marchisio2013.pdf}
\end{sourceblock}
1 和 ze(j) 同样在这种情况下，可以定义关键点 ze(i) 2，并且 NDF 的结果表示为

\begin{sourceblock}
\includegraphics[page=320,pagebox=mediabox,trim={155.44bp 506.52bp 65.07bp 161.37bp},clip,max width=0.94\textwidth,max height=0.38\textheight]{Marchisio2013.pdf}
\end{sourceblock}
其中 M1 和 M2 是用于两个内坐标的部分或间隔的数量。 Nij 的演化方程得出

\begin{sourceblock}
\includegraphics[page=320,pagebox=mediabox,trim={187.80bp 430.69bp 65.07bp 233.55bp},clip,max width=0.94\textwidth,max height=0.38\textheight]{Marchisio2013.pdf}
\end{sourceblock}
类方法仅适用于解决简单问题的双变量 PBE，通常将其扩展到多变量问题是相当复杂的。下面讨论一些选定的例子来说明这些困难。

\paragraph*{点流程}
在两个内坐标的情况下，不连续点过程可以呈现不同的形式，具体取决于所考虑的现象（例如聚集/合并与断裂）和内坐标的性质。让我们通过一个例子来说明主要问题：聚合。此外，让我们考虑聚集混合问题，即两种不同类型粒子的聚集。在这种特殊情况下，内坐标在聚合过程中是相加的，如 Matsoukas 等人所述。 （2006、2009）。在这些条件下，PBE 为

\begin{sourceblock}
\includegraphics[page=320,pagebox=mediabox,trim={130.52bp 204.62bp 127.37bp 434.74bp},clip,max width=0.94\textwidth,max height=0.38\textheight]{Marchisio2013.pdf}
\end{sourceblock}
1 和 I(j) 2 间隔为两个内坐标，并应用等式中的定义。 (7.40)，我们得到 通过对 I(i) 进行积分

\begin{sourceblock}
\includegraphics[page=320,pagebox=mediabox,trim={98.66bp 98.44bp 90.62bp 535.02bp},clip,max width=0.94\textwidth,max height=0.38\textheight]{Marchisio2013.pdf}
\end{sourceblock}
{\small 对于单变量情况，当粒子在聚集事件后形成时，必须将其分配给相邻单元（双变量情况下为四个）。当使用固定枢轴技术将粒子分配给单元时，得到以下表达式：\par}
dNij

\begin{sourceblock}
\includegraphics[page=321,pagebox=mediabox,trim={132.77bp 550.69bp 129.92bp 68.98bp},clip,max width=0.94\textwidth,max height=0.38\textheight]{Marchisio2013.pdf}
\end{sourceblock}
其中后两个求和分别对索引对 (k, l) 和 (m, n) 进行，M1 和 M2 是两个内坐标离散化中使用的区间数，并且 βklmn = β

\begin{sourceblock}
\includegraphics[page=321,pagebox=mediabox,trim={231.87bp 496.43bp 189.91bp 179.14bp},clip,max width=0.94\textwidth,max height=0.38\textheight]{Marchisio2013.pdf}
\end{sourceblock}
。重分配系数 ηij 是通过叠加 NDF 的一些矩守恒来计算的。例如，如果使用分布的前三个矩（即 m0,0、m1,0 和 m0,1），则获得以下表达式：

\begin{sourceblock}
\includegraphics[page=321,pagebox=mediabox,trim={74.31bp 313.19bp 67.56bp 245.35bp},clip,max width=0.94\textwidth,max height=0.38\textheight]{Marchisio2013.pdf}
\end{sourceblock}
\begin{sourceblock}
\includegraphics[page=321,pagebox=mediabox,trim={229.25bp 290.99bp 246.02bp 395.11bp},clip,max width=0.94\textwidth,max height=0.38\textheight]{Marchisio2013.pdf}
\end{sourceblock}
0 Ψk 1Ψl 2n dΨ。方程中的系数。 （7.41）是通过施加 m0,0、m1,0 和 m0,1 守恒来计算的，但理论上可以容纳任何其他矩集。类似地，通过保持相同矩集 m0,0、m1,0 和 m0,1 的守恒，也可以采用单元平均技术来计算系数： 其中矩定义为 mk,l =
dNij

\begin{sourceblock}
\includegraphics[page=321,pagebox=mediabox,trim={71.65bp 56.06bp 68.57bp 487.19bp},clip,max width=0.94\textwidth,max height=0.38\textheight]{Marchisio2013.pdf}
\end{sourceblock}
其中 H(xi) 是 Heaviside 阶跃函数，聚合源项为

\begin{sourceblock}
\includegraphics[page=322,pagebox=mediabox,trim={137.48bp 578.27bp 163.26bp 84.30bp},clip,max width=0.94\textwidth,max height=0.38\textheight]{Marchisio2013.pdf}
\end{sourceblock}
βklmnNkmNln, Bij =

\begin{sourceblock}
\includegraphics[page=322,pagebox=mediabox,trim={231.95bp 548.70bp 137.62bp 133.01bp},clip,max width=0.94\textwidth,max height=0.38\textheight]{Marchisio2013.pdf}
\end{sourceblock}
其他参数由 ζ 定义

\begin{sourceblock}
\includegraphics[page=322,pagebox=mediabox,trim={118.03bp 501.51bp 84.63bp 161.06bp},clip,max width=0.94\textwidth,max height=0.38\textheight]{Marchisio2013.pdf}
\end{sourceblock}
Mij,1 =

\begin{sourceblock}
\includegraphics[page=322,pagebox=mediabox,trim={116.34bp 448.69bp 82.96bp 213.88bp},clip,max width=0.94\textwidth,max height=0.38\textheight]{Marchisio2013.pdf}
\end{sourceblock}
Mij,2 =
以及

\begin{sourceblock}
\includegraphics[page=322,pagebox=mediabox,trim={173.01bp 392.49bp 170.55bp 264.69bp},clip,max width=0.94\textwidth,max height=0.38\textheight]{Marchisio2013.pdf}
\end{sourceblock}
该方法可以很容易地扩展到破损的情况，具体细节请读者参阅专业文献。重要的是要报告通过使用例如扩展网格来减少和优化计算成本的许多努力，如 Chiney 等人的工作所示。 （2010）。

\paragraph*{相空间漂移项}
在纯漂移的情况下，源项可以写为

\begin{sourceblock}
\includegraphics[page=322,pagebox=mediabox,trim={136.90bp 244.68bp 65.07bp 421.78bp},clip,max width=0.94\textwidth,max height=0.38\textheight]{Marchisio2013.pdf}
\end{sourceblock}
与对一个内部坐标所做的类似，这项可以写为

\begin{sourceblock}
\includegraphics[page=322,pagebox=mediabox,trim={217.19bp 193.71bp 65.07bp 488.73bp},clip,max width=0.94\textwidth,max height=0.38\textheight]{Marchisio2013.pdf}
\end{sourceblock}
其中 f ij 1 和 f ij 2 分别是由于第一和第二内坐标的连续变化率而产生的粒子净通量。与有限体积方案的类比是显而易见的，如果使用一阶迎风方案，我们可以写成

\begin{sourceblock}
\includegraphics[page=322,pagebox=mediabox,trim={177.11bp 98.72bp 65.07bp 557.91bp},clip,max width=0.94\textwidth,max height=0.38\textheight]{Marchisio2013.pdf}
\end{sourceblock}
{\small 其中 ai、bi、ci 和 di 是根据 Puel 等人的解释从网格计算的系数。 （2003a）。显然该格式受数值扩散的影响，也可以使用更高阶的格式。\par}
\paragraph*{相空间扩散项}
在纯扩散的情况下，源项可以写为

\begin{sourceblock}
\includegraphics[page=323,pagebox=mediabox,trim={147.33bp 535.34bp 67.56bp 97.74bp},clip,max width=0.94\textwidth,max height=0.38\textheight]{Marchisio2013.pdf}
\end{sourceblock}
与对一个内部坐标所做的类似，这项可以写为

\begin{sourceblock}
\includegraphics[page=323,pagebox=mediabox,trim={243.91bp 487.51bp 67.56bp 194.93bp},clip,max width=0.94\textwidth,max height=0.38\textheight]{Marchisio2013.pdf}
\end{sourceblock}
S ij = dij
其中 dij 1 和 dij 2 是粒子的扩散通量，可以使用有限体积方案中使用的标准方法进行描述。

\subsection*{7.2.3 碰撞KE}\phantomsection\addcontentsline{toc}{subsection}{7.2.3 碰撞KE}
均匀系统的动力学方程（KE）可以通过将方程中出现的所有源项分组在一起来编写。 (7.1) 右侧：

\begin{sourceblock}
\includegraphics[page=323,pagebox=mediabox,trim={161.65bp 356.48bp 67.56bp 315.27bp},clip,max width=0.94\textwidth,max height=0.38\textheight]{Marchisio2013.pdf}
\end{sourceblock}
其中由不连续事件引起的源项用 C 表示（参见第 6 章）。在这里，我们将注意力限制在没有外力和相空间扩散的系统上，16 使得 ̇ 和 D 为零。由此产生的碰撞 KE 描述了具有恒定尺寸和质量的粒子之间的动量交换，其中 NDF 是粒子速度 ε = (v1, v2, v3) 的函数。玻尔兹曼硬球碰撞积分为

\begin{sourceblock}
\includegraphics[page=323,pagebox=mediabox,trim={144.64bp 252.59bp 67.56bp 413.87bp},clip,max width=0.94\textwidth,max height=0.38\textheight]{Marchisio2013.pdf}
\end{sourceblock}
式中， ψ 和 ψ* 分别为测试粒子和场粒子碰撞前的粒子速度， ψ′ 和 ψ′ ＊ 为碰撞后粒子速度， B 为碰撞核， g = ψ*− ψ 为场粒子相对于测试粒子的相对速度， x 为沿瞬时碰撞过程中连接两个粒子中心的方向从场指向测试粒子的单位向量。碰撞粒子 17 的线性动量和动能守恒导致碰撞前和碰撞后速度之间存在以下关系：

\begin{sourceblock}
\includegraphics[page=323,pagebox=mediabox,trim={210.75bp 131.76bp 67.56bp 535.51bp},clip,max width=0.94\textwidth,max height=0.38\textheight]{Marchisio2013.pdf}
\end{sourceblock}
16 对于 KE，相空间扩散通常与由与外部相的随机碰撞引起的布朗运动相关。相比之下，C 是由于同类型粒子之间的自碰撞造成的。将布朗运动表示为扩散项意味着与外相的碰撞发生的时间尺度远小于自碰撞的时间尺度。

\footnotetext[17]{弹性碰撞期间动能守恒，而非弹性碰撞期间动能的一小部分会不可逆地损失。恢复系数量化了这两种情况，更多细节可以在第 6 章中找到。}
碰撞核取决于粒子之间相互作用的类型，这里我们关注硬球碰撞核：B(g, x) = d2 p |g · x|。例如，另一种核是第 6 章中描述的麦克斯韦核。求解碰撞 KE 的分段方法和类方法通常称为离散速度方法（DVM）。这些方法基于将速度空间离散化为由有限数量的点构成的网格的简单思想。现有方法的特点是不同的网格结构（Aristov，2001）。例如，格子玻尔兹曼方法将速度空间离散成具有恒定晶格尺寸的规则立方晶格（Li-Shi，2000），而其他方法则采用不同的离散化方案（Monaco＆Preziosi，1990）。通过使用与 PBE 类似的方法，可以将 Ni 定义为速度为 xii 的粒子的数密度，并且离散化的 KE 变为

\begin{sourceblock}
\includegraphics[page=324,pagebox=mediabox,trim={221.38bp 472.56bp 65.07bp 208.44bp},clip,max width=0.94\textwidth,max height=0.38\textheight]{Marchisio2013.pdf}
\end{sourceblock}
其中 Gi 和 Li 分别表示速度等于 xii 的粒子的增益率和损失率。在这种离散化表示中，当速度为 xii 的粒子与另一个速度为 xij 的粒子碰撞时，它会丢失，并且碰撞会产生一对不同的碰撞后速度 (ψk, ψl)。类似地，两个速度为 Ψk 和 Ψl 的粒子碰撞可能会产生一个速度为 Ψi 的粒子。这些过程产生离散碰撞算子的表达式：

\begin{sourceblock}
\includegraphics[page=324,pagebox=mediabox,trim={179.27bp 359.43bp 65.07bp 308.42bp},clip,max width=0.94\textwidth,max height=0.38\textheight]{Marchisio2013.pdf}
\end{sourceblock}
其中 Aij kl 和 Akl ij 是粒子速度类别之间的碰撞频率或转移概率，并且仅对可能的对 (k, l) 执行第二次求和。离散化方案必须与原始KE一致。例如，由于在两个粒子 i 和 j 的一般碰撞中，产生粒子 k 和 l，由于动量守恒 xii + xij = xik + xil 和动能 xi2，四个相应的速度不是任意的。

\begin{sourceblock}
\includegraphics[page=324,pagebox=mediabox,trim={103.30bp 272.41bp 327.02bp 411.50bp},clip,max width=0.94\textwidth,max height=0.38\textheight]{Marchisio2013.pdf}
\end{sourceblock}
升。很容易证明（例如，请参见 Chapman \& Cowling (1961) 的图形推导），对于满足这两个条件的碰撞，碰撞前和碰撞后相对速度矢量的长度是相同的， | xi i − xi j | = |ψk − ψl|，并且由于角动量守恒，非旋转粒子的碰撞必须发生在平面内。因此，离散化的结构必须设计为产生尊重质量、动量和能量守恒性质的转移概率。此外，非常希望离散化方法能够收敛到连续体 KE，以实现越来越精细的离散化（Gorsch，1999）。当 KE 中包含完整碰撞积分时，这一点尤其重要，而当碰撞积分被动力学模型（即 BGK（Bhatnagar 等人，1954）和 ES-BGK（Holway，1966）模型）替代时，这不是问题。例如，当粒子被表示为具有有限半径的点质量时，属于 j 类的粒子只有在属于 i 类的粒子恰好位于横截面积为 18 S c 且高度为 | xi i − xi j | 的圆柱体中时才能够与 i 类粒子发生碰撞。 dt。由此产生的转移概率是 Akl

\begin{sourceblock}
\includegraphics[page=324,pagebox=mediabox,trim={354.50bp 117.09bp 68.07bp 566.85bp},clip,max width=0.94\textwidth,max height=0.38\textheight]{Marchisio2013.pdf}
\end{sourceblock}
ij。假设在碰撞中所有可能的对都以相同的概率获得，则结果是转移概率是非负常数（值为零对应于无法实现的碰撞），其相对于

\footnotetext[18]{对于硬球，该横截面积对应于半径等于 dp 的圆的表面。}
索引 i、j 和 k、l：Akl ij = Akl ji = Alk ij ≥0。作为离散公式的结果，我们还确保遵守微观可逆性关系，这意味着粒子 i 和 j 的碰撞产生粒子 k 和 l 的概率与逆碰撞相同。利用微观可逆性关系并对指数 j、k 和 l 而非对对进行求和，离散动力学方程变为

\begin{sourceblock}
\includegraphics[page=325,pagebox=mediabox,trim={179.68bp 540.27bp 67.56bp 127.66bp},clip,max width=0.94\textwidth,max height=0.38\textheight]{Marchisio2013.pdf}
\end{sourceblock}
其中 i = 1, 2, .。。添加了 、M 和系数 1/2，以避免对同一对粒子进行两次计数。正如已经预期的那样，确保离散化方案尊重一些基本属性非常重要，并且可以证明等式中的方案。 (7.52) 具有原始连续 KE 的许多性质，例如守恒定律和具有麦克斯韦平衡的 H 定理（Aristov，2001）。等式的比较(7.52) 与等式。式（7.13）显示了两种离散化之间的异同。例如，在粒子聚集过程中，碰撞后形成单个粒子，因此仅涉及三个索引（两个碰撞前粒子的索引以及子粒子或碰撞后粒子的索引）。相反，在粒子碰撞期间（仅交换动量和动能），事件使粒子数量保持不变，因此必须采用四个指数。在不同的简化假设下开发了不同的技术，在下文中我们将讨论两种不同的情况，对应于各向同性情况和一般的非各向同性情况。在 PBE 分段方法的并行中，针对这两种不同情况说明的方法与为单变量和多元 PBE 求解而开发的数值方法有许多相似之处。

\paragraph*{各向同性情况}
在许多问题中，我们对粒子速度矢量的分布不感兴趣，但在某些条件下，粒子速度大小的分布可以充分描述系统。其中一种情况是各向同性假设成立。在这个简化的假设下，速度矢量的方向失去了任何重要性，因此 可以用标量 ψ = | ψ| 代替。 = (v2 1 + v2 2 + v2 3)1/2 对应于粒子速度大小（或粒子速度）。一般来说，如果 NDF 可以用速度矢量来表示，则速度分布函数 n(xi) 是已知的。以速度大小表示的 NDF（表示为 nxi(xi)）可以通过对所有可能的方向进行积分来明确地写出：

\begin{sourceblock}
\includegraphics[page=325,pagebox=mediabox,trim={171.45bp 143.99bp 67.56bp 522.21bp},clip,max width=0.94\textwidth,max height=0.38\textheight]{Marchisio2013.pdf}
\end{sourceblock}
其中 ε = (xi sin θ cos φ, ϋ sin θ sin φ, ϋ cos θ)，θ 和 φ 分别是球坐标的极角和方位角。后者用于表示碰撞积分中出现的三维积分（例如，参见第 6 章）。方程（7.53）始终有效，但可以在各向同性假设下进一步简化，因为 n(xi) 在半径 xi 恒定的球体上是恒定的。这导致

\begin{sourceblock}
\includegraphics[page=325,pagebox=mediabox,trim={202.17bp 51.49bp 67.56bp 633.42bp},clip,max width=0.94\textwidth,max height=0.38\textheight]{Marchisio2013.pdf}
\end{sourceblock}
其中 | xi |由于速度矢量的方向并不重要，因此在 NDF 中被限制为等于 ψ。从今以后，我们将滥用该符号，表示 n(ψ) = n(ψ)|ψ|=ψ。当初始 NDF 各向同性时，方程 (7.54) 在平衡状态下成立。为了说明动力学方程的离散化通常是如何进行的，必须引入一些额外的简化。将碰撞积分写为很方便

\begin{sourceblock}
\includegraphics[page=326,pagebox=mediabox,trim={129.19bp 531.38bp 65.07bp 134.82bp},clip,max width=0.94\textwidth,max height=0.38\textheight]{Marchisio2013.pdf}
\end{sourceblock}
其中积分现在扩展到整个立体角，但函数 S 定义为

\begin{sourceblock}
\includegraphics[page=326,pagebox=mediabox,trim={205.66bp 476.29bp 204.40bp 187.54bp},clip,max width=0.94\textwidth,max height=0.38\textheight]{Marchisio2013.pdf}
\end{sourceblock}
通过写出关于立体角的第二个积分，我们得到

\begin{sourceblock}
\includegraphics[page=326,pagebox=mediabox,trim={93.89bp 402.36bp 65.07bp 243.14bp},clip,max width=0.94\textwidth,max height=0.38\textheight]{Marchisio2013.pdf}
\end{sourceblock}
p|ψ*cos(αy)− ψ cos(αx)|，其中 αx 是 ψ 和 x 形成的角度，αy 是 ψ* 和 x 形成的角度。如果将立体角写成极坐标，dx = sin(αx)dαx dβx 且 dx*= sin(αy)dαy dβy，则碰撞动力学方程为

\begin{sourceblock}
\includegraphics[page=326,pagebox=mediabox,trim={117.70bp 316.75bp 65.07bp 334.36bp},clip,max width=0.94\textwidth,max height=0.38\textheight]{Marchisio2013.pdf}
\end{sourceblock}
引入变量 x = cos(αx) 和 y = cos(αy)，并回顾由于各向同性，NDF 不是 αx、βx、αy 和 βy 的函数，方程： (7.57) 简化为

\begin{sourceblock}
\includegraphics[page=326,pagebox=mediabox,trim={76.17bp 249.06bp 65.07bp 417.14bp},clip,max width=0.94\textwidth,max height=0.38\textheight]{Marchisio2013.pdf}
\end{sourceblock}
此外，由于各向同性，碰撞算子是对称的，导致

\begin{sourceblock}
\includegraphics[page=326,pagebox=mediabox,trim={92.84bp 195.95bp 65.07bp 470.25bp},clip,max width=0.94\textwidth,max height=0.38\textheight]{Marchisio2013.pdf}
\end{sourceblock}
通过分离增益项和损失项，我们得到

\begin{sourceblock}
\includegraphics[page=326,pagebox=mediabox,trim={127.78bp 114.46bp 65.07bp 523.35bp},clip,max width=0.94\textwidth,max height=0.38\textheight]{Marchisio2013.pdf}
\end{sourceblock}
其中碰撞后速度的大小（参见方程（7.49））由下式定义

\begin{sourceblock}
\includegraphics[page=326,pagebox=mediabox,trim={193.03bp 51.15bp 65.07bp 605.45bp},clip,max width=0.94\textwidth,max height=0.38\textheight]{Marchisio2013.pdf}
\end{sourceblock}
有关此推导的详细信息可以在 Aristov (2001) 中针对硬球情况的情况中找到，而 Asinari (2010) 中报告了其对通用内核的扩展。与单变量 PBE 的类方法和截面方法一样，粒子速度幅值空间 Σ 在大小为 ΔΣ:19 的 M 个区间上均匀离散化 Σ1 < Σ2 < Σ3 < · · · < ΣM−1 < ΣM < ΣM+1。对于速度大小等于 Ni = n(xii) 的粒子的数密度，可得出以下方程：

\begin{sourceblock}
\includegraphics[page=327,pagebox=mediabox,trim={133.57bp 490.93bp 67.56bp 137.08bp},clip,max width=0.94\textwidth,max height=0.38\textheight]{Marchisio2013.pdf}
\end{sourceblock}
其中关于场粒子的积分已被索引 j 的求和所取代。在等式中。 (7.62)式中， xii 对应于测试粒子 xi 的速度， xij 对应于场粒子 xi* 的速度，并且 n(xii) 和 n(xi j) 在最后一项中很容易被 Ni 和 Nj 代替。显而易见的是，由于 Ni 和 Nj 可以从方程 (2) 的最终积分中取出。 （7.62），方程中定义的损失项。 (7.50) 是

\begin{sourceblock}
\includegraphics[page=327,pagebox=mediabox,trim={187.82bp 382.73bp 67.56bp 292.94bp},clip,max width=0.94\textwidth,max height=0.38\textheight]{Marchisio2013.pdf}
\end{sourceblock}
其中

\begin{sourceblock}
\includegraphics[page=327,pagebox=mediabox,trim={115.35bp 333.87bp 67.56bp 326.79bp},clip,max width=0.94\textwidth,max height=0.38\textheight]{Marchisio2013.pdf}
\end{sourceblock}
相反，增益项的处理需要更多步骤。该术语量化了一对具有离散速度 xii 和 xij 的粒子的形成速率，这是由于另一对具有速度 xi' 和 xi' * 的粒子碰撞而产生的结果。让我们使用以下不等式分配索引 k 和 l 来指示后面的粒子：

\begin{sourceblock}
\includegraphics[page=327,pagebox=mediabox,trim={161.73bp 243.10bp 67.56bp 421.16bp},clip,max width=0.94\textwidth,max height=0.38\textheight]{Marchisio2013.pdf}
\end{sourceblock}
对于给定的一对离散粒子，存在许多遵循方程的粒子对。 (7.65) 对于不同的碰撞参数值−1 ≤x ≤1 和−1 ≤y ≤1。这些对沿着碰撞后速度平面中半径等于 (xi2 j)1/2 的圆排列，并且根据该圆与速度平面中线性离散化生成的网格之间的交集，可以识别四种不同的情况。当用 Nk 和 Nl 替换 n(xi′) 和 n(xi′ i + xi2 ∗) 时，增益项变为

\begin{sourceblock}
\includegraphics[page=327,pagebox=mediabox,trim={176.50bp 124.21bp 67.56bp 551.41bp},clip,max width=0.94\textwidth,max height=0.38\textheight]{Marchisio2013.pdf}
\end{sourceblock}
19 值得注意的是，在这个公式中，粒子速度区间的大小 Δψ 出现在离散方程中，这与群体平衡中发生的情况形成鲜明对比。这仅仅是因为在群体平衡的情况下，Ni 被定义为 n(xi) 在第 i 个区间内的积分，因此代表内坐标位于 xi 和 xi + dψ 之间的粒子的有限数量： Ni =

\begin{sourceblock}
\includegraphics[page=327,pagebox=mediabox,trim={310.48bp 78.08bp 165.85bp 609.00bp},clip,max width=0.94\textwidth,max height=0.38\textheight]{Marchisio2013.pdf}
\end{sourceblock}
{\small Ψi n(Ψ)dΨ。在 KE 的情况下，Ni 被定义为枢轴处分布的实际值：Ni = n(xii)。由于后者对应于将 n(xi) 作为积分中的常数，因此两种方法本质上相同，最多为 Δxi 因子。\par}
其中索引 k 和 l 的总和由碰撞后速度平面中的圆定义，并且

\begin{sourceblock}
\includegraphics[page=328,pagebox=mediabox,trim={219.34bp 590.27bp 65.07bp 75.41bp},clip,max width=0.94\textwidth,max height=0.38\textheight]{Marchisio2013.pdf}
\end{sourceblock}
αikl =
一旦索引（i，j）和（k，l）固定，积分域Ωkl，以及方程的结果。 （7.67），可以通过识别六种不同的可能情况来分析计算（Aristov，2001）。最终离散化后的KE为

\begin{sourceblock}
\includegraphics[page=328,pagebox=mediabox,trim={155.03bp 507.97bp 65.07bp 152.24bp},clip,max width=0.94\textwidth,max height=0.38\textheight]{Marchisio2013.pdf}
\end{sourceblock}
当 NDF 各向同性时有效。

\paragraph*{各向异性情况}
当各向同性假设不成立时，需要对速度矢量的三个分量进行离散化。相对于等式中的另一种更方便的表示法。 (7.50) 使用索引向量 i = (i1, i2, i3) 并将离散速度定义为

\begin{sourceblock}
\includegraphics[page=328,pagebox=mediabox,trim={199.40bp 388.52bp 65.07bp 298.18bp},clip,max width=0.94\textwidth,max height=0.38\textheight]{Marchisio2013.pdf}
\end{sourceblock}
其中索引向量在正整数值和负整数值上变化，因此每个方向上的粒子速度在最小（负）值和最大（正）值之间变化。选择间隔 Δ xi 的宽度，使得动力学的省略部分小于数值离散的近似值。请注意，方程中的离散化。 （7.69）对应于相空间中的线性网格，因此可能需要大量网格点来获得精确的解。对于各向同性情况，我们将速度矢量等于 xii 的粒子的数密度定义为 Ni = n(xii)。将离散化应用于动力学方程并用矢量索引 j = ( j1, j2, j3) 上的求和代替相对于场粒子的积分，属于第 i 类的粒子的变化率变为

\begin{sourceblock}
\includegraphics[page=328,pagebox=mediabox,trim={91.94bp 223.09bp 65.07bp 440.45bp},clip,max width=0.94\textwidth,max height=0.38\textheight]{Marchisio2013.pdf}
\end{sourceblock}
其中，速度差矢量 g 现在表示 xii 和 xij 之间的差，矢量索引 j 的总和对应于三个索引 j1、j2 和 j3 的三重总和。通过认识到 Ni = n(xii) 和 Nj = n(xij) 可以从方程 (2) 的最终积分中提取出来。 (7.70)，损失项变为

\begin{sourceblock}
\includegraphics[page=328,pagebox=mediabox,trim={181.71bp 98.72bp 65.07bp 537.86bp},clip,max width=0.94\textwidth,max height=0.38\textheight]{Marchisio2013.pdf}
\end{sourceblock}
p|xii − xij|。增益项（也称为逆碰撞积分）可以进一步简化，回想一下碰撞后和碰撞前速度之间的关系可以写为 ψ′ = 1 2(ψ + ψ*) + 1 2(g · x)x 和 ψ′ ∗= 1 2(ψ + ψ*) −1 2(g · x)x。如果将这些关系离散化并用相应的索引表示，我们得到 k = 1
{\small 2(i + j) + r\par}
l = 1 2(i + j)−r，其中r是整数或半整数索引。增益项的近似表达式为

\begin{sourceblock}
\includegraphics[page=329,pagebox=mediabox,trim={169.81bp 580.40bp 67.56bp 83.14bp},clip,max width=0.94\textwidth,max height=0.38\textheight]{Marchisio2013.pdf}
\end{sourceblock}
碰撞积分的剩余部分可以通过识别不同的情况来分析计算，从而将离散化的KE简化为类似于式（1）的形式。 （7.52）。详细描述可以在 Aristov (2001) 中找到。

\section*{7.3 矩方法}\phantomsection\addcontentsline{toc}{section}{7.3 矩方法}
MOM 背后的基本思想是求解 NDF 矩的演化方程，而不是通过求解实际的 GPBE 来确定 NDF 的演化。齐次系统的第 k 矩的一般演化方程（k = (k1,...,kM) 是矩指数向量，包含与 xi 的每个分量有关的矩的阶数），可以通过将矩变换应用于方程（1）来轻松获得。 （7.1）：

\begin{sourceblock}
\includegraphics[page=329,pagebox=mediabox,trim={148.31bp 355.83bp 67.56bp 279.40bp},clip,max width=0.94\textwidth,max height=0.38\textheight]{Marchisio2013.pdf}
\end{sourceblock}
替代但等效的矩符号为 mk1,...,kM = mk = m(k1, ..., kM) = m(k)。 MOM 具有许多令人感兴趣的优点。第一个优点来自于问题维数的减少。事实上，通过求解矩方程而不是 NDF，内坐标就被积分出来了。我们将看到，这通常是以必须检索的信息严重丢失为代价的（例如，通过重建 NDF）。另一个重要的优点是，矩对应于具有有意义的物理解释的量，因此至少在原则上是可以直接测量的。这是至关重要的一点，因为在许多应用中，NDF 不是直接测量的，而是通过积分测量来推断的。这些推理方法解决了反演背后的病态问题（通常有多个解决方案），并在 NDF 的确定中引入了不可忽略的不确定性。由于测量的积分属性与 NDF 的矩直接且唯一相关，因此在许多应用中，根据矩本身来制定和解决问题被证明是一种更加一致和方便的方法。使用 MOM 时会出现两个主要问题。第一个与要跟踪的时刻数量有关。理论上，无限矩集的知识相当于访问底层的 NDF，因此将原始无限集截断并投影到有限集意味着信息丢失。关键是要评估这种损失是否可以接受。这一点与第二个问题严格相关，即不可能将源项仅写为矩的函数。已经开发了不同的方法来克服这个闭合问题，并且通常它们的准确性取决于可用矩的数量。因此，要跟踪的时刻数量和采用的闭合问题通常会一起解决。下面，将针对单变量和双变量 PBE 以及 KE 分别讨论这些问题。

\subsection*{7.3.1 单变量 PBE}\phantomsection\addcontentsline{toc}{subsection}{7.3.1 单变量 PBE}
让我们只考虑一个内部坐标。均匀系统的 PBE 为 Eq. (7.3)，由此得到的 k 阶矩的演化方程为
dmk

\begin{sourceblock}
\includegraphics[page=330,pagebox=mediabox,trim={231.15bp 561.70bp 65.07bp 119.30bp},clip,max width=0.94\textwidth,max height=0.38\textheight]{Marchisio2013.pdf}
\end{sourceblock}
\begin{sourceblock}
\includegraphics[page=330,pagebox=mediabox,trim={67.52bp 498.58bp 65.07bp 147.77bp},clip,max width=0.94\textwidth,max height=0.38\textheight]{Marchisio2013.pdf}
\end{sourceblock}
0 Σkn dΣ，对于粒子速度，积分为 mk ≡ +∞ −∞Σkn dΣ。在单变量分布的情况下，零阶矩 m0 是总粒子数密度，表示每单位体积的粒子数。例如，如果内坐标是颗粒尺寸，则 m1 是颗粒总长度密度，或者换句话说，是每单位体积颗粒的总长度。因此很明显，这两个量的比率代表了数均平均粒径：d10 = m1/m0。二阶矩 m2 通过表面形状因子 kA 与每单位体积的总颗粒表面积成正比。三阶矩 m3 通过体积形状因子 kV 与单位体积的总颗粒体积（或者换句话说，总分散相体积分数）相关。另一个非常常见的平均粒径是表面积平均粒径，通常称为索特平均直径，d32 = m3/m2，而 d43 = m4/m3 量化了按颗粒体积而不是颗粒表面积加权的平均粒径。如果内坐标是颗粒体积，则 m0 又是总颗粒数密度，而 m1 是总分散相体积分数。

\paragraph*{具有相空间漂移的矩方程}
在漂移的情况下，方程的源项。 (7.75) 可以利用分部积分规则计算：

\begin{sourceblock}
\includegraphics[page=330,pagebox=mediabox,trim={166.33bp 213.65bp 65.07bp 452.56bp},clip,max width=0.94\textwidth,max height=0.38\textheight]{Marchisio2013.pdf}
\end{sourceblock}
等式右侧的第一项。 (7.77) 对于 Ψ → ∞ 始终为零；事实上，在这个极限下，n(t, xi) 趋向于零的速度比 xik ̇xi(xi) 趋于无穷大的速度要快。对于 xi = 0，当 k > 0 时，第一项为空，而对于 k = 0，我们得到

\begin{sourceblock}
\includegraphics[page=330,pagebox=mediabox,trim={215.72bp 137.12bp 65.07bp 542.22bp},clip,max width=0.94\textwidth,max height=0.38\textheight]{Marchisio2013.pdf}
\end{sourceblock}
它表示 xi = 0 处颗粒通过边界的净通量。对于 ̇ xi(xi) > 0 的颗粒生长，该项可以很容易计算，并且通常与成核率相关：20

\begin{sourceblock}
\includegraphics[page=330,pagebox=mediabox,trim={223.64bp 72.22bp 65.07bp 608.78bp},clip,max width=0.94\textwidth,max height=0.38\textheight]{Marchisio2013.pdf}
\end{sourceblock}
\footnotetext[20]{通常，假设成核仅发生在 xi = 0 时，因此 J(t, 0) 是总体成核率。}
更复杂的是当 ̇xi(0) < 0 时的情况，例如在颗粒溶解或液滴蒸发期间。在这种情况下，方程中的第一项。 (7.78) 与成核速率无关，而是与 NDF 的逐点值相关，因此 n(t, 0) 必须从分布的已知矩推断出来。21 在速率独立于内坐标的正漂移（即生长）的情况下，很容易证明

\begin{sourceblock}
\includegraphics[page=331,pagebox=mediabox,trim={219.18bp 548.41bp 67.56bp 130.92bp},clip,max width=0.94\textwidth,max height=0.38\textheight]{Marchisio2013.pdf}
\end{sourceblock}
而如果增长率相对于内坐标遵循幂律关系 ̇ ̇ ̇ ̇ ̇ 0 ̇α，我们发现

\begin{sourceblock}
\includegraphics[page=331,pagebox=mediabox,trim={213.04bp 498.11bp 67.56bp 181.22bp},clip,max width=0.94\textwidth,max height=0.38\textheight]{Marchisio2013.pdf}
\end{sourceblock}
得出特定 α 值的一组未封闭方程。对于更复杂的增长率，没有可用的解析解，必须寻求数值闭包。

\paragraph*{具有相空间扩散的矩方程}
相空间中的扩散主要与速度相关，因此对于此项，我们将在区间 (−∞, +∞) 上进行积分。在纯扩散的情况下，方程的源项。 (7.75) 可以通过两次分部积分来计算：

\begin{sourceblock}
\includegraphics[page=331,pagebox=mediabox,trim={105.25bp 326.55bp 67.56bp 294.98bp},clip,max width=0.94\textwidth,max height=0.38\textheight]{Marchisio2013.pdf}
\end{sourceblock}
等式右侧的前两项。 (7.82) 对于|Ψ| 始终为空→ 无穷大；事实上，在这个极限下，NDF 趋向于零的速度比 xikD 趋于无穷大的速度要快。由此我们得到

\begin{sourceblock}
\includegraphics[page=331,pagebox=mediabox,trim={184.91bp 264.94bp 67.56bp 401.51bp},clip,max width=0.94\textwidth,max height=0.38\textheight]{Marchisio2013.pdf}
\end{sourceblock}
它充当 k ≥ 2 的源项。对于 D = D0 + D1Σ2 的情况（回想一下，对于所有 Σ，D ≥0），我们找到扩散项的闭合表达式：

dmk

\begin{sourceblock}
\includegraphics[page=331,pagebox=mediabox,trim={182.43bp 205.17bp 67.56bp 475.82bp},clip,max width=0.94\textwidth,max height=0.38\textheight]{Marchisio2013.pdf}
\end{sourceblock}
否则，必须为方程式提供矩闭合。 （7.83）。

\paragraph*{具有点过程的矩方程}
在聚合的情况下，当内部坐标 xi 在一个聚合事件期间纯粹是相加时，PBE 具有等式 1 中给出的源项。 (7.9) 则 k 阶矩的演化方程为

\begin{sourceblock}
\includegraphics[page=331,pagebox=mediabox,trim={70.34bp 80.85bp 67.56bp 576.76bp},clip,max width=0.94\textwidth,max height=0.38\textheight]{Marchisio2013.pdf}
\end{sourceblock}
\footnotetext[21]{这对于蒸发喷雾来说尤其重要。对于相关的理论和数值问题，以及对该主题的深入讨论，我们建议读者参考专业文献。}
在聚合核 β(xi, xi′) = β0 恒定的情况下，可以得到

\begin{sourceblock}
\includegraphics[page=332,pagebox=mediabox,trim={149.31bp 591.61bp 65.07bp 75.83bp},clip,max width=0.94\textwidth,max height=0.38\textheight]{Marchisio2013.pdf}
\end{sourceblock}
因此，对于 k = 0，

\begin{sourceblock}
\includegraphics[page=332,pagebox=mediabox,trim={218.73bp 545.92bp 65.07bp 128.70bp},clip,max width=0.94\textwidth,max height=0.38\textheight]{Marchisio2013.pdf}
\end{sourceblock}
对于 k = 1，

\begin{sourceblock}
\includegraphics[page=332,pagebox=mediabox,trim={233.21bp 505.53bp 65.07bp 175.47bp},clip,max width=0.94\textwidth,max height=0.38\textheight]{Marchisio2013.pdf}
\end{sourceblock}
它说明了纯加性内坐标聚合期间一阶矩的守恒。相反，高阶矩由

\begin{sourceblock}
\includegraphics[page=332,pagebox=mediabox,trim={199.68bp 434.88bp 65.07bp 228.20bp},clip,max width=0.94\textwidth,max height=0.38\textheight]{Marchisio2013.pdf}
\end{sourceblock}
在纯粹破损的情况下，PBE 为

\begin{sourceblock}
\includegraphics[page=332,pagebox=mediabox,trim={165.16bp 377.95bp 65.07bp 288.51bp},clip,max width=0.94\textwidth,max height=0.38\textheight]{Marchisio2013.pdf}
\end{sourceblock}
k 阶矩的演化方程为

\begin{sourceblock}
\includegraphics[page=332,pagebox=mediabox,trim={127.65bp 320.13bp 65.07bp 346.07bp},clip,max width=0.94\textwidth,max height=0.38\textheight]{Marchisio2013.pdf}
\end{sourceblock}
对于恒定破损核的某些特定情况，可以使用解析解，a(xi) = a0。例如，在具有子分布的对称碎片的情况下

\begin{sourceblock}
\includegraphics[page=332,pagebox=mediabox,trim={202.36bp 265.20bp 65.07bp 415.79bp},clip,max width=0.94\textwidth,max height=0.38\textheight]{Marchisio2013.pdf}
\end{sourceblock}
mk 的演化方程为

\begin{sourceblock}
\includegraphics[page=332,pagebox=mediabox,trim={202.13bp 207.30bp 65.07bp 471.98bp},clip,max width=0.94\textwidth,max height=0.38\textheight]{Marchisio2013.pdf}
\end{sourceblock}
这再次遵守一阶矩的守恒约束。在断裂成两个以 x 和 1 −x 为特征的碎片的情况下，

\begin{sourceblock}
\includegraphics[page=332,pagebox=mediabox,trim={174.71bp 140.42bp 65.07bp 532.61bp},clip,max width=0.94\textwidth,max height=0.38\textheight]{Marchisio2013.pdf}
\end{sourceblock}
mk 的最终演化方程为

\begin{sourceblock}
\includegraphics[page=332,pagebox=mediabox,trim={179.90bp 84.01bp 65.07bp 587.75bp},clip,max width=0.94\textwidth,max height=0.38\textheight]{Marchisio2013.pdf}
\end{sourceblock}
{\small 在尺寸相关核的情况下，没有可用的解析解并且需要数值闭包。\par}
\paragraph*{PBE 的矩方程}
经历成核、生长、扩散、聚集和破碎的均匀系统的一般演化方程为

\begin{sourceblock}
\includegraphics[page=333,pagebox=mediabox,trim={89.09bp 508.54bp 67.56bp 104.42bp},clip,max width=0.94\textwidth,max height=0.38\textheight]{Marchisio2013.pdf}
\end{sourceblock}
矩闭包是一种将源项写在等式右侧的算法。 （7.96），明确或隐含地，就矩而言：

\begin{sourceblock}
\includegraphics[page=333,pagebox=mediabox,trim={193.66bp 373.82bp 67.56bp 235.71bp},clip,max width=0.94\textwidth,max height=0.38\textheight]{Marchisio2013.pdf}
\end{sourceblock}
我们已经展示了这组微分方程在一些简单情况下如何闭合（即，对于零尺寸的成核，具有恒定的增长率以及恒定的聚集和破碎内核）。一般来说，当系统不封闭时，必须通过稳定的数值方法来寻求封闭。

\paragraph*{矩方程的闭包}
多年来，已经开发了几种方法，我们从插值闭包矩方法（MOMIC）开始对这些方法进行简要回顾（Frenklach，2002；Frenklach＆Harris，1987）。该方法基于非常简单且有效的想法，即可以通过使用例如拉格朗日（或样条）插值从已知的矩集内插（或外推）未知矩。通过绘制分布矩与矩阶的关系可以清楚地说明这一概念。从图 7.1 中可以看出，了解整数矩足以（通过插值）估计未知的非整数矩。该方法仅对于可以根据某些（已知或未知）矩来简化矩的源项的情况有效。对于简单问题来说几乎总是如此，但在许多（并非非常罕见）的情况下这是不可能的。我们现在用两个基于聚合问题的示例来说明这些概念。聚合的主要困难是聚合核对内部坐标的依赖，特别是其非加性特征。例如，在连续体体系22中，当 xi 表示粒子质量时，聚合核为

\begin{sourceblock}
\includegraphics[page=333,pagebox=mediabox,trim={162.32bp 73.06bp 67.56bp 597.97bp},clip,max width=0.94\textwidth,max height=0.38\textheight]{Marchisio2013.pdf}
\end{sourceblock}
\footnotetext[22]{不同的聚集/凝固机制和相应的内核将在第 5 章中讨论。}
\begin{sourceblock}
\includegraphics[page=334,pagebox=mediabox,trim={58.00bp 439.80bp 58.23bp 59.00bp},clip,width=\textwidth,max height=0.72\textheight]{Marchisio2013.pdf}
\par\vspace{0.45\baselineskip}
\small 图7.1。正态分布的矩与矩阶。如果仅知道整数矩（符号），则简单插值（连续线）允许人们估计非整数矩。
\end{sourceblock}
聚合的源项可以写为

\begin{sourceblock}
\includegraphics[page=334,pagebox=mediabox,trim={147.41bp 338.37bp 186.88bp 327.83bp},clip,max width=0.94\textwidth,max height=0.38\textheight]{Marchisio2013.pdf}
\end{sourceblock}
\begin{sourceblock}
\includegraphics[page=334,pagebox=mediabox,trim={147.41bp 280.54bp 65.07bp 367.79bp},clip,max width=0.94\textwidth,max height=0.38\textheight]{Marchisio2013.pdf}
\end{sourceblock}
其中 k = 2, 3, .。。。如果方程。 （7.98）代入式（7.98） (7.99)，我们得到

\begin{sourceblock}
\includegraphics[page=334,pagebox=mediabox,trim={193.94bp 226.21bp 65.08bp 445.53bp},clip,max width=0.94\textwidth,max height=0.38\textheight]{Marchisio2013.pdf}
\end{sourceblock}
以及

\begin{sourceblock}
\includegraphics[page=334,pagebox=mediabox,trim={113.01bp 177.04bp 65.08bp 486.03bp},clip,max width=0.94\textwidth,max height=0.38\textheight]{Marchisio2013.pdf}
\end{sourceblock}
其中 k = 2, 3, .。。产生一组未封闭的方程。为了获得闭合，源项中出现的分数阶矩必须闭合。在 MOMIC 中，这是通过整数矩之间的插值来完成的。当使用例如拉格朗日多项式对约简矩的对数 log(μk) = log(mk) −log(m0) 以及对可用于正分数矩的整个矩集进行插值时，插值会更有效，log(μp) = Lp(log(μ0), log(μ1), log(μ2), ..., log(μN)) for p > 0，其中 N 是最大数跟踪的矩，并仅从前三个矩推断负分数矩，log(μp) = Lp(log(μ0), log(μ1), log(μ2)) for p < 0。当无法将源项写为某些已知或未知矩的显式函数时，MOMIC 会变得非常麻烦。例如，对于自由分子体系中的聚集，聚集核为

\begin{sourceblock}
\includegraphics[page=335,pagebox=mediabox,trim={165.44bp 586.20bp 67.57bp 81.11bp},clip,max width=0.94\textwidth,max height=0.38\textheight]{Marchisio2013.pdf}
\end{sourceblock}
显而易见，指数 1/2 的存在不允许我们将源项写为显式求和。在这种情况下，一个有趣的替代方案是估计以下内核的源项中出现的结果总和：

\begin{sourceblock}
\includegraphics[page=335,pagebox=mediabox,trim={171.57bp 496.53bp 67.57bp 175.59bp},clip,max width=0.94\textwidth,max height=0.38\textheight]{Marchisio2013.pdf}
\end{sourceblock}
对于 l = 0, 1, 2, .。。，然后再次插值以获得 l = 1/2 的解。正如预期的那样，在这些情况下，该方法失去了效率和简单性，而它不适用于内核（或连续变化率）包含对内部坐标的指数依赖性，或无法通过多项式在内坐标方面精确近似的任何其他类型的函数依赖性（例如指数断裂内核、指数增长率等）的情况（不幸的是，这种情况并不罕见！）。我们建议对有关插值有效性的数学细节感兴趣的读者参阅 Frenklach (2002)。参见 Alexiadis 等人。 (2004)、Barrett \& Jheeta (1996) 和 Diemer \& Olson (2002b)，用于相关插值方法的描述。克服封闭问题的第二类方法是对 NDF 做出函数假设。最简单的是假设 NDF 由以内坐标平均值（例如 m1/m0）为中心的 delta 函数组成，或者换句话说，假设粒子总体是单分散的。利用这种方法，可以很容易地计算出缺失的矩（例如 mk = m1−k 0 mk 1），如 Kruis 等人所示。 （1993）。一个非常流行的替代方案是对 NDF 使用不太近似的函数形式，例如，在对数正态分布的情况下。三个低阶矩的知识足以计算对数正态分布的三个参数。 （详细信息请参见第 3 章。）闭合 NDF 可用于计算所有缺失矩以及源项和汇项（Bensberg 等，1999；Friedlander \& Wang，1966；Lee，1983；Swift \& Friedlander，1964）。对于明显呈现双峰（或多峰）分布的问题，粒子群可以分为不同的组（例如种子和生长的晶体或核和聚集体等），由以不同内坐标平均值（例如数均和体积平均值）为中心的δ函数描述，如Wu等人所解释的。 (1988)，或通过对数正态分布(Megaridis \& Dobbins, 1990)。一种非常有效的技术是有限尺寸域完整试函数矩集法 (FCMOM) (Strumendo \& Arastoopour, 2008; Strumendo \& Arastoopour, 2009, 2010)，其中 NDF 是通过使用完整的正交函数系统的级数展开从矩重构的。对不同闭包的细节以及分布重建问题感兴趣的读者应该查阅 Diemer \& Olson (2002a, 2002b)。然而，应该强调的是，大多数提出的单变量 PBE 方法都存在一些重要的缺点，主要与其适用范围和数值精度有关。 McGraw (1997) 提出了一种基于正交近似的方法，似乎可以克服这些限制。这种方法称为矩求积法 (QMOM)，在 7.4 节中进行了描述。

\subsection*{7.3.2 双变量和多变量 PBE}\phantomsection\addcontentsline{toc}{subsection}{7.3.2 双变量和多变量 PBE}
正如 Hulburt \& Katz (1964) 所解释的，MOM 也适用于双变量 PBE。在这种情况下，NDF 的矩定义为 mk ≡

\begin{sourceblock}
\includegraphics[page=336,pagebox=mediabox,trim={296.56bp 593.47bp 176.22bp 85.88bp},clip,max width=0.94\textwidth,max height=0.38\textheight]{Marchisio2013.pdf}
\end{sourceblock}
1 Ψk2 2 n dΞ，其中 Ψ1 和 Ψ2 是两个内坐标，k = (k1, k2) 是包含矩相对于两个内坐标的阶数的指数向量。矩可以具有特定的物理意义，对于单变量情况，m0,0 是总粒子数密度。同样，不同矩的比值可以用来表征总体的积分性质，例如m1,0/m0,0是第一个内坐标ψ1的数均值，而m0,1/m0,0是第二个内坐标ψ2的数均值。有关双变量矩和可测量量之间关系的详细信息和示例，读者可以参考 Rosner 等人的工作。 （2003）。对于单变量情况，必须求解矩的演化方程，

\begin{sourceblock}
\includegraphics[page=336,pagebox=mediabox,trim={114.84bp 443.73bp 65.08bp 223.20bp},clip,max width=0.94\textwidth,max height=0.38\textheight]{Marchisio2013.pdf}
\end{sourceblock}
属于 MOM 模拟期间跟踪的矩集。两个内坐标的存在使矩数量的选择和要包含在矩集中的矩的性质变得复杂。事实上，如果对于单变量情况，最自然的选择是选择从 m0 开始到一定顺序的完整矩集，23 对于双变量情况，不存在直接且琐碎的选择。一种可能的策略是通过按全局顺序对矩进行分类来定义矩集，γ = k1 + k2。根据这种分类，全局零阶矩（γ=0）只有一个，即m0,0，但一阶矩（γ=1）有两个，即m1,0和m0,1，二阶矩（γ=2）有3个，即m2,0、m1,1和m0,2，高阶矩以此类推。为了使选择包含在集合中的矩的过程合理化，通常将它们组织在三角形表中，如第 3.2.1 节的表 3.6 所示。只要直接使用方程式的可能性即可。对于（7.104）式而言，对于单变量情况，主要限制在于难以找到解析解和开发简单且稳定的闭包。解析解非常罕见，仅适用于连续变化率和独立于内坐标的内核。在这些简化假设下，对于成核、生长、扩散、24 聚集和断裂（具有加性内坐标）的情况，双变量 PBE 变为

\begin{sourceblock}
\includegraphics[page=336,pagebox=mediabox,trim={74.34bp 149.07bp 65.08bp 501.77bp},clip,max width=0.94\textwidth,max height=0.38\textheight]{Marchisio2013.pdf}
\end{sourceblock}
23 对于简单的应用，可以选择（m0、m1、m2、m3），但对于复杂的问题，可能需要更多矩，因此最高阶可以固定为 5 或 6。这对于 NDF 重建也有一些重要的影响；事实上，复杂 NDF 的精确重建需要许多矩的知识。

\footnotetext[24]{我们提醒读者，扩散通常与无界域上的速度相空间相关。如果在有界域上使用扩散，则必须为扩散通量提供适当的边界条件。在等式中。 (7.105) 我们假设 NDF 不会因相空间边界的扩散而改变。}
因此 k 阶混合矩的输运方程为

\begin{sourceblock}
\includegraphics[page=337,pagebox=mediabox,trim={95.14bp 535.16bp 67.57bp 81.35bp},clip,max width=0.94\textwidth,max height=0.38\textheight]{Marchisio2013.pdf}
\end{sourceblock}
对于恒定内核和对称破损的单变量情况（参见方程（7.92）），在破损的加性内坐标假设下，

\begin{sourceblock}
\includegraphics[page=337,pagebox=mediabox,trim={173.09bp 485.72bp 67.57bp 193.21bp},clip,max width=0.94\textwidth,max height=0.38\textheight]{Marchisio2013.pdf}
\end{sourceblock}
矩演化方程组（见方程（7.106））是封闭的并且存在一些解析解。对于低阶矩我们可以写

\begin{sourceblock}
\includegraphics[page=337,pagebox=mediabox,trim={182.63bp 382.18bp 67.57bp 242.64bp},clip,max width=0.94\textwidth,max height=0.38\textheight]{Marchisio2013.pdf}
\end{sourceblock}
对于高阶矩可以写出类似的方程。在所有其他情况下，必须使用已知时刻来恢复闭合。通过类似的形式主义，这些概念可以扩展到多元系统；然而，对于双变量系统，不存在解决矩闭合问题的通用程序。此外，随着内坐标数量的增加，保证给定的矩闭合是可实现的变得越来越困难（更多细节请参见第3章）。

\subsection*{7.3.3 碰撞KE}\phantomsection\addcontentsline{toc}{subsection}{7.3.3 碰撞KE}
在速度分布的情况下，NDF 的矩代表重要的量。例如零阶矩，m(0) = m(0, 0, 0) =

\begin{sourceblock}
\includegraphics[page=337,pagebox=mediabox,trim={351.99bp 260.16bp 135.35bp 428.04bp},clip,max width=0.94\textwidth,max height=0.38\textheight]{Marchisio2013.pdf}
\end{sourceblock}
n dψ，表示粒子数密度N。如果将该量乘以单个粒子的质量Mp，则得到分散相质量密度ρp。相对于粒子速度的一阶矩与分散相动量密度有关： ρpUp = Mp

\begin{sourceblock}
\includegraphics[page=337,pagebox=mediabox,trim={115.24bp 198.78bp 336.08bp 475.85bp},clip,max width=0.94\textwidth,max height=0.38\textheight]{Marchisio2013.pdf}
\end{sourceblock}
Ψn Ψ。 NDF 的二阶矩与总颗粒应力张量相关：P = Mp ( xi −Up )⊗( xi −Up )n d ψ。众所周知，该张量的迹与总分散相动能有关，总分散相动能可写为 ρpep = 1 2 Mp ϋ2n dϋ，其中 ep 为分散相比能。由于 Ψ2 = Ψ2 1 + Ψ2 2 + Ψ2 3，所以上式可以写为 ρpep = 1 2 Mp[m(2, 0, 0) + m(0, 2, 0) + m(0, 0, 2)]。当使用 MOM 解决与碰撞 KE 相关的问题时，必须定义要跟踪的矩集。在低阶矩的情况下，选择非常简单，因为所有低阶矩都具有清晰且重要的物理意义，并且应该包含在任何描述中。当该方法扩展到没有直接物理意义的高阶矩时，就会出现应该考虑哪些矩来准确描述颗粒系统的问题。25对此的一个答案

\footnotetext[25]{与矩数相关的另一个重要问题将在稍后讨论，并且已在第 3 章中讨论过。}
问题是按照特定顺序包含所有时刻。例如，如果仅考虑一阶矩，则将使用四个矩（即 m0,0,0、m1,0,0、m0,1,0 和 m0,0,1）。如果考虑二阶矩，则应将六个矩添加到集合中，从而总共十个。增加矩集中包含的矩的顺序，20, 35, 56, 84, 120, 165, 220, 286, 364, .。。将使用时刻。另一种答案可以基于对一组扩展矩的选择（Grad，1949a，1958）。均相系统的动力学方程由方程给出。 （7.47）。 NDF 零阶矩的演化方程为空，这是因为碰撞积分不会改变粒子数量，或者更明确地说，

\begin{sourceblock}
\includegraphics[page=338,pagebox=mediabox,trim={379.87bp 528.41bp 107.47bp 159.79bp},clip,max width=0.94\textwidth,max height=0.38\textheight]{Marchisio2013.pdf}
\end{sourceblock}
C dΣ = 0。如果粒子速度的变化率（即粒子加速度）是粒子速度的线性函数（即 ̇ xi = a + b Σ），则一阶矩的演化方程为

\begin{sourceblock}
\includegraphics[page=338,pagebox=mediabox,trim={193.18bp 407.32bp 65.08bp 217.52bp},clip,max width=0.94\textwidth,max height=0.38\textheight]{Marchisio2013.pdf}
\end{sourceblock}
其中我们使用了碰撞积分的一阶矩相同为零的事实：

\begin{sourceblock}
\includegraphics[page=338,pagebox=mediabox,trim={109.28bp 385.85bp 378.07bp 302.35bp},clip,max width=0.94\textwidth,max height=0.38\textheight]{Marchisio2013.pdf}
\end{sourceblock}
ψC dψ = 0。由于现在碰撞积分的矩不为零，因此高阶矩的演化方程变得复杂。例如，具有各向同性扩散的 KE 通用矩的演化方程为
dmi, j,k dt = ia1mi−1, j,k + ja2mi, j−1,k + ka3mi, j,k−1 + (i + j + k)bmi, j,k

\begin{sourceblock}
\includegraphics[page=338,pagebox=mediabox,trim={123.59bp 275.91bp 65.08bp 385.00bp},clip,max width=0.94\textwidth,max height=0.38\textheight]{Marchisio2013.pdf}
\end{sourceblock}
1Ψ j 2Ψk 3C dΨ。与将 MOM 应用于 PBE 一样，方程： （7.110）通常是非封闭的，因为碰撞积分的矩不能用矩本身来写。27 需要提醒读者的是，方程（7.110）通常是非封闭的。 (7.110) 是方程的空间齐次版本。 （6.131）（加上由于粒子速度的连续变化率而导致的漂移项）。对于 KE，在某些限制情况下可以使用简单的矩闭合。例如，在麦克斯韦粒子的情况下，碰撞积分的阶矩 γ = i + j + k 可以明确地用高达 γ 阶的矩来写，如方程 1 所示。第 6 章的方程（6.109）。此外，当使用 BGK 近似对碰撞积分进行线性化时，克服了与碰撞积分相关的封闭问题（参见方程（6.120）和方程（6.121））。在所有其他情况下，需要某种形式的关闭。要对将 MOM 应用于碰撞 KE 所生成的方程组进行全面而彻底的讨论，读者可以参考 Truesdell \& Muncaster (1980) 的经典著作。下面将简要讨论 Grad 方法提供的闭包。其中 Πi, j,k =
26 可以通过设置 b = −1/τ 和 D = [m(2, 0, 0) + m(0, 2, 0) + m(0, 0, 2)]/[τm(0, 0, 0)] 来构建保存总能量的各向同性扩散模型，其中 τ 是弛豫时间尺度。

\footnotetext[27]{此外，在处理非均匀系统时，对流项的存在会产生一些额外的闭合问题。这些问题将在第 8 章中讨论。}
\paragraph*{Grad 的 13 时刻闭幕式}
Grad 的 13 矩理论基于分散相质量密度 ρp、分散相动量密度 ρpUp 的三个分量、总分散相应力张量 P 的六个分量以及分散相能量通量 q 的三个分量的演化方程的解，这些方程由 q = 1 定义

\begin{sourceblock}
\includegraphics[page=339,pagebox=mediabox,trim={65.03bp 531.29bp 67.55bp 121.09bp},clip,max width=0.94\textwidth,max height=0.38\textheight]{Marchisio2013.pdf}
\end{sourceblock}
i | xi −Up|2(xii −Upi)] f ∗，其中重复索引意味着求和，并且选择系数以便 n[13] 再现 13 矩集。请注意，该速度分布函数是根据关于平均速度的速度波动来定义的 (xi′ = xi −Up)。这种选择确保，像原始 KE 和矩方程一样，n[13] 和所得闭包将是伽利略不变的（即系数的值不依赖于平均速度）。我们建议读者参阅 Struchtrup (2005)，了解有关此方法的详细信息以及对其扩展到 26 个矩的描述。 Grad 方法的一个重要限制是远离平衡，尽管可以很好地描述矩的演化，但 NDF 的相应假设常常被证明是非物理的，因为它会导致负 NDF（如 Icardi 等人（2012）中关于齐次各向同性玻尔兹曼方程的说明）。这个问题与 n[13] 定义中的系数以线性形式出现有关。当从 n[13] 计算的 13 个矩设置为等于已知矩时，将产生线性系统，并且不能保证所得系数对于所有 xi 都将保持 n[13] ≥0。任何基于线性展开的矩闭合都会出现同样的问题（例如 Strumendo \& Arastoopour, 2008, 2009, 2010）。相比之下，7.4 节中描述的方法使用非线性重建，可以保证可实现的 NDF。 （有关更多详细信息，请参阅第 3 章。）Grad 方法的另一个关键点是应跟踪的时刻数量。理想情况下，在求解 KE 时（以及在 GPBE 的背景下），最好提前知道这个数字，即在运行多相过程的特定模拟之前。此外，最好将该数字定义为所涉及的无量纲数字的函数（参见第 1 章）。不幸的是，这样一个简单的标准很难定义，因此最好的选择是用不断增加的矩来模拟该过程，然后当矩数量的进一步增加不会在最终结果中产生任何显著差异时停止。

\paragraph*{格子玻尔兹曼方法}
格子玻尔兹曼方法 (LBM) 可以表示为多元 MOM (Lalle-mand \& Luo, 2003)，其本质与 Grad 方法类似。在LBM中，晶格速度cα是固定的，使得在均匀晶格上向量ci Δt（其中Δt是时间步长）对应于晶格向量。虽然这一点对于同质系统并不重要，但对于非齐次系统中用于空间传输的数值算法至关重要。每个晶格速度都有一个非负权重 fα，NDF 的假设形式为 n(t, xi) = - α fα(t)δ(xi −cα)，其中多维狄拉克 δ 函数定义为 δ(xi −cα) = (3 i=1 δ(xii −ci,α)。与 Grad 方法一样，LBM 中的 N 个权重是通过强制与由下式给出的 N 个矩的 KE 保持一致来确定的： m(k1, k2, k3) = -
{\small 1αck2 2αck3 3α fα。由于分量 c1α、c2α 和 c3α 是固定的，因此矩可以写成未知权重的线性系统：\par}
{\small α ck1\par}
\begin{sourceblock}
\includegraphics[page=340,pagebox=mediabox,trim={195.78bp 597.17bp 65.07bp 59.44bp},clip,max width=0.94\textwidth,max height=0.38\textheight]{Marchisio2013.pdf}
\end{sourceblock}
其中系数矩阵 A 取决于 cα 的分量。等式右侧使用的 N 个矩的集合。 （7.111）包括所有二阶矩和选定的高阶矩（Lallemand＆Luo，2003），并且被选择使得A是可逆的。平衡权重 f * α 可以由等式确定。 （7.111）使用高斯矩，但实际上采用平衡分布的线性化形式（Lallemand＆Luo，2003）。与 Grad 方法中的系数一样，LBM 中的权重是矩的线性函数（即 f = A−1m），因此，对于可实现矩集 m，其中一些权重将为负。还值得注意的是，权重将取决于平均速度的值，因此 LBM 本质上不是伽利略不变量（Lallemand \& Luo，2003）。事实上，只有当平均速度的大小相对于晶格速度的大小较小时（即低马赫数）才可以使用 LBM。实际上，LBM 用于低马赫数非均匀流，其中权重 fα 保持接近 f * α。事实上，如前所述，LBM 的真正优势在于通过选择晶格速度而产生的高效平流方案，而不是其捕获 KE 非平衡解的能力。

\section*{7.4 基于求积的矩方法}\phantomsection\addcontentsline{toc}{section}{7.4 基于求积的矩方法}
基于求积的矩方法 (QBMM) 构成了一类非常成功的矩方法，它通过使用正交近似克服了闭合问题。 M 个内坐标的 N 个节点的一般（多元）求积需要了解 N 个权重、wα 和 N 个节点（或节点）， Σα = Σ1,α, Σ2,α,。。。 , ΣM,α，其中 1 ≤α ≤N。请注意，每个节点 xiα 都是长度为 M 的向量。在我们的表示法中，节点的第一个索引定义内部坐标，而第二个索引定义节点的编号。 N 点、M 维求积通过以下等式克服了封闭问题（始终以积分形式出现）：

\begin{sourceblock}
\includegraphics[page=340,pagebox=mediabox,trim={75.74bp 244.22bp 357.43bp 436.76bp},clip,max width=0.94\textwidth,max height=0.38\textheight]{Marchisio2013.pdf}
\end{sourceblock}
α=1 wαg(xiα)。读者会将此表达式视为近似积分的经典求积公式。有两种主要方法可用于模拟特定矩集（以下称为矩集）的演化。在第一种方法（称为矩求积法 (QMOM)）中，求解矩的演化方程，并通过求积近似估计源项（出现在方程 (7.73) 的右侧，并在方程 (7.74) 中详细说明）。正交近似又是通过采用矩反演算法从属于矩集的可用矩来确定的。 （可用的不同矩反演算法将在第 3 章中讨论。）在第二种方法（称为矩的直接求积法 (DQMOM)）中，通过确保属于矩集的矩遵守正确的演化方程来求解求积近似本身的演化。 QMOM/DQMOM的基本思想是，GPBE中出现的NDF由N个多维狄拉克δ函数之和表示：n(t, ϋ) = -N α=1 wα(t)δ(ϋ −ϋα(t))。该方程定义了关于 NDF 函数形式的主要假设，这是 QBMM 的基础。读者可能已经注意到，QMOM 与 LBM 非常相似。然而，存在一些非常重要的差异。首先，QMOM 中的节点不是恒定的，而是允许变化的。其次，QMOM 中的矩演化方程非线性地依赖于节点，因此，矩反演问题比 LBM 的困难得多。正如我们将在下面看到的，这种额外的复杂性有一线希望（详细信息请参阅第 3 章）：所有可实现的矩集都将导致非负权重，因此，QMOM 将始终代表可实现的 NDF。下面将首先针对单变量 PBE 讨论 QBMM，然后针对双变量和多变量 PBE，最后针对碰撞 KE。如第 3 章所述，可以通过在 n(t, xi) 的定义中使用核密度函数（即有限分布）来扩展 QBMM 的精度。然而，为了简单起见，我们将在下面的讨论中假设狄拉克δ函数。

\subsection*{7.4.1 单变量 PBE}\phantomsection\addcontentsline{toc}{subsection}{7.4.1 单变量 PBE}
QBMM 中的闭合问题通过使用求积近似来克服，该近似具有计算积分非常准确的巨大优势。事实上，它们通常表现出比相同阶的其他插值方案高得多的精度。通过确定与 NDF 正交的多项式序列来构建正交近似。该理论的主要结果（详细信息请参见第 3 章）指出，如果求积近似的节点是与 NDF 正交的 N 阶多项式的根，则求积近似为

\begin{sourceblock}
\includegraphics[page=341,pagebox=mediabox,trim={74.74bp 388.47bp 356.03bp 292.51bp},clip,max width=0.94\textwidth,max height=0.38\textheight]{Marchisio2013.pdf}
\end{sourceblock}
如果 f(xi) 是 2N -1 阶或更小的多项式，则 α=1 f(xiα)wα 是精确的（Dette \& Studden，1997）。在所有其他情况下，f(xi) 越接近多项式，近似值就越准确。在解释如何根据 NDF 的矩以及 QMOM 和 DQMOM 之间的异同来计算求积近似的权重与节点之前，让我们讨论一下方程 1 中源项 S k 的估计。 （7.76）。该项的估计存在一定的误差，该误差取决于正交近似的阶数。直观上，人们可能认为随着节点数量的增加，误差会减小到非常小的值；然而，误差对节点数量的依赖性很难先验预测。此外，提醒读者增加节点数量意味着需要跟踪更多时刻是有用的。一般来说，对于单变量系统，可以通过跟踪 2N 个矩来构造 N 阶求积。这通常会影响计算时间（由于需要求解更多方程，因此计算时间当然会增加）和方法的稳定性（由于与矩反演算法相关的问题）。由于正交近似的精度很大程度上取决于应用，因此让我们针对具体示例详细讨论其中一些问题。在标准成核、（正）生长、聚集和断裂的情况下，将 QBMM 应用于方程式中的源项。 (7.96) 收益率

\begin{sourceblock}
\includegraphics[page=341,pagebox=mediabox,trim={76.18bp 110.34bp 67.57bp 522.39bp},clip,max width=0.94\textwidth,max height=0.38\textheight]{Marchisio2013.pdf}
\end{sourceblock}
其中 ̇Σα = ̇Σ(Σα)、Dα = D(Σα)、βα,γ = β(Σα, Σγ) 和 aα = a(Σα)，子分布函数的矩为 ̄bk

\begin{sourceblock}
\includegraphics[page=341,pagebox=mediabox,trim={223.95bp 74.76bp 244.42bp 602.13bp},clip,max width=0.94\textwidth,max height=0.38\textheight]{Marchisio2013.pdf}
\end{sourceblock}
{\small Ψkb(Ψ|Ψα)dΨ。表 7.1 报告了子分布函数及其相应矩变换的几个示例。当破损产生两个碎片时，分布函数由两个delta组成\par}
\begin{center}\small 表7.1。公共子分布函数及其矩变换。\end{center}
\begin{sourceblock}
\includegraphics[page=342,pagebox=mediabox,trim={58.00bp 450.47bp 58.23bp 72.20bp},clip,width=\textwidth,max height=0.72\textheight]{Marchisio2013.pdf}
\end{sourceblock}
函数以表征片段的内部坐标值为中心。如果断裂导致形成两个相等的碎片，则两个 δ 函数重叠，而在侵蚀的情况下，一个 δ 函数以 Ψ = Ψ0（即单体）为中心，另一个 δ 函数以 Ψ′ - Υ0 为中心。当形成内坐标比为 1 比 4 的两个碎片时，分布函数由以 xi′(1/5)1/3 和 xi′(4/5)1/3 为中心的两个 delta 函数组成。对于简单的聚集和破损问题，QMBB 具有出色的准确性（Marchisio 等人，2003b，2003c；Marchisio 和 Fox，2005）。在这些情况下，使用少量节点（即 2 ≤ N ≤ 5），因此求解属于矩集的矩演化的少量方程，足以使总体误差保持在可接受的范围内。通常，不仅对于跟踪的时刻，而且对于高阶时刻，误差仍然非常小。事实上，通过外推正交近似，可以成功地使用正交近似来估计未包括在原始矩集中的矩。图 7.2 对此进行了说明，其中将 QBMM 获得的预测与精确解进行了比较。报告的预测对应于 Marchisio 等人中案例 4 的稳态解。 （2003c）。精度根据全局误差函数进行分析，该函数通过使用 N = 2、3 和 4 的 QBMM（因此通过跟踪 4、6 和 8 个矩）量化在预测不同矩阶的 k 阶矩时所犯的误差。从图 7.2 可以看出，不仅预测求积所涉及的矩的误差是可以接受的，而且预测高阶矩的误差也相当小。当使用更多节点时，源项的近似精度更高，并且 NDF 的表示更真实。图 7.3 对此进行了说明，其中将使用截面方法预测的 NDF 与正交近似的节点位置进行了比较，以垂直条示意性表示。可以看出，源项正在驱动 NDF 达到截面法所描述的形状。当仅使用两个节点时，它们的位置是为了捕获 NDF 的主要特征并表示 NDF 的四个矩。当添加更多节点时，它们用于表示更多矩，对应于分布的一条尾部。对这个例子的细节感兴趣的读者可以参考原始作品（Marchisio et al., 2006）。作为一般规则，可以安全地说，对于正常内核和连续变化率，

\begin{sourceblock}
\includegraphics[page=343,pagebox=mediabox,trim={58.00bp 408.90bp 58.23bp 59.08bp},clip,width=\textwidth,max height=0.72\textheight]{Marchisio2013.pdf}
\par\vspace{0.45\baselineskip}
\small 图7.2。全局误差函数 E(k) 与矩阶 k 的关系，用于通过使用 QBMM（从上到下：N = 2、3 和 4）相对于精确解获得有关经历聚集和破碎的颗粒系统的预测。改编自马尔基西奥等人。 （2003c）。
\end{sourceblock}
\begin{sourceblock}
\includegraphics[page=343,pagebox=mediabox,trim={58.00bp 201.94bp 58.23bp 351.18bp},clip,width=\textwidth,max height=0.72\textheight]{Marchisio2013.pdf}
\par\vspace{0.45\baselineskip}
\small 图7.3。对于标准聚合和破损问题，通过截面方法（虚线）和 QBMM（竖线）预测 NDF（从左到右、从上到下）N = 2、3、4 和 5。改编自马尔基西奥等人。 （2006）。
\end{sourceblock}
由此产生的源项是一个平滑函数，其近似精度随着正交近似的节点数量的增加而不断增加。由于计算时间和稳定性问题，正交近似的节点数量 (N) 以及相应的跟踪矩数量 (2N) 不能无限增加。事实上，随着 N 的增加，必须跟踪更大的矩集，从而导致计算时间更长。此外，随着矩数量的增加，数值问题就会出现。这些问题与矩反演算法有关，该算法有时无法计算大 N 值的正交近似。对于大多数已调查的案例，合理的准确度达到 2 ≤ N ≤ 5。对于经历聚集和破损问题的颗粒系统的更多分析和讨论，读者可以参考 Alopaeus 等人的工作。 (2006)、Diemer \& Ehrman (2005)、Gerber \& Mousavi (2007)、Prat \& Ducoste (2006)、Silva 等人。 (2010)、Upadhyay 和 Ezekoye (2003、2005) 和 Wang 等人。 （2005a，2005b）。在内坐标连续变化的标准速率的情况下也可以得出同样的结论。这些速率相对于内坐标的大多数函数依赖性可以通过正交近似轻松处理，正如许多作品中所证明的那样（例如，参见 McGraw (1997) 和 Wright 等人 (2000, 2001a)）。

非常关键的是源项在相空间中高度局域化的情况（例如，参见 Grosch 等人 (2007)）。事实上，在处理高度局部化的现象时，可能会面临在需要节点的相空间区域中没有足够节点的风险。一个典型的例子是粒子消失（由于精细溶解循环），其中负汇项是针对相空间的特定（通常是有限的）值范围制定的。如果没有节点定位在该窗口中，则忽略该现象。此外，由于新节点的位置在其数量增加时不容易预测，因此当节点数量变化时近似源项的误差的行为是不可预测的。在这种情况下，简单地添加更多节点（并求解更多矩演化方程）并不是一个好的策略。相反，使用第 3 章中描述的带有核密度函数的扩展高斯求积可能会产生更好的结果，因为求积的阶数可以独立于传输矩的数量而改变。另一个相关情况是，当矩定义中出现的两个积分极限之一具有非零 NDF 时，QBMM 往往表现不佳。一个典型的例子是，当 NDF 被定义为介于零和无穷大之间的内坐标值（因此矩也被定义在零和无穷大之间），并且粒子受到内坐标负变化率的影响。从给定的初始条件开始，随着粒子减小其内坐标值，NDF 趋向于零，这是积分的两个极限之一。当 NDF 为零时，n(0) 不为空，一些粒子逃离相空间，消失，并在两个边界之一生成非零通量（参见方程（7.77））。这就是粒子溶解或液滴蒸发过程中发生的情况：当粒子达到有限尺寸时，它们离开相空间，导致 NDF 的零阶矩发生变化。为了正确预测这些现象，QBMM 必须与能够高精度估计 n(0) 的重建方法相结合。 Strumendo 和 Arastoopour (2008, 2010)、Massot 等人提出了一些有趣的解决方案。 （2010）和袁等人。 (2012)，其中最后一个计算效率更高、更稳健。下面详细讨论 QBMM 的两个最流行的版本。对最近的阐述和新发展感兴趣的读者可以参考专业文献（Aamir et al., 2009, 2010; Attarakih et al., 2009; Bove et al., 2005; Gimbun et al., 2009; Nagy, 2009; Qamar et al., 2011a; Su et al., 2007, 2008）。

\paragraph*{QMOM}
{\small 当采用这种方法时，k = 0, 1, 2, ... 的矩 mk 的演化方程。。。 , 2N -1 被求解，并且源项在每个时间步被评估\par}
求助于正交近似。所包含的矩的数量理论上可以是任意整数；然而，通常只跟踪 NDF 的前 2N -1 个矩，因为这些足以计算 N 阶的正交近似。当然，需要一个初始条件，并且这可以从 NDF 的初始条件（众所周知）轻松计算出来：28 mk(0) =

\begin{sourceblock}
\includegraphics[page=345,pagebox=mediabox,trim={369.59bp 588.18bp 117.75bp 100.02bp},clip,max width=0.94\textwidth,max height=0.38\textheight]{Marchisio2013.pdf}
\end{sourceblock}
Ψkn(0, Ψ)dΨ。矩演化方程的积分需要在源项的每个时间步进行评估，为此需要能够提供正交近似的权重 wα 和节点 xiα 的矩反演算法。计算正交近似的直接方法是通过其矩的定义：mk = -N

\begin{sourceblock}
\includegraphics[page=345,pagebox=mediabox,trim={99.74bp 516.44bp 358.26bp 167.60bp},clip,max width=0.94\textwidth,max height=0.38\textheight]{Marchisio2013.pdf}
\end{sourceblock}
α，k = 0, 1, 2, .。。 , 2N -1。事实上，N 阶的正交近似由其 N 个权重 wα 和 N 个节点 xiα 定义，并且可以通过取其前 2N 个矩 m0,... 来计算。。。，m2N−1。然而，这不是最好的方法，因为需要非常好的初始猜测。更聪明的方法是根据矩 mk 编写与 NDF 正交的多项式的递归关系。一旦将这种关系写成矩阵形式，就很容易证明多项式的根对应于雅可比矩阵的特征值。 （有关单变量矩反演算法的完整讨论，请参阅第 3.1.2 和 3.1.3 节。）矩反演算法要求矩集可实现。因此，用于求解矩演化方程的 ODE 求解器必须确保，给定可实现矩 mk(0)，矩集 mk(t) 始终是可实现的。实际上，对于齐次系统，源项使用大多数 ODE 求解器生成可实现的矩。然而，对于非均匀系统来说情况并非如此，对于非均匀系统，平流项可能存在问题。我们将在第 8 章中回到这个主题。

\paragraph*{DQMOM}
无需求解矩的演化方程，而是可以直接跟踪正交近似中权重与节点的演化（Marchisio \& Fox, 2005；McGraw \& Wright, 2003）。权重与节点的演化方程可以通过将 NDF 的 δ 函数表示形式正式代入方程 1 来导出。 （7.3）。如果权重与节点是时间的连续函数，则此过程产生

\begin{sourceblock}
\includegraphics[page=345,pagebox=mediabox,trim={137.66bp 206.91bp 67.57bp 456.17bp},clip,max width=0.94\textwidth,max height=0.38\textheight]{Marchisio2013.pdf}
\end{sourceblock}
其中 δ′(Σ) 是广义函数 δ(Σ) 的一阶导数，如 Arfken (1985) 所解释的。如果引入加权节点（或加权节点） ςα = wαψα，则经过一些操作后，方程： (7.113) 可以重写为

\begin{sourceblock}
\includegraphics[page=345,pagebox=mediabox,trim={105.40bp 117.26bp 67.57bp 545.82bp},clip,max width=0.94\textwidth,max height=0.38\textheight]{Marchisio2013.pdf}
\end{sourceblock}
28 例如，在典型的气相或液相粒子形成过程中，在 t = 0 时所有矩都相同。实际上，在这些情况下，在模拟开始时通常不存在粒子，因为它们是通过成核或初始机制产生的。在其他情况下（例如在晶种结晶过程或工业凝固器中），初始时刻对应于最初存在的晶种/单体/颗粒的 NDF。

现在我们将 aα 和 bα 分别定义为权重 wα 和加权节点 ςα 的源项：

\begin{sourceblock}
\includegraphics[page=346,pagebox=mediabox,trim={200.56bp 591.29bp 65.08bp 82.96bp},clip,max width=0.94\textwidth,max height=0.38\textheight]{Marchisio2013.pdf}
\end{sourceblock}
使用这些定义，方程。 (7.114) 可以重写为更简单的形式：

\begin{sourceblock}
\includegraphics[page=346,pagebox=mediabox,trim={131.30bp 536.02bp 65.08bp 130.92bp},clip,max width=0.94\textwidth,max height=0.38\textheight]{Marchisio2013.pdf}
\end{sourceblock}
该方程现在可用于通过应用矩变换来确定未知函数 aα 和 bα。一般来说，DQMOM 可以应用于任何希望控制的独立矩集，唯一的必要条件是矩的数量必须等于未知函数的数量（即 aα 和 bα，或总共 2N 个未知数）。鉴于此

\begin{sourceblock}
\includegraphics[page=346,pagebox=mediabox,trim={214.28bp 468.23bp 68.06bp 210.24bp},clip,max width=0.94\textwidth,max height=0.38\textheight]{Marchisio2013.pdf}
\end{sourceblock}
α，方程的矩变换。 (7.116) 产量

\begin{sourceblock}
\includegraphics[page=346,pagebox=mediabox,trim={175.86bp 415.34bp 65.08bp 251.60bp},clip,max width=0.94\textwidth,max height=0.38\textheight]{Marchisio2013.pdf}
\end{sourceblock}
其中 k = k1, k2, .。。，k2N。从等式可以看出。根据方程（7.117），N 个权重 wα 和 N 个加权节点 ςα 的输运方程的源项可以通过涉及 2N 个整数矩的线性系统来定义。方程中的线性系统。 (7.117) 可以写成矩阵形式，Aα = d，其中 α = [a1 · · · aN b1 · · · bN]T = [aT bT]T，d = [S k1 · · · S k2N−1]T，矩阵 A 的分量为

\begin{sourceblock}
\includegraphics[page=346,pagebox=mediabox,trim={192.57bp 307.30bp 65.08bp 354.03bp},clip,max width=0.94\textwidth,max height=0.38\textheight]{Marchisio2013.pdf}
\end{sourceblock}
aij =
如果（如在 QMOM 中）选择前 2N 个整数矩（即 k = 0, ..., 2N -1），则线性系统的矩阵为

\begin{sourceblock}
\includegraphics[page=346,pagebox=mediabox,trim={82.62bp 188.43bp 65.08bp 428.26bp},clip,max width=0.94\textwidth,max height=0.38\textheight]{Marchisio2013.pdf}
\end{sourceblock}
请注意，A 不依赖于权重 wα。此外，如果节点 xiα 是唯一的，则 A 对于任意 N 都将是满秩的。对于这种情况，权重 wα 和加权节点 ςα 的传输方程的源项就是 α = A−1d。 DQMOM 的过程可以总结如下。正交近似的权重与节点的演化方程，如方程所示。 （7.115），通过计算逆线性系统的源项来求解。使用以下初始条件对方程进行积分： wα(0) = w0

\begin{sourceblock}
\includegraphics[page=346,pagebox=mediabox,trim={67.52bp 86.35bp 65.07bp 586.10bp},clip,max width=0.94\textwidth,max height=0.38\textheight]{Marchisio2013.pdf}
\end{sourceblock}
选择前 2N 个矩（如第 3 章所述），并在选择另一个非序列矩集时直接求解非线性系统。许多实际情况的特点是 NDF 矩的初始条件为零，并且作为权重 wα 和加权节点 ςα 的结果。因此，在第一个时间步中，局部节点 xiα 未定义，矩阵 A 不可逆，因此无法找到源项 a 和 b。一种可能的解决方案是在初始条件中引入“种子”，以便以非常小的权重和加权节点值开始模拟。尽管由于初始值的任意选择，这种方法可能看起来有问题，但已发现它非常稳定且有效。另一种可能性是求解一个非常小的时间步长 δt 的矩演化方程以找到 mk(δt)（因为在这种情况下 mk(0) = 0，mk(δt) 将由成核确定），然后将矩反演算法应用于 mk(δt)。在内坐标连续变化的情况下源项（即漂移项）的计算值得更详细地讨论，因为通过一些数学段落，可以进一步简化。在这种情况下，源项简化为 S (ϋ) = −∂ ̇ϋn/∂ϋ，并且 S k 可以表示为

\begin{sourceblock}
\includegraphics[page=347,pagebox=mediabox,trim={169.51bp 417.68bp 67.57bp 249.76bp},clip,max width=0.94\textwidth,max height=0.38\textheight]{Marchisio2013.pdf}
\end{sourceblock}
将 NDF 表示为以节点为中心的 delta 函数的总和，我们得到 S k = k -N

\begin{sourceblock}
\includegraphics[page=347,pagebox=mediabox,trim={128.61bp 382.67bp 332.62bp 301.44bp},clip,max width=0.94\textwidth,max height=0.38\textheight]{Marchisio2013.pdf}
\end{sourceblock}
α ̇xiαwα，它是整数矩 mk 的源项的 N 点求积近似。在这种情况下，定义 DQMOM 传输方程中源项的线性系统简化为

\begin{sourceblock}
\includegraphics[page=347,pagebox=mediabox,trim={221.98bp 324.65bp 67.57bp 343.42bp},clip,max width=0.94\textwidth,max height=0.38\textheight]{Marchisio2013.pdf}
\end{sourceblock}
其中 b*= b −Pg。对角权重矩阵由 P = diag(w1, . . . , wN) 和 g = [ ̇Σ1 · · · ̇ΣN]T 定义。方程的解。 (7.121) 是 a = 0 且 b = Pg，以便直接从漂移项找到加权节点的源项，同时权重保持不变。对于单变量情况，QMOM 和 DQMOM 产生完全相同的结果（如果权重与节点是时间的连续函数），因此它们可以互换使用。开发 DQMOM 的主要动机是处理多元 PBE，而 QMOM 由于缺乏稳健的矩反演算法而存在问题。

\subsection*{7.4.2 双变量和多变量 PBE}\phantomsection\addcontentsline{toc}{subsection}{7.4.2 双变量和多变量 PBE}
通过扩展上一节的概念，可以导出二元系统的近似值（不再是真正的高斯求积）。由此产生的 N 点求积近似将一般二元矩的定义转换为 (Wright et al., 2001b) mk = -N

\begin{sourceblock}
\includegraphics[page=347,pagebox=mediabox,trim={213.29bp 121.92bp 242.21bp 561.18bp},clip,max width=0.94\textwidth,max height=0.38\textheight]{Marchisio2013.pdf}
\end{sourceblock}
1,αψk2 2,α，其中，对于单变量情况，wα是第α个节点的权重，ψ1,α是第α个节点对于第一个内坐标的位置，ψ2,α是第α个节点对于第二个内坐标的位置。使用这种正交近似，可以以良好的精度估计 k = (k1, k2) 阶矩的源项。例如，在成核、分子生长、扩散、聚集和断裂的情况下（在漂移率和内核取决于内坐标的一般情况下），源项的近似表达式为

\begin{sourceblock}
\includegraphics[page=348,pagebox=mediabox,trim={98.02bp 402.84bp 65.08bp 91.70bp},clip,max width=0.94\textwidth,max height=0.38\textheight]{Marchisio2013.pdf}
\end{sourceblock}
其中，聚合核根据第 α 和 γ 节点的内坐标值进行评估，βα,γ = β(xi1,α, xi1,γ; xi2,α, ϋ2,γ)，并且在第 α 节点处评估破损核 aα = (Σ1,α, Σ2,α)。对于子分布函数的矩，可得到以下通用表达式： b

\begin{sourceblock}
\includegraphics[page=348,pagebox=mediabox,trim={246.98bp 332.64bp 203.34bp 344.25bp},clip,max width=0.94\textwidth,max height=0.38\textheight]{Marchisio2013.pdf}
\end{sourceblock}
1 Σk2 2 b(Σ1,α, Σ2,α|Σ)dΣ。在对称破坏的情况下，这已经在方程中讨论过。 (7.107), 表达式为 b

\begin{sourceblock}
\includegraphics[page=348,pagebox=mediabox,trim={72.49bp 306.03bp 355.50bp 375.46bp},clip,max width=0.94\textwidth,max height=0.38\textheight]{Marchisio2013.pdf}
\end{sourceblock}
1,αψk2 2,α。有关将 QBMM 用于双变量颗粒系统的更多详细信息，请参阅 Buffo 等人的文章。 （2012）和祖卡等人。 （2007）。与单变量系统一样，有两种方法可以解决双变量求积问题：(i) 通过求解一组双变量矩 (QMOM) 的方程并使用矩反演算法计算求积近似值，29 或 (ii) 通过求解近似值的权重和加权节点/节点方程 (DQMOM) 来直接跟踪求积的演化。与单变量问题的主要区别在于选择要表示和跟踪的时刻的自由度较高。事实上，如果对于单变量分布，矩的选择是自动定义的（为了获得真正的 N 阶正交近似，我们必须使用前 2N 个矩），对于双变量和多元分布，理论上可以考虑任何矩集，如 McGraw 等人的工作所示。 (1998)、McGraw \& Wright (2003)、Rosner 等人。 (2003)，以及 Yoon 和 McGraw (2004a，2004b)。尽管矩的选择受到所采用的方法（即 QMOM 与 DQMOM）以及使用 QMOM 时所采用的矩反演算法（即强力 QMOM、张量积 QMOM、CQMOM）的影响，但可以做出一些一般性评论。其中一些评论将在下一节讨论 QMOM 和 DQMOM 的双变量和多变量应用时报告，而定义和使用最佳矩集的机会将在第 3.2.1 节中讨论。

\footnotetext[29]{双变量和多变量问题的矩反演算法在 3.2 节中讨论。}
\paragraph*{QMOM}
QMOM 扩展到双变量和多变量系统是相对简单的。矩集 mk，在方程中定义。 （2.2），其演化由方程决定。 （7.73），可以通过估计方程中的源项来求解。 (7.74) 具有 N 点求积近似。方程式 1 给出了一个示例，说明该方法如何适用于经历内坐标连续变化以及聚集、断裂和成核的简单双变量系统。 （7.122）。当然，需要矩的初始条件，这可以根据初始 NDF 的知识来计算，就像在单变量情况下一样。我们还应该注意到，即使对于时间上不连续的源项（这种情况很少见），QMOM 也是有效的，而 DQMOM 不能用于这种情况。 QMOM 收敛是通过使用从第 3 章详细讨论的矩反演算法中找到的正交近似来实现的。

\paragraph*{DQMOM}
DQMOM 到双变量和多变量系统的扩展使用 N 阶的正交近似。双变量分布的简单情况使用

\begin{sourceblock}
\includegraphics[page=349,pagebox=mediabox,trim={163.30bp 392.37bp 67.57bp 274.58bp},clip,max width=0.94\textwidth,max height=0.38\textheight]{Marchisio2013.pdf}
\end{sourceblock}
DQMOM 方程通过代入方程得到。 (7.123) 代入方程。 （7.104）。如果权重与节点是时间的连续函数，30 我们得到

\begin{sourceblock}
\includegraphics[page=349,pagebox=mediabox,trim={75.53bp 278.22bp 67.57bp 343.35bp},clip,max width=0.94\textwidth,max height=0.38\textheight]{Marchisio2013.pdf}
\end{sourceblock}
未知函数 a = (a1, . . . , aN)、b1 = (b1,1, . . . , b1,N) 和 b2 = (b2,1, . . . , b2,N) 是 DQMOM 方程的源项：

\begin{sourceblock}
\includegraphics[page=349,pagebox=mediabox,trim={161.45bp 215.45bp 67.57bp 458.56bp},clip,max width=0.94\textwidth,max height=0.38\textheight]{Marchisio2013.pdf}
\end{sourceblock}
对于 1 ≤α ≤N，其中加权节点由 ς1,α = wαΣ1,α 和 ς2,α = wαΣ2,α 定义。 DQMOM 源项可以使用矩变换根据阶矩 k = \{k1, k2\} 找到： mk1,k2 =

\begin{sourceblock}
\includegraphics[page=349,pagebox=mediabox,trim={234.87bp 176.19bp 237.90bp 503.17bp},clip,max width=0.94\textwidth,max height=0.38\textheight]{Marchisio2013.pdf}
\end{sourceblock}
1 Σk2 2 nd Σ。如果将矩变换应用于方程： （7.124），方程的右侧。 (7.125) 通过求解线性系统得到

\begin{sourceblock}
\includegraphics[page=349,pagebox=mediabox,trim={84.18bp 124.45bp 67.57bp 542.49bp},clip,max width=0.94\textwidth,max height=0.38\textheight]{Marchisio2013.pdf}
\end{sourceblock}
这里源项 S k1,k2 由等式右侧定义。 (7.104): S k1,k2 =

\begin{sourceblock}
\includegraphics[page=349,pagebox=mediabox,trim={73.25bp 92.95bp 407.75bp 593.26bp},clip,max width=0.94\textwidth,max height=0.38\textheight]{Marchisio2013.pdf}
\end{sourceblock}
1 Ψk2 2 S dΨ。在此再次强调，源项是通过使用来评估的，这一点很重要

\footnotetext[30]{矩可能在时间上连续，而权重和/或节点则不然。然而，对于输运方程为双曲线的非齐次系统，最有可能出现不连续解。}
N 点求积方案，因此有时用上标 (N) 报告。对于双变量情况，有 3N 个未知数 \{aα, b1,α, b2,α\}，因此方程 1 必须使用相同数量的矩。 （7.126）。由此产生的 3N 方程线性系统可以写成矩阵形式：Aα = d，其中 α = [a1 · · · aN b1,1 · · · b1,N b2,1 · · · b2,N]T = [aT bT 1 bT 2]T。有趣的是，这里线性系统的系数矩阵 A 与强力 QMOM 中使用的系数矩阵相同，并且是针对方程 1 中的一般多变量情况定义的。 （3.44）。对于双变量系统（即 M = 2）和 N = 1，混合矩的自然选择是 \{k1, k2\} = \{0, 0\}, \{1, 0\}, \{0, 1\}。因此，DQMOM 方程简化为

\begin{sourceblock}
\includegraphics[page=350,pagebox=mediabox,trim={160.91bp 509.24bp 65.08bp 164.77bp},clip,max width=0.94\textwidth,max height=0.38\textheight]{Marchisio2013.pdf}
\end{sourceblock}
与单变量情况一样，用于近似右侧源项的单点求积通常不准确。因此，通常只有 N > 1 的情况才会令人感兴趣。对于 N = 2，NDF 由两个正交点描述。由于相空间中的每个点 α 由三个参数 \{wα, Ψ1,α, Ψ2,α\} 确定，因此需要六个混合矩才能解决问题。没有独特的方式来选择时刻。如果选择六个最低阶矩 \{k1, k2\} = \{0, 0\}、\{1, 0\}、\{0, 1\}、\{2, 0\}、\{1, 1\} 和 \{0, 2\}，则系数矩阵为

\begin{sourceblock}
\includegraphics[page=350,pagebox=mediabox,trim={136.14bp 342.01bp 65.08bp 279.04bp},clip,max width=0.94\textwidth,max height=0.38\textheight]{Marchisio2013.pdf}
\end{sourceblock}
然而，这一选择是一个糟糕的选择。对于 Ψ1,α 和 Ψ2,α 的任何值，方程式中的矩阵： (7.128) 是单数。这是因为协方差（其中 pα = wα/(w1 + w2)）、σ12 = p1Σ1,1Σ2,1 + p2Σ1,2Σ2,2 −(p1Σ1,1 + p2Σ1,2)(p1Σ2,1 + p2Σ2,2) 与方差相关，σ2 1 = p1Σ2 1,1 + p2Σ2 1,2 −(p1Σ1,1 + p2Σ1,2)2 和 σ2 2 = p1Σ2 2,1 + p2Σ2 2,2 −(p1Σ2,1 + p2Σ2,2)2，由 σ2 12 = σ2 1σ2 2。从几何角度来看，这种等式源于相空间中的两点必然共线。因此，为了控制协方差，需要N≥3。双变量和多变量系统矩的选择问题将在下面和第 3 章中进一步讨论。我们将通过注意到方程的多变量版本来结束对 DQMOM 的讨论。 (7.126) 通过检查如下：

\begin{sourceblock}
\includegraphics[page=350,pagebox=mediabox,trim={132.39bp 124.78bp 65.08bp 478.98bp},clip,max width=0.94\textwidth,max height=0.38\textheight]{Marchisio2013.pdf}
\end{sourceblock}
n=1 ψkn n dψ 线性系统的系数矩阵为方程 1。 （3.44）。选择一组独立的低阶矩 \{k1,...。。，kM\}，方程。 (7.129) 生成一个线性系统，可用于查找 N(M + 1) 个未知源项 \{aα, b1,α, .。。 , bM,α\} 出现在 \{wα, ς1,α, 的 DQMOM 方程中。。。 , ςM,α\} （参见方程（7.125））。下面报告了有关矩集选择的所有相关问题的简要总结，而详细讨论可以在第 3 章中找到。最后但并非最不重要的一点是，对于在相空间有限域上定义的多元 NDF，必须仔细检查可实现性问题（例如，参见 Wang \& Fox (2003) 的工作）。

\paragraph*{矩集的选择}
我们现在简要讨论为 QBMM 的通用应用选择 NDF 矩的问题。首先，对于一个具体问题，需要建立集合内矩的总数以及矩集合的最高阶。这通常由要使用的正交近似的阶数决定。正交近似的节点数量 N 越高，矩集合 N(M + 1) 中包含的矩就越多，因此该集合中包含的矩的阶数越高。不存在选择 N 的先验标准，因此，根据源项的性质，正交近似的阶数必须通过与替代的、更复杂的方法进行比较来确定。矩集的选择也会影响方法的整体稳定性，因此这方面也应该仔细考虑。我们将首先分析双变量情况，然后将结果扩展到多变量情况。讨论将基于矩集的选择对系数矩阵 A 的影响，正如已经提到的，系数矩阵 A 在 QMOM 和 DQMOM 中都起着重要作用。此外，尽管某些矩反演算法的矩集是通过固定正交近似的节点数来自动定义的，但值得讨论的是不同方法需要什么矩集，以及这对最终系数矩阵有何影响。作为一般规则，我们可以从等式推断。根据方程 (7.128)，为了在二阶矩级别正确表示 NDF，节点数量 N 必须大于内坐标 M 的数量。当使用单变量 NDF (M = 1) 时，N 必须至少等于 2 才能预测方差。对于二元 NDF (M = 2)，N 必须至少等于 3 才能正确表示二阶矩（即方差和协方差）。但请注意，如果在等式 1 中替换 (k1, k2) = (1, 1)，则可以对二元情况使用 N = 2。 (7.128) 通过高阶矩（例如 (2, 2)）。事实上，作为一般规则，即使对于 M > 2 的多变量情况，通常也可以找到 N ≤ M 的任何值的矩集，其中 A 是满秩的。如果我们决定用两个节点（即 N = 2）描述二元分布，则问题是用六个变量定义的，即每个内部坐标有两个权重和两个节点，因此必须选择六个矩来生成线性系统。如果我们考虑直到五阶的所有可能的整数矩，我们将有 21 个矩，因此要跟踪的六个矩的多达 54 264 个可能的选择。当然，并非所有这些选择都是有效的。其中 373 个结果形成奇异矩阵。如果我们考虑到我们希望高精度评估的初阶矩的物理意义，则可能的矩组合的数量会大大减少。前三个混合矩实际上与粒子总数 (m0,0) 和内坐标平均值（m1,0 和 m0,1）有关，这些属性通常非常令人感兴趣，需要尽可能准确地了解。 如果我们将分析限制在包含前三个矩的矩集上，我们只需要在 18 个矩的组中选择剩余的三个矩，并且导致非奇异系统的选择减少到九个。所有这些都可以通过选择单个内坐标的任意一组纯矩来获得。为了获得非奇异矩阵 A，选择两个内坐标的混合矩或纯矩就足够了。然而，人们可能还想非常准确地知道二阶纯矩的值（即 m2,0 和 m0,2），它们与两个内坐标的方差有关。 31 在这种情况下，我们只需要从 0 到 5 之间的 21 个阶矩的原始集合中的剩余 16 个矩中选择一个。可以选择的第一个可能的矩是m3,0。有了这组矩，该方法就可以应用，并且数值稳定性令人满意。然而，值得注意的是，在问题的表述中仅使用纯矩存在一些缺点，因为没有正确考虑内坐标之间的相关性。如果有人有兴趣获得有关二元分布函数形状的信息，这一点可能很重要。由于双节点公式对于双变量问题并不令人满意，让我们考虑使用三个节点时涉及的问题，N = 3。在这种情况下，必须跟踪一组九个整数矩，通常从全局五阶矩中选择。32在这种情况下有 293 930 个可能的选择，其中只有 36 个导致非奇异矩阵，A。如果我们再次将分析限制在包括前五个纯矩的集合上((k1, k2) = (0, 0), (1, 0), (0, 1), (2, 0) 和 (0,2))，我们仅通过在矩集 m1,1 中包含两个全局三阶矩和一个全局四阶矩来获得非奇异矩阵，得到 m3,0, m2,1, m4,0 或 m0,3, m1,2 和 m0,4，或者m2,1、m1,2 和 m2,2。值得注意的是，在矩集中可以找到一些规则结构。当二元分布的矩显示在三角表中时，这些结构尤其明显，如 Zucca (2005) 所示（例如，参见表 7.2）。

我们已经讨论了导致可逆矩阵 A 的一组矩的选择。这是一个必要条件，并且可以通过条件和数值稳定性来驱动矩的选择的进一步细化。如果问题变得病态，计算的准确性可能会急剧降低，从而导致数值不稳定。在计算过程中矩阵 A 距离奇异点越远，跟踪和非跟踪时刻的预测就越准确。如果节点彼此接近，则行列式变得接近于零。随着节点数量的增加，发现两个节点彼此接近的概率也随之增加。另一个考虑因素涉及用于计算权重与节点源项的混合矩的顺序。从这个角度来看，应该避免包含不需要知道的高阶矩，并且评估使用分数阶矩以减少行列式的程度的可能性是很有趣的（例如，参见 Zucca (2005) 的工作）。还值得注意的是，当

\begin{center}\small 表7.2。报告二元 NDF 高达全局五阶的不同矩的三角表\end{center}
\begin{sourceblock}
\includegraphics[page=352,pagebox=mediabox,trim={58.00bp 77.23bp 58.23bp 505.27bp},clip,width=\textwidth,max height=0.72\textheight]{Marchisio2013.pdf}
\end{sourceblock}
\footnotetext[32]{矩的全局阶数被定义为相对于每个内部坐标的阶数之和。例如，m0,1 是全局一阶，而 m2,3 是全局五阶。}
两个内坐标相差数量级时，条件数可以高于内坐标相同数量级的情况。作为一般规则，我们可以说，以大小差异较大为特征的内坐标会导致矩阵缩放不良，应避免使用它们。如有必要，可以通过重新调整和标准化来解决这个问题。迄今为止得出的所有结论都是基于数值问题的分析，但应该记住，正如已经指出的那样，跟踪时刻的预测精度通常比集合中未包含的其他时刻的预测精度要好得多。因此，重要时刻应尽可能包含在所选时刻集中。图 7.4 对此进行了说明，其中考虑了简单的聚合和破坏问题。在所考虑的操作条件下，达到了 NDF 的稳态解。图 7.4 将使用 Monte Carlo 方法获得的三阶矩 m3,0 的预测与使用 DQMOM 在 N = 3 和不同矩集下获得的预测进行了比较。所有考虑的集合都包含 m0,0、m1,0、m0,1、m2,0、m1,1 和 m0,2，并且用全局阶数为 3、4 和 5 的三个矩来饱和自由度。从图7.4中可以看出，只有在矩集中包含全局三阶矩才能正确预测m3,0。扩展到多变量情况由于许多因素而变得复杂，并且，特别是对于高维情况，有必要确定可靠的过程。不同领域的应用表明，应该通过跟踪属于第 3.2.1 节中描述的最佳矩集 DQMOM 的矩的演化来进行模拟。总之，当使用最佳矩集时，模拟通常运行顺利并且没有太多问题。当考虑齐次系统时，可以通过求解矩方程 (QMOM) 或正交近似方程 (DQMOM) 来跟踪最佳矩集的演化。

\begin{sourceblock}
\includegraphics[page=353,pagebox=mediabox,trim={58.00bp 117.00bp 58.23bp 365.72bp},clip,width=\textwidth,max height=0.72\textheight]{Marchisio2013.pdf}
\par\vspace{0.45\baselineskip}
\small 图7.4。对于标准聚合与破坏问题，N = 3 的 QBMM 对 m3,0（用其初始值标准化）的演化进行预测，在给定前六个（m0,0、m1,0、m0,1、m2,0、m1,1 和 m0,2）的情况下，从全局三阶、全局四阶和全局五阶矩中选择最后三个时刻。细线，QBMM；粗线：蒙特卡罗方法。改编自 Zucca 等人。 （2007）。
\end{sourceblock}
当采用第一个过程时，应使用 CQMOM 作为矩反演算法（Yuan \& Fox，2011）。当采用第二个过程时，直接跟踪M维、N点求积近似的演化。在空间均匀系统的情况下，这两种方法是等效的。

\subsection*{7.4.3 碰撞KE}\phantomsection\addcontentsline{toc}{subsection}{7.4.3 碰撞KE}
QBMM 在齐次碰撞 KE 中的通用应用非常简单。回想一下，速度 NDF 矩的通用定义是 mijk = m(i, j, k) =

\begin{sourceblock}
\includegraphics[page=354,pagebox=mediabox,trim={75.74bp 476.77bp 65.08bp 172.70bp},clip,max width=0.94\textwidth,max height=0.38\textheight]{Marchisio2013.pdf}
\end{sourceblock}
在强调漂移和不连续事件项的存在时，得到以下方程：

\begin{sourceblock}
\includegraphics[page=354,pagebox=mediabox,trim={83.87bp 408.10bp 65.08bp 259.34bp},clip,max width=0.94\textwidth,max height=0.38\textheight]{Marchisio2013.pdf}
\end{sourceblock}
其中 Πijk 是碰撞积分的矩变换。 KE 的过程与 PBE 的过程相同。首先定义一个矩集。然后，根据这些矩，可以（使用第 3 章中讨论的任何矩反演算法）确定表示 NDF 矩的 N 阶正交近似：mijk = -N 1,αψ j 2,αψk 3,α。通过对方程的源项应用正交近似。 (7.131) 得到如下演化方程：

\begin{sourceblock}
\includegraphics[page=354,pagebox=mediabox,trim={105.09bp 336.60bp 354.06bp 347.44bp},clip,max width=0.94\textwidth,max height=0.38\textheight]{Marchisio2013.pdf}
\end{sourceblock}
\begin{sourceblock}
\includegraphics[page=354,pagebox=mediabox,trim={83.31bp 284.27bp 65.08bp 382.68bp},clip,max width=0.94\textwidth,max height=0.38\textheight]{Marchisio2013.pdf}
\end{sourceblock}
产生代表矩集演化的一组闭微分方程。对于空间均匀的情况，碰撞源项通常表示为两项之和（例如，参见方程（6.181）），被简化为单个项：

\begin{sourceblock}
\includegraphics[page=354,pagebox=mediabox,trim={176.37bp 193.33bp 65.08bp 473.61bp},clip,max width=0.94\textwidth,max height=0.38\textheight]{Marchisio2013.pdf}
\end{sourceblock}
Πijk = 6
函数 I(0) ijk 是从硬球碰撞积分中分析得出的，对确切形式感兴趣的读者请参阅第 6.33 章的表 6.1–6.3 一个关键问题是用 QBMM 描述平衡分布。事实上，由于驱动 NDF 及其矩达到麦克斯韦平衡的非线性碰撞源项是近似的，因此平衡通常无法完美描述。所涉及的误差通常非常小，并且随着节点数量的增加而减小，但可以通过使用一些简单的修正来轻松克服。 Icardi 等人报告了有关各向同性玻尔兹曼方程测试用例的这些修正的详细信息。 （2012）。

\footnotetext[33]{这里不讨论用更简单的方法（例如麦克斯韦粒子或 BGK）表示碰撞积分，因为，正如已经指出的，在这些情况下不会出现闭合问题。}
\section*{7.5 蒙特卡罗方法}\phantomsection\addcontentsline{toc}{section}{7.5 蒙特卡罗方法}
本章到目前为止讨论的矩方法几乎总是涉及某种类型的闭包。因此，为了验证闭包的准确性，有必要将时刻的预测与更复杂的求解方法进行比较。对于单变量 PBE，第 7.2 节中描述的类方法和分段方法是验证矩方法的绝佳选择。然而，我们已经看到，对于多元 PBE，分段方法的实现变得相当繁琐，必须非常小心才能避免错误。对于多元系统，分段方法的一个有吸引力的替代方法是一类称为蒙特卡罗 (MC) 方法的模拟技术。蒙特卡罗（和直接模拟蒙特卡罗 (DSMC)）方法基于对所研究的多相系统行为的人工实现（Kendal，1950）。正如已经提到的，我们感兴趣的是预测由粒子群构成的多相系统的时间演化，这些粒子通常也是空间分布的，受一些连续和不连续事件的控制。在 MC 方法中，每个粒子的演化都会被跟踪。连续事件确定性地表示为内坐标向量 ψ 的连续变化（以特定速率 ̇ ψ 变化），而不连续事件则表示为以特定且规定的概率发生的随机事件。与先前提出的方法中使用的确定性方法的联系是清楚的（Ramkrishna，1981）：底层的多粒子联合 PDF（参见第 4 章）或类似的主密度方程（参见 Ramkrishna（2000）的第 2 章）不是直接平均（导致 GPBE）而是直接求解。该解是通过对随机微分方程采用适当的方法获得的，从而产生许多复杂的解，每个解代表一个实现。这些实现代表必须进行平均的不同样本路径，从而消除低概率事件并最终产生平均值，可以与使用其他确定性方法（例如 CM、MOM、QMOM、DQMOM 等）获得的结果进行比较。这通常是通过多次重复相同的模拟（对于时间相关的问题）或通过在不同时间对单个模拟求平均值（对于稳态问题）来完成的。在本节中，将简要描述可用于验证矩闭合的不同 MC 方法，重点讨论空间均匀系统，尽管扩展到非均匀系统很简单。这将在本节末尾进行快速分析，但更多详细信息可在文献中找到。所有这些方法都基于对粒子初始状态的了解。这在统计意义上是已知的，意味着多相系统的初始 NDF 是指定的。然后在计算域或框中随机创建一组 N 个粒子（表示给定的 NDF），并直接跟踪其演化。不同方法之间的主要区别与模拟过程中时间前进的方式有关。在时间驱动的模拟中，选择时间步长 Δt，然后定义作用于粒子群并在时间间隔内实现的事件的数量和类型（Collins \& Knudsen，1970；Liffman，1992；Spielman \& Levenspiel，1965； Zeitlin 和 Tavlarides，1972a，1972b）。在事件驱动的模拟中，选择发生一个事件，然后相应地提前时间（Shah 等，1976，1977）。该方法不需要显式的时间离散化，并且在模拟过程中计算的时间步长会根据各种过程的速率进行自我调整。方法之间的另一个重要区别与随着时间的推移属于计算域的粒子总数 N 发生的情况有关。事实上，一些不连续的事件（例如
聚集和破碎）改变了域中的粒子数量。例如，在聚集事件之后，N -1 个粒子保留在域中，而在破碎事件之后，必须跟踪 N + 1 个粒子。随着模拟的进行，粒子总数发生变化，而计算域中粒子的总体积（或总质量）保持不变。因此，这些方法称为恒量蒙特卡罗方法。然而，使用这种方法不能无限期地运行模拟，因为一段时间后，如果破损占主导地位，则要跟踪的粒子数量太大，而如果聚集占主导地位，则粒子数量太低而无法保证可接受的统计精度。为了克服这些限制，在计算域中的粒子数增加或减少其初始值两倍的每个时间步，N分别减半或加倍（Kruis等人，2000；Liffman，1992；Maisels等人，2004a，2004b）。这个概念的逻辑扩展是引入计算域的连续重新填充，通过从域中随机选择一个粒子并将其副本放置在聚合腾出的位置，并在破损的情况下用新形成的粒子覆盖旧粒子（Khalili et al., 2010; Lee et al., 2008; Lee \& Matsoukas, 2000; Lin et al., 2002; Matsoukas et al., 2006、2009；史密斯和马苏卡斯，1998）。在聚集、破碎、生长和沉降的情况下对不同的方法进行了比较，结果表明，当它们使用相当数量的粒子运行时，所有方法都以相当的误差水平计算 NDF 的演化。然而，尽管方法的计算成本和总体精度很大程度上取决于细节，但事件驱动方法通常执行速度更快，并且具有更高的精度（Zhao et al., 2007）。让我们非常简要地讨论这些方法的理论基础。由于真实系统由非常大的粒子数密度构成（例如，从毫米和微米粒子的每立方米<109到纳米粒子的每立方米<1018），在MC模拟中仅考虑测试或概念粒子的子系统（例如104-106），对应于真实域的一小部分。当还考虑计算域内的粒子运动时，会实现周期性边界条件，以便逃出域的粒子立即被相同的粒子替换。潜在的隐含假设是子系统的行为复制了整个系统。如果 n0(Σ) = n(0, Σ) 是初始 NDF，则可以使用相应的归一化累积分布轻松确定 N 个粒子总体的初始状态： N0(Σ) = (1/m(0))

\begin{sourceblock}
\includegraphics[page=356,pagebox=mediabox,trim={268.36bp 238.61bp 181.64bp 436.68bp},clip,max width=0.94\textwidth,max height=0.38\textheight]{Marchisio2013.pdf}
\end{sourceblock}
0 n0 dψ，其中M是内坐标向量的维数。事实上，通过生成均匀分布在 0 和 1 之间的随机数，然后识别导致等于该特定随机数的归一化累积分布值的内坐标集，就生成了属于该分布的一个粒子。然后可以执行该过程 N 次，以产生属于初始分布的 N 个粒子。在每个瞬间，系统的状态由内坐标向量 xii(t) 的值来识别，表征每个单独的粒子（其中 i = 1, 2, ..., N）。我们现在假设粒子群在两个后续的不连续事件之间连续改变其状态。没有不连续事件发生的这个时间间隔称为静止间隔 T。从时刻 t 开始的粒子群将在静止期间演化，直到达到时刻 t +T，或者换句话说，直到遇到新的不连续事件，根据以下常微分方程组：

\begin{sourceblock}
\includegraphics[page=356,pagebox=mediabox,trim={171.32bp 50.16bp 65.08bp 626.19bp},clip,max width=0.94\textwidth,max height=0.38\textheight]{Marchisio2013.pdf}
\end{sourceblock}
假设 xii(t) 的初始条件已知。静止间隔是一个随机变量，在继续进行之前，必须量化相应的概率。接下来，仅考虑一个一阶（或单粒子）事件（即破裂）以及一个二阶点过程（例如碰撞和/或聚集）。扩展到任何其他类型的颗粒过程都很简单，所以让我们首先考虑仅破碎。破损频率由破损内核定义。根据这个量，可以估计其中一个粒子在无穷小的时间间隔 dτ 内分裂的概率，从而导致

\begin{sourceblock}
\includegraphics[page=357,pagebox=mediabox,trim={287.98bp 504.16bp 67.57bp 162.79bp},clip,max width=0.94\textwidth,max height=0.38\textheight]{Marchisio2013.pdf}
\end{sourceblock}
Pr\{一个粒子在间隔 dτ 内破裂\} =
方程（7.135）允许我们推导出以下静止间隔概率的累积分布表达式：

\begin{sourceblock}
\includegraphics[page=357,pagebox=mediabox,trim={222.13bp 436.35bp 67.57bp 225.05bp},clip,max width=0.94\textwidth,max height=0.38\textheight]{Marchisio2013.pdf}
\end{sourceblock}
FT(τ) = 1 −exp
累积分布FT(τ)表示静止时间T小于τ的概率。假设内坐标向量的连续变化的影响可以忽略不计，对于该概率的计算，并且通过将碰撞和聚合也作为附加的不连续事件包括在内，得到以下表达式：

\begin{sourceblock}
\includegraphics[page=357,pagebox=mediabox,trim={189.60bp 340.48bp 67.57bp 317.86bp},clip,max width=0.94\textwidth,max height=0.38\textheight]{Marchisio2013.pdf}
\end{sourceblock}
FT(τ) = 1 −exp
通过首先计算概率分布 fT = dFT/dτ，然后计算其相对于 τ 的一阶矩，可以很容易地找到以粒子初始状态为条件的相应预期静止区间，从而得到

\begin{sourceblock}
\includegraphics[page=357,pagebox=mediabox,trim={147.09bp 265.93bp 67.57bp 404.56bp},clip,max width=0.94\textwidth,max height=0.38\textheight]{Marchisio2013.pdf}
\end{sourceblock}
其中βV是计算框构成的体积中的碰撞率。在时间驱动的 MC 模拟中，规定了特定的时间间隔 Δt。该时间间隔必须小于 〈τ〉，并且通常每个时间步长都不同。一旦时间步长被任意固定，就依次分析每个粒子，并决定该粒子是否会经历不连续事件。描述碰撞、聚合和破坏所使用的策略略有不同，因此让我们分别讨论它们。单个粒子（由索引 i 标识）在时间间隔 Δt 内破裂的概率为

\begin{sourceblock}
\includegraphics[page=357,pagebox=mediabox,trim={190.40bp 140.88bp 67.57bp 532.65bp},clip,max width=0.94\textwidth,max height=0.38\textheight]{Marchisio2013.pdf}
\end{sourceblock}
其中 〈τi〉= 1/b(xii(t))。一旦选择了 0 到 1 之间均匀分布的随机数，粒子破碎的特征就是通过方程 1 中计算的概率来表征的粒子。 (7.139) 大于随机数。然后，选定的粒子将分裂，形成新的碎片，这些碎片必须容纳在模拟框中。34

\footnotetext[34]{在简单子分布函数（例如对称二元破坏）的情况下，此过程非常简单，而对于更复杂的子分布函数（例如 beta、uniform 等），可能必须选择额外的随机数}
在碰撞或聚集的情况下，事情会稍微复杂一些，因为当识别出聚集的粒子时，也必须识别出与它们相互作用的粒子。在碰撞或聚集的情况下，不连续事件仅发生在那些表现出用方程式计算的概率的粒子上。 (7.139) 大于生成的随机数。对于此计算，特征时间为 〈τi〉= 1/-N j=i+1 β(Σi(t), Σj(t))。一旦找到碰撞或聚集粒子 i 的索引，就必须识别另一个粒子的索引。这是通过将粒子 i 和 j 相遇的一般概率定义为

\begin{sourceblock}
\includegraphics[page=358,pagebox=mediabox,trim={193.88bp 504.96bp 65.08bp 164.83bp},clip,max width=0.94\textwidth,max height=0.38\textheight]{Marchisio2013.pdf}
\end{sourceblock}
现在生成另一个均匀分布在 0 和 1 之间的随机数 r，并且通过满足以下关系从所有可能的遭遇中选择遭遇：

\begin{sourceblock}
\includegraphics[page=358,pagebox=mediabox,trim={205.90bp 428.97bp 65.08bp 237.97bp},clip,max width=0.94\textwidth,max height=0.38\textheight]{Marchisio2013.pdf}
\end{sourceblock}
在这种情况下，这意味着粒子 k 将遇到粒子 i。在聚合的情况下，将产生一个内坐标矢量满足基本守恒定律（即质量和动量守恒）的单个粒子，而在粒子碰撞的情况下，将生成两个粒子，最终速度根据线动量和角动量以及动能守恒计算。特别是关于后一种情况（粒子碰撞而不聚集）及其与玻尔兹曼和动力学方程的关系的更多细节可以在专业文献中找到（Bird，1963，1994）。在事件驱动的 MC 模拟中，时间步的幅度 Δt = τ 通过从随机变量中采样来确定，该随机变量的累积概率分布由式（1）给出。 （7.137）。当新的不连续事件发生的时间已经确定后，我们必须从所有可能的事件中选择该事件，其过程与刚才描述的过程非常相似。事件的选择由逆方法控制：对可能的事件进行排序，从0到1之间的均匀分布中选取一个随机数，最终选择与随机数最接近概率的事件。从 MC 方法首次应用于一般聚集和破碎问题（Hsu \& Tsao，1992；Laurenzi \& Diamond，1999；Tandon \& Rosner，1999；Van Peborgh Gooch \& Hounslow，1996）开始，研究了一些具体问题，例如液-液分散（Balmelli \& Steiner，2000；Ghalehchian，2002；Goodson \& Goodson \& Kraft，2004）、聚合过程（Saliakas 等，2008；Tong，2001）、火焰中纳米颗粒的形成（Balthasar 和 Kraft，2003；Celnik 等，2007）和研磨（Mishra，2000）。特别关注使用 MC 方法模拟造粒（Braumann 等人，2010；Cameron 等人，2005；Gantt 和 Gatzke，2006；Oullion 等人，2009；Wildeboer 等人，2005）、喷雾干燥（Terrazas-Velarde 等人，2009）和气溶胶过程。 （Kalani 和 Christofides，2002 年；Rosner 等人，2003 年）。结晶和沉淀过程（Haseltine et al., 2005）、胶体系统的自组装（Ethayaraja et al., 2006）和静电沉淀（Zhao \& Cheng, 2008）过程也通过MC方法进行了研究。就新方法论的发展而言，值得一提的是在估计潜在误差源（Li et al., 2000）和减少统计噪声（Zhao et al., 2009）方面的研究工作。同样重要的是引用常量蒙特卡罗方法对破坏的扩展（Friesen \& Dabros，2003）以及 Irizarry 对整体算法的加速（2008a，2008b）。还特别关注多组分系统模拟性能的提高（Laurenzi 等，2002；Zhao 等，2010）以及其他问题的解决（Vikhansky \& Kraft，2004；Zhao 等，2005a，2005b；Zhao 和 Cheng，2009a，2009b）。 还值得一提的是，由于 MC 模拟中产生和可用的大量信息，并且由于易于扩展到多变量问题，MC 方法通常用于精确求解 GPBE，以测试、比较和验证其他方法（例如，参见 Alexopoulos \& Kiparissides (2007) 和 Meimaroglou \& Kiparissides (2007) 的工作）。然而，应该提到的是，由于 NDF 的高阶矩受分布尾部的影响很大，因此需要大量的粒子才能非常准确地估计它们的演化。在这种情况下，引用 MC 模拟和 QBMM 之间的比较（Marchisio 等人，2003b）特别有趣，该比较表明仅使用三个和四个节点求积的 QMOM 和 DQMOM 计算如何与使用大量粒子（即 N > 20 000 个粒子）进行的 MC 模拟一样准确。空间分布式系统的扩展非常直观。空间分布的计算域被划分为多个盒子，每个盒子中有许多粒子。由于平流和扩散而导致的粒子从一个盒子到另一个盒子的运动现在被表示为一个额外的随机事件，其规定的概率由对流和扩散通量决定。

\section*{7.6 均匀 PBE 示例}\phantomsection\addcontentsline{toc}{section}{7.6 均匀 PBE 示例}
我们通过讨论两个示例 PBE 来结束本章，其中涉及 (i) 单变量同质连续系统和 (ii) 双变量同质批处理系统。对于这两种情况，我们对两种求解方法进行了比较，即第一个系统的 QMOM 和 DQMOM 之间以及第二个系统的 CQMOM 和蒙特卡洛方法之间的比较。但请注意，这些示例可以使用 EQMOM 和 ECQMOM 以更高的精度进行计算（Yuan 等人，2012），如第 3.3 节所述。我们首先讨论假设 PBE 同质性的有效性。

\subsection*{7.6.1 关于空间均匀 PBE 的一些说明}\phantomsection\addcontentsline{toc}{subsection}{7.6.1 关于空间均匀 PBE 的一些说明}
在查看同质系统示例之前，分析何时可以将空间均匀性假设用于 PBE 是有用的。对于 PBE，内坐标不包括粒子速度。例如，考虑以下 PBE：

\begin{sourceblock}
\includegraphics[page=359,pagebox=mediabox,trim={91.20bp 102.98bp 67.57bp 518.09bp},clip,max width=0.94\textwidth,max height=0.38\textheight]{Marchisio2013.pdf}
\end{sourceblock}
{\small 该 PBE 以通用形式编写，包含表示累积、物理空间平流、相空间平流、相空间扩散以及二阶、一阶和零阶点过程的术语。 （有关这些过程的更多详细信息，请参阅第 5 章。）让我们\par}
现在假设 PBE 承认稳态解；在分离出物理空间平流项时，我们有（使用指数的符号）

\begin{sourceblock}
\includegraphics[page=360,pagebox=mediabox,trim={91.36bp 558.66bp 65.08bp 87.52bp},clip,max width=0.94\textwidth,max height=0.38\textheight]{Marchisio2013.pdf}
\end{sourceblock}
通过适当选择特征长度和时间尺度，可以使方程（7.143）变得无量纲。对于特定的系统来说，很容易确定特征长度L和特征速度U，从而得到特征平流时间τM=L/U。一旦达到稳定状态，就可以直接计算内坐标的平均值：

\begin{sourceblock}
\includegraphics[page=360,pagebox=mediabox,trim={131.94bp 453.06bp 65.08bp 212.48bp},clip,max width=0.94\textwidth,max height=0.38\textheight]{Marchisio2013.pdf}
\end{sourceblock}
通过使用这些值，可以估计内坐标的特征变化率 ̇ xi(xi)，进而可用于定义每个内坐标的相空间平流的特征时间尺度： τ ̇ xi = [ xi1/ ̇ xi1(xi1), .。。 , ΨM/ ̇ΨM(ΨM)]。类似地，对于扩散，特征时间尺度很容易定义：τD,ij = ψiψ j/Dij。对于点过程的源项，也可以定义一些特征时间尺度。对于分散相的形成速率，时间尺度定义为 τJ = m0,...,0/Jf(xi)。对于二阶点过程，时间尺度为 τII = 1/[β(Ψ, Ψ)m0,...,0]，而 τI = 1/b(Ψ) 用于一阶过程。通过引入归一化量 X = x/L、ze = Ψ/Ψ、V = Uf/U、F(ze) = n(Ψ)Ψ1Ψ2。。。 ψM/m0,...,0, D*= D/D(ψ), ̇ze= ̇ ψ/ ̇ ψ(ψ), J*= Jf/Jf(ψ), β*= β/β(ψ, ψ), b*= b/b(ψ)，代入式 (1) (7.143)，我们得到

\begin{sourceblock}
\includegraphics[page=360,pagebox=mediabox,trim={90.87bp 251.66bp 65.08bp 383.66bp},clip,max width=0.94\textwidth,max height=0.38\textheight]{Marchisio2013.pdf}
\end{sourceblock}
这里需要强调的是，在等式中。 (7.145) 归一化的 NDF F 和流体速度 V 对空间坐标有明确的依赖性，但也对内坐标的变化率 ̇ζ、分散相 J* 的形成速率以及内核 β* 和 b* 存在隐式依赖性。事实上，如第 5 章所述，这些速率和内核取决于系统中逐点变化的流属性。我们开始分析方程。 （7.145）注意到，由于所有项都已标准化，因此不同项的重要性由时间尺度的值决定。例如，如果 τM 远大于 τ ̇xii、τD,ij、τJ、τII 和 τI，则粒子平流比其他粒子过程慢得多，并且左侧的项相对于右侧的项可以忽略不计。相反，当 τM 远小于 τ ̇xii、τD,ij、τJ、τII 和 τI 时，粒子平流比所涉及的颗粒过程快得多，并且左侧的项比右侧的项更重要。

{\small 有趣的是，提醒读者，例如，如果等式： （7.145）描述了分散在液体中并悬浮在搅拌罐中的颗粒的这些特定颗粒过程，然后它们的速率将随点而显著变化。事实上，聚合和\par}
搅拌槽中悬浮颗粒的破碎率和相应的颗粒在叶轮周围非常高，而远离叶轮则非常低。当 τM 趋于零时，平流变得更快，并且通过乘以方程： (7.145) 由 τM 很容易证明 NDF 的空间梯度必定趋于零。因此，我们的第一个结论是，随着 τM 趋于零，NDF 在整个计算域中变得更加均匀。相反，当其他颗粒时间尺度 τ ̇xii、τD,ij、τJ、τII 和 τI 减小时，相应的颗粒过程变得更快。此外，如果这些特征时间尺度之一趋于零，并且如果速率表现出对空间坐标的强烈依赖性，则所得的 NDF 将表现出相同的空间依赖性。因此，我们的第二个结论是，如果 τ ̇xii、τD,ij、τJ、τII 或 τI 趋于零，则 NDF 在域中将变得越来越不均匀。我们对方程的解释。 （7.145）的要点是存在两组具有相反效果的过程。一方面，平流混合粒子并平滑 NDF 的所有梯度，使其在空间上越来越均匀。另一方面，颗粒过程的速率很大程度上取决于空间坐标，在 NDF 中创建并维持这些梯度。一个典型的例子是搅拌罐中的颗粒聚集和破碎（Marchisio 等，2006）。在这种情况下，通常会达到稳定状态，其中颗粒在高剪切速率区域不断破裂并在更多停滞区域聚集。当颗粒的特征混合时间远短于聚集和破碎的特征时间时，颗粒平流比聚集和破碎更快，因此NDF在整个容器中相当均匀。相反，当聚集和破碎的特征时间尺度比混合时间快得多时，混合速度不够快，无法均匀化和平滑生成的梯度，因此容器中各个点得到的 NDF 有很大不同。这为确定系统是否可以作为空间均匀系统进行分析提供了标准。然而，虽然 NDF 的同质性允许我们用单个量来描述系统，但这并不意味着整个容器的速率是恒定的！35 如果我们最终处于第一种情况，我们可以简化方程： (7.142) 使用体积平均 NDF： ˆn = (1/V)

\begin{sourceblock}
\includegraphics[page=361,pagebox=mediabox,trim={178.68bp 301.26bp 308.67bp 386.94bp},clip,max width=0.94\textwidth,max height=0.38\textheight]{Marchisio2013.pdf}
\end{sourceblock}
n dx，其中 V 是所研究系统的总体积。完美混合的情况意味着 NDF 的空间变化非常小，以至于整个容器中的 n ≈ˆn。通过将体积平均过程应用于方程。 (7.142) 在完全混合的条件下我们得到

\begin{sourceblock}
\includegraphics[page=361,pagebox=mediabox,trim={80.72bp 202.00bp 67.57bp 445.83bp},clip,max width=0.94\textwidth,max height=0.38\textheight]{Marchisio2013.pdf}
\end{sourceblock}
其中 ˆ ̇xi、ˆD、ˆJf、ˆβ 和 ˆb 分别是内坐标的体积平均变化率、分散相的形成速率以及二阶和一阶过程的核。方程（7.146）是在将系统描述为空间均匀时必须求解的体积平均 PBE。仅当已知 ˆ ̇ xi、ˆD、ˆJf、ˆβ 和 ˆb 时，才能求解该方程。请注意，这些数量很难估计，因为它们包含隐藏的空间依赖性。再次提及搅拌中的聚集和破碎
35 这种误解在结晶领域尤其常见。对于结晶和沉淀过程，完美混合系统的假设被标记为混合悬浮液、混合产物去除（MSMPR）。通过该模型，结晶器使用空间均匀的 NDF 进行建模，通常称为晶体尺寸分布 (CSD)。然而，CSD 在容器中保持恒定这一事实并不意味着晶体成核、分子生长、聚集和破裂的速率是恒定的。

在储罐中，介观尺度聚集和破损率随空间变化，必须在方程的体积平均率中考虑在内。 （7.146）。例如，正如第 5 章所强调的， ̇ 和 Jf 可能取决于过饱和度以及局部流动和湍流，而容器中的点与点之间的差异可能相当大。估计其体积平均值的一种方法是明确考虑空间依赖性。例如，在质量传递的情况下，使 ̇ xi 强烈依赖于位置的参数可能是湍流耗散率 ε(t, x)，因此 ˆ ̇ xi = (1/V)

\begin{sourceblock}
\includegraphics[page=362,pagebox=mediabox,trim={315.93bp 553.19bp 140.68bp 125.79bp},clip,max width=0.94\textwidth,max height=0.38\textheight]{Marchisio2013.pdf}
\end{sourceblock}
̇xi fε(ε)dε，其中 f(t, ε) 包含湍流耗散率空间分布的信息。这里，fε(ε)dε 表示以湍流耗散率 ε 为特征的容器的体积分数，并且体积分布函数积分为 1，

\begin{sourceblock}
\includegraphics[page=362,pagebox=mediabox,trim={369.28bp 526.54bp 118.07bp 161.65bp},clip,max width=0.94\textwidth,max height=0.38\textheight]{Marchisio2013.pdf}
\end{sourceblock}
fε(ε)dε ≡1，而其一阶矩表示容器内的体积平均湍流耗散率，ˆε =

\begin{sourceblock}
\includegraphics[page=362,pagebox=mediabox,trim={113.14bp 500.77bp 374.21bp 187.43bp},clip,max width=0.94\textwidth,max height=0.38\textheight]{Marchisio2013.pdf}
\end{sourceblock}
εfε(ε)dε。由于 ˆ ̇ xi 和 ε 之间的关系是非线性的，因此不等式 ˆ ̇ ̇ ̇ xi(ˆ ε) 在大多数情况下成立。类似地，对于二阶点过程的体积平均核 \textasciicircum{}β = (1/V)

\begin{sourceblock}
\includegraphics[page=362,pagebox=mediabox,trim={182.71bp 459.82bp 261.39bp 214.75bp},clip,max width=0.94\textwidth,max height=0.38\textheight]{Marchisio2013.pdf}
\end{sourceblock}
β dx 以及一阶过程的体积平均核 \textasciicircum{}b = (1/V) b dx，对于体积平均 PBE 中出现的所有其他速率和核也类似。例如，在聚集和破碎的情况下，可能包含显著空间依赖性的关键参数是剪切速率 Gf，如方程式 1 所示。 （5.170）、（5.136）和（5.136）。在这些情况下，一旦确定了容器中的体积剪切速率分布 fGf，体积平均聚集和破碎内核就可以计算为 \textasciicircum{}β =

\begin{sourceblock}
\includegraphics[page=362,pagebox=mediabox,trim={244.49bp 400.04bp 142.11bp 288.16bp},clip,max width=0.94\textwidth,max height=0.38\textheight]{Marchisio2013.pdf}
\end{sourceblock}
β(Gf) fGf dGf 且 ˆb = b(G) fGf dGf。同样对于涉及的其他变量，一旦识别出包含对空间坐标的依赖性的参数，就可以确定该参数的体积分布并用于计算相应的体积平均速率和核。总之，上面讨论的内容对于 PBE 控制的多相过程具有重要意义。

1. NDF 在整个域中几乎恒定的事实通常意味着系统可以通过求解方程 1 中报告的体积平均 PBE 来用体积平均 NDF 来描述。 （7.146）。然而，这并不意味着方程中出现的速率是已知的。

2. 如果内坐标的变化率、分散相的形成率以及二阶和一阶过程的核根据空间坐标没有表现出强烈的变化，则它们的体积平均值出现在式（1）中。 （7.146）可以很容易地估计。

3. 在所有其他情况下，必须确定速率和核中出现的负责空间依赖性的参数，并且必须通过确定其体积分布来量化它们在整个容器中的变化。后一个量通常可以通过模拟来测量或估计。一旦知道这些分布，它们就可以用于估计方程式中出现的体积平均速率和内核。 （7.146）。

4. 此外，当解决反问题时，即在完美混合的假设下通过实验与模型预测的比较（或者换句话说，使用方程（7.146））来确定体积平均速率和内核时，必须记住，体积平均速率和内核并不是真正的“动力学”参数，而是包含一些流体动力学因素。这就是为什么在将这些速率和内核传输到不同的操作条件或从一个系统到另一个系统时，可以观察到非常差的性能。

\subsection*{7.6.2 QMOM 与 DQMOM 的比较}\phantomsection\addcontentsline{toc}{subsection}{7.6.2 QMOM 与 DQMOM 的比较}
考虑一个具有连续入口和出口的简单均匀系统。相应的单变量 PBE 为

\begin{sourceblock}
\includegraphics[page=363,pagebox=mediabox,trim={201.74bp 571.73bp 67.57bp 101.73bp},clip,max width=0.94\textwidth,max height=0.38\textheight]{Marchisio2013.pdf}
\end{sourceblock}
其中 n(t, xi) 是 NDF，ψ 是内坐标（例如颗粒质量），nin(t, xi) 是入口 NDF，τR 是颗粒的平均停留时间，S (xi) 是代表不同颗粒过程的源项。当应用矩变换时，获得矩 mk(t) 的方程：

\begin{sourceblock}
\includegraphics[page=363,pagebox=mediabox,trim={197.20bp 489.27bp 67.57bp 184.56bp},clip,max width=0.94\textwidth,max height=0.38\textheight]{Marchisio2013.pdf}
\end{sourceblock}
矩可以从方程中找到。 (7.148) 使用 k = 0 的 QMOM，.。。 , 2N −1 和给定的初始条件 mk(0) = mo k。源项通过使用 N 个节点 xiα 和 N 个权重 wα 的正交近似来封闭。后者是用单变量矩反演算法（即 PD 或 Wheeler 算法）计算的。让我们定义一个新的源项，将方程的整个右侧组合在一起。 （7.148）：

\begin{sourceblock}
\includegraphics[page=363,pagebox=mediabox,trim={196.00bp 391.83bp 67.57bp 282.00bp},clip,max width=0.94\textwidth,max height=0.38\textheight]{Marchisio2013.pdf}
\end{sourceblock}
如果我们使用 DQMOM 求解问题，则 N 个权重 wα 和 N 个加权节点 ζα = wαψα 的传输方程为

\begin{sourceblock}
\includegraphics[page=363,pagebox=mediabox,trim={225.06bp 306.43bp 224.73bp 343.61bp},clip,max width=0.94\textwidth,max height=0.38\textheight]{Marchisio2013.pdf}
\end{sourceblock}
其中 α = 1，.。。 、N 以及初始条件 wα(0) 和 xiα(0) 通过应用单变量矩反演算法根据 mo k 计算得出。源项a*

\begin{sourceblock}
\includegraphics[page=363,pagebox=mediabox,trim={226.57bp 262.27bp 128.14bp 408.48bp},clip,max width=0.94\textwidth,max height=0.38\textheight]{Marchisio2013.pdf}
\end{sourceblock}
α 根据线性系统 Aα*= d* 计算，其中 α*= [a* * 0, .。。 , S * 2N−1]T。如果我们现在写出方程式中所示的源项。 (7.149) 并利用线性系统的性质，我们得到 α*= A−1(din −dout + d) = αin −αout + α，其中 αin = [ain 1, .。。 , bin N]T 是表示进入系统的粒子的源项的一部分，αout = [aout

\begin{sourceblock}
\includegraphics[page=363,pagebox=mediabox,trim={95.84bp 214.06bp 287.27bp 458.85bp},clip,max width=0.94\textwidth,max height=0.38\textheight]{Marchisio2013.pdf}
\end{sourceblock}
1 , .。。，约 N ]T 是代表粒子离开系统的部分，α = [a1, .。。 , aN, b1, .。。 , bN]T 是代表颗粒过程的部分。如方程式所示。 (7.149), din = [min 0/τR, .。。 , min 2N−1/τR]T, doout = [mout 0 /τR, .。。，mout 2N−1/τR]T，并且 d = [S 0, .。。，S 2N−1]T。利用矩阵 A 仅取决于节点值的事实，可以进一步简化：

\begin{sourceblock}
\includegraphics[page=363,pagebox=mediabox,trim={202.65bp 108.77bp 67.57bp 539.11bp},clip,max width=0.94\textwidth,max height=0.38\textheight]{Marchisio2013.pdf}
\end{sourceblock}
其中 α = 1，.。。，N。使用任一方程解决问题。 （7.148）或等式。 （7.150），应该得到相同的解决方案。 36

\footnotetext[36]{有趣的是，如果源项 αin 是通过使用矩阵 A 计算的，则矩阵 A 不是由节点 xiα 确定，而是由节点 xiin 确定}
{\small α 对应于入口矩（或者换句话说，由\par}
作为一个例子，考虑以质量为内坐标和矩 min k = mk(0) = 1 且 k = 0 的单分散粒子群。。。，2N−1。这群粒子被连续地送入一个系统，其中粒子经历聚集和具有恒定内核的对称二元破碎。描述矩演化的方程是

\begin{sourceblock}
\includegraphics[page=364,pagebox=mediabox,trim={83.53bp 541.58bp 75.90bp 125.36bp},clip,max width=0.94\textwidth,max height=0.38\textheight]{Marchisio2013.pdf}
\end{sourceblock}
当 k = 0 时，.。。 , 2N −1 且其中 βo = 1 且 ao = 0.02 是恒定的聚集和破碎核。本例中使用二元对称破坏，与 Marchisio 等人的案例 1 类似。 （2003c）。如果使用 N = 2 的 QMOM，则会求解四个方程，对应于图 7.5 中所示的 NDF 的前四个矩。可见这是一个以聚合为主的情况；事实上，m0 减少，而 m2 和 m3 增加，并且 m1 保持不变，因为颗粒质量在聚集和破碎过程中保持不变。两个节点和两个权值对应的演化如图7.6所示。如果使用 DQMOM37，则等式： (7.150) 得到解决，导致与图 7.5 和 7.6 中所示的完全相同的演变。

\subsection*{7.6.3 CQMOM 与蒙特卡罗的比较}\phantomsection\addcontentsline{toc}{subsection}{7.6.3 CQMOM 与蒙特卡罗的比较}
现在让我们考虑以两个内坐标为特征的粒子群，即尺寸 Ψ1 和成分 Ψ2。组成被解释为颗粒内一种成分的质量。在聚集过程中，如果 xi′ 1 和 ̃ψ1 是导致尺寸为 xi1 的颗粒的聚集颗粒的尺寸，并且 xi′ 1 和 ̃ xi1 是导致形成成分为 xi1 的颗粒的聚集颗粒的组成，则以下关系成立： xi3 1 = ψ′ 1 3 + ̃ψ3 1 和 xi2 = ψ′ 2 + ̃ ψ2。颗粒被限制在均匀区域内并发生聚集和破碎。 NDF n(t, Ψ1, Ψ2) 的演变由下式描述

\begin{sourceblock}
\includegraphics[page=364,pagebox=mediabox,trim={97.06bp 215.51bp 65.07bp 392.79bp},clip,max width=0.94\textwidth,max height=0.38\textheight]{Marchisio2013.pdf}
\end{sourceblock}
将矩反演算法应用于入口矩集 min k )，将得到以下方程组：

\begin{sourceblock}
\includegraphics[page=364,pagebox=mediabox,trim={94.24bp 136.89bp 199.98bp 508.27bp},clip,max width=0.94\textwidth,max height=0.38\textheight]{Marchisio2013.pdf}
\end{sourceblock}
其中通过应用矩求逆算法赢得 k。然而，这组方程不会产生正确的矩演化，因为只有 m0 和 m1 会被正确预测，而所有其他矩（即 k ≥ 2 的 mk ）将不可避免地是错误的。 37 使用 DQMOM 时，我们求解节点和权重的方程，并且为了计算源项，必须求解线性系统。如果两个或多个节点相同，则该线性系统是奇异的，并且，特别是对于 N ≥ 3，以完全单分散分布开始模拟可能会导致问题。例如，对于正在调查的案例，原始矩集为 mk(0) = 1（对于 k = 1）。。。 、 3，但对于 DQMOM，则使用以下矩集：m0 = m1 = 1、m2 = 1.009 和 m3 = 1.0027。这些矩代表了稍微多分散的粒子群，实际上相当于完全单分散的粒子，但导致了更稳定的计算。

\begin{sourceblock}
\includegraphics[page=365,pagebox=mediabox,trim={58.00bp 464.80bp 58.23bp 59.88bp},clip,width=\textwidth,max height=0.72\textheight]{Marchisio2013.pdf}
\par\vspace{0.45\baselineskip}
\small 图7.5。使用 QMOM 和 DQMOM 进行聚合和破坏的连续同质系统的前四个矩的演化。
\end{sourceblock}
\begin{sourceblock}
\includegraphics[page=365,pagebox=mediabox,trim={58.00bp 250.08bp 58.23bp 272.32bp},clip,width=\textwidth,max height=0.72\textheight]{Marchisio2013.pdf}
\par\vspace{0.45\baselineskip}
\small 图7.6。使用 QMOM 和 DQMOM 进行聚合和破坏的连续同质系统的权重与节点的演化。
\end{sourceblock}
其中所有术语都是根据前面章节中解释的规则导出的。从方程。 (7.151) 混合矩 mk,l 的演化方程可以很容易地导出：

\begin{sourceblock}
\includegraphics[page=365,pagebox=mediabox,trim={77.89bp 114.53bp 67.57bp 513.93bp},clip,max width=0.94\textwidth,max height=0.38\textheight]{Marchisio2013.pdf}
\end{sourceblock}
当方程。 (7.152) 用 CQMOM 求解，NDF 近似为

\begin{sourceblock}
\includegraphics[page=365,pagebox=mediabox,trim={143.86bp 51.10bp 146.39bp 615.84bp},clip,max width=0.94\textwidth,max height=0.38\textheight]{Marchisio2013.pdf}
\end{sourceblock}
其中 N1 是为颗粒尺寸选择的节点数，N2 是为颗粒成分选择的节点数，从而得到 N1N2 节点的正交近似值。将此近似应用于方程。 (7.152) 产量

\begin{sourceblock}
\includegraphics[page=366,pagebox=mediabox,trim={109.29bp 506.88bp 65.08bp 105.93bp},clip,max width=0.94\textwidth,max height=0.38\textheight]{Marchisio2013.pdf}
\end{sourceblock}
其中 βα,γ,δ,σ = β(Σ1,α, Σ2;α,γ, Σ1,δ, Σ2;δ,σ), aα,γ = a(Σ1,α, Σ2;α,γ)，子分布函数的矩为

\begin{sourceblock}
\includegraphics[page=366,pagebox=mediabox,trim={161.64bp 445.66bp 158.94bp 225.78bp},clip,max width=0.94\textwidth,max height=0.38\textheight]{Marchisio2013.pdf}
\end{sourceblock}
对于这个具体示例，考虑恒定聚合和破损内核的情况，即 βα,γ,δ,σ = β0, aα,γ = a0 和 beta 子分布函数

\begin{sourceblock}
\includegraphics[page=366,pagebox=mediabox,trim={130.12bp 367.06bp 127.42bp 294.19bp},clip,max width=0.94\textwidth,max height=0.38\textheight]{Marchisio2013.pdf}
\end{sourceblock}
导致

\begin{sourceblock}
\includegraphics[page=366,pagebox=mediabox,trim={137.34bp 324.82bp 134.64bp 349.80bp},clip,max width=0.94\textwidth,max height=0.38\textheight]{Marchisio2013.pdf}
\end{sourceblock}
作为初始条件，我们采用第 3.2.3 节中的二变量高斯分布，其中 μ1 = 4 × 10−3 m、μ2 = 2.4 × 10−7 mol、σ1 = 0.15 和 σ2 = 0.15。现在让我们使用 CQMOM 和 MC 方法研究初始分布在聚合和破坏过程中的演化。 MC 模拟是使用第 7.5 节中报告的常数 MC 方法进行的，名义上有 10 000 个

\begin{sourceblock}
\includegraphics[page=366,pagebox=mediabox,trim={58.00bp 81.03bp 58.23bp 452.48bp},clip,width=\textwidth,max height=0.72\textheight]{Marchisio2013.pdf}
\par\vspace{0.45\baselineskip}
\small 图7.7。双变量分布 (a) 如蒙特卡罗模拟所示，在模拟开始时和 (b) 当达到稳定状态时。
\end{sourceblock}
\begin{sourceblock}
\includegraphics[page=367,pagebox=mediabox,trim={58.00bp 474.80bp 58.23bp 60.02bp},clip,width=\textwidth,max height=0.72\textheight]{Marchisio2013.pdf}
\par\vspace{0.45\baselineskip}
\small 图7.8。蒙特卡罗模拟中表示的二元分布与 CQMOM 表示的比较，在模拟的中间时刻，(a) N1 = 3、N2 = 1 和 (b) N1 = 3、N2 = 2。
\end{sourceblock}
\begin{sourceblock}
\includegraphics[page=367,pagebox=mediabox,trim={58.00bp 179.75bp 58.23bp 274.50bp},clip,width=\textwidth,max height=0.72\textheight]{Marchisio2013.pdf}
\par\vspace{0.45\baselineskip}
\small 图7.9。蒙特卡洛模拟和 CQMOM 预测的低阶矩演化的比较：黑实线，MC 方法；白色圆圈，CQMOM，N1 = 3，N2 = 1；和白色三角形，CQMOM，其中 N1 = 3 且 N2 = 2（以通用矩 Mk,l 为单位：mk−3 moll）。
\end{sourceblock}
颗粒。更多细节可以在 Buffo 等人的工作中找到。 （2013）。图 7.7 报告了模拟开始时和达到稳定状态时的二元分布（如 MC 模拟所示）。可以看出，一开始分布是沿着两个内坐标展开的。然而，一旦开始聚集和破碎，粒子群就会均匀化并呈现立方体形状，代表成分和尺寸之间的关系（即 Ψ2 ∝Ψ3
1）。看看如何用 CQMOM 表示分布是很有趣的。图 7.8 报告了 MC 模拟中中间时间的双变量分布，并将其与 CQMOM 中的隐式假设进行比较。可以看出，当 N1 = 3 且 N2 = 1 时，粒子群分为三组，每组都以不同的节点值 Ψ1 和 Ψ2 为特征。相反，当 N1 = 3 且 N2 = 2 时，对于 Ψ1 的每个节点，识别出 Ψ2 的两个节点。 MC 模拟和 CQMOM 预测的矩演化的比较如图 7.9 所示。可以看出，m3,0 和 m0,1 是恒定的，因为它们在聚集和破坏过程中是守恒量，而所有其他时刻均按预期演化。有趣的是，当使用 N1 = 3 和 N2 = 1 总共三个节点时，以及当使用 N1 = 3 和 N2 = 2 总共六个节点时，MC 和 CQMOM 之间的差异非常小。事实上，在这两种情况下，CQMOM 提供的求积似乎可以相当准确地评估跟踪时刻的源项。用于方程 1 中报告的 ODE 集时间积分的数值方案也发挥着至关重要的作用。 （7.153）。事实上，这种数值方案一定不能引入额外的误差，这些误差会破坏矩集，将可实现的矩集转变为不可实现的矩集。对于空间均匀系统模拟中所涉及的常微分方程的数值积分来说，这通常不是问题，而在处理第 8 章中描述的非均匀系统时，问题就变得很严重。对矩可实现性问题感兴趣的读者应该参考第 3 章。一个单独的问题是求积近似的可实现性。例如，对于所研究的情况，由于 xi1 和 xi2（代表颗粒的尺寸和化学成分）均为正，因此必须确保采用可实现的正交来对相空间进行采样。在这种情况下，可实现的正交具有产生合理值的正节点。对于这个特定问题，事实证明，只有在跟踪矩集中包含一些混合矩（例如 m2,1）时，所探索的相空间部分才完全可实现。这意味着不包括 m2,1 的选项 N1 = 2 和 N2 = 1 可能会导致问题（例如，当涉及传质和化学反应时），而选项 N1 = 3 和 N2 = 1 是更好的选择。

\chapter{非均匀系统的矩方法}
在本章中，我们讨论将矩方法应用于空间非均匀系统的具体问题。我们特别关注平流、扩散和自由传输引起的矩集的空间传输。在第七章中，我们彻底讨论了与相空间输运相关的问题，在这里，我们将仅讨论这些影响空间输运的术语。在第 8.1 节中，简要回顾了空间非均匀系统出现的主要建模问题。在接下来的章节中，我们分别讨论（i）第8.3节中的非齐次群体平衡方程（PBE）（即内部坐标不包括或影响速度），（ii）第8.4节中的非齐次动力学方程（KE）（即唯一的内部坐标是速度），以及最后（iii）第8.5节中的完全非齐次广义群体平衡方程（GPBE）的矩方法。为每种类型的系统提供了具体的例子和相应的离散化公式，以便读者充分理解模拟非齐次系统时出现的问题。贯穿整章的一个重要主题是可实现矩集的问题，以及空间传输如何影响可实现性。因此，为了获得数值问题的明确示例，我们在第 8.2 节中引入了基于动理学的有限体积方法 (KBFVM)。尽管如此，读者应该记住，这些数值问题对于矩传输来说是通用的，并且会在所有空间离散方法中出现。最后，我们在本章中集中讨论基于正交矩方法（QBMM）所需的矩集的空间传输。正如第 7 章中所见，QBMM 在处理 GPBE 中的相空间传输项方面具有许多优势。正如我们将在本章中发现的，QBMM 也是解决空间传输带来的可实现性问题的有效工具。

\section*{8.1 空间建模问题概述}\phantomsection\addcontentsline{toc}{section}{8.1 空间建模问题概述}
当对非齐次系统使用矩方法时，矩集由于平流、扩散和自由传输而在物理空间中传输。由于矩输运方程是从 NDF 的传递方程导出的，因此矩传递问题与 NDF 的传递问题密切相关。将 NDF 表示为 n(t, x, ψ)，空间传输过程涉及对于内坐标 ψ 固定值的 n(t, x, ψ) 的变化。例如，具有恒定速度 u 的平流，数学上表示为

\begin{sourceblock}
\includegraphics[page=370,pagebox=mediabox,trim={214.28bp 591.30bp 65.07bp 83.05bp},clip,max width=0.94\textwidth,max height=0.38\textheight]{Marchisio2013.pdf}
\end{sourceblock}
只需将 NDF 从其原始值 n(0, x, Ψ) 移动到新值 n(t, x, Ψ)，此时精确解为 n(t, x, Ψ) = n(0, x −ut, Ψ)。请注意，精确解意味着 NDF 在物理空间中移动而不改变其在相空间中的形状。此外，由于矩 mk(t, x) 是 NDF 的线性变换，我们从等式中发现： (8.1) 矩的平流式为

\begin{sourceblock}
\includegraphics[page=370,pagebox=mediabox,trim={207.69bp 498.70bp 65.07bp 175.66bp},clip,max width=0.94\textwidth,max height=0.38\textheight]{Marchisio2013.pdf}
\end{sourceblock}
其精确解为 mk(t, x) = mk(0, x −ut)，这又是精确 NDF 的线性变换。不幸的是，方程的数值解。 （8.1）和（8.2）永远不会精确，正是这些数值误差导致了矩方法的严重问题。值得强调的是，这些错误与方程的推导无关。 (8.2) 来自等式。 (8.1)，或者使用 QBMM 中使用的闭包来封闭相空间项。相反，它们是平流数值近似的直接结果。熟悉数值平流算法的读者将意识到一些最相关的问题，例如数值扩散（在一阶方案中）和由于振荡而导致的最大/最小界限的损失（在高阶方案中）。与用于独立标量平流的经典 CFD 算法不同，矩集中的标量（即矩）是相关的，因此不能独立截断（或限制），否则会面临矩集不再可实现的风险。在本节的其余部分中，我们将介绍使用矩方法与空间传输相关的主要建模问题。首先，我们讨论 NDF 和矩集（与上面讨论的数值误差有关）的可实现性。其次，我们介绍了非齐次 KE 中发生的粒子轨迹穿越 (PTC) 现象（并且由 NDF 精确捕获），并描述了它如何导致矩输运方程中的闭合问题。接下来，我们研究与 GPBE 中空间和相空间传输之间的耦合相关的问题（即由于速度和内坐标（例如粒子体积）之间的相关性）。最后，我们介绍了用于求解与 QBMM 相关的矩输运方程的 KBFVM，并简要讨论了如何使用它们来确保可实现性以及捕获 PTC 和处理耦合矩。

\subsection*{8.1.1 可实现性}\phantomsection\addcontentsline{toc}{subsection}{8.1.1 可实现性}
正如上面的讨论所表明的，我们关心两种类型的可实现性：（i）NDF 的可实现性和（ii）时刻的可实现性。理论上，由于可实现的 NDF 生成可实现的矩，因此我们可以将我们的讨论限制为类型 (i) 可实现性。然而，在矩方法中，我们将在不知道真正的 NDF 的情况下传输矩，因此类型 (ii) 可实现性将是最重要的问题。其中一些问题已经在第 3.1.4 节中讨论过，特别关注评估和恢复矩集可实现性的算法。根据定义，如果 NDF 对于内坐标的所有可能值都是非负的，则 NDF 是可实现的：对于所有 ψ ∈Ω ψ，n(t, x, ψ) ≥0。因此，NDF 空间传输的主要问题是离散传输算法绝不能产生负 NDF：如果对于任何 ε n(t, x, ϋ) < 0，则 NDF 是不可实现的。对于具有一阶精度的经典平流方案，确保 NDF 保持可实现相对简单；然而，高阶方案的情况并非如此。由于一阶方案表现出过度的数值扩散，因此只有可实现的高阶平流方案才对计算 NDF 感兴趣。实际上，由于内坐标的数量和 NDF 的高维性，GPBE 的直接数值求解通常是不可行的。因此，开发数值方案时主要关注的是相应矩集的可实现性。

正如前面章节中所解释的，矩集是矩的向量。例如，对于单变量 QMOM，我们使用矩集 M = (m0, m1, ..., m2N−1)，其中 N 是正交节点的数量。如果将矩集的分量用作 2N 维矩空间的坐标，我们可以绘制对应于特定 NDF 的矩空间中所有点的集合。一般来说，矩空间是凸的，因此如果 M1 和 M2 是矩空间中的两个可实现的“点”，那么对于 0 ≤ α ≤ 1 由 M3 = αM1 + (1 − α)M2 定义的“点”就是一个可实现的矩集。更一般地，可实现矩集的任何凸组合1都是可实现的。因此，证明数值算法仅生成可实现的矩集的一种便捷方法是表明它使用前一时间步的矩集（例如不同空间位置处的矩集）的凸组合来更新矩。一般来说，矩空间的边界可能非常复杂，使得检查特定的多元矩集 M 是否可实现变得困难。对于有界区间上 NDF 的单变量矩的特殊情况，可以使用 M 的非线性变换来定义规范矩集 P (Dette \& Studden, 1997)，其中矩空间是单位立方体：P ∈[0, 1]2N。 （规范矩与方程（3.4）中定义的 Hankel-Hadamard 行列式相关。）因此，对于单变量矩集，可以通过用规范矩来制定可实现的空间传输方案（Kah 等人，2012）。单变量矩集的另一种可能性是使用第 3.1.4 节中描述的矩校正算法。对于多元矩集，已知可能性的范围要有限得多，并且大多数都依赖于上面提到的凸矩组合。为了描述单变量 PBE 的平流，有必要使用方程 1 进行平流。 (8.2) 长度为 2N 的矩集 M，如果该矩集对应于非负 NDF n(t, x, xi)，则它对于所有 t 和 x 都是可实现的。对于多元 PBE，适用完全相同的条件，但现在 QBMM 中使用的矩集的维数要大得多。实际上，在使用 QBMM 求解 PBE 时，可实现性如此重要的原因是，只有可实现的矩集才能与第 3 章中描述的矩反演算法一起使用。因此，如果空间传输在某个空间位置 x 生成不可实现的矩集，则无法应用 QBMM 来评估矩源项，并且计算将中止。简而言之，当且仅当保证离散空间传递方案产生可实现的矩集时，使用 QBMM 求解矩输运方程的整个过程才是稳健的！因此，我们的立场是，虽然有时可以进行临时修正来恢复可实现的矩，但最好的方法是仅采用可以在数学上证明矩集总是可实现的数值算法。此外，对于空间传输，我们将需要高阶平流方案，以最大限度地减少数值扩散。此后，我们将把具有这些属性的算法称为可实现的高阶空间传输方案。

\footnotetext[1]{凸组合是系数 αn ≥0 且满足 - n αn = 1 的线性组合。}
另一种与空间传输不直接相关的可实现性问题是可实现求积的构造。对于在凸域 Ω 上定义的单变量 NDF，高斯求积算法保证返回一组 N 节点，对于所有 α ε N 具有属性 ψα ε Ω。换句话说，节点对应于来自 NDF 的可实现样本。对于多元求积，即使用于构造求积的矩集是可实现的，也不能保证所有节点都位于 Ω ψ 内（当然，除非 Ω ψ 是整条实线）。换句话说，可实现的多元矩集并不总是生成可实现的节点。当 Ω 是有限域时，这一点非常重要，因为 Ω 之外的点通常没有物理意义。例如，如果内坐标是化学质量分数，则它们必须为非负且总和为 1。采用负质量分数或大于1的质量分数来评估化学反应速率将产生无意义的值。因此，使用不可实现的节点来评估相空间源项（如第 7 章所述）可能会导致较大的正交误差，从而损害模拟。因此，用于从可实现矩集构造多元求积的矩反演算法必须验证所得求积是否可实现。这当然与用于构建它的矩集有关（如第 7.6.3 节中所述）。然而，我们重申，空间传输不是决定矩反演算法的因素，因此本章的重点将放在可实现的矩集上。

\subsection*{8.1.2 粒子轨迹穿越}\phantomsection\addcontentsline{toc}{subsection}{8.1.2 粒子轨迹穿越}
粒子轨迹穿越 (PTC) 的发生与无碰撞 KE（以及扩展的 GPBE）中的自由传输项相关，并导致难以用欧拉求解器捕获的多速度状态（Sachdev 等人，2007 年；Saurel 等人，1994 年）。在一维速度相空间中，NDF n(t, x, v) 的最简单 KE 为

\begin{sourceblock}
\includegraphics[page=372,pagebox=mediabox,trim={218.30bp 283.67bp 65.07bp 390.68bp},clip,max width=0.94\textwidth,max height=0.38\textheight]{Marchisio2013.pdf}
\end{sourceblock}
其中第二项称为免费运输。该 KE 的解析解为 n(t, x, v) = n(0, x −vt, v)，与方程 1 的关键区别在于： （8.1）式中速度 v 既是内坐标又是传输速度。这意味着 NDF 不会简单地在物理空间中传输而不改变形状，而是自由传输会改变 NDF 的空间位置及其形状。这导致了 PTC 的反直觉现象，其中密度“数据包”可以在物理空间中的单个点交叉。例如，考虑初始 NDF n(0, x, v) = δ(x + 1)δ(v −1) + δ(x −1)δ(v + 1)，它对应于两个数据包，一个位于 x = −1，速度 v = 1，另一个位于 x = 1，速度 v = −1。矩的初始值为 mk(0, x) = δ(x + 1) + (−1)kδ(x −1)。使用方程的解析解。 (8.3), 那么我们有

\begin{sourceblock}
\includegraphics[page=372,pagebox=mediabox,trim={137.86bp 110.22bp 65.08bp 562.66bp},clip,max width=0.94\textwidth,max height=0.38\textheight]{Marchisio2013.pdf}
\end{sourceblock}
其中最终表达式遵循对速度 v 应用狄拉克 δ 函数。这导致了以下矩的解析解：mk(t, x) = δ(x −t + 1) + (−1)kδ(x +t −1)。从方程式中的解决方案。根据方程 (8.4)，我们可以观察到，在 t = 1 时，出现 PTC，其中 n(1, x, v) = δ(x)[δ(v −1) + δ(v + 1)] 且 mk(1, x) = [1 + (−1)k]δ(x)。用言语来说，

\begin{sourceblock}
\includegraphics[page=373,pagebox=mediabox,trim={58.00bp 405.80bp 58.23bp 59.00bp},clip,width=\textwidth,max height=0.72\textheight]{Marchisio2013.pdf}
\par\vspace{0.45\baselineskip}
\small 图8.1。粒子轨迹在 2D 域中交叉的示例。完全发展的颗粒数密度 m00 与平均颗粒速度矢量（m10/m00、m01/m00）一起显示。这个例子是使用 QBMM 与 KBFVM (Yuan \& Fox, 2011) 和可实现的二阶通量 (Vikas et al., 2011a) 进行计算的。请注意，平均粒子速度未定义在密度恰好为零的区域中。由于少量的数值扩散，两个射流外部的密度非常小，但非零。
\end{sourceblock}
在 t = 1 时，NDF 集中在 x = 0 处，并具有两个相等的部分，一个以速度 v = 1 向右移动，一个以速度 v = −1 向左移动。同样，在 t = 1 处，仅在 x = 0 处矩才非零。本质上，两个数据包在 t = 1 处交叉，并且方程的精确解为： (8.3) 重现了这种行为。将 PTC 扩展到更高的空间维度很简单，并且可以在文献中找到许多示例（例如，参见 Desjardins 等人 (2008)）。在图 8.1 中，我们展示了两个射流交叉的 2D 示例。 PTC 出现的一个重要问题是矩输运方程是否可以预测它。为了回答这个问题，我们回到方程。 (8.3) 并在相空间上积分以找到矩输运方程：

\begin{sourceblock}
\includegraphics[page=373,pagebox=mediabox,trim={208.13bp 165.64bp 67.56bp 508.71bp},clip,max width=0.94\textwidth,max height=0.38\textheight]{Marchisio2013.pdf}
\end{sourceblock}
我们立即观察到 KE 中的自由传输项导致了矩输运方程中的闭包问题：如果我们截断 m2N−1 处的矩集，那么为了求解方程： （8.5），我们必须为 m2N 提供一个矩闭包。对于方程中给出的例子。 （8.4），很容易验证矩的解析表达式满足方程（8.4）。 （8.5）。事实上，在这个例子中，我们可以观察到，w1(t, x) = δ(x −t + 1)、w2(t, x) = δ(x + t + 1)、ψ1(t, x) = 1 和 ψ2(t, x) = −1 的两点求积精确地再现了任意 k 的矩 mk(t, x)。因此，从矩集（m0，m1，m2，m3）开始，两点求积就是m4的最优闭包。两点求积所需的矩输运方程为

\begin{sourceblock}
\includegraphics[page=374,pagebox=mediabox,trim={215.14bp 542.51bp 65.07bp 57.97bp},clip,max width=0.94\textwidth,max height=0.38\textheight]{Marchisio2013.pdf}
\end{sourceblock}
其中未闭合矩的最佳选择是 m4 = w1ψ4 1 + w2ψ4 2。如本例所示，在一维问题中，未知矩为 m2N，但在更高维度中，其他未知矩将出现在空间通量中，并使用求积来闭合。总之，使用 QBMM，KE 中的自由传输项通过首先从传输矩集合构造 N 节点求积，然后使用 N 个权重和 N 个节点来估计空间通量中出现的未知矩来闭合。方程中系统的保守通量形式。 （8.6），

\begin{sourceblock}
\includegraphics[page=374,pagebox=mediabox,trim={209.73bp 420.06bp 65.07bp 254.29bp},clip,max width=0.94\textwidth,max height=0.38\textheight]{Marchisio2013.pdf}
\end{sourceblock}
其中F是通量函数，自然就引出了系统是否是双曲线的问题。换句话说，如果我们考虑通量函数 ∂F/∂M 的雅可比行列式，它的特征值都是实数吗？ Chalons 等人研究了这个问题。 (2012)，其中表明系统是弱双曲的，即它具有与节点相对应的 N 对重复（实）特征值。经典的弱双曲系统是无压气体动力学方程：

\begin{sourceblock}
\includegraphics[page=374,pagebox=mediabox,trim={214.29bp 285.03bp 65.07bp 363.31bp},clip,max width=0.94\textwidth,max height=0.38\textheight]{Marchisio2013.pdf}
\end{sourceblock}
众所周知，它会在密度 ρ 中产生 Delta 冲击（Bouchut 等人，2003 年）。有趣的是，方程。 (8.8) 是自由传输的 QBMM 近似，其中 N = 1（即 m0 = ρ、m1 = ρu、w1 = ρ 和 xi1 = u，因此 m2 = w1u2）。如果我们尝试解决式（1）中的 PTC 问题： (8.4) 使用 N = 1，我们会发现对于所有 1 ≤ t，矩将保持“卡”在 x = 0 处：m0(t ≥1, x) = 2δ(x)，m1(t ≥1, x) = 0 ⇒ρ(t ≥1, x) = 2δ(x)，以及 u(t ≥1, x) = 0，这对应于密度中 delta 冲击的定义。然而，我们在上面已经看到，当 N = 2 时，不会发生 Delta 冲击，并且对于本例中的 N ≥ 2 也是如此。一般来说，方程的解。只要给定点 x 处的 PTC 数量大于可用节点 N 的数量，方程（8.7）就会在某些权重 wα(t, x) 中表现出 Delta 冲击。总之，尽管方程（8.7）的弱双曲性质， （8.7）仅针对 1D 相空间（即 KE 中的一个速度分量）得到严格证明，经验强烈表明完整的 3D 系统也是弱双曲的。这一观察结果表明，用于求解 QBMM 封闭的矩输运方程的数值方案必须能够处理矩中的局部 Delta 冲击。一类这样的数值方案由第 8.2 节中介绍的基于动理学的有限体积求解器组成。最后一点，我们应该提到 Chalons 等人的工作。 （2012）使用扩展高斯求积（参见第 3.3.2 节）和基于动理学的有限体积求解器来关闭方程。 (8.7) 表明具有 2N +1 矩的系统是完全双曲线的，因此不会表现出 delta 冲击。正如我们将在下面看到的（以及我们在上面的图 8.1 中看到的），基于动力学的求解器可以使用可实现的高阶空间传输方案进行设计，从而既确保自由传输项的最佳 QBMM 闭包，又保证传输矩始终是可实现的。此外，此类方法可以扩展到涉及主动（例如影响平流速度的速度或内坐标）和被动内坐标的 GPBE。2

\subsection*{8.1.3 主动内坐标与被动内坐标之间的耦合}\phantomsection\addcontentsline{toc}{subsection}{8.1.3 主动内坐标与被动内坐标之间的耦合}
在 GPBE 中，内坐标包括主动变量和被动变量。例如，考虑粒子速度 (v) 和表面积 (s) 的 1D NDF：n(t, x, v, s)，其中两个内坐标为 xi = (v, s)。 NDF 的自由传输项与等式 1 中的相同。 （8.3）。然而，我们现在有由 mk,l = 定义的双变量矩

\begin{sourceblock}
\includegraphics[page=375,pagebox=mediabox,trim={91.49bp 477.70bp 395.85bp 210.50bp},clip,max width=0.94\textwidth,max height=0.38\textheight]{Marchisio2013.pdf}
\end{sourceblock}
vksln dv ds，因此矩输运方程为

\begin{sourceblock}
\includegraphics[page=375,pagebox=mediabox,trim={204.44bp 431.88bp 67.56bp 242.21bp},clip,max width=0.94\textwidth,max height=0.38\textheight]{Marchisio2013.pdf}
\end{sourceblock}
和以前一样，方程的精确解。 (8.3) 为 n(t, x, v, s) = n(0, x −vt, v, s)，从中我们可以观察到 s 确实是被动的，因为它只是随速度 v 一起携带。或者换句话说，GPBE 中的自由传输项可以针对 s 的每个值单独求解。然而，矩输运方程是耦合的，因为它们取决于 k 和 l。尽管如此，我们可以观察到，如果我们考虑一组 k 固定但 l 自由的矩，则可以使用 QBMM 来表示未闭合矩。例如，当 N = 2 时，矩输运方程为

\begin{sourceblock}
\includegraphics[page=375,pagebox=mediabox,trim={208.96bp 231.18bp 67.56bp 368.31bp},clip,max width=0.94\textwidth,max height=0.38\textheight]{Marchisio2013.pdf}
\end{sourceblock}
并且，对于每个l，我们可以使用矩集(m0,l, m1,l, m2,l, m3,l)找到两个权重和两个节点(w1,l, w2,l, Ψ1,l和Ψ2,l)，由此可以闭合未知矩：m4,l = w1,l

\begin{sourceblock}
\includegraphics[page=375,pagebox=mediabox,trim={81.00bp 187.16bp 343.60bp 493.04bp},clip,max width=0.94\textwidth,max height=0.38\textheight]{Marchisio2013.pdf}
\end{sourceblock}
4.同样，由于方程中大小矩的被动性质，求解算法的这种分解是可能的。 （8.10）。然而，一般来说，GBPE 的右侧将包含未闭源项，需要 2D 求积近似才能收敛。例如，如果构造一个 N = 4 个节点的 2D 正交，则必须传输表 3.5 中的 12 个混合矩。自由传输项将对应于方程。 (8.10) 其中 l = 0, 1, 2, 3。然而，对于 l = 2, 3，仅传输 k = 0, 1 的四个矩（即 m0,2、m1,2、m0,3 和 m1,3），因此我们将无法找到两个权重和

\footnotetext[2]{我们将使用术语“主动”来指代速度内坐标，因为它的值直接影响空间传输。其他内坐标（例如尺寸、质量等）如果只是通过自由传输项在空间中平流传播，则被称为“被动”。然而，如果平流速度被建模为内部坐标（例如大小）的函数，则此类坐标也称为活动坐标。}
l = 2, 3 时有两个节点。相反，我们本质上有两个选择：(1) 使用第 3.2 节中描述的强力 QMOM 计算 2D、四节点求积，或 (2) 使用给定表面积的条件速度矩来构造 CQMOM 求积。如第 3.2 节所示，这种情况的 CQMOM 需要 10 个混合矩 m0,0、m0,1、m0,2、m0,3、m1,0、m1,1、m2,0、m2,1、m3,0 和 m3,1，其中自由传输方程为

\begin{sourceblock}
\includegraphics[page=376,pagebox=mediabox,trim={100.04bp 462.48bp 65.07bp 137.01bp},clip,max width=0.94\textwidth,max height=0.38\textheight]{Marchisio2013.pdf}
\end{sourceblock}
读者可以验证方程中的系统。 (8.11) 有四个未闭矩：m1,2、m1,3、m4,0 和 m4,1。矩的 CQMOM 闭包的形式为 mk,l = w1(w11vk 11 + w12vk 12)sl 1 + w2(w21vk 21 + w22vk 22)sl 2，其中权重与节点根据十个传输矩确定。 （根据定义，w11 + w12 = 1 和 w21 + w22 = 1，因此正好有 10 个未知数。）因此，例如，m4,0 的闭包为 m4,0 = w1(w11v4 11 + w12v4 12) + w2(w21v4 21 + w22v4 22)，一般来说，它与使用矩集找到的闭包不同直接 (m0,0, m1,0, m2,0, m3,0)。仅在 NDF 可分离的情况下，n(v, s) = n(v)n(s)，m4,0 的两个正交近似值才会相同。换句话说，只有当主动和被动内坐标解耦时，所有可能的正交重建才会相同。作为一个例子，让我们考虑方程 1 中 NDF 的推广。 （8.4）：

\begin{sourceblock}
\includegraphics[page=376,pagebox=mediabox,trim={94.19bp 307.64bp 65.07bp 380.02bp},clip,max width=0.94\textwidth,max height=0.38\textheight]{Marchisio2013.pdf}
\end{sourceblock}
可以用两点求积来精确表示，其中 w1(t, x) = δ(x −t + 1)、v1(t, x) = 1、s1(t, x) = 1、w2(t, x) = δ(x + t + 1)、v2(t, x) = −1 和 s2(t, x) = 2。任意 l 的矩集为

\begin{sourceblock}
\includegraphics[page=376,pagebox=mediabox,trim={174.97bp 197.20bp 65.07bp 438.79bp},clip,max width=0.94\textwidth,max height=0.38\textheight]{Marchisio2013.pdf}
\end{sourceblock}
产生权重与节点 w1,l = δ(x −t + 1)、w2,l = δ(x + t + 1)、 xi1,l = 1 和 xi2,l = 2。因此，未知矩 m4,l(t, x) = δ(x −t + 1) + (2)lδ(x + t + 1) 完全闭合。对于 CQMOM，我们从四个时刻开始

\begin{sourceblock}
\includegraphics[page=376,pagebox=mediabox,trim={178.74bp 93.26bp 65.07bp 548.62bp},clip,max width=0.94\textwidth,max height=0.38\textheight]{Marchisio2013.pdf}
\end{sourceblock}
{\small 并计算权重与节点 w1 = δ(x −t + 1)、w2 = δ(x + t + 1)、s1 = 1 和 s2 = 2，这也是精确的。为了计算条件权重与节点，我们使用六个矩\par}
\begin{sourceblock}
\includegraphics[page=377,pagebox=mediabox,trim={176.24bp 552.54bp 67.57bp 59.44bp},clip,max width=0.94\textwidth,max height=0.38\textheight]{Marchisio2013.pdf}
\end{sourceblock}
k = 1, 2, 3 的条件矩 〈vk〉α 通过求解线性系统找到

\begin{sourceblock}
\includegraphics[page=377,pagebox=mediabox,trim={120.11bp 495.20bp 67.56bp 171.97bp},clip,max width=0.94\textwidth,max height=0.38\textheight]{Marchisio2013.pdf}
\end{sourceblock}
由此得出 〈v〉1 = 〈v2〉1 = 〈v3〉1 = 1 和 〈v〉2 = −1、〈v2〉2 = 1 和 〈v3〉2 = −1，因此条件权重与节点为 w11 = w12 = w21 = w22 = 1/2, v11 = v12 = 1，且 v21 = v22 = −1。然后，该结果可用于估计输运方程中的未知矩（方程（8.11））。例如，m4,0 的 CQMOM 闭包为 m4,0 = w1v4 11 + w2v4 21 = δ(x −t + 1) + δ(x + t + 1)，这是精确的。我们应该注意到，如果 Yuan \& Fox (2011) 中描述的自适应 Wheeler 算法与条件矩一起使用，那么条件权重与节点只会返回一个节点：w11 = 1, v11 = 1 和 w21 = 1, v21 = -1。换句话说，如 Yuan \& Fox (2011) 所示，当 NDF 由加权 Dirac delta 函数（即 N 点分布）组成时，CQMOM 中使用的数值算法将找到精确的 NDF。总之，CQMOM 求积对于本例来说是精确的，因为 NDF 可以由两个多元 delta 函数表示（方程（8.12））。对于 NDF 连续的情况，CQMOM 求积将提供 NDF 的近似值，可用于估计矩输运方程中的未知矩。实际上，由于相空间源项，主动和被动内坐标将强耦合。例如，在流体-颗粒流中，颗粒速度将由于流体阻力而改变，并且阻力的大小与颗粒表面积成反比。因此，具有小 s 的粒子将具有与具有大 s 的粒子非常不同的速度 v。因为我们希望使用 GPBE 捕获的正是内坐标之间的这种相关性，所以用于处理矩输运方程中的空间传递的求解方法必须能够处理耦合混合矩组，例如方程 1 中的耦合混合矩组。 （8.11）。正如我们在本节中所示，CQMOM 是构建高度相关混合矩集求积的有效工具。此外，CQMOM 可以轻松地与基于动力学的方法 (KBFVM) 结合，处理矩输运方程中的空间传递项，同时保持可实现的矩集。在第 8.5 节中，我们表明 EQMOM 也可用于具有尺寸相关空间通量的二变量求积。

\subsection*{8.1.4 QMOM 与 DQMOM}\phantomsection\addcontentsline{toc}{subsection}{8.1.4 QMOM 与 DQMOM}
第 3 章介绍了矩求积法 (QMOM) 和直接求积矩方法 (DQMOM)，作为求解齐次 GPBE 的等效方法。事实上，DQMOM 是由 Marchisio \& Fox (2005) 推导出来的，主要目的是求解空间非齐次多元矩输运方程。与单变量情况不同，在单变量情况下，矩反演算法是为给定的一组矩唯一定义的，而多变量情况下的 QMOM 则更难以应用。使用 DQMOM，一旦选择了矩集（例如 Fox (2009b) 中的最佳矩集），并且初始化了权重与节点，（原则上）就不再需要应用矩反演算法，因为权重与节点是直接求解的。在实践中，DQMOM 在模拟过程中在空间和时间上的每个点上节点保持连续且不同的条件下工作良好。事实上，DQMOM 方程的推导（参见 Marchisio \& Fox (2005) 和 Fox et al. (2008)）要求节点的空间和时间导数是有限的（以便定义狄拉克 δ 函数的导数）并且系数矩阵 A（例如参见方程（7.118））是非奇异的。如果这些条件中的任何一个失败，那么 DQMOM 模拟就会受到影响。因此，考虑 DQMOM 可能失败的情况是有启发性的。

在查看 DQMOM 可能失败的具体示例之前，我们应该注意到使用 QMOM 直接求解矩输运方程非常稳健。在守恒方程中，矩输运方程（以及 NDF）以保守形式编写：

\begin{sourceblock}
\includegraphics[page=378,pagebox=mediabox,trim={195.48bp 426.30bp 65.07bp 248.04bp},clip,max width=0.94\textwidth,max height=0.38\textheight]{Marchisio2013.pdf}
\end{sourceblock}
其中 F 是通量函数，S 是源项。3 因此，用输运理论的术语来说，矩是守恒变量（当 S = 0 时）。相反，权重与节点是原始变量，其传输方程（即 DQMOM）不以保守形式出现。从气体动力学中众所周知，方程形式的方程。 (8.17) 可以表现出原始变量的冲击（例如，在欧拉方程中，气体速度在空间中可以是不连续的）。因此，我们可以预期，取决于方程中 F 的性质。根据方程（8.17），即使矩集 M 被明确定义，DQMOM 中的权重与节点也可能是不连续的。 CFD 社区众所周知，有限体积方法 (FVM) 非常适合求解方程 1。 （8.17）。然而，正如下面第 8.2 节所讨论的，在求解矩集时，使用可实现的 FVM 非常重要。考虑到有关 DQMOM 以及保守变量与原始变量的背景信息，我们现在可以讨论 DQMOM 可能失败的情况。

(a) 如果矩输运方程（方程（8.17））是纯双曲线的（即 NDF 或 KE 的纯平流），则当节点存在激波时，DQMOM 将失效，节点（也可能是权重）不连续。如附录 C 中所述，当特征（对应于节点）在物理空间中交叉时就会发生这种情况，因此不是唯一的。对于这种情况，使用 QMOM 和多元矩反演算法似乎是目前唯一可行的解决方案。这种情况最常发生在 KE 或 GPBE 中，其中速度是内坐标。

(b) 尽管矩输运方程中的平流项是纯双曲线的，即使没有激波且矩在空间中连续，但当用标准 FVM 求解时，DQMOM 中使用的非保守形式的输运方程可能仍然无法守恒矩，因为不可避免地存在数值误差和数值扩散。事实上，仅

\footnotetext[3]{严格来说，保守形式是指没有源项的守恒方程。当求解矩输运方程时，经常采用算子分裂，其中传递步骤用 S = 0 求解。}
权重与节点以及这些变量的线性组合（即零阶矩和一阶矩）将通过 DQMOM 方程的有限体积解来守恒。所有高阶矩（权重与节点的非线性组合）不会被数值通量函数严格守恒。这种数值误差可能会产生无法实现的矩集，特别是当在数值通量函数中使用高阶空间重建时。处理此问题的一种方法是直接根据以矩（守恒量）形式编写的可实现数值通量函数来处理物理空间中的矩平流，而不是使用 DQMOM 方程的数值通量函数。 DQMOM 的这种变体称为完全保守 DQMOM (DQMOM-FC)，在附录 D 中详细讨论。

(c) 如果矩输运方程包含空间扩散项，但扩散率非常小（因此空间梯度很难求解），则扩散的 DQMOM 校正项（详细信息请参见 Marchisio \& Fox (2005)）可能不准确。对于这种情况，DQMOM 将不会与 QMOM 完全相同，除非节点中的空间梯度在计算网格上完全解析。当使用 FVM 求解 DQMOM 方程并且数值扩散与真实扩散率具有相同数量级（或大于！）时，可能会发生这种情况。因此，经典的 DQMOM 方程更适合与在模型方程中提供显著“湍流扩散性”的湍流模型一起使用，而空间扩散非常小的情况应使用 DQMOM-FC 进行处理，如上一点所述。

(d) 如果节点的初始条件在流域内或入口处不连续，则当输运方程是纯双曲线时，DQMOM 可能会失败（参见附录 C）。例如，当域的上半部分具有一个 NDF 而下半部分具有不同的 NDF 时，通常会发生这种情况。请注意，在这种情况下，矩 M 在整个域中并不连续，但这并不妨碍我们通过求解方程 1 来获得校正解。 （8.17）使用可实现的FVM。显然，当解域中的流入NDF与初始NDF不同时，也会出现同样的问题。一般来说，流域的内部必须使用与表示流入相同的节点进行初始化，以便 DQMOM 正常工作。文献中报道的许多 DQMOM 困难都是由于不连续的初始/边界条件导致的交叉特性。

(e) 退化矩集可能会出现另一个更微妙的问题。后者是指表示所有矩所需的节点数量是有限的矩集（或空间或时间的变化）。这个问题最明显的例子是当初始条件对应于由单个加权狄拉克δ函数组成的NDF时。对于这种情况，矩 M 被明确定义（如方程（8.17）），但在只需要其中两个来重建 NDF 的意义上是退化的。例如，如果由于源项，矩在流中的某些点变得非退化，则具有自适应 Wheeler 算法的 QMOM（Yuan \& Fox，2011）将自动调整 N 的值，以增加节点的数量。相比之下，DQMOM 是用固定的 N 值制定的，因此必须采取一些“聪明”的措施来处理这种情况（例如，使 N 自适应）。

(f) 另一个更微妙的问题是在矩的空间扩散的情况下出现的所谓“X 问题”。考虑具有零平流但非零扩散的单变量 PBE：

\begin{sourceblock}
\includegraphics[page=380,pagebox=mediabox,trim={192.51bp 572.38bp 163.96bp 100.90bp},clip,max width=0.94\textwidth,max height=0.38\textheight]{Marchisio2013.pdf}
\end{sourceblock}
为简单起见，令 m0(t, x) = 1、m1(t, x) = 0、m2(0, x) = x2 和 m3(t, x) = 0，以便 x = 0 处的方差最初为零。二阶矩的精确解为 m2(t, x) = 2t + x2。相应的权重与节点为 w1(t, x) = w2(t, x) = 1/2 和 Σ1(t, x) = −Σ2(t, x) =

\begin{sourceblock}
\includegraphics[page=380,pagebox=mediabox,trim={294.01bp 528.78bp 190.22bp 159.42bp},clip,max width=0.94\textwidth,max height=0.38\textheight]{Marchisio2013.pdf}
\end{sourceblock}
2 t + x2（0 ≤ t）。因此，在 t = 0 时，节点形成一个“X”，当 0 < t 时，该“X”分为两部分。显然，对于高阶矩也可能发生同样的情况，并且必须修改 DQMOM 来处理这种情况。实际上，这种由扩散引起的简并矩数量的变化会发生在其中一个内坐标由源项生成的问题中，例如化学反应的产物。

毫无疑问，还有其他一些 DQMOM 失败但 QMOM 有效的例子。然而，作为一般规则，空间通量的扩散分量越占主导地位，节点就越有可能在空间中表现出平滑的变化，并且 DQMOM 将适用。在 DQMOM 适用的情况下，可实现矩集的问题很容易解决，因为唯一的要求是权重是非负的。尽管如此，在多变量情况下仍然有必要验证节点在物理上是可实现的，即它们必须在内坐标的相空间中取值（参见第 8.1.1 节）。使用 DQMOM（与 QMOM 一样），无法保证解将代表多元 NDF 的可实现正交。

\section*{8.2 基于动理学的有限体积方法}\phantomsection\addcontentsline{toc}{section}{8.2 基于动理学的有限体积方法}
在本节中，我们将介绍用于矩输运方程的 KBFVM，并展示如何制定它们以保证更新的矩集是可实现的。更详细的介绍可以在附录 B 中找到。为了清楚起见，我们从 PBE 开始，然后讨论 KE 产生的其他复杂性，最后讨论 GPBE。由于平流项是无法实现矩集的主要原因，因此我们将讨论限制在该项上，并重点关注减少数值扩散所需的高阶空间重建。本节中的材料大致遵循 Vikas 等人的工作。 (2011a, 2012)，读者可参阅此处和附录 B 了解更多详细信息。用于矩输运方程的 KBFVM 背后的主要思想是，每个矩集背后都存在一个相应的 NDF。因此，矩不被视为独立标量，而是被视为单个 NDF 的线性变换。考虑到这个想法，我们可以根据矩或 NDF 自由地编写相同的方案。然后，为了保证矩集是可实现的，我们对 NDF 强制执行一个可实现性条件，该条件规定“数值”NDF 对于方案中使用的参数（例如时间步长和网格间距）必须是非负的。对于任意的NDF，只有一阶空间重构（即NDF在每个网格单元中是恒定的）才能保证矩集的可实现性。然而，我们将看到，通过 QBMM，NDF 具有加权狄拉克δ函数表示，使得可以对权重使用高阶空间重构，同时仍然保证数值 NDF 是非负的，从而保证矩集是可实现的。下面，我们使用涉及一维空间平流或自由传输的简单案例演示 KBFVM 背后的基本思想。更多详细信息请参见附录 B。

\subsection*{8.2.1 应用于PBE}\phantomsection\addcontentsline{toc}{subsection}{8.2.1 应用于PBE}
为了简单起见，这里只讨论一维情况，我们用 n(t, x, Ψ) 表示单变量 NDF，其中 Ψ 是被动内坐标。仅包含平流项的 PBE 为

\begin{sourceblock}
\includegraphics[page=381,pagebox=mediabox,trim={215.53bp 510.84bp 67.56bp 163.51bp},clip,max width=0.94\textwidth,max height=0.38\textheight]{Marchisio2013.pdf}
\end{sourceblock}
其中 u(t, x) 是已知速度。从方程找到的矩输运方程。 (8.18) 是

\begin{sourceblock}
\includegraphics[page=381,pagebox=mediabox,trim={209.60bp 460.36bp 67.56bp 214.00bp},clip,max width=0.94\textwidth,max height=0.38\textheight]{Marchisio2013.pdf}
\end{sourceblock}
将 QBMM4 中使用的矩集表示为 M = [m0 m1 m2。。。 m2N−1]T 且矩通量为 F(u, M) = uM，矩输运方程可写为

\begin{sourceblock}
\includegraphics[page=381,pagebox=mediabox,trim={65.03bp 363.07bp 67.56bp 278.28bp},clip,max width=0.94\textwidth,max height=0.38\textheight]{Marchisio2013.pdf}
\end{sourceblock}
Kn dΣ，其中 K = [1 Σ Σ2 · · · Σ2N−1]T。从等式开始（8.20），我们现在开发一个可实现的高阶方案，在给定时间 t 处的可实现矩集的情况下更新时间 t + Δt 处的矩集。在间距为 Δx 的均匀网格上，我们将时间 t 处的体积平均矩定义为

\begin{sourceblock}
\includegraphics[page=381,pagebox=mediabox,trim={186.13bp 290.28bp 67.56bp 375.40bp},clip,max width=0.94\textwidth,max height=0.38\textheight]{Marchisio2013.pdf}
\end{sourceblock}
其中 xi 是像元中心，t = m Δt 是当前时间。如果使用单阶段欧拉显式时间积分方案，则更新后的矩集可以写为

\begin{sourceblock}
\includegraphics[page=381,pagebox=mediabox,trim={130.27bp 208.10bp 67.56bp 441.68bp},clip,max width=0.94\textwidth,max height=0.38\textheight]{Marchisio2013.pdf}
\end{sourceblock}
其中 λ = Δt/Δx，G 是数值通量函数，定义为

\begin{sourceblock}
\includegraphics[page=381,pagebox=mediabox,trim={148.43bp 158.15bp 67.56bp 514.44bp},clip,max width=0.94\textwidth,max height=0.38\textheight]{Marchisio2013.pdf}
\end{sourceblock}
2(ul + |ul|) 和 u− r = 1 l = 1 且 u+ 2(ur −|ur|)。下标i−1/2(i+1/2)表示网格单元i−1和i(i和i+1)之间的界面。下标l(r)表示从左(右)接近界面的限制。例如，点 xi−1/2,l 位于网格单元 i -1 中但靠近与网格单元 i 的界面，点 xi−1/2,r 位于同一界面附近但位于网格单元 i 中。同样，um i−1/2,l 表示网格中时间 t 的速度

\footnotetext[4]{为简单起见，我们使用 QMOM 来说明数值方法。第 8.3 节提供了使用 EQMOM 的示例。}
单元 i -1，但位于网格单元 i -1 和 i 之间的界面。读者应注意等式中。 (8.23) 数值通量函数要求我们估计位于单元界面左侧和右侧的 NDF nl 和 nr。这些估计必须根据所有网格单元中的已知矩集（即 Mm i ）构建。在引入数值 NDF hi(xi) 时，方程中设置的更新矩： (8.22) 可以写成 Mm+1

\begin{sourceblock}
\includegraphics[page=382,pagebox=mediabox,trim={118.42bp 564.57bp 343.60bp 111.97bp},clip,max width=0.94\textwidth,max height=0.38\textheight]{Marchisio2013.pdf}
\end{sourceblock}
基 d δ，其中

\begin{sourceblock}
\includegraphics[page=382,pagebox=mediabox,trim={101.98bp 525.02bp 65.07bp 142.56bp},clip,max width=0.94\textwidth,max height=0.38\textheight]{Marchisio2013.pdf}
\end{sourceblock}
读者应该注意到，该表达式中的 xi 依赖性来自重构的 NDF：纳米

\begin{sourceblock}
\includegraphics[page=382,pagebox=mediabox,trim={152.24bp 492.66bp 146.71bp 190.91bp},clip,max width=0.94\textwidth,max height=0.38\textheight]{Marchisio2013.pdf}
\end{sourceblock}
我。假设时间级别 m 处的矩对应于非负密度函数（即 nm i 对于所有 i 都是非负的），方程右侧的前三项(8.24) 是非负的，最后两项是非正的。为了使时间级别 m + 1 的时刻对应于非负密度函数，hi(xi) 对于所有 xi 都必须是非负的。使用 NDF 的一阶空间重建，细胞界面处的 NDF 与细胞中心处的 NDF 相同， nm

\begin{sourceblock}
\includegraphics[page=382,pagebox=mediabox,trim={89.69bp 398.73bp 65.07bp 262.64bp},clip,max width=0.94\textwidth,max height=0.38\textheight]{Marchisio2013.pdf}
\end{sourceblock}
因此更新后的矩集是

\begin{sourceblock}
\includegraphics[page=382,pagebox=mediabox,trim={97.89bp 359.23bp 65.07bp 316.84bp},clip,max width=0.94\textwidth,max height=0.38\textheight]{Marchisio2013.pdf}
\end{sourceblock}
hi(xi) 的非负性在以下可实现条件下得到保证：

\begin{sourceblock}
\includegraphics[page=382,pagebox=mediabox,trim={98.00bp 306.03bp 65.07bp 360.30bp},clip,max width=0.94\textwidth,max height=0.38\textheight]{Marchisio2013.pdf}
\end{sourceblock}
读者应该注意到，即使对速度使用高阶空间重构，时间步长 Δt 上的这种可实现性条件也适用。还值得注意的是，方程。即使速度不连续（例如冲击时），方程（8.26）也适用。对于方程中 NDF 的二阶和一般情况下的任何高阶空间重建。 (8.24)，不能保证 hi(xi) 的非负性。然而，如果将 QBMM 与下面讨论的特殊重构一起使用，则可以保证 hi(xi) 的非负性。使用具有 N 节点求积的 QBMM，5 NDF 为

\begin{sourceblock}
\includegraphics[page=382,pagebox=mediabox,trim={173.85bp 178.48bp 65.07bp 488.45bp},clip,max width=0.94\textwidth,max height=0.38\textheight]{Marchisio2013.pdf}
\end{sourceblock}
其中 wα 是权重，Σα 是节点。对于这个重建的 NDF，由 Vikas 等人提出以下定理。 (2011a) 适用。

定理 8.1 令 N, p ∈N 且 α ∈\{1, 2, .。。，N\}。对于体积平均矩集 Mm

\begin{sourceblock}
\includegraphics[page=382,pagebox=mediabox,trim={291.93bp 116.85bp 104.33bp 567.07bp},clip,max width=0.94\textwidth,max height=0.38\textheight]{Marchisio2013.pdf}
\end{sourceblock}
i，令重建的权重值满足 wm i−1/2,α,r ≥0 ∀i, α。如果设计使用单级欧拉显式时间积分方案的有限体积方案，该方案对权重使用 p 阶空间重构，对节点使用一阶空间重构，则方程 1 中数值 NDF 的非负性。 (8.24) 总是可以在时间步长 (Δt ∈R+) 的显式约束下得到保证。

\footnotetext[5]{QMOM 和EQMOM 都可以通过使用正交节点写成这种离散形式。}
证明如下。使用方程式（8.24），无论使用哪种方案，数值密度函数都可以写为

\begin{sourceblock}
\includegraphics[page=383,pagebox=mediabox,trim={149.80bp 590.12bp 67.56bp 93.48bp},clip,max width=0.94\textwidth,max height=0.38\textheight]{Marchisio2013.pdf}
\end{sourceblock}
在上面的表达式中，方程中的非负项。 (8.24) 已归入 ε+ 下。对于 N 节点求积，使用等式： (8.28)，hm i 的表达式变为

\begin{sourceblock}
\includegraphics[page=383,pagebox=mediabox,trim={130.11bp 494.88bp 67.56bp 149.97bp},clip,max width=0.94\textwidth,max height=0.38\textheight]{Marchisio2013.pdf}
\end{sourceblock}
如果节点采用一阶空间重建，则节点的界面值将与“体积平均”值相同：6 Σm

\begin{sourceblock}
\includegraphics[page=383,pagebox=mediabox,trim={326.06bp 456.48bp 92.13bp 227.30bp},clip,max width=0.94\textwidth,max height=0.38\textheight]{Marchisio2013.pdf}
\end{sourceblock}
i−1/2,α,r。将其代入方程式。 (8.30)，数值 NDF 变为

\begin{sourceblock}
\includegraphics[page=383,pagebox=mediabox,trim={95.47bp 405.68bp 67.56bp 261.27bp},clip,max width=0.94\textwidth,max height=0.38\textheight]{Marchisio2013.pdf}
\end{sourceblock}
对于 Δt ∈R+ 满足可实现性条件，

\begin{sourceblock}
\includegraphics[page=383,pagebox=mediabox,trim={152.86bp 346.63bp 67.56bp 316.71bp},clip,max width=0.94\textwidth,max height=0.38\textheight]{Marchisio2013.pdf}
\end{sourceblock}
hi 对于所有 Ψ 都是非负的。证明到此结束。当权重采用高阶空间重构时，更新后的矩集为

\begin{sourceblock}
\includegraphics[page=383,pagebox=mediabox,trim={92.81bp 210.90bp 67.56bp 403.55bp},clip,max width=0.94\textwidth,max height=0.38\textheight]{Marchisio2013.pdf}
\end{sourceblock}
其中 公里
6 将权重与节点表示为“体积平均值”在技术上是不正确的，因为只有矩是体积平均值。根据体积平均矩的定义，我们可以写出每个矩

\begin{sourceblock}
\includegraphics[page=383,pagebox=mediabox,trim={156.40bp 149.69bp 158.93bp 520.45bp},clip,max width=0.94\textwidth,max height=0.38\textheight]{Marchisio2013.pdf}
\end{sourceblock}
由于对节点的非线性依赖，该表达式的右侧一般不能用真实的体积平均重量和节点代替：

\begin{sourceblock}
\includegraphics[page=383,pagebox=mediabox,trim={78.67bp 95.41bp 80.00bp 575.82bp},clip,max width=0.94\textwidth,max height=0.38\textheight]{Marchisio2013.pdf}
\end{sourceblock}
{\small 然而，对于特殊的重建，其中 Ψα(m Δt, x) = Ψm\par}
{\small i,α 在每个网格单元中是恒定的，将权重与节点称为体积平均值是正确的。然而，矩和节点之间的非线性依赖性使得在保留体积平均矩的值的同时对节点使用高阶空间重建具有挑战性。\par}
\begin{sourceblock}
\includegraphics[page=384,pagebox=mediabox,trim={100.91bp 623.57bp 353.10bp 60.00bp},clip,max width=0.94\textwidth,max height=0.38\textheight]{Marchisio2013.pdf}
\end{sourceblock}
其中 wm i,α 是从体积平均矩集 Mm i 中找到的。因此，单元 i 中的重构矩为

\begin{sourceblock}
\includegraphics[page=384,pagebox=mediabox,trim={81.36bp 580.32bp 65.07bp 86.62bp},clip,max width=0.94\textwidth,max height=0.38\textheight]{Marchisio2013.pdf}
\end{sourceblock}
请注意，由于节点在区间 xi −Δx/2 ≤x ≤xi + Δx/2 内恒定，因此重构的矩集自动满足式（1）。 (8.21) 对于任何 ai,α 值。实际上，ai,α 是使用具有流量限制器的相邻单元中的 wm i,α 值来估计的（Vikas 等人，2011a）。方程中公式的推广。将 (8.33) 转换为 2D 和 3D 空间域以及非结构化网格非常简单，如 Vikas 等人所述。 （2011a，2012）。为了提高时间精度，可以使用二阶、两级、强保稳龙格库塔（RK2SSP）方案（Vikas et al., 2011a）。更新后的矩集可从以下位置找到

\begin{sourceblock}
\includegraphics[page=384,pagebox=mediabox,trim={201.68bp 448.77bp 65.07bp 225.40bp},clip,max width=0.94\textwidth,max height=0.38\textheight]{Marchisio2013.pdf}
\end{sourceblock}
其中第一阶段矩集定义为

\begin{sourceblock}
\includegraphics[page=384,pagebox=mediabox,trim={67.52bp 401.40bp 65.07bp 265.54bp},clip,max width=0.94\textwidth,max height=0.38\textheight]{Marchisio2013.pdf}
\end{sourceblock}
(8.36) 和第二阶段矩由

\begin{sourceblock}
\includegraphics[page=384,pagebox=mediabox,trim={67.88bp 333.83bp 65.07bp 323.66bp},clip,max width=0.94\textwidth,max height=0.38\textheight]{Marchisio2013.pdf}
\end{sourceblock}
第二阶段中出现的权重与节点是根据 M(1) i 计算的。读者可以注意到，每个阶段使用的公式与方程式完全相同。 （8.33），因此是可实现的。令人惊讶的是，标准的 Runge-Kutta 方案无法实现（Vikas 等人，2011a），因此在设计高阶时间步进方案时必须小心。方程的扩展。将（8.33）（或方程（8.35））转换为多元 PBE n(t, x, xi) 相当简单（参见附录 B 的 B.3 节）。首先，需要一个多元求积：

\begin{sourceblock}
\includegraphics[page=384,pagebox=mediabox,trim={173.46bp 205.59bp 65.07bp 461.35bp},clip,max width=0.94\textwidth,max height=0.38\textheight]{Marchisio2013.pdf}
\end{sourceblock}
体积平均矩集 Mm i 现在包含找到 wm 所需的多元矩

\begin{sourceblock}
\includegraphics[page=384,pagebox=mediabox,trim={105.29bp 175.84bp 345.27bp 507.73bp},clip,max width=0.94\textwidth,max height=0.38\textheight]{Marchisio2013.pdf}
\end{sourceblock}
我，α。和以前一样，仅对权重采用高阶空间重建，因此等式：除了 Km i,α 的定义之外（8.33）保持不变，Km i,α 现在包括用于构造多元求积的所有混合矩。如第 3 章所述，传输矩集的一个方便选择是最优矩集。如果使用 CQMOM 来构造求积，则权重与节点将取决于调节中使用的变量的顺序（即排列）。为了对权重使用高阶空间重建，必须在所有网格单元中使用相同的排列（有关详细信息，请参见 Yuan \& Fox (2011) 和附录 B 的 B.3 节）。此外，当应用于等式时，CQMOM 的每个排列。 （8.33），将产生不同的更新矩集（Mm+1
{\small i )p, 其中 p ∈1, .。。 、P 和 P 是排列数（例如，对于 2D 相空间，P = 2；对于 3D 相空间，P = 6）。最终更新的矩集定义为\par}
\begin{sourceblock}
\includegraphics[page=385,pagebox=mediabox,trim={196.95bp 608.32bp 67.56bp 58.62bp},clip,max width=0.94\textwidth,max height=0.38\textheight]{Marchisio2013.pdf}
\end{sourceblock}
想要了解更多详细信息的读者请参阅 Yuan \& Fox (2011) 和附录 B 的 B.3 节。

\subsection*{8.2.2 KE 的应用}\phantomsection\addcontentsline{toc}{subsection}{8.2.2 KE 的应用}
我们再次仅考虑 1D KE，并用 n(t, x, v) 表示单变量 NDF，其中 v 是速度。仅包含自由传输项的 KE 为

\begin{sourceblock}
\includegraphics[page=385,pagebox=mediabox,trim={215.81bp 488.62bp 67.56bp 185.73bp},clip,max width=0.94\textwidth,max height=0.38\textheight]{Marchisio2013.pdf}
\end{sourceblock}
矩输运方程为

\begin{sourceblock}
\includegraphics[page=385,pagebox=mediabox,trim={208.13bp 442.42bp 67.56bp 231.93bp},clip,max width=0.94\textwidth,max height=0.38\textheight]{Marchisio2013.pdf}
\end{sourceblock}
将矩通量表示为 F(n)，其中 n 是根据传递矩集 M 重建的 NDF，矩输运方程可写为

\begin{sourceblock}
\includegraphics[page=385,pagebox=mediabox,trim={209.45bp 384.27bp 67.56bp 290.07bp},clip,max width=0.94\textwidth,max height=0.38\textheight]{Marchisio2013.pdf}
\end{sourceblock}
为了方便起见，将矩通量写成分裂形式，即 F(n) = Q−(n) + Q+(n)，其中

\begin{sourceblock}
\includegraphics[page=385,pagebox=mediabox,trim={152.73bp 330.20bp 67.56bp 336.00bp},clip,max width=0.94\textwidth,max height=0.38\textheight]{Marchisio2013.pdf}
\end{sourceblock}
且 K = [1 v v2 · · · v2N−1]T。最后，为了适应 NDF 的任意重建（例如高斯 EQMOM），我们定义了两个新的半矩集：

\begin{sourceblock}
\includegraphics[page=385,pagebox=mediabox,trim={167.11bp 269.72bp 67.56bp 396.48bp},clip,max width=0.94\textwidth,max height=0.38\textheight]{Marchisio2013.pdf}
\end{sourceblock}
\begin{sourceblock}
\includegraphics[page=385,pagebox=mediabox,trim={65.02bp 226.33bp 67.55bp 440.61bp},clip,max width=0.94\textwidth,max height=0.38\textheight]{Marchisio2013.pdf}
\end{sourceblock}
α.8 请注意，在这些求积中，M 可以取任何正值，因此，如果 n 被重构为 N 个狄拉克 delta 函数，则应使用自适应求积算法将未使用的权重设置为零（详细信息请参见 Yuan \& Fox (2011)）。从等式开始根据方程（8.42），我们可以开发一个可实现的高阶方案来更新时间 t + Δt 处的传输矩集（给定时间 t 处的可实现矩集）。如果使用单阶段欧拉显式时间积分方案，则更新后的矩集可以写为

\begin{sourceblock}
\includegraphics[page=385,pagebox=mediabox,trim={97.70bp 123.32bp 67.56bp 553.10bp},clip,max width=0.94\textwidth,max height=0.38\textheight]{Marchisio2013.pdf}
\end{sourceblock}
7 一般而言，传输矩集 M 可以包含比半矩集 M− 和 M+ 更少的矩。实际上，选择 M− 和 M+ 的长度，使得从半矩集构造的高斯求积包括出现在 M 的传输方程中的最高阶矩，例如米2N。这意味着 M = N + 1。

\footnotetext[8]{在第 8.4 节中，采用直接使用 EQMOM 和矩的替代定义来确定权重与节点。在这两种公式中，总是可以将节点集分为正速度和负速度。}
其中数值通量函数由 G(nl, nr) = Q−(nr) + Q+(nl) 定义。然后，使用等式中的定义。 （8.43），我们可以写

\begin{sourceblock}
\includegraphics[page=386,pagebox=mediabox,trim={139.44bp 574.76bp 65.07bp 100.65bp},clip,max width=0.94\textwidth,max height=0.38\textheight]{Marchisio2013.pdf}
\end{sourceblock}
其中 K− α,r (K+ α,r (v+ α,l) 的定义与 K 相同，但使用节点 v− α,l 进行评估）。同样，我们有9

\begin{sourceblock}
\includegraphics[page=386,pagebox=mediabox,trim={173.22bp 511.12bp 65.07bp 161.89bp},clip,max width=0.94\textwidth,max height=0.38\textheight]{Marchisio2013.pdf}
\end{sourceblock}
因此方程。 (8.45) 可以重写为

\begin{sourceblock}
\includegraphics[page=386,pagebox=mediabox,trim={80.51bp 411.02bp 65.07bp 211.50bp},clip,max width=0.94\textwidth,max height=0.38\textheight]{Marchisio2013.pdf}
\end{sourceblock}
等式右侧的求和。 （8.48）可以写成以下两个数值NDF：

\begin{sourceblock}
\includegraphics[page=386,pagebox=mediabox,trim={129.54bp 340.18bp 65.07bp 326.90bp},clip,max width=0.94\textwidth,max height=0.38\textheight]{Marchisio2013.pdf}
\end{sourceblock}
根据可实现性条件，这些 NDF 对于 v 的所有值都必须是非负的。如前所述，方程中的数值 NDF 为：仅当我们使用 QBMM 并对节点施加以下可实现性条件时，方程 (8.49) 才为非负：

\begin{sourceblock}
\includegraphics[page=386,pagebox=mediabox,trim={137.65bp 263.78bp 65.07bp 419.82bp},clip,max width=0.94\textwidth,max height=0.38\textheight]{Marchisio2013.pdf}
\end{sourceblock}
然后我们有

\begin{sourceblock}
\includegraphics[page=386,pagebox=mediabox,trim={89.00bp 173.62bp 65.07bp 460.07bp},clip,max width=0.94\textwidth,max height=0.38\textheight]{Marchisio2013.pdf}
\end{sourceblock}
因此 Δt 的可实现条件为

\begin{sourceblock}
\includegraphics[page=386,pagebox=mediabox,trim={96.07bp 116.42bp 65.07bp 544.47bp},clip,max width=0.94\textwidth,max height=0.38\textheight]{Marchisio2013.pdf}
\end{sourceblock}
该结果概括了 Vikas 等人的定理 1。 (2011a) 到 M− 和 M+ 定义中使用的 NDF 的任意重构。请注意，如果对权重使用一阶空间重建，则 wm+

\begin{sourceblock}
\includegraphics[page=386,pagebox=mediabox,trim={217.57bp 71.32bp 149.45bp 611.91bp},clip,max width=0.94\textwidth,max height=0.38\textheight]{Marchisio2013.pdf}
\end{sourceblock}
{\small α,i−1/2,r，则方程： (8.52)\par}
\footnotetext[9]{这些方程是恒等式，因为 M ≥N，因此达到 2M -1 阶的半矩设置是精确的。}
简化为基于速度节点最大幅度的 CFL 条件。上面讨论中引入的一个微妙但非常重要的点是，即使重构的 NDF 可能是连续的，我们仍然使用半矩集 M− 和 M+ 的求积来计算等式 1 中出现的权重与节点。 （8.48）。此步骤确保方程中的可实现性条件。对于有限的 Δt 值，方程 (8.52) 得到了很好的定义。综上所述，从传输矩集 Mm 开始的整体时间步长过程如下：

\begin{sourceblock}
\includegraphics[page=387,pagebox=mediabox,trim={125.89bp 529.74bp 128.43bp 151.78bp},clip,max width=0.94\textwidth,max height=0.38\textheight]{Marchisio2013.pdf}
\end{sourceblock}
(i) 使用传输矩集重建 NDF（例如高斯 EQMOM）。

(ii) 使用重建的 NDF 来分析计算半矩集。

(iii) 使用半矩集构造两个具有 M 个节点的高斯求积。

(iv) 使用根据半矩集计算的权重与节点推进传输矩集。

步骤(iv)中设置的更新矩是从10找到的

\begin{sourceblock}
\includegraphics[page=387,pagebox=mediabox,trim={97.36bp 343.83bp 67.56bp 289.85bp},clip,max width=0.94\textwidth,max height=0.38\textheight]{Marchisio2013.pdf}
\end{sourceblock}
如果方程中的可实现性条件。 (8.52) 满足，Mm+1 i 由式(8.52)找到。 (8.53) 保证可以实现。使用上述 PBE 方法（方程（8.35）），可以直接扩展方程： （8.53）到二阶时间步进。在多个空间维度上对多个速度分量 v 的扩展有点复杂。如 Yuan \& Fox (2011) 和附录 B 的 B.3 节中所述，当使用 CQMOM 构造多元求积时，必须以一致的方式使用所有排列，以获得正确的动能通量。然而，对于 CQMOM 的每个排列，方程中的基本时间步进公式。 (8.53) 用于通过修改 Km α,i 的定义来更新传输矩集，以将多元矩包含在最佳矩集中。对多维自由运输的更多细节感兴趣的读者可以查阅 Yuan \& Fox (2011) 和附录 B 的 B.3 节。

\subsection*{8.2.3 GPBE 的应用}\phantomsection\addcontentsline{toc}{subsection}{8.2.3 GPBE 的应用}
KBFVM 在 GPBE 中的应用与上面讨论的 KE 情况非常相似，但现在 NDF 包含具有已知（或计算的）平流的活动标量 n(t, x, xi)

10 正如 Yuan \& Fox (2011) 和附录 B 的 B.3 节中所讨论的，在多变量情况下(8.53) 不等于

\begin{sourceblock}
\includegraphics[page=387,pagebox=mediabox,trim={73.43bp 85.80bp 75.96bp 587.78bp},clip,max width=0.94\textwidth,max height=0.38\textheight]{Marchisio2013.pdf}
\end{sourceblock}
{\small 因为CQMOM并不能控制Mm中出现的所有最优时刻\par}
{\small 我。或者，换句话说，只有方程式中给出的形式。 （8.53）保证可以实现。因此，最好保留式（1）的右侧。 （8.53）在权重与节点方面，即使对于单变量情况也是如此。\par}
速度 u(t, x, Ψ)，或者它是速度 v 和被动标量 Ψ 的二元变量。我们从主动标量的情况开始，其 GPBE 的形式为等式： （8.18）。单步欧拉显式时间积分方案与式（1）相同。 (8.22)，但是数值通量函数现在是

\begin{sourceblock}
\includegraphics[page=388,pagebox=mediabox,trim={152.57bp 573.88bp 65.07bp 98.03bp},clip,max width=0.94\textwidth,max height=0.38\textheight]{Marchisio2013.pdf}
\end{sourceblock}
l (xi) = 1 2(ul(多项式) + |ul(多项式)|) 且 u− r (多项式) = 1 2(ur(多项式) −|ur(多项式)|)。请注意，平流速度是针对每个 xi 值单独定义的，但其他符号与方程 1 中的相同。 （8.23）。对权重和平流速度 u(m Δt, x, xi) 的 x 相关性使用高阶空间重建，更新后的矩集为 u+

\begin{sourceblock}
\includegraphics[page=388,pagebox=mediabox,trim={107.16bp 461.41bp 65.07bp 188.91bp},clip,max width=0.94\textwidth,max height=0.38\textheight]{Marchisio2013.pdf}
\end{sourceblock}
i,α 的定义如式 (1) 所示。 （8.33）。读者会注意到，这个表达式为 Mm+1，其中 Km

\begin{sourceblock}
\includegraphics[page=388,pagebox=mediabox,trim={67.52bp 380.61bp 65.06bp 256.10bp},clip,max width=0.94\textwidth,max height=0.38\textheight]{Marchisio2013.pdf}
\end{sourceblock}
这类似于等式。 (8.32) 除了现在平流速度取决于节点之外。接下来我们考虑 n(t, x, v, s) 的一维域中自由传输的情况，其中 s 是被动内部坐标（例如表面积）。附录 B 中讨论了 GPBE 的更一般情况。方程 1 中给出了具有自由传递的矩输运方程。 （8.9）。传递的二元矩集将再次用 M 表示，紧凑的矩输运方程用 Eq.1 表示。 （8.42） 矩通量如方程（8.42）所示。 （8.43），但 n 表示重建的二元 NDF。矩通量的分裂形式再次由方程给出。 (8.43) 但现在必须使用二元内坐标来定义向量 K：K = [1 v s v2 vs s2 · · · vksl · · · ]T，其中 K 的长度由最佳矩集确定。半矩集使用速度定义：

\begin{sourceblock}
\includegraphics[page=388,pagebox=mediabox,trim={139.65bp 211.94bp 65.07bp 454.27bp},clip,max width=0.94\textwidth,max height=0.38\textheight]{Marchisio2013.pdf}
\end{sourceblock}
其中，根据定义，M = M−+ M+。与KE一样，半矩集可以与QBMM一起使用来构造两组权重与节点：11

\begin{sourceblock}
\includegraphics[page=388,pagebox=mediabox,trim={123.79bp 141.11bp 65.07bp 517.87bp},clip,max width=0.94\textwidth,max height=0.38\textheight]{Marchisio2013.pdf}
\end{sourceblock}
α 用于定义空间通量。单阶段欧拉显式时间积分方案的进行与等式 1 完全相同。 （8.45）和方程的结果。 （8.48）。然而，在这种情况下，两个二元数值 NDF 是

\begin{sourceblock}
\includegraphics[page=388,pagebox=mediabox,trim={100.45bp 80.73bp 65.07bp 586.30bp},clip,max width=0.94\textwidth,max height=0.38\textheight]{Marchisio2013.pdf}
\end{sourceblock}
\footnotetext[11]{正如第 8.4 节中所做的那样，也可以直接从 EQMOM 求积获得权重与节点。}
并且对于 v 和 s 的所有值来说，它们都必须是非负的。使用 NDF 的 QBMM 重构，并在方程 (1) 中施加可实现性条件。 (8.50) 以及可实现性条件 sm−

\begin{sourceblock}
\includegraphics[page=389,pagebox=mediabox,trim={127.97bp 481.84bp 67.56bp 83.72bp},clip,max width=0.94\textwidth,max height=0.38\textheight]{Marchisio2013.pdf}
\end{sourceblock}
因此，Δt 的可实现条件与等式 1 中的相同。 （8.52）。综上所述，从传输矩集 Mm 开始的四阶段时间步长过程如下：

\begin{sourceblock}
\includegraphics[page=389,pagebox=mediabox,trim={108.17bp 422.73bp 110.70bp 258.79bp},clip,max width=0.94\textwidth,max height=0.38\textheight]{Marchisio2013.pdf}
\end{sourceblock}
其中更新后的矩集 Mm+1 是从等式中找到的。 (8.53) 使用上面给出的 K 的定义。对高阶时间步长和多元主动和/或被动内部坐标的扩展的评论出现在方程之后。 (8.53) 也适用于 GPBE。

\section*{8.3 非均匀PBE}\phantomsection\addcontentsline{toc}{section}{8.3 非均匀PBE}
在本节中，我们考虑单变量 NDF n(t, x, Ψ) 的平流和扩散，其中 Ψ ε R+ 是被动标量（即速度 u(t, x) 不依赖于 Ψ）。例如，ψ 可以表示斯托克斯数非常小的液体中细颗粒的颗粒质量。假设扩散系数 Γ(xi) 是 xi 的函数。例如，使用斯托克斯-爱因斯坦公式计算方程中报告的扩散率。 (5.116) 第 187 页，Γ(ϋ) = Γ0/ϋ，其中 ϋ 是粒子质量。此案例的 PBE 内容为

\begin{sourceblock}
\includegraphics[page=389,pagebox=mediabox,trim={191.69bp 220.47bp 67.56bp 453.87bp},clip,max width=0.94\textwidth,max height=0.38\textheight]{Marchisio2013.pdf}
\end{sourceblock}
其中速度 u 和初始 NDF n(0, x, xi) 已知。对于本节中考虑的示例，我们不会在方程式中添加项。 （8.61）用于相空间过程，例如生长和聚集。当采用算子分裂时，这些项与空间传输项分开处理，如第 7 章所述。然而，为了有实际的例子，我们将假设初始 NDF 足够复杂，足以代表具有聚合和破坏的系统。

\subsection*{8.3.1 矩输运方程}\phantomsection\addcontentsline{toc}{subsection}{8.3.1 矩输运方程}
\begin{sourceblock}
\includegraphics[page=389,pagebox=mediabox,trim={307.92bp 104.44bp 179.42bp 583.75bp},clip,max width=0.94\textwidth,max height=0.38\textheight]{Marchisio2013.pdf}
\end{sourceblock}
Ψkn(t, x, Ψ) d Υ，从式（1）中找到的矩输运方程。 (8.61) 是将 NDF 的整数矩定义为 mk(t, x) =

\begin{sourceblock}
\includegraphics[page=389,pagebox=mediabox,trim={166.47bp 50.16bp 67.56bp 616.32bp},clip,max width=0.94\textwidth,max height=0.38\textheight]{Marchisio2013.pdf}
\end{sourceblock}
Σkn(0, x, Σ)dΣ。读者可以注意到，由于扩散系数与 xi 相关，方程： (8.62) 不是封闭的。如果 Γ(xi) 不依赖于 t 或 x（例如由于温度的空间梯度），我们可以重写方程 1 中的扩散项。 (8.62) 使用 必须在初始条件下求解 mk(0, x) =

\begin{sourceblock}
\includegraphics[page=390,pagebox=mediabox,trim={163.85bp 540.57bp 155.61bp 111.85bp},clip,max width=0.94\textwidth,max height=0.38\textheight]{Marchisio2013.pdf}
\end{sourceblock}
Ψk Г(Ψ)n d Ψ = Г0mk−1，本质上是第 8.3.4 节数值示例中将使用的形式 12。总之，我们将在第 8.3.4 节的数值示例中考虑矩输运方程的两种变体。第一个例子将使用闭矩系统，其中扩散率不依赖于 Ψ：此外，采用斯托克斯-爱因斯坦公式，我们可以写出

\begin{sourceblock}
\includegraphics[page=390,pagebox=mediabox,trim={190.92bp 446.77bp 65.07bp 227.59bp},clip,max width=0.94\textwidth,max height=0.38\textheight]{Marchisio2013.pdf}
\end{sourceblock}
u 和 Γ 是已知常数。第二个例子将考虑非闭合矩系统

\begin{sourceblock}
\includegraphics[page=390,pagebox=mediabox,trim={184.97bp 386.62bp 65.07bp 287.73bp},clip,max width=0.94\textwidth,max height=0.38\textheight]{Marchisio2013.pdf}
\end{sourceblock}
其中 u 和 Γ0 是已知常数。因为它是封闭的，方程。 (8.63) 可以使用标准 FVM 求解。因此，在第 8.3.2 节中，我们使用等式： （8.63）来说明高阶方案可能由于平流项和扩散项而产生无法实现的矩。在第 8.3.3 节中，我们为等式 2 实现了可实现的高阶方案。 (8.63) 和方程。 （8.64）使用 gamma EQMOM 从时刻重建 NDF。

\subsection*{8.3.2 矩的标准有限体积格式}\phantomsection\addcontentsline{toc}{subsection}{8.3.2 矩的标准有限体积格式}
在本节中，我们开发一个高阶方案来求解方程。 （8.63）在间距为 Δx 的均匀网格上。用 xi 表示网格点，体积平均矩为

\begin{sourceblock}
\includegraphics[page=390,pagebox=mediabox,trim={185.78bp 220.84bp 183.34bp 444.79bp},clip,max width=0.94\textwidth,max height=0.38\textheight]{Marchisio2013.pdf}
\end{sourceblock}
其中 xi−1/2 和 xi+1/2 分别表示以 xi 为中心的网格单元的左界面和右界面的位置。方程的体积平均。 (8.63) 导致

\begin{sourceblock}
\includegraphics[page=390,pagebox=mediabox,trim={104.19bp 151.39bp 65.07bp 514.89bp},clip,max width=0.94\textwidth,max height=0.38\textheight]{Marchisio2013.pdf}
\end{sourceblock}
其中 mk,i±1/2(t) 是界面处的重构矩，而 Г ∂mk/∂x|i+1/2 和 Г ∂mk/∂x|i−1/2 是重构梯度。在标准方案中，重建值是从相邻网格单元中的体积平均矩找到的。下面我们考虑具有一阶、二阶和四阶精度的方案。

\footnotetext[12]{由于斯托克斯-爱因斯坦公式对于有限尺寸粒子有效，因此 NDF 仅当 Ψ0 ≤ Ψ < ∞ 时才应为非零。为了方便起见，我们将允许 0 ≤ xi < ∞ 并将 ψ0 添加到斯托克斯-爱因斯坦公式的分母上，以确保扩散率在 ψ = 0 时保持有限。}
\paragraph*{一阶方案}
对方程中的平流项使用一阶迎风重建。 （8.65），我们有

\begin{sourceblock}
\includegraphics[page=391,pagebox=mediabox,trim={129.46bp 570.67bp 133.18bp 92.38bp},clip,max width=0.94\textwidth,max height=0.38\textheight]{Marchisio2013.pdf}
\end{sourceblock}
其中

\begin{sourceblock}
\includegraphics[page=391,pagebox=mediabox,trim={149.35bp 527.76bp 67.56bp 146.86bp},clip,max width=0.94\textwidth,max height=0.38\textheight]{Marchisio2013.pdf}
\end{sourceblock}
使用中心差异，梯度项近似为

\begin{sourceblock}
\includegraphics[page=391,pagebox=mediabox,trim={185.43bp 478.59bp 185.57bp 188.57bp},clip,max width=0.94\textwidth,max height=0.38\textheight]{Marchisio2013.pdf}
\end{sourceblock}
在等式中实现这些近似。 (8.65) 产量

\begin{sourceblock}
\includegraphics[page=391,pagebox=mediabox,trim={82.05bp 419.33bp 67.56bp 246.86bp},clip,max width=0.94\textwidth,max height=0.38\textheight]{Marchisio2013.pdf}
\end{sourceblock}
该耦合方程组必须通过从初始条件 mk,i(0) 开始及时积分来求解。为了稳定性，可以使用例如 Crank-Nicolson (CN) 方案隐式处理扩散项的时间积分。如果我们用 mn k,i 表示时间 t = n Δt 时的体积平均矩，则求解方程式的半隐式方案。 (8.67) 是

\begin{sourceblock}
\includegraphics[page=391,pagebox=mediabox,trim={87.57bp 311.69bp 67.56bp 345.92bp},clip,max width=0.94\textwidth,max height=0.38\textheight]{Marchisio2013.pdf}
\end{sourceblock}
其中 λ1 = Γ Δt/(Δx2) 且 λ2 = |u|Δt/Δx。方程中的方案。 (8.68) α = 1 时完全隐含扩散，α = 1/2 时是 CN 方案。当 α = 0 时，整体方案是显式欧拉方案。方程中的系统。 (8.68) 可以重写为 mn+1 k,i 的三对角线性系统：

\begin{sourceblock}
\includegraphics[page=391,pagebox=mediabox,trim={89.36bp 229.85bp 67.56bp 453.75bp},clip,max width=0.94\textwidth,max height=0.38\textheight]{Marchisio2013.pdf}
\end{sourceblock}
其中系数定义为

\begin{sourceblock}
\includegraphics[page=391,pagebox=mediabox,trim={194.26bp 159.97bp 67.56bp 495.11bp},clip,max width=0.94\textwidth,max height=0.38\textheight]{Marchisio2013.pdf}
\end{sourceblock}
写方程。 (8.69) 由于 Amn+1 = Bmn，更新矩的形式解为 mn+1 = A−1Bmn。定义矩阵 P = A−1B，这是方程的众所周知的稳定性条件。 (8.69) 的结论是 P 的所有特征值的大小均小于 1。对于矩方法来说，同样重要的条件是矩集合的可实现性。等式可实现性的充分条件13。 （8.69）的要点是 P 的所有分量都是非负的。因为

\footnotetext[13]{时刻 n + 1 处的矩是时刻 n 处矩的加权和，其中系数是 P 的行。如果所有分量都为正，则由于矩空间是凸的，因此保证矩是可实现的。}
矩阵 A−1 具有非负分量，则 B 的分量为非负就足够了。因此，可实现矩的充分条件是 ad > 0，从而得出

\begin{sourceblock}
\includegraphics[page=392,pagebox=mediabox,trim={193.28bp 580.68bp 65.07bp 86.25bp},clip,max width=0.94\textwidth,max height=0.38\textheight]{Marchisio2013.pdf}
\end{sourceblock}
当一阶方案是半隐式且 0 < α < 1 时，此可实现性条件是纯平流和纯扩散的条件的混合。

\paragraph*{二阶方案}
对方程中的平流项使用二阶迎风重建。 （8.65），我们有

\begin{sourceblock}
\includegraphics[page=392,pagebox=mediabox,trim={122.18bp 473.77bp 119.73bp 193.17bp},clip,max width=0.94\textwidth,max height=0.38\textheight]{Marchisio2013.pdf}
\end{sourceblock}
重建的斜率由下式定义： S k,i(t) = minmod\%mk,i+1(t) −mk,i(t), mk,i(t) −mk,i−1(t)\&,
其中 minmod 斜率限制器是

\begin{sourceblock}
\includegraphics[page=392,pagebox=mediabox,trim={214.60bp 383.63bp 271.13bp 282.40bp},clip,max width=0.94\textwidth,max height=0.38\textheight]{Marchisio2013.pdf}
\end{sourceblock}
如果 ab > 0，则sign(a)min(|a|, |b|)，否则为 0。最小模（a，b）=
在扩散项中，梯度再次近似为

\begin{sourceblock}
\includegraphics[page=392,pagebox=mediabox,trim={189.17bp 329.73bp 185.52bp 337.44bp},clip,max width=0.94\textwidth,max height=0.38\textheight]{Marchisio2013.pdf}
\end{sourceblock}
以及

\begin{sourceblock}
\includegraphics[page=392,pagebox=mediabox,trim={187.93bp 290.29bp 183.08bp 376.88bp},clip,max width=0.94\textwidth,max height=0.38\textheight]{Marchisio2013.pdf}
\end{sourceblock}
在等式中实现这些近似。 (8.65) 产量

\begin{sourceblock}
\includegraphics[page=392,pagebox=mediabox,trim={79.50bp 220.44bp 65.07bp 429.39bp},clip,max width=0.94\textwidth,max height=0.38\textheight]{Marchisio2013.pdf}
\end{sourceblock}
该耦合方程组必须通过从初始条件 mk,i(0) 开始及时积分来求解。使用显式时间步进算法，方程： (8.72) 导致

\begin{sourceblock}
\includegraphics[page=392,pagebox=mediabox,trim={78.78bp 119.34bp 65.07bp 517.10bp},clip,max width=0.94\textwidth,max height=0.38\textheight]{Marchisio2013.pdf}
\end{sourceblock}
为了检查可实现性，我们可以将条款重新排列为

\begin{sourceblock}
\includegraphics[page=392,pagebox=mediabox,trim={117.17bp 57.27bp 115.22bp 605.22bp},clip,max width=0.94\textwidth,max height=0.38\textheight]{Marchisio2013.pdf}
\end{sourceblock}
当 Δt 满足可实现性条件时，右侧第一项将为非负数

\begin{sourceblock}
\includegraphics[page=393,pagebox=mediabox,trim={202.53bp 590.31bp 205.05bp 76.63bp},clip,max width=0.94\textwidth,max height=0.38\textheight]{Marchisio2013.pdf}
\end{sourceblock}
接下来的两项始终是非负的。然而，最后一项可以是正数也可以是负数，具体取决于斜率，因此我们不能保证等式： （8.73）将产生可实现的时刻。事实上，经验表明，高于一阶的标准方案几乎总是会在计算域中的某个点产生无法实现的矩。事实上，如果采用四阶方案，即使是扩散项也会产生无法实现的矩。 14 由于无法在矩反演算法中使用不可实现的矩来查找评估源项所需的权重与节点，因此标准方案对于求解矩输运方程来说是糟糕的选择。

\subsection*{8.3.3 可实现的有限体积矩格式}\phantomsection\addcontentsline{toc}{subsection}{8.3.3 可实现的有限体积矩格式}
在本节中，我们讨论等式的可实现方案。 （8.61）基于第 8.2 节和附录 B 中介绍的 KBFVM。为了实现这些方法，我们必须从矩重建 NDF。为此，我们将采用具有 N = 2 个节点的 gamma EQMOM（方程（3.84）），其中重构矩由下式给出

\begin{sourceblock}
\includegraphics[page=393,pagebox=mediabox,trim={65.02bp 351.81bp 67.56bp 295.79bp},clip,max width=0.94\textwidth,max height=0.38\textheight]{Marchisio2013.pdf}
\end{sourceblock}
αβ,i 对应于使用相关拉盖尔多项式 P(a)(t) 的高斯-拉盖尔求积，其中 a = ( Σn

\begin{sourceblock}
\includegraphics[page=393,pagebox=mediabox,trim={85.21bp 314.64bp 233.45bp 357.76bp},clip,max width=0.94\textwidth,max height=0.38\textheight]{Marchisio2013.pdf}
\end{sourceblock}
i ) −1。使用 gamma EQMOM 的矩反演算法找到两个权重 wn i：\{mn 0,i, mn 1,i, mn 2,i, mn 3,i, mn 4,i\} =⇒\{wn 1,i, wn 2,i, Σn 1,i, Σn 2,i, σn i \}。高斯-拉盖尔求积的定义为 limσ→0 Ψαβ = Ψα。此外，方程中的第二个求和。 (8.74) 对应于 gamma PDF 的矩，

\begin{sourceblock}
\includegraphics[page=393,pagebox=mediabox,trim={200.77bp 243.38bp 203.31bp 424.51bp},clip,max width=0.94\textwidth,max height=0.38\textheight]{Marchisio2013.pdf}
\end{sourceblock}
对于 2M -1 阶的矩 m(α) k 是精确的。因此，为了方程。 (8.74) 准确地说，对于 mn 4,i 之前的时刻，我们必须使 M ≥3。对于下面具有恒定扩散率的示例，我们将设置 M = 3，而对于与尺寸相关的扩散率，我们使用 M = 20。总之，将针对矩集 \{mn 0,i, mn 1,i, mn 2,i, mn 3,i, mn 4,i\} 编写可实现的方案。在每个时间步 n 和每个网格单元 i 中，我们使用 gamma EQMOM 来查找 (wn 1,i, wn 2,i, Ψn 1,i, Ψn 2,i, σn i )，从中我们定义式 (2) 中的矩重构。 （8.74）。后者将用于定义空间通量，特别是在网格单元之间的界面处开发准高阶通量重建。最后，为了方便起见，我们令Wαβ=wαwαβ。

\paragraph*{一阶方案}
{\small 对流的一阶方案与第 8.3.2 节中描述的完全相同。唯一的区别（这不是绝对必要的）是我们将替换\par}
\footnotetext[14]{在四阶扩散方案中，i + 2 和 i -2 处的矩显示为负系数。}
矩由方程式（8.74）。扩散的二阶方案 15 在附录 B 的 B.6 节中导出。使用显式时间步进方案并结合平流和扩散产量方案

\begin{sourceblock}
\includegraphics[page=394,pagebox=mediabox,trim={67.86bp 524.75bp 65.07bp 107.47bp},clip,max width=0.94\textwidth,max height=0.38\textheight]{Marchisio2013.pdf}
\end{sourceblock}
其中

\begin{sourceblock}
\includegraphics[page=394,pagebox=mediabox,trim={135.51bp 473.49bp 65.07bp 187.66bp},clip,max width=0.94\textwidth,max height=0.38\textheight]{Marchisio2013.pdf}
\end{sourceblock}
添加 Stokes-Einstein 扩散系数中的常数 Υ0 来处理消失尺寸的极限，从而使扩散系数保持有限。16 方程的可实现条件。 (8.75) 是

\begin{sourceblock}
\includegraphics[page=394,pagebox=mediabox,trim={123.99bp 387.32bp 121.54bp 271.13bp},clip,max width=0.94\textwidth,max height=0.38\textheight]{Marchisio2013.pdf}
\end{sourceblock}
请注意，不可能用斯托克斯-爱因斯坦扩散系数编写扩散项的隐式形式，因为时间 n + 1 时的权重与节点以高度非线性方式依赖于 mn+1 k,i。另外，我们应该注意到，通过增加 M 可以更准确地表示斯托克斯-爱因斯坦（或任何与 xi 相关的）扩散系数。读者应该记住，增加 M 不会改变需要求解的矩的数量，因此计算成本对 M 的值相对不敏感。

\paragraph*{准二阶方案}
平流的准二阶方案使用恒定节点和带有权重限制器的斜率估计器来重建界面通量。重建的权重由下式给出

\begin{sourceblock}
\includegraphics[page=394,pagebox=mediabox,trim={192.55bp 168.02bp 65.07bp 482.37bp},clip,max width=0.94\textwidth,max height=0.38\textheight]{Marchisio2013.pdf}
\end{sourceblock}
重建的斜率根据 Wαβ 定义：

\begin{sourceblock}
\includegraphics[page=394,pagebox=mediabox,trim={147.60bp 126.10bp 145.15bp 550.01bp},clip,max width=0.94\textwidth,max height=0.38\textheight]{Marchisio2013.pdf}
\end{sourceblock}
由于斜率限制器的特性，重建的权重遵循最大/最小原则，并且不能大于/小于它们在单元格 i −1、i 和 i 中的值

\footnotetext[15]{在均匀网格上，扩散项始终至少为二阶。 16由于 0 < Σ0 可以设置得非常小，引入它只是让我们可以在区间 [0, ∞) 而不是 [Σ0, ∞) 上定义 NDF。}
i + 1。因此，我们保证方程中的重构权重。 (8.77) 绝不会为负。扩散项的二阶有限体积近似与方程 1 中的相同。 （8.75）。因此可实现的准二阶方案为

\begin{sourceblock}
\includegraphics[page=395,pagebox=mediabox,trim={87.15bp 490.45bp 67.56bp 106.10bp},clip,max width=0.94\textwidth,max height=0.38\textheight]{Marchisio2013.pdf}
\end{sourceblock}
方程的可实现条件。 (8.78) 中涉及权重的所有系数必须是非负的。这个条件产生

\begin{sourceblock}
\includegraphics[page=395,pagebox=mediabox,trim={126.03bp 422.93bp 67.56bp 244.00bp},clip,max width=0.94\textwidth,max height=0.38\textheight]{Marchisio2013.pdf}
\end{sourceblock}
请注意，只有等式右侧的第一项。 (8.79) 是正数，可以通过减小 Δt 使其任意大。

\paragraph*{准三阶方案}
用于平流的准三阶 MUSCL 方案使用恒定节点和带有权重限制器的斜率估计器来重建界面通量。重建的权重定义为

\begin{sourceblock}
\includegraphics[page=395,pagebox=mediabox,trim={157.34bp 264.98bp 67.56bp 385.28bp},clip,max width=0.94\textwidth,max height=0.38\textheight]{Marchisio2013.pdf}
\end{sourceblock}
斜率函数的定义是

\begin{sourceblock}
\includegraphics[page=395,pagebox=mediabox,trim={168.06bp 204.34bp 67.56bp 453.26bp},clip,max width=0.94\textwidth,max height=0.38\textheight]{Marchisio2013.pdf}
\end{sourceblock}
其中差异由下式定义

\begin{sourceblock}
\includegraphics[page=395,pagebox=mediabox,trim={187.97bp 144.91bp 67.56bp 522.21bp},clip,max width=0.94\textwidth,max height=0.38\textheight]{Marchisio2013.pdf}
\end{sourceblock}
对于三阶方案，方程中的函数 φ−(a, b) 和 φ+(a, b)。 (8.81) 是针对四种不同情况定义的：

\begin{sourceblock}
\includegraphics[page=395,pagebox=mediabox,trim={75.55bp 57.42bp 121.45bp 593.71bp},clip,max width=0.94\textwidth,max height=0.38\textheight]{Marchisio2013.pdf}
\end{sourceblock}
(ii) 如果 a ≥ 0 且 b < 0：

\begin{sourceblock}
\includegraphics[page=396,pagebox=mediabox,trim={134.11bp 590.20bp 106.29bp 73.62bp},clip,max width=0.94\textwidth,max height=0.38\textheight]{Marchisio2013.pdf}
\end{sourceblock}
(iii) 如果 a、b < 0：

\begin{sourceblock}
\includegraphics[page=396,pagebox=mediabox,trim={146.78bp 528.24bp 118.96bp 135.59bp},clip,max width=0.94\textwidth,max height=0.38\textheight]{Marchisio2013.pdf}
\end{sourceblock}
(iv) 如果 a < 0 且 b ≥0：

\begin{sourceblock}
\includegraphics[page=396,pagebox=mediabox,trim={134.11bp 466.27bp 106.29bp 197.55bp},clip,max width=0.94\textwidth,max height=0.38\textheight]{Marchisio2013.pdf}
\end{sourceblock}
由平流项引起的更新可以表示为

\begin{sourceblock}
\includegraphics[page=396,pagebox=mediabox,trim={83.74bp 364.25bp 65.07bp 267.97bp},clip,max width=0.94\textwidth,max height=0.38\textheight]{Marchisio2013.pdf}
\end{sourceblock}
只要涉及权重的系数都不为负，就可以实现。实际上，当准三阶方案的可实现性条件失败时，有必要降低给定单元 i 的平流方案的阶数。附录 B 的 B.6 节描述了扩散项的准四阶方案。扩散引起的更新具有以下形式

\begin{sourceblock}
\includegraphics[page=396,pagebox=mediabox,trim={73.74bp 185.98bp 65.07bp 405.75bp},clip,max width=0.94\textwidth,max height=0.38\textheight]{Marchisio2013.pdf}
\end{sourceblock}
其中

\begin{sourceblock}
\includegraphics[page=396,pagebox=mediabox,trim={127.37bp 144.40bp 65.07bp 530.21bp},clip,max width=0.94\textwidth,max height=0.38\textheight]{Marchisio2013.pdf}
\end{sourceblock}
以及

\begin{sourceblock}
\includegraphics[page=396,pagebox=mediabox,trim={134.95bp 109.57bp 65.07bp 565.04bp},clip,max width=0.94\textwidth,max height=0.38\textheight]{Marchisio2013.pdf}
\end{sourceblock}
方程中的扩散更新。只要所有系数都是非负的，（8.84）就可以实现。这导致了两种情况。

\begin{sourceblock}
\includegraphics[page=396,pagebox=mediabox,trim={192.13bp 61.52bp 190.57bp 621.89bp},clip,max width=0.94\textwidth,max height=0.38\textheight]{Marchisio2013.pdf}
\end{sourceblock}
{\small 1. 如果 Wn αβ,i+1/2 + 0.5Dn\par}
{\small αβ,i−1/2 是负数，然后恢复到二阶扩散方案，如式（1）所示。 （8.78）。\par}
2. 否则，修复仿真参数，使得方程中的第一个系数。 (8.84) 为正。17
在实践中，第一种情况很少发生，因此方程（1）中的准四阶格式。 （8.84）几乎总是用于扩散。通过结合等式中的方案。 (8.83)和(8.84)，我们找到可实现的准三阶格式：

\begin{sourceblock}
\includegraphics[page=397,pagebox=mediabox,trim={75.06bp 438.80bp 67.56bp 124.84bp},clip,max width=0.94\textwidth,max height=0.38\textheight]{Marchisio2013.pdf}
\end{sourceblock}
方程的可实现条件。 (8.87) 的要点是方括号中的所有系数都必须是非负的。这又导致了两种情况。

\begin{sourceblock}
\includegraphics[page=397,pagebox=mediabox,trim={190.48bp 391.55bp 191.37bp 291.87bp},clip,max width=0.94\textwidth,max height=0.38\textheight]{Marchisio2013.pdf}
\end{sourceblock}
1. 如果 Wn αβ,i+1/2 + 0.5Dn αβ,i−1/2 中的任一个为负，则在单元格 i 中恢复到二阶扩散方案，如方程 1 所示。 （8.78）。

2. 否则，检查是否

\begin{sourceblock}
\includegraphics[page=397,pagebox=mediabox,trim={125.24bp 308.05bp 67.56bp 330.79bp},clip,max width=0.94\textwidth,max height=0.38\textheight]{Marchisio2013.pdf}
\end{sourceblock}
如果不是，则降低单元 i 中平流和扩散方案的阶数。

在实践中，第一种情况仅发生在少数孤立点，因此几乎总是使用准四阶扩散方案。

\paragraph*{整体有限体积方案}
上面讨论的所有有限体积方案都可以写成标准通量形式：

\begin{sourceblock}
\includegraphics[page=397,pagebox=mediabox,trim={77.15bp 185.65bp 79.67bp 479.82bp},clip,max width=0.94\textwidth,max height=0.38\textheight]{Marchisio2013.pdf}
\end{sourceblock}
(8.89) 其中 Hk,n αβ,i+1/2 和 Gk,n αβ,i+1/2 分别是扩散和对流数值通量函数。这些函数可以写成以下形式：

\begin{sourceblock}
\includegraphics[page=397,pagebox=mediabox,trim={131.82bp 80.08bp 67.56bp 545.63bp},clip,max width=0.94\textwidth,max height=0.38\textheight]{Marchisio2013.pdf}
\end{sourceblock}
\footnotetext[17]{如下所述，通常最好使用稳定性条件来固定 Δt 并降低方案的阶数以确保可实现性。}
右侧术语的定义取决于用于扩散和平流的方案的阶数。对于扩散，我们有

\begin{sourceblock}
\includegraphics[page=398,pagebox=mediabox,trim={98.72bp 582.98bp 65.07bp 80.61bp},clip,max width=0.94\textwidth,max height=0.38\textheight]{Marchisio2013.pdf}
\end{sourceblock}
以及

\begin{sourceblock}
\includegraphics[page=398,pagebox=mediabox,trim={123.96bp 539.07bp 65.07bp 124.53bp},clip,max width=0.94\textwidth,max height=0.38\textheight]{Marchisio2013.pdf}
\end{sourceblock}
而对于平流我们有

\begin{sourceblock}
\includegraphics[page=398,pagebox=mediabox,trim={126.34bp 480.07bp 65.07bp 172.26bp},clip,max width=0.94\textwidth,max height=0.38\textheight]{Marchisio2013.pdf}
\end{sourceblock}
以及

\begin{sourceblock}
\includegraphics[page=398,pagebox=mediabox,trim={125.87bp 421.87bp 65.07bp 230.48bp},clip,max width=0.94\textwidth,max height=0.38\textheight]{Marchisio2013.pdf}
\end{sourceblock}
时间步长 Δt 和所用方法的顺序确定如下。应用总体 CFL 条件来固定给定网格上的时间步长：

\begin{sourceblock}
\includegraphics[page=398,pagebox=mediabox,trim={195.05bp 359.08bp 192.60bp 307.85bp},clip,max width=0.94\textwidth,max height=0.38\textheight]{Marchisio2013.pdf}
\end{sourceblock}
其中选择CFL=1/4是为了确保准三阶平流稳定。请注意，通过选择 Δt，一阶平流和二阶扩散方案始终是稳定且可实现的。接下来，计算 λ1 = Γ Δt/(Δx2) 和 λ2 = |u|Δt/Δx。对于准四阶扩散方案，我们检查是否满足可实现性条件

\begin{sourceblock}
\includegraphics[page=398,pagebox=mediabox,trim={90.05bp 271.87bp 83.69bp 401.45bp},clip,max width=0.94\textwidth,max height=0.38\textheight]{Marchisio2013.pdf}
\end{sourceblock}
对于每个 i 都成立，如果不成立，则在单元 i 中使用二阶扩散方案。接下来，方程中的可实现性条件。对于每个 i 检查（8.88），如果失败，则平流方案在单元 i 中降低到准二阶。最后，重新检查准二阶单元的可实现性条件。 （8.88），如果再次失败，平流和扩散方案将设置为最低阶（如前所述，总是可以实现的）。以这种方式，时间步保持恒定，并且局部改变有限体积方案的阶数以确保矩的可实现性。在实践中，大多数单元在每个时间步都使用最高阶方案，特别是在扩散显著平滑了轮廓之后。

\subsection*{8.3.4 非均匀 PBE 的示例结果}\phantomsection\addcontentsline{toc}{subsection}{8.3.4 非均匀 PBE 的示例结果}
上一节中开发的可实现方案应用于图 8.2 所示的矩分布。归一化矩通过将其值除以 x = 0 处的初始值来定义：(1, 10, 200, 6000, 240 000)。对于具有恒定扩散率的情况，我们设置 M = 20，而对于具有 Stokes-Einstein 扩散率的情况，我们使用 Σ0 = 1，并设置 M = 100 以改善体积相关扩散通量的表示。

\begin{sourceblock}
\includegraphics[page=399,pagebox=mediabox,trim={58.00bp 420.17bp 58.23bp 59.00bp},clip,width=\textwidth,max height=0.72\textheight]{Marchisio2013.pdf}
\par\vspace{0.45\baselineskip}
\small 图8.2。 t = 0 时 PBE 的归一化矩。
\end{sourceblock}
\begin{sourceblock}
\includegraphics[page=399,pagebox=mediabox,trim={58.00bp 197.32bp 58.23bp 296.32bp},clip,width=\textwidth,max height=0.72\textheight]{Marchisio2013.pdf}
\par\vspace{0.45\baselineskip}
\small 图8.3。 t = 2 时 PBE 的归一化矩是在 u = 0 和恒定扩散率 Γ = 10−3 时发现的。
\end{sourceblock}
\paragraph*{恒定扩散率}
图 8.3 显示了时间 t = 2 时 u = 0 且 Γ = 10−3 的纯扩散结果。如果没有平流，我们可以观察到从方波初始条件开始，所有矩以相同的速率扩散。对于这种情况，用于表示 NDF 的所有节点均独立于 x 和 t。当 t > 0 时，四阶扩散方案几乎在域中的所有地方使用。图 8.4 显示了在时间 t = 2 时 u = 1 且 Γ = 10−3 时的平流和扩散结果。由于边界条件是周期性的，因此 t = 2 对应于方波返回到其原始位置所需的时间。

\begin{sourceblock}
\includegraphics[page=400,pagebox=mediabox,trim={58.00bp 493.17bp 58.23bp 59.00bp},clip,width=\textwidth,max height=0.72\textheight]{Marchisio2013.pdf}
\par\vspace{0.45\baselineskip}
\small 图8.4。 t = 2 时 PBE 的归一化矩是在 u = 1 和恒定扩散率 Γ = 10−3 的情况下发现的。
\end{sourceblock}
\begin{sourceblock}
\includegraphics[page=400,pagebox=mediabox,trim={58.00bp 252.71bp 58.23bp 243.19bp},clip,width=\textwidth,max height=0.72\textheight]{Marchisio2013.pdf}
\par\vspace{0.45\baselineskip}
\small 图8.5。 t = 2 时的 PBE 归一化矩是在 u = 0 和斯托克斯-爱因斯坦扩散系数 Γ = 10−3 时发现的。
\end{sourceblock}
与前面的情况一样，用于表示 NDF 的所有节点均独立于 x 和 t。当 t > 0 时，高阶平流和扩散方案几乎在域中随处使用。

\paragraph*{斯托克斯-爱因斯坦扩散率}
图 8.5 显示了在时间 t = 2 时 u = 0 且 Γ = 10−3 时斯托克斯-爱因斯坦扩散的结果。我们可以观察到低阶矩比高阶矩扩散得更快。对于这种情况，用于表示 NDF 的节点取决于 x 和 t。事实上，在方波右侧，t > 0 时的节点较小，因为小颗粒比大颗粒扩散得更快。相反，在方波内部，节点较大。时间 t = 2 时 u = 1 且 Γ = 10−3 时斯托克斯-爱因斯坦扩散的结果如图 8.6 所示。对于平流，矩轮廓的形状

\begin{sourceblock}
\includegraphics[page=401,pagebox=mediabox,trim={58.00bp 429.17bp 58.23bp 59.00bp},clip,width=\textwidth,max height=0.72\textheight]{Marchisio2013.pdf}
\par\vspace{0.45\baselineskip}
\small 图8.6。 t = 2 时 PBE 的归一化矩是在 u = 1 和斯托克斯-爱因斯坦扩散系数 Γ = 10−3 时发现的。
\end{sourceblock}
\begin{sourceblock}
\includegraphics[page=401,pagebox=mediabox,trim={58.00bp 179.10bp 58.23bp 308.90bp},clip,width=\textwidth,max height=0.72\textheight]{Marchisio2013.pdf}
\par\vspace{0.45\baselineskip}
\small 图8.7。 t = 2 时 PBE 的归一化矩是在 u = 1 和斯托克斯-爱因斯坦扩散系数 Γ = 10−3 时发现的。 CFL 编号设置为 1。
\end{sourceblock}
对 NDF 的 EQMOM 重建敏感，该重建与第二个求积一起使用来估计与尺寸相关的扩散系数。一般来说，通过增加 n（EQMOM 节点的数量）可以得到更好的结果，但代价是传输更多的矩。在图 8.7 中，我们重复与图 8.6 中相同的模拟，但 CFL 数设置为 1。可以看出，方波前部的扩散性低于后侧，这是由于使用 CFL 数高于其稳定极限的高阶平流方案而产生的数值伪影。这可以通过切换来验证

\begin{sourceblock}
\includegraphics[page=402,pagebox=mediabox,trim={58.00bp 430.17bp 58.23bp 59.00bp},clip,width=\textwidth,max height=0.72\textheight]{Marchisio2013.pdf}
\par\vspace{0.45\baselineskip}
\small 图8.8。 t = 2 时 PBE 的归一化矩是在 u = 1 和斯托克斯-爱因斯坦扩散系数 Γ = 10−3 时发现的。采用CFL数为1的一阶平流方案。
\end{sourceblock}
为一阶平流格式，如图 8.8 所示。请注意，在所有情况下，时刻仍然是可实现的，这是通过时间步长和可实现方案的选择来保证的。

\section*{8.4 非均匀KE}\phantomsection\addcontentsline{toc}{section}{8.4 非均匀KE}
在本节中，我们将研究速度 NDF n(t, x, v) 的动力学方程，其中 v = (u, v) 是二分量速度矢量（即速度相空间是二维的）。为了显示不同数量的粒子-粒子碰撞的动力学，我们将使用 BGK 碰撞模型。 （有关碰撞模型的更多详细信息，请参阅第 6 章。）这种情况的非齐次动力学方程为

\begin{sourceblock}
\includegraphics[page=402,pagebox=mediabox,trim={188.73bp 199.83bp 65.07bp 473.29bp},clip,max width=0.94\textwidth,max height=0.38\textheight]{Marchisio2013.pdf}
\end{sourceblock}
其中 τ 是碰撞时间常数（与克努森数成正比），而 θ 是根据双变量速度矩定义的颗粒温度：

\begin{sourceblock}
\includegraphics[page=402,pagebox=mediabox,trim={166.05bp 123.28bp 163.60bp 538.56bp},clip,max width=0.94\textwidth,max height=0.38\textheight]{Marchisio2013.pdf}
\end{sourceblock}
在这里，我们使用速度相空间中第二方向和第三方向之间的统计对称性来用二元矩 mj,k 来表示颗粒温度。 NDF n* 是具有守恒性质的平衡（麦克斯韦）分布：m*
{\small 0,0 = m0,0，m* 1,0 = m1,0，m* 0,1 = m0,1，且 θ*= θ。从物理上讲，这些等式源于碰撞过程中质量、动量和能量守恒的事实。\par}
\subsection*{8.4.1 矩输运方程}\phantomsection\addcontentsline{toc}{subsection}{8.4.1 矩输运方程}
从方程中找到的双变量矩的传输方程。 (8.95) 是

\begin{sourceblock}
\includegraphics[page=403,pagebox=mediabox,trim={161.98bp 574.39bp 67.56bp 97.46bp},clip,max width=0.94\textwidth,max height=0.38\textheight]{Marchisio2013.pdf}
\end{sourceblock}
其中平衡矩 m* j,k 是守恒量 m0,0、m1,0、m0,1 和 θ 的已知函数。方程式中的碰撞项。 (8.96) 是封闭的，因此唯一未封闭的项是涉及 mj+1,k 的空间通量。为了结束这项，我们将使用第 3.3.4 节中描述的多元 EQMOM 和高斯函数（详细信息请参见方程（3.134））：

\begin{sourceblock}
\includegraphics[page=403,pagebox=mediabox,trim={155.70bp 481.75bp 67.56bp 186.14bp},clip,max width=0.94\textwidth,max height=0.38\textheight]{Marchisio2013.pdf}
\end{sourceblock}
其中 δσ(x, y) 由方程式定义。 （3.88）。 （另请参见第 98 页的方程（3.141）。）方程的矩反演算法。 （8.97）在等式之后描述。 (3.140)，第 98 页，需要以下一组 13 个矩：

\begin{sourceblock}
\includegraphics[page=403,pagebox=mediabox,trim={177.71bp 369.75bp 180.24bp 264.12bp},clip,max width=0.94\textwidth,max height=0.38\textheight]{Marchisio2013.pdf}
\end{sourceblock}
然而，根据对称性，如果初始 NDF 是高斯分布，则四个矩 m0,1、m1,1、m0,3 和 m1,3 为零，在这种情况下，仅需要九个矩即可求解方程 1。 （8.96）。然而，求解所有 13 个矩（以下称为传输矩集）并使用已知的零矩来检查算法的数值精度是很方便的。

\subsection*{8.4.2 矩方程的算子分裂}\phantomsection\addcontentsline{toc}{subsection}{8.4.2 矩方程的算子分裂}
由于碰撞时间 τ 是任意的并且可能导致方程僵化，因此最好求解方程 1。 （8.96）使用算子分裂，即首先求解

\begin{sourceblock}
\includegraphics[page=403,pagebox=mediabox,trim={203.60bp 211.20bp 67.56bp 462.45bp},clip,max width=0.94\textwidth,max height=0.38\textheight]{Marchisio2013.pdf}
\end{sourceblock}
随后

\begin{sourceblock}
\includegraphics[page=403,pagebox=mediabox,trim={182.88bp 169.42bp 67.56bp 502.43bp},clip,max width=0.94\textwidth,max height=0.38\textheight]{Marchisio2013.pdf}
\end{sourceblock}
方程式中的碰撞项。 (8.99) 可以通过设置半解析求解

\begin{sourceblock}
\includegraphics[page=403,pagebox=mediabox,trim={185.55bp 135.01bp 188.09bp 550.31bp},clip,max width=0.94\textwidth,max height=0.38\textheight]{Marchisio2013.pdf}
\end{sourceblock}
这导致近似

\begin{sourceblock}
\includegraphics[page=403,pagebox=mediabox,trim={127.63bp 91.82bp 130.16bp 584.72bp},clip,max width=0.94\textwidth,max height=0.38\textheight]{Marchisio2013.pdf}
\end{sourceblock}
{\small 根据均匀网格上的有限体积矩编写，该公式变为\par}
\begin{sourceblock}
\includegraphics[page=403,pagebox=mediabox,trim={168.83bp 50.23bp 67.57bp 625.64bp},clip,max width=0.94\textwidth,max height=0.38\textheight]{Marchisio2013.pdf}
\end{sourceblock}
\begin{sourceblock}
\includegraphics[page=404,pagebox=mediabox,trim={98.67bp 590.39bp 175.11bp 60.00bp},clip,max width=0.94\textwidth,max height=0.38\textheight]{Marchisio2013.pdf}
\end{sourceblock}
where κn
矩 m† j,k;i 是通过求解方程得出的。 (8.98) 从 mn j,k;i 开始，平衡 j,k;i （其中 j, k = 0, 1, 2）.18 因为方程： (8.100) 是可实现矩的凸和，对于 Δt ≥0 总是可实现的。因此，在实践中，时间步长Δt可以仅根据空间通量的可实现条件来选择。仍有待描述方程中的空间通量如何。 （8.98）在可实现的有限体积方案中处理速度矩。 j,k;i 是 m† 矩 m* 的已知函数 在离开碰撞项之前，我们应该向读者指出方程的右侧。式（8.99）可以用更复杂的碰撞模型代替，例如第6章中描述的玻尔兹曼硬球模型。一般来说，碰撞模型在传递矩方面不会封闭，因此矩更新将比式（8.99）中的更复杂。 （8.100）。例如，对于玻尔兹曼碰撞项，右侧将涉及乘积 NDF n(t, x, v1)n(t, x, v2) 的二重积分（参见第 217 页上的方程（6.9））。通过使用等式中的多高斯闭包。 (8.96) 中，二重积分可以写成高斯-埃尔米特求积节点上的二重求和（更多详细信息，请参见第 3.3.1 节）。对于此处考虑的情况，高斯-埃尔米特求积将用于表示方程 1 中给出的二元矩源项中的积分。第 232.19 页上的 (6.68) 请注意，表 6.1–6.3 中给出了计算硬球碰撞源项所需的表达式 I(0) l1l2l3 （其中 ω = 1）。

\subsection*{8.4.3 双变量速度矩的可实现的有限体积格式}\phantomsection\addcontentsline{toc}{subsection}{8.4.3 双变量速度矩的可实现的有限体积格式}
方程的可实现方案。 (8.98) 可以写成标准通量形式：

\begin{sourceblock}
\includegraphics[page=404,pagebox=mediabox,trim={67.52bp 230.16bp 65.08bp 399.42bp},clip,max width=0.94\textwidth,max height=0.38\textheight]{Marchisio2013.pdf}
\end{sourceblock}
(i) 从传输矩集 mn j,k;i 开始，应用矩反演算法求出四个高斯分布的参数：wα、uα、vα、σ1α 和 σ2α。请注意，这里我们对 α = 1, 进行一次求和。。。，4，代替等式中使用的双重求和。 (3.134) 第 96 页。因此，例如，σ1α 对于所有 α 都是相同的，并且 u1 = u2、u3 = u4、σ21 = σ22 和 σ23 = σ24，因为速度分量 v 取决于 u 的值。实施矩反演算法所需的具体步骤在方程 1 后描述。 (3.140) 第 98 页。

18 回想一下，碰撞项具有守恒性质，因此 m† 0,0;i、m† 1,0;i、m† 0,1;i 和 m† 2,0;i + 2m† 0,2;i 在碰撞步骤中不会改变。因此，在没有空间通量的情况下，κ(t, x) 是恒定的。

\footnotetext[19]{速度相空间第三方向上的高斯-埃尔米特求积与速度分量 v 相同，因此可以形成为第二方向上的求积权重与节点的张量积。然而，在实践中，重写方程中的项在计算上更加有效。 （6.68）利用对称性来消除对第三速度分量的依赖性。}
(ii) 给定 M1 和 M2，找到相应的高斯-埃尔米特求积权重并

\begin{sourceblock}
\includegraphics[page=405,pagebox=mediabox,trim={139.01bp 610.88bp 271.14bp 70.89bp},clip,max width=0.94\textwidth,max height=0.38\textheight]{Marchisio2013.pdf}
\end{sourceblock}
节点 (w(1) β2 )，我们假设它是标准化的，使得权重总和为 1，方差为 1。

(iii) 对于每个 α，求积权重与节点由下式定义：

\begin{sourceblock}
\includegraphics[page=405,pagebox=mediabox,trim={212.19bp 520.88bp 67.57bp 124.67bp},clip,max width=0.94\textwidth,max height=0.38\textheight]{Marchisio2013.pdf}
\end{sourceblock}
请注意，由于高斯-埃尔米特求积权重与节点仅取决于点数，因此如果 M1 和 M2 恒定（通常是这种情况），则只需计算一次。在附录 B 的 B.4 节中，我们提出了一种使用等式 1 中定义的半矩来计算通量正交的替代方法。 (B.37) 第 431 页，带有 CQMOM 的第二次应用。实际上，CQMOM 的精度和在没有过多舍入误差的情况下可以确定的正交节点的数量之间存在权衡。因此，如果需要许多通量节点才能准确计算空间通量，则即使半矩是精确的，也可能很难使用 CQMOM 找到所需的精度。相比之下，使用高斯-埃尔米特求积法，M1 和 M2 可以做得非常大，而不会损失任何精度，因为它们不需要矩反演算法。在实践中，这意味着我们应该能够在不经过半矩的情况下实现非常准确的通量表示。例如，当 M = 14 时，最小的高斯-埃尔米特权重约为 10−8，并且添加更多节点对数值通量函数几乎没有任何影响。对于这里考虑的情况，我们发现使用 M1 = 14 和 M2 = 4 可以准确地模拟传递矩集。但是，一般来说，所需的节点数量将取决于传递矩集中矩的最高阶。对于方程中的自由传输问题。 (8.98)，可实现的数值通量函数为

\begin{sourceblock}
\includegraphics[page=405,pagebox=mediabox,trim={96.70bp 235.72bp 67.57bp 419.32bp},clip,max width=0.94\textwidth,max height=0.38\textheight]{Marchisio2013.pdf}
\end{sourceblock}
权重 Wn αβ1β2,i+1/2,l 和 Wn αβ1β2,i+1/2,r 取决于方案的顺序。例如，一阶方案使用

\begin{sourceblock}
\includegraphics[page=405,pagebox=mediabox,trim={194.81bp 176.31bp 67.57bp 490.57bp},clip,max width=0.94\textwidth,max height=0.38\textheight]{Marchisio2013.pdf}
\end{sourceblock}
而准二阶和准三阶方案使用类似于方程的形式。分别为（8.77）和（8.80）。由方程得出的可实现性条件。 (8.101) 是

\begin{sourceblock}
\includegraphics[page=405,pagebox=mediabox,trim={87.92bp 123.95bp 67.57bp 551.76bp},clip,max width=0.94\textwidth,max height=0.38\textheight]{Marchisio2013.pdf}
\end{sourceblock}
对于 (α, β1, β2, i) 的每个值。实际上，时间步长 Δt 是使用 CFL 条件确定的，该条件基于非零权重的一阶方案（即当 0 < Wn

\begin{sourceblock}
\includegraphics[page=405,pagebox=mediabox,trim={73.32bp 50.16bp 67.57bp 606.74bp},clip,max width=0.94\textwidth,max height=0.38\textheight]{Marchisio2013.pdf}
\end{sourceblock}
\begin{sourceblock}
\includegraphics[page=406,pagebox=mediabox,trim={58.00bp 443.79bp 58.23bp 59.00bp},clip,width=\textwidth,max height=0.72\textheight]{Marchisio2013.pdf}
\par\vspace{0.45\baselineskip}
\small 图8.9。黎曼激波的非零初始矩。
\end{sourceblock}
且 CFL < 1。当方程 1 时，使用准三阶/二阶对流方案。 (8.105) 满足固定的 λ1 = Δt/Δx，否则对于给定的 (α, β1, β2, i) 值采用低阶重构方案。对于一阶方案，方程。 (8.106) 足以保证更新后的矩集 m† j,k;i 的可实现性。总之，可实现的方案 j,k;i 并使用等式从 mn j,k;i 开始计算 m†。 （8.101）。该结果与等式一起使用。 (8.100) 求更新矩 mn+1 j,k;i。只要等式。满足(8.105)和(8.106)，则保证更新后的矩是可实现的。

\subsection*{8.4.4 非均匀 KE 的示例结果}\phantomsection\addcontentsline{toc}{subsection}{8.4.4 非均匀 KE 的示例结果}
在这里，我们应用有限体积方案来模拟非齐次动力学方程的两个不同示例。第一个例子是具有不同碰撞时间 τ 值的非平衡黎曼激波问题。第二个示例是具有不同碰撞时间的两个一维交叉射流。实际上，碰撞时间由数密度 M00 控制，在这些示例中我们将其相对于统一进行归一化。因此，读者可以将 τ 的不同值解释为非归一化数密度的不同值。如上所述，对于多高斯求积，我们使用 M1 = 14 和 M2 = 4 以及统一的 CFL 数来计算空间通量。

\paragraph*{黎曼激波}
黎曼激波问题中非零矩的初始条件如图 8.9 所示。在域的左半部分，归一化数密度为 M00 = 1，在右半部分为 0.1。最初，域两侧的 RMS 速度均为 1，平均速度为零。矩被初始化为麦克斯韦矩，因此初始条件处于局部平衡。随着时间的增加，从 x = 0 开始的冲击波向右移动，而通货紧缩波向左移动。在 τ = 0 的极限下，解与气体动力学欧拉方程相同。时间 t = 0.5 时 τ = 100 时刻的示例结果如图 8.10 所示。对于这种情况，碰撞非常弱，因此几乎没有动能从 u 分量转移到

\begin{sourceblock}
\includegraphics[page=407,pagebox=mediabox,trim={58.00bp 445.17bp 58.23bp 59.00bp},clip,width=\textwidth,max height=0.72\textheight]{Marchisio2013.pdf}
\par\vspace{0.45\baselineskip}
\small 图8.10。黎曼激波的非零矩，t = 0.5 时 τ = 100。
\end{sourceblock}
\begin{sourceblock}
\includegraphics[page=407,pagebox=mediabox,trim={58.00bp 228.38bp 58.23bp 281.64bp},clip,width=\textwidth,max height=0.72\textheight]{Marchisio2013.pdf}
\par\vspace{0.45\baselineskip}
\small 图8.11。黎曼激波的统计数据，t = 0.5 时 τ = 100。键表示将 ex 和 ey 绘制为 e 的线，对于键中的 q 和 η 也同样如此。
\end{sourceblock}
速度波动的 v 分量。这可以从图 8.11 中看出，其中 v 统计量保持在平衡值附近。例如 σ21 = σ22 ≈1，表明所有 v 速度波动都包含在高斯速度分布中。然而，σ1 远小于 1，表明速度的 u 分量非常不平衡。同样，对于所有 x，v 分量中的能量为 ey ≈1，v 偏度为 qy ≈0，v 峰度为 ηy ≈3。相反，速度 u 分量的值远未达到平衡状态。正如黎曼激波问题所预期的那样，平均 u 速度 Up 不为零。请注意，权重和 u 速度节点远离平衡值。相比之下，v 速度节点非常接近于零，正如高斯速度分布所预期的那样。时间 t = 0.5 时 τ = 1 时刻的示例结果如图 8.12 所示。对于这种情况，碰撞仍然是

\begin{sourceblock}
\includegraphics[page=408,pagebox=mediabox,trim={58.00bp 440.17bp 58.23bp 59.00bp},clip,width=\textwidth,max height=0.72\textheight]{Marchisio2013.pdf}
\par\vspace{0.45\baselineskip}
\small 图8.12。黎曼激波的非零矩，t = 0.5 时 τ = 1。
\end{sourceblock}
\begin{sourceblock}
\includegraphics[page=408,pagebox=mediabox,trim={58.00bp 211.03bp 58.23bp 287.69bp},clip,width=\textwidth,max height=0.72\textheight]{Marchisio2013.pdf}
\par\vspace{0.45\baselineskip}
\small 图8.13。黎曼激波的统计，t = 0.5 时 τ = 1。请参见图 8.11 中不同行的关键信息。
\end{sourceblock}
相对较弱，因此没有太多动能从速度波动的 u 分量转移到 v 分量。这可以从图 8.13 中看出，其中 v 统计数据仅略微偏离平衡值。 σ21 和 σ22 也接近一致，表明 v 速度波动接近高斯速度分布。同样，σ1 明显小于 1，表明速度的 u 分量不平衡。正如预期的那样，平均 u 速度 Up 不为零，并且权重和 u 速度节点远离平衡值。然而，现在 v 速度节点离零更远，表明偏离高斯速度分布。时间 t = 0.5 时 τ = 0.01 时刻的示例结果如图 8.14 所示。对于

\begin{sourceblock}
\includegraphics[page=409,pagebox=mediabox,trim={58.00bp 438.17bp 58.23bp 59.00bp},clip,width=\textwidth,max height=0.72\textheight]{Marchisio2013.pdf}
\par\vspace{0.45\baselineskip}
\small 图8.14。黎曼激波的非零矩，t = 0.5 时 τ = 0.01。
\end{sourceblock}
\begin{sourceblock}
\includegraphics[page=409,pagebox=mediabox,trim={58.00bp 200.56bp 58.23bp 297.97bp},clip,width=\textwidth,max height=0.72\textheight]{Marchisio2013.pdf}
\par\vspace{0.45\baselineskip}
\small 图8.15。 t = 0.5 时 τ = 0.01 的黎曼激波统计。请参见图 8.11 中不同行的关键信息。
\end{sourceblock}
在这种情况下，碰撞相对较强，因此动能从速度波动的 u 分量到 v 分量有显著的转移。这可以在图 8.15 中看到，其中 v 统计量接近 u 分量的值。在这种情况下，σ21 和 σ22 接近于 1，但对 x 具有复杂的依赖性。 σ1 也是如此。在 τ = 0 的极限下，σ 的所有分量都等于 1。平均 u 速度 Up 再次非零，并且权重与节点表现出对 x 的复杂依赖性。总体而言，可实现方案能够处理整个 τ 值范围，同时确保矩始终可实现。

\paragraph*{穿越喷气式飞机}
交叉射流问题中非零矩的初始条件如图 8.16 所示。初始数密度是均匀的，M00 = 1，但在域的左半部分，平均速度为 Up = 1，而在右半部分，平均速度为 -1。最初，域两侧的 RMS 速度均为 0.1，矩初始化为麦克斯韦矩。随着时间的增加，左半侧的密度向右移动，右半侧的密度向左移动，在 x = 0 处产生一维交叉射流。根据 τ 的值，两个射流将松弛到平衡速度分布函数。时间 t = 1 时 τ = 100 的结果如图 8.17 和 8.18 所示。对于这个 τ 值，碰撞非常弱，因此松弛到平衡的速度非常慢。这是显而易见的

\begin{sourceblock}
\includegraphics[page=410,pagebox=mediabox,trim={58.00bp 298.04bp 58.23bp 200.75bp},clip,width=\textwidth,max height=0.72\textheight]{Marchisio2013.pdf}
\par\vspace{0.45\baselineskip}
\small 图8.16。交叉射流的非零初始矩。
\end{sourceblock}
\begin{sourceblock}
\includegraphics[page=410,pagebox=mediabox,trim={58.00bp 69.44bp 58.23bp 430.15bp},clip,width=\textwidth,max height=0.72\textheight]{Marchisio2013.pdf}
\par\vspace{0.45\baselineskip}
\small 图8.17。 t = 1 时 τ = 100 的交叉射流的非零矩。
\end{sourceblock}
\begin{sourceblock}
\includegraphics[page=411,pagebox=mediabox,trim={58.00bp 440.16bp 58.23bp 59.00bp},clip,width=\textwidth,max height=0.72\textheight]{Marchisio2013.pdf}
\par\vspace{0.45\baselineskip}
\small 图8.18。 t = 1 时 τ = 100 交叉射流的统计数据。
\end{sourceblock}
\begin{sourceblock}
\includegraphics[page=411,pagebox=mediabox,trim={58.00bp 214.01bp 58.23bp 285.38bp},clip,width=\textwidth,max height=0.72\textheight]{Marchisio2013.pdf}
\par\vspace{0.45\baselineskip}
\small 图8.19。 t = 1 时 τ = 1 的交叉射流的非零矩。
\end{sourceblock}
来自能量 ey、偏度 qy 和峰度 ηy 的 v 分量统计，它们非常接近其初始值。相比之下，u 分量的相同统计数据表现出与平衡的强烈偏差（正如交叉射流所预期的那样）。 σ 各分量的行为也很有趣。在域的中心，σ1 几乎为零，表明速度的 u 分量几乎具有以两个速度节点 uα 为中心的双狄拉克分布。相反，分量 σ21 和 σ22 彼此偏离并且偏离统一，正如速度 v 分量不平衡时所预期的那样。由于射流交叉，区域中心的平均密度 ρ 较高，平均速度 Up 表现出复杂的空间依赖性。 t = 1 时 τ = 1 的结果如图 8.19 和 8.20 所示。对于这个 τ 值，碰撞是显著的，因此松弛到平衡的速度相当快。能源的 v 分量统计

\begin{sourceblock}
\includegraphics[page=412,pagebox=mediabox,trim={58.00bp 440.16bp 58.23bp 59.00bp},clip,width=\textwidth,max height=0.72\textheight]{Marchisio2013.pdf}
\par\vspace{0.45\baselineskip}
\small 图8.20。 t = 1 时与 τ = 1 交叉的射流的统计数据。参见图 8.18 中不同线的关键部分。
\end{sourceblock}
\begin{sourceblock}
\includegraphics[page=412,pagebox=mediabox,trim={58.00bp 202.06bp 58.23bp 297.34bp},clip,width=\textwidth,max height=0.72\textheight]{Marchisio2013.pdf}
\par\vspace{0.45\baselineskip}
\small 图8.21。 t = 1 时 τ = 0.01 的交叉射流的非零矩。
\end{sourceblock}
ey、偏度 qy 和峰度 ηy 显著偏离平衡值，而对于 u 分量，这些统计数据再次表现出与平衡值的强烈偏差。在域的中心，σ1 再次接近于零，但现在分量 σ21 和 σ22 与平衡值的偏差非常大。另请注意，节点 vα 现在在域的很大一部分中远离零。正如预期的那样，区域中心的平均密度 ρ 较高，平均速度 Up 从左到右单调递减。时间 t = 1 时 τ = 0.01 的结果如图 8.21 和 8.22 所示。对于这个 τ 值，碰撞非常强烈，因此松弛到平衡的速度很快。能量 ey、偏度 qy 和峰度 ηy 的 v 分量统计数据在各处都接近平衡值，除了两个射流交叉点之外。对于 u 分量，

\begin{sourceblock}
\includegraphics[page=413,pagebox=mediabox,trim={58.00bp 440.17bp 58.23bp 59.00bp},clip,width=\textwidth,max height=0.72\textheight]{Marchisio2013.pdf}
\par\vspace{0.45\baselineskip}
\small 图8.22。 t = 1 时 τ = 0.01 交叉射流的统计数据。参见图 8.18 中不同线的关键部分。
\end{sourceblock}
这些统计数据也接近与 v 分量的平衡。除射流交叉点外，节点 vα 在大部分域中接近于零。平均密度 ρ 在域中心接近方波，而平均速度 Up 在域中心接近零。总的来说，我们可以从交叉射流示例中得出结论，与碰撞项耦合的可实现方案可以毫无困难地捕获速度相空间中的复杂物理场，同时始终产生可实现的矩集。

\section*{8.5 非均匀GPBE}\phantomsection\addcontentsline{toc}{section}{8.5 非均匀GPBE}
\subsection*{8.5.1 GPBE 的类别}\phantomsection\addcontentsline{toc}{subsection}{8.5.1 GPBE 的类别}
在广义 PBE (GPBE) 中，相空间变量包括粒子速度和粒子属性（例如尺寸或体积）。一般来说，GPBE 主要分为两类，这取决于粒子速度的表示方式：

(1) 粒子速度由已知的与标量相关的平流速度 u(xi) 给出，并且
(2) 粒子速度由联合速度标量 NDF 求得。

第一类是第 8.3 节中讨论的平流扩散问题的概括，其中开发的大部分材料可以重新表述以开发可实现的 GPBE 方案。第二类是第 8.4 节中考虑的动力学方程的推广，在文献中被称为半动力学模型（Laurent 等人，2004 年；Laurent 和 Massot，2001 年）。在下文中，我们将分别对待每一类，尽管读者无疑会注意到数值方案之间的许多相似之处。第三类 GPBE，位于上面列出的两类之间，也可以被识别，其中标量相关的对流速度具有参数形式，例如

\begin{sourceblock}
\includegraphics[page=413,pagebox=mediabox,trim={162.00bp 52.15bp 67.57bp 633.42bp},clip,max width=0.94\textwidth,max height=0.38\textheight]{Marchisio2013.pdf}
\end{sourceblock}
特别是，左侧可以解释为以颗粒尺寸为条件的颗粒速度，右侧是条件速度的局部（在物理空间中）二阶近似。然而，为了准确地再现矩通量，方程中的近似值： (8.107) 必须接近真实速度。或者，换句话说，粒子速度的条件方差必须很小，以便条件速度分布函数紧密地以 u(t, x, xi) 为中心。例如，当粒子斯托克斯数低于 PTC 开始发生的临界值时（即惯性非常小的粒子），情况就是如此。

\paragraph*{速度参数的确定}
方程中的流量相关参数 u0、u1 和 u2 (8.107) 可以通过引入速度 v 和尺寸 xi 之间的交叉矩来确定。这是通过注意到当联合速度标量 NDF 为时对流速度的参数形式有效来完成的

\begin{sourceblock}
\includegraphics[page=414,pagebox=mediabox,trim={173.15bp 454.18bp 65.08bp 223.27bp},clip,max width=0.94\textwidth,max height=0.38\textheight]{Marchisio2013.pdf}
\end{sourceblock}
其中 n(xi) 是 NDF 的大小。因此，联合整数矩由下式给出

\begin{sourceblock}
\includegraphics[page=414,pagebox=mediabox,trim={168.48bp 412.37bp 65.08bp 260.78bp},clip,max width=0.94\textwidth,max height=0.38\textheight]{Marchisio2013.pdf}
\end{sourceblock}
前三个交叉时刻是

\begin{sourceblock}
\includegraphics[page=414,pagebox=mediabox,trim={183.46bp 343.63bp 65.08bp 313.19bp},clip,max width=0.94\textwidth,max height=0.38\textheight]{Marchisio2013.pdf}
\end{sourceblock}
通过传输三个附加矩 m1,0、m1,1 和 m1,2，该线性系统可以求解方程式中出现的三个未知参数。 （8.107）：

\begin{sourceblock}
\includegraphics[page=414,pagebox=mediabox,trim={176.51bp 270.50bp 65.08bp 385.60bp},clip,max width=0.94\textwidth,max height=0.38\textheight]{Marchisio2013.pdf}
\end{sourceblock}
右侧逆矩阵中出现的矩是大小 NDF 的纯矩。因此，只要NDF的大小不是由一个或两个δ函数组成，逆矩阵就可以很好地定义。否则，在逆条件条件较差的情况下，方程中的近似阶数为： （8.107）可以简单地通过将u2（以及，如果需要，u0）设置为零，并求解一组减少的参数来减少。20请注意，这是在空间和时间上局部完成的，因为参数随着混合矩而演变。然后就足以为扩展矩集\{m0,0, m0,1, m0,2, m0,3, m0,4, m1,0, m1,1, m1,2\}导出与GPBE一致的可实现方案。我们在第 8.5.4 节中考虑具有斯托克斯阻力的粒子系统示例。

20 更准确地说，定义

\begin{sourceblock}
\includegraphics[page=414,pagebox=mediabox,trim={195.39bp 110.79bp 192.94bp 554.66bp},clip,max width=0.94\textwidth,max height=0.38\textheight]{Marchisio2013.pdf}
\end{sourceblock}
如果 A 的条件数足够小以进行精确反演，我们使用式（1）。 （8.111）。否则我们设置u2 = 0，重新定义系数矩阵为

\begin{sourceblock}
\includegraphics[page=414,pagebox=mediabox,trim={207.23bp 72.27bp 204.78bp 600.19bp},clip,max width=0.94\textwidth,max height=0.38\textheight]{Marchisio2013.pdf}
\end{sourceblock}
{\small 并求解 u0 和 u1 的简化系统。然而，如果这个新的 A 的条件也很差（即 xi 的方差可以忽略不计），那么我们设置 u0 = 0 和 u1 = m1,1/m0,2。\par}
\paragraph*{使用 EQMOM 确定速度参数}
如果使用 EQMOM 来重建大小 NDF，则 xi ∈[0, 1] 的联合速度标量 NDF 可以写成分段线性形式：21

\begin{sourceblock}
\includegraphics[page=415,pagebox=mediabox,trim={171.75bp 555.63bp 67.57bp 121.50bp},clip,max width=0.94\textwidth,max height=0.38\textheight]{Marchisio2013.pdf}
\end{sourceblock}
其中参数 wα、σ 和 Σα 在第 3.3 节中定义，条件速度建模为 U(Σ) = u0g0(Σ) + u1g1(Σ) + u2g2(Σ) + u3g3(Σ)，函数 gn 定义为

\begin{sourceblock}
\includegraphics[page=415,pagebox=mediabox,trim={152.54bp 403.62bp 67.57bp 185.56bp},clip,max width=0.94\textwidth,max height=0.38\textheight]{Marchisio2013.pdf}
\end{sourceblock}
如果 x ∈[x1, x2]，则指示函数 I[x1,x2](x) 为 1，否则为零。请注意，g0(0) = 1、g1(Σ1) = 1、g2(Σ2) = 1 和 g3(1) = 1。一般来说，对于 beta EQMOM，0 < Σ1 < Σ2 < 1；然而，如果 Ψ1 = Ψ2，则 u1 = u2。在 σ = 0 的简并情况下，只有 g1 和 g2 非零，因此 u0 和 u3 没有定义。四个速度参数（u0、u1、u2 和 u3）由四个混合矩确定。混合整数矩由下式给出

\begin{sourceblock}
\includegraphics[page=415,pagebox=mediabox,trim={85.08bp 288.36bp 67.57bp 375.22bp},clip,max width=0.94\textwidth,max height=0.38\textheight]{Marchisio2013.pdf}
\end{sourceblock}
使用 j = 1 和 k = 0, 1, 2, 3 找到速度参数。这些矩可以写为

\begin{sourceblock}
\includegraphics[page=415,pagebox=mediabox,trim={172.67bp 243.57bp 67.57bp 440.25bp},clip,max width=0.94\textwidth,max height=0.38\textheight]{Marchisio2013.pdf}
\end{sourceblock}
其中

\begin{sourceblock}
\includegraphics[page=415,pagebox=mediabox,trim={101.61bp 190.29bp 67.57bp 473.29bp},clip,max width=0.94\textwidth,max height=0.38\textheight]{Marchisio2013.pdf}
\end{sourceblock}
wαβ 和 xiαβ 是第二个高斯-雅可比求积的权重与节点。请注意，M 必须选择足够大，以便准确计算 Gk n。速度参数的线性系统为

\begin{sourceblock}
\includegraphics[page=415,pagebox=mediabox,trim={175.00bp 84.81bp 67.57bp 556.02bp},clip,max width=0.94\textwidth,max height=0.38\textheight]{Marchisio2013.pdf}
\end{sourceblock}
\footnotetext[21]{为了清楚起见，我们考虑为 Ψ ε[0, 1] 定义的双节点 beta EQMOM 的情况。毫无疑问，更通用的形式是可能的。}
如果 Σ1 < Σ2 且 σ > 0，左侧的系数矩阵将是可逆的。在 σ = 0 的极限情况下，系数矩阵为

\begin{sourceblock}
\includegraphics[page=416,pagebox=mediabox,trim={203.21bp 552.00bp 200.76bp 89.74bp},clip,max width=0.94\textwidth,max height=0.38\textheight]{Marchisio2013.pdf}
\end{sourceblock}
只要 Ψ1 < Ψ2 则其等级为 2；否则为秩1。在秩1的情况下，只能确定一个速度参数（u1），而在秩2的情况下，可以找到u1和u2。这些可以通过删除方程式中系统的最后两行来找到（对于 σ > 0）。 （8.117）。因此，等式。求解方程（8.117）即可找到四个速度参数。然后就足以为扩展矩集\{m0,0, m0,1, m0,2, m0,3, m0,4, m1,0, m1,1, m1,2, m1,3\}导出一个与GPBE一致的可实现方案。有趣的是，通过 EQMOM 重建，比方程 1 传输的矩更多。 （8.108），这表明使用方程可以获得更好的表示。 （8.112）。然而，由于需要调整算法以解决等式可能存在的秩不足问题。 （8.117），其数值实现比较困难。

\subsection*{8.5.2 已知标量相关速度的空间传输}\phantomsection\addcontentsline{toc}{subsection}{8.5.2 已知标量相关速度的空间传输}
设 Ψ 表示一个非负标量，NDF n(t, x, Ψ)。推广方程。 （8.61），我们假设 GPBE 是

\begin{sourceblock}
\includegraphics[page=416,pagebox=mediabox,trim={188.54bp 331.21bp 65.08bp 343.15bp},clip,max width=0.94\textwidth,max height=0.38\textheight]{Marchisio2013.pdf}
\end{sourceblock}
其中平流速度 u(xi) 是 xi 的已知函数。这种情况下的矩输运方程是

\begin{sourceblock}
\includegraphics[page=416,pagebox=mediabox,trim={139.06bp 264.73bp 65.08bp 400.96bp},clip,max width=0.94\textwidth,max height=0.38\textheight]{Marchisio2013.pdf}
\end{sourceblock}
很明显，空间通量需要闭合。正如第 8.3.3 节中所做的那样，基于求积的闭包可用于此目的：

\begin{sourceblock}
\includegraphics[page=416,pagebox=mediabox,trim={160.21bp 188.55bp 157.76bp 475.03bp},clip,max width=0.94\textwidth,max height=0.38\textheight]{Marchisio2013.pdf}
\end{sourceblock}
其中 Wαβ 和 xiαβ 取决于高斯-拉盖尔求积节点。对于这种情况，M 的值必须选择得足够大，以代表所需精度的磁通积分。实际上，这取决于 u(xi) 的确切函数形式。求解方程的可实现方案。 （8.119）与式（8.119）相同。 (8.89)，但 Gk,n 的定义略有不同

\begin{sourceblock}
\includegraphics[page=416,pagebox=mediabox,trim={73.81bp 85.75bp 83.52bp 570.06bp},clip,max width=0.94\textwidth,max height=0.38\textheight]{Marchisio2013.pdf}
\end{sourceblock}
{\small αβ,i+1/2,l 和 Wn\par}
{\small αβ,i+1/2,r 与等式中给出的相同。 （8.93）和（8.94）。方程中的可实现条件。 (8.88) 必须修改为\par}
\begin{sourceblock}
\includegraphics[page=417,pagebox=mediabox,trim={80.46bp 587.20bp 67.57bp 56.44bp},clip,max width=0.94\textwidth,max height=0.38\textheight]{Marchisio2013.pdf}
\end{sourceblock}
与平流扩散情况一样，时间步长通过设置总体 CFL 条件来固定：

\begin{sourceblock}
\includegraphics[page=417,pagebox=mediabox,trim={167.75bp 518.67bp 170.28bp 137.49bp},clip,max width=0.94\textwidth,max height=0.38\textheight]{Marchisio2013.pdf}
\end{sourceblock}
利用该 Δt 值，网格单元 i 中的对流方案的阶数由等式确定。 （8.121）。

\subsection*{8.5.3 已知标量相关速度的示例结果}\phantomsection\addcontentsline{toc}{subsection}{8.5.3 已知标量相关速度的示例结果}
作为与标量相关的速度的示例，我们将考虑速度是 xi 的线性函数的两种情况。为简单起见，我们假设在时间 0 ψ 可以用区间 [0, 1] 上的 beta 分布来表示，矩的空间分布如图 8.23 所示。正如附录 C 中所讨论的，矩的演变取决于速度是大小的增函数还是减函数。在这里，我们考虑一个速度为 u(Ψ) = −1 + 2Ψ 的分离系统，定义为当 Ψ > 0.5 时，粒子上升，而当 Ψ < 0.5 时，粒子下降；以及 u(xi) = 1 −2xi 的混合系统。从图 8.23 的初始条件中，我们可以观察到，最初较小的粒子位于 x = 0 点的下方，而较大的粒子位于该点的上方。因此，对于分离系统，粒子将远离 x = 0，从而产生空隙；而对于混合系统，域的上半部分和下半部分上的粒子将向彼此移动，从而导致粒子混合。对于这些模拟，我们将使用 n = 2 的 beta EQMOM 来重建 NDF，并将扩散率 Γ 设置为零。因此，如附录 C 所示，GPBE 是纯双曲的，其交叉特征为

\begin{sourceblock}
\includegraphics[page=417,pagebox=mediabox,trim={58.00bp 69.07bp 58.23bp 445.90bp},clip,width=\textwidth,max height=0.72\textheight]{Marchisio2013.pdf}
\par\vspace{0.45\baselineskip}
\small 图8.23。标量相关速度情况的初始时刻。
\end{sourceblock}
\begin{sourceblock}
\includegraphics[page=418,pagebox=mediabox,trim={58.00bp 445.17bp 58.23bp 68.74bp},clip,width=\textwidth,max height=0.72\textheight]{Marchisio2013.pdf}
\par\vspace{0.45\baselineskip}
\small 图8.24。分离系统在 t = 0.2 时的矩。
\end{sourceblock}
混合箱。为了获得高精度的空间通量，我们在第二个高斯求积中使用 M = 100 并将 CFL 数设置为 1。因此，所有正交误差均归因于第一个 beta EQMOM 正交。请注意，与附录 C 不同，初始 NDF 被清晰地分成两部分（即 xi = 0.5 上方和下方），这里 EQMOM 表示在整个区间 [0, 1] 上使用 beta 分布，这意味着 x = 0 上方和下方最初将存在所有可能尺寸的粒子（尽管平均尺寸非常不同，如图 8.23 所示）。

\paragraph*{隔离系统}
时间 t = 0.2 时分离系统的样本结果如图 8.24 和 8.25 所示。正如预期的那样，当粒子平流远离中点时，m0 在 x = 0 处具有最小值。基础 GPBE 的双曲线性质在图 8.24 中的矩曲线中清晰可见，我们看到粒子相对于图 8.23 中的初始位置精确地上升/下降了 0.2 个空间单位。由于分离系统不存在交叉特征，因此矩分布的行为符合预期，较小的颗粒下降，较大的颗粒上升。然而，图 8.25 中所示的 beta EQMOM 的权重与节点表现出复杂的空间依赖性，扩展参数 σ 也是如此。请注意，平均粒径 m1/m0 随着 x 从零单调增加到 1，因此 beta EQMOM 能够处理两种极端情况，其中 NDF 在区间边界处几乎是 delta 函数。最后，图 8.25 显示了 α = 1 时前五个通量节点 xiαβ 和通量权重 Wαβ 的空间分布。从通量节点可以看出，通过对通量使用 M = 100 的第二高斯求积，可以很好地对尺寸相空间进行采样。

\paragraph*{混合系统}
{\small 时间 t = 0.2 时混合系统的样本结果如图 8.26 和 8.27 所示。正如预期的那样，m0 在 x = 0 处具有最大值，因为粒子从\par}
\begin{sourceblock}
\includegraphics[page=419,pagebox=mediabox,trim={58.00bp 344.16bp 58.23bp 59.00bp},clip,width=\textwidth,max height=0.72\textheight]{Marchisio2013.pdf}
\par\vspace{0.45\baselineskip}
\small 图8.25。 t = 0.2 时分离系统的统计数据。
\end{sourceblock}
\begin{sourceblock}
\includegraphics[page=419,pagebox=mediabox,trim={58.00bp 80.67bp 58.23bp 413.58bp},clip,width=\textwidth,max height=0.72\textheight]{Marchisio2013.pdf}
\par\vspace{0.45\baselineskip}
\small 图8.26。混合系统在 t = 0.2 时的矩。
\end{sourceblock}
\begin{sourceblock}
\includegraphics[page=420,pagebox=mediabox,trim={58.00bp 296.17bp 58.23bp 59.00bp},clip,width=\textwidth,max height=0.72\textheight]{Marchisio2013.pdf}
\par\vspace{0.45\baselineskip}
\small 图8.27。 t = 0.2 时混合系统的统计数据。
\end{sourceblock}
中点。基础 GPBE 的双曲线性质在图 8.26 中的矩曲线中再次清晰可见，我们看到粒子相对于图 8.23 中的初始位置恰好上升/下降了 0.2 个空间单位。在此示例中，交叉特性是一个重要特征，并且当较大/较小颗粒在空间中重叠时，矩分布表现出混合行为。这里，图 8.27 中所示的 beta EQMOM 的权重与节点再次表现出复杂的空间依赖性，扩展参数 σ 也是如此。此外，由于特性的混合行为，平均粒径 m1/m0 首先从 1 减小，然后在 x < 0 时增大到 0.5。在图 8.27 中，α = 1 时前五个通量节点 xiαβ 和通量权重 Wαβ 的空间分布表现出对 x 的非平凡依赖性。总体而言，可实现的方案在处理由于初始条件和速度对颗粒尺寸的依赖性而交叉的特性方面没有困难。如前所述，矩输运方程解的精度由 beta EQMOM 求积控制，并且可以通过使用较大的 n 值（即 EQMOM 节点数）来提高精度。然而，尽管它们的空间轮廓复杂，矩集总是可以实现的。

\subsection*{8.5.4 标量条件速度的空间传输}\phantomsection\addcontentsline{toc}{subsection}{8.5.4 标量条件速度的空间传输}
在本节中，我们考虑粒子速度由方程式给出的中间情况。 (8.107) 具有与流量相关的参数 u0、u1 和 u2，或通过方程 (8.107) 中给出的 EQMOM 表示。 （8.112）。作为起点，我们将采用以下 GPBE 来表示关节速度大小 NDF n(t, x, v, xi)：

\begin{sourceblock}
\includegraphics[page=421,pagebox=mediabox,trim={192.78bp 531.89bp 67.57bp 142.36bp},clip,max width=0.94\textwidth,max height=0.38\textheight]{Marchisio2013.pdf}
\end{sourceblock}
其中 τ(ψ) = τ0ψ2/3 是归一化体积 ψ = 1 的粒子的阻力特征时间尺度。在此 GPBE 中，ψ 是粒子体积，阻力项隐式假设流体速度为零。因此，阻力项将具有以与尺寸相关的衰减率衰减初始粒子速度的效果。对应于方程的矩输运方程。 (8.122) 是

\begin{sourceblock}
\includegraphics[page=421,pagebox=mediabox,trim={184.72bp 418.34bp 67.57bp 254.43bp},clip,max width=0.94\textwidth,max height=0.38\textheight]{Marchisio2013.pdf}
\end{sourceblock}
因此，扩展矩集由 22 决定

\begin{sourceblock}
\includegraphics[page=421,pagebox=mediabox,trim={187.50bp 151.13bp 187.17bp 318.87bp},clip,max width=0.94\textwidth,max height=0.38\textheight]{Marchisio2013.pdf}
\end{sourceblock}
请注意，空间通量项和阻力项都需要矩闭合。这些闭包将取决于 NDF 是否由等式 1 表示。 （8.108）或等式。 （8.112）。当 NDF 由式（1）表示时： (8.108)，不使用 m1,3 的传输方程。

\footnotetext[22]{矩集的选择并不是唯一的。然而，矩 m1,1 对应于动量密度，它是一个整体守恒量，因此是包含在扩展矩集中的自然选择。}
\paragraph*{方程中的速度参数。 (8.108)}
作为一个例子，我们将使用一个两节点 beta EQMOM 来表示 n(xi)，因此，混合矩可以用封闭形式表示为

\begin{sourceblock}
\includegraphics[page=422,pagebox=mediabox,trim={131.10bp 556.52bp 65.08bp 111.37bp},clip,max width=0.94\textwidth,max height=0.38\textheight]{Marchisio2013.pdf}
\end{sourceblock}
其中 Wαβ 和 xiαβ 取决于高斯-雅可比求积节点。函数 u0(t, x)、u1(t, x) 和 u2(t, x) 是使用方程 (1) 从传输矩集中找到的。 （8.111）。方程式中阻力项中出现的分数矩。 (8.123) 可以使用等式进行评估。 (8.124) 其中 k 设置为适当的分数。如果 NDF 在 xi = 0 处具有非零值，则矩 m1,−2/3 可能是无穷大，因此数值方法必须考虑这种可能性。方程中的系统。使用算子分裂将通量项与阻力项分开，可以方便地求解方程（8.123）。第一步，通过求解来更新通量项

\begin{sourceblock}
\includegraphics[page=422,pagebox=mediabox,trim={206.09bp 429.73bp 65.08bp 243.92bp},clip,max width=0.94\textwidth,max height=0.38\textheight]{Marchisio2013.pdf}
\end{sourceblock}
求解方程的可实现方案。 （8.125）与式（8.125）相同。 （8.89）（没有扩散通量），但数值通量的定义略有不同：

\begin{sourceblock}
\includegraphics[page=422,pagebox=mediabox,trim={92.48bp 324.19bp 65.08bp 296.54bp},clip,max width=0.94\textwidth,max height=0.38\textheight]{Marchisio2013.pdf}
\end{sourceblock}
\begin{sourceblock}
\includegraphics[page=422,pagebox=mediabox,trim={166.81bp 254.07bp 65.08bp 399.82bp},clip,max width=0.94\textwidth,max height=0.38\textheight]{Marchisio2013.pdf}
\end{sourceblock}
其中右侧仅显示体积平均传递矩。注意，参数的这种指定与可实现的方案一致，其中速度节点在每个网格单元中是恒定的。方程中的可实现条件。 (8.88) 必须修改为

\begin{sourceblock}
\includegraphics[page=422,pagebox=mediabox,trim={85.42bp 175.46bp 80.11bp 497.68bp},clip,max width=0.94\textwidth,max height=0.38\textheight]{Marchisio2013.pdf}
\end{sourceblock}
(8.128) 其中 λ1 = Δt/Δx。通过设置总体 CFL 条件来固定时间步长：

\begin{sourceblock}
\includegraphics[page=422,pagebox=mediabox,trim={192.76bp 105.62bp 190.31bp 546.84bp},clip,max width=0.94\textwidth,max height=0.38\textheight]{Marchisio2013.pdf}
\end{sourceblock}
\footnotetext[23]{通过计算方程 23 的体积平均值，找到每个单元中混合矩的有限体积表示。 （8.110）。右侧显示 u0m0,0 等产品的平均成交量。因此，如果 u0 在单元中是分段常数，则体积平均乘积取决于体积平均矩，这是有限体积方案中的守恒量。}
对于 CFL ≤1。利用该 Δt 值，网格单元 i 中的对流方案的阶数由等式确定。 （8.128）。在平流步骤结束时，更新的矩集由 m* j,k;i 表示，并且保证是可实现的。接下来，使用具有固定权重与节点的速度参数的等效 ODE 系统求解阻力项。 ODE 系统的发现起始于

\begin{sourceblock}
\includegraphics[page=423,pagebox=mediabox,trim={198.70bp 546.39bp 200.04bp 127.10bp},clip,max width=0.94\textwidth,max height=0.38\textheight]{Marchisio2013.pdf}
\end{sourceblock}
必须用初始条件 m1,k;i(0) = m* 1,k;i 来求解。通过替换以权重与节点表示的矩，可以得到表达式

\begin{sourceblock}
\includegraphics[page=423,pagebox=mediabox,trim={123.65bp 478.10bp 126.18bp 189.78bp},clip,max width=0.94\textwidth,max height=0.38\textheight]{Marchisio2013.pdf}
\end{sourceblock}
我们为每个网格单元 i 找到一个 ODE 系统：

\begin{sourceblock}
\includegraphics[page=423,pagebox=mediabox,trim={77.21bp 380.79bp 67.57bp 246.62bp},clip,max width=0.94\textwidth,max height=0.38\textheight]{Marchisio2013.pdf}
\end{sourceblock}
这与以下相同

\begin{sourceblock}
\includegraphics[page=423,pagebox=mediabox,trim={95.45bp 317.35bp 67.57bp 342.20bp},clip,max width=0.94\textwidth,max height=0.38\textheight]{Marchisio2013.pdf}
\end{sourceblock}
纯尺寸矩 m0,k;i 在拖动步骤中与时间无关，因此在给定从方程 1 找到的初始条件的情况下，可以针对参数 u0;i(t)、u1;i(t) 和 u2;i(t) 求解该常数系数线性 ODE 系统。 (8.127) 和矩 m* j,k;i。当 k = 0, 1, 2 时，由方程得出的线性 ODE 系统： （8.130）可以写成（为了清楚起见，删除 i 依赖性）24

\begin{sourceblock}
\includegraphics[page=423,pagebox=mediabox,trim={197.68bp 214.12bp 67.57bp 441.98bp},clip,max width=0.94\textwidth,max height=0.38\textheight]{Marchisio2013.pdf}
\end{sourceblock}
其中系数矩阵根据大小矩定义为

\begin{sourceblock}
\includegraphics[page=423,pagebox=mediabox,trim={125.41bp 154.48bp 67.57bp 501.62bp},clip,max width=0.94\textwidth,max height=0.38\textheight]{Marchisio2013.pdf}
\end{sourceblock}
请注意，在阻力评估期间，尺寸矩 m0,k 是常数 25，因此 A 和 B 是常数矩阵。此外，这些系数矩阵都是对称且正定的，这意味着一阶 ODE 系统的特征值
24 如果流体速度非零，则方程的右侧为（8.131）将有一个额外的非齐次项，具体取决于流体速度。如果流体速度 uf 与时间无关，则只需将右侧的 u0 替换为 u0 -uf 即可找到等效表示。

\footnotetext[25]{尺寸矩是恒定的，因为权重与节点是恒定的，这与推导方程 25 所用的假设一致。 （8.130）。}
等式。 (8.131) 是实值且为负值。因此，参数 u0(t)、u1(t) 和 u2(t) 将随着 t 的增加呈指数衰减至零。26 使用隐式欧拉方案更新参数向量（并重写等式（8.111）中的定义以重写右侧），我们发现对于网格单元 i：

\begin{sourceblock}
\includegraphics[page=424,pagebox=mediabox,trim={183.04bp 544.77bp 65.08bp 109.13bp},clip,max width=0.94\textwidth,max height=0.38\textheight]{Marchisio2013.pdf}
\end{sourceblock}
其中 m* j,k;i 是对流步骤后单元 i 中的体积平均矩。矩阵 A*

\begin{sourceblock}
\includegraphics[page=424,pagebox=mediabox,trim={207.24bp 513.51bp 245.23bp 169.91bp},clip,max width=0.94\textwidth,max height=0.38\textheight]{Marchisio2013.pdf}
\end{sourceblock}
i 和 B* i 使用 W* αβ;i 进行评估，它们是使用 EQMOM 与矩 m* 0,k;i 找到的权重与节点。时间步长 Δt 与对流步骤中使用的时间步长相同。当 m0,−2/3 →∞ 时，会出现一个重要的特殊情况，即 NDF 的液滴由于蒸发等原因而接近零体积。在这种情况下，等式右侧的逆矩阵。 (8.133) 减少至27

\begin{sourceblock}
\includegraphics[page=424,pagebox=mediabox,trim={67.52bp 392.62bp 65.06bp 230.09bp},clip,max width=0.94\textwidth,max height=0.38\textheight]{Marchisio2013.pdf}
\end{sourceblock}
0;当 m0,−2/3 = ∞ 时，i = 0。在实践中，等式。 （8.134）可以用在等式中。 (8.133) 当 m0,−2/3 相对于 m0,0 非常大时。因为 m0,−2/3 是根据高斯求积权重与节点 (W*

\begin{sourceblock}
\includegraphics[page=424,pagebox=mediabox,trim={79.14bp 342.76bp 371.47bp 340.67bp},clip,max width=0.94\textwidth,max height=0.38\textheight]{Marchisio2013.pdf}
\end{sourceblock}
αβ;i)，它永远不会是无限的。方程精确解中的矩阵指数。 (8.131) 可能很难实现，特别是在 B 是病态的情况下。物理上一致且数值上稳健的曳力的替代公式可以按如下方式实现。

(i) 在平流步骤之后，使用矩 m* 0,k;i 计算权重 W* αβ,i 和节点 xi*

\begin{sourceblock}
\includegraphics[page=424,pagebox=mediabox,trim={135.79bp 278.46bp 341.10bp 407.34bp},clip,max width=0.94\textwidth,max height=0.38\textheight]{Marchisio2013.pdf}
\end{sourceblock}
(ii) 构造方程中的矩阵 A。 (8.127) 使用矩 m* 0,k;i，并计算平流产生的速度参数：

\begin{sourceblock}
\includegraphics[page=424,pagebox=mediabox,trim={194.60bp 184.89bp 174.10bp 449.67bp},clip,max width=0.94\textwidth,max height=0.38\textheight]{Marchisio2013.pdf}
\end{sourceblock}
26 精确解是 u(t) = exp A−1m(0)，其中矩阵 C = A−1B 是正定的。因为 m0,−2/3 是根据求积求得的（B 中出现的所有分数矩也是如此），所以它应该始终是有界的，尽管它可能非常大。因此，原则上，精确解可用于更新速度系数。 27 若用线性函数表示u(xi)，则对应的矩阵为

\begin{sourceblock}
\includegraphics[page=424,pagebox=mediabox,trim={173.11bp 122.25bp 170.67bp 551.07bp},clip,max width=0.94\textwidth,max height=0.38\textheight]{Marchisio2013.pdf}
\end{sourceblock}
且极限表达式为

\begin{sourceblock}
\includegraphics[page=424,pagebox=mediabox,trim={163.02bp 87.02bp 160.57bp 586.15bp},clip,max width=0.94\textwidth,max height=0.38\textheight]{Marchisio2013.pdf}
\end{sourceblock}
{\small 在这种情况下，un+1\par}
{\small 0;i = 0 且 un+1 1;i = m* 1,1;i/(m* 0,2;i + m* 0,4/3;iτ)。 28根据τ0的定义，τ的单位为长度平方，因此A和B的分量具有一致的单位。\par}
回想一下，A 可能是病态的（甚至是奇异的），因此在这一步中处理退化情况时必须小心。

(iii) 使用速度参数创建样本速度对应的集合 -

\begin{sourceblock}
\includegraphics[page=425,pagebox=mediabox,trim={89.93bp 571.63bp 67.56bp 101.00bp},clip,max width=0.94\textwidth,max height=0.38\textheight]{Marchisio2013.pdf}
\end{sourceblock}
(iv) 根据阻力定律更新样本速度：

\begin{sourceblock}
\includegraphics[page=425,pagebox=mediabox,trim={145.01bp 503.93bp 121.43bp 147.61bp},clip,max width=0.94\textwidth,max height=0.38\textheight]{Marchisio2013.pdf}
\end{sourceblock}
此步骤中的标量指数始终是明确定义的，即使在 xi*

\begin{sourceblock}
\includegraphics[page=425,pagebox=mediabox,trim={94.28bp 466.67bp 362.48bp 218.36bp},clip,max width=0.94\textwidth,max height=0.38\textheight]{Marchisio2013.pdf}
\end{sourceblock}
(v) 使用新的样本速度来计算样本混合矩：

\begin{sourceblock}
\includegraphics[page=425,pagebox=mediabox,trim={190.26bp 410.07bp 167.88bp 257.82bp},clip,max width=0.94\textwidth,max height=0.38\textheight]{Marchisio2013.pdf}
\end{sourceblock}
(vi) 使用样本混合矩和 A 更新速度参数

\begin{sourceblock}
\includegraphics[page=425,pagebox=mediabox,trim={89.93bp 332.53bp 191.96bp 313.08bp},clip,max width=0.94\textwidth,max height=0.38\textheight]{Marchisio2013.pdf}
\end{sourceblock}
该公式对于存在其他力（例如浮力、虚拟质量、升力等）的情况也很方便。对于这种情况，只需在步骤 (iv) 中将额外的力添加到样本速度的更新中即可。更新速度参数后，拖动步骤结束时的传输矩可从

\begin{sourceblock}
\includegraphics[page=425,pagebox=mediabox,trim={80.10bp 237.38bp 67.57bp 421.50bp},clip,max width=0.94\textwidth,max height=0.38\textheight]{Marchisio2013.pdf}
\end{sourceblock}
因为权重在拖动步骤期间是恒定的，并且在平流步骤之后始终为非负值，所以更新的矩在方程中设置。 （8.135）保证可以实现。在第 8.5.5 节中，我们将上述算法应用于 τ0 小到足以避免 PTC 的情况。

\paragraph*{方程中的速度参数。 (8.112)}
我们将再次使用二节点 beta EQMOM 来表示 n(xi)，因此混合矩可以用闭合形式表示为

\begin{sourceblock}
\includegraphics[page=425,pagebox=mediabox,trim={171.77bp 84.23bp 67.57bp 583.66bp},clip,max width=0.94\textwidth,max height=0.38\textheight]{Marchisio2013.pdf}
\end{sourceblock}
{\small 条件速度由 U(t, x, xi) = u0(t, x)g0(ψ) + u1(t, x)g1(xi) + u2(t, x)g2(ψ) + u3(t, x)g3(ψ) 定义，其中 Wαβ 和 ψαβ 取决于高斯-雅可比求积\par}
节点。函数 uk(t, x) 是使用方程式从传输矩集中找到的。 （8.117）。函数 gk(xi) 通过 EQMOM 节点 xiα(t, x) 取决于 t 和 x。为了避免繁琐的符号，我们假设 gk(xi) 的 t 和 x 依赖性与其参数相同。例如，gk(Σn αβ,i) 在时间步 n 和网格单元 i 处计算。方程中的系统。使用算子分裂将通量项与阻力项分开，可以方便地求解方程（8.123）。第一步，通过求解方程更新通量项。 （8.125）。求解方程的可实现方案。 （8.125）与式（8.125）相同。 (8.89)（无扩散通量），但数值通量 G j,k;n 的定义略有不同

\begin{sourceblock}
\includegraphics[page=426,pagebox=mediabox,trim={67.52bp 453.40bp 65.06bp 157.04bp},clip,max width=0.94\textwidth,max height=0.38\textheight]{Marchisio2013.pdf}
\end{sourceblock}
k;i 被假定为每个有限体积单元中的分段常数。从等式开始（8.136），这导致

\begin{sourceblock}
\includegraphics[page=426,pagebox=mediabox,trim={155.16bp 411.83bp 65.08bp 271.06bp},clip,max width=0.94\textwidth,max height=0.38\textheight]{Marchisio2013.pdf}
\end{sourceblock}
其中，遵循等式： (8.116),

\begin{sourceblock}
\includegraphics[page=426,pagebox=mediabox,trim={174.34bp 359.06bp 65.08bp 308.83bp},clip,max width=0.94\textwidth,max height=0.38\textheight]{Marchisio2013.pdf}
\end{sourceblock}
式中取 k = 0, 1, 2, 3 （8.138）导致线性系统：29

\begin{sourceblock}
\includegraphics[page=426,pagebox=mediabox,trim={159.06bp 272.92bp 65.08bp 360.59bp},clip,max width=0.94\textwidth,max height=0.38\textheight]{Marchisio2013.pdf}
\end{sourceblock}
其中，只有从传输时刻找到的组件才会出现在右侧。注意，速度参数的这种指定与可实现的方案一致，其中速度节点在每个网格单元中是恒定的。方程中的可实现条件。 (8.88) 必须修改为

\begin{sourceblock}
\includegraphics[page=426,pagebox=mediabox,trim={86.92bp 194.89bp 65.08bp 478.25bp},clip,max width=0.94\textwidth,max height=0.38\textheight]{Marchisio2013.pdf}
\end{sourceblock}
其中 λ1 = Δt/Δx。通过设置总体 CFL 条件来固定时间步长：

\begin{sourceblock}
\includegraphics[page=426,pagebox=mediabox,trim={199.71bp 133.86bp 197.26bp 520.89bp},clip,max width=0.94\textwidth,max height=0.38\textheight]{Marchisio2013.pdf}
\end{sourceblock}
对于 CFL ≤1。利用该 Δt 值，网格单元 i 中的对流方案的阶数由等式确定。 （8.141）。在平流步骤结束时，更新的矩集由 m* j,k;i 表示，并且保证是可实现的。接下来，如下评估阻力项。

\footnotetext[29]{系数矩阵的性质在等式之后讨论。 （8.117），并且数值方法必须考虑可能的奇异情况。}
(i) 使用矩 m* 0,k;i 计算权重 W* αβ,i 和节点 xi*

\begin{sourceblock}
\includegraphics[page=427,pagebox=mediabox,trim={362.19bp 624.40bp 114.69bp 61.41bp},clip,max width=0.94\textwidth,max height=0.38\textheight]{Marchisio2013.pdf}
\end{sourceblock}
(ii) 求解方程式中的线性系统。 (8.140) 的速度参数由

\begin{sourceblock}
\includegraphics[page=427,pagebox=mediabox,trim={89.93bp 544.79bp 193.98bp 92.97bp},clip,max width=0.94\textwidth,max height=0.38\textheight]{Marchisio2013.pdf}
\end{sourceblock}
其中

\begin{sourceblock}
\includegraphics[page=427,pagebox=mediabox,trim={195.91bp 472.06bp 173.53bp 163.64bp},clip,max width=0.94\textwidth,max height=0.38\textheight]{Marchisio2013.pdf}
\end{sourceblock}
\begin{sourceblock}
\includegraphics[page=427,pagebox=mediabox,trim={346.65bp 449.34bp 106.00bp 234.09bp},clip,max width=0.94\textwidth,max height=0.38\textheight]{Marchisio2013.pdf}
\end{sourceblock}
分量 Gk;* m;i 是根据方程计算的。 (8.139) 使用 W* αβ,i。回想一下，A 可能是病态的（即当节点彼此接近或域的边界时），因此在此步骤中处理退化情况时必须小心。

(iii) 使用速度参数创建样本速度对应的集合 -

\begin{sourceblock}
\includegraphics[page=427,pagebox=mediabox,trim={89.93bp 370.56bp 67.56bp 301.52bp},clip,max width=0.94\textwidth,max height=0.38\textheight]{Marchisio2013.pdf}
\end{sourceblock}
(iv) 根据阻力定律更新样本速度：

\begin{sourceblock}
\includegraphics[page=427,pagebox=mediabox,trim={142.79bp 305.53bp 119.22bp 348.30bp},clip,max width=0.94\textwidth,max height=0.38\textheight]{Marchisio2013.pdf}
\end{sourceblock}
此步骤中的标量指数始终是明确定义的，即使在 xi*

\begin{sourceblock}
\includegraphics[page=427,pagebox=mediabox,trim={94.28bp 271.59bp 362.48bp 413.44bp},clip,max width=0.94\textwidth,max height=0.38\textheight]{Marchisio2013.pdf}
\end{sourceblock}
(v) 使用新的样本速度来计算样本混合矩：

\begin{sourceblock}
\includegraphics[page=427,pagebox=mediabox,trim={190.26bp 213.91bp 167.88bp 453.98bp},clip,max width=0.94\textwidth,max height=0.38\textheight]{Marchisio2013.pdf}
\end{sourceblock}
(vi) 使用样本混合矩和 A 更新速度参数

\begin{sourceblock}
\includegraphics[page=427,pagebox=mediabox,trim={89.93bp 121.72bp 191.96bp 510.17bp},clip,max width=0.94\textwidth,max height=0.38\textheight]{Marchisio2013.pdf}
\end{sourceblock}
该公式对于存在其他力（例如浮力、虚拟质量、升力等）的情况也很方便。对于这种情况，只需在步骤 (iv) 中将额外的力添加到样本速度的更新中即可。由于函数 gk 在拖动步骤期间是恒定的（即它们仅取决于大小矩），因此也可以构造线性 ODE 系统，例如方程 1。 （8.131）。然而，在实践中，上述算法更方便，并且提供了更大的灵活性来合并其他力。更新速度参数后，拖动步骤结束时的传输矩可从

\begin{sourceblock}
\includegraphics[page=428,pagebox=mediabox,trim={102.71bp 552.87bp 65.08bp 87.52bp},clip,max width=0.94\textwidth,max height=0.38\textheight]{Marchisio2013.pdf}
\end{sourceblock}
其中Un+1

\subsection*{8.5.5 标量条件速度的示例结果}\phantomsection\addcontentsline{toc}{subsection}{8.5.5 标量条件速度的示例结果}
在本节中，我们将使用方程 1 中的 NDF 表示，展示上一节中介绍的标量条件速度算法得出的结果。 （8.108）。所有这些计算都是使用 beta EQMOM 完成的，第一个求积使用 n = 2 个节点，第二个求积使用 M = 50 个节点。初始 NDF 具有两种形式之一，其中 xi ∈[0, 1] 是粒子体积。在第一种情况下，称为“小分布”，NDF 由两个 beta 函数表示，一个接近 xi = 1/3，一个接近 xi = 2/3，具有较小的分布参数 σ。这种情况的初始时刻如图 8.28 所示。从图中可以看出，仅在空间域 x ε(−0.8, −0.2) 内，矩才非零。图 8.29 给出了根据初始 NDF 计算的一些统计数据。请注意，初始条件速度 u(0, x, xi) = 1 与颗粒尺寸无关，并且体积的平均值和标准偏差不取决于 x。

\begin{sourceblock}
\includegraphics[page=428,pagebox=mediabox,trim={58.00bp 69.07bp 58.23bp 354.64bp},clip,width=\textwidth,max height=0.72\textheight]{Marchisio2013.pdf}
\par\vspace{0.45\baselineskip}
\small 图8.28。具有小扩展的标量条件速度情况的初始时刻。
\end{sourceblock}
\begin{sourceblock}
\includegraphics[page=429,pagebox=mediabox,trim={58.00bp 262.80bp 58.23bp 59.00bp},clip,width=\textwidth,max height=0.72\textheight]{Marchisio2013.pdf}
\par\vspace{0.45\baselineskip}
\small 图8.29。小范围标量条件速度情况的初始统计。
\end{sourceblock}
对于这种情况， ε = 0 处的 NDF 为空。第二种情况，称为“大分布”，在区间 xi ∈[0, 1] 上具有均匀的初始 NDF。对于这种情况， ε = 0 处的 NDF 是一致的。矩的初始空间分布如图 8.30 所示，选定的统计数据如图 8.31 所示。小价差和大价差情况之间的主要区别是成交量标准差。根据定义，n = 2 的 beta EQMOM 可以准确地重建两个初始 NDF。除了这些情况之外，我们还成功地测试了具有更大分布的算法，其中 xi = 0 处的精确 NDF 是无限的（m0,−2/3 的精确值也是如此）。作为示例结果，我们考虑 τ0 的三个值：10、1 和 0.1。请注意，由于初始平均粒子速度为 1（且流体速度为零），因此初始斯托克斯数等于 τ0。如前所述，标量条件速度模型仅对小斯托克斯数有效。然而，如果我们避免 PTC，该模型原则上可以用于更大的斯托克斯数。由于我们的模拟域具有周期性边界条件，因此我们可以通过在粒子退出域之前停止模拟来避免 PTC（即 t = 1.2，τ0 = ∞）。此外，由于粒子速度因阻力而减小，斯托克斯数将随着

\begin{sourceblock}
\includegraphics[page=430,pagebox=mediabox,trim={58.00bp 359.80bp 58.23bp 67.14bp},clip,width=\textwidth,max height=0.72\textheight]{Marchisio2013.pdf}
\par\vspace{0.45\baselineskip}
\small 图8.30。具有大扩展的标量条件速度情况的初始时刻。
\end{sourceblock}
时间。事实上，对于 τ0 ≤1，粒子将由于阻力而在域结束之前停止。然而，我们将看到 τ0 > 0.1 的情况可能会导致与 PTC 无关的其他问题。

\paragraph*{价差小}
我们首先查看 τ0 = 10（即非常弱的阻力）的情况，其中 t = 1 时的矩如图 8.32 所示，样本统计数据如图 8.33 所示。从矩图中可以明显看出，初始 NDF 主要只是在空间中以接近 1 的平均速度传输，这与初始条件一致。从高阶矩可以看出，空间分布的左侧有更多的小粒子，而右侧有更多的大粒子。从图 8.33 中的平均体积图也可以明显看出这一点。质量平均速度（由 m1,1/m0,1 定义）在分布左侧略低，在右侧略高。从体积标准差的分布中，我们可以观察到，作为颗粒体积函数的速度差异导致大颗粒和小颗粒分离，小颗粒落后于领先的大颗粒。最后，图 8.33 显示了五个 x 值的条件速度 u(xi) = u0 + u1ψ + u2ψ2。由于分布的展度很小，因此在这种情况下，基本上不存在体积为 Ψ = 0 或 Ψ = 1 的粒子。然而，函数 u(xi) 可以外推到这些点。请注意，由于阻力非常弱（斯托克斯数高），因此条件速度非常缓慢地向零移动（即向流体速度移动）。 t = 3 时 τ0 = 1（即中等阻力）的结果如图 8.34 和 8.35 所示。如上所述，此时 NDF 为

\begin{sourceblock}
\includegraphics[page=431,pagebox=mediabox,trim={58.00bp 267.80bp 58.23bp 59.00bp},clip,width=\textwidth,max height=0.72\textheight]{Marchisio2013.pdf}
\par\vspace{0.45\baselineskip}
\small 图8.31。大范围标量条件速度情况的初始统计。
\end{sourceblock}
由于速度接近于零，因此基本上冻结在太空中。矩的空间分布显示了大颗粒和小颗粒的非常强烈的空间分离，图 8.35 中颗粒体积的平均值和标准偏差图也是如此。尽管质量平均速度很小，但其形状遵循平均体积的形状。另请注意，平均体积曲线的最大值和最小值大致对应于小扩散情况下使用的两个 beta 函数的峰值，因此颗粒在 x = -0.45 和 x = 0.65 处几乎完全分离。图 8.35 中的条件速度曲线显示出对 x 的明显依赖性。事实上，对于 x = 0.8，条件速度在 xi = 0 时为负值，这与阻力只会使粒子单调减慢至零这一事实并不相符。尽管如此，u(xi) 的曲线与用于从方程 (1) 中找到它的矩 mj,k 一致。 （8.127）。事实上，重建的速度以方程的形式。 (8.107) 没有施加可实现性约束来限制 u(xi) 的最大/最小值。当斯托克斯数不小时，这可能会导致较大的速度值，即使该表达式再现了正确的速度标量矩。在实践中，由于较大的速度会导致较小的时间步长，

\begin{sourceblock}
\includegraphics[page=432,pagebox=mediabox,trim={58.00bp 394.05bp 58.23bp 59.00bp},clip,width=\textwidth,max height=0.72\textheight]{Marchisio2013.pdf}
\par\vspace{0.45\baselineskip}
\small 图8.32。具有小扩展的标量条件速度情况的矩，在时间 t = 1 时 τ0 = 10。
\end{sourceblock}
\begin{sourceblock}
\includegraphics[page=432,pagebox=mediabox,trim={58.00bp 65.50bp 58.23bp 345.80bp},clip,width=\textwidth,max height=0.72\textheight]{Marchisio2013.pdf}
\par\vspace{0.45\baselineskip}
\small 图8.33。小范围标量条件速度情况的统计数据，时间 t = 1 时 τ0 = 10。显示了五个空间位置的条件速度：x = -0.8、-0.4、0、0.4 和 0.8。
\end{sourceblock}
\begin{sourceblock}
\includegraphics[page=433,pagebox=mediabox,trim={58.00bp 364.17bp 58.23bp 66.54bp},clip,width=\textwidth,max height=0.72\textheight]{Marchisio2013.pdf}
\par\vspace{0.45\baselineskip}
\small 图8.34。具有小扩展的标量条件速度情况的矩，在时间 t = 3 时 τ0 = 1。
\end{sourceblock}
这种可实现性错误可能会对模拟产生负面影响。对于较小的斯托克斯数（即在模型的有效性范围内），u(xi)与流体速度的偏差较小，因此模拟时间步长由流体CFL数固定。 t = 1 时 τ0 = 0.1（即强阻力）的结果如图 8.36 和 8.37 所示。对于这种情况，由于强大的阻力，分布很快就会冻结。然而，从统计数据中我们仍然可以观察到许多与之前案例相同的趋势。请注意，这种情况在条件速度模型有效的斯托克斯数的预期范围内。因此，条件速度曲线与流体速度紧密耦合，因此图 8.37 中无法实现的速度非常小。

\paragraph*{价差大}
接下来我们看一下传播范围较大的案例。对于 t = 1 时 τ0 = 10（即非常弱的阻力），结果如图 8.38 和 8.39 所示。与小扩散的情况相比，我们现在从平均体积图中观察到，分离后剩余颗粒的平均体积（和体积标准偏差）接近于零。图 8.39 中 x = −0.4 处的相应条件速度与零速度存在较大偏差。这是因为 x = −0.4 处的 NDF 完全集中在接近零的 xi 值处，因此用于计算 u(xi) 的矩仅包含有关该相空间区域的信息。对于较大的 xi 值，曲线只需使用 u(xi) 的二次模型进行推断。无法实现速度的另一个明显迹象是图 8.39 中的统计数据变化超过 x = 0.8（即跨越周期边界）。由于最大初始速度为 1，因此实际系统中的信息不能以超过 1 的速度传播。作为

\begin{sourceblock}
\includegraphics[page=434,pagebox=mediabox,trim={58.00bp 291.16bp 58.23bp 59.00bp},clip,width=\textwidth,max height=0.72\textheight]{Marchisio2013.pdf}
\par\vspace{0.45\baselineskip}
\small 图8.35。小范围标量条件速度情况的统计数据，时间 t = 3 时 τ0 = 1。显示了五个空间位置的条件速度：x = -0.8、-0.4、0、0.4 和 0.8。
\end{sourceblock}
如上所述，出现这种无法实现的结果是因为条件速度模型不适用于大斯托克斯数。事实上，在更仔细地分析模拟时，发现无法实现的速度是在空间传输步骤中出现的（即由于方程（8.125）和（8.127））。理论上，有可能为方程中的系数开发一致的限制器。 （8.127）。然而，真正的困难在于，除非斯托克斯数很小，否则模型无法捕捉正确的物理现象。 t = 3 时 τ0 = 1（即中等阻力）的结果如图 8.40 和 8.41 所示。此处，由于 NDF 分布范围较大，尺寸偏析非常明显。同样，在模拟结束时，由于阻力的影响，NDF 几乎冻结在原地。然而，在模拟的初始步骤中，u(xi) 的广泛传播和无法实现的重建导致平流速度明显大于单位速度，从而导致相对较小的时间步长。斯托克斯数较小的情况如图 8.42 和 8.43 所示。同样，虽然分布很快就固定在适当的位置，但大小隔离非常显著。总体而言，在斯托克斯数较小的情况下，从时刻重构 u(xi) 时，速度可实现性是一个小问题。这表明，当在其有效性范围内使用时，

\begin{sourceblock}
\includegraphics[page=435,pagebox=mediabox,trim={58.00bp 394.05bp 58.23bp 59.00bp},clip,width=\textwidth,max height=0.72\textheight]{Marchisio2013.pdf}
\par\vspace{0.45\baselineskip}
\small 图8.36。具有小扩展的标量条件速度情况下的矩，在时间 t = 1 时 τ0 = 0.1。
\end{sourceblock}
\begin{sourceblock}
\includegraphics[page=435,pagebox=mediabox,trim={58.00bp 66.51bp 58.23bp 348.15bp},clip,width=\textwidth,max height=0.72\textheight]{Marchisio2013.pdf}
\par\vspace{0.45\baselineskip}
\small 图8.37。小范围标量条件速度情况的统计数据，时间 t = 1 时 τ0 = 0.1。显示了五个空间位置的条件速度：x = -0.8、-0.4、0、0.4 和 0.8。
\end{sourceblock}
\begin{sourceblock}
\includegraphics[page=436,pagebox=mediabox,trim={58.00bp 364.17bp 58.23bp 66.54bp},clip,width=\textwidth,max height=0.72\textheight]{Marchisio2013.pdf}
\par\vspace{0.45\baselineskip}
\small 图8.38。具有大扩展的标量条件速度情况的矩，在时间 t = 1 时 τ0 = 10。
\end{sourceblock}
标量条件速度是模拟多分散型多相流的一种数值高效选择，只需添加计算 u(xi) 所需的一小组速度矩即可。例如，分散的气泡流通常具有非常小的斯托克斯数，因此可以用这种方式处理。对于较大的斯托克斯数，更好的解决方案是使用下面描述的速度标量 NDF 方法。

\subsection*{8.5.6 速度标量 NDF 的空间传输}\phantomsection\addcontentsline{toc}{subsection}{8.5.6 速度标量 NDF 的空间传输}
作为最后一种情况，我们将再次考虑由方程 2 控制的联合速度标量 NDF。 （8.122），但没有指定标量条件速度的函数形式。矩输运方程再次由方程给出。 （8.122）。然而，我们现在将使用标量条件多元 EQMOM 来重建联合 NDF：

\begin{sourceblock}
\includegraphics[page=436,pagebox=mediabox,trim={162.66bp 117.73bp 65.08bp 550.17bp},clip,max width=0.94\textwidth,max height=0.38\textheight]{Marchisio2013.pdf}
\end{sourceblock}
对于标量，核密度函数将定义为

\begin{sourceblock}
\includegraphics[page=436,pagebox=mediabox,trim={158.88bp 50.16bp 65.08bp 616.32bp},clip,max width=0.94\textwidth,max height=0.38\textheight]{Marchisio2013.pdf}
\end{sourceblock}
\begin{sourceblock}
\includegraphics[page=437,pagebox=mediabox,trim={58.00bp 353.16bp 58.23bp 59.00bp},clip,width=\textwidth,max height=0.72\textheight]{Marchisio2013.pdf}
\par\vspace{0.45\baselineskip}
\small 图8.39。标量条件速度情况下的统计数据，在时间 t = 1 时，τ0 = 10，分布范围较大。显示了五个空间位置的条件速度：x = −0.8、−0.4、0、0.4 和 0.8。
\end{sourceblock}
对应于 gamma EQMOM。对于速度，将使用高斯核密度函数：

\begin{sourceblock}
\includegraphics[page=437,pagebox=mediabox,trim={159.07bp 234.20bp 67.57bp 428.79bp},clip,max width=0.94\textwidth,max height=0.38\textheight]{Marchisio2013.pdf}
\end{sourceblock}
对于 8.4 节中的双变量速度矩，必须传输 13 个双变量矩才能找到方程 1 中的参数。 (8.143):\{m0,0, m0,1, m0,2, m0,3, m0,4, m1,0, m1,1, m2,0, m2,1, m3,0, m3,1, m4,0, m4,1\}，其中矩由 mj,k = 定义

\begin{sourceblock}
\includegraphics[page=437,pagebox=mediabox,trim={331.88bp 198.24bp 155.46bp 489.96bp},clip,max width=0.94\textwidth,max height=0.38\textheight]{Marchisio2013.pdf}
\end{sourceblock}
vjΣkn(v, Σ)dv dΣ。权重 (wαβ)、节点 (xiα, vαβ) 和宽度参数 (σ1, σ1α) 是使用 ECQMOM 从 13 个二元矩确定的，如第 3.3.4 节中所述。为清楚起见，针对所考虑的案例，双变量 ECQMOM 所需的步骤如下。

(i) 使用 gamma EQMOM 和大小矩 (m0,0, m0,1, m0,2, m0,3, m0,4) 来找到

\begin{sourceblock}
\includegraphics[page=437,pagebox=mediabox,trim={89.93bp 117.09bp 247.39bp 569.25bp},clip,max width=0.94\textwidth,max height=0.38\textheight]{Marchisio2013.pdf}
\end{sourceblock}
(ii) 通过求解线性系统找到标量条件速度矩 μk α

\begin{sourceblock}
\includegraphics[page=437,pagebox=mediabox,trim={122.40bp 53.96bp 100.02bp 611.34bp},clip,max width=0.94\textwidth,max height=0.38\textheight]{Marchisio2013.pdf}
\end{sourceblock}
\begin{sourceblock}
\includegraphics[page=438,pagebox=mediabox,trim={58.00bp 372.16bp 58.23bp 66.43bp},clip,width=\textwidth,max height=0.72\textheight]{Marchisio2013.pdf}
\par\vspace{0.45\baselineskip}
\small 图8.40。具有小扩展的标量条件速度情况的矩，在时间 t = 3 时 τ0 = 1。
\end{sourceblock}
(iii) 对于每个 α = 1, 2，使用具有条件速度矩的高斯 EQMOM

\begin{sourceblock}
\includegraphics[page=438,pagebox=mediabox,trim={92.42bp 301.02bp 120.98bp 383.63bp},clip,max width=0.94\textwidth,max height=0.38\textheight]{Marchisio2013.pdf}
\end{sourceblock}
(iv) 最终的正交权重由下式给出： wαβ = wαw*

\begin{sourceblock}
\includegraphics[page=438,pagebox=mediabox,trim={324.78bp 269.79bp 155.79bp 416.20bp},clip,max width=0.94\textwidth,max height=0.38\textheight]{Marchisio2013.pdf}
\end{sourceblock}
正如对双变量速度矩所做的那样，我们将通过替换方程中的双和来简化符号。 (8.143) 通过索引 α = 1 的单个和，.。。 , 4. 然后我们将 ECQMOM 求积表示为 wα、Σα、vα、σ1α 和 σ2α。作为最后的求积步骤，我们为尺寸变量引入高斯-拉盖尔求积，为速度变量引入高斯-埃尔米特求积。令 mn j,k;i 表示时间步 n 处网格单元 i 中的体积平均二元矩，矩的双正交表示为

\begin{sourceblock}
\includegraphics[page=438,pagebox=mediabox,trim={158.47bp 138.70bp 156.02bp 528.62bp},clip,max width=0.94\textwidth,max height=0.38\textheight]{Marchisio2013.pdf}
\end{sourceblock}
如式中所示。根据方程（8.102），该表达式中的权重与节点来自（i）ECQ-MOM，（ii）高斯-拉盖尔求积，以及（iii）高斯-埃尔米特求积。同样，M1 和 M2 可以根据需要选择尽可能大的值，以便最小化空间通量和阻力项中的正交误差。

{\small 矩输运方程的数值求解再次通过使用算子分裂首先更新空间通量然后更新阻力来完成。可实现的方案\par}
\begin{sourceblock}
\includegraphics[page=439,pagebox=mediabox,trim={58.00bp 361.16bp 58.23bp 59.00bp},clip,width=\textwidth,max height=0.72\textheight]{Marchisio2013.pdf}
\par\vspace{0.45\baselineskip}
\small 图8.41。标量条件速度情况的统计数据，在时间 t = 3 时，τ0 = 1，分布范围较大。显示了五个空间位置的条件速度：x = −0.8、−0.4、0、0.4 和 0.8。
\end{sourceblock}
用于求解方程。 （8.125）与式（8.125）相同。 (8.89)（无扩散通量），但数值通量具有以下定义：

\begin{sourceblock}
\includegraphics[page=439,pagebox=mediabox,trim={98.92bp 235.93bp 67.57bp 419.11bp},clip,max width=0.94\textwidth,max height=0.38\textheight]{Marchisio2013.pdf}
\end{sourceblock}
权重Wn αβ1β2,i+1/2,l 和Wn αβ1β2,i+1/2,r 的公式与等式1之后描述的相同。 （8.103）。本案例的可实现条件为

\begin{sourceblock}
\includegraphics[page=439,pagebox=mediabox,trim={87.20bp 178.89bp 67.57bp 496.81bp},clip,max width=0.94\textwidth,max height=0.38\textheight]{Marchisio2013.pdf}
\end{sourceblock}
时间步长是固定的，如方程式所示。 (8.106) 使用

\begin{sourceblock}
\includegraphics[page=439,pagebox=mediabox,trim={195.18bp 120.21bp 67.57bp 541.90bp},clip,max width=0.94\textwidth,max height=0.38\textheight]{Marchisio2013.pdf}
\end{sourceblock}
对流步骤结束时的双变量矩用 m* j,k;i 表示，并且如果方程 1 成立，则保证可实现。 (8.147) 是满意的。在对流步骤之后，从 m* 开始构造一个新的双正交表示
{\small j，k；i。这个新的正交，参数为 W*\par}
\begin{sourceblock}
\includegraphics[page=439,pagebox=mediabox,trim={120.38bp 50.16bp 139.86bp 621.70bp},clip,max width=0.94\textwidth,max height=0.38\textheight]{Marchisio2013.pdf}
\end{sourceblock}
{\small αβ2,i 不变。\par}
\begin{sourceblock}
\includegraphics[page=440,pagebox=mediabox,trim={58.00bp 364.17bp 58.23bp 66.54bp},clip,width=\textwidth,max height=0.72\textheight]{Marchisio2013.pdf}
\par\vspace{0.45\baselineskip}
\small 图8.42。具有大扩展的标量条件速度情况的矩，在时间 t = 1 时 τ0 = 0.1。
\end{sourceblock}
相反，每个网格单元 i 中的速度节点是通过求解 ODE 找到的：

\begin{sourceblock}
\includegraphics[page=440,pagebox=mediabox,trim={166.50bp 260.79bp 163.35bp 398.89bp},clip,max width=0.94\textwidth,max height=0.38\textheight]{Marchisio2013.pdf}
\end{sourceblock}
初始条件为 V* αβ1,i。权重 w* αβ2,i 通过高斯-拉盖尔求积求得，并针对所有 α 和 i 值相对于 β2 求和至一。因此，方括号中的项是用于表示 n(xi) 的每个伽玛核密度函数的 1/ψ2 期望值。使用隐式欧拉方案，更新后的速度节点为

\begin{sourceblock}
\includegraphics[page=440,pagebox=mediabox,trim={158.13bp 150.56bp 65.08bp 509.13bp},clip,max width=0.94\textwidth,max height=0.38\textheight]{Marchisio2013.pdf}
\end{sourceblock}
与平流步骤中的 Δt 相同。最后，更新后的矩集由

\begin{sourceblock}
\includegraphics[page=440,pagebox=mediabox,trim={158.47bp 78.16bp 65.08bp 589.15bp},clip,max width=0.94\textwidth,max height=0.38\textheight]{Marchisio2013.pdf}
\end{sourceblock}
{\small 通过构造，这个矩集合保证是可实现的。\par}
\begin{sourceblock}
\includegraphics[page=441,pagebox=mediabox,trim={58.00bp 350.17bp 58.23bp 59.00bp},clip,width=\textwidth,max height=0.72\textheight]{Marchisio2013.pdf}
\par\vspace{0.45\baselineskip}
\small 图8.43。标量条件速度情况下的统计数据，在时间 t = 1 时具有较大分布，τ0 = 0.1。显示了五个空间位置的条件速度：x = -0.8、-0.4、0、0.4 和 0.8。
\end{sourceblock}
\section*{8.6 结束语}\phantomsection\addcontentsline{toc}{section}{8.6 结束语}
在本章中，我们描述了使用矩方法求解具有多个相空间变量的非齐次动力学方程所需的数值方法和算法。使用矩方法时的一个关键点是矩集的可实现性问题。我们已经看到，某些数值方案（例如矩的高阶空间重构）将不可避免地导致无法实现的矩。因此，我们强调可实现的有限体积方案，我们可以明确地证明矩集将在可实现性约束下实现。由于对可实现方案的开发的研究相对较新，读者可以预期在不久的将来会得到进一步的改进，因此本章中提出的一些方法可能会被更准确的方案所取代。使用本书中描述的矩方法时的第二个关键点是从可实现矩集开始重建 NDF 的算法。一般来说，解的精度很大程度上取决于 NDF 重建的质量，特别是对于其中数值通量与重建的 NDF 密切相关的非齐次系统。由于其灵活性、鲁棒性和计算效率，在本书中我们重点关注用于 NDF 重建的基于求积的矩方法。例如，对于单变量分布，通过增加第一个求积中的节点数量，通常可以使使用 EQMOM 找到的重构 NDF 变得足够准确。此外，我们已经看到，由于连续 NDF 和高斯求积节点的离散表示之间的对偶性，EQMOM 中的第二个求积在开发数值方案时具有很大的灵活性。然而，当使用多元 NDF 时，重建 NDF 的任务要困难得多。在本书中，我们重点讨论了现有的多元求积方法，但读者可以再次期待，随着对多元求积的研究兴趣的不断增加，将会带来进一步的改进。

\appendix
\titleformat{\chapter}[display]{\bfseries\Huge\centering}{附录\thechapter}{0.75em}{}
\chapter{矩反演算法}
本附录报告了第 3 章中讨论的矩反演算法的 Matlab 脚本。

\section*{A.1 单变量求积}\phantomsection\addcontentsline{toc}{section}{A.1 单变量求积}
\subsection*{A.1.1 PD算法}\phantomsection\addcontentsline{toc}{subsection}{A.1.1 PD算法}
下面报告了从 NDF 矩开始计算正交近似的节点 xi 和权重 w 的 Matlab 脚本。该脚本使用包含 N 个矩的向量 m 来计算矩阵 P、向量 zeta 和雅可比矩阵 jacobi 的系数。然后利用Matlab函数eig计算特征值和特征向量(eval, evect)，最终计算并返回节点和权重。

\begin{sourceblock}
\includegraphics[page=443,pagebox=mediabox,trim={58.00bp 52.85bp 58.23bp 413.18bp},clip,width=\textwidth,max height=0.66\textheight]{Marchisio2013.pdf}
\end{sourceblock}
\begin{sourceblock}
\includegraphics[page=444,pagebox=mediabox,trim={58.00bp 283.99bp 58.23bp 62.50bp},clip,width=\textwidth,max height=0.66\textheight]{Marchisio2013.pdf}
\end{sourceblock}
\subsection*{A.1.2 自适应Wheeler算法}\phantomsection\addcontentsline{toc}{subsection}{A.1.2 自适应Wheeler算法}
下面报告了使用 Wheeler 算法根据一组已知矩 m 计算 N 阶正交近似的 Matlab 脚本。该脚本计算中间系数 sigma 和雅可比矩阵，并且对于 PD 算法，根据矩阵的特征值和特征向量确定正交近似的节点和权重。

\begin{sourceblock}
\includegraphics[page=444,pagebox=mediabox,trim={58.00bp 52.84bp 58.23bp 524.76bp},clip,width=\textwidth,max height=0.66\textheight]{Marchisio2013.pdf}
\end{sourceblock}
\begin{sourceblock}
\includegraphics[page=445,pagebox=mediabox,trim={58.00bp 268.05bp 58.23bp 62.50bp},clip,width=\textwidth,max height=0.66\textheight]{Marchisio2013.pdf}
\end{sourceblock}
\section*{A.2 矩校正算法}\phantomsection\addcontentsline{toc}{section}{A.2 矩校正算法}
\subsection*{A.2.1 McGraw修正算法}\phantomsection\addcontentsline{toc}{subsection}{A.2.1 McGraw修正算法}
下面报告了一个 Matlab 脚本，其中包含用于校正不可实现的矩集的算法的实现。该脚本首先分析矩集，然后，如果该矩集不可实现，则脚本会识别出必须更改最少的矩才能使该集可实现。迭代该过程，直到矩集变得可实现或达到最大迭代次数。通过确保二阶差分的正性来验证可实现性；然而，可以更有效地使用其他更严格的条件。

\begin{sourceblock}
\includegraphics[page=445,pagebox=mediabox,trim={58.00bp 52.84bp 58.23bp 604.46bp},clip,width=\textwidth,max height=0.66\textheight]{Marchisio2013.pdf}
\end{sourceblock}
\begin{sourceblock}
\includegraphics[page=446,pagebox=mediabox,trim={58.00bp 475.27bp 58.23bp 62.50bp},clip,width=\textwidth,max height=0.66\textheight]{Marchisio2013.pdf}
\end{sourceblock}
\begin{sourceblock}
\includegraphics[page=446,pagebox=mediabox,trim={67.52bp 396.57bp 324.86bp 175.07bp},clip,max width=0.94\textwidth,max height=0.38\textheight]{Marchisio2013.pdf}
\end{sourceblock}
\begin{sourceblock}
\includegraphics[page=446,pagebox=mediabox,trim={58.00bp 411.51bp 58.23bp 237.83bp},clip,width=\textwidth,max height=0.66\textheight]{Marchisio2013.pdf}
\end{sourceblock}
\begin{sourceblock}
\includegraphics[page=446,pagebox=mediabox,trim={58.00bp 371.65bp 58.23bp 301.59bp},clip,width=\textwidth,max height=0.66\textheight]{Marchisio2013.pdf}
\end{sourceblock}
\begin{sourceblock}
\includegraphics[page=446,pagebox=mediabox,trim={58.00bp 108.64bp 58.23bp 341.43bp},clip,width=\textwidth,max height=0.66\textheight]{Marchisio2013.pdf}
\end{sourceblock}
\begin{sourceblock}
\includegraphics[page=446,pagebox=mediabox,trim={58.00bp 52.86bp 58.23bp 604.44bp},clip,width=\textwidth,max height=0.66\textheight]{Marchisio2013.pdf}
\end{sourceblock}
\begin{sourceblock}
\includegraphics[page=447,pagebox=mediabox,trim={65.02bp 627.69bp 422.54bp 63.50bp},clip,max width=0.94\textwidth,max height=0.38\textheight]{Marchisio2013.pdf}
\end{sourceblock}
\begin{sourceblock}
\includegraphics[page=447,pagebox=mediabox,trim={58.00bp 578.87bp 58.23bp 70.46bp},clip,width=\textwidth,max height=0.66\textheight]{Marchisio2013.pdf}
\end{sourceblock}
\subsection*{A.2.2 赖特校正算法}\phantomsection\addcontentsline{toc}{subsection}{A.2.2 赖特校正算法}
下面报告了一个实现 Wright 版本校正算法的简单 Matlab 脚本。

\begin{sourceblock}
\includegraphics[page=447,pagebox=mediabox,trim={58.00bp 52.85bp 58.23bp 205.96bp},clip,width=\textwidth,max height=0.66\textheight]{Marchisio2013.pdf}
\end{sourceblock}
\begin{sourceblock}
\includegraphics[page=448,pagebox=mediabox,trim={58.00bp 586.85bp 58.23bp 62.50bp},clip,width=\textwidth,max height=0.66\textheight]{Marchisio2013.pdf}
\end{sourceblock}
\section*{A.3 多元求积}\phantomsection\addcontentsline{toc}{section}{A.3 多元求积}
\subsection*{A.3.1 穷举式 QMOM}\phantomsection\addcontentsline{toc}{subsection}{A.3.1 穷举式 QMOM}
下面报告了一个针对二元问题 (M = 2) 实现强力 QMOM 的 Matlab 脚本。该脚本需要求积近似中使用的节点总数 N 作为输入。此外，由于理论上可以使用此方法，可以使用任何大小为 N(M + 1) 的矩集，因此必须指定为反演选择的集。这是通过指数矩阵 k 完成的，该矩阵具有 M = 2 列和 N(M + 1) 行，报告相对于所用矩阶数的两个内坐标的索引。然后，N(M + 1) 个矩通过在不同位置包含不同矩的向量 m，遵循指数矩阵 k 中指示的顺序。正如预期的那样，开始搜索需要一个良好的初步猜测。在矩阵变量 xi i 中，节点的初始猜测被组织成两个不同的行（对应于两个内部坐标）和 N 列（对应于不同的节点），而 w i 是包含权重的 N 个初始猜测的向量。该脚本计算求积并返回 N 个权重 w、N 个节点（再次组织在矩阵 xi 中）以及收敛到最终解所需的迭代次数 (iter)。该脚本采用迭代牛顿-拉夫森法来求解非线性系统，收敛准则要求整体归一化残差小于 tol，默认值为 1 × 10−6。为了帮助收敛，欠松弛因子 u rel 可以取比 1 小很多；典型值范围为 1 至 0.001。迭代循环在收敛时或达到最大迭代次数 iter max 时关闭。扩展到 M > 2 很简单，但收敛通常更难实现。

\begin{sourceblock}
\includegraphics[page=448,pagebox=mediabox,trim={58.00bp 52.85bp 58.23bp 453.02bp},clip,width=\textwidth,max height=0.66\textheight]{Marchisio2013.pdf}
\end{sourceblock}
\begin{sourceblock}
\includegraphics[page=449,pagebox=mediabox,trim={58.00bp 355.72bp 58.23bp 62.50bp},clip,width=\textwidth,max height=0.66\textheight]{Marchisio2013.pdf}
\end{sourceblock}
\begin{sourceblock}
\includegraphics[page=449,pagebox=mediabox,trim={58.00bp 68.80bp 58.23bp 357.38bp},clip,width=\textwidth,max height=0.66\textheight]{Marchisio2013.pdf}
\end{sourceblock}
\begin{sourceblock}
\includegraphics[page=449,pagebox=mediabox,trim={79.67bp 53.86bp 309.01bp 629.36bp},clip,max width=0.94\textwidth,max height=0.38\textheight]{Marchisio2013.pdf}
\end{sourceblock}
\begin{sourceblock}
\includegraphics[page=450,pagebox=mediabox,trim={82.16bp 627.69bp 394.42bp 63.50bp},clip,max width=0.94\textwidth,max height=0.38\textheight]{Marchisio2013.pdf}
\end{sourceblock}
\begin{sourceblock}
\includegraphics[page=450,pagebox=mediabox,trim={58.00bp 483.23bp 58.23bp 70.46bp},clip,width=\textwidth,max height=0.66\textheight]{Marchisio2013.pdf}
\end{sourceblock}
\begin{sourceblock}
\includegraphics[page=450,pagebox=mediabox,trim={96.80bp 468.30bp 324.84bp 214.92bp},clip,max width=0.94\textwidth,max height=0.38\textheight]{Marchisio2013.pdf}
\end{sourceblock}
\begin{sourceblock}
\includegraphics[page=450,pagebox=mediabox,trim={58.00bp 387.59bp 58.23bp 229.86bp},clip,width=\textwidth,max height=0.66\textheight]{Marchisio2013.pdf}
\end{sourceblock}
\subsection*{A.3.2 张量积 QMOM}\phantomsection\addcontentsline{toc}{subsection}{A.3.2 张量积 QMOM}
下面报告了一个 Matlab 脚本，该脚本针对本节中描述的简单双变量情况实现了张量积 QMOM。所需的输入是第一个 (N1) 和第二个 (N2) 内坐标的节点数。由于在上述公式中，用于计算正交近似的矩是由方法本身定义的，因此不需要指数矩阵。使用的矩通过矩阵变量 m 传递，其元素由两个索引定义。第一个表示矩相对于第一个内坐标的顺序（索引 1 表示矩 0，索引 2 表示矩顺序 1，等等），而第二个表示矩相对于第二个内坐标的顺序。最终矩阵与表 3.8 中报告的非常相似。该脚本以通常的形式返回正交近似值：权重存储在大小为 N = N1N2 的权重向量 w 中，而节点存储在具有两行（对应于第一个和第二个内部坐标）和 N = N1N2 列（对应于不同节点）的矩阵中。

\begin{sourceblock}
\includegraphics[page=450,pagebox=mediabox,trim={58.00bp 51.38bp 58.23bp 526.23bp},clip,width=\textwidth,max height=0.66\textheight]{Marchisio2013.pdf}
\end{sourceblock}
\begin{sourceblock}
\includegraphics[page=451,pagebox=mediabox,trim={58.00bp 52.86bp 58.23bp 62.50bp},clip,width=\textwidth,max height=0.66\textheight]{Marchisio2013.pdf}
\end{sourceblock}
\subsection*{A.3.3 CQMOM}\phantomsection\addcontentsline{toc}{subsection}{A.3.3 CQMOM}
下面报告了用于确定双变量情况的 CQMOM 近似值的最小 Matlab 脚本（Yuan \& Fox，2011）。该脚本需要指定第一个 (N1) 和第二个 (N2) 内部坐标所需的节点数。计算中使用的矩必须以矩阵形式提供。包含矩 m 的矩阵由两个索引定义；第一个表示矩相对于第一个内部坐标的顺序（索引 1 表示矩 0，索引 2 表示矩顺序 1，等等），而第二个表示矩相对于第二个内部坐标的顺序。数据结构类似于表 3.9-3.11。该脚本计算 N = N1N2 权重与节点，并将它们存储在向量 w 和矩阵 xi 中，它们的结构与前面的脚本相同。该过程基于相对于第二内坐标的矩的计算，以第一内坐标的值为条件。通过为脚本提供矩矩阵的转置，可以简单地执行以第二内坐标为条件的矩构造的求积计算。

\begin{sourceblock}
\includegraphics[page=452,pagebox=mediabox,trim={58.00bp 116.61bp 58.23bp 269.71bp},clip,width=\textwidth,max height=0.66\textheight]{Marchisio2013.pdf}
\end{sourceblock}
\begin{sourceblock}
\includegraphics[page=452,pagebox=mediabox,trim={82.16bp 101.68bp 357.81bp 581.55bp},clip,max width=0.94\textwidth,max height=0.38\textheight]{Marchisio2013.pdf}
\end{sourceblock}
\begin{sourceblock}
\includegraphics[page=452,pagebox=mediabox,trim={58.00bp 52.85bp 58.23bp 596.48bp},clip,width=\textwidth,max height=0.66\textheight]{Marchisio2013.pdf}
\end{sourceblock}
\begin{sourceblock}
\includegraphics[page=453,pagebox=mediabox,trim={79.67bp 627.69bp 374.94bp 63.50bp},clip,max width=0.94\textwidth,max height=0.38\textheight]{Marchisio2013.pdf}
\end{sourceblock}
\begin{sourceblock}
\includegraphics[page=453,pagebox=mediabox,trim={58.00bp 363.69bp 58.23bp 70.46bp},clip,width=\textwidth,max height=0.66\textheight]{Marchisio2013.pdf}
\end{sourceblock}
\section*{A.4 EQMOM}\phantomsection\addcontentsline{toc}{section}{A.4 EQMOM}
在本节中，我们提供了计算扩展正交的示例代码，如第 3.3.2 节中所述。这些代码旨在说明寻找 σ 所需的迭代方法的一般结构，并且无疑可以改进。

\subsection*{A.4.1 Beta EQMOM}\phantomsection\addcontentsline{toc}{subsection}{A.4.1 Beta EQMOM}
\begin{sourceblock}
\includegraphics[page=453,pagebox=mediabox,trim={58.00bp 52.85bp 58.23bp 461.00bp},clip,width=\textwidth,max height=0.66\textheight]{Marchisio2013.pdf}
\end{sourceblock}
\begin{sourceblock}
\includegraphics[page=454,pagebox=mediabox,trim={58.00bp 52.86bp 58.23bp 62.50bp},clip,width=\textwidth,max height=0.66\textheight]{Marchisio2013.pdf}
\end{sourceblock}
\begin{sourceblock}
\includegraphics[page=455,pagebox=mediabox,trim={58.00bp 52.86bp 58.23bp 62.50bp},clip,width=\textwidth,max height=0.66\textheight]{Marchisio2013.pdf}
\end{sourceblock}
\begin{sourceblock}
\includegraphics[page=456,pagebox=mediabox,trim={58.00bp 395.57bp 58.23bp 62.50bp},clip,width=\textwidth,max height=0.66\textheight]{Marchisio2013.pdf}
\end{sourceblock}
\subsection*{A.4.2 伽马EQMOM}\phantomsection\addcontentsline{toc}{subsection}{A.4.2 伽马EQMOM}
\begin{sourceblock}
\includegraphics[page=456,pagebox=mediabox,trim={58.00bp 60.82bp 58.23bp 341.44bp},clip,width=\textwidth,max height=0.66\textheight]{Marchisio2013.pdf}
\end{sourceblock}
\begin{sourceblock}
\includegraphics[page=456,pagebox=mediabox,trim={82.16bp 53.85bp 204.03bp 637.34bp},clip,max width=0.94\textwidth,max height=0.38\textheight]{Marchisio2013.pdf}
\end{sourceblock}
\begin{sourceblock}
\includegraphics[page=457,pagebox=mediabox,trim={65.02bp 627.69bp 367.61bp 63.50bp},clip,max width=0.94\textwidth,max height=0.38\textheight]{Marchisio2013.pdf}
\end{sourceblock}
\begin{sourceblock}
\includegraphics[page=457,pagebox=mediabox,trim={58.00bp 52.86bp 58.23bp 70.46bp},clip,width=\textwidth,max height=0.66\textheight]{Marchisio2013.pdf}
\end{sourceblock}
\begin{sourceblock}
\includegraphics[page=458,pagebox=mediabox,trim={58.00bp 124.59bp 58.23bp 62.50bp},clip,width=\textwidth,max height=0.66\textheight]{Marchisio2013.pdf}
\end{sourceblock}
\subsection*{A.4.3 高斯 EQMOM}\phantomsection\addcontentsline{toc}{subsection}{A.4.3 高斯 EQMOM}
\begin{sourceblock}
\includegraphics[page=458,pagebox=mediabox,trim={58.00bp 52.84bp 58.23bp 612.43bp},clip,width=\textwidth,max height=0.66\textheight]{Marchisio2013.pdf}
\end{sourceblock}
\begin{sourceblock}
\includegraphics[page=459,pagebox=mediabox,trim={58.00bp 52.86bp 58.23bp 62.50bp},clip,width=\textwidth,max height=0.66\textheight]{Marchisio2013.pdf}
\end{sourceblock}
\begin{sourceblock}
\includegraphics[page=460,pagebox=mediabox,trim={58.00bp 451.36bp 58.23bp 62.50bp},clip,width=\textwidth,max height=0.66\textheight]{Marchisio2013.pdf}
\end{sourceblock}
\chapter{基于动理学的有限体积方法}
在本附录中，我们讨论使用基于动理学的有限体积方法 (KBFVM) 在多个空间维度上从 GPBE 中找到的矩输运方程的解。与第 8.2 节一样，讨论的重点是使用单阶段欧拉显式时间积分方案的时间推进。对有限体积方法的更多细节感兴趣的读者可以参考Leveque (2002)。

\section*{B.1 GPBE 的空间依赖性}\phantomsection\addcontentsline{toc}{section}{B.1 GPBE 的空间依赖性}
下面，我们将 NDF 表示为 n(t, x, v, Ψ)，其中速度内坐标为 v，其余内坐标为 Ψ。要考虑的 GPBE 的一般形式是

\begin{sourceblock}
\includegraphics[page=461,pagebox=mediabox,trim={126.23bp 256.37bp 67.55bp 416.91bp},clip,max width=0.94\textwidth,max height=0.38\textheight]{Marchisio2013.pdf}
\end{sourceblock}
其中左侧第二项和第三项分别是具有已知速度 u(t, x, xi) 和自由传输的多维平流，右侧最后一项 S (t, x, v, xi) 是相空间源项。涉及非负二阶张量 λ(t, x, v, xi) 的项是混合平流项 1，描述速度相空间中的梯度如何导致物理空间中的通量。该术语出现在气体-颗粒流的湍流模型中，用于描述未解决的气体速度波动对颗粒速度分布函数的影响。涉及非负二阶张量 Г(t, x, v, xi) 的项对物理空间中的多维扩散进行建模，例如由于布朗运动。为了方便起见，我们在方程中包含了所有可能的项。 (B.1)，但在应用中几乎总是出现仅出现术语子集的情况。

1 混合平流通常写成以下形式(Reeks, 1983, 1991, 1992; Zaichik, 1997; Zaichik et al., 2008)

\begin{sourceblock}
\includegraphics[page=461,pagebox=mediabox,trim={126.88bp 78.40bp 128.22bp 594.78bp},clip,max width=0.94\textwidth,max height=0.38\textheight]{Marchisio2013.pdf}
\end{sourceblock}
{\small 其中 μ 是非负二阶张量。该表达式右侧的第二个和第三个贡献可以添加到相空间源项 S，因此等式： (B.1) 是等效形式。\par}
例如，在 PBE 中，速度内坐标不存在，因此只会出现平流、扩散和源项。同样，对于 KE，仅会出现自由传输、混合平流和源项。为了清楚和简单起见，我们讨论每个方向上具有均匀间距的 3D 结构化笛卡尔网格的有限体积方法 (FVM)：(Δx、Δy、Δz)。此外，我们仅考虑各向同性扩散和各向同性混合平流，因此 Г(t, x, v, Ψ) = Г(t, x, v, Ψ)δij 且 λ(t, x, v, Ψ) = λ(t, x, v, Ψ)δij。 Vikas 等人描述了可实现的高阶 FVM 到非结构化网格的扩展。 （2011a）。这里我们的重点是确保在采用高阶空间重建时传输矩集的可实现性，因此结构化网格的更简单的公式就足够了。此外，通过使用运算符拆分方程中的每一项。 (B.1) 可以在时间上单独推进，因此足以单独描述每一项的可实现 FVM。因此，在 B.3 节中，我们描述了多维平流的可实现 FVM：

\begin{sourceblock}
\includegraphics[page=462,pagebox=mediabox,trim={214.92bp 451.50bp 65.07bp 222.04bp},clip,max width=0.94\textwidth,max height=0.38\textheight]{Marchisio2013.pdf}
\end{sourceblock}
对于每个方向 i = 1, 2, 3。在 B.4 节中，我们描述了自由传输的可实现 FVM：

\begin{sourceblock}
\includegraphics[page=462,pagebox=mediabox,trim={213.13bp 397.05bp 65.07bp 276.48bp},clip,max width=0.94\textwidth,max height=0.38\textheight]{Marchisio2013.pdf}
\end{sourceblock}
在 B.5 节中，我们对混合平流做了同样的事情：

\begin{sourceblock}
\includegraphics[page=462,pagebox=mediabox,trim={210.28bp 349.08bp 65.07bp 323.08bp},clip,max width=0.94\textwidth,max height=0.38\textheight]{Marchisio2013.pdf}
\end{sourceblock}
在 B.6 节中，我们这样做是为了扩散：

\begin{sourceblock}
\includegraphics[page=462,pagebox=mediabox,trim={211.53bp 298.53bp 65.07bp 375.00bp},clip,max width=0.94\textwidth,max height=0.38\textheight]{Marchisio2013.pdf}
\end{sourceblock}
读者应该注意，在等式中。 (B.2)–(B.5) 空间导数出现在右侧，因此有必要为每种情况定义一个可实现的高阶 FVM。相反，方程中的源项 S。 (B.1) 不包含空间导数2，因此在每个有限体积网格单元中是局部的。换句话说，通过运算符拆分源项，可以为每个网格单元生成一个（刚性）常微分方程 (ODE)。如第 7 章所述，可以使用 QBMM 处理源项，其中权重与节点是根据多元矩计算的。因此，有必要设计一个可实现的 ODE 求解器，以与等式 1 中的源项一起使用。 (B.1)。在本附录的其余部分中，我们假设存在这样的 ODE 求解器，因此我们将注意力集中在空间传输求解器上。尽管如此，读者应该查阅 Yuan \& Fox (2011)，了解如何使用具有多种排列的 CQMOM 以统计一致的方式评估源项的详细信息。简而言之，当使用 CQMOM 时
2 源项可以包含内坐标的导数。例如，速度相空间中的扩散可写为

\begin{sourceblock}
\includegraphics[page=462,pagebox=mediabox,trim={213.92bp 85.60bp 210.28bp 590.52bp},clip,max width=0.94\textwidth,max height=0.38\textheight]{Marchisio2013.pdf}
\end{sourceblock}
{\small 其中 B(t, x, v) 是非负扩散张量。由于该项不涉及空间导数，因此包含在源项 S 中。\par}
如第 3.2.3 节所述，在定义条件 PDF 时，用户可以选择哪个内部坐标排在第一位、第二位等，我们将每个唯一的选择称为 CQMOM 的排列。除非 CQMOM 的所有排列都能准确再现基础 NDF，否则权重与节点的结果集将会不同。尽管如此，每个排列都会生成 NDF 的有效重建！因此，如果我们假设有 P 个排列，那么每个排列的源项都会以不同的方式闭合。让 Sp 表示 CQMOM 第 p 次排列的源项闭包。由于源项 Sp 引起的传递矩集 M 的变化可以表示为

\begin{sourceblock}
\includegraphics[page=463,pagebox=mediabox,trim={184.27bp 498.74bp 67.55bp 164.84bp},clip,max width=0.94\textwidth,max height=0.38\textheight]{Marchisio2013.pdf}
\end{sourceblock}
其中 K 是对应于传输矩集的矩向量（示例见方程（B.9））。换句话说，在定义（齐次）矩方程的闭包时，所有 P 排列都被同等对待。为了完整起见，我们应该注意到，如果可以应用第 3.2.1 节中描述的强力 QMOM，则 P = 1，因为最大可能的矩数由正交重建控制。不幸的是，强力 QMOM 不能保证收敛，因此不适合与 FVM 一起使用。

\section*{B.2 可实现的 FVM}\phantomsection\addcontentsline{toc}{section}{B.2 可实现的 FVM}
在结构化笛卡尔网格上，以点 (i Δx, j Δy, k Δz) 为中心的网格单元可以用 Ωijk 表示，其体积用 |Ωijk| 表示。 = Δx Δy Δz。时间 t = m Δt,3 时 Ωijk 中的体积平均 NDF 表示为 nm ijk(v, xi)，定义为 (Leveque, 2002)

\begin{sourceblock}
\includegraphics[page=463,pagebox=mediabox,trim={174.69bp 289.58bp 67.55bp 376.11bp},clip,max width=0.94\textwidth,max height=0.38\textheight]{Marchisio2013.pdf}
\end{sourceblock}
传递矩集 M(t, x) 与 NDF 的关系为

\begin{sourceblock}
\includegraphics[page=463,pagebox=mediabox,trim={197.99bp 233.34bp 67.55bp 432.47bp},clip,max width=0.94\textwidth,max height=0.38\textheight]{Marchisio2013.pdf}
\end{sourceblock}
其中 K 是由下式定义的矩向量

\begin{sourceblock}
\includegraphics[page=463,pagebox=mediabox,trim={120.70bp 183.61bp 67.55bp 491.38bp},clip,max width=0.94\textwidth,max height=0.38\textheight]{Marchisio2013.pdf}
\end{sourceblock}
K 的长度由多元传递矩集决定。乘法方程式(B.7) 通过 K 并在内坐标上积分得出体积平均传递矩集：

\begin{sourceblock}
\includegraphics[page=463,pagebox=mediabox,trim={144.11bp 105.02bp 67.56bp 560.66bp},clip,max width=0.94\textwidth,max height=0.38\textheight]{Marchisio2013.pdf}
\end{sourceblock}
\footnotetext[3]{仅当时间步长恒定时，写 t = m Δt 才是正确的。一般来说，情况不会如此。然而，我们会在所有情况下通过写 t = m Δt 来滥用这种表示法，并隐含地理解，当 Δt 在每个时间步长上变化时，t 的公式会更加复杂。}
一般来说，体积平均传递矩集合与体积平均NDF相关。因此，为了证明Mm ijk 是可实现矩集，只需证明nm ijk 是非负的即可。这种使用 NDF 来展示矩集的可实现性的想法是基于动力学的 FVM 的基础。 ijk，nm 的值 给定 n0 ijk（0 < m），由方程 (1) 的有限体积解确定。 (B.1)。因此，作为第一步，我们需要考虑方程的体积平均形式。 (B.1) 或等价地，方程中各项的体积平均形式。 (B.2)–(B.5)。使用单阶段欧拉显式时间积分方案（Leveque，2002），更新后的 NDF 的有限体积表达式具有形式4

\begin{sourceblock}
\includegraphics[page=464,pagebox=mediabox,trim={116.06bp 446.31bp 65.07bp 174.69bp},clip,max width=0.94\textwidth,max height=0.38\textheight]{Marchisio2013.pdf}
\end{sourceblock}
其中 G(nl, nr) 是必须为每个空间传输项构建的数值通量函数。方程中需要观察的重要一点。 (B.11) 是数值通量函数中出现的 NDF 对应于单元面。例如， nm i+1/2 jk,l 是位于 (i + 1/2)Δx 处的单元面左侧 5 的 NDF，nm i−1/2 jk,r 是位于 (i -1/2)Δx 处的单元面右侧的 NDF。在 FVM 中，仅知道体积平均值 nm ijk，因此使用空间重建来查找面值。例如，对于一阶空间重建，我们将单元中的所有值设置为 nm ijk。另一方面，对于二阶空间重建，我们使用线性近似：

\begin{sourceblock}
\includegraphics[page=464,pagebox=mediabox,trim={124.14bp 303.59bp 65.07bp 372.89bp},clip,max width=0.94\textwidth,max height=0.38\textheight]{Marchisio2013.pdf}
\end{sourceblock}
其中 xijk 表示像元中心，斜率向量 am ijk 是使用通量限制器根据相邻像元中 nm ijk 的值计算的，以确保 nm(x, v, xi) 对于 x ∈Ωijk 始终为非负。在可实现的高阶 FVM 中，数值通量函数的定义方式必须保证 nm+1 ijk (v, xi) 由方程 1 计算得出。 (B.11) 总是非负的。一般来说，通过一阶空间重建，可以直接构造可实现的数值通量函数。然而，对于高阶空间重建，情况通常并非如此。在 QBMM 中，体积平均 NDF nm ijk 是使用矩反演算法从体积平均传输矩集 Mm ijk 重建的（详细信息请参阅第 3 章）。生成的 NDF 具有形式 6

\begin{sourceblock}
\includegraphics[page=464,pagebox=mediabox,trim={120.88bp 131.23bp 65.07bp 535.71bp},clip,max width=0.94\textwidth,max height=0.38\textheight]{Marchisio2013.pdf}
\end{sourceblock}
4 Vikas 等人提出了对非结构化网格的扩展。 （2011a）。对于结构化和非结构化网格来说，关键点在于，考虑到网格单元中的体积平均 NDF，有必要在单元面上重建 NDF。

\footnotetext[5]{我们使用与一维中相同的约定来定义左和右：xi 的值增大对应于右，值减小对应于左。}
\footnotetext[6]{即使当多元 EQMOM 与核密度函数一起使用时，也会采用 Dirac delta 函数（双求积）表示来闭合 GPBE 中的项。更多详细信息，请参见第 3.3.4 节。}
\begin{sourceblock}
\includegraphics[page=465,pagebox=mediabox,trim={216.59bp 624.19bp 245.90bp 60.00bp},clip,max width=0.94\textwidth,max height=0.38\textheight]{Marchisio2013.pdf}
\end{sourceblock}
其中 wm ijk,α 是权重，(vm ijk,α) 是对应于网格单元 Ωijk 中体积平均矩的节点。此后，为了避免混淆，我们将将此步骤称为正交重建。正交重建在每个网格单元中单独完成，而空间重建使用来自（可能）所有网格单元的信息。如 Vikas 等人所示。 （2011a），如果通过使用通过体积平均矩的空间重建找到的单元面矩的正交重建来找到单元面的NDF，则不能保证单元面的NDF是可实现的。同样，如果权重与节点-sas 的空间重建与方程一起使用。 (B.13) 来近似单元面处的 NDF，那么当使用高阶空间重建时，后者将无法实现（示例参见第 8.2.1 节）。如果 CQMOM 用于多元正交重建，则另一个需要记住的重要技术点是方程： (B.13) 对应于 CQMOM 的一种排列。换句话说，CQMOM 的每个 P 排列都会生成 nm ijk(v, xi) 的不同表示。然而，下面给出的从 Mm ijk 开始计算 nm+1 ijk (v, xi) 的公式可以应用于每个排列，以便在一个时间步结束时，每个排列将产生一个不同的更新矩集，用 (Mm+1 ijk )p 表示。然后我们可以将排列平均更新矩集定义为

\begin{sourceblock}
\includegraphics[page=465,pagebox=mediabox,trim={195.07bp 370.51bp 67.56bp 296.44bp},clip,max width=0.94\textwidth,max height=0.38\textheight]{Marchisio2013.pdf}
\end{sourceblock}
这是下一次迭代的起点。正如我们将在 B.4 节中看到的，对于速度内坐标，CQMOM 排列的使用略有不同。尽管如此，使用所有有效 CQMOM 排列的总体思路保持不变。继维卡斯等人之后。 (2011a)中，为了保证NDF的空间重建是可实现的，我们将对求积节点采用一阶空间重建，对求积权重采用高阶空间重建。形式上，这可以写成

\begin{sourceblock}
\includegraphics[page=465,pagebox=mediabox,trim={101.19bp 214.67bp 67.56bp 452.26bp},clip,max width=0.94\textwidth,max height=0.38\textheight]{Marchisio2013.pdf}
\end{sourceblock}
其中权重 0 ≤wm α(x) 是使用所有网格单元中 wm ijk,α 的已知值通过通量限制器重建的。例如，对于一阶空间重建，网格单元 Ωijk 内的 NDF 表示为

\begin{sourceblock}
\includegraphics[page=465,pagebox=mediabox,trim={74.31bp 106.74bp 67.56bp 550.62bp},clip,max width=0.94\textwidth,max height=0.38\textheight]{Marchisio2013.pdf}
\end{sourceblock}
其中斜率向量 am ijk 是使用通量限制器根据相邻单元中的 wm ijk 值计算的，以确保 nm(x, v, xi) 始终为非负。请注意，等式。 (B.15) 和 (B.16) 对于任意单元格均有效。对于笛卡尔单元的特殊情况，12 NDF 位于方程右侧出现的单元面。 (B.11) 成为

\begin{sourceblock}
\includegraphics[page=466,pagebox=mediabox,trim={93.19bp 231.99bp 65.07bp 58.62bp},clip,max width=0.94\textwidth,max height=0.38\textheight]{Marchisio2013.pdf}
\end{sourceblock}
其中斜率向量的分量写为a = [a1 a2 a3]T。通过构造，方程中出现的所有权重。 (B.17) 是非负的，因此单元面处的重构 NDF 都是可实现的。总之，可实现的 FVM 求解每个时间步和网格单元的体积平均传输矩集 Mm ijk，并使用每个时间步的这些矩集来重建单元面上的 NDF：

\begin{sourceblock}
\includegraphics[page=466,pagebox=mediabox,trim={121.52bp 133.58bp 98.26bp 542.83bp},clip,max width=0.94\textwidth,max height=0.38\textheight]{Marchisio2013.pdf}
\end{sourceblock}
ijk

在下面的部分中，我们将描述如何使用单元表面处的重建 NDF 来为方程 1 中的每个空间传输项定义数值通量函数 G(nl, nr)。 (B.2)–(B.5) 保证从等式中找到更新的 NDF nm+1 ijk。 (B.11)（因此体积平均传递矩集 Mm+1
{\small ijk ) 是可以实现的。使用 RK2SSP 将可实现的 FVM 扩展到时间二阶（Vikas 等人，2011a）在第 8.2 节中进行了描述。\par}
\section*{B.3 平流}\phantomsection\addcontentsline{toc}{section}{B.3 平流}
为了清楚起见，我们将描述方程式中的平流项。 (B.2) 在每个方向上使用算子分裂。形式上，该过程可以表示为

\begin{sourceblock}
\includegraphics[page=467,pagebox=mediabox,trim={164.55bp 556.46bp 146.26bp 124.20bp},clip,max width=0.94\textwidth,max height=0.38\textheight]{Marchisio2013.pdf}
\end{sourceblock}
ijk

我们使用单个全局时间步长 Δt 连续解决每个方向的平流问题。以 x 方向与 xi 无关的速度为例，NDF 和传递矩集的平流方程为

\begin{sourceblock}
\includegraphics[page=467,pagebox=mediabox,trim={157.35bp 469.32bp 67.56bp 205.02bp},clip,max width=0.94\textwidth,max height=0.38\textheight]{Marchisio2013.pdf}
\end{sourceblock}
由于矩输运方程是封闭的，人们可能会尝试直接为 M 设计一个高阶方案。然而，对于这样的方案很难确保可实现性，因此更好的方法是直接使用 NDF。此外，当平流速度取决于 ε 时，矩输运方程不是封闭的，因此使用 NDF 将得到更通用的公式。现在考虑已知平流速度 u(t, x, xi) 取决于 xi 的情况。7 对于这种情况，矩输运方程不是封闭的，但 NDF 传递方程仍然由方程 1 给出。 （B.18）。对于 x 方向，NDF 的有限体积公式为

\begin{sourceblock}
\includegraphics[page=467,pagebox=mediabox,trim={118.48bp 331.86bp 67.56bp 344.22bp},clip,max width=0.94\textwidth,max height=0.38\textheight]{Marchisio2013.pdf}
\end{sourceblock}
其中 λx = Δt/Δx，数值通量函数定义为

\begin{sourceblock}
\includegraphics[page=467,pagebox=mediabox,trim={194.89bp 278.10bp 67.56bp 405.89bp},clip,max width=0.94\textwidth,max height=0.38\textheight]{Marchisio2013.pdf}
\end{sourceblock}
其中

\begin{sourceblock}
\includegraphics[page=467,pagebox=mediabox,trim={94.75bp 216.64bp 67.56bp 457.52bp},clip,max width=0.94\textwidth,max height=0.38\textheight]{Marchisio2013.pdf}
\end{sourceblock}
请注意，如果所有 x 的速度 u(t, x, ψ) 已知，则 u1,l 和 u1,r 对应于在固定 ε 值的单元面上计算的速度分量 u1(t, x, ψ)。否则，u1,l 和 u1,r 必须根据相应网格单元中速度的体积平均值重建。在整个附录中，我们表示 x 相关变量，例如 u1(t, x, xi)，在单元面（时间步 m）处通过 um 1,i+1/2jk 和 um 1,i−1/2jk 进行评估。附加下标 l 或 r 用于指示细胞表面的哪一侧是感兴趣的。

7 如第 5.2 节所述，对 xi 的典型依赖关系可写为

\begin{sourceblock}
\includegraphics[page=467,pagebox=mediabox,trim={191.95bp 82.68bp 194.48bp 604.72bp},clip,max width=0.94\textwidth,max height=0.38\textheight]{Marchisio2013.pdf}
\end{sourceblock}
{\small 其中 uf、u* 和 g 是已知函数。例如，在流体-粒子系统中，uf 可以是流体速度，g 可以是斯托克斯数，因此 u* 表示对非常小的斯托克斯数有效的速度修正。\par}
使用方程式重建的 NDF。 (B.15)，方程中更新的 NDF。 (B.19) 变为8

\begin{sourceblock}
\includegraphics[page=468,pagebox=mediabox,trim={78.45bp 516.83bp 65.07bp 81.69bp},clip,max width=0.94\textwidth,max height=0.38\textheight]{Marchisio2013.pdf}
\end{sourceblock}
方程中的最后两项。 (B.22) 是非负的，因此 n(1) ijk 将可实现，如果

\begin{sourceblock}
\includegraphics[page=468,pagebox=mediabox,trim={104.64bp 468.68bp 102.86bp 214.91bp},clip,max width=0.94\textwidth,max height=0.38\textheight]{Marchisio2013.pdf}
\end{sourceblock}
对于 i jk 和 α 的所有值，产生 Δt 的第一个可实现性条件：

\begin{sourceblock}
\includegraphics[page=468,pagebox=mediabox,trim={86.43bp 410.34bp 65.07bp 252.19bp},clip,max width=0.94\textwidth,max height=0.38\textheight]{Marchisio2013.pdf}
\end{sourceblock}
请注意，如果 wm ijk,α = 0，则 wm i+1/2jk,α,l = wm i−1/2jk,α,r = 0，并且对于所有 0 ≤Δt 都满足可实现性条件。如果方程中的可实现性条件。满足(B.23)，则由式(B.23)可求出矩集M(1)ijk。 (B.22) 使用定义 9

\begin{sourceblock}
\includegraphics[page=468,pagebox=mediabox,trim={205.75bp 317.29bp 65.07bp 348.52bp},clip,max width=0.94\textwidth,max height=0.38\textheight]{Marchisio2013.pdf}
\end{sourceblock}
其中（如方程（B.9））K 是由传输矩集确定的矩矢量。使用 M(1) ijk 进行正交重建，然后生成权重 w(1) ijk,α 和节点 xi(1)

\begin{sourceblock}
\includegraphics[page=468,pagebox=mediabox,trim={409.17bp 282.55bp 65.07bp 403.00bp},clip,max width=0.94\textwidth,max height=0.38\textheight]{Marchisio2013.pdf}
\end{sourceblock}
由此可以计算 NDF n(2) ijk：

\begin{sourceblock}
\includegraphics[page=468,pagebox=mediabox,trim={78.70bp 157.26bp 65.07bp 441.25bp},clip,max width=0.94\textwidth,max height=0.38\textheight]{Marchisio2013.pdf}
\end{sourceblock}
因此 Δt 的第二个可实现条件是

\begin{sourceblock}
\includegraphics[page=468,pagebox=mediabox,trim={86.44bp 94.32bp 65.07bp 565.11bp},clip,max width=0.94\textwidth,max height=0.38\textheight]{Marchisio2013.pdf}
\end{sourceblock}
\footnotetext[8]{对于平流，不需要速度内坐标，因此我们用 Σ 表示所有内坐标。 9由于n(1) i jk(xi)涉及加权狄拉克δ函数，矩集M(1)}
{\small i jk 涉及求和并且可以分析计算。\par}
如果这个条件成立，则n(2) ijk 可实现。 z 方向也有类似的公式：

\begin{sourceblock}
\includegraphics[page=469,pagebox=mediabox,trim={203.26bp 584.64bp 67.56bp 81.18bp},clip,max width=0.94\textwidth,max height=0.38\textheight]{Marchisio2013.pdf}
\end{sourceblock}
从中我们找到权重 w(2) ijk,α 和节点 xi(2) ijk,α，然后

\begin{sourceblock}
\includegraphics[page=469,pagebox=mediabox,trim={74.18bp 459.07bp 67.56bp 139.45bp},clip,max width=0.94\textwidth,max height=0.38\textheight]{Marchisio2013.pdf}
\end{sourceblock}
第三个可实现条件是

\begin{sourceblock}
\includegraphics[page=469,pagebox=mediabox,trim={83.95bp 405.96bp 67.56bp 253.46bp},clip,max width=0.94\textwidth,max height=0.38\textheight]{Marchisio2013.pdf}
\end{sourceblock}
全局可实现性条件（决定时间步长）是

\begin{sourceblock}
\includegraphics[page=469,pagebox=mediabox,trim={197.53bp 373.22bp 67.56bp 313.83bp},clip,max width=0.94\textwidth,max height=0.38\textheight]{Marchisio2013.pdf}
\end{sourceblock}
如果这个条件成立，那么可实现的更新矩集可以从

\begin{sourceblock}
\includegraphics[page=469,pagebox=mediabox,trim={198.44bp 326.83bp 67.56bp 338.98bp},clip,max width=0.94\textwidth,max height=0.38\textheight]{Marchisio2013.pdf}
\end{sourceblock}
从而完成多维平流的一个时间步骤。在 CQMOM 用于正交重建的情况下，导致方程 1 的步骤：对于每个排列，遵循（B.31）以找到（Mm+1 ijk ）p，然后应用方程。 (B.14) 求 Mm+1 ijk。同样有效的替代方法是使用所有 CQMOM 排列的总和来定义中间矩集 M(1) ijk，然后对 M(2) ijk 和 Mm+1 ijk 执行相同的操作。总之，理解方程式很重要。 （B.22）是构建可实现的多维平流高阶方案的关键。由此产生的方案被称为基于动力学的方案，因为数值通量函数和可实现性条件是根据 NDF 制定的，尽管我们最终只想知道传递矩集的值。因此，虽然在矩方法的背景下直接使用 NDF 似乎过于严格，但它使我们能够设计一个用于多维平流的高阶方案，该方案对于所有可能的传递矩集都是可实现的。

\section*{B.4 自由输运}\phantomsection\addcontentsline{toc}{section}{B.4 自由输运}
为了清楚起见，我们将再次描述方程式中的自由传输项。 (B.3) 在每个方向上使用算子分裂。以 x 方向为例，NDF 的自由传输方程为

\begin{sourceblock}
\includegraphics[page=469,pagebox=mediabox,trim={211.90bp 50.16bp 67.56bp 626.13bp},clip,max width=0.94\textwidth,max height=0.38\textheight]{Marchisio2013.pdf}
\end{sourceblock}
其中 v1 是内部坐标。10 为简单起见，在本节中我们假设 NDF 仅包含速度：n(t, x, v)。然而，由于其他内部坐标对于自由运输而言是被动的，因此包含额外的内部坐标将相对简单。方程中 NDF 的有限体积公式。 (B.18) 与式 (B.18) 相同。 (B.19) 但数值通量函数定义为

\begin{sourceblock}
\includegraphics[page=470,pagebox=mediabox,trim={201.15bp 555.36bp 65.07bp 128.67bp},clip,max width=0.94\textwidth,max height=0.38\textheight]{Marchisio2013.pdf}
\end{sourceblock}
其中

\begin{sourceblock}
\includegraphics[page=470,pagebox=mediabox,trim={158.32bp 514.53bp 65.07bp 160.09bp},clip,max width=0.94\textwidth,max height=0.38\textheight]{Marchisio2013.pdf}
\end{sourceblock}
请注意，用于重建 nl 和 nr 的数值 NDF nm ijk(v) 现在取决于自由传输速度分量 v1。再次，我们可以使用 nm ijk(v) 的正交重构来重写式（1）中的方案。 (B.19) 作为

\begin{sourceblock}
\includegraphics[page=470,pagebox=mediabox,trim={103.23bp 355.52bp 65.07bp 243.00bp},clip,max width=0.94\textwidth,max height=0.38\textheight]{Marchisio2013.pdf}
\end{sourceblock}
权重 wm ijk,α 和速度节点 vm ijk,α 使用体积平均矩集 Mm ijk 的正交重建找到，并且

\begin{sourceblock}
\includegraphics[page=470,pagebox=mediabox,trim={98.02bp 288.76bp 65.07bp 382.98bp},clip,max width=0.94\textwidth,max height=0.38\textheight]{Marchisio2013.pdf}
\end{sourceblock}
由此产生的可实现性条件类似于等式。 (B.23)，但速度节点 v+ 1,ijk,α 和 v− 1,ijk,α 分别代替 u+ 1,i+1/2jk,l 和 u− 1,i−1/2jk,r。事实上，方程的净效应。 (B.36) 是将 v1 节点分为两组：(1) N+ ijk 节点，0 ≤ vm 1,ijk,α < 0，其中 N = N− ijk + N+ 1,ijk,α 和 (2) N− ijk 节点，其中 vm ijk，正（负）节点的数量可以在空间和时间上变化。在真正的 NDF 是 N 点分布函数的情况下，这种空间通量的分解是精确的。 11 在许多系统中（例如涉及碰撞的系统），速度 NDF 将是连续的，因此不能用 N 点分布函数很好地表示。在这种情况下，为了更准确地表示空间通量，可以使用第 3.3.4 节中描述的多元 EQMOM 从传输矩集 Mm ijk 重建 NDF nm ijk(v)。请注意，与“常规”求积不同，12 该 NDF 将是 v 的连续（已知）函数。尽管如此，为了评估方程 1 中的空间通量， (B.35)，使用 nm ijk(v) 的矩构造正则求积仍然是有利的。作为
10 自由传输可以被视为多维平流的特例，其中 u(t, x, Ψ) = Ψ。然而，由于缺乏与 xi 无关的速度分量，我们需要单独考虑它。

\footnotetext[11]{N 点分布函数具有 N 加权狄拉克 δ 函数的形式，位于相空间中具有正权重的不同点。}
\footnotetext[12]{“正则”是指 NDF 可以用狄拉克 δ 函数之和表示的求积。}
在第 3.3.4 节中讨论过，从 nm ijk(v) 的矩构造的求积的阶数是任意的。换句话说，它可以不同于可以从传输矩集 Mm ijk 构建的规则求积的阶数。这种额外的自由度使我们能够在不增加传输矩集大小的情况下实现更高的空间通量精度。当然，在 nm ijk(v) 是 N 点分布函数的限制下，多元 EQMOM 简化为常规求积（即 CQMOM），在这种情况下(B.35) 保持不变。使用 EQMOM，我们可以将计算空间通量所需的速度矩表示为两个贡献的总和：

\begin{sourceblock}
\includegraphics[page=471,pagebox=mediabox,trim={170.61bp 470.23bp 67.56bp 167.84bp},clip,max width=0.94\textwidth,max height=0.38\textheight]{Marchisio2013.pdf}
\end{sourceblock}
其中 K* 是计算 N* 节点正则多元求积所需的矩向量，并且 N ≤ N*（即向量 K* 比 K 长，但 K 的所有分量都在 K* 中）。在等式中。 (B.37)， nm ijk(v) 是使用多元 EQMOM 从传输矩集 Mm ijk 重建的。正如第 3.3.4 节中所讨论的，由于 nm ijk(v) 具有已知的函数形式（例如多元高斯），因此方程 1 中的半矩集为： (B.37) 可以通过分析进行评估。一旦知道这些半矩集，就可以应用 N* 节点多元求积来找到两组权重与节点。总体流程如下：13

\begin{sourceblock}
\includegraphics[page=471,pagebox=mediabox,trim={115.75bp 329.40bp 97.47bp 320.13bp},clip,max width=0.94\textwidth,max height=0.38\textheight]{Marchisio2013.pdf}
\end{sourceblock}
ijk

其中 β = 1, 2, .。。，N*。通过构造，向量 v+ ijk,β（或 v− ijk,β）的 x 分量为正（或负）。我们现在可以重新表达方程。 (B.35) 作为两个贡献的总和：

\begin{sourceblock}
\includegraphics[page=471,pagebox=mediabox,trim={211.37bp 274.74bp 67.56bp 408.94bp},clip,max width=0.94\textwidth,max height=0.38\textheight]{Marchisio2013.pdf}
\end{sourceblock}
其中

\begin{sourceblock}
\includegraphics[page=471,pagebox=mediabox,trim={65.02bp 216.76bp 67.56bp 449.64bp},clip,max width=0.94\textwidth,max height=0.38\textheight]{Marchisio2013.pdf}
\end{sourceblock}
(B.39) and

\begin{sourceblock}
\includegraphics[page=471,pagebox=mediabox,trim={65.03bp 155.12bp 67.55bp 511.28bp},clip,max width=0.94\textwidth,max height=0.38\textheight]{Marchisio2013.pdf}
\end{sourceblock}
(B.40) 通过构造，方程式中的最终条款。 (B.39) 和 (B.40) 是非负的。因此，Δt 的第一个可实现条件是

\begin{sourceblock}
\includegraphics[page=471,pagebox=mediabox,trim={137.34bp 82.08bp 67.56bp 580.46bp},clip,max width=0.94\textwidth,max height=0.38\textheight]{Marchisio2013.pdf}
\end{sourceblock}
\footnotetext[13]{正如第 8.5 节中所做的那样，也可以直接使用 EQMOM 求积来分割正负速度。}
如果满足这个条件，则 n(1) ijk 是可实现的，因此矩集 M(1) ijk 也是可实现的。 14 结合中间步骤，从 Mm ijk 计算自由传输的 M(1) ijk 的过程为

\begin{sourceblock}
\includegraphics[page=472,pagebox=mediabox,trim={171.85bp 561.59bp 148.59bp 111.98bp},clip,max width=0.94\textwidth,max height=0.38\textheight]{Marchisio2013.pdf}
\end{sourceblock}
ijk

此时，仍然存在一个技术问题：如何定义多元求积？如果使用强力 QMOM，则只有一种选择，因为所有传输的矩都是用这种方法再现的。然而，由于暴力 QMOM 在数值上是病态的，所以更可能的选择是使用 CQMOM（详细信息请参见第 3.2.3 节）。正如 Yuan \& Fox (2011) 中所述，当 CQMOM 用于空间传输时，在选择排列时必须小心。例如，对于 v（即 3D 相空间），CQMOM 有 P = 6 种排列，但只有其中两种给出了 x 方向上的正确空间通量（即排列 123 和 132）。 15 为了适应这两种 CQMOM 排列，上述更新过程修改如下：

\begin{sourceblock}
\includegraphics[page=472,pagebox=mediabox,trim={124.42bp 373.09bp 112.56bp 270.12bp},clip,max width=0.94\textwidth,max height=0.38\textheight]{Marchisio2013.pdf}
\end{sourceblock}
请注意，在求 M(1) ijk 时，相同的 CQMOM 排列同时用于常规求积和扩展求积。 y 方向的公式与上面给出的 x 方向的公式类似。唯一的区别是，现在我们必须计算 y 方向上的半矩集，定义为

\begin{sourceblock}
\includegraphics[page=472,pagebox=mediabox,trim={173.93bp 241.04bp 65.07bp 397.03bp},clip,max width=0.94\textwidth,max height=0.38\textheight]{Marchisio2013.pdf}
\end{sourceblock}
计算 M(2) ijk 的过程可概括如下：

\begin{sourceblock}
\includegraphics[page=472,pagebox=mediabox,trim={151.04bp 182.13bp 148.59bp 491.46bp},clip,max width=0.94\textwidth,max height=0.38\textheight]{Marchisio2013.pdf}
\end{sourceblock}
Δt 的第二个可实现条件是

\begin{sourceblock}
\includegraphics[page=472,pagebox=mediabox,trim={139.83bp 116.92bp 65.07bp 545.63bp},clip,max width=0.94\textwidth,max height=0.38\textheight]{Marchisio2013.pdf}
\end{sourceblock}
14 形式上，矩集 M(1) i jk 是根据等式 1 中的定义找到的。 （B.24）。然而，考虑到 n(1) i jk 是一个 2N* 点的分布函数，方程中的积分： (B.24) 成为求和。 15排列根据条件变量进行编号。例如，排列 123 从 v1 方向开始 CQMOM，然后是 v2 方向，最后是 v3 方向。换句话说，v2 以 v1 为条件，v3 以 v2 和 v1 为条件。

当 CQMOM 用于多元求积时，我们必须适应两个 CQMOM 排列：

\begin{sourceblock}
\includegraphics[page=473,pagebox=mediabox,trim={112.52bp 556.01bp 115.04bp 87.21bp},clip,max width=0.94\textwidth,max height=0.38\textheight]{Marchisio2013.pdf}
\end{sourceblock}
如果第二个可实现性条件成立，则 M(2) ijk 可实现。 z方向的处理以同样的方式进行，从M(2) ijk 开始。 Δt 的第三个可实现条件是

\begin{sourceblock}
\includegraphics[page=473,pagebox=mediabox,trim={137.34bp 463.24bp 67.56bp 199.31bp},clip,max width=0.94\textwidth,max height=0.38\textheight]{Marchisio2013.pdf}
\end{sourceblock}
因此全局时间步必须满足 Δt ≤min(Δt1, Δt2, Δt3)。当 CQMOM 用于多元求积时，我们必须再次适应两个 CQMOM 排列：

\begin{sourceblock}
\includegraphics[page=473,pagebox=mediabox,trim={102.88bp 361.76bp 105.40bp 281.45bp},clip,max width=0.94\textwidth,max height=0.38\textheight]{Marchisio2013.pdf}
\end{sourceblock}
如果 CQMOM 的所有排列都满足全局可实现性条件，则更新的传输矩集 Mm+1 ijk 将是可实现的。16 为了结束本节，我们讨论在应用自由传输的可实现高阶方案时出现的三个技术点。第一点是如何选择 N* 的值，假设 N = n3 用于多元 EQMOM。为了回答这个问题，我们首先注意到任何一个方向上的最高阶传递矩都是 2n 阶（对应于任何一个方向上的 2n+1 个纯矩）。 EQMOM 中使用“额外”（偶数阶）矩来确定核密度函数的分布（详细信息请参见第 3.3.2 节）。当自由传输应用于 2n 阶矩时，通量函数涉及 2n + 1 阶矩。因此，为了使用半矩集准确预测自由传输的通量函数，我们必须有 N*≥(n+1)3。因此，N* 的明显选择是 N*= (n + 1)3。第二点涉及 N* 权重的高阶空间重建。由于 CQMOM 使用自适应一维正交例程来控制正交误差（详细信息请参阅 Yuan \& Fox (2011)），因此相邻单元中的 CQMOM 重建有时可能会有所不同，因为并非所有单元都具有 N* 非零权重。对权重使用高阶空间重建要求相邻单元中的 CQMOM 重建以完全相同的方式完成。如果不这样做，则只能使用一阶空间重建。因此，为了使用高阶空间重建，必须将索引添加到唯一定义 CQMOM 重建的 CQMOM 例程中（示例参见 Yuan \& Fox (2011)）。通过比较相邻单元中的索引，可以确定空间重建的正确局部顺序。

\footnotetext[16]{换句话说，给定方向的可实现性条件必须满足该方向使用的 CQMOM 的两种排列。}
最后一点涉及添加被动内坐标，使 NDF 为 n(t, x, v, xi)。当使用 CQMOM 重建多元求积近似时，（原则上）对添加被动内坐标没有限制。然而，在实践中，传输矩的数量类似于 (1 + 3 + M)n3+M，其中 M 是 ϋ 的长度，因此即使 M 值中等，构建求积也需要大量矩。然而，即使 M = 1，在评估自由传输数值通量时如何选择 CQMOM 排列的问题也会出现。例如，在 x 方向上，我们可以使用 xi1 启动 CQMOM，然后选择 123 或 132 作为速度变量，或者我们可以将 xi1 放在调节顺序中的其他位置。目前，这个问题还没有得到明确的解决，但是，从我们在第 8.1.3 节中的观察来看，为了正确预测 Ψ 的空间通量，在计算自由传输空间通量时最好将被动变量放在 CQMOM 条件序列的开头。

\section*{B.5 混合平流}\phantomsection\addcontentsline{toc}{section}{B.5 混合平流}
方程中的混合平流项。 (B.4) 涉及 NDF 的空间梯度和相空间梯度，因此它比本附录 B.3 节中描述的纯平流项更复杂。尽管如此，通过处理时刻的输运方程，我们可以证明它可以被视为具有“特殊”属性的平流项。将 NDF 表示为 n(t, x, v, xi)，混合平流情况下 x 方向的 NDF 输运方程为

\begin{sourceblock}
\includegraphics[page=474,pagebox=mediabox,trim={220.81bp 332.01bp 65.07bp 340.50bp},clip,max width=0.94\textwidth,max height=0.38\textheight]{Marchisio2013.pdf}
\end{sourceblock}
其中 λ(t, x, v, xi) 为非负数。相应的矩输运方程为

\begin{sourceblock}
\includegraphics[page=474,pagebox=mediabox,trim={139.13bp 282.57bp 65.07bp 383.59bp},clip,max width=0.94\textwidth,max height=0.38\textheight]{Marchisio2013.pdf}
\end{sourceblock}
其中右侧的最终形式是使用分部积分找到的。更进一步，我们注意到矩向量 K 的分量形式为 vk1 1 vk2 2 vk3 3 ΨkΨ，其中指数向量 k = (k1, k2, k3, kΨ) 包含对应于传输矩集 M 的非负整数。因此，方程 1 中出现的 K 的 v1 导数。 (B.46) 的形式为

\begin{sourceblock}
\includegraphics[page=474,pagebox=mediabox,trim={171.63bp 193.01bp 169.19bp 480.45bp},clip,max width=0.94\textwidth,max height=0.38\textheight]{Marchisio2013.pdf}
\end{sourceblock}
通过消除相空间积分，方程。 (B.46) 可以重写为 NDF 的 k1 相关平流项：

\begin{sourceblock}
\includegraphics[page=474,pagebox=mediabox,trim={209.23bp 144.24bp 65.07bp 522.79bp},clip,max width=0.94\textwidth,max height=0.38\textheight]{Marchisio2013.pdf}
\end{sourceblock}
形式上，我们可以将平流速度定义为 u1(t, x, v, xi) = k1λ/v1，因此方程： (B.47) 简化为 NDF 的平流方程：

\begin{sourceblock}
\includegraphics[page=474,pagebox=mediabox,trim={213.28bp 82.57bp 65.07bp 591.68bp},clip,max width=0.94\textwidth,max height=0.38\textheight]{Marchisio2013.pdf}
\end{sourceblock}
{\small 请注意，平流速度线性取决于 k1，因此最高速度将与 M 中出现的最高阶矩相关。令 k*\par}
\footnotetext[1]{= max(k1) 是最大的}
出现在传输矩集中的 k1 值（即，对于常规求积，k* 1 = 2n -1，对于扩展求积，k* 1 = 2n -1）。然后我们可以将最大平流速度定义为 u* 1 = k* 1λ/v1，并且正是这个速度将出现在可实现性条件中。请注意，由于 λ 为非负数，因此 u1 的符号由 v1 确定。另外，由于 u1 与 v1 成反比，因此当 v1 = 0 时，平流速度将无穷大，除非 λ 也为零。一般来说，λ的函数形式对于确定算法的行为非常重要，我们将在介绍混合平流方案后回到这个主题。方程中的平流问题。 (B.48) 的形式与方程 1 中讨论的形式完全相同。 （B.3），但具有平流速度的特定定义。例如，对应于等式。 （B.21），我们有

\begin{sourceblock}
\includegraphics[page=475,pagebox=mediabox,trim={140.51bp 472.87bp 67.56bp 200.60bp},clip,max width=0.94\textwidth,max height=0.38\textheight]{Marchisio2013.pdf}
\end{sourceblock}
然后，更新后的数值 NDF n(1) ijk 可以写为两个贡献之和，如式（1）所示。 （B.38），其中

\begin{sourceblock}
\includegraphics[page=475,pagebox=mediabox,trim={88.43bp 363.52bp 67.56bp 267.86bp},clip,max width=0.94\textwidth,max height=0.38\textheight]{Marchisio2013.pdf}
\end{sourceblock}
以及

\begin{sourceblock}
\includegraphics[page=475,pagebox=mediabox,trim={85.42bp 269.54bp 67.56bp 361.84bp},clip,max width=0.94\textwidth,max height=0.38\textheight]{Marchisio2013.pdf}
\end{sourceblock}
方程式中的最终条款。 （B.50）和等式。 (B.51) 是非负的，因此如果 Δt 的以下可实现性条件成立，则 n(1) ijk 将可实现：

\begin{sourceblock}
\includegraphics[page=475,pagebox=mediabox,trim={73.23bp 196.91bp 67.56bp 465.63bp},clip,max width=0.94\textwidth,max height=0.38\textheight]{Marchisio2013.pdf}
\end{sourceblock}
接下来是更新后的矩集，定义为

\begin{sourceblock}
\includegraphics[page=475,pagebox=mediabox,trim={189.36bp 143.84bp 67.56bp 521.97bp},clip,max width=0.94\textwidth,max height=0.38\textheight]{Marchisio2013.pdf}
\end{sourceblock}
也是可以实现的。如前所述，正整数 k*1 出现在可实现性条件中，导致当传递矩集中包含高阶矩时，时间步更小。同样，我们可以观察到时间步长将由 |v1| 控制。节点最接近于零。

{\small 与我们在自由传输的情况下所做的类似，通过混合平流，我们定义了计算等式中使用的权重与节点所需的半矩集。 (B.50) 和 (B.51) 为\par}
\begin{sourceblock}
\includegraphics[page=476,pagebox=mediabox,trim={161.61bp 583.48bp 65.07bp 63.78bp},clip,max width=0.94\textwidth,max height=0.38\textheight]{Marchisio2013.pdf}
\end{sourceblock}
在此定义中，nm ijk(v, xi) 是使用来自传输矩集 Mm ijk 的扩展多元求积来重建的。由于 nm ijk(v, xi) 具有已知的函数形式，因此半矩在方程 1 中设定。 （B.54）可以进行分析评估。一旦知道这些半矩集，就应用 N 节点多元求积来找到两组权重与节点。整体流程如下：

\begin{sourceblock}
\includegraphics[page=476,pagebox=mediabox,trim={106.83bp 466.38bp 83.56bp 183.15bp},clip,max width=0.94\textwidth,max height=0.38\textheight]{Marchisio2013.pdf}
\end{sourceblock}
ijk

请注意，对于混合平流，扩展求积和常规求积都使用 N 个节点。这是可能的，因为对于混合平流，空间通量中出现的最高阶矩由于除以 v1 而具有较低阶（即，对于自由传输而不是乘以 v1）。然而，只有当 λ 不依赖于 v 时，这一观察结果才成立。y 和 z 方向上的混合对流公式与上面针对 x 方向导出的公式类似。同样，本附录 B.4 节中有关使用 CQMOM 进行免费运输的评论也适用于混合平流。混合平流的一个关键技术点是 λ 对 v 的依赖性，特别是当 v 的任何一个分量趋于零时。例如，方程中的 λ(v)。 (B.52) 具有三种可能的限制行为：

\begin{sourceblock}
\includegraphics[page=476,pagebox=mediabox,trim={164.37bp 276.60bp 295.90bp 375.74bp},clip,max width=0.94\textwidth,max height=0.38\textheight]{Marchisio2013.pdf}
\end{sourceblock}
0 自由传输，恒定平流，∞ 有效（椭圆）压力。林 v1→0
前两个对应于自由传输和平流，其中时间步长由具有最大幅度的速度节点控制。另一方面，第三种情况是由具有最小幅度的速度节点控制的。表现出这些行为的二阶张量 λ 的简单形式是

\begin{sourceblock}
\includegraphics[page=476,pagebox=mediabox,trim={168.65bp 171.08bp 65.07bp 485.01bp},clip,max width=0.94\textwidth,max height=0.38\textheight]{Marchisio2013.pdf}
\end{sourceblock}
其中 0 ≤ν ≤2 且 0 ≤λ0(xi) 的维度为 (m/s)2−ν，并且取决于粒子斯托克斯数等。在混合平流对气体-粒子流的建议应用中（Pandya \& Mashayek, 2003b; Reeks, 1977, 1983, 1993; Zaichik, 1997, 1999; Zaichik et al., 2008），λ 被认为与 v 无关，因此它表现为有效压力。因此，为了从方程计算出时间步长。 (B.52) 为非零，则 λ 在 v− 的空间位置处必须为零。

{\small 1,ijk,β 或 v+ 1,ijk,β 为零。由于这些节点是根据半矩集 M− 计算得出的\par}
{\small 1,ijk 和 M+ 1,ijk 分别为，仅当 EQMOM 中使用的核密度函数具有零宽度时，它们才可以为零。 （参见第 3.3.2 节\par}
有关 EQMOM 的更多详细信息。）作为示例，让我们考虑宽度为 σ 且平均值为 u1 的高斯核密度函数。为了计算 v+ 1，我们需要正尾中的矩，定义为

\begin{sourceblock}
\includegraphics[page=477,pagebox=mediabox,trim={169.35bp 569.58bp 175.60bp 99.18bp},clip,max width=0.94\textwidth,max height=0.38\textheight]{Marchisio2013.pdf}
\end{sourceblock}
当 u1 <0 时会发生最坏的情况（即对于常规求积，狄拉克 δ 函数将位于负 v1 速度）。在这个极限下，我们发现

\begin{sourceblock}
\includegraphics[page=477,pagebox=mediabox,trim={214.23bp 491.44bp 216.76bp 181.05bp},clip,max width=0.94\textwidth,max height=0.38\textheight]{Marchisio2013.pdf}
\end{sourceblock}
因此σ2 控制了当u1 <0 时速度节点距原点的距离。对于极限 0 ≪u1 中的 v− 1 也是如此。对于多元 EQMOM，矩阵 Σ 起着 σ2 的作用。这一观察结果表明，一个好的 λ 模型是

\begin{sourceblock}
\includegraphics[page=477,pagebox=mediabox,trim={199.66bp 424.15bp 67.56bp 263.68bp},clip,max width=0.94\textwidth,max height=0.38\textheight]{Marchisio2013.pdf}
\end{sourceblock}
请注意，由于 Σ 的单位为 (m/s)2，因此非负函数 h(xi) 将是无量纲的。有了这个模型，方程中的可实现性条件为 λ。当 h(xi) 有限时，(B.52) 总是会在 Δt 上产生非零上限。17 物理上，Σ 在纯粒子轨迹穿越的极限中为零，其中真正的 NDF 是狄拉克 δ 函数的总和。另一方面，当 Σ 达到最大值时，NDF 呈高斯分布。因此，由于混合平流与“随机”粒子运动相关，方程中的模型(B.56) 也具有物理意义。尽管如此，更新公式中潜在的奇异行为使得混合平流的处理存在问题。总之，在本节中，我们证明了混合平流的可实现方案与平流和自由传输的方案非常相似。然而，该方案的数值特性很大程度上取决于 λ 使用的函数形式。为了使混合平流的“标准”模型得到良好的调节，λ 的函数形式必须补偿距离原点太近的速度节点。对于多元 EQMOM，方程中给出的模型。 (B.56) 可用于此目的。然而，混合平流的奇异行为的可能性使其成为问题，并且可能需要替代方案来正确处理这种行为。

\section*{B.6 扩散}\phantomsection\addcontentsline{toc}{section}{B.6 扩散}
对于扩散系数为 Γ(t, x, xi) 的多维扩散，我们将再次在每个空间方向上使用算子分裂（Vikas et al., 2012）。在x方向上，扩散方程为

\begin{sourceblock}
\includegraphics[page=477,pagebox=mediabox,trim={211.44bp 115.03bp 67.56bp 559.32bp},clip,max width=0.94\textwidth,max height=0.38\textheight]{Marchisio2013.pdf}
\end{sourceblock}
17 在气体-颗粒流中，h 取决于颗粒斯托克斯数，如下所示
{\small h ∝ 1 St(St + 1)。\par}
{\small 因此，当斯托克斯数非零时，h 将是有限的。\par}
形式上，相应的有限体积公式为

\begin{sourceblock}
\includegraphics[page=478,pagebox=mediabox,trim={142.83bp 588.54bp 65.07bp 75.68bp},clip,max width=0.94\textwidth,max height=0.38\textheight]{Marchisio2013.pdf}
\end{sourceblock}
其中 NDF 的空间梯度必须根据正交表示来定义：

\begin{sourceblock}
\includegraphics[page=478,pagebox=mediabox,trim={88.98bp 510.73bp 65.07bp 151.81bp},clip,max width=0.94\textwidth,max height=0.38\textheight]{Marchisio2013.pdf}
\end{sourceblock}
对于空间重建的每个阶，有必要定义方程式中出现的空间梯度项。 （B.59）。对于二阶空间重建，界面处的权重定义为

\begin{sourceblock}
\includegraphics[page=478,pagebox=mediabox,trim={184.79bp 408.72bp 65.07bp 241.67bp},clip,max width=0.94\textwidth,max height=0.38\textheight]{Marchisio2013.pdf}
\end{sourceblock}
以及接口处的 delta 函数

\begin{sourceblock}
\includegraphics[page=478,pagebox=mediabox,trim={147.02bp 332.41bp 65.07bp 317.98bp},clip,max width=0.94\textwidth,max height=0.38\textheight]{Marchisio2013.pdf}
\end{sourceblock}
梯度定义为

\begin{sourceblock}
\includegraphics[page=478,pagebox=mediabox,trim={177.53bp 250.42bp 65.07bp 393.11bp},clip,max width=0.94\textwidth,max height=0.38\textheight]{Marchisio2013.pdf}
\end{sourceblock}
以及

\begin{sourceblock}
\includegraphics[page=478,pagebox=mediabox,trim={139.89bp 179.95bp 65.07bp 463.58bp},clip,max width=0.94\textwidth,max height=0.38\textheight]{Marchisio2013.pdf}
\end{sourceblock}
在方程式中使用这些公式。 （B.58）导致

\begin{sourceblock}
\includegraphics[page=478,pagebox=mediabox,trim={126.70bp 54.73bp 65.07bp 543.78bp},clip,max width=0.94\textwidth,max height=0.38\textheight]{Marchisio2013.pdf}
\end{sourceblock}
其中 γx = Δt/(Δx)2。该 NDF 是可以实现的，前提是

\begin{sourceblock}
\includegraphics[page=479,pagebox=mediabox,trim={157.96bp 590.67bp 67.56bp 71.87bp},clip,max width=0.94\textwidth,max height=0.38\textheight]{Marchisio2013.pdf}
\end{sourceblock}
y 和 z 方向的表达式类似于等式： （B.64）。读者可以注意到，当 Г(t, x) 不依赖于 ψ 时，方程： (B.64) 得出更新矩集的封闭的、可实现的公式：

\begin{sourceblock}
\includegraphics[page=479,pagebox=mediabox,trim={76.65bp 527.12bp 67.56bp 148.95bp},clip,max width=0.94\textwidth,max height=0.38\textheight]{Marchisio2013.pdf}
\end{sourceblock}
事实上，该表达式与直接使用使用传递矩集编写的扩散项时发现的二阶表达式相同。一般来说，非结构化网格上的二阶空间重构也将是可以实现的。也许由于这个原因，扩散项通常不被认为是不可实现矩集的潜在来源。对于高阶空间重建，只有利用相邻单元的节点来构建方案，NDF才能实现。换句话说，我们必须继续使用方程式。 （B.61）和（B.63）。然后我们可以使用高阶方案来近似权重的空间导数。例如，对于四阶方案，我们将使用

\begin{sourceblock}
\includegraphics[page=479,pagebox=mediabox,trim={121.95bp 357.39bp 67.56bp 293.00bp},clip,max width=0.94\textwidth,max height=0.38\textheight]{Marchisio2013.pdf}
\end{sourceblock}
以及

\begin{sourceblock}
\includegraphics[page=479,pagebox=mediabox,trim={99.17bp 284.35bp 67.56bp 356.57bp},clip,max width=0.94\textwidth,max height=0.38\textheight]{Marchisio2013.pdf}
\end{sourceblock}
请注意，如果所有单元中的节点都是恒定的，则总体方案将具有与权重的空间重建相同的顺序。另一方面，如果权重恒定，则节点的总体方案是二阶的。应用上面的公式，更新后的数值 NDF 可以写为

\begin{sourceblock}
\includegraphics[page=479,pagebox=mediabox,trim={101.42bp 93.55bp 67.56bp 474.73bp},clip,max width=0.94\textwidth,max height=0.38\textheight]{Marchisio2013.pdf}
\end{sourceblock}
{\small 如果每个乘以 delta 函数的系数对于所有网格单元和所有节点都是非负的，则该 NDF 是可实现的。如前所述，y 和 z 方向的表达式类似于等式： （B.69）。\par}
由于 x 方向上的扩散而引起的矩集变化的最终方程为

\begin{sourceblock}
\includegraphics[page=480,pagebox=mediabox,trim={99.51bp 486.56bp 65.07bp 81.71bp},clip,max width=0.94\textwidth,max height=0.38\textheight]{Marchisio2013.pdf}
\end{sourceblock}
y 和 z 方向具有类似的表达式。如前所述，如果采用 CQMOM 进行多元求积重建，则应使用每个 CQMOM 排列来更新矩集，然后对结果进行平均以找到 Mm+1 ijk。

\chapter{双曲方程的矩方法}
在文献中（Mazzei，2011）有人认为，如果不添加“扩散”项，基于求积的矩方法（QBMM）就不能用于表示双曲动力学方程的解。该论点认为，扩散项对于表示源自系统空间不同区域的两个 NDF 之间的“微混合”是必要的。本附录的目的是提供一个分析反例，清楚地表明双曲动力学方程的精确解（即没有扩散项）可以在很大程度上产生“完全混合”（或更准确地说，均匀的）NDF。正如我们将看到的，该解的出现完全是由于平流速度的大小依赖性。我们还使用解析解来说明使用 DQMOM 和 QMOM 求解具有空间非光滑解的动力学方程（即当矩的空间导数不连续时）时出现的问题。最后，值得注意的是，本附录中导出的解析解可用于与模型动理学方程的数值解进行比较，以验证数值算法的准确性。

\section*{C.1 模型动理学方程}\phantomsection\addcontentsline{toc}{section}{C.1 模型动理学方程}
为了清楚地说明相关问题，我们考虑具有一个空间维度的 NDF，其平流速度取决于单个内部坐标 ε（我们将其称为“尺寸”）。相应的动力学方程为

\begin{sourceblock}
\includegraphics[page=481,pagebox=mediabox,trim={203.36bp 153.36bp 67.55bp 520.90bp},clip,max width=0.94\textwidth,max height=0.38\textheight]{Marchisio2013.pdf}
\end{sourceblock}
其中 f(t, x, xi) 是 NDF，平流速度 u(xi) 假设已知。请注意，虽然可以将 u(xi) 推广为取决于 t 和 x，而不改变方程的基本数学属性。 (C.1)，这样做会使解析解变得非常复杂，因此只会模糊其基本属性。在多相流的应用中，u(xi) 对应于所有尺寸为 xi 的颗粒的平均速度（即条件速度），因此根据这个定义，u(xi) 是单值的（即对于每个 xi 都有一个速度）。这绝不排除粒子轨迹穿越的可能性，因为后者是瞬时速度的属性，而不是其条件平均值的属性。

\section*{C.2 分离初始条件的解析解}\phantomsection\addcontentsline{toc}{section}{C.2 分离初始条件的解析解}
给定初始 NDF f(0, x, xi)，方程： (C.1)有解析解

\begin{sourceblock}
\includegraphics[page=482,pagebox=mediabox,trim={191.95bp 582.13bp 65.07bp 105.70bp},clip,max width=0.94\textwidth,max height=0.38\textheight]{Marchisio2013.pdf}
\end{sourceblock}
可以通过代入等式来验证（C.2）代入等式。 (C.1)。该解对于任意初始 NDF 都是精确的。这里我们将考虑空间隔离的初始 NDF：

\begin{sourceblock}
\includegraphics[page=482,pagebox=mediabox,trim={160.21bp 498.25bp 65.07bp 154.08bp},clip,max width=0.94\textwidth,max height=0.38\textheight]{Marchisio2013.pdf}
\end{sourceblock}
其中 0 < xi*< 1 表示截止粒径。换句话说，所有较小的粒子最初都位于 x = 0 的下方，而较大的粒子则位于它们上方。 NDF f1 和 f2 是针对等式中指示的 xi 范围定义的。 (C.3)，但具有任意形状（即它们可以是平滑函数，甚至是 delta 函数）。使用方程式。 (C.2)和(C.3)，初始条件分离的解析解可写为

\begin{sourceblock}
\includegraphics[page=482,pagebox=mediabox,trim={124.69bp 402.48bp 65.07bp 284.57bp},clip,max width=0.94\textwidth,max height=0.38\textheight]{Marchisio2013.pdf}
\end{sourceblock}
其中 IS (x) 是区间 S 的指示函数：

\begin{sourceblock}
\includegraphics[page=482,pagebox=mediabox,trim={199.60bp 343.39bp 65.07bp 320.43bp},clip,max width=0.94\textwidth,max height=0.38\textheight]{Marchisio2013.pdf}
\end{sourceblock}
请注意方程中的解析解。 (C.4) 沿着 x = u(xi)t 定义的 (t, x, xi)- 空间中的表面是不连续的。因此，正如我们接下来将看到的，u(xi) 的函数形式将决定 NDF 是否会随着 t 的增加而进一步混合或分离。

\subsection*{C.2.1 分离解}\phantomsection\addcontentsline{toc}{subsection}{C.2.1 分离解}
我们首先考虑 u(xi) 是 xi 的单调递增函数的情况。具体来说，我们令 u(xi) = xi − xi*。方程的特征。 (C.1) 然后由 x(t, Ψ) = x(0) + (Ψ − Ψ*)t 定义，因此，使用等式： (C.3)，我们看到

\begin{sourceblock}
\includegraphics[page=482,pagebox=mediabox,trim={162.19bp 186.90bp 65.07bp 476.92bp},clip,max width=0.94\textwidth,max height=0.38\textheight]{Marchisio2013.pdf}
\end{sourceblock}
换句话说，随着 t 的增加，从负 x 开始的特征变得更加消极，而从正 x 开始的特征变得更加积极。在所有情况下，我们都可以从等式中看到。 (C.6) 特征永远不会相互交叉（即，如果 Ψ1 Ψ2，则 x(t, Ψ1) x(t, Ψ2) 对于任何 t > 0）。从 x(0) = 0 开始且 xi = xi* 的特征具有 x(t, xi*) = 0 的属性，并且对应于两个初始解之间的接触点（或不连续性），分别为左 NDF f1(xi) 和右 NDF f2(xi)。由于特征不交叉，解析解可以写成空间中的两个半区间：

\begin{sourceblock}
\includegraphics[page=482,pagebox=mediabox,trim={137.06bp 51.65bp 65.07bp 611.39bp},clip,max width=0.94\textwidth,max height=0.38\textheight]{Marchisio2013.pdf}
\end{sourceblock}
因此，NDF 在任何时候都保持“未混合”，并且在 x = ( xi − xi ) t 线上不连续地变化。请注意，对于较大的 t，只有从 xi* 开始的特征（即速度为零）将保留在系统中。结果，随着时间的增加，初始 NDF 根据尺寸变得越来越分离。最后，请注意，只要 u(ψ) 单调递增，改变对 ψ 的函数依赖性就不会改变方程 (1) 的定性行为。 (C.6) 或等式。 （C.7）。

\subsection*{C.2.2 混合解}\phantomsection\addcontentsline{toc}{subsection}{C.2.2 混合解}
接下来我们考虑 u(xi) 是 xi 的单调递减函数的情况。具体来说，我们令 u(xi) = xi*−xi。方程的特征。 (C.1) 由 x(t, Ψ) = x(0) + (Ψ*−Ψ)t 定义，因此，使用等式： (C.3)，我们看到

\begin{sourceblock}
\includegraphics[page=483,pagebox=mediabox,trim={159.69bp 440.40bp 67.55bp 223.43bp},clip,max width=0.94\textwidth,max height=0.38\textheight]{Marchisio2013.pdf}
\end{sourceblock}
换句话说，随着 t 的增加，从负 x 开始的特征最终会变为正，而从正 x 开始的特征最终会变为负。因此，我们可以从等式得出结论。 (C.8) 对应于 Ψ < Ψ* 的特征将在 t 的某个值处交叉对应于 Ψ*< Ψ 的所有特征。1 从 x(0) = 0 开始且 Ψ = Ψ* 的特征具有属性 x(t, Ψ*) = 0，并且再次对应于两个初始 NDF 之间的接触点。然而，现在有两个不连续点对应于从 x(0) = 0 开始的两个特征，分别为 xi = 0 和 xi = 1。第一个表示 NDF f1 渗透到最初由 f2 占据的空间区域，第二个表示相反。通过考虑交叉特性，解析解可以写为

\begin{sourceblock}
\includegraphics[page=483,pagebox=mediabox,trim={118.76bp 237.77bp 67.55bp 400.42bp},clip,max width=0.94\textwidth,max height=0.38\textheight]{Marchisio2013.pdf}
\end{sourceblock}
从这个解析解中，可以清楚地看出 xi*−1 < x/t < xi* 的 NDF 是两个原始 NDF 的组合。此外，当 t 增加到无穷大时，我们发现

\begin{sourceblock}
\includegraphics[page=483,pagebox=mediabox,trim={133.55bp 166.77bp 67.56bp 506.61bp},clip,max width=0.94\textwidth,max height=0.38\textheight]{Marchisio2013.pdf}
\end{sourceblock}
它对应于整个实数线上的初始 NDF 分布（即“充分混合”的系统）。再次注意，只要 u(ψ) 单调递减，改变对 ψ 的函数依赖就不会改变方程 (1) 的定性行为。 (C.8) 或等式。 （C.10）。如前所述，混合解是在不添加扩散的情况下从全双曲动力学方程获得的。事实上，导致混合的关键技术点是双曲方程解中特征的交叉。

\footnotetext[1]{确切的交叉时间为 tc = (x2(0) −x1(0))/(ψ2 −ψ1)，其中 ψ1 < ψ*< ψ2。}
\section*{C.3 矩和求积近似}\phantomsection\addcontentsline{toc}{section}{C.3 矩和求积近似}
给定解析解，使用定义计算整数矩很简单

\begin{sourceblock}
\includegraphics[page=484,pagebox=mediabox,trim={191.03bp 567.87bp 65.07bp 98.34bp},clip,max width=0.94\textwidth,max height=0.38\textheight]{Marchisio2013.pdf}
\end{sourceblock}
同样，可以从这些时刻找到正交近似。例如，对于单节点求积，权重与节点由下式定义

\begin{sourceblock}
\includegraphics[page=484,pagebox=mediabox,trim={147.37bp 507.34bp 65.07bp 166.13bp},clip,max width=0.94\textwidth,max height=0.38\textheight]{Marchisio2013.pdf}
\end{sourceblock}
请注意，由于动力学方程的双曲形式，矩、权重与节点不必是 x 的平滑函数。为简单起见，我们将考虑两种情况的时刻。第一个使用由 f1 = f2 = I[0,1](Ψ) 定义的连续初始 NDF，第二个使用由 f1 = δ(Ψ−Ψ*/2) 和 f2 = δ(Ψ -3Ψ*/2) 定义的离散 NDF。

\subsection*{C.3.1 分离解的矩}\phantomsection\addcontentsline{toc}{subsection}{C.3.1 分离解的矩}
\paragraph*{连续无纤维膜}
具有连续初始 NDF 的分离解的矩可以通过方程 1 进行分析计算。 （C.7）：

\begin{sourceblock}
\includegraphics[page=484,pagebox=mediabox,trim={89.87bp 337.22bp 65.07bp 326.61bp},clip,max width=0.94\textwidth,max height=0.38\textheight]{Marchisio2013.pdf}
\end{sourceblock}
示例如图 C.1 所示。单节点正交近似由下式给出

\begin{sourceblock}
\includegraphics[page=484,pagebox=mediabox,trim={81.85bp 246.47bp 141.48bp 386.54bp},clip,max width=0.94\textwidth,max height=0.38\textheight]{Marchisio2013.pdf}
\end{sourceblock}
（C.14）和二节点求积如图 C.2 所示。对于这种情况，当 t > 0 时，x = 0 处的上方和下方的极限产生 xi1(t, 0) = xi*，因此节点在所有 x 中都是连续的。但请注意，当 t > 0 时，mk(t, 0) = 0，并且矩的空间梯度在 x = 0 处不连续。当使用多个节点表示 t > 0 时的矩时，所有权重在 x = 0 处均为零。然而，由于真正的 NDF 接近于以 ψ = ψ* 为中心的 δ 函数（对于 x 接近于零），因此所有节点都平滑地逼近 ψ*，因为 x →0 可以这样实现：见图C.2。

\paragraph*{离散NDF}
对于离散初始条件，矩为

\begin{sourceblock}
\includegraphics[page=484,pagebox=mediabox,trim={153.93bp 52.26bp 65.07bp 597.55bp},clip,max width=0.94\textwidth,max height=0.38\textheight]{Marchisio2013.pdf}
\end{sourceblock}
\begin{sourceblock}
\includegraphics[page=485,pagebox=mediabox,trim={58.00bp 466.80bp 58.23bp 59.00bp},clip,width=\textwidth,max height=0.72\textheight]{Marchisio2013.pdf}
\par\vspace{0.45\baselineskip}
\small 图C.1。方程中分离解的矩示例。 (C.13) 其中 xi*= 1/2。
\end{sourceblock}
\begin{sourceblock}
\includegraphics[page=485,pagebox=mediabox,trim={58.00bp 235.11bp 58.23bp 290.75bp},clip,width=\textwidth,max height=0.72\textheight]{Marchisio2013.pdf}
\par\vspace{0.45\baselineskip}
\small 图C.2。从图 C.1 中的四个时刻找到的权重与节点。
\end{sourceblock}
读者可以注意到，非零区域的上限和下限对应于根据等式计算的特征。 (C.6) 从 x(0) = 0 开始。这种情况的权重与节点具有简单的形式：

w1 = 1, Σ1 = Σ*/2 (−∞< x/t ≤−Σ*/2), w1 = 0, Σ1 未定义 (−Σ*/2 < x/t ≤Σ*/2), w1 = 1, Σ1 = 3Σ*/2 (Σ*/2 ≤x/t < Infini)。

\begin{sourceblock}
\includegraphics[page=485,pagebox=mediabox,trim={397.94bp 104.89bp 67.56bp 583.31bp},clip,max width=0.94\textwidth,max height=0.38\textheight]{Marchisio2013.pdf}
\end{sourceblock}
{\small 由于隔离，该系统创建了一个无法定义节点的真空区域。在该区域之外，权重与节点是恒定的。\par}
\subsection*{C.3.2 混合解的矩}\phantomsection\addcontentsline{toc}{subsection}{C.3.2 混合解的矩}
\paragraph*{连续无纤维膜}
分离解的矩可以通过方程进行分析计算。 （C.10）：

\begin{sourceblock}
\includegraphics[page=486,pagebox=mediabox,trim={74.21bp 502.50bp 65.07bp 120.99bp},clip,max width=0.94\textwidth,max height=0.38\textheight]{Marchisio2013.pdf}
\end{sourceblock}
示例如图C.3所示，相应的正交权重与节点如图C.4所示。对于这种情况，x = 0 处的矩只是在所有 t > 0 的情况下从 f1 和 f2 找到的矩之和。还可以看出，虽然矩

\begin{sourceblock}
\includegraphics[page=486,pagebox=mediabox,trim={58.00bp 271.06bp 58.23bp 257.90bp},clip,width=\textwidth,max height=0.72\textheight]{Marchisio2013.pdf}
\par\vspace{0.45\baselineskip}
\small 图C.3。方程式中混合解的矩示例。 (C.17) 其中 xi*= 1/2。
\end{sourceblock}
\begin{sourceblock}
\includegraphics[page=486,pagebox=mediabox,trim={58.00bp 69.07bp 58.23bp 446.66bp},clip,width=\textwidth,max height=0.72\textheight]{Marchisio2013.pdf}
\par\vspace{0.45\baselineskip}
\small 图C.4。从图 C.3 中的四个时刻找到的权重与节点。
\end{sourceblock}
在图 C.3 中的 x 中是连续的，但它们的空间梯度在 x/t = xi*−1、x = 0 和 x/t = xi* 处不连续。图 C.4 中的权重与节点也是如此。

\paragraph*{离散NDF}
对于离散初始条件，矩为

\begin{sourceblock}
\includegraphics[page=487,pagebox=mediabox,trim={131.98bp 511.28bp 67.56bp 138.52bp},clip,max width=0.94\textwidth,max height=0.38\textheight]{Marchisio2013.pdf}
\end{sourceblock}
读者可以注意到，中心区域的上限和下限对应于根据等式计算的特征。 (C.8) 从 x(0) = 0 开始。这种情况下的权重与节点具有简单的形式：

\begin{sourceblock}
\includegraphics[page=487,pagebox=mediabox,trim={101.27bp 414.86bp 67.56bp 246.70bp},clip,max width=0.94\textwidth,max height=0.38\textheight]{Marchisio2013.pdf}
\end{sourceblock}
但请注意，节点的数量在特征交叉的点处发生变化。如前所述，交叉特性导致初始 NDF 的“混合”。

\section*{C.4 QBMM的应用}\phantomsection\addcontentsline{toc}{section}{C.4 QBMM的应用}
QBMM 在方程中的应用当 u(Ψ) 取决于 Ψ 时，(C.1) 需要闭包。尽管如此，生成的矩方程（用于 QMOM 或 EQMOM）以及权重与节点的传输方程（用于 DQMOM）仍然是双曲线的。根据双曲守恒定律，矩是守恒变量（由 f 的线性运算产生），而权重与节点是原始变量。由于矩守恒对于矩反演算法的稳定性很重要，因此数值算法必须保证矩守恒。对于双曲系统，使用有限体积方法 (FVM)（或更具体地说，可实现的 FVM）最容易实现这一点。另一个重要的考虑因素是用于闭合函数的矩闭合的准确性，如下所述。

\subsection*{C.4.1 矩输运方程}\phantomsection\addcontentsline{toc}{subsection}{C.4.1 矩输运方程}
从方程找到的矩输运方程。 (C.1) 是

\begin{sourceblock}
\includegraphics[page=487,pagebox=mediabox,trim={212.40bp 112.06bp 67.56bp 562.29bp},clip,max width=0.94\textwidth,max height=0.38\textheight]{Marchisio2013.pdf}
\end{sourceblock}
其中通量函数针对上述两种情况定义为

\begin{sourceblock}
\includegraphics[page=487,pagebox=mediabox,trim={118.37bp 51.31bp 67.56bp 611.39bp},clip,max width=0.94\textwidth,max height=0.38\textheight]{Marchisio2013.pdf}
\end{sourceblock}
为了求解矩输运方程，必须为 mk+1 提供闭包。使用 QMOM，前 2N 个矩可用于找到 N 个权重 wα 和 N 个节点 xiα。然后可以使用以下命令关闭通量函数

\begin{sourceblock}
\includegraphics[page=488,pagebox=mediabox,trim={203.53bp 562.86bp 65.07bp 104.08bp},clip,max width=0.94\textwidth,max height=0.38\textheight]{Marchisio2013.pdf}
\end{sourceblock}
或者，我们可以使用 EQMOM 并传输 2N + 1 个矩，从中可以找到 N 个权重 wα、N 个节点 xiα 和 σ。在这种情况下，m2N+2 的闭包可以从重构的 NDF 中找到，如 3.3 节所述。正如双曲系统通常的情况一样，读者可以注意到方程： （C.20）以保守形式编写，可以使用 KBFVM 进行数值求解（参见附录 B 的 B.3 节）。对于离散 NDF，QMOM 闭包将是精确的，其中 N = 1（方程（C.15）中的分离解）或 N = 2（方程（C.19）中的混合解）。请注意，对于混合解决方案，自适应惠勒算法将从 N = 1 切换到 N = 2 在 x/t = | xi * / 2 | 处，这对应于交叉特性。此外，在正确制定的基于动理学的有限体积算法中，通过反转矩得到的节点将保持其原始值，即± xi *，并且权重应该非常接近统一（仅由于数值扩散而不同）。对于连续 NDF，由于关闭 Fk 时产生的错误，QMOM 和 EQMOM 闭包都不会准确。然而，通过增加 N，应该可以接近图 C.1 和 C.3 中所示的矩的演变。在所有情况下，用 KBFVM 找到的闭矩方程组的解将表现良好（例如，矩将是可实现的），并且没有理由在方程的右侧添加扩散项。 （C.20）。

\subsection*{C.4.2 权重与节点的传输方程}\phantomsection\addcontentsline{toc}{subsection}{C.4.2 权重与节点的传输方程}
在 DQMOM 中，权重与节点是通过求解从方程 1 导出的传输方程来找到的。 (C.20) 将矩定义为

\begin{sourceblock}
\includegraphics[page=488,pagebox=mediabox,trim={171.90bp 254.67bp 65.07bp 412.27bp},clip,max width=0.94\textwidth,max height=0.38\textheight]{Marchisio2013.pdf}
\end{sourceblock}
对于方程。 (C.1)，这导致

\begin{sourceblock}
\includegraphics[page=488,pagebox=mediabox,trim={205.16bp 213.13bp 65.07bp 461.22bp},clip,max width=0.94\textwidth,max height=0.38\textheight]{Marchisio2013.pdf}
\end{sourceblock}
以及

\begin{sourceblock}
\includegraphics[page=488,pagebox=mediabox,trim={206.20bp 178.54bp 65.07bp 495.81bp},clip,max width=0.94\textwidth,max height=0.38\textheight]{Marchisio2013.pdf}
\end{sourceblock}
该双曲方程组有精确解：

\begin{sourceblock}
\includegraphics[page=488,pagebox=mediabox,trim={67.52bp 124.68bp 65.07bp 535.95bp},clip,max width=0.94\textwidth,max height=0.38\textheight]{Marchisio2013.pdf}
\end{sourceblock}
α 是常数，对应于式（1）中的离散 NDF。 （C.15）和等式。 （C.19）。如前所述，方程中的解决方案。 (C.26) 等价于原始双曲方程的特征，并且对于混合解，特征是交叉的。这一观察结果意味着描述时刻所需的节点数量可以在空间和时间上不连续地变化。此外，因为等式。 (C.25) 是一阶的，精确解中只有一个积分常数可用。因此，将给定 α 的 xiα(0, x) 值设置为 x 的函数将导致在特征交叉处得到多值解。相反， xiα(0, x) 必须独立于 x，并且 α 的范围应足够大以捕获所有 x.2 的 f(0, x)。2 同时在方程的右侧添加扩散项。 (C.24) 和 (C.25) 将使 DQMOM 方程呈抛物线形（因此允许平滑特征交叉点处生成的冲击），这种做法是非物理的，并且不尊重方程 1 的双曲性质。 （C.20）。与数值求解 DQMOM 方程相关的另一个困难是很难保证矩守恒。只有一阶平流格式可以保证守恒，但数值扩散太大，无法获得准确的解。因此，作者认为，双曲矩系统最好使用方程来求解。 （C.20）与 Fk 的适当 QBMM 闭包（例如基于 EQMOM 的闭包）结合可实现的高阶方案（Vikas 等人，2011a）。

2 这正是描述双曲系统 DQMOM 的原始论文中采用的方法（Fan 等人，2004 年）。当使用这种方法时，使用具有固定节点 xi* α 的权重 wα(0, x) 来近似 f(0, x)。

\chapter{全守恒直接求积矩方法}
在本附录中，我们讨论了 DQMOM 的矩守恒问题，并提出了一种称为 DQMOM 完全保守 (DQMOM-FC) 的变体，可用于部分克服该问题。为了说明这个问题，我们考虑单变量 NDF n(t, x, xi) 的一维空间非均匀 PBE，其具有恒定速度和恒定扩散系数，如第 8.3 节所述。然而，正如下面将变得更加清楚的那样，扩展到其他更复杂的问题是很简单的。

\section*{D.1 非均匀 PBE}\phantomsection\addcontentsline{toc}{section}{D.1 非均匀 PBE}
让我们考虑简单的空间非均匀 PBE：

\begin{sourceblock}
\includegraphics[page=490,pagebox=mediabox,trim={173.28bp 270.06bp 65.07bp 404.29bp},clip,max width=0.94\textwidth,max height=0.38\textheight]{Marchisio2013.pdf}
\end{sourceblock}
其中粒子速度 u 和扩散系数 Γ 是常数（即它们不依赖于 t、x 和 xi）。这并不是一个重大限制，因为使用第 8 章中介绍的方法，本附录中介绍的材料可以轻松扩展到更一般的情况。此外，我们假设初始 NDF n(0, x, xi) 以及相空间过程（例如成核、生长、聚集和断裂时效）的源项 S(t, x, xi) 已知。当采用算子分裂时，源项通常与第 8 章中所述的空间传输项分开处理，因此可以视为空。然而，为了与我们最初的 DQMOM 推导（Marchisio \& Fox，2005）的完整性和一致性，我们将在讨论中考虑这个术语。

\section*{D.2 标准 DQMOM}\phantomsection\addcontentsline{toc}{section}{D.2 标准 DQMOM}
当应用标准 DQMOM 时，权重 w′ 的传输方程

\begin{sourceblock}
\includegraphics[page=490,pagebox=mediabox,trim={111.24bp 50.73bp 65.07bp 600.92bp},clip,max width=0.94\textwidth,max height=0.38\textheight]{Marchisio2013.pdf}
\end{sourceblock}
以及

\begin{sourceblock}
\includegraphics[page=491,pagebox=mediabox,trim={175.94bp 600.36bp 67.55bp 72.45bp},clip,max width=0.94\textwidth,max height=0.38\textheight]{Marchisio2013.pdf}
\end{sourceblock}
其中 α = 1，.。。 , N，其中使用撇号 '，其原因将在下面变得更加清楚，表示使用标准 DQMOM 获得的解决方案。还必须定义权重和加权节点的初始条件，这通常是通过计算与初始 NDF 相关的矩来完成的，mk(0, x) =

\begin{sourceblock}
\includegraphics[page=491,pagebox=mediabox,trim={276.41bp 556.06bp 210.93bp 132.13bp},clip,max width=0.94\textwidth,max height=0.38\textheight]{Marchisio2013.pdf}
\end{sourceblock}
◆kn(0, x, ◆)d◆，然后应用第 3 章中介绍的矩反演算法之一。（更多详细信息请参阅第 7.4 节。）源项 a′ α 和 b′ α 通过选择一组 2N 矩（为简单起见，这里我们使用 2N 整数矩）并求解以下线性系统来计算（Marchisio \& Fox，2005）：

\begin{sourceblock}
\includegraphics[page=491,pagebox=mediabox,trim={65.03bp 444.27bp 67.55bp 204.48bp},clip,max width=0.94\textwidth,max height=0.38\textheight]{Marchisio2013.pdf}
\end{sourceblock}
\begin{sourceblock}
\includegraphics[page=491,pagebox=mediabox,trim={169.08bp 398.91bp 67.55bp 268.02bp},clip,max width=0.94\textwidth,max height=0.38\textheight]{Marchisio2013.pdf}
\end{sourceblock}
重要的是要提醒读者，为了正确传输大于 1 的阶矩（即 k ≥ 2），该校正项是必要的。

\paragraph*{有限体积离散化}
正如第 8 章所讨论的，当采用具有均匀网格间距 Δx 的 FVM 时，定义了以 xi 为中心的单元中的以下体积平均量：

\begin{sourceblock}
\includegraphics[page=491,pagebox=mediabox,trim={179.58bp 283.77bp 67.55bp 381.91bp},clip,max width=0.94\textwidth,max height=0.38\textheight]{Marchisio2013.pdf}
\end{sourceblock}
以及

\begin{sourceblock}
\includegraphics[page=491,pagebox=mediabox,trim={181.51bp 233.40bp 67.55bp 432.27bp},clip,max width=0.94\textwidth,max height=0.38\textheight]{Marchisio2013.pdf}
\end{sourceblock}
其中 xi −Δx/2 和 xi + Δx/2 表示以 xi 为中心的单元左右界面的位置。该体积平均程序在方程中的应用。 (D.2) 和 (D.3) 得出以下表达式：

\begin{sourceblock}
\includegraphics[page=491,pagebox=mediabox,trim={198.84bp 164.86bp 67.55bp 506.24bp},clip,max width=0.94\textwidth,max height=0.38\textheight]{Marchisio2013.pdf}
\end{sourceblock}
以及

\begin{sourceblock}
\includegraphics[page=491,pagebox=mediabox,trim={198.29bp 117.51bp 67.55bp 553.59bp},clip,max width=0.94\textwidth,max height=0.38\textheight]{Marchisio2013.pdf}
\end{sourceblock}
其中方程右侧的第一项。 (D.8)和(D.9)表示由于物理空间中的平流而导致的以xi为中心的单元中的权重和加权节点的变化率，可以写为

\begin{sourceblock}
\includegraphics[page=491,pagebox=mediabox,trim={167.12bp 54.69bp 67.56bp 617.16bp},clip,max width=0.94\textwidth,max height=0.38\textheight]{Marchisio2013.pdf}
\end{sourceblock}
Here, w′

\begin{sourceblock}
\includegraphics[page=492,pagebox=mediabox,trim={67.52bp 602.98bp 65.07bp 71.77bp},clip,max width=0.94\textwidth,max height=0.38\textheight]{Marchisio2013.pdf}
\end{sourceblock}
\begin{sourceblock}
\includegraphics[page=492,pagebox=mediabox,trim={171.54bp 557.00bp 65.07bp 114.84bp},clip,max width=0.94\textwidth,max height=0.38\textheight]{Marchisio2013.pdf}
\end{sourceblock}
\begin{sourceblock}
\includegraphics[page=492,pagebox=mediabox,trim={67.52bp 521.49bp 203.13bp 150.54bp},clip,max width=0.94\textwidth,max height=0.38\textheight]{Marchisio2013.pdf}
\end{sourceblock}
k,i+1(t)\}。等式右侧的第二项。 (D.8) 和 (D.9) 是由于扩散引起的，可以使用可实现的有限体积方案以相同的方式计算：

\begin{sourceblock}
\includegraphics[page=492,pagebox=mediabox,trim={159.74bp 463.91bp 65.07bp 207.83bp},clip,max width=0.94\textwidth,max height=0.38\textheight]{Marchisio2013.pdf}
\end{sourceblock}
以及

\begin{sourceblock}
\includegraphics[page=492,pagebox=mediabox,trim={160.56bp 417.35bp 65.07bp 254.40bp},clip,max width=0.94\textwidth,max height=0.38\textheight]{Marchisio2013.pdf}
\end{sourceblock}
有关如何在标准有限体积方案中根据相邻单元的体积平均值计算重建值的更多详细信息，请参见第 8 章。 (D.8) 和 (D.9)，a′ α,i 和 b′ α,i 表示因相空间过程而导致以 xi 为中心的单元中的权重和加权节点的变化率，并由式（1）中的线性系统计算得出。 (D.4)，通过在以 xi 为中心的单元中使用体积平均源和校正项 Sk,i 和 Ck,i。这可以通过使用体积平均权重 w′ α,i 和体积平均节点 ψ′ 评估 Sk 和 Ck 来轻松完成

\begin{sourceblock}
\includegraphics[page=492,pagebox=mediabox,trim={144.45bp 312.54bp 299.27bp 371.25bp},clip,max width=0.94\textwidth,max height=0.38\textheight]{Marchisio2013.pdf}
\end{sourceblock}
α，i。时间积分通常使用显式或隐式方案来执行，这将在下面讨论。使用此过程，所得数值方案仅保留权重和加权节点及其线性组合。事实上，由于所进行的近似（例如，在评估方程（D.5）中的梯度时）以及由于离散化误差（例如数值扩散），FVM 应用于方程（D.5）。 (D.2) 和 (D.3) 导致矩守恒误差。1 显而易见，NDF 中只有两个矩会守恒，即零阶矩和一阶矩（即 m0 和 m1）。通过使用 FVM 和求积近似，这两个矩可以写为 m′

\begin{sourceblock}
\includegraphics[page=492,pagebox=mediabox,trim={67.52bp 202.82bp 66.06bp 461.27bp},clip,max width=0.94\textwidth,max height=0.38\textheight]{Marchisio2013.pdf}
\end{sourceblock}
αi。有趣的是，仅两个守恒矩的阶数（即 k = 0, 1）对应于零校正项，这可以通过检查等式轻松确认。 （D.5）。这意味着对于零阶矩和一阶矩，物理空间扩散项（实数或数值）的矩的正确传递（以及因此守恒）始终得到保证，但对于高阶矩需要一个校正项。事实上，对于所有其他涉及的量，这些量不是这些原始变量的线性组合，特别是大于 1 的阶矩，无法保证守恒，因为数值扩散引起的误差不会抵消。此外，如附录 C 中所讨论的，在存在冲击以及不连续权重与节点的情况下，不能使用 DQMOM 处理 Γ = 0（其中 Ck = 0）发现的双曲系统。

\footnotetext[1]{在 FVM 中，保守意味着对于封闭系统，所有网格单元 i 的传输项之和为零。}
\section*{D.3 DQMOM-FC}\phantomsection\addcontentsline{toc}{section}{D.3 DQMOM-FC}
对于 Γ > 0，可以通过正确处理物理空间中的平流和扩散项来开发保守方案（用于物理空间中的传输）。关于重写等式中的 PBE。 (D.1) 作为

\begin{sourceblock}
\includegraphics[page=493,pagebox=mediabox,trim={156.04bp 548.57bp 67.57bp 125.68bp},clip,max width=0.94\textwidth,max height=0.38\textheight]{Marchisio2013.pdf}
\end{sourceblock}
我们可以应用矩变换来获得

\begin{sourceblock}
\includegraphics[page=493,pagebox=mediabox,trim={115.38bp 490.94bp 67.57bp 174.63bp},clip,max width=0.94\textwidth,max height=0.38\textheight]{Marchisio2013.pdf}
\end{sourceblock}
ψkn dψ 和 Sk 是对应于相空间过程的源项的 k 阶矩。物理空间 2 中对应于平流的项是其中 mk =

\begin{sourceblock}
\includegraphics[page=493,pagebox=mediabox,trim={170.42bp 436.79bp 67.57bp 230.24bp},clip,max width=0.94\textwidth,max height=0.38\textheight]{Marchisio2013.pdf}
\end{sourceblock}
物理空间中的扩散术语是

\begin{sourceblock}
\includegraphics[page=493,pagebox=mediabox,trim={171.16bp 386.35bp 67.57bp 280.58bp},clip,max width=0.94\textwidth,max height=0.38\textheight]{Marchisio2013.pdf}
\end{sourceblock}
如果我们现在将 DQMOM 方法3应用于方程。 (D.15)，得到以下方程：

\begin{sourceblock}
\includegraphics[page=493,pagebox=mediabox,trim={199.10bp 325.80bp 67.56bp 347.01bp},clip,max width=0.94\textwidth,max height=0.38\textheight]{Marchisio2013.pdf}
\end{sourceblock}
以及

\begin{sourceblock}
\includegraphics[page=493,pagebox=mediabox,trim={200.75bp 275.11bp 67.57bp 397.69bp},clip,max width=0.94\textwidth,max height=0.38\textheight]{Marchisio2013.pdf}
\end{sourceblock}
其中双撇号 '' 表示使用 DQMOM-FC 获得的解。右侧的源项是根据以下线性系统计算得出的：

\begin{sourceblock}
\includegraphics[page=493,pagebox=mediabox,trim={156.46bp 174.33bp 67.56bp 458.40bp},clip,max width=0.94\textwidth,max height=0.38\textheight]{Marchisio2013.pdf}
\end{sourceblock}
2 可以直接证明，如果 u 不是常数而是取决于内坐标，则可以通过使用正交近似来克服由此产生的闭包问题：

\begin{sourceblock}
\includegraphics[page=493,pagebox=mediabox,trim={194.33bp 116.66bp 196.86bp 553.48bp},clip,max width=0.94\textwidth,max height=0.38\textheight]{Marchisio2013.pdf}
\end{sourceblock}
类似的近似可用于处理与尺寸相关的扩散系数 Γ(xi)，如第 8 章所述。 3 更准确地说，我们将定义替换为

\begin{sourceblock}
\includegraphics[page=493,pagebox=mediabox,trim={148.50bp 65.14bp 151.63bp 607.68bp},clip,max width=0.94\textwidth,max height=0.38\textheight]{Marchisio2013.pdf}
\end{sourceblock}
{\small 进入方程的左侧。 (D.15)，并取关于 t 的偏导数。\par}
以及

\begin{sourceblock}
\includegraphics[page=494,pagebox=mediabox,trim={165.73bp 584.24bp 65.07bp 82.70bp},clip,max width=0.94\textwidth,max height=0.38\textheight]{Marchisio2013.pdf}
\end{sourceblock}
当 k = 0 时，.。。 , 2N -1。请注意，由于没有使用空间导数，因此不需要校正项来确保当 Γ > 0 时大于 1 阶的矩在空间中守恒。这一观察结果表明，可以通过直接从方程 1 开始导出保守的有限体积方案。 （D.15）。

\paragraph*{有限体积离散化}
我们可以从方程开始推导出体积平均矩演化方程。 （D.15）：

dmk,i

\begin{sourceblock}
\includegraphics[page=494,pagebox=mediabox,trim={189.18bp 449.59bp 65.07bp 231.40bp},clip,max width=0.94\textwidth,max height=0.38\textheight]{Marchisio2013.pdf}
\end{sourceblock}
其中体积平均矩和右侧定义为

\begin{sourceblock}
\includegraphics[page=494,pagebox=mediabox,trim={181.48bp 394.27bp 65.07bp 271.42bp},clip,max width=0.94\textwidth,max height=0.38\textheight]{Marchisio2013.pdf}
\end{sourceblock}
以及

\begin{sourceblock}
\includegraphics[page=494,pagebox=mediabox,trim={179.04bp 333.82bp 65.07bp 331.86bp},clip,max width=0.94\textwidth,max height=0.38\textheight]{Marchisio2013.pdf}
\end{sourceblock}
根据定义，对于给定时间 t，每个 i 的权重和加权节点与体积平均矩 4 的关系为：

\begin{sourceblock}
\includegraphics[page=494,pagebox=mediabox,trim={184.34bp 261.11bp 65.07bp 405.82bp},clip,max width=0.94\textwidth,max height=0.38\textheight]{Marchisio2013.pdf}
\end{sourceblock}
\begin{sourceblock}
\includegraphics[page=494,pagebox=mediabox,trim={327.11bp 238.23bp 124.83bp 445.20bp},clip,max width=0.94\textwidth,max height=0.38\textheight]{Marchisio2013.pdf}
\end{sourceblock}
对于 k = 0，.。。 , 2N -1。这里，值得注意的是 w′′ α,i 不等于从 wα(t, x) 和 ςα(t, x) 中找到的体积平均值（即它们不是方程（D.6）和（D.7）中定义的变量）。再次应用 DQMOM 方法，权重和加权节点的演化方程为

\begin{sourceblock}
\includegraphics[page=494,pagebox=mediabox,trim={199.49bp 169.95bp 65.07bp 501.15bp},clip,max width=0.94\textwidth,max height=0.38\textheight]{Marchisio2013.pdf}
\end{sourceblock}
以及

\begin{sourceblock}
\includegraphics[page=494,pagebox=mediabox,trim={100.11bp 91.32bp 65.07bp 558.58bp},clip,max width=0.94\textwidth,max height=0.38\textheight]{Marchisio2013.pdf}
\end{sourceblock}
where a′′

\footnotetext[4]{这些量可以从矩 mk,i(t) 中找到，例如使用第 3 章中描述的矩反演算法。}
\begin{sourceblock}
\includegraphics[page=495,pagebox=mediabox,trim={65.03bp 599.63bp 67.56bp 73.36bp},clip,max width=0.94\textwidth,max height=0.38\textheight]{Marchisio2013.pdf}
\end{sourceblock}
\begin{sourceblock}
\includegraphics[page=495,pagebox=mediabox,trim={168.06bp 557.89bp 67.57bp 116.37bp},clip,max width=0.94\textwidth,max height=0.38\textheight]{Marchisio2013.pdf}
\end{sourceblock}
其中 mk,i±1/2 是从体积平均矩 \{mk,i−1(t), mk,i(t), mk,i+1(t)\} 重建的细胞界面处的矩。我们建议对这种重构的可实现性的众多含义感兴趣的读者参阅第 8 章。类似地，扩散项 ad α,i 和 bd α,i 是用相同的方式从方程 1 中计算出来的。 （D.21）通过使用可实现的有限体积方案来表示 Dk,i：

\begin{sourceblock}
\includegraphics[page=495,pagebox=mediabox,trim={153.94bp 457.71bp 67.56bp 216.45bp},clip,max width=0.94\textwidth,max height=0.38\textheight]{Marchisio2013.pdf}
\end{sourceblock}
其中，与第 8 章的符号一致，mk,i+1、mk,i 和 mk,i−1 分别是单元 i + 1、i 和 i −1 中的体积平均矩。使用 DQMOM-FC，固定 t 时会自动遵守矩守恒，因为用于表示空间传输的有限体积方案是保守的。5 然而，方程 1 的数值解的整体守恒性质。 (D.27) 和 (D.28) 将取决于时间积分方案，如下所述。

\section*{D.4 时间积分}\phantomsection\addcontentsline{toc}{section}{D.4 时间积分}
现在让我们详细讨论时间积分过程中的矩守恒问题。与第 8 章一致，相空间过程的源项被设置为零，因此本讨论中仅考虑物理空间中的传输项。当方程。 (D.23) 使用显式欧拉方案进行积分，以 xi 为中心的单元在 (n + 1)Δt 时刻的 k 阶体积平均矩直接由 k 阶体积平均矩在 n Δt 时刻计算得出，公式如下：

\begin{sourceblock}
\includegraphics[page=495,pagebox=mediabox,trim={184.77bp 240.99bp 67.56bp 435.08bp},clip,max width=0.94\textwidth,max height=0.38\textheight]{Marchisio2013.pdf}
\end{sourceblock}
如果使用一阶迎风空间差分格式，则单元界面处的重构矩 mk,i±1/2 等于迎风值，并且在 u 为正的情况下，得到以下方程：

\begin{sourceblock}
\includegraphics[page=495,pagebox=mediabox,trim={84.41bp 165.03bp 67.56bp 501.92bp},clip,max width=0.94\textwidth,max height=0.38\textheight]{Marchisio2013.pdf}
\end{sourceblock}
这种方法（即 FVM 中的 QMOM）始终是保守的，这意味着以下关系始终成立（在相空间过程的源项为空的情况下）：

\begin{sourceblock}
\includegraphics[page=495,pagebox=mediabox,trim={171.93bp 108.75bp 157.96bp 574.85bp},clip,max width=0.94\textwidth,max height=0.38\textheight]{Marchisio2013.pdf}
\end{sourceblock}
{\small 其中两个最终等式是系统关闭后得出的。\par}
\footnotetext[5]{换句话说，方程 1 的所有网格单元 i 的总和。由于封闭系统的平流（方程（D.29））和扩散（方程（D.30）），（D.23）在右侧产生零贡献。}
对于标准 DQMOM，当采用时间上的显式欧拉格式和空间上的一阶迎风差分格式时，以 xi 为中心的单元在时间 (n + 1)Δt 处的体积平均权重为

\begin{sourceblock}
\includegraphics[page=496,pagebox=mediabox,trim={87.94bp 569.63bp 65.07bp 97.30bp},clip,max width=0.94\textwidth,max height=0.38\textheight]{Marchisio2013.pdf}
\end{sourceblock}
其中 α = 1，.。。，N，而体积平均加权节点是

\begin{sourceblock}
\includegraphics[page=496,pagebox=mediabox,trim={91.67bp 516.49bp 65.07bp 150.45bp},clip,max width=0.94\textwidth,max height=0.38\textheight]{Marchisio2013.pdf}
\end{sourceblock}
其中 α = 1，.。。，N。如果在n Δt时刻，权重和加权节点是准确的，即关系

\begin{sourceblock}
\includegraphics[page=496,pagebox=mediabox,trim={198.72bp 455.51bp 65.07bp 211.43bp},clip,max width=0.94\textwidth,max height=0.38\textheight]{Marchisio2013.pdf}
\end{sourceblock}
对于 k = 0 有效，.。。 , 2N −1, 则直接求解式(n + 1)Δt计算出时刻(n + 1)Δt的矩。 (D.32)（即使用 QMOM）可能与根据等式 1 确定的权重和加权节点计算出的矩不同。 （D.33）和（D.34）（即使用 DQMOM）。事实上，很容易证明根据方程计算出的矩： (D.32)，mn+1

\begin{sourceblock}
\includegraphics[page=496,pagebox=mediabox,trim={250.69bp 389.46bp 185.04bp 293.92bp},clip,max width=0.94\textwidth,max height=0.38\textheight]{Marchisio2013.pdf}
\end{sourceblock}
k,i 以及根据 w′n+1 计算的值

\begin{sourceblock}
\includegraphics[page=496,pagebox=mediabox,trim={187.86bp 349.86bp 65.07bp 317.08bp},clip,max width=0.94\textwidth,max height=0.38\textheight]{Marchisio2013.pdf}
\end{sourceblock}
仅当 k = 0 和 k = 1 时相同，而当 k ≥2 时，两个解将导致不同的 mn+1 k,i 和 m′n+1 k,i 值。此外，在 FVM 中使用 QMOM 传输矩时，时间积分始终是保守的。相反，在 FVM 中使用标准 DQMOM 传输权重和加权节点时，时间积分仅在 k = 0 和 k = 1 时刻是保守的，或者换句话说，

\begin{sourceblock}
\includegraphics[page=496,pagebox=mediabox,trim={222.43bp 258.93bp 203.49bp 425.08bp},clip,max width=0.94\textwidth,max height=0.38\textheight]{Marchisio2013.pdf}
\end{sourceblock}
仅适用于 k = 0 和 k = 1。有趣的是，在使用标准 DQMOM 与 QMOM 计算矩时，估计在时间 (n + 1)Δt 处发生的守恒误差，并与 QMOM 进行比较，从时间 n Δt 开始，具有相同的矩集。例如，在无扩散的情况下（即 Γ = 0），对于 k = 2，时间步 n 处的误差由下式给出

\begin{sourceblock}
\includegraphics[page=496,pagebox=mediabox,trim={67.52bp 133.41bp 65.07bp 507.72bp},clip,max width=0.94\textwidth,max height=0.38\textheight]{Marchisio2013.pdf}
\end{sourceblock}
α,i−1。方程 (D.37) 清楚地表明，对于空间非均匀系统，仅当节点空间均匀时，QMOM 和标准 DQMOM 预测才相同。如果节点不均匀（或者更糟糕的是，如果节点由于例如冲击而不连续），则误差将会很大。请注意，无量纲量 u Δt/Δx 是 CFL 数，在模拟过程中通常保持恒定。因此，减少时间步长不会降低模拟结束时的总误差，因为完成模拟需要更多的时间步长，而总误差是总和 -
{\small n Errn.\par}
当 DQMOM-FC 与显式欧拉时间方案一起使用时，方程

\begin{sourceblock}
\includegraphics[page=497,pagebox=mediabox,trim={200.77bp 604.90bp 67.56bp 78.85bp},clip,max width=0.94\textwidth,max height=0.38\textheight]{Marchisio2013.pdf}
\end{sourceblock}
其中 α = 1，.。。，N 为权重，

\begin{sourceblock}
\includegraphics[page=497,pagebox=mediabox,trim={202.42bp 565.84bp 67.56bp 117.91bp},clip,max width=0.94\textwidth,max height=0.38\textheight]{Marchisio2013.pdf}
\end{sourceblock}
其中 α = 1，.。。，N为加权节点，用于时间积分。对于 k = 0，.。。 , 2N −1，项 An α,i 和 Bn α,i 与方程 (2) 中的矩源项有关。 (D.31) 由线性系统

\begin{sourceblock}
\includegraphics[page=497,pagebox=mediabox,trim={65.02bp 465.46bp 67.56bp 181.35bp},clip,max width=0.94\textwidth,max height=0.38\textheight]{Marchisio2013.pdf}
\end{sourceblock}
\begin{sourceblock}
\includegraphics[page=497,pagebox=mediabox,trim={98.28bp 417.19bp 67.56bp 249.75bp},clip,max width=0.94\textwidth,max height=0.38\textheight]{Marchisio2013.pdf}
\end{sourceblock}
如果在时间 n Δt 时，权重和加权节点恰好代表矩集，

\begin{sourceblock}
\includegraphics[page=497,pagebox=mediabox,trim={138.12bp 360.85bp 67.56bp 306.09bp},clip,max width=0.94\textwidth,max height=0.38\textheight]{Marchisio2013.pdf}
\end{sourceblock}
即 QMOM 的矩设置与 DQMOM-FC 的矩设置相同，很明显，方程（D.31）与式（D.31）一致。 (D.41) 仅当 k = 0 时，

\begin{sourceblock}
\includegraphics[page=497,pagebox=mediabox,trim={157.14bp 292.56bp 67.56bp 374.38bp},clip,max width=0.94\textwidth,max height=0.38\textheight]{Marchisio2013.pdf}
\end{sourceblock}
and k = 1,

\begin{sourceblock}
\includegraphics[page=497,pagebox=mediabox,trim={146.45bp 244.56bp 67.56bp 422.37bp},clip,max width=0.94\textwidth,max height=0.38\textheight]{Marchisio2013.pdf}
\end{sourceblock}
这意味着用方程计算矩的演化。 （D.31）（即使用 QMOM）或通过权重和加权节点与方程。 (D.38) 和 (D.39) 仅当 k = 0 和 k = 1 时才会产生相同的矩。事实上，经过一些繁琐的代数计算并通过 Δt 的幂展开，可以证明

\begin{sourceblock}
\includegraphics[page=497,pagebox=mediabox,trim={86.65bp 155.39bp 67.56bp 507.15bp},clip,max width=0.94\textwidth,max height=0.38\textheight]{Marchisio2013.pdf}
\end{sourceblock}
因此，当权重与节点随时间平滑变化时，使用 DQMOM-FC 的显式欧拉方案对于大于 1 阶的矩具有 Δt2 阶的守恒误差。请注意，误差项的系数具有与等式 1 中的校正项类似的形式。 （D.5）。使用等式中 An α,i 和 Bn α,i 的定义。 （D.38）和（D.39），方程中的系数。 (D.45) 可以重写为

\begin{sourceblock}
\includegraphics[page=497,pagebox=mediabox,trim={78.85bp 50.16bp 67.56bp 611.64bp},clip,max width=0.94\textwidth,max height=0.38\textheight]{Marchisio2013.pdf}
\end{sourceblock}
该形式清楚地表明，DQMOM-FC 中由于时间积分而产生的矩守恒误差与节点的空间梯度不成正比（如在 DQMOM 中），而是与 (xi′′n+1

\begin{sourceblock}
\includegraphics[page=498,pagebox=mediabox,trim={196.55bp 600.43bp 252.58bp 83.72bp},clip,max width=0.94\textwidth,max height=0.38\textheight]{Marchisio2013.pdf}
\end{sourceblock}
α,i)2，因此只有当节点随时间平滑变化时，它才会与 Δt2 成正比。例如，双曲系统（节点不连续）和具有小 Γ 的系统（节点的空间梯度非常大）可能由于平流而具有较大的矩守恒误差。另一个有问题的系统是包含与时间相关的空区的系统：即 wn α,i = 0 但 wn+1 α,i > 0。因此，对于所有此类系统，最好直接使用方程 1 中给出的矩输运方程。 （D.15）。

\chapter*{参考文献}\addcontentsline{toc}{chapter}{参考文献}
\markboth{参考文献}{参考文献}
\clearpage
\begin{sourceblock}
\includegraphics[page=499,pagebox=mediabox,trim={58.00bp 50.16bp 58.23bp 185.00bp},clip,width=\textwidth,max height=0.82\textheight]{Marchisio2013.pdf}
\end{sourceblock}
\clearpage
\begin{sourceblock}
\includegraphics[page=500,pagebox=mediabox,trim={58.00bp 50.16bp 58.23bp 55.00bp},clip,width=\textwidth,max height=0.82\textheight]{Marchisio2013.pdf}
\end{sourceblock}
\clearpage
\begin{sourceblock}
\includegraphics[page=501,pagebox=mediabox,trim={58.00bp 50.16bp 58.23bp 55.00bp},clip,width=\textwidth,max height=0.82\textheight]{Marchisio2013.pdf}
\end{sourceblock}
\clearpage
\begin{sourceblock}
\includegraphics[page=502,pagebox=mediabox,trim={58.00bp 50.16bp 58.23bp 55.00bp},clip,width=\textwidth,max height=0.82\textheight]{Marchisio2013.pdf}
\end{sourceblock}
\clearpage
\begin{sourceblock}
\includegraphics[page=503,pagebox=mediabox,trim={58.00bp 50.16bp 58.23bp 55.00bp},clip,width=\textwidth,max height=0.82\textheight]{Marchisio2013.pdf}
\end{sourceblock}
\clearpage
\begin{sourceblock}
\includegraphics[page=504,pagebox=mediabox,trim={58.00bp 50.16bp 58.23bp 55.00bp},clip,width=\textwidth,max height=0.82\textheight]{Marchisio2013.pdf}
\end{sourceblock}
\clearpage
\begin{sourceblock}
\includegraphics[page=505,pagebox=mediabox,trim={58.00bp 50.16bp 58.23bp 55.00bp},clip,width=\textwidth,max height=0.82\textheight]{Marchisio2013.pdf}
\end{sourceblock}
\clearpage
\begin{sourceblock}
\includegraphics[page=506,pagebox=mediabox,trim={58.00bp 50.16bp 58.23bp 55.00bp},clip,width=\textwidth,max height=0.82\textheight]{Marchisio2013.pdf}
\end{sourceblock}
\clearpage
\begin{sourceblock}
\includegraphics[page=507,pagebox=mediabox,trim={58.00bp 50.16bp 58.23bp 55.00bp},clip,width=\textwidth,max height=0.82\textheight]{Marchisio2013.pdf}
\end{sourceblock}
\clearpage
\begin{sourceblock}
\includegraphics[page=508,pagebox=mediabox,trim={58.00bp 50.16bp 58.23bp 55.00bp},clip,width=\textwidth,max height=0.82\textheight]{Marchisio2013.pdf}
\end{sourceblock}
\clearpage
\begin{sourceblock}
\includegraphics[page=509,pagebox=mediabox,trim={58.00bp 50.16bp 58.23bp 55.00bp},clip,width=\textwidth,max height=0.82\textheight]{Marchisio2013.pdf}
\end{sourceblock}
\clearpage
\begin{sourceblock}
\includegraphics[page=510,pagebox=mediabox,trim={58.00bp 50.16bp 58.23bp 55.00bp},clip,width=\textwidth,max height=0.82\textheight]{Marchisio2013.pdf}
\end{sourceblock}
\clearpage
\begin{sourceblock}
\includegraphics[page=511,pagebox=mediabox,trim={58.00bp 50.16bp 58.23bp 55.00bp},clip,width=\textwidth,max height=0.82\textheight]{Marchisio2013.pdf}
\end{sourceblock}
\clearpage
\begin{sourceblock}
\includegraphics[page=512,pagebox=mediabox,trim={58.00bp 50.16bp 58.23bp 55.00bp},clip,width=\textwidth,max height=0.82\textheight]{Marchisio2013.pdf}
\end{sourceblock}
\clearpage
\begin{sourceblock}
\includegraphics[page=513,pagebox=mediabox,trim={58.00bp 50.16bp 58.23bp 55.00bp},clip,width=\textwidth,max height=0.82\textheight]{Marchisio2013.pdf}
\end{sourceblock}
\clearpage
\begin{sourceblock}
\includegraphics[page=514,pagebox=mediabox,trim={58.00bp 50.16bp 58.23bp 55.00bp},clip,width=\textwidth,max height=0.82\textheight]{Marchisio2013.pdf}
\end{sourceblock}
\clearpage
\begin{sourceblock}
\includegraphics[page=515,pagebox=mediabox,trim={58.00bp 50.16bp 58.23bp 55.00bp},clip,width=\textwidth,max height=0.82\textheight]{Marchisio2013.pdf}
\end{sourceblock}
\clearpage
\begin{sourceblock}
\includegraphics[page=516,pagebox=mediabox,trim={58.00bp 50.16bp 58.23bp 55.00bp},clip,width=\textwidth,max height=0.82\textheight]{Marchisio2013.pdf}
\end{sourceblock}
\clearpage
\begin{sourceblock}
\includegraphics[page=517,pagebox=mediabox,trim={58.00bp 50.16bp 58.23bp 55.00bp},clip,width=\textwidth,max height=0.82\textheight]{Marchisio2013.pdf}
\end{sourceblock}
\clearpage
\begin{sourceblock}
\includegraphics[page=518,pagebox=mediabox,trim={58.00bp 50.16bp 58.23bp 55.00bp},clip,width=\textwidth,max height=0.82\textheight]{Marchisio2013.pdf}
\end{sourceblock}
\clearpage
\begin{sourceblock}
\includegraphics[page=519,pagebox=mediabox,trim={58.00bp 50.16bp 58.23bp 55.00bp},clip,width=\textwidth,max height=0.82\textheight]{Marchisio2013.pdf}
\end{sourceblock}
\clearpage
\begin{sourceblock}
\includegraphics[page=520,pagebox=mediabox,trim={58.00bp 50.16bp 58.23bp 55.00bp},clip,width=\textwidth,max height=0.82\textheight]{Marchisio2013.pdf}
\end{sourceblock}
\clearpage
\begin{sourceblock}
\includegraphics[page=521,pagebox=mediabox,trim={58.00bp 50.16bp 58.23bp 55.00bp},clip,width=\textwidth,max height=0.82\textheight]{Marchisio2013.pdf}
\end{sourceblock}
\clearpage
\begin{sourceblock}
\includegraphics[page=522,pagebox=mediabox,trim={58.00bp 50.16bp 58.23bp 55.00bp},clip,width=\textwidth,max height=0.82\textheight]{Marchisio2013.pdf}
\end{sourceblock}
\clearpage
\begin{sourceblock}
\includegraphics[page=523,pagebox=mediabox,trim={58.00bp 50.16bp 58.23bp 55.00bp},clip,width=\textwidth,max height=0.82\textheight]{Marchisio2013.pdf}
\end{sourceblock}
\clearpage
\begin{sourceblock}
\includegraphics[page=524,pagebox=mediabox,trim={58.00bp 50.16bp 58.23bp 55.00bp},clip,width=\textwidth,max height=0.82\textheight]{Marchisio2013.pdf}
\end{sourceblock}
\clearpage
\begin{sourceblock}
\includegraphics[page=525,pagebox=mediabox,trim={58.00bp 50.16bp 58.23bp 55.00bp},clip,width=\textwidth,max height=0.82\textheight]{Marchisio2013.pdf}
\end{sourceblock}
\clearpage
\begin{sourceblock}
\includegraphics[page=526,pagebox=mediabox,trim={58.00bp 50.16bp 58.23bp 55.00bp},clip,width=\textwidth,max height=0.82\textheight]{Marchisio2013.pdf}
\end{sourceblock}
\clearpage
\begin{sourceblock}
\includegraphics[page=527,pagebox=mediabox,trim={58.00bp 50.16bp 58.23bp 55.00bp},clip,width=\textwidth,max height=0.82\textheight]{Marchisio2013.pdf}
\end{sourceblock}

\end{document}
